% généré par scripts/03_catalogue.py — ne pas modifier

\annee{340}

\begin{notice}{1}{340}{ICUR n.s. I, 1419 · EDCS-39000407 · EDB31183}
\label{not:1}
\nf{Datation}{340 — directe : nom consulaire lu}
\nf{Genre}{Épitaphe ; ensembles A, B, C ; 92 lettres transmises ; 36 \% restituées.}
\nf{Références}{ICUR I 1419 = ILCV 4647.}
\nf{Provenance}{Rattachement selon l’EDB : \emph{Incertae originis} (le volume I des ICVR réunit les inscriptions dont la provenance est perdue). Conservation selon l’EDB : \emph{Musei Vaticani: Lapidario Cristiano}.}
\nf{Support}{Plaque de marbre. 35 × 65 cm (h. × l.). Inscription gravée. Reproductions signalées par l’EDB : dessin, ICVR I (de Rossi) (droits réservés ; non reproduites). Fiche EDB : Filippo Piazzolla, première version 2015.}
\ntext{\lnum{1}[- - - coniu]gi dulcissim(a)e et fideli\\
\lnum{2}[quae vixit] annis XXVIII m(ensibus) II qu(a)e\\
\lnum{3}[deposita est - - - N]onas Aug(ustas) A⸢c⸣ind⸢y⸣no\\
\lnum{4}[et Procu]lo conss(ulibus)}
\nf{Témoins}{Export EDCS de 2020 : \emph{[3 coniu]gi dulcissim(a)e et fideli / [quae vixit] annis XXVIII m(ensibus) II qu(a)e / [deposita est 3 N]onas Aug(ustas) Aindno / [et Procu]lo conss(ulibus)}. EDB (EDB31183) : texte concordant aux conventions de notation près (indice de similarité 0,925).}
\nf{Apparat}{⸢c⸣ : Q gravé ; ⸢y⸣ : I gravé}
\nf{Traduction}{[…] à son épouse très douce et fidèle, qui vécut 28 ans 2 mois, qui fut déposée […] (jour des) nones d’août, sous le consulat d’Acindynus et de Proculus.}
\nf{Langue}{L’éditeur signale 2 graphies corrigées : Q gravé pour c ; I gravé pour y (graphie de y).}
\nf{Chronologie}{Noms consulaires : Acindynus (lu) et Proculus (restitué). Dates du jour : déposition : … Nonas Augustas (partiellement restitué). Âge déclaré : 28 a. 2 m. (partiellement restitué).}
\nf{Bibliographie}{ICUR I 1419 ; ILCV 4647. Bibliographie complémentaire selon l’EDB : ILCV, 4647 ; Ferrua 1979b, 36, n. 1419 ; Lega 2000, 105.}
\nf{Statut}{Lecture secondaire (EDCS) concordante avec l’EDB : lecture stabilisée entre deux bases, non collationnée sur l’ICVR ni sur le monument. L’EDCS signale une photographie (non consultée).}
\end{notice}

\begin{notice}{2}{340}{ICUR n.s. III, 8417 · EDCS-35802005 · EDB25692}
\label{not:2}
\nf{Datation}{340 — directe : lecture partielle d’un nom consulaire}
\nf{Genre}{Épitaphe ; ensembles A, B, C ; 46 lettres transmises ; 41 \% restituées.}
\nf{Références}{ICUR III 8417.}
\nf{Provenance}{Rattachement selon l’EDB : \emph{Incertum locum in coem. Domitillae}.}
\nf{Support}{Plaque de marbre. Inscription gravée. Reproductions signalées par l’EDB : photographie PCAS ×2 (droits réservés ; non reproduites). Fiche EDB : Carlo Carletti / Antonio Enrico Felle, première version 2003.}
\ntext{\lnum{1}q]ui vixit [- - -]\\
\lnum{2}[- - -] pater fec[it - - -]\\
\lnum{3}[Acin]dino et Procu[lo conss(ulibus)}
\nf{Témoins}{EDB (EDB25692) : texte concordant aux conventions de notation près (indice de similarité 0,981).}
\nf{Traduction}{[…] qui vécut […] ; […] son père a fait (ce monument) […] ; sous le consulat d’Acindynus et de Proculus.}
\nf{Langue}{1 abréviation développée par l’éditeur (dont conss(ulibus)).}
\nf{Chronologie}{Noms consulaires : Acindynus (lu en partie) et Proculus (lu en partie).}
\nf{Bibliographie}{ICUR III 8417.}
\nf{Statut}{Lecture secondaire (EDCS) concordante avec l’EDB : lecture stabilisée entre deux bases, non collationnée sur l’ICVR ni sur le monument. L’EDCS signale une photographie (non consultée).}
\end{notice}

\begin{notice}{3}{340}{ICUR n.s. IV, 11089 · EDCS-37200371 · EDB39548}
\label{not:3}
\nf{Datation}{340 — directe : lecture partielle d’un nom consulaire}
\nf{Genre}{Épitaphe ; ensembles A, B, C ; 58 lettres transmises ; 64 \% restituées.}
\nf{Références}{ICUR IV 11089.}
\nf{Provenance}{Rattachement selon l’EDB : \emph{Coem. Callisti area subdialis et incerti loci}. Conservation selon l’EDB : \emph{Coem. Callisti area subdialis et incerti loci}.}
\nf{Support}{Plaque de marbre. 18 × 22 × 4,5 cm (h. × l. × ép.). Inscription gravée. Fiche EDB : Marida Pierno, première version 2015.}
\ntext{\lnum{1}ben]e meren[ti in pace - - -]\\
\lnum{2}[- - - perii]t pri(die) Idu[s - - -]\\
\lnum{3}[- - - Acin]dino e[t Proculo conss(ulibus)}
\nf{Témoins}{EDB (EDB39548 ; écarts à examiner) : \emph{------ / [ben]emeren[ti in pace] / [peri]t pri(die) idu[s ---] / [---]dino e[t ---] / ------}.}
\nf{Traduction}{[…] bien méritant, en paix […] ; […] mourut la veille des ides […] ; sous le consulat d’Acindynus et de Proculus.}
\nf{Langue}{2 abréviations développées par l’éditeur (dont pri(die), conss(ulibus)).}
\nf{Chronologie}{Noms consulaires : Acindynus (lu en partie) et Proculus (restitué).}
\nf{Commentaire}{La formule in pace est attestée seulement par restitution.}
\nf{Bibliographie}{ICUR IV 11089.}
\nf{Statut}{Lecture secondaire (EDCS) ; écarts avec l’EDB à arbitrer sur l’ICVR. Plus de la moitié des lettres sont restituées.}
\end{notice}

\annee{341}

\begin{notice}{4}{341}{ICUR n.s. I, 1420 · EDCS-34501143 · EDB31184}
\label{not:4}
\nf{Datation}{341 — directe : nom consulaire lu}
\nf{Genre}{Épitaphe ; ensembles A, B, C ; 189 lettres transmises.}
\nf{Références}{ICUR I 1420 = ILCV 2816 (em) = ICVaticano p 299 = ICaRoma 89 = AE 1997, 166.}
\nf{Provenance}{Rattachement selon l’EDB : \emph{Incertae originis} (le volume I des ICVR réunit les inscriptions dont la provenance est perdue). Conservation selon l’EDB : \emph{Musei Vaticani: Lapidario Cristiano}.}
\nf{Support}{Plaque de marbre. 49 × 186 cm (h. × l.). Inscription gravée. Signes décrits par l’EDB : chrisme. Reproductions signalées par l’EDB : dessin, ICVR I (de Rossi) et reproduction, Perret 1851 (droits réservés ; non reproduites). Fiche EDB : Filippo Piazzolla, première version 2015.}
\ntext{\lnum{1}Depossio Iuniani pri(die) Idus Apriles Marcellino et Probino conss(ulibus)\\
\lnum{2}qui ⸢v⸣ixit annis XL in pace dec⸢e⸣ssit et amator paupe\{ro\}rum ⸢v⸣ixit\\
\lnum{3}cum ⸢v⸣(i)rg(i)nia annis XV bene merenti ⸢v⸣irginia sua Bictora\\
\lnum{4}bene merenti fecit amatrix paupe\{ro\}rum et operaria}
\nf{Témoins}{Export EDCS de 2020 : \emph{Depossio Iuniani pri(die) Idus Apriles Marcellino et Probino conss(ulibus) / qui ixit annis XL in pace decssit et amator paupe\{ro\}rum ixit / cum (i)rg(i)nia annis XV bene merenti irginia sua Bictora / bene merenti fecit amatrix paupe\{ro\}rum et operaria}. EDB (EDB31184) : même texte ; les écarts relevés dans la liste de résultats tiennent à la notation (relecture de la fiche détaillée). Relecture de la fiche détaillée de l’EDB : Même texte. L'EDB garde la graphie de la pierre (bixit, birginia), l'EDCS la corrige (<v=B>). L'EDB décrit un chrisme final. Supports bibliographiques : Di Stefano Manzella 1997 (avec photo), Carletti 1986a n. 89.}
\nf{Apparat}{⸢v⸣ : B gravé ; ⸢e⸣ : I gravé ; ⸢v⸣ : B gravé ; ⸢v⸣ : B gravé ; ⸢v⸣ : B gravé ; \{ro\} : lettres superflues ; \{ro\} : lettres superflues}
\nf{Traduction}{Déposition de Iunianus, la veille des ides d’avril, sous le consulat de Marcellinus et de Probinus ; il vécut 40 ans ; il est mort en paix ; ami des pauvres, il vécut avec sa vierge (= épouse épousée vierge) 15 ans ; à lui, bien méritant, sa vierge Victora a fait (ce monument), bien méritant ; amie des pauvres et travailleuse.}
\nf{Langue}{L’éditeur signale 5 graphies corrigées : B gravé pour v (bêtacisme) ×4 ; I gravé pour e (confusion e/i). Des lettres gravées sont jugées superflues par l’éditeur (entre accolades). 4 abréviations développées par l’éditeur (dont pri(die), conss(ulibus), (i)rg, (i)nia).}
\nf{Chronologie}{Noms consulaires : Marcellinus (lu) et Probinus (lu). Dates du jour : déposition : pridie Idus Apriles = 12 avril 341 (lu). Âge déclaré : 40 a. (lu).}
\nf{Commentaire}{La formule in pace est conservée. Commentaire développé : dossier p.~\pageref{dos:4}.}
\nf{Bibliographie}{ICUR I 1420 ; ILCV 2816 (em) ; ICVaticano p 299 ; ICaRoma 89 ; AE 1997, 166. Bibliographie complémentaire selon l’EDB : Di Stefano Manzella 1997, 299-200, 3.7.9, con foto (C. Lega) ; Lega 2000, 92 ; Carletti 1986a, 101-102, n. 89 ; ILCV, 2816.}
\nf{Statut}{Lecture secondaire (EDCS) concordante avec l’EDB : lecture stabilisée entre deux bases, non collationnée sur l’ICVR ni sur le monument. L’EDCS signale une photographie (non consultée).}
\end{notice}

\annee{342}

\begin{notice}{5}{342}{ICUR n.s. I, 1251 · EDCS-39101961 · EDB30459}
\label{not:5}
\nf{Datation}{342 — directe : lecture partielle d’un nom consulaire}
\nf{Genre}{Épitaphe ; ensembles A, B, C ; 98 lettres transmises ; 55 \% restituées.}
\nf{Références}{ICUR I 1251 = Louvre 966.}
\nf{Provenance}{Rattachement selon l’EDB : \emph{Incertae originis} (le volume I des ICVR réunit les inscriptions dont la provenance est perdue). Conservation selon l’EDB : \emph{Musei Capitolini}.}
\nf{Support}{Plaque de marbre. 26 × 30 cm (h. × l.). Inscription gravée. Fiche EDB : Filippo Piazzolla, première version 2014.}
\ntext{\lnum{1}vixit ann(os) - - -]XVIII m(enses) [- - -]\\
\lnum{2}[- - -] Non(as) Sept[embr(es) - - -]\\
\lnum{3}[dd(ominis) nn(ostris) Co]nstanti[o III]\\
\lnum{4}[et Const]ante II Augg(ustis) [cons(ulibus)]\\
\lnum{5}[The]odotus [posuit}
\nf{Témoins}{EDB (EDB30459 ; écarts à examiner) : \emph{------ / [--- vixit annis] XVIII m[ensibus ---] / [--- dep(osit---) ---] ṇon(as) sept[embres ---] / [dd(ominis) nn(ostris) Co]nstanti[o III] / [et Const]ante II Augg̣[ustibus cons(ulibu)s] / [--- The]odotus [--- posuit] / ------ [texte tronqué dans la liste de résultats]}.}
\nf{Traduction}{[…] vécut […]18 ans, […] mois ; […] nones de septembre […] ; sous le consulat de nos seigneurs Constance (Auguste) pour la 3e fois et Constant (Auguste) pour la 2e fois ; Theodotus a posé (ce monument).}
\nf{Langue}{8 abréviations développées par l’éditeur (dont ann(os), m(enses), Non(as), Septembr(es), dd(ominis), nn(ostris)).}
\nf{Chronologie}{Noms consulaires : Constantius (lu en partie) et Constans (lu en partie). Dates du jour : événement indéterminé : … Nonas Septembres (partiellement restitué).}
\nf{Commentaire}{Vocabulaire religieux : anthroponyme (niveau 6) — theodotus (anthroponyme).}
\nf{Bibliographie}{ICUR I 1251 ; Louvre 966. Bibliographie complémentaire selon l’EDB : Molisani 1973, 110 ; Ferrua 1979b, 30, n. 1251.}
\nf{Statut}{Lecture secondaire (EDCS) ; écarts avec l’EDB à arbitrer sur l’ICVR. Plus de la moitié des lettres sont restituées. L’EDCS signale une photographie (non consultée).}
\end{notice}

\begin{notice}{6}{342}{ICUR n.s. IV, 9556 · EDCS-36101085 · EDB20872}
\label{not:6}
\nf{Datation}{342 — directe : lecture partielle d’un nom consulaire}
\nf{Genre}{Épitaphe ; ensembles A, B, C ; 152 lettres transmises ; 66 \% restituées.}
\nf{Références}{ICUR IV 9556 = ILCV +2846.}
\nf{Provenance}{Rattachement selon l’EDB : \emph{Coem. Callisti pars inferior}.}
\nf{Support}{Plaque de marbre. Inscription gravée. Reproductions signalées par l’EDB : photographie PCAS (droits réservés ; non reproduites). Fiche EDB : Antonio Enrico Felle, première version 2003.}
\ntext{\lnum{1}[Aur(elius)] Paulinus Aur(eliae) Ta[eo]dore [filiae dul]\\
\lnum{2}[ciss]im(a)e fecit qu(a)e pr[aec]essit in [pace dd(ominis) nn(ostris)]\\
\lnum{3}[Con]statio Aug(usto) III et [Constante Aug(usto) II conss(ulibus)\\
\lnum{4}[- - - D]ec(embres) et ⸢v⸣ixit an(nos) LI[- - -]}
\nf{Témoins}{Export EDCS de 2020 : \emph{[Aur(elius)] Paulinus Aur(eliae) Ta[eo]dore [filiae dul]/[ciss]im(a)e fecit qu(a)e pr[aec]essit in [pace dd(ominis) nn(ostris)] / [Con]statio Aug(usto) III et [Constante Aug(usto) II conss(ulibus) / [3 D]ec(embres) et ixit an(nos) LI[3]}. EDB (EDB20872) : texte concordant aux conventions de notation près (indice de similarité 0,98).}
\nf{Apparat}{⸢v⸣ : B gravé}
\nf{Traduction}{Aurelius Paulinus à Aurelia Theodora, sa fille très douce, a fait (ce monument) ; elle (l’)a précédé dans la paix, sous le consulat de nos seigneurs Constance Auguste pour la 3e fois et Constant Auguste pour la 2e fois, […] décembre ; elle vécut 51 (?) ans […].}
\nf{Langue}{L’éditeur signale une graphie corrigée : B gravé pour v (bêtacisme).}
\nf{Chronologie}{Noms consulaires : Constantius (lu en partie) et Constans (restitué).}
\nf{Commentaire}{La formule in pace est attestée seulement par restitution.}
\nf{Bibliographie}{ICUR IV 9556 ; ILCV +2846.}
\nf{Statut}{Lecture secondaire (EDCS) concordante avec l’EDB : lecture stabilisée entre deux bases, non collationnée sur l’ICVR ni sur le monument. Plus de la moitié des lettres sont restituées. L’EDCS signale une photographie (non consultée).}
\end{notice}

\annee{343}

\begin{notice}{7}{343}{ICUR n.s. I, 264 · EDCS-38600270 · EDB27782}
\label{not:7}
\nf{Datation}{343 — directe : nom consulaire lu}
\nf{Genre}{Épitaphe ; ensembles A, B, C ; 86 lettres transmises ; 12 \% restituées.}
\nf{Références}{ICUR I 264 = ILCV 4393.}
\nf{Provenance}{Rattachement selon l’EDB : \emph{Incertae originis} (le volume I des ICVR réunit les inscriptions dont la provenance est perdue).}
\nf{Support}{Support non précisé par l’EDB. Reproductions signalées par l’EDB : dessin, ICVR I (de Rossi) (droits réservés ; non reproduites). Fiche EDB : Filippo Piazzolla, première version 2014.}
\ntext{\lnum{1}Placi⸢d⸣o ⸢Fl⸣(avio) Romulo consulibus III Nona[s Iulias die]\\
\lnum{2}Martis decessit Leo annorum XXX PO[- - -]\\
\lnum{3}SV vixit ann(os) nn(u)mm(ero ?) X}
\nf{Témoins}{Export EDCS de 2020 : \emph{Placio (avio) Romulo consulibus III Nona[s Iulias die] / Martis decessit Leo annorum XXX PO[3] / SV vixit ann(os) nn(u)mm(ero ?) X}. EDB (EDB27782) : texte concordant aux conventions de notation près (indice de similarité 0,948).}
\nf{Apparat}{⸢d⸣ : T gravé ; ⸢Fl⸣ : IT gravé}
\nf{Traduction}{Sous le consulat de Placidus et de Flavius Romulus, le 3e jour avant les nones de juillet, un mardi, mourut Leo, âgé de 30 ans […] ; […] vécut en tout 10 (?) ans [texte corrompu].}
\nf{Langue}{L’éditeur signale 2 graphies corrigées : T gravé pour d ; IT gravé pour Fl.}
\nf{Chronologie}{Noms consulaires : Placidus (lu) et Romulus (lu). Dates du jour : événement indéterminé : III Nonas Iulias = 5 juillet 343 (partiellement restitué). Âge déclaré : 30 a. (lu).}
\nf{Bibliographie}{ICUR I 264 ; ILCV 4393. Bibliographie complémentaire selon l’EDB : ILCV, 4393.}
\nf{Statut}{Lecture secondaire (EDCS) concordante avec l’EDB : lecture stabilisée entre deux bases, non collationnée sur l’ICVR ni sur le monument. L’EDCS signale une photographie (non consultée).}
\end{notice}

\begin{notice}{8}{343}{ICUR n.s. I, 2080 · EDCS-25402137 · EDB15588}
\label{not:8}
\nf{Datation}{343 — directe : nom consulaire lu}
\nf{Genre}{Épitaphe ; ensembles A, B, C ; 37 lettres transmises.}
\nf{Références}{ICUR I 2080.}
\nf{Provenance}{Rattachement selon l’EDB : \emph{Incertae originis} (le volume I des ICVR réunit les inscriptions dont la provenance est perdue). Conservation selon l’EDB : \emph{Musei Vaticani: Lapidario Cristiano}.}
\nf{Support}{Sarcophage de marbre. 15 × 20 cm (h. × l.). Inscription gravée. Fiche EDB : Antonio Enrico Felle, première version 2003.}
\ntext{\lnum{1}κατάθεσις\\
\lnum{2}Ὀκταβίλλης\\
\lnum{3}πρὸ θʹ καλ(ανδῶν) σε(π)τ(εμβρίων)\\
\lnum{4}(ὑπατείᾳ) Ῥωμουλι}
\nf{Témoins}{Export EDCS de 2020 : \emph{texte grec absent}. EDB (EDB15588 ; écarts à examiner) : \emph{κατάθεσις / Ὀκταβίλλης / πρὸ θ´ καλ(ανδῶν) σε(π)τ(εμβρίων) / ⟨ὑπατείᾳ⟩ Ῥωμουλι [texte tronqué dans la liste de résultats]}.}
\nf{Traduction}{(grec) Déposition d’Octavilla, le 9e jour avant les calendes de septembre, sous le consulat de Romulus.}
\nf{Langue}{Le texte est rédigé en grec, en tout ou en partie ; l’export de 2020 ne le transmettait pas. 4 abréviations développées par l’éditeur (dont (ανδῶν), (π), (εμβρίων), (ὑπατείᾳ)).}
\nf{Chronologie}{Noms consulaires : Romulus (lu).}
\nf{Bibliographie}{ICUR I 2080. Bibliographie complémentaire selon l’EDB : Deichmann 1967, 70, no. 84 (with photo at taf. 24) ; CLRE, 633.}
\nf{Statut}{Lecture secondaire (EDCS) ; écarts avec l’EDB à arbitrer sur l’ICVR.}
\end{notice}

\begin{notice}{9}{343}{ICUR n.s. III, 8720 · EDCS-35900965 · EDB18778}
\label{not:9}
\nf{Datation}{343 — directe : nom consulaire lu}
\nf{Genre}{Épitaphe ; ensembles A, B, C ; 68 lettres transmises.}
\nf{Références}{ICUR III 8720 = ILCV 2571 (em) = ILCV 2819 (em).}
\nf{Provenance}{Rattachement selon l’EDB : \emph{Tituli qui in coem. Callisti reperti traduntur}.}
\nf{Support}{Plaque de marbre. Inscription gravée. Signes décrits par l’EDB : chrisme. Reproductions signalées par l’EDB : dessin, ICVR I (de Rossi) (droits réservés ; non reproduites). Fiche EDB : Carlo Carletti, première version 2003.}
\ntext{\lnum{1}Placido et Romulo (!) XVII ⸢K⸣(alendas) Dec(embres)\\
\lnum{2}Ianuaria in pace quae\\
\lnum{3}vixit annos XVIIII m(enses)\\
\lnum{4}VIII d(ies) XIIII}
\nf{Témoins}{Export EDCS de 2020 : \emph{Placido et Romulo (!) XVII (alendas) Dec(embres) / Ianuaria in pace quae / vixit annos XVIIII m(enses) / VIII d(ies) XIIII}. EDB (EDB18778) : texte concordant aux conventions de notation près (indice de similarité 0,988).}
\nf{Apparat}{⸢K⸣ : C gravé ; (!) : forme signalée par l’éditeur}
\nf{Traduction}{Sous (le consulat de) Placidus et Romulus, le 17e jour avant les calendes de décembre, Ianuaria, en paix, qui vécut 19 ans 8 mois 14 jours.}
\nf{Langue}{L’éditeur signale une graphie corrigée : C gravé pour K. L’éditeur signale au moins une forme anormale (sic). 4 abréviations développées par l’éditeur (dont (alendas), Dec(embres), m(enses), d(ies)).}
\nf{Chronologie}{Noms consulaires : Placidus (lu) et Romulus (lu). Dates du jour : événement indéterminé : XVII Kalendas Decembres = 15 novembre 343 (lu). Âge déclaré : 19 a. 8 m. 14 j. (lu).}
\nf{Commentaire}{La formule in pace est conservée.}
\nf{Bibliographie}{ICUR III 8720 ; ILCV 2571 (em) ; ILCV 2819 (em).}
\nf{Statut}{Lecture secondaire (EDCS) concordante avec l’EDB : lecture stabilisée entre deux bases, non collationnée sur l’ICVR ni sur le monument. L’EDCS signale une photographie (non consultée).}
\end{notice}

\annee{344}

\begin{notice}{10}{344}{ICUR n.s. IV, 11752 · EDCS-37400847 · EDB699}
\label{not:10}
\nf{Datation}{344 — directe : lecture partielle d’un nom consulaire}
\nf{Genre}{Épitaphe ; ensembles A, B, C ; 188 lettres transmises ; 66 \% restituées.}
\nf{Références}{ICUR IV 11752.}
\nf{Provenance}{Rattachement selon l’EDB : \emph{Coem. anonymum ad viam Appiam}. Conservation selon l’EDB : \emph{Coem. anonymum ad viam Appiam}.}
\nf{Support}{Plaque de marbre. Inscription gravée. Fiche EDB : Maria Cristina Grisanzio, première version 2003.}
\ntext{\lnum{1}[Her]mogenia [quae vixit a]nnis XXVI m[ensibus - - - diebus - - -]\\
\lnum{2}[recessit - - - Non(as) Mai(as) deposi]ta Nonis s(upra) s(criptis) u[xor Primi - - -]\\
\lnum{3}[hic requi]escit in pace [- - -] Primus co[mpari su(a)e fecit]\\
\lnum{4}[bene] merenti in pac[e Leonti]o et Salu[stio vv(iris) cc(larissimis) co(n)ss(ulibus)]}
\nf{Témoins}{EDB (EDB699 ; écarts à examiner) : \emph{[Her]mogenia [quae vixit a]nnis XXVI m[ensibus --- diebus ---] / [--- deposi]ta nonis s(upra) s(criptis) u[xor Primi ---] / [hic requi]escit in pace --- Primus co[mpari suae fecit] / [bene]merenti in pac[e Leonti]o et Salu[stio viris clarissimis co(n)s(ulibu)s] [texte tronqué dans la liste de résultats]}.}
\nf{Traduction}{Hermogenia, qui vécut 26 ans […] mois […] jours ; elle s’en alla […] (avant) les nones de mai ; déposée aux nones susdites ; épouse de Primus […] ; elle repose ici en paix ; […] Primus à sa compagne a fait (ce monument), bien méritante, en paix ; sous le consulat de Leontius et Sallustius, clarissimes.}
\nf{Langue}{9 abréviations développées par l’éditeur (dont Non(as), Mai(as), s(upra), s(criptis), su(a)e, vv(iris)).}
\nf{Chronologie}{Noms consulaires : Leontius (restitué) et Sallustius (absent). Dates du jour : décès : … Nonas Maias (restitué). Âge déclaré : 26 a. (partiellement restitué).}
\nf{Commentaire}{La formule in pace est conservée. Vocabulaire religieux : repos et sommeil de la mort — requiescit (partiellement restitué).}
\nf{Bibliographie}{ICUR IV 11752.}
\nf{Statut}{Lecture secondaire (EDCS) ; écarts avec l’EDB à arbitrer sur l’ICVR. Plus de la moitié des lettres sont restituées.}
\end{notice}

\annee{345}

\begin{notice}{11}{345}{ICUR n.s. I, 1104 · EDCS-38700644 · EDB30041}
\label{not:11}
\nf{Datation}{345 — directe : nom consulaire lu}
\nf{Genre}{Épitaphe ; ensembles A, B, C ; 33 lettres transmises ; 58 \% restituées.}
\nf{Références}{ICUR I 1104.}
\nf{Provenance}{Rattachement selon l’EDB : \emph{Incertae originis} (le volume I des ICVR réunit les inscriptions dont la provenance est perdue). Conservation selon l’EDB : \emph{s. Saba}.}
\nf{Support}{Plaque de marbre. 45 × 58 × 4,5 cm (h. × l. × ép.). Inscription gravée. Fiche EDB : Filippo Piazzolla, première version 2014.}
\ntext{\lnum{1}] Augustas Amanti[o et Albino cons(ulibus)}
\nf{Témoins}{EDB (EDB30041 ; écarts à examiner) : \emph{------ / ạụg̣ụṣṭạṣ Ạṃạṇṭị[o et Albino cons(ulibu)s ---] [texte tronqué dans la liste de résultats]}.}
\nf{Traduction}{[…] août, sous le consulat d’Amantius et d’Albinus.}
\nf{Langue}{1 abréviation développée par l’éditeur (dont cons(ulibus)).}
\nf{Chronologie}{Noms consulaires : Amantius (lu) et Albinus (restitué).}
\nf{Bibliographie}{ICUR I 1104.}
\nf{Statut}{Lecture secondaire (EDCS) ; écarts avec l’EDB à arbitrer sur l’ICVR. Plus de la moitié des lettres sont restituées.}
\end{notice}

\annee{352}

\begin{notice}{12}{352}{ICUR n.s. II, 4798 · EDCS-36101279}
\label{not:12}
\nf{Datation}{352 — directe : nom consulaire lu}
\nf{Genre}{Épitaphe ; ensembles A, B, C ; 83 lettres transmises ; 24 \% restituées.}
\nf{Références}{ICUR II 4798.}
\nf{Provenance}{Rome (indication de l’EDCS) ; aucun rattachement cimétérial transmis.}
\nf{Support}{Non documenté dans les sources consultées.}
\ntext{\lnum{1}]teni sorori castissim[ae - - -]\\
\lnum{2}[- - -]b(res) Amantio et Albino con[s(ulibus) - - -]\\
\lnum{3}[- - -] XXIII d(e)p(osita) XIIII K(alendas) Aug(ustas) [Decentio]\\
\lnum{4}[et P]aulo conss(ulibus)}
\nf{Témoins}{Aucune fiche EDB appariée (numéro ICUR absent de la liste publique de l’EDB).}
\nf{Traduction}{[…] à sa sœur très chaste […] ; […] (mois) sous le consulat d’Amantius et d’Albinus […] ; […] 23 […] déposée le 14e jour avant les calendes d’août, sous le consulat de Décence et Paulus.}
\nf{Langue}{7 abréviations développées par l’éditeur (dont b(res), cons(ulibus), d(e)p, (osita), K(alendas), Aug(ustas)).}
\nf{Chronologie}{Noms consulaires : Amantius (lu), Albinus (lu), Decentius (restitué) et Paulus (lu en partie). Dates du jour : déposition : XIIII Kalendas Augustas = 19 juillet 352 (lu).}
\nf{Bibliographie}{ICUR II 4798.}
\nf{Statut}{Lecture secondaire (EDCS) sans second témoin numérique. L’EDCS signale une photographie (non consultée).}
\end{notice}

\annee{345}

\begin{notice}{13}{345}{ICUR n.s. II, 4799 · EDCS-36101280}
\label{not:13}
\nf{Datation}{345 — directe : nom consulaire lu}
\nf{Genre}{Épitaphe ; ensembles A, B, C ; 85 lettres transmises ; 21 \% restituées.}
\nf{Références}{ICUR II 4799.}
\nf{Provenance}{Rome (indication de l’EDCS) ; aucun rattachement cimétérial transmis.}
\nf{Support}{Non documenté dans les sources consultées.}
\ntext{\lnum{1}qua]e vi⸢x⸣it an(n)is XXV[- - -]\\
\lnum{2}[- - - f]ecit compar[i - - -]\\
\lnum{3}[- - -]i filiae bene mer[enti - - -]\\
\lnum{4}[- - - e]t deposita est [- - -]\\
\lnum{5}[- - -]di cons(ulato) Amantio [et]\\
\lnum{6}[Albino con]ss(ulibus)}
\nf{Témoins}{Export EDCS de 2020 : \emph{qua]e viit an(n)is XXV[3] / [3 f]ecit compar[i 3] / [3]i filiae bene mer[enti 3] / [3 e]t deposita est [3] / [3]di cons(ulato) Amantio [et] / [Albino con]ss(ulibus)}. Aucune fiche EDB appariée (numéro ICUR absent de la liste publique de l’EDB).}
\nf{Apparat}{⸢x⸣ : S gravé}
\nf{Traduction}{[…] qui vécut 25 ans […] ; […] a fait à sa compagne […] ; […] à sa fille bien méritante […] ; […] et elle a été déposée […] ; […] sous le consulat d’Amantius et d’Albinus.}
\nf{Langue}{L’éditeur signale une graphie corrigée : S gravé pour x (s pour x).}
\nf{Chronologie}{Noms consulaires : Amantius (lu) et Albinus (restitué).}
\nf{Bibliographie}{ICUR II 4799.}
\nf{Statut}{Lecture secondaire (EDCS) sans second témoin numérique.}
\end{notice}

\begin{notice}{14}{345}{ICUR n.s. IV, 11753 · EDCS-37400848 · EDB700}
\label{not:14}
\nf{Datation}{345 — directe : nom consulaire lu}
\nf{Genre}{Épitaphe ; ensembles A, B, C ; 63 lettres transmises ; 75 \% restituées.}
\nf{Références}{ICUR IV 11753.}
\nf{Provenance}{Rattachement selon l’EDB : \emph{Coem. anonymum ad viam Appiam}. Conservation selon l’EDB : \emph{Coem. anonymum ad viam Appiam}.}
\nf{Support}{Plaque de marbre. Inscription gravée. Reproductions signalées par l’EDB : photographie PCAS (droits réservés ; non reproduites). Fiche EDB : Maria Cristina Grisanzio, première version 2003.}
\ntext{\lnum{1}Dorm[itio - - - vixit]\\
\lnum{2}annis X[- - -]\\
\lnum{3}Amanti[o et Albino vv(iris) cc(larissimis) conss(ulibus)]}
\nf{Témoins}{EDB (EDB700 ; écarts à examiner) : \emph{dorm[itio --- q(---) vixit] / annis X [---] / Amanti[o et Albino viris clarissimis co(n)s(ulibu)s] [texte tronqué dans la liste de résultats]}.}
\nf{Traduction}{Endormissement (dormitio) de […] ; il/elle vécut […] ans […] ; sous le consulat d’Amantius et d’Albinus, clarissimes.}
\nf{Langue}{3 abréviations développées par l’éditeur (dont vv(iris), cc(larissimis), conss(ulibus)).}
\nf{Chronologie}{Noms consulaires : Amantius (lu) et Albinus (restitué).}
\nf{Commentaire}{Vocabulaire religieux : repos et sommeil de la mort — Dormitio (partiellement restitué).}
\nf{Bibliographie}{ICUR IV 11753.}
\nf{Statut}{Lecture secondaire (EDCS) ; écarts avec l’EDB à arbitrer sur l’ICVR. Plus de la moitié des lettres sont restituées. L’EDCS signale une photographie (non consultée).}
\end{notice}

\annee{346}

\begin{notice}{15}{346}{ICUR n.s. I, 1930 · EDCS-20000191 · EDB27212}
\label{not:15}
\nf{Datation}{346 — directe : nom consulaire lu}
\nf{Genre}{Épitaphe ; ensembles A, B, C ; 128 lettres transmises ; 11 \% restituées.}
\nf{Références}{CIL VI 37122 (p 4820) = ILCV 162 (add) = ICUR I 1930 = AE 1907, 133.}
\nf{Provenance}{Rattachement selon l’EDB : \emph{Incertae originis} (le volume I des ICVR réunit les inscriptions dont la provenance est perdue). Conservation selon l’EDB : \emph{Museo Archeologico Nazionale}.}
\nf{Support}{Sarcophage de marbre. 17 × 130 × 15,5 cm (h. × l. × ép.). Inscription gravée. Fiche EDB : Filippo Piazzolla, première version 2015.}
\ntext{\lnum{1}]\\
\lnum{2}Cassia Pisonis c(larissima) f(emina) vixit annis XXX diebus XCIII nata Sabino et Rufino [co(n)s(ulibus)]\\
\lnum{3}XVII Kal(endas) Iulias reddit pos(t) consulatum Amanti et [Albi]ni XVII Kal(endas) Octobres}
\nf{Témoins}{EDB (EDB27212) : texte concordant aux conventions de notation près (indice de similarité 0,983).}
\nf{Traduction}{[…] Cassia, fille de Piso, clarissime, vécut 30 ans et 93 jours ; née sous le consulat de Sabinus et Rufinus, le 17e jour avant les calendes de juillet ; elle rendit (l’âme) après le consulat d’Amantius et d’Albinus, le 17e jour avant les calendes d’octobre.}
\nf{Langue}{7 abréviations développées par l’éditeur (dont c(larissima), f(emina), co(n)s, (ulibus), Kal(endas), pos(t)).}
\nf{Chronologie}{Noms consulaires : Amantius (lu), Albinus (restitué) et Rufinus (lu). Datation par post-consulat. Dates du jour : naissance : XVII Kalendas Iulias = 15 juin 346 (lu) ; décès : XVII Kalendas Octobres = 15 septembre 346 (lu). Âge déclaré : 30 a. 93 j. (lu).}
\nf{Commentaire}{Rang sénatorial d’une personne commémorée (cassia pisonis clarissima femina vixi). Commentaire développé : dossier p.~\pageref{dos:15}.}
\nf{Bibliographie}{CIL VI 37122 (p 4820) ; ILCV 162 (add) ; ICUR I 1930 ; AE 1907, 133. Bibliographie complémentaire selon l’EDB : CIL VI, 37122, cfr. p. 4820 ; ILCV, 162, cfr. II p. 507 ; Ferrua 1986a, 50, n. 1930.}
\nf{Statut}{Lecture secondaire (EDCS) concordante avec l’EDB : lecture stabilisée entre deux bases, non collationnée sur l’ICVR ni sur le monument. L’EDCS signale une photographie (non consultée).}
\end{notice}

\begin{notice}{16}{346}{ICUR n.s. I, 3126 · EDCS-39900588 · EDB33080}
\label{not:16}
\nf{Datation}{346 — directe : nom consulaire lu}
\nf{Genre}{Épitaphe ; ensembles A, B, C ; 82 lettres transmises.}
\nf{Références}{CIL X *949,05 = ICUR I 3126 = ILCV 2602 (em).}
\nf{Provenance}{Rattachement selon l’EDB : \emph{Incertae originis} (le volume I des ICVR réunit les inscriptions dont la provenance est perdue). Conservation selon l’EDB : \emph{Velletri (Roma)   ; Museo Municipale}.}
\nf{Support}{Plaque de marbre. 23 × 72 cm (h. × l.). Inscription gravée. Signes décrits par l’EDB : chrisme. Reproductions signalées par l’EDB : dessin, ICVR I (de Rossi) et reproduction, Perret 1851 (droits réservés ; non reproduites). Fiche EDB : Filippo Piazzolla, première version 2015.}
\ntext{\lnum{1}Bonoso bene merenti in pace\\
\lnum{2}qui vixit annis II m(ensibus) III d(iebus) XX\\
\lnum{3}dep(ositus) prid(ie) Idus Sept(embres) post\\
\lnum{4}cons\{s\}(ulatu) Amanti et Albini}
\nf{Témoins}{EDB (EDB33080 ; écarts à examiner) : \emph{⟨:a⟩ / ((monogramma Christi)) / ⟨:b⟩ / Bonoso benemerenti in pace / qui vixit annis II m(ensibus) III d(iebus) XX / dep(ositus) prid(ie) idus sept(embres) post / cons(ulibu)s Amanti et Albini ((monogramma Christi)) [texte tronqué dans la liste de résultats]}.}
\nf{Apparat}{\{s\} : lettres superflues}
\nf{Traduction}{À Bonosus, bien méritant, en paix, qui vécut 2 ans 3 mois 20 jours ; déposé la veille des ides de septembre, après le consulat d’Amantius et d’Albinus.}
\nf{Langue}{Des lettres gravées sont jugées superflues par l’éditeur (entre accolades). 6 abréviations développées par l’éditeur (dont m(ensibus), d(iebus), dep(ositus), prid(ie), Sept(embres), (ulatu)).}
\nf{Chronologie}{Noms consulaires : Amantius (lu) et Albinus (lu). Datation par post-consulat. Dates du jour : déposition : pridie Idus Septembres = 12 septembre 346 (lu). Âge déclaré : 2 a. 3 m. 20 j. (lu).}
\nf{Commentaire}{La formule in pace est conservée.}
\nf{Bibliographie}{CIL X *949,05 ; ICUR I 3126 ; ILCV 2602 (em). Bibliographie complémentaire selon l’EDB : CIL X, 949\textbackslash{}*,5 ; ILCV, 2602.}
\nf{Statut}{Lecture secondaire (EDCS) ; écarts avec l’EDB à arbitrer sur l’ICVR. L’EDCS signale une photographie (non consultée).}
\end{notice}

\begin{notice}{17}{346}{ICUR n.s. I, 1421 · EDCS-39000408 · EDB31185}
\label{not:17}
\nf{Datation}{346 — directe : nom consulaire lu}
\nf{Genre}{Épitaphe ; ensembles A, B, C ; 93 lettres transmises.}
\nf{Références}{ICUR I 1421 = ILCV 2627.}
\nf{Provenance}{Rattachement selon l’EDB : \emph{Incertae originis} (le volume I des ICVR réunit les inscriptions dont la provenance est perdue). Conservation selon l’EDB : \emph{Musei Vaticani: Lapidario Cristiano}.}
\nf{Support}{Plaque de marbre. 20 × 78 cm (h. × l.). Inscription gravée. Reproductions signalées par l’EDB : dessin, ICVR I (de Rossi) (droits réservés ; non reproduites). Fiche EDB : Filippo Piazzolla, première version 2015.}
\ntext{\lnum{1}B(ene) m(erenti) in p(ace) Cheretio q(ui) v(ixit) ann(os) V m(enses) V d(ies) XV\\
\lnum{2}Cheretius et Severa parentes\\
\lnum{3}d(e)p(ositus) VI Kal(endas) Septembres\\
\lnum{4}pos(t) consulatum Amanti et Albini}
\nf{Témoins}{EDB (EDB31185) : texte concordant aux conventions de notation près (indice de similarité 0,944).}
\nf{Traduction}{Au bien méritant, en paix, Cheretius, qui vécut 5 ans 5 mois 15 jours ; Cheretius et Severa, ses parents ; déposé le 6e jour avant les calendes de septembre, après le consulat d’Amantius et d’Albinus.}
\nf{Langue}{12 abréviations développées par l’éditeur (dont B(ene), m(erenti), p(ace), q(ui), v(ixit), ann(os)).}
\nf{Chronologie}{Noms consulaires : Amantius (lu) et Albinus (lu). Datation par post-consulat. Dates du jour : déposition : VI Kalendas Septembres = 27 août 346 (lu). Âge déclaré : 5 a. 5 m. 15 j. (lu).}
\nf{Commentaire}{La formule in pace est transmise sous forme abrégée, sans restitution de lacune ; son développement reste éditorial.}
\nf{Bibliographie}{ICUR I 1421 ; ILCV 2627. Bibliographie complémentaire selon l’EDB : ILCV, 2627 ; Lega 2000, 51 ; Ferrua 1964, 108 ; Ferrua 1979b, 36, n. 1421.}
\nf{Statut}{Lecture secondaire (EDCS) concordante avec l’EDB : lecture stabilisée entre deux bases, non collationnée sur l’ICVR ni sur le monument. L’EDCS signale une photographie (non consultée).}
\end{notice}

\begin{notice}{18}{346}{ICUR n.s. IV, 11754 · EDCS-35700937 · EDB965}
\label{not:18}
\nf{Datation}{346 — directe : nom consulaire lu}
\nf{Genre}{Épitaphe ; ensembles A, B, C ; 61 lettres transmises ; 31 \% restituées.}
\nf{Références}{ICUR IV 11754 = ILCV 1126b.}
\nf{Provenance}{Rattachement selon l’EDB : \emph{Coem. anonymum ad viam Appiam}. Conservation selon l’EDB : \emph{Coem. anonymum ad viam Appiam}.}
\nf{Support}{Paroi de enduit. Inscription peinte. Fiche EDB : Maria Cristina Grisanzio, première version 2003.}
\ntext{\lnum{1}]\\
\lnum{2}iterum\\
\lnum{3}con[- - - ecclesiae(?)]\\
\lnum{4}catholica[e]\\
\lnum{5}d(epositio) V K(alendas) Ma[- - -]\\
\lnum{6}post con[s(ulatum)]\\
\lnum{7}Aman[ti]\\
\lnum{8}et Albini}
\nf{Témoins}{EDB (EDB965) : texte concordant aux conventions de notation près (indice de similarité 0,92).}
\nf{Traduction}{[…] pour la seconde fois […] de l’Église (?) catholique ; déposition le 5e jour avant les calendes de mai (ou mars) ; après le consulat d’Amantius et d’Albinus.}
\nf{Langue}{4 abréviations développées par l’éditeur (dont ecclesiae(?), d(epositio), K(alendas), cons(ulatum)).}
\nf{Chronologie}{Noms consulaires : Amantius (lu en partie) et Albinus (lu). Datation par post-consulat.}
\nf{Commentaire}{Vocabulaire religieux : vocabulaire ecclésial — ecclesiae (entièrement restitué).}
\nf{Bibliographie}{ICUR IV 11754 ; ILCV 1126b.}
\nf{Statut}{Lecture secondaire (EDCS) concordante avec l’EDB : lecture stabilisée entre deux bases, non collationnée sur l’ICVR ni sur le monument.}
\end{notice}

\annee{347}

\begin{notice}{19}{347}{ICUR n.s. IV, 10851 · EDCS-19000686 · EDB39161}
\label{not:19}
\nf{Datation}{347 — directe : nom consulaire lu}
\nf{Genre}{Épitaphe ; ensembles A, B, C ; 200 lettres transmises ; 26 \% restituées.}
\nf{Références}{CIL VI 9237 = ILCV 589 = ICUR IV 10851.}
\nf{Provenance}{Rattachement selon l’EDB : \emph{Coem. Callisti pars superior}. Conservation selon l’EDB : \emph{Coem. Callisti pars superior}.}
\nf{Support}{Sarcophage de marbre. 31 × 88 cm (h. × l.). Inscription gravée. Fiche EDB : Marida Pierno, première version 2015.}
\ntext{\lnum{1}[In ho]c sarcofa⸢g⸣o conditus\\
\lnum{2}[est S]antulus Fl(avius) Aur(elius) Casarius\\
\lnum{3}[- - - A]ureli et [Vic]torin(a)e\\
\lnum{4}[c(larissimiae)] m(emoriae) f(eminae) venit [ante annos d]uos\\
\lnum{5}Preniste cu[m compare sua Lea(?)]\\
\lnum{6}et Lea vixit ann[os - - -]\\
\lnum{7}et dies XXI decess[it]\\
\lnum{8}autem VI Idus Iulias\\
\lnum{9}Amantio et Albino c(on)s(ulibus)\\
\lnum{10}Romae vixit d(ies) XXII[}
\nf{Témoins}{Export EDCS de 2020 : \emph{[In ho]c sarcofao conditus / [est S]antulus Fl(avius) Aur(elius) Casarius / [3 A]ureli et [Vic]torin(a)e / [c(larissimiae)] m(emoriae) f(eminae) venit [ante annos d]uos / Preniste cu[m compare sua Lea(?)] // et Lea vixit ann[os 3] / et dies XXI decess[it] / autem VI Idus Iulias / Amantio et Albino c(on)s(ulibus) / Romae vixit d(ies) XXII[}. EDB (EDB39161 ; divergence forte) : \emph{[in ho]c sarcofago conditus / [est S]anctulus Fl(avius) Aur(elius) casarius / [--- A]urelii et [Vic]torine / [---] m(onumentum?) f(ecerunt?) ++++[-c.10-] / Preniste cu[m ----]uos}. L'EDB (EDB39161 = ICVR IV 10851 = CIL VI 9237) restitue beaucoup moins que l'EDCS : là où l'EDCS donne \emph{[c(larissimae)] m(emoriae) f(eminae) venit [ante annos d]uos … cu[m compare sua Lea(?)]}, l'EDB lit \emph{m(onumentum?) f(ecerunt?) ++++ … cu[m ---]uos}, et \emph{dies XX} au lieu de \emph{XXI}. Le support est un sarcophage de marbre (31 × 88 cm). Ni le nom de Lea, ni la qualité clarissime, ni l'indication « il y a deux ans » ne peuvent être exploités avant contrôle de l'ICVR.}
\nf{Apparat}{⸢g⸣ : C gravé}
\nf{Traduction}{Dans ce sarcophage a été placé Santulus Flavius Aurelius Casarius, […] (fils ?) d’Aurelius et de Victorina, femme clarissime de bonne mémoire ; il était venu il y a deux ans de Préneste avec sa compagne Lea ; // et Lea vécut […] ans et 21 jours ; elle mourut le 6e jour avant les ides de juillet, sous le consulat d’Amantius et d’Albinus ; à Rome elle vécut 22 jours (?).}
\nf{Langue}{L’éditeur signale une graphie corrigée : C gravé pour g (c pour g).}
\nf{Chronologie}{Noms consulaires : Amantius (lu) et Albinus (lu). Le décès de Lea est daté par le consulat d'Amantius et Albinus (345) ; l'année 347 de la base paraît déduite de l'indication restituée \emph{[ante annos d]uos}. Cette déduction repose sur une restitution que l'EDB ne retient pas : l'année de classement 347 est conservée sous réserve, et le texte est daté au plus tôt de 345. Dates du jour : décès : VI Idus Iulias = 10 juillet 347 (lu). Âge déclaré : 21 j. (partiellement restitué).}
\nf{Commentaire}{Rang sénatorial d’une personne commémorée (ius casarius … aureli et victorinae clarissimiae memo).}
\nf{Bibliographie}{CIL VI 9237 ; ILCV 589 ; ICUR IV 10851. Bibliographie complémentaire selon l’EDB : CIL VI, 9237 ; ILCV, 589 ; Bisconti 2000a, 209 ; Dresken-Weiland 2003, 220.}
\nf{Statut}{Lecture secondaire (EDCS) ; divergence forte avec l’EDB : le texte ne doit pas être exploité sans contrôle de l’édition imprimée.}
\end{notice}

\begin{notice}{20}{347}{ICUR n.s. X, 26312 · EDCS-01000568 · EDB10655}
\label{not:20}
\nf{Datation}{347 — directe : nom consulaire lu}
\nf{Genre}{Épitaphe ; ensembles A, B, C ; 136 lettres transmises ; 11 \% restituées.}
\nf{Références}{CIL VI 41424a = ICUR X 26312 = ICaRoma 70 = RAC-1924-82 = AE 1988, 34.}
\nf{Provenance}{Rattachement selon l’EDB : \emph{Pars infima coem. Pamphili}. Conservation selon l’EDB : \emph{Pars infima coem. Pamphili}.}
\nf{Support}{Plaque de marbre. 84,5 × 114,5 × 3,4 cm (h. × l. × ép.). Inscription gravée. Reproductions signalées par l’EDB : photographie PCAS (droits réservés ; non reproduites). Fiche EDB : Antonio Enrico Felle, première version 2014.}
\ntext{\lnum{1}In hoc loco positus est Fl(avius) Gaian[us ex]\\
\lnum{2}praesidibus qui vixit annis n(umero) [- - -]\\
\lnum{3}defunctus est autem Rufino e[t Eusebio]\\
\lnum{4}conss(ulibus) die VIIII Kal(endas) Mart(ias) depositus [est]\\
\lnum{5}VIII Kal(endas) Mart(ias) in pace}
\nf{Témoins}{EDB (EDB10655) : texte concordant aux conventions de notation près (indice de similarité 1).}
\nf{Traduction}{En ce lieu a été placé Flavius Gaianus, ancien gouverneur, qui vécut […] ans ; il mourut sous le consulat de Rufinus et d’Eusebius, le 9e jour avant les calendes de mars ; il fut déposé le 8e jour avant les calendes de mars, en paix.}
\nf{Langue}{7 abréviations développées par l’éditeur (dont Fl(avius), n(umero), conss(ulibus), Kal(endas), Mart(ias)).}
\nf{Chronologie}{Noms consulaires : Rufinus (lu) et Eusebius (restitué). Dates du jour : événement indéterminé : die VIIII Kalendas Martias = 21 février 347 (lu) ; déposition : VIII Kalendas Martias = 22 février 347 (lu). Âge déclaré : 9 j. (lu).}
\nf{Commentaire}{La formule in pace est conservée. Fonction ou métier mentionné : praesidibus.}
\nf{Bibliographie}{CIL VI 41424a ; ICUR X 26312 ; ICaRoma 70 ; RAC-1924-82 ; AE 1988, 34. Bibliographie complémentaire selon l’EDB : Carletti 1986a, 82-83, n. 70 ; CIL VI, 41424a.}
\nf{Statut}{Lecture secondaire (EDCS) concordante avec l’EDB : lecture stabilisée entre deux bases, non collationnée sur l’ICVR ni sur le monument. L’EDCS signale une photographie (non consultée).}
\end{notice}

\begin{notice}{21}{347}{ICUR n.s. I, 4003 · EDCS-34801295 · EDB12186}
\label{not:21}
\nf{Datation}{347 — directe : lecture partielle d’un nom consulaire}
\nf{Genre}{Épitaphe ; ensembles A, B, C ; 115 lettres transmises ; 56 \% restituées.}
\nf{Références}{ICUR I 4003.}
\nf{Provenance}{Rattachement selon l’EDB : \emph{Incertae originis} (le volume I des ICVR réunit les inscriptions dont la provenance est perdue).}
\nf{Support}{Plaque de marbre. Inscription gravée. Fiche EDB : Antonio Enrico Felle, première version 2003.}
\ntext{\lnum{1}] quae fecit cum m[arito ann(os) - - -]\\
\lnum{2}[quae vixit a]nn(os) XXXV m(enses) II d(ies) X d[ep(osita) - - -]\\
\lnum{3}[- - - Ἡ]λιοδώρα ἐτῶν τρ[ιάκοντα πέντε]\\
\lnum{4}[ὑπατείᾳ Ῥουφίνου καὶ Εὐ]σεβίου μηνὶ μαΐ[ου τό - - -]}
\nf{Témoins}{Export EDCS de 2020 : \emph{] quae fecit cum m[arito ann(os) 3] / [quae vixit a]nn(os) XXXV m(enses) II d(ies) X d[ep(osita) 3] //}. EDB (EDB12186) : texte concordant aux conventions de notation près (indice de similarité 0,94).}
\nf{Traduction}{[…] qui a vécu avec son mari […] ans ; qui vécut 35 ans 2 mois 10 jours ; déposée […] // (grec) […] Héliodôra, de trente-cinq ans, sous le consulat de Rufinus et d’Eusebius, au mois de mai […].}
\nf{Langue}{Le texte est rédigé en grec, en tout ou en partie ; l’export de 2020 ne le transmettait pas. 5 abréviations développées par l’éditeur (dont ann(os), m(enses), d(ies), dep(osita)).}
\nf{Chronologie}{Noms consulaires : Rufinus (restitué) et Eusebius (lu en partie). Âge déclaré : 35 a. 2 m. 10 j. (partiellement restitué).}
\nf{Bibliographie}{ICUR I 4003.}
\nf{Statut}{Lecture secondaire (EDCS) concordante avec l’EDB : lecture stabilisée entre deux bases, non collationnée sur l’ICVR ni sur le monument. Plus de la moitié des lettres sont restituées.}
\end{notice}

\annee{348}

\begin{notice}{22}{348}{ICUR n.s. I, 887 · EDCS-24100635 · EDB29417}
\label{not:22}
\nf{Datation}{348 — directe : nom consulaire lu}
\nf{Genre}{Épitaphe ; ensembles A, B, C ; 273 lettres transmises.}
\nf{Références}{CIL VI 33865 = CLE 659 = ICUR I 887 = ILCV 748 (add).}
\nf{Provenance}{Rattachement selon l’EDB : \emph{Incertae originis} (le volume I des ICVR réunit les inscriptions dont la provenance est perdue).}
\nf{Support}{Texte métrique. Signes décrits par l’EDB : chrisme et chrisme radié. Reproductions signalées par l’EDB : dessin, ICVR I (de Rossi) (droits réservés ; non reproduites). Fiche EDB : Filippo Piazzolla, première version 2014.}
\ntext{\lnum{1}Felix vita viri felicior exitus ipse\\
\lnum{2}Caiani semper crescit per saecula nomen\\
\lnum{3}nescit fama mori sed semper vi⸢v⸣it ubique\\
\lnum{4}advenit hospes Romanus princeps in urbem\\
\lnum{5}qui fuit hic primum iuris consultor amicus\\
\lnum{6}quiescit in pace\{m\} depositus diem quar\\
\lnum{7}tum Nonas Augu\\
\lnum{8}stas Flavio Filippo\\
\lnum{9}et Flavio Sallea con\\
\lnum{10}sulibus pater Sabba\\
\lnum{11}tius fecit}
\nf{Témoins}{Export EDCS de 2020 : \emph{Felix vita viri felicior exitus ipse / Caiani semper crescit per saecula nomen / nescit fama mori sed semper viit ubique / advenit hospes Romanus princeps in urbem / qui fuit hic primum iuris consultor amicus / quiescit in pace\{m\} depositus diem quar/tum Nonas Augu/stas Flavio Filippo / et Flavio Sallea con/sulibus pater Sabba/tius fecit}. EDB (EDB29417) : texte concordant aux conventions de notation près (indice de similarité 0,976).}
\nf{Apparat}{⸢v⸣ : B gravé ; \{m\} : lettres superflues}
\nf{Traduction}{(vers) Heureuse fut la vie de l’homme, plus heureuse encore sa fin : le nom de Caianus ne cesse de croître au fil des siècles ; la renommée ne sait mourir, mais vit partout. Il vint, hôte romain, dans la ville du prince (?), lui qui fut ici d’abord jurisconsulte et ami. Il repose en paix ; déposé le 4e jour avant les nones d’août, sous le consulat de Flavius Philippus et de Flavius Salia ; son père Sabbatius a fait (ce monument).}
\nf{Langue}{L’éditeur signale une graphie corrigée : B gravé pour v (bêtacisme). Des lettres gravées sont jugées superflues par l’éditeur (entre accolades).}
\nf{Chronologie}{Noms consulaires : Philippus (lu). Dates du jour : déposition : diem quartum Nonas Augustas = 2 août 348 (lu).}
\nf{Commentaire}{La formule in pace est conservée. Vocabulaire religieux : repos et sommeil de la mort — quiescit (lu). Commentaire développé : dossier p.~\pageref{dos:22}.}
\nf{Bibliographie}{CIL VI 33865 ; CLE 659 ; ICUR I 887 ; ILCV 748 (add). Bibliographie complémentaire selon l’EDB : CLE I (nn. 1-859), 659 ; CIL VI, 33865 ; ILCV, 748 ; PLRE I, 378 ; Caldelli 2001, 21, nt. 2 ; Tantillo 2006, 275, nt. 29 ; Krummrey 1965, 329, nt. 41.}
\nf{Statut}{Lecture secondaire (EDCS) concordante avec l’EDB : lecture stabilisée entre deux bases, non collationnée sur l’ICVR ni sur le monument. L’EDCS signale une photographie (non consultée).}
\end{notice}

\begin{notice}{23}{348}{ICUR n.s. I, 81 · EDCS-38600097 · EDB27454}
\label{not:23}
\nf{Datation}{348 — directe : nom consulaire lu}
\nf{Genre}{Épitaphe ; ensembles A, B, C ; 94 lettres transmises ; 16 \% restituées.}
\nf{Références}{ICUR I 81 = ILCV 2940.}
\nf{Provenance}{Rattachement selon l’EDB : \emph{Incertae originis} (le volume I des ICVR réunit les inscriptions dont la provenance est perdue). Conservation selon l’EDB : \emph{s. Cecilia in Trastevere}.}
\nf{Support}{Plaque de marbre. 70 × 185 × 2,5 cm (h. × l. × ép.). Inscription gravée. Reproductions signalées par l’EDB : dessin, ICVR I (de Rossi) (droits réservés ; non reproduites). Fiche EDB : Filippo Piazzolla, première version 2013.}
\ntext{\lnum{1}Artorius Amachius Aurel[iae - - -]\\
\lnum{2}uxori dulcissim(a)e sibi et E[- - -]\\
\lnum{3}deposita XIII Kal(endas) Dece⸢m⸣[bres - - -]\\
\lnum{4}Filippo et Sa(l)ia con[ss(ulibus)]}
\nf{Témoins}{Export EDCS de 2020 : \emph{Artorius Amachius Aurel[iae 3] / uxori dulcissim(a)e sibi et E[3] / deposita XIII Kal(endas) Dece[bres 3] / Filippo et Sa(l)ia con[ss(ulibus)]}. EDB (EDB27454) : texte concordant aux conventions de notation près (indice de similarité 0,933).}
\nf{Apparat}{⸢m⸣ : N gravé}
\nf{Traduction}{Artorius Amachius à Aurelia […], son épouse très douce, et pour lui-même et E[…] ; déposée le 13e jour avant les calendes de décembre […] ; sous le consulat de Philippus et de Salia.}
\nf{Langue}{L’éditeur signale une graphie corrigée : N gravé pour m (n pour m).}
\nf{Chronologie}{Noms consulaires : Philippus (lu) et Salia (lu). Dates du jour : déposition : XIII Kalendas Decembres = 19 novembre 348 (partiellement restitué).}
\nf{Bibliographie}{ICUR I 81 ; ILCV 2940.}
\nf{Statut}{Lecture secondaire (EDCS) concordante avec l’EDB : lecture stabilisée entre deux bases, non collationnée sur l’ICVR ni sur le monument. L’EDCS signale une photographie (non consultée).}
\end{notice}

\begin{notice}{24}{348}{ICUR n.s. I, 318 · EDCS-38600315 · EDB27940}
\label{not:24}
\nf{Datation}{348 — directe : nom consulaire lu}
\nf{Genre}{Épitaphe ; ensembles A, B, C ; 62 lettres transmises ; 42 \% restituées.}
\nf{Références}{ICUR I 318 = ILCV 3797a.}
\nf{Provenance}{Rattachement selon l’EDB : \emph{Incertae originis} (le volume I des ICVR réunit les inscriptions dont la provenance est perdue).}
\nf{Support}{Support non précisé par l’EDB. Signes décrits par l’EDB : croix. Reproductions signalées par l’EDB : dessin, ICVR I (de Rossi) (droits réservés ; non reproduites). Fiche EDB : Filippo Piazzolla, première version 2014.}
\ntext{\lnum{1}[- - -]terio Avi[- - -]\\
\lnum{2}[- - - bi]som[u]m PEI[- - -]\\
\lnum{3}[- - - se] vi⸢v⸣o sibi et [- - -]\\
\lnum{4}[- - - Id]us Ian(uarias) Philip[po et Salia cons(ulibus)}
\nf{Témoins}{Export EDCS de 2020 : \emph{[3]terio Avi[3] / [3 bi]som[u]m PEI[3] / [3 se] vio sibi et [3] / [3 Id]us Ian(uarias) Philip[po et Salia cons(ulibus)}. EDB (EDB27940 ; écarts à examiner) : \emph{[---]terio ((crux quadrata)) Avi̲[---] / [--- bi]som[u]m ((crux quadrata)) pe+̲[---] / [--- se] vibo sibi ((crux quadrata)) et [ ---] / [--- id]us ian(uarias) Philip[po et Salia cons(ulibu)s] / ------ [texte tronqué dans la liste de résultats]}.}
\nf{Apparat}{⸢v⸣ : B gravé}
\nf{Traduction}{[…] Avi[…] ; […] tombe double (bisomum) […] ; de son vivant, pour lui-même et […] ; […] ides de janvier, sous le consulat de Philippus et de Salia.}
\nf{Langue}{L’éditeur signale une graphie corrigée : B gravé pour v (bêtacisme).}
\nf{Chronologie}{Noms consulaires : Philippus (lu) et Salia (restitué). Dates du jour : événement indéterminé : … Idus Ianuarias (partiellement restitué).}
\nf{Bibliographie}{ICUR I 318 ; ILCV 3797a. Bibliographie complémentaire selon l’EDB : ILCV, 3797A.}
\nf{Statut}{Lecture secondaire (EDCS) ; écarts avec l’EDB à arbitrer sur l’ICVR. L’EDCS signale une photographie (non consultée).}
\end{notice}

\begin{notice}{25}{348}{ICUR n.s. II, 4800 · EDCS-36101281}
\label{not:25}
\nf{Datation}{348 — directe : lecture partielle d’un nom consulaire}
\nf{Genre}{Épitaphe ; ensembles A, B, C ; 66 lettres transmises ; 48 \% restituées.}
\nf{Références}{ICUR II 4800.}
\nf{Provenance}{Rome (indication de l’EDCS) ; aucun rattachement cimétérial transmis.}
\nf{Support}{Non documenté dans les sources consultées.}
\ntext{\lnum{1}Domino [- - - bene me]\\
\lnum{2}renti in[- - - qui]\\
\lnum{3}vixit ann[os - - -]\\
\lnum{4}d⸢e⸣pos\{s\}itus [- - -]\\
\lnum{5}Fili[ppo et Sallia cons(ulibus ?)]}
\nf{Témoins}{Export EDCS de 2020 : \emph{Domino [3 bene me]/renti in[3 qui] / vixit ann[os 3] / dpos\{s\}itus [3] / Fili[ppo et Sallia cons(ulibus ?)]}. Aucune fiche EDB appariée (numéro ICUR absent de la liste publique de l’EDB).}
\nf{Apparat}{⸢e⸣ : I gravé ; \{s\} : lettres superflues}
\nf{Traduction}{Au seigneur (?) […] bien méritant […] qui vécut […] ans ; déposé […] ; sous le consulat de Philippus et de Salia.}
\nf{Langue}{L’éditeur signale une graphie corrigée : I gravé pour e (confusion e/i). Des lettres gravées sont jugées superflues par l’éditeur (entre accolades). 1 abréviation développée par l’éditeur (dont cons(ulibus ?)).}
\nf{Chronologie}{Noms consulaires : Philippus (lu en partie) et Salia (restitué).}
\nf{Bibliographie}{ICUR II 4800.}
\nf{Statut}{Lecture secondaire (EDCS) sans second témoin numérique.}
\end{notice}

\begin{notice}{26}{348}{ICUR n.s. II, 5711 · EDCS-38302409}
\label{not:26}
\nf{Datation}{348 — directe : nom consulaire lu}
\nf{Genre}{Épitaphe ; ensembles A, B, C ; 154 lettres transmises.}
\nf{Références}{ICUR II 5711 = ILCV 4326 = Arctos-2012-225 = AE 2012, 197.}
\nf{Provenance}{Rome (indication de l’EDCS) ; aucun rattachement cimétérial transmis.}
\nf{Support}{Non documenté dans les sources consultées.}
\ntext{\lnum{1}Mir⸢ae⸣ bonitatis atq(ue) sanctitati(s)\\
\lnum{2}Esu⸢v⸣iae Ianuariae coniugi qu(a)e vixit mecum\\
\lnum{3}annos XVIII m(enses) V d(ies) XXV Aradius Melissus maritus\\
\lnum{4}et sibi fecit deposita in pace III Kal(endas) Apr(iles) Filippo et Sallia co(n)ss(ulibus)}
\nf{Témoins}{Export EDCS de 2020 : \emph{Mir bonitatis atq(ue) sanctitati(s) / Esuiae Ianuariae coniugi qu(a)e vixit mecum / annos XVIII m(enses) V d(ies) XXV Aradius Melissus maritus / et sibi fecit deposita in pace III Kal(endas) Apr(iles) Filippo et Sallia co(n)ss(ulibus)}. Aucune fiche EDB appariée (numéro ICUR absent de la liste publique de l’EDB).}
\nf{Apparat}{⸢ae⸣ : I gravé ; ⸢v⸣ : B gravé}
\nf{Traduction}{D’une bonté et d’une sainteté admirables : à Esuvia Ianuaria, son épouse, qui vécut avec moi 18 ans 5 mois 25 jours, Aradius Melissus, son mari, a fait (ce monument), pour lui aussi ; déposée en paix le 3e jour avant les calendes d’avril, sous le consulat de Philippus et de Salia.}
\nf{Langue}{L’éditeur signale 2 graphies corrigées : I gravé pour ae ; B gravé pour v (bêtacisme).}
\nf{Chronologie}{Noms consulaires : Philippus (lu) et Salia (lu). Dates du jour : déposition : III Kalendas Apriles = 30 mars 348 (lu).}
\nf{Commentaire}{La formule in pace est conservée. Vocabulaire religieux : épithète du défunt (non doctrinal) — sanctitatis (lu). Commentaire développé : dossier p.~\pageref{dos:26}.}
\nf{Bibliographie}{ICUR II 5711 ; ILCV 4326 ; Arctos-2012-225 ; AE 2012, 197.}
\nf{Statut}{Lecture secondaire (EDCS) sans second témoin numérique. L’EDCS signale une photographie (non consultée).}
\end{notice}

\begin{notice}{27}{348}{ICUR n.s. III, 6498 · EDCS-34801298 · EDB20721}
\label{not:27}
\nf{Datation}{348 — directe : nom consulaire lu}
\nf{Genre}{Épitaphe ; ensembles A, B, C ; 68 lettres transmises ; 21 \% restituées.}
\nf{Références}{ICUR III 6498 = ILCV 3382 (em).}
\nf{Provenance}{Rattachement selon l’EDB : \emph{Coem. Domitillae pars inferior}. Conservation selon l’EDB : \emph{Coem. Domitillae pars inferior}.}
\nf{Support}{Enduit de mortier de mortier. Inscription tracée dans le mortier frais. Reproductions signalées par l’EDB : photographie PCAS (droits réservés ; non reproduites). Fiche EDB : Carlo Carletti / Antonio Enrico Felle, première version 2003.}
\ntext{\lnum{1}] moritur IIII Idus Iulias\\
\lnum{2}[depositus est - - - Id]us Filippo et Sal(l)ia co(n)ss(ulibus)\\
\lnum{3}\{i\}scri(p)s⸢i⸣(t) Donatus}
\nf{Témoins}{Export EDCS de 2020 : \emph{] moritur IIII Idus Iulias / [depositus est 3 Id]us Filippo et Sal(l)ia co(n)ss(ulibus) / \{i\}scri(p)s(t) Donatus}. EDB (EDB20721 ; écarts à examiner) : \emph{[---] moritur IIII idus iulias / 1[---]us1 Filippo et Sal(l)ia / co(n)s(ulibu)s / iscrise Donatus [texte tronqué dans la liste de résultats]}.}
\nf{Apparat}{⸢i⸣ : E gravé ; \{i\} : lettres superflues}
\nf{Traduction}{[…] meurt le 4e jour avant les ides de juillet ; il a été déposé […] ides, sous le consulat de Philippus et de Salia ; Donatus a écrit (l’inscription).}
\nf{Langue}{L’éditeur signale une graphie corrigée : E gravé pour i (confusion e/i). Des lettres gravées sont jugées superflues par l’éditeur (entre accolades). 5 abréviations développées par l’éditeur (dont Sal(l)ia, co(n)ss, (ulibus), scri(p)s, (t)).}
\nf{Chronologie}{Noms consulaires : Philippus (lu) et Salia (lu). Dates du jour : décès : IIII Idus Iulias = 12 juillet 348 (lu).}
\nf{Bibliographie}{ICUR III 6498 ; ILCV 3382 (em).}
\nf{Statut}{Lecture secondaire (EDCS) ; écarts avec l’EDB à arbitrer sur l’ICVR. L’EDCS signale une photographie (non consultée).}
\end{notice}

\begin{notice}{28}{348}{ICUR n.s. IV, 11755 · EDCS-37400849 · EDB966}
\label{not:28}
\nf{Datation}{348 — directe : nom consulaire lu}
\nf{Genre}{Épitaphe ; ensembles A, B, C ; 104 lettres transmises ; 36 \% restituées.}
\nf{Références}{ICUR IV 11755.}
\nf{Provenance}{Rattachement selon l’EDB : \emph{Coem. anonymum ad viam Appiam}. Conservation selon l’EDB : \emph{Coem. anonymum ad viam Appiam}.}
\nf{Support}{Plaque de marbre. Inscription gravée. Fiche EDB : Maria Cristina Grisanzio, première version 2003.}
\ntext{\lnum{1}con]iugi bene mer[enti - - -]\\
\lnum{2}[qui(?) vixit anno]s septem mens⸢e⸣s octo (h)o[ras - - -]\\
\lnum{3}[- - - qui(?)] de(ces)sit in pace die III Kal(endas) F[ebruarias - - -]\\
\lnum{4}[Filippo] et Salli⸢a⸣ conss(ulibus)}
\nf{Témoins}{Export EDCS de 2020 : \emph{con]iugi bene mer[enti 3] / [qui(?) vixit anno]s septem menss octo (h)o[ras 3] / [3 qui(?)] de(ces)sit in pace die III Kal(endas) F[ebruarias 3] / [Filippo] et Salli conss(ulibus)}. EDB (EDB966 ; écarts à examiner) : \emph{[--- con]iugi benemer[enti] / [cum q(---) vixit anni]s ṣẹp̣ṭẹm mensis octo o[ras] / --- quae] de(ces)sit in pace die III kaḷ(endas) [---] / [Philippo] et Sallie cons(ulibu)s [texte tronqué dans la liste de résultats]}.}
\nf{Apparat}{⸢e⸣ : I gravé ; ⸢a⸣ : E gravé}
\nf{Traduction}{À son conjoint bien méritant […] qui vécut 7 ans (?) 8 mois […] heures […] ; qui mourut en paix le 3e jour avant les calendes de février […] ; sous le consulat de Philippus et de Salia.}
\nf{Langue}{L’éditeur signale 2 graphies corrigées : I gravé pour e (confusion e/i) ; E gravé pour a (e pour a).}
\nf{Chronologie}{Noms consulaires : Philippus (restitué) et Salia (lu). Dates du jour : décès : die III Kalendas Februarias = 30 janvier 348 (partiellement restitué). Âge déclaré : 7 a. 8 m. 3 j. (partiellement restitué).}
\nf{Commentaire}{La formule in pace est conservée.}
\nf{Bibliographie}{ICUR IV 11755.}
\nf{Statut}{Lecture secondaire (EDCS) ; écarts avec l’EDB à arbitrer sur l’ICVR.}
\end{notice}

\begin{notice}{29}{348}{ICUR n.s. IV, 11756 · EDCS-37400850 · EDB6728}
\label{not:29}
\nf{Datation}{348 — directe : nom consulaire lu}
\nf{Genre}{Épitaphe ; ensembles A, B, C ; 95 lettres transmises ; 38 \% restituées.}
\nf{Références}{ICUR IV 11756.}
\nf{Provenance}{Rattachement selon l’EDB : \emph{Coem. anonymum ad viam Appiam}. Conservation selon l’EDB : \emph{Coem. anonymum ad viam Appiam}.}
\nf{Support}{Plaque de marbre. Inscription gravée. Fiche EDB : Maria Cristina Grisanzio, première version 2003.}
\ntext{\lnum{1}] Rufino marito suo bene m[erenti]\\
\lnum{2}[qui vixit] ann(os) n(umero) LXXV et m(enses) n(umero) III d[ep(ositus) - - -]\\
\lnum{3}[Kalenda]s Maias Filippo e(t) Sallia\{e\} con[sulibus]}
\nf{Témoins}{EDB (EDB6728) : texte concordant aux conventions de notation près (indice de similarité 0,924).}
\nf{Apparat}{\{e\} : lettres superflues}
\nf{Traduction}{[…] à Rufinus, son mari bien méritant, qui vécut 75 ans et 3 mois ; déposé […] calendes de mai, sous le consulat de Philippus et de Salia.}
\nf{Langue}{Des lettres gravées sont jugées superflues par l’éditeur (entre accolades). 6 abréviations développées par l’éditeur (dont ann(os), n(umero), m(enses), dep(ositus), e(t)).}
\nf{Chronologie}{Noms consulaires : Philippus (lu) et Salia (lu). Ne sont pas des consuls : Rufino marito suo (époux défunt (dédicataire)). Dates du jour : déposition : Kalendas Maias = 1 mai 348 (partiellement restitué). Âge déclaré : 75 a. 3 m. (partiellement restitué).}
\nf{Bibliographie}{ICUR IV 11756.}
\nf{Statut}{Lecture secondaire (EDCS) concordante avec l’EDB : lecture stabilisée entre deux bases, non collationnée sur l’ICVR ni sur le monument.}
\end{notice}

\annee{349}

\begin{notice}{30}{349}{ICUR n.s. III, 8418 · EDCS-35802006 · EDB25693}
\label{not:30}
\nf{Datation}{349 — directe : lecture partielle d’un nom consulaire}
\nf{Genre}{Épitaphe ; ensembles A, B, C ; 43 lettres transmises ; 63 \% restituées.}
\nf{Références}{ICUR III 8418.}
\nf{Provenance}{Rattachement selon l’EDB : \emph{Incertum locum in coem. Domitillae}. Conservation selon l’EDB : \emph{coem. Domitillae, aula supra coemeterium}.}
\nf{Support}{Plaque de marbre. Inscription gravée. Reproductions signalées par l’EDB : photographie PCAS (droits réservés ; non reproduites). Fiche EDB : Carlo Carletti / Antonio Enrico Felle, première version 2003.}
\ntext{\lnum{1}Fili[te - - -]\\
\lnum{2}et d(ies) XX [- - -]\\
\lnum{3}Sept(embres) Lim[enio et Catulino conss(ulibus)}
\nf{Témoins}{EDB (EDB25693) : texte concordant aux conventions de notation près (indice de similarité 0,979).}
\nf{Traduction}{[…] et 20 jours […] ; […] septembre, sous le consulat de Limenius et de Catullinus.}
\nf{Langue}{3 abréviations développées par l’éditeur (dont d(ies), Sept(embres), conss(ulibus)).}
\nf{Chronologie}{Noms consulaires : Limenius (lu en partie) et Catullinus (restitué).}
\nf{Bibliographie}{ICUR III 8418.}
\nf{Statut}{Lecture secondaire (EDCS) concordante avec l’EDB : lecture stabilisée entre deux bases, non collationnée sur l’ICVR ni sur le monument. Plus de la moitié des lettres sont restituées. L’EDCS signale une photographie (non consultée).}
\end{notice}

\begin{notice}{31}{349}{ICUR n.s. III, 8419 · EDCS-35802007 · EDB25694}
\label{not:31}
\nf{Datation}{349 — directe : lecture partielle d’un nom consulaire}
\nf{Genre}{Épitaphe ; ensembles A, B, C ; 42 lettres transmises ; 52 \% restituées.}
\nf{Références}{ICUR III 8419.}
\nf{Provenance}{Rattachement selon l’EDB : \emph{Incertum locum in coem. Domitillae}.}
\nf{Support}{Plaque de marbre. Inscription gravée. Fiche EDB : Carlo Carletti / Antonio Enrico Felle, première version 2003.}
\ntext{\lnum{1}]m[- - -]\\
\lnum{2}[- - -] deces(s)it [- - -]\\
\lnum{3}[- - - biso]mum et [- - -]\\
\lnum{4}[- - - cons(ulatu) Limeni et C]atulini}
\nf{Témoins}{EDB (EDB25694) : texte concordant aux conventions de notation près (indice de similarité 0,948).}
\nf{Traduction}{[…] mourut […] ; […] tombe double et […] ; […] sous le consulat de Limenius et de Catullinus.}
\nf{Langue}{2 abréviations développées par l’éditeur (dont deces(s)it, cons(ulatu)).}
\nf{Chronologie}{Noms consulaires : Limenius (restitué) et Catullinus (lu en partie).}
\nf{Bibliographie}{ICUR III 8419.}
\nf{Statut}{Lecture secondaire (EDCS) concordante avec l’EDB : lecture stabilisée entre deux bases, non collationnée sur l’ICVR ni sur le monument. Plus de la moitié des lettres sont restituées.}
\end{notice}

\begin{notice}{32}{349}{ICUR n.s. IV, 11090 · EDCS-37200372 · EDB39549}
\label{not:32}
\nf{Datation}{349 — directe : lecture partielle d’un nom consulaire}
\nf{Genre}{Épitaphe ; ensembles A, B, C ; 80 lettres transmises ; 60 \% restituées.}
\nf{Références}{ICUR IV 11090.}
\nf{Provenance}{Rattachement selon l’EDB : \emph{Coem. Callisti area subdialis et incerti loci}. Conservation selon l’EDB : \emph{Coem. Callisti area subdialis et incerti loci}.}
\nf{Support}{Plaque de marbre. 36 × 38 × 2 cm (h. × l. × ép.). Inscription gravée. Fiche EDB : Marida Pierno, première version 2015.}
\ntext{\lnum{1}an]nos [- - -] m[ens(es) - - -]\\
\lnum{2}[- - - cum] marit⸢o⸣ fe[cit - - -] dies [- - -]\\
\lnum{3}[- - - Li]menio et Ac[one Cat]ulino [vv(iris) cc(larissimis) conss(ulibus)}
\nf{Témoins}{Export EDCS de 2020 : \emph{an]nos [3] m[ens(es) 3] / [3 cum] marit fe[cit 3] dies [3] / [3 Li]menio et Ac[one Cat]ulino [vv(iris) cc(larissimis) conss(ulibus)}. EDB (EDB39549 ; écarts à examiner) : \emph{------ / [--- vixit] / [an]nos [---] m[---] / [cum] maritu fe[cit ---] dies [---] / [Ulpio Li]menio et Ac[onio Cat]ulino [cons(ulibu)s]}.}
\nf{Apparat}{⸢o⸣ : V gravé}
\nf{Traduction}{[…] ans […] mois ; […] a vécu avec son mari […] jours […] ; sous le consulat de Limenius et d’Aco Catullinus, clarissimes.}
\nf{Langue}{L’éditeur signale une graphie corrigée : V gravé pour o (confusion o/u).}
\nf{Chronologie}{Noms consulaires : Limenius (lu en partie) et Catullinus (lu en partie).}
\nf{Bibliographie}{ICUR IV 11090. Bibliographie complémentaire selon l’EDB : CLRE, 635.}
\nf{Statut}{Lecture secondaire (EDCS) ; écarts avec l’EDB à arbitrer sur l’ICVR. Plus de la moitié des lettres sont restituées.}
\end{notice}

\begin{notice}{33}{349}{ICUR n.s. IV, 12524 · EDCS-37801588 · EDB42440}
\label{not:33}
\nf{Datation}{349 — directe : nom consulaire lu}
\nf{Genre}{Épitaphe ; ensembles A, B, C ; 119 lettres transmises.}
\nf{Références}{ICUR IV 12524 = ILCV 2795b.}
\nf{Provenance}{Rattachement selon l’EDB : \emph{inter vias Appiam et Ardeatinam}.}
\nf{Support}{La fiche EDB ne figure pas dans l’état 2019 des données EpiDoc ; support non documenté.}
\ntext{\lnum{1}Limenio et Catulino conss(ulibus) III\\
\lnum{2}Idus Ianuarias defunctus est\\
\lnum{3}Eu\{u\}odius qui vixit annos LXV\\
\lnum{4}menses tres et dies XI bene me\\
\lnum{5}renti in pace fecit coniux}
\nf{Témoins}{EDB (EDB42440) : texte concordant aux conventions de notation près (indice de similarité 0,997).}
\nf{Apparat}{\{u\} : lettres superflues}
\nf{Traduction}{Sous le consulat de Limenius et de Catullinus, le 3e jour avant les ides de janvier, mourut Euodius, qui vécut 65 ans 3 mois et 11 jours ; à lui, bien méritant, en paix, son épouse a fait (ce monument).}
\nf{Langue}{Des lettres gravées sont jugées superflues par l’éditeur (entre accolades). 1 abréviation développée par l’éditeur (dont conss(ulibus)).}
\nf{Chronologie}{Noms consulaires : Limenius (lu) et Catullinus (lu). Ne sont pas des consuls : Eu\{u\}odius qui vixit annos LXV (défunt). Dates du jour : événement indéterminé : Idus Ianuarias = 13 janvier 349 (lu). Âge déclaré : 65 a. 3 m. 11 j. (lu).}
\nf{Commentaire}{La formule in pace est conservée.}
\nf{Bibliographie}{ICUR IV 12524 ; ILCV 2795b. Bibliographie complémentaire selon l’EDB : ILCV, 2795 B.}
\nf{Statut}{Lecture secondaire (EDCS) concordante avec l’EDB : lecture stabilisée entre deux bases, non collationnée sur l’ICVR ni sur le monument.}
\end{notice}

\annee{350}

\begin{notice}{34}{350}{ICUR n.s. I, 2596 · EDCS-03900021 · EDB32543}
\label{not:34}
\nf{Datation}{350 — directe : nom consulaire lu}
\nf{Genre}{Épitaphe ; ensembles A, B, C ; 131 lettres transmises.}
\nf{Références}{Anagni 41 = ICUR I 2596 = ILCV 2940a.}
\nf{Provenance}{Rattachement selon l’EDB : \emph{Incertae originis} (le volume I des ICVR réunit les inscriptions dont la provenance est perdue). Conservation selon l’EDB : \emph{Anagni (FR)   ; monastero delle suore Cistercensi della Carità}.}
\nf{Support}{Plaque de marbre. 28 × 77 cm (h. × l.). Inscription gravée. Signes décrits par l’EDB : chrisme. Reproductions signalées par l’EDB : dessin, ICVR I (de Rossi), reproduction, Perret 1851 et reproduction, Solin et alii 1996 (droits réservés ; non reproduites). Fiche EDB : Filippo Piazzolla, première version 2015.}
\ntext{\lnum{1}Gau\\
\lnum{2}de\\
\lnum{3}n\\
\lnum{4}ti\\
\lnum{5}us\\
\lnum{6}Felicissima in pace uxor qu(a)e\\
\lnum{7}vi\{c\}xit annis ⸢v⸣iginti no⸢v⸣e(m)\\
\lnum{8}me(n)s⸢e⸣s no⸢v⸣e(m) dies no⸢v⸣e(m) deposita\\
\lnum{9}VII Idus Iulias pos(t) consulatu(m)\\
\lnum{10}Limeni et Catulini Gauden\\
\lnum{11}tius}
\nf{Témoins}{Export EDCS de 2020 : \emph{Gau/de/n/ti/us // Felicissima in pace uxor qu(a)e / vi\{c\}xit annis iginti noe(m) / me(n)ss noe(m) dies noe(m) deposita / VII Idus Iulias pos(t) consulatu(m) / Limeni et Catulini Gauden/tius}. EDB (EDB32543 ; écarts à examiner) : \emph{⟨:a⟩ /((monogramma Christi)) / Gau/de/n/t⌜i⌝/us /⟨:b⟩ /Felicissima in pace uxor que / vicxit annis biginti nobe / mesis nobe dies no⌜b⌝e deposita / VII idus iulias pos⟨t⟩ consulatu / Limeni et Catulini Gauden/tius [texte tronqué dans la liste de résultats]}.}
\nf{Apparat}{⸢v⸣ : B gravé ; ⸢v⸣ : B gravé ; ⸢e⸣ : I gravé ; ⸢v⸣ : B gravé ; ⸢v⸣ : R gravé ; \{c\} : lettres superflues}
\nf{Traduction}{Gaudentius // Felicissima, en paix, épouse, qui vécut 29 ans 9 mois 9 jours ; déposée le 7e jour avant les ides de juillet, après le consulat de Limenius et de Catullinus ; Gaudentius.}
\nf{Langue}{L’éditeur signale 5 graphies corrigées : B gravé pour v (bêtacisme) ×3 ; I gravé pour e (confusion e/i) ; R gravé pour v. Des lettres gravées sont jugées superflues par l’éditeur (entre accolades). 7 abréviations développées par l’éditeur (dont qu(a)e, e(m), me(n)s, pos(t)).}
\nf{Chronologie}{Noms consulaires : Limenius (lu) et Catullinus (lu). Datation par post-consulat. Dates du jour : déposition : VII Idus Iulias = 9 juillet 350 (lu). Âge déclaré : 29 a. 9 m. 9 j. (lu).}
\nf{Commentaire}{La formule in pace est conservée.}
\nf{Bibliographie}{Anagni 41 ; ICUR I 2596 ; ILCV 2940a. Bibliographie complémentaire selon l’EDB : ILCV, 2940A ; Solin et alii 1996, 47-48 n. 41 con foto.}
\nf{Statut}{Lecture secondaire (EDCS) ; écarts avec l’EDB à arbitrer sur l’ICVR. L’EDCS signale une photographie (non consultée).}
\end{notice}

\begin{notice}{35}{350}{ICUR n.s. I, 1422 · EDCS-39000409 · EDB31186}
\label{not:35}
\nf{Datation}{350 — directe : nom consulaire lu}
\nf{Genre}{Épitaphe ; ensembles A, B, C ; 83 lettres transmises.}
\nf{Références}{ICUR I 1422 = ILCV 2628 (em).}
\nf{Provenance}{Rattachement selon l’EDB : \emph{Incertae originis} (le volume I des ICVR réunit les inscriptions dont la provenance est perdue). Conservation selon l’EDB : \emph{Musei Vaticani: Lapidario Cristiano}.}
\nf{Support}{Plaque de marbre. 29 × 74 cm (h. × l.). Inscription gravée. Reproductions signalées par l’EDB : dessin, ICVR I (de Rossi) (droits réservés ; non reproduites). Fiche EDB : Filippo Piazzolla, première version 2015.}
\ntext{\lnum{1}Successae filiae dul\\
\lnum{2}cissimae parentes q(uae) v(ixit) an\\
\lnum{3}nos duos m(enses) XI d(ies) II b(ene) m(erenti) in pace\\
\lnum{4}d(e)p(osita) XVIII Kal(endas) Sept(embres) Ni⸢g⸣riniano}
\nf{Témoins}{Export EDCS de 2020 : \emph{Successae filiae dul/cissimae parentes q(uae) v(ixit) an/nos duos m(enses) XI d(ies) II b(ene) m(erenti) in pace / d(e)p(osita) XVIII Kal(endas) Sept(embres) Niriniano}. EDB (EDB31186) : texte concordant aux conventions de notation près (indice de similarité 0,945).}
\nf{Apparat}{⸢g⸣ : C gravé}
\nf{Traduction}{À Successa, leur fille très douce, ses parents ; elle vécut 2 ans 11 mois 2 jours ; bien méritante, en paix ; déposée le 18e jour avant les calendes de septembre, sous (le consulat de) Nigrinianus.}
\nf{Langue}{L’éditeur signale une graphie corrigée : C gravé pour g (c pour g).}
\nf{Chronologie}{Noms consulaires : Nigrinianus (lu). Dates du jour : déposition : XVIII Kalendas Septembres = 15 août 350 (lu). Âge déclaré : 2 a. 11 m. 2 j. (lu).}
\nf{Commentaire}{La formule in pace est conservée.}
\nf{Bibliographie}{ICUR I 1422 ; ILCV 2628 (em). Bibliographie complémentaire selon l’EDB : Lega 2000, 278 ; ILCV, 2628.}
\nf{Statut}{Lecture secondaire (EDCS) concordante avec l’EDB : lecture stabilisée entre deux bases, non collationnée sur l’ICVR ni sur le monument. L’EDCS signale une photographie (non consultée).}
\end{notice}

\annee{351}

\begin{notice}{36}{351}{ICUR n.s. I, 19 · EDCS-18100517 · EDB27137}
\label{not:36}
\nf{Datation}{351 — indirecte}
\nf{Genre}{Épitaphe ; ensembles A, B, C ; 63 lettres transmises.}
\nf{Références}{CIL VI 1712 (p 3173, 4741) = ILCV 1850 = ICUR I 19 = AE 2006, 93.}
\nf{Provenance}{Rattachement selon l’EDB : \emph{Incertae originis} (le volume I des ICVR réunit les inscriptions dont la provenance est perdue).}
\nf{Support}{Plaque de marbre. Inscription gravée. Fiche EDB : Anita Rocco, première version 2013.}
\ntext{\lnum{1}Clo\\
\lnum{2}di\\
\lnum{3}us\\
\lnum{4}A\\
\lnum{5}del\\
\lnum{6}fi\\
\lnum{7}us\\
\lnum{8}v(ir) c(larissimus)\\
\lnum{9}ex\\
\lnum{10}prae\\
\lnum{11}fe\\
\lnum{12}ctis\\
\lnum{13}ur\\
\lnum{14}bis\\
\lnum{15}u\\
\lnum{16}xo\\
\lnum{17}ri\\
\lnum{18}in\\
\lnum{19}co⸢m⸣\\
\lnum{20}pa\\
\lnum{21}ra\\
\lnum{22}bi\\
\lnum{23}li\\
\lnum{24}et\\
\lnum{25}si\\
\lnum{26}bi\\
\lnum{27}fe\\
\lnum{28}cit}
\nf{Témoins}{Export EDCS de 2020 : \emph{Clo/di/us / A/del/fi/us / v(ir) c(larissimus) / ex / prae/fe/ctis / ur/bis / u/xo/ri / in/co/pa/ra/bi/li / et / si/bi / fe/cit}. EDB (EDB27137) : texte concordant aux conventions de notation près (indice de similarité 0,942).}
\nf{Apparat}{⸢m⸣ : N gravé}
\nf{Traduction}{Clodius Adelfius, clarissime, ancien préfet de la Ville, à son épouse incomparable et pour lui-même a fait (ce monument).}
\nf{Langue}{L’éditeur signale une graphie corrigée : N gravé pour m (n pour m).}
\nf{Chronologie}{Aucun élément consulaire n'est gravé. L'année 351 est celle de la base ; l'EDB donne « 351 ; 354 ». Elle repose sur l'identification de Clodius Adelfius avec le préfet de la Ville de 351, non sur le texte : 351 est un terme postérieur (après la préfecture), non l'année du décès. La notice est exclue des séries chronologiques fines.}
\nf{Commentaire}{Rang sénatorial d’une personne commémorée (clodius adelfius vir clarissimus ex p). Fonction ou métier mentionné : praefectis. Commentaire développé : dossier p.~\pageref{dos:36}.}
\nf{Bibliographie}{CIL VI 1712 (p 3173, 4741) ; ILCV 1850 ; ICUR I 19 ; AE 2006, 93. Bibliographie complémentaire selon l’EDB : CIL VI, 1712 ; ILCV, 1850.}
\nf{Statut}{Lecture secondaire (EDCS) concordante avec l’EDB : lecture stabilisée entre deux bases, non collationnée sur l’ICVR ni sur le monument. L’EDCS signale une photographie (non consultée).}
\end{notice}

\annee{352}

\begin{notice}{37}{352}{ICUR n.s. I, 2081 · EDCS-39300305 · EDB32043}
\label{not:37}
\nf{Datation}{352 — directe : nom consulaire lu}
\nf{Genre}{Épitaphe ; ensembles A, B, C ; 62 lettres transmises ; 31 \% restituées.}
\nf{Références}{ICUR I 2081.}
\nf{Provenance}{Rattachement selon l’EDB : \emph{Incertae originis} (le volume I des ICVR réunit les inscriptions dont la provenance est perdue). Conservation selon l’EDB : \emph{Musei Vaticani: Galleria Lapidaria}.}
\nf{Support}{Plaque de marbre. 38 × 39 cm (h. × l.). Inscription gravée. Reproductions signalées par l’EDB : dessin, ICVR I (de Rossi) (droits réservés ; non reproduites). Fiche EDB : Filippo Piazzolla, première version 2015.}
\ntext{\lnum{1}[Locus] Anicitionis e⸢t⸣ [- - -]\\
\lnum{2}[qui vi]xit ann(os) XXVI[- - -]\\
\lnum{3}[dep(ositus) - - -] Aug(ustas) De⸢c⸣en(tio) et Paulo conss(ulibus)}
\nf{Témoins}{Export EDCS de 2020 : \emph{[Locus] Anicitionis e [3] / [qui vi]xit ann(os) XXVI[3] / [dep(ositus) 3] Aug(ustas) Deen(tio) et Paulo conss(ulibus)}. EDB (EDB32043 ; écarts à examiner) : \emph{[locus] Anicitionis e⌜t⌝ [---] / [qui vi]xit ann(is) XXVI [---] / [ dep(osit---) ] aug(ustas) De⌜c⌝en⟨tio⟩ et Paul̲o̲ c̲o̲n̲s̲(ulibu)s̲ [---] [texte tronqué dans la liste de résultats]}.}
\nf{Apparat}{⸢t⸣ : I gravé ; ⸢c⸣ : I gravé}
\nf{Traduction}{Emplacement d’Anicitio (?) et […] qui vécut 26 ans […] ; déposé […] août, sous le consulat de Décence et de Paulus.}
\nf{Langue}{L’éditeur signale 2 graphies corrigées : I gravé pour t ; I gravé pour c.}
\nf{Chronologie}{Noms consulaires : Decentius (lu) et Paulus (lu).}
\nf{Bibliographie}{ICUR I 2081. Bibliographie complémentaire selon l’EDB : Lega 2000.}
\nf{Statut}{Lecture secondaire (EDCS) ; écarts avec l’EDB à arbitrer sur l’ICVR. L’EDCS signale une photographie (non consultée).}
\end{notice}

\begin{notice}{38}{352}{ICUR n.s. II, 4241 · EDCS-33900148 · EDB17657}
\label{not:38}
\nf{Datation}{352 — directe : nom consulaire lu}
\nf{Genre}{Épitaphe ; ensembles A, B, C ; 103 lettres transmises.}
\nf{Références}{ICUR II 4241 = ILCV 3252a (em).}
\nf{Provenance}{Rattachement selon l’EDB : \emph{Coem. sub divo apud basilicam S.Petri apostoli}. Conservation selon l’EDB : \emph{Musei Vaticani: Lapidario Cristiano}.}
\nf{Support}{Plaque de marbre. Inscription gravée. Signes décrits par l’EDB : chrisme. Reproductions signalées par l’EDB : dessin, ICVR I (de Rossi) (droits réservés ; non reproduites). Fiche EDB : Antonio Enrico Felle, première version 2003.}
\ntext{\lnum{1}Privata qu(a)e vixit annis LIII me(n)ses(!)\\
\lnum{2}VIII die(bu)s II deposita est IIII Nonas\\
\lnum{3}Febr(uarias) bene merenti in pace Chr(isti)\\
\lnum{4}Decentio et Paulo conss(ulibus)}
\nf{Témoins}{EDB (EDB17657) : texte concordant aux conventions de notation près (indice de similarité 0,972).}
\nf{Apparat}{(!) : forme signalée par l’éditeur}
\nf{Traduction}{Privata, qui vécut 53 ans 8 mois 2 jours ; déposée le 4e jour avant les nones de février ; bien méritante, en paix ☧ ; sous le consulat de Décence et Paulus.}
\nf{Langue}{Cinq autres développements : qu(a)e, me(n)ses, die(bu)s, Febr(uarias), conss(ulibus). Chr(isti) est conservé comme transmission éditoriale ; son statut de mot abrégé n’est pas tenu pour établi.}
\nf{Chronologie}{Noms consulaires : Decentius (lu) et Paulus (lu). Dates du jour : déposition : IIII Nonas Februarias = 2 février 352 (lu). Âge déclaré : 53 a. 8 m. 2 j. (lu).}
\nf{Commentaire}{Après \emph{in pace}, l'EDCS transmet \emph{Chr(isti)} ; l'EDB (EDB17657, notice d'A. E. Felle) transcrit \emph{in pace ((monogramma Christi))} et décrit un chrisme dans les données matérielles. L'estampage universitaire ILCV 3252A (Ohio State University) montre lui aussi un christogramme placé après IN PACE. Deux témoins indépendants de l'EDCS concordent donc : la pierre porte un monogramme, non le mot \emph{Christi} écrit en lettres. Le texte est classé parmi les signes (catégorie S) et non parmi les mentions lexicales du Christ. La question de savoir si le monogramme fonctionne ici comme abréviation (\emph{in pace Chr(isti)}) ou comme signe juxtaposé ne peut être tranchée sans photographie : elle est discutée avec les n° 140, 151 et 311 (chapitre sur Dieu et le Christ). Commentaire développé : dossier p.~\pageref{dos:38}.}
\nf{Bibliographie}{ICUR II 4241 ; ILCV 3252a (em).}
\nf{Statut}{Lecture secondaire (EDCS) concordante avec l’EDB : lecture stabilisée entre deux bases, non collationnée sur l’ICVR ni sur le monument. L’EDCS signale une photographie (non consultée).}
\end{notice}

\begin{notice}{39}{352}{ICUR n.s. III, 8140 · EDCS-35100216 · EDB24663}
\label{not:39}
\nf{Datation}{352 — directe : nom consulaire lu}
\nf{Genre}{Épitaphe ; ensembles A, B, C ; 34 lettres transmises ; 47 \% restituées.}
\nf{Références}{ICUR III 8140.}
\nf{Provenance}{Rattachement selon l’EDB : \emph{Basilica ss. Nerei et Achillei in coem. Domitillae}.}
\nf{Support}{Plaque de marbre. Inscription gravée. Fiche EDB : Carlo Carletti / Antonio Enrico Felle, première version 2003.}
\ntext{\lnum{1}Decentio] et Paulo [conss(ulibus)]\\
\lnum{2}[- - - A]fro[disi}
\nf{Témoins}{EDB (EDB24663 ; écarts à examiner) : \emph{[--- Decentio] et Paulo [cons(ulibu)s ---] / [--- A]fro[disi---] [texte tronqué dans la liste de résultats]}.}
\nf{Traduction}{Sous le consulat de Décence et de Paulus […] Afrodisius (?).}
\nf{Langue}{1 abréviation développée par l’éditeur (dont conss(ulibus)).}
\nf{Chronologie}{Noms consulaires : Decentius (restitué) et Paulus (lu).}
\nf{Bibliographie}{ICUR III 8140.}
\nf{Statut}{Lecture secondaire (EDCS) ; écarts avec l’EDB à arbitrer sur l’ICVR. L’EDCS signale une photographie (non consultée).}
\end{notice}

\annee{353}

\begin{notice}{40}{353}{ICUR n.s. I, 1252 · EDCS-39101962 · EDB30464}
\label{not:40}
\nf{Datation}{353 — directe : nom consulaire lu ; datation contestée}
\nf{Genre}{Épitaphe ; ensembles A, B, C ; 111 lettres transmises ; 23 \% restituées.}
\nf{Références}{ICUR I 1252.}
\nf{Provenance}{Rattachement selon l’EDB : \emph{Incertae originis} (le volume I des ICVR réunit les inscriptions dont la provenance est perdue). Conservation selon l’EDB : \emph{Musei Vaticani: Lapidario Cristiano}.}
\nf{Support}{Plaque de marbre. 42 × 52 cm (h. × l.). Inscription gravée. Reproductions signalées par l’EDB : dessin, ICVR I (de Rossi) (droits réservés ; non reproduites). Fiche EDB : Filippo Piazzolla, première version 2015.}
\ntext{\lnum{1}bene me]renti qui ⸢v⸣ix⸢i⸣t cum co\\
\lnum{2}[niuge sua ann(os) - - -]XXI me(n)s(es) III depositus\\
\lnum{3}[annorum - - -] me(n)s(es) III d(ies) VIII depositus in\\
\lnum{4}[pace - - -] M⸢ag⸣n\{n\}entio et Decen\\
\lnum{5}[ti]o cc(on)ss(ulibus)}
\nf{Témoins}{Export EDCS de 2020 : \emph{bene me]renti qui ixt cum co/[niuge sua ann(os) 3]XXI me(n)s(es) III depositus / [annorum 3] me(n)s(es) III d(ies) VIII depositus in / [pace 3] Mn\{n\}entio et Decen/[ti]o cc(on)ss(ulibus)}. EDB (EDB30464 ; écarts à examiner) : \emph{[--- beneme]erenti qui bixet cum co/[iuge sua ann(os) ---]XXI me(nse)s III depositus / [--- annorum ---] me(n)s(ium) III d(ierum) VIII depositus in [pace ---] MVANNENTIO et Decen/[ti]o cc(on)s(ulibu)s / ((avis)) [texte tronqué dans la liste de résultats]}.}
\nf{Apparat}{⸢v⸣ : B gravé ; ⸢i⸣ : E gravé ; ⸢ag⸣ : VA gravé ; \{n\} : lettres superflues}
\nf{Traduction}{[…] bien méritant, qui vécut avec son épouse […]21 ans 3 mois ; déposé ; âgé de […] ans 3 mois 8 jours ; déposé en paix […] sous le consulat de Magnence et de Décence.}
\nf{Langue}{L’éditeur signale 3 graphies corrigées : B gravé pour v (bêtacisme) ; E gravé pour i (confusion e/i) ; VA gravé pour ag. Des lettres gravées sont jugées superflues par l’éditeur (entre accolades). 8 abréviations développées par l’éditeur (dont ann(os), me(n)s, (es), d(ies)).}
\nf{Chronologie}{Noms consulaires : Magnentius (lu) et Decentius (lu en partie). La formule \emph{M⸢a⸣nnentio et Decen[ti]o cc(on)s(ulibus)} nomme Magnence et Décence. La lecture retenue y voit le collège de 353 (Magnence II et Décence II, sans itération gravée). L'EDB, citant la CLRE (p. 636), date le texte de 352 en supposant que le lapicide a combiné la formule de 352 (\emph{Magno Decentio}) avec le nom de Magnence gravé l'année précédente. Deux interprétations publiées s'opposent ; l'année de classement 353 est conservée et la datation est déclarée contestée (352/353).}
\nf{Commentaire}{La formule in pace est attestée seulement par restitution (répartie sur deux lignes).}
\nf{Bibliographie}{ICUR I 1252. Bibliographie complémentaire selon l’EDB : ILCV, 2967 adn ; Ferrua 1979b, 30, n. 1252 ; Molisani 1973, 110 ; CLRE, 636.}
\nf{Statut}{Lecture secondaire (EDCS) ; écarts avec l’EDB à arbitrer sur l’ICVR. L’EDCS signale une photographie (non consultée).}
\end{notice}

\annee{355}

\begin{notice}{41}{355}{ICUR n.s. I, 2769 · EDCS-39900318 · EDB32698}
\label{not:41}
\nf{Datation}{355 — directe : nom consulaire lu}
\nf{Genre}{Épitaphe ; ensembles A, B, C ; 51 lettres transmises.}
\nf{Références}{ICUR I 2769 = ILCV 2978a.}
\nf{Provenance}{Rattachement selon l’EDB : \emph{Incertae originis} (le volume I des ICVR réunit les inscriptions dont la provenance est perdue). Conservation selon l’EDB : \emph{Fossombrone (PU)   ; Lapidario Diocesano}.}
\nf{Support}{Plaque de marbre. 21,5 × 15 cm (h. × l.). Inscription gravée. Reproductions signalées par l’EDB : dessin, ICVR I (de Rossi) (droits réservés ; non reproduites). Fiche EDB : Filippo Piazzolla, première version 2018.}
\ntext{\lnum{1}Marasi\\
\lnum{2}us d(epositus) V Id(us)\\
\lnum{3}Oct(obres) in p(ace)\\
\lnum{4}qui vixit\\
\lnum{5}an(nos) XXXVII\\
\lnum{6}Arb(itione) et Lol(liano)\\
\lnum{7}Marasius d(epositus)}
\nf{Témoins}{EDB (EDB32698 ; écarts à examiner) : \emph{⟨:a⟩ /Marasi/us d(epositus) V id(us) / oct(obres) in p(ace) / qui vixit / an(nis) XXXVII / Arb(itio) et Lol(liano) / ⟨:b⟩ /Marasius d(epositus) [texte tronqué dans la liste de résultats]}.}
\nf{Traduction}{Marasius, déposé le 5e jour avant les ides d’octobre, en paix, qui vécut 37 ans ; sous (le consulat d’)Arbitio et Lollianus. // Marasius déposé.}
\nf{Langue}{8 abréviations développées par l’éditeur (dont d(epositus), Id(us), Oct(obres), p(ace), an(nos), Arb(itione)).}
\nf{Chronologie}{Noms consulaires : Arbitio (lu) et Lollianus (lu). Dates du jour : déposition : V Idus Octobres = 11 octobre 355 (lu). Âge déclaré : 37 a. (lu).}
\nf{Commentaire}{La formule in pace est transmise sous forme abrégée, sans restitution de lacune ; son développement reste éditorial.}
\nf{Bibliographie}{ICUR I 2769 ; ILCV 2978a. Bibliographie complémentaire selon l’EDB : ILCV, 2975 adn.}
\nf{Statut}{Lecture secondaire (EDCS) ; écarts avec l’EDB à arbitrer sur l’ICVR. L’EDCS signale une photographie (non consultée).}
\end{notice}

\begin{notice}{42}{355}{ICUR n.s. II, 4268 · EDCS-33900175 · EDB8228}
\label{not:42}
\nf{Datation}{355 — directe : nom consulaire lu}
\nf{Genre}{Épitaphe ; ensembles A, B, C ; 134 lettres transmises.}
\nf{Références}{ICUR II 4268 = ILCV 3099 (em).}
\nf{Provenance}{Rattachement selon l’EDB : \emph{Coem. Octavillae ad s.Pancratium}. Conservation selon l’EDB : \emph{Musei Vaticani: Lapidario Cristiano}.}
\nf{Support}{Plaque de marbre. Inscription gravée. Reproductions signalées par l’EDB : dessin, ICVR I (de Rossi) (droits réservés ; non reproduites). Fiche EDB : Antonio Enrico Felle, première version 2003.}
\ntext{\lnum{1}Statilia Tigris filia dul\\
\lnum{2}cissima qu(a)e vixit ann(os)\\
\lnum{3}XXII d(ies) XVIIII et cu(m) marito\\
\lnum{4}fecit ann(os) IIII m(enses) III d(ies) VII\\
\lnum{5}Kal(endas) Iulias(!) quiescet in pace\\
\lnum{6}Flaviis Arbitione et Lolliano cons(ulibus)}
\nf{Témoins}{EDB (EDB8228) : texte concordant aux conventions de notation près (indice de similarité 1).}
\nf{Apparat}{(!) : forme signalée par l’éditeur}
\nf{Traduction}{Statilia Tigris, fille très douce, qui vécut 22 ans 19 jours, et vécut avec son mari 4 ans 3 mois 7 jours ; aux calendes de juillet elle repose en paix ; sous le consulat des Flavii Arbitio et Lollianus.}
\nf{Langue}{L’éditeur signale au moins une forme anormale (sic). 10 abréviations développées par l’éditeur (dont qu(a)e, ann(os), d(ies), cu(m), m(enses)).}
\nf{Chronologie}{Noms consulaires : Arbitio (lu) et Lollianus (lu). Dates du jour : événement indéterminé : Kalendas Iulias = 1 juillet 355 (lu). Âge déclaré : 22 a. 3 m. 19 j. (lu).}
\nf{Commentaire}{La formule in pace est conservée. Vocabulaire religieux : repos et sommeil de la mort — quiescet (lu).}
\nf{Bibliographie}{ICUR II 4268 ; ILCV 3099 (em).}
\nf{Statut}{Lecture secondaire (EDCS) concordante avec l’EDB : lecture stabilisée entre deux bases, non collationnée sur l’ICVR ni sur le monument. L’EDCS signale une photographie (non consultée).}
\end{notice}

\begin{notice}{43}{355}{ICUR n.s. III, 8420 · EDCS-31900329 · EDB25695}
\label{not:43}
\nf{Datation}{355 — directe : nom consulaire lu}
\nf{Genre}{Épitaphe ; ensembles A, B, C ; 237 lettres transmises ; 53 \% restituées.}
\nf{Références}{ICUR III 8420 = CLE 660 = ILCV 2451 (em).}
\nf{Provenance}{Rattachement selon l’EDB : \emph{Incertum locum in coem. Domitillae}. Conservation selon l’EDB : \emph{coem. Domitillae, aula supra coemeterium}.}
\nf{Support}{Plaque de marbre. Inscription gravée. Texte métrique. Signes décrits par l’EDB : chrisme et chrisme avec alpha et oméga. Reproductions signalées par l’EDB : photographie PCAS ×2 et dessin, ICVR I (de Rossi) (droits réservés ; non reproduites). Fiche EDB : Carlo Carletti / Antonio Enrico Felle, première version 2003.}
\ntext{\lnum{1}In nomine [Christi - - -]\\
\lnum{2}tristis Anastasio Cons[tantia carmina scribit(?)]\\
\lnum{3}coniunx qui lucem t[enebris mutavit amaris(?)]\\
\lnum{4}vita quater denis et q[uinis constitit annis(?)]\\
\lnum{5}quam cito praerept[us dilectae uxoris amori(?)]\\
\lnum{6}fletus duodecimum [Ianus cum sumeret ortum(?)]\\
\lnum{7}conditus Arbitio co[nsul cum duceret annum]\\
\lnum{8}[}
\nf{Témoins}{EDB (EDB25695) : texte concordant aux conventions de notation près (indice de similarité 0,936).}
\nf{Traduction}{(vers) Au nom du Christ […]. Constantia, épouse affligée, écrit ces vers pour Anastasius, qui a échangé la lumière contre d’amères ténèbres ; sa vie s’arrêta à quatre fois dix et cinq ans (= 45 ans) ; combien vite il fut ravi à l’amour de son épouse chérie ! Il fut pleuré quand Janus prenait son douzième lever (= 12 janvier ?) ; il fut enseveli l’année que conduisait le consul Arbitio […].}
\nf{Langue}{5 abréviations développées par l’éditeur (dont scribit(?), amaris(?), annis(?), amori(?), ortum(?)).}
\nf{Chronologie}{Noms consulaires : Arbitio (lu).}
\nf{Commentaire}{Vocabulaire religieux : désignation du Christ — Christi (entièrement restitué). Commentaire développé : dossier p.~\pageref{dos:43}.}
\nf{Bibliographie}{ICUR III 8420 ; CLE 660 ; ILCV 2451 (em). Bibliographie complémentaire selon l’EDB : Mazzoleni 2015, 165 ; CLE I (nn. 1-859), 660 ; ILCV, 2451.}
\nf{Statut}{Lecture secondaire (EDCS) concordante avec l’EDB : lecture stabilisée entre deux bases, non collationnée sur l’ICVR ni sur le monument. Plus de la moitié des lettres sont restituées. L’EDCS signale une photographie (non consultée).}
\end{notice}

\begin{notice}{44}{355}{ICUR n.s. IV, 11091 · EDCS-37200373 · EDB39550}
\label{not:44}
\nf{Datation}{355 — directe : nom consulaire lu}
\nf{Genre}{Épitaphe ; ensembles A, B, C ; 40 lettres transmises ; 57 \% restituées.}
\nf{Références}{ICUR IV 11091.}
\nf{Provenance}{Rattachement selon l’EDB : \emph{Coem. Callisti area subdialis et incerti loci}. Conservation selon l’EDB : \emph{Coem. Callisti area subdialis et incerti loci}.}
\nf{Support}{Plaque de marbre. 26 × 46 × 3,5 cm (h. × l. × ép.). Inscription gravée. Fiche EDB : Marida Pierno, première version 2015.}
\ntext{\lnum{1}De(positio) Cel[eri]ni v(iri) c(larissimi) A[- - -]\\
\lnum{2}[- - -]T Arb⸢i⸣t(ione) e[t Lolliano conss(ulibus)}
\nf{Témoins}{Export EDCS de 2020 : \emph{De(positio) Cel[eri]ni v(iri) c(larissimi) A[3] / [3]T Arbt(ione) e[t Lolliano conss(ulibus)}. EDB (EDB39550 ; écarts à examiner) : \emph{d(e)p(ositio) Cel[eri]ni v(iri) c(larissimi) a[---] /[---]t Arbut(o) e[t Lolliano]}.}
\nf{Apparat}{⸢i⸣ : V gravé}
\nf{Traduction}{Déposition de Celerinus, clarissime, A[…] ; […] sous le consulat d’Arbitio et de Lollianus.}
\nf{Langue}{L’éditeur signale une graphie corrigée : V gravé pour i.}
\nf{Chronologie}{Noms consulaires : Arbitio (lu) et Lollianus (restitué).}
\nf{Commentaire}{Rang sénatorial d’une personne commémorée (depositio celerini viri clarissimi a… …).}
\nf{Bibliographie}{ICUR IV 11091.}
\nf{Statut}{Lecture secondaire (EDCS) ; écarts avec l’EDB à arbitrer sur l’ICVR. Plus de la moitié des lettres sont restituées.}
\end{notice}

\annee{356}

\begin{notice}{45}{356}{ICUR n.s. II, 4802 · EDCS-36101283}
\label{not:45}
\nf{Datation}{356 — directe : nom consulaire lu}
\nf{Genre}{Épitaphe ; ensembles A, B, C ; 85 lettres transmises ; 16 \% restituées.}
\nf{Références}{ICUR II 4802 = ILCV +1735.}
\nf{Provenance}{Rome (indication de l’EDCS) ; aucun rattachement cimétérial transmis.}
\nf{Support}{Non documenté dans les sources consultées.}
\ntext{\lnum{1}]CAT I[- - -] Iustina [- - -]\\
\lnum{2}[- - -]XVIII d[- - -]\\
\lnum{3}[- - - vir]ginium su(u)m in pace\\
\lnum{4}vix\{s\}it annos XXV m(enses) IIII\\
\lnum{5}[- - - consulatu Ar]bitioni et Lolliani cons(ulum)}
\nf{Témoins}{Aucune fiche EDB appariée (numéro ICUR absent de la liste publique de l’EDB).}
\nf{Apparat}{\{s\} : lettres superflues}
\nf{Traduction}{[…] Iustina […] 18 […] ; […] son ‘vierge’ (= époux épousé vierge), en paix ; elle vécut 25 ans 4 mois ; […] sous le consulat d’Arbitio et de Lollianus.}
\nf{Langue}{Des lettres gravées sont jugées superflues par l’éditeur (entre accolades). 3 abréviations développées par l’éditeur (dont su(u)m, m(enses), cons(ulum)).}
\nf{Chronologie}{Noms consulaires : Arbitio (lu en partie) et Lollianus (lu). Âge déclaré : 25 a. 4 m. (partiellement restitué).}
\nf{Commentaire}{La formule in pace est conservée.}
\nf{Bibliographie}{ICUR II 4802 ; ILCV +1735.}
\nf{Statut}{Lecture secondaire (EDCS) sans second témoin numérique.}
\end{notice}

\begin{notice}{46}{356}{ICUR n.s. III, 6499 · EDCS-34801299 · EDB19876}
\label{not:46}
\nf{Datation}{356 — directe : lecture partielle d’un nom consulaire}
\nf{Genre}{Épitaphe ; ensembles A, B, C ; 135 lettres transmises ; 13 \% restituées.}
\nf{Références}{ICUR III 6499 = ILCV 3663 (em).}
\nf{Provenance}{Rattachement selon l’EDB : \emph{Coem. Domitillae pars inferior}. Conservation selon l’EDB : \emph{Coem. Domitillae pars inferior}.}
\nf{Support}{Plaque de marbre. 72 × 60 × 2,5 cm (h. × l. × ép.). Inscription gravée. Signes décrits par l’EDB : chrisme. Reproductions signalées par l’EDB : photographie PCAS ×6 (droits réservés ; non reproduites). Fiche EDB : Carlo Carletti / Antonio Enrico Felle, première version 2016.}
\ntext{\lnum{1}Marcus Karisiae\\
\lnum{2}co⸢m⸣pari carissimae\\
\lnum{3}feci nobis domum\\
\lnum{4}aete[r]nale(m) deposita\\
\lnum{5}pr[id]ie Kal(endas) Martias bene\\
\lnum{6}[me]renti in pace\\
\lnum{7}[Co]nstantio Aug(usto) VI[II] et\\
\lnum{8}[Iu]liano Caes(are) con[s(ulibus)]}
\nf{Témoins}{Export EDCS de 2020 : \emph{Marcus Karisiae / copari carissimae / feci nobis domum / aete[r]nale(m) deposita / pr[id]ie Kal(endas) Martias bene / [me]renti in pace / [Co]nstantio Aug(usto) VI[II] et / [Iu]liano Caes(are) con[s(ulibus)]}. EDB (EDB19876) : texte concordant aux conventions de notation près (indice de similarité 0,929).}
\nf{Apparat}{⸢m⸣ : N gravé}
\nf{Traduction}{Marcus à Karisia, sa compagne très chère : j’ai fait pour nous une demeure éternelle ; déposée la veille des calendes de mars ; bien méritante, en paix ; sous le consulat de Constance Auguste pour la 8e fois et de Julien César.}
\nf{Langue}{L’éditeur signale une graphie corrigée : N gravé pour m (n pour m).}
\nf{Chronologie}{Noms consulaires : Constantius (lu en partie) et Iulianus (lu en partie). Dates du jour : déposition : pridie Kalendas Martias = 29 février 356 (partiellement restitué).}
\nf{Commentaire}{La formule in pace est conservée.}
\nf{Bibliographie}{ICUR III 6499 ; ILCV 3663 (em).}
\nf{Statut}{Lecture secondaire (EDCS) concordante avec l’EDB : lecture stabilisée entre deux bases, non collationnée sur l’ICVR ni sur le monument. L’EDCS signale une photographie (non consultée).}
\end{notice}

\annee{357}

\begin{notice}{47}{357}{ICUR n.s. I, 3166 · EDCS-40300313 · EDB33117}
\label{not:47}
\nf{Datation}{357 — directe : nom consulaire lu}
\nf{Genre}{Épitaphe ; ensembles A, B, C ; 106 lettres transmises ; 24 \% restituées.}
\nf{Références}{ICUR I 3166.}
\nf{Provenance}{Rattachement selon l’EDB : \emph{Incertae originis} (le volume I des ICVR réunit les inscriptions dont la provenance est perdue).}
\nf{Support}{Support non précisé par l’EDB. Reproductions signalées par l’EDB : dessin, ICVR I (de Rossi) (droits réservés ; non reproduites). Fiche EDB : Filippo Piazzolla, première version 2015.}
\ntext{\lnum{1}] Quintiano bene merenti qui ⸢v⸣ixit annos\\
\lnum{2}[- - -] depositus VII Idus Septemb(res) in pace Legitima\\
\lnum{3}[coniux fecit Const(antio) IIII] et Iul(iano) II cons(ulibus)}
\nf{Témoins}{Export EDCS de 2020 : \emph{] Quintiano bene merenti qui ixit annos / [3] depositus VII Idus Septemb(res) in pace Legitima / [coniux fecit Const(antio) IIII] et Iul(iano) II cons(ulibus)}. EDB (EDB33117) : texte concordant aux conventions de notation près (indice de similarité 0,964).}
\nf{Apparat}{⸢v⸣ : B gravé}
\nf{Traduction}{[…] à Quintianus, bien méritant, qui vécut […] ans ; déposé le 7e jour avant les ides de septembre, en paix ; Legitima, son épouse, a fait (ce monument) ; sous le consulat de Constance pour la 4e (sic) fois et de Julien pour la 2e fois.}
\nf{Langue}{L’éditeur signale une graphie corrigée : B gravé pour v (bêtacisme).}
\nf{Chronologie}{Noms consulaires : Constantius (restitué) et Iulianus (lu). Dates du jour : déposition : VII Idus Septembres = 7 septembre 357 (lu).}
\nf{Commentaire}{La formule in pace est conservée.}
\nf{Bibliographie}{ICUR I 3166. Bibliographie complémentaire selon l’EDB : CLRE, 638.}
\nf{Statut}{Lecture secondaire (EDCS) concordante avec l’EDB : lecture stabilisée entre deux bases, non collationnée sur l’ICVR ni sur le monument. L’EDCS signale une photographie (non consultée).}
\end{notice}

\begin{notice}{48}{357}{ICUR n.s. III, 8722 · EDCS-35900967 · EDB20990}
\label{not:48}
\nf{Datation}{357 — directe : initiales consulaires}
\nf{Genre}{Épitaphe ; ensembles A, B, C ; 55 lettres transmises.}
\nf{Références}{ICUR III 8722.}
\nf{Provenance}{Rattachement selon l’EDB : \emph{Tituli qui in coem. Callisti reperti traduntur}.}
\nf{Support}{Plaque de marbre. Inscription gravée. Fiche EDB : Carlo Carletti, première version 2003.}
\ntext{\lnum{1}Locus Secundi et Socrati(a)es(!)\\
\lnum{2}secundos d(e)p(ositi) (d)ie \{s\} VI Idus Febr(uarias) Co(nstantio ?) VIIII I(uliano ?) II}
\nf{Témoins}{EDB (EDB20990) : texte concordant aux conventions de notation près (indice de similarité 0,943).}
\nf{Apparat}{\{s\} : lettres superflues ; (!) : forme signalée par l’éditeur}
\nf{Traduction}{Emplacement de Secundus et de Socratia ; Secundus déposé le 6e jour avant les ides de février, sous (le consulat de) Constance pour la 9e fois et Julien pour la 2e fois.}
\nf{Langue}{Des lettres gravées sont jugées superflues par l’éditeur (entre accolades). L’éditeur signale au moins une forme anormale (sic). 7 abréviations développées par l’éditeur (dont Socrati(a)es, d(e)p, (ositi), (d)ie, Febr(uarias), Co(nstantio ?)).}
\nf{Chronologie}{Noms consulaires : Constantius (restitué) et Iulianus (restitué). Dates du jour : déposition : VI Idus Februarias = 8 février 357 (lu).}
\nf{Bibliographie}{ICUR III 8722. Bibliographie complémentaire selon l’EDB : CLRE, 637-638.}
\nf{Statut}{Lecture secondaire (EDCS) concordante avec l’EDB : lecture stabilisée entre deux bases, non collationnée sur l’ICVR ni sur le monument.}
\end{notice}

\annee{358}

\begin{notice}{49}{358}{ICUR n.s. VII, 17438 · EDCS-35200265 · EDB18146}
\label{not:49}
\nf{Datation}{358 — directe : nom consulaire lu}
\nf{Genre}{Épitaphe ; ensembles A, B, C ; 92 lettres transmises.}
\nf{Références}{CIL X *951,6 = ICUR VII 17438 = Epigraphica-1986-186.}
\nf{Provenance}{Rattachement selon l’EDB : \emph{Coem. Cyriacae ad s.Laurentium}.}
\nf{Support}{Support non précisé par l’EDB. Signes décrits par l’EDB : chrisme. Reproductions signalées par l’EDB : dessin, ICVR I (de Rossi) (droits réservés ; non reproduites). Fiche EDB : Miriam Ramunni, première version 2003.}
\ntext{\lnum{1}Exsuperia(!) coniux castissima\\
\lnum{2}mihique vixit\\
\lnum{3}annos XXXIIII deposita XI Kal(endas) Aug(ustas)\\
\lnum{4}Datiano et Cereale conss(ulibus)\\
\lnum{5}in pace}
\nf{Témoins}{EDB (EDB18146) : texte concordant aux conventions de notation près (indice de similarité 0,991).}
\nf{Apparat}{(!) : forme signalée par l’éditeur}
\nf{Traduction}{Exsuperia, épouse très chaste, et qui a vécu pour moi 34 ans ; déposée le 11e jour avant les calendes d’août, sous le consulat de Datianus et de Cerealis ; en paix.}
\nf{Langue}{L’éditeur signale au moins une forme anormale (sic). 4 abréviations développées par l’éditeur (dont Exsuperia(!), Kal(endas), Aug(ustas), conss(ulibus)).}
\nf{Chronologie}{Noms consulaires : Datianus (lu) et Cerealis (lu). Dates du jour : déposition : XI Kalendas Augustas = 22 juillet 358 (lu). Âge déclaré : 34 a. (lu).}
\nf{Commentaire}{La formule in pace est conservée.}
\nf{Bibliographie}{CIL X *951,6 ; ICUR VII 17438 ; Epigraphica-1986-186.}
\nf{Statut}{Lecture secondaire (EDCS) concordante avec l’EDB : lecture stabilisée entre deux bases, non collationnée sur l’ICVR ni sur le monument. L’EDCS signale une photographie (non consultée).}
\end{notice}

\begin{notice}{50}{358}{ICUR n.s. I, 1423 · EDCS-38801380 · EDB20348}
\label{not:50}
\nf{Datation}{358 — directe : nom consulaire lu}
\nf{Genre}{Épitaphe ; ensembles A, B, C ; 35 lettres transmises.}
\nf{Références}{ICUR I 1423 = ICUR VII 17437 = ILCV +2941.}
\nf{Provenance}{Rattachement selon l’EDB : \emph{Coem. Cyriacae ad s.Laurentium}. Conservation selon l’EDB : \emph{Musei Vaticani: Lapidario Cristiano}.}
\nf{Support}{Plaque de marbre. 59 × 168 × 3 cm (h. × l. × ép.). Inscription gravée. Signes décrits par l’EDB : chrisme. Fiche EDB : Miriam Ramunni, première version 2014.}
\ntext{\lnum{1}Quinta d(e)p(osita)\\
\lnum{2}Nonis Iuniis Datiano et Cere\\
\lnum{3}al⸢i⸣}
\nf{Témoins}{Export EDCS de 2020 : \emph{Quinta d(e)p(osita) / Nonis Iuniis Datiano et Cere/al}. EDB (EDB20348) : texte concordant aux conventions de notation près (indice de similarité 0,953).}
\nf{Apparat}{⸢i⸣ : E gravé}
\nf{Traduction}{Quinta, déposée aux nones de juin, sous (le consulat de) Datianus et Cerealis.}
\nf{Langue}{L’éditeur signale une graphie corrigée : E gravé pour i (confusion e/i).}
\nf{Chronologie}{Noms consulaires : Datianus (lu) et Cerealis (lu). Dates du jour : déposition : Nonis Iuniis = 5 juin 358 (lu).}
\nf{Bibliographie}{ICUR I 1423 ; ICUR VII 17437 ; ILCV +2941. Bibliographie complémentaire selon l’EDB : ILCV, 2941 ; ICVR I, 1423 ; Lega 2000, 260.}
\nf{Statut}{Lecture secondaire (EDCS) concordante avec l’EDB : lecture stabilisée entre deux bases, non collationnée sur l’ICVR ni sur le monument. L’EDCS signale une photographie (non consultée).}
\end{notice}

\begin{notice}{51}{358}{ICUR n.s. I, 1847 · EDCS-39200298 · EDB31763}
\label{not:51}
\nf{Datation}{358 — formule transmise par l’EDB (partiellement restituée)}
\nf{Genre}{Épitaphe ; ensembles A, B, C ; 25 lettres transmises ; 20 \% restituées.}
\nf{Références}{ICUR I 1847,04.}
\nf{Provenance}{Rattachement selon l’EDB : \emph{Incertae originis} (le volume I des ICVR réunit les inscriptions dont la provenance est perdue). Conservation selon l’EDB : \emph{Musei Vaticani: Lapidario Cristiano}.}
\nf{Support}{Plaque de marbre. 24 × 20 cm (h. × l.). Inscription gravée. Fiche EDB : Filippo Piazzolla, première version 2015.}
\ntext{\lnum{1}an]n(os) IIII\\
\lnum{2}[- - - N]on(as) Febr(uarias)\\
\lnum{3}[- - - bene] merenti [}
\nf{Témoins}{EDB (EDB31763 ; divergence forte) : \emph{------ / [---]n(---) IIII / [--- n]on(--) febr(uar--) / [--- Datiano et C]aereale [consulibus? ---] / ------}.}
\nf{Traduction}{[…] 4 ans […] ; […] nones de février […] ; bien méritant(e).}
\nf{Langue}{3 abréviations développées par l’éditeur (dont ann(os), Non(as), Febr(uarias)).}
\nf{Chronologie}{Le texte EDCS ne conserve aucun nom consulaire. La fiche EDB (EDB31763, d'après Lega 2000) transmet une ligne supplémentaire : \emph{[--- Datiano et C]aereale [consulibus? ---]}. Seule la fin du nom de Cerealis est lue ; l'année 358 repose désormais sur cette ligne, non plus sur la seule base. Dates du jour : événement indéterminé : … Nonas Februarias (partiellement restitué).}
\nf{Bibliographie}{ICUR I 1847,04. Bibliographie complémentaire selon l’EDB : Ferrua 1979b, 53 ; Lega 2000.}
\nf{Statut}{Lecture secondaire (EDCS) ; divergence forte avec l’EDB : le texte ne doit pas être exploité sans contrôle de l’édition imprimée.}
\end{notice}

\begin{notice}{52}{358}{ICUR n.s. II, 6022 · EDCS-34200648 · EDB17462}
\label{not:52}
\nf{Datation}{358 — directe : consulat restitué ; datation contestée}
\nf{Genre}{Épitaphe ; ensembles A, B, C ; 32 lettres transmises ; 47 \% restituées.}
\nf{Références}{ICUR II 6022.}
\nf{Provenance}{Rattachement selon l’EDB : \emph{Coem. Commodillae}.}
\nf{Support}{Plaque de marbre. Inscription gravée. Fiche EDB : Carlo Carletti, première version 2003.}
\ntext{\lnum{1}dep(ositus ?)] XIII Kal(endas) [- - -]\\
\lnum{2}[- - -] conss(ulatu) Da[tiani et Cerealis}
\nf{Témoins}{EDB (EDB17462 ; divergence forte) : \emph{[--- dep(osit---)] XIII kal(endas) [---] / [---] cons(ulibu)s Da[---] [texte tronqué dans la liste de résultats]}.}
\nf{Traduction}{[…] déposé le 13e jour avant les calendes de […] ; […] sous le consulat de Datianus et de Cerealis.}
\nf{Langue}{3 abréviations développées par l’éditeur (dont dep(ositus ?), Kal(endas), conss(ulatu)).}
\nf{Chronologie}{Noms consulaires : Datianus (restitué) et Cerealis (restitué). Seules les lettres \emph{Da[} sont lues après \emph{conss(ulatu)}. Silvagni restitue \emph{Da[tiani et Cerealis]} (358) ; la CLRE (p. 631), citée par l'EDB, rappelle que Dalmatius, consul en 333 avec Zenophilus, commence aussi par DA. La pierre est perdue et l'EDB signale une transcription différente de l'ICVR. L'année 358 est conservée comme année de classement ; si 333 était retenu, la notice sortirait du corpus.}
\nf{Bibliographie}{ICUR II 6022. Bibliographie complémentaire selon l’EDB : CLRE, 631.}
\nf{Statut}{Lecture secondaire (EDCS) ; divergence forte avec l’EDB : le texte ne doit pas être exploité sans contrôle de l’édition imprimée.}
\end{notice}

\begin{notice}{53}{358}{ICUR n.s. IV, 12420 · EDCS-37600944 · EDB41720}
\label{not:53}
\nf{Datation}{358 — directe : nom consulaire lu}
\nf{Genre}{Épitaphe ; ensembles A, B, C ; 80 lettres transmises ; 29 \% restituées.}
\nf{Références}{ICUR IV 12420 = ILCV +2690.}
\nf{Provenance}{Rattachement selon l’EDB : \emph{basilica Marci}.}
\nf{Support}{La fiche EDB ne figure pas dans l’état 2019 des données EpiDoc ; support non documenté.}
\ntext{\lnum{1}[Datiano] et C\{a\}eriale conss(ulibus)\\
\lnum{2}[dep(osita) - - -] qu(a)e vixit annos XXXII et\\
\lnum{3}[- - - in p]ace q(uae) vixit cum marito\\
\lnum{4}[ann(os) - - -] m(enses) VIII}
\nf{Témoins}{EDB (EDB41720 ; écarts à examiner) : \emph{[---? Datiano] et Caeriale cons(ulibu)s / [deposita ---] que vixit annos XXXII et / [menses --- in p]ace q(uae) vixit cum marito / [annos --- et] m(enses) VIII}.}
\nf{Apparat}{\{a\} : lettres superflues}
\nf{Traduction}{Sous le consulat de Datianus et de Cerealis […] ; déposée […], qui vécut 32 ans et […] ; en paix ; qui vécut avec son mari […] ans […] 8 mois.}
\nf{Langue}{Des lettres gravées sont jugées superflues par l’éditeur (entre accolades). 6 abréviations développées par l’éditeur (dont conss(ulibus), dep(osita), qu(a)e, q(uae), ann(os), m(enses)).}
\nf{Chronologie}{Noms consulaires : Datianus (restitué) et Cerealis (lu). Âge déclaré : 32 a. (partiellement restitué).}
\nf{Commentaire}{La formule in pace est attestée seulement par restitution.}
\nf{Bibliographie}{ICUR IV 12420 ; ILCV +2690. Bibliographie complémentaire selon l’EDB : ILCV, 2690 adn.}
\nf{Statut}{Lecture secondaire (EDCS) ; écarts avec l’EDB à arbitrer sur l’ICVR.}
\end{notice}

\annee{359}

\begin{notice}{54}{359}{ICUR n.s. I, 1424 · EDCS-39000410 · EDB31187}
\label{not:54}
\nf{Datation}{359 — directe : nom consulaire lu}
\nf{Genre}{Épitaphe ; ensembles A, B, C ; 75 lettres transmises.}
\nf{Références}{ICUR I 1424.}
\nf{Provenance}{Rattachement selon l’EDB : \emph{Incertae originis} (le volume I des ICVR réunit les inscriptions dont la provenance est perdue). Conservation selon l’EDB : \emph{Musei Vaticani: Lapidario Cristiano}.}
\nf{Support}{Plaque de marbre. 24 × 33 cm (h. × l.). Inscription gravée. Fiche EDB : Filippo Piazzolla, première version 2015.}
\ntext{\lnum{1}Eusebio et (H)ypatio cc(onsulibus)\\
\lnum{2}V Idus Iunias deces\\
\lnum{3}s⸢i⸣t Prima quae\\
\lnum{4}vixit ann⸢o⸣s XXXV\\
\lnum{5}m(enses) II deposita in pac(e)}
\nf{Témoins}{Export EDCS de 2020 : \emph{Eusebio et (H)ypatio cc(onsulibus) / V Idus Iunias deces/st Prima quae / vixit anns XXXV / m(enses) II deposita in pac(e)}. EDB (EDB31187 ; écarts à examiner) : \emph{Eusevio et Ypatio cc(onsulibus) / V idus iunias deces/set Prima quae / vixit annus XXXV / m(enses) II deposita in pac(e) / ((ornamenta duo)) [texte tronqué dans la liste de résultats]}.}
\nf{Apparat}{⸢i⸣ : E gravé ; ⸢o⸣ : V gravé}
\nf{Traduction}{Sous le consulat d’Eusebius et d’Hypatius, le 5e jour avant les ides de juin, mourut Prima, qui vécut 35 ans 2 mois ; déposée en paix.}
\nf{Langue}{L’éditeur signale 2 graphies corrigées : E gravé pour i (confusion e/i) ; V gravé pour o (confusion o/u).}
\nf{Chronologie}{Noms consulaires : Eusebius (lu) et Hypatius (lu). Dates du jour : événement indéterminé : V Idus Iunias = 9 juin 359 (lu). Âge déclaré : 35 a. 2 m. (lu).}
\nf{Commentaire}{La formule in pace est transmise sous forme abrégée, sans restitution de lacune ; son développement reste éditorial.}
\nf{Bibliographie}{ICUR I 1424. Bibliographie complémentaire selon l’EDB : ILCV, 2807A adn. ; Lega 2000, 120.}
\nf{Statut}{Lecture secondaire (EDCS) ; écarts avec l’EDB à arbitrer sur l’ICVR. L’EDCS signale une photographie (non consultée).}
\end{notice}

\begin{notice}{55}{359}{ICUR n.s. IV, 11757 · EDCS-37400851 · EDB6729}
\label{not:55}
\nf{Datation}{359 — directe : nom consulaire lu}
\nf{Genre}{Épitaphe ; ensembles A, B, C ; 165 lettres transmises ; 12 \% restituées.}
\nf{Références}{ICUR IV 11757 = ILCV 4399 (em).}
\nf{Provenance}{Rattachement selon l’EDB : \emph{Coem. anonymum ad viam Appiam}. Conservation selon l’EDB : \emph{Coem. anonymum ad viam Appiam}.}
\nf{Support}{Plaque de marbre. Inscription gravée. Signes décrits par l’EDB : chrisme et chrisme avec alpha et oméga. Reproductions signalées par l’EDB : photographie PCAS (droits réservés ; non reproduites). Fiche EDB : Maria Cristina Grisanzio, première version 2003.}
\ntext{\lnum{1}Aqundino et Proculo co[nsulibus olim na]tus\\
\lnum{2}XI Kal(endas) Iulias die ⸢V⸣eneris qu[i v]ix\{s\}it ann⸢o⸣s XVIIII\\
\lnum{3}et mens⸢e⸣s III dies VI depositus est V Kal(endas) Octobr⸢e⸣s Marcellinus\\
\lnum{4}depositus in pace\{m\} Marcellinus se [viv]u(s) fecit}
\nf{Témoins}{Export EDCS de 2020 : \emph{Aqundino et Proculo co[nsulibus olim na]tus / XI Kal(endas) Iulias die eneris qu[i v]ix\{s\}it anns XVIIII / et menss III dies VI depositus est V Kal(endas) Octobrs Marcellinus / depositus in pace\{m\} Marcellinus se [viv]u(s) fecit}. EDB (EDB6729 ; écarts à examiner) : \emph{Aqundino et Proculo co[nsulibus olim nat]us / XI kal(endas) iulias die Beneris, qu[i v]ixsit annus XVIIII / et mensis III dies VI, depositus est V kal(endas) octobris Marcellinus / depositus in pace\{m\} Marcellinus se [viv]u fecit / ((monogramma Christi cum alpha et omega)) [texte tronqué dans la liste de résultats]}.}
\nf{Apparat}{⸢V⸣ : B gravé ; ⸢o⸣ : V gravé ; ⸢e⸣ : I gravé ; ⸢e⸣ : I gravé ; \{s\} : lettres superflues ; \{m\} : lettres superflues}
\nf{Traduction}{Né jadis sous le consulat d’Acindynus et de Proculus, le 11e jour avant les calendes de juillet, un vendredi ; il vécut 19 ans 3 mois et 6 jours ; il a été déposé le 5e jour avant les calendes d’octobre ; Marcellinus déposé en paix ; Marcellinus a fait (ce tombeau) de son vivant.}
\nf{Langue}{L’éditeur signale 4 graphies corrigées : B gravé pour V (bêtacisme) ; V gravé pour o (confusion o/u) ; I gravé pour e (confusion e/i) ×2. Des lettres gravées sont jugées superflues par l’éditeur (entre accolades). 3 abréviations développées par l’éditeur (dont Kal(endas), vivu(s)).}
\nf{Chronologie}{Noms consulaires : Acindynus (absent) et Proculus (lu). Ne sont pas des consuls : depositus est … Marcellinus (défunt (la date consulaire Acindyno et Proculo, 340, est celle de la naissance)). Le texte associe une naissance sous Acindynus et Proculus (340), le 21 juin, et une déposition le 5e jour avant les calendes d'octobre. Le calcul de 19 ans 3 mois 6 jours conduit au 27 septembre 359, conformément au quantième gravé. Le jour de naissance est dit \emph{die Veneris} ; le calcul julien donne un samedi pour le 21 juin 340 : l'écart d'un jour est conservé sans correction. Le nom \emph{Marcellinus} est celui du défunt, non d'un consul. Dates du jour : naissance : XI Kalendas Iulias = 21 juin 359 (lu) ; déposition : V Kalendas Octobres = 27 septembre 359 (lu). Âge déclaré : 19 a. 3 m. 6 j. (partiellement restitué).}
\nf{Commentaire}{La formule in pace est conservée.}
\nf{Bibliographie}{ICUR IV 11757 ; ILCV 4399 (em).}
\nf{Statut}{Lecture secondaire (EDCS) ; écarts avec l’EDB à arbitrer sur l’ICVR. L’EDCS signale une photographie (non consultée).}
\end{notice}

\annee{360}

\begin{notice}{56}{360}{ICUR n.s. II, 5858 · EDCS-33100921 · EDB4428}
\label{not:56}
\nf{Datation}{360 — directe : nom consulaire lu}
\nf{Genre}{Épitaphe ; ensembles A, B, C ; 91 lettres transmises.}
\nf{Références}{ICUR II 5858 = ICUR II 5932 = ICUR V 13105 = ILCV 2883 (add, em).}
\nf{Provenance}{Rattachement selon l’EDB : \emph{Coem. subterraneum ad Catacumbas}. Conservation selon l’EDB : \emph{s. Paolo f.l.m.}.}
\nf{Support}{Plaque de marbre. 158 × 50 cm (h. × l.). Inscription gravée. Reproductions signalées par l’EDB : dessin, ICVR I (de Rossi) (droits réservés ; non reproduites). Fiche EDB : Carlo Carletti, première version 2015.}
\ntext{\lnum{1}Pelecrinus(!) in pace cum uxore\{m\} sua\{m\} Sil\\
\lnum{2}vana\{m\} qui(!) exi⸢v⸣it in pace XIIII Kal(endas) Ian(uarias)\\
\lnum{3}d(ominis) n(ostris) Const(antio) Aug(usto) X et Iul(iano) Caesare III conss(ulibus)}
\nf{Témoins}{Export EDCS de 2020 : \emph{Pelecrinus(!) in pace cum uxore\{m\} sua\{m\} Sil/vana\{m\} qui(!) exiit in pace XIIII Kal(endas) Ian(uarias) / d(ominis) n(ostris) Const(antio) Aug(usto) X et Iul(iano) Caesare III conss(ulibus)}. EDB (EDB4428) : texte concordant aux conventions de notation près (indice de similarité 0,966).}
\nf{Apparat}{⸢v⸣ : B gravé ; \{m\} : lettres superflues ; \{m\} : lettres superflues ; \{m\} : lettres superflues ; (!) : forme signalée par l’éditeur}
\nf{Traduction}{Pelegrinus, en paix, avec son épouse Silvana ; il s’en alla en paix le 14e jour avant les calendes de janvier, sous le consulat de nos seigneurs Constance Auguste pour la 10e fois et Julien César pour la 3e fois.}
\nf{Langue}{L’éditeur signale une graphie corrigée : B gravé pour v (bêtacisme). Des lettres gravées sont jugées superflues par l’éditeur (entre accolades). L’éditeur signale au moins une forme anormale (sic). 10 abréviations développées par l’éditeur (dont Pelecrinus(!), qui(!), Kal(endas), Ian(uarias), d(ominis), n(ostris)).}
\nf{Chronologie}{Noms consulaires : Constantius (lu) et Iulianus (lu). Dates du jour : décès : XIIII Kalendas Ianuarias = 19 décembre 360 (lu).}
\nf{Commentaire}{La formule in pace est conservée.}
\nf{Bibliographie}{ICUR II 5858 ; ICUR II 5932 ; ICUR V 13105 ; ILCV 2883 (add, em). Bibliographie complémentaire selon l’EDB : ICVR II, 5858 ; ICVR II, 5932.}
\nf{Statut}{Lecture secondaire (EDCS) concordante avec l’EDB : lecture stabilisée entre deux bases, non collationnée sur l’ICVR ni sur le monument. L’EDCS signale une photographie (non consultée).}
\end{notice}

\begin{notice}{57}{360}{ICUR n.s. III, 8141 · EDCS-35100217 · EDB24664}
\label{not:57}
\nf{Datation}{360 — directe : nom consulaire lu}
\nf{Genre}{Épitaphe ; ensembles A, B, C ; 93 lettres transmises ; 23 \% restituées.}
\nf{Références}{ICUR III 8141 = ILCV +2090.}
\nf{Provenance}{Rattachement selon l’EDB : \emph{Basilica ss. Nerei et Achillei in coem. Domitillae}. Conservation selon l’EDB : \emph{Basilica ss. Nerei et Achillei in coem. Domitillae}.}
\nf{Support}{Plaque de marbre. 44 × 133 × 9 cm (h. × l. × ép.). Inscription gravée. Reproductions signalées par l’EDB : photographie PCAS (droits réservés ; non reproduites). Fiche EDB : Carlo Carletti / Antonio Enrico Felle, première version 2003.}
\ntext{\lnum{1}[DD(ominis)] nn(ostris) Constantio Aug(usto) X\\
\lnum{2}[et] Iul(iano) III conss(ulibus) dep(ositus) pri(di)e\\
\lnum{3}[N]on(as) Maias F[l(avius)] Car(a)ba\\
\lnum{4}sinus qui vixit ann[os - - -]\\
\lnum{5}[me]n(ses) XI in pace}
\nf{Témoins}{EDB (EDB24664) : texte concordant aux conventions de notation près (indice de similarité 0,947).}
\nf{Traduction}{Sous le consulat de nos seigneurs Constance Auguste pour la 10e fois et Julien pour la 3e fois, déposé la veille des nones de mai, Flavius Carabasinus, qui vécut […] ans 11 mois, en paix.}
\nf{Langue}{11 abréviations développées par l’éditeur (dont DD(ominis), nn(ostris), Aug(usto), Iul(iano), conss(ulibus), dep(ositus)).}
\nf{Chronologie}{Noms consulaires : Constantius (lu) et Iulianus (lu). Dates du jour : déposition : Nonas Maias = 7 mai 360 (partiellement restitué). Âge déclaré : 11 m. (partiellement restitué).}
\nf{Commentaire}{La formule in pace est conservée.}
\nf{Bibliographie}{ICUR III 8141 ; ILCV +2090. Bibliographie complémentaire selon l’EDB : IC I Supplementum, 1495 ; ILCV, 2090 adn.}
\nf{Statut}{Lecture secondaire (EDCS) concordante avec l’EDB : lecture stabilisée entre deux bases, non collationnée sur l’ICVR ni sur le monument. L’EDCS signale une photographie (non consultée).}
\end{notice}

\begin{notice}{58}{360}{ICUR n.s. IV, 11093 · EDCS-37200375 · EDB39552}
\label{not:58}
\nf{Datation}{360 — directe : nom consulaire lu}
\nf{Genre}{Épitaphe ; ensembles A, B, C ; 63 lettres transmises.}
\nf{Références}{ICUR IV 11093 = ILCV +2807a.}
\nf{Provenance}{Rattachement selon l’EDB : \emph{Coem. Callisti area subdialis et incerti loci}. Conservation selon l’EDB : \emph{Musei Vaticani: Lapidario Cristiano}.}
\nf{Support}{Plaque de marbre. 27 × 123 cm (h. × l.). Inscription gravée. Fiche EDB : Marida Pierno, première version 2015.}
\ntext{\lnum{1}DD(ominis) n(ostris) Cons(tan)tio Aug(usto) X et Iuliano C(a)esar(e) III\\
\lnum{2}conss(ulibus) VIIII Kal(endas) Ma(ia)s decessit Ianuaras}
\nf{Témoins}{EDB (EDB39552) : texte concordant aux conventions de notation près (indice de similarité 0,993).}
\nf{Traduction}{Sous le consulat de nos seigneurs Constance Auguste pour la 10e fois et Julien César pour la 3e fois, le 9e jour avant les calendes de mai, mourut Ianuara.}
\nf{Langue}{9 abréviations développées par l’éditeur (dont DD(ominis), n(ostris), Cons(tan)tio, Aug(usto), C(a)esar, (e)).}
\nf{Chronologie}{Noms consulaires : Constantius (lu) et Iulianus (lu). Dates du jour : événement indéterminé : VIIII Kalendas Maias = 23 avril 360 (lu).}
\nf{Bibliographie}{ICUR IV 11093 ; ILCV +2807a. Bibliographie complémentaire selon l’EDB : Lega 2000.}
\nf{Statut}{Lecture secondaire (EDCS) concordante avec l’EDB : lecture stabilisée entre deux bases, non collationnée sur l’ICVR ni sur le monument. L’EDCS signale une photographie (non consultée).}
\end{notice}

\annee{361}

\begin{notice}{59}{361}{ICUR n.s. III, 8142 · EDCS-35100218 · EDB24665}
\label{not:59}
\nf{Datation}{361 — directe : nom consulaire lu}
\nf{Genre}{Épitaphe ; ensembles A, B, C ; 122 lettres transmises.}
\nf{Références}{CIL VI 32011 (p 4803) = ICUR III 8142 = ILCV 288 (add, em).}
\nf{Provenance}{Rattachement selon l’EDB : \emph{Basilica ss. Nerei et Achillei in coem. Domitillae}. Conservation selon l’EDB : \emph{Basilica ss. Nerei et Achillei in coem. Domitillae}.}
\nf{Support}{Plaque de marbre. 40 × 117 cm (h. × l.). Inscription gravée. Reproductions signalées par l’EDB : photographie PCAS (droits réservés ; non reproduites). Fiche EDB : Carlo Carletti / Antonio Enrico Felle, première version 2018.}
\ntext{\lnum{1}Dulcissimo filio Mariniano\\
\lnum{2}qui vixit annos XI mens(es) IIII\\
\lnum{3}dies XIII dep(ositus) XI Kal(endas) Aug(ustas) qu(i)es⸢c⸣et\\
\lnum{4}in pace Tauro et Florentio con(sulibus)\\
\lnum{5}Fl(avius) Concordius v(ir) p(erfectissimus) filio fecit}
\nf{Témoins}{Export EDCS de 2020 : \emph{Dulcissimo filio Mariniano / qui vixit annos XI mens(es) IIII / dies XIII dep(ositus) XI Kal(endas) Aug(ustas) qu(i)eset / in pace Tauro et Florentio con(sulibus) / Fl(avius) Concordius v(ir) p(erfectissimus) filio fecit}. EDB (EDB24665) : texte concordant aux conventions de notation près (indice de similarité 0,937).}
\nf{Apparat}{⸢c⸣ : QV gravé}
\nf{Traduction}{À son fils très doux Marinianus, qui vécut 11 ans 4 mois 13 jours ; déposé le 11e jour avant les calendes d’août ; il repose en paix ; sous le consulat de Taurus et de Florentius ; Flavius Concordius, perfectissime, a fait (ce monument) pour son fils.}
\nf{Langue}{L’éditeur signale une graphie corrigée : QV gravé pour c.}
\nf{Chronologie}{Noms consulaires : Taurus (lu) et Florentius (lu). Dates du jour : déposition : XI Kalendas Augustas = 22 juillet 361 (lu). Âge déclaré : 11 a. 4 m. 13 j. (lu).}
\nf{Commentaire}{La formule in pace est conservée. Vocabulaire religieux : repos et sommeil de la mort — quiescet (lu).}
\nf{Bibliographie}{CIL VI 32011 (p 4803) ; ICUR III 8142 ; ILCV 288 (add, em). Bibliographie complémentaire selon l’EDB : CIL VI, 32011 ; CIL VI, 288 ; IC I Supplementum, 1499 ; ILCV, 288.}
\nf{Statut}{Lecture secondaire (EDCS) concordante avec l’EDB : lecture stabilisée entre deux bases, non collationnée sur l’ICVR ni sur le monument. L’EDCS signale une photographie (non consultée).}
\end{notice}

\begin{notice}{60}{361}{ICUR n.s. I, 2979 · EDCS-39900423 · EDB32878}
\label{not:60}
\nf{Datation}{361 — directe : nom consulaire lu}
\nf{Genre}{Épitaphe ; ensembles A, B, C ; 85 lettres transmises.}
\nf{Références}{ICUR I 2979 = ILCV 4368 = Bencivenni-2018, 00038,1.}
\nf{Provenance}{Rattachement selon l’EDB : \emph{Incertae originis} (le volume I des ICVR réunit les inscriptions dont la provenance est perdue). Conservation selon l’EDB : \emph{Ravenna   ; Museo Nazionale}.}
\nf{Support}{Plaque de marbre. 25 × 38 × 2 cm (h. × l. × ép.). Inscription gravée. Signes décrits par l’EDB : chrisme. Fiche EDB : Filippo Piazzolla, première version 2015.}
\ntext{\lnum{1}(H)ic positus e⸢st⸣ Leo qui\\
\lnum{2}vixit annos XXVI de\\
\lnum{3}positus VIII Idus O\\
\lnum{4}ctobr⸢e⸣s natus est\\
\lnum{5}Tauro et Florentio\\
\lnum{6}cc(on)ss(ulibus)}
\nf{Témoins}{Export EDCS de 2020 : \emph{(H)ic positus e Leo qui / vixit annos XXVI de/positus VIII Idus O/ctobrs natus est / Tauro et Florentio / cc(on)ss(ulibus)}. EDB (EDB32878) : texte concordant aux conventions de notation près (indice de similarité 0,981).}
\nf{Apparat}{⸢st⸣ : TS gravé ; ⸢e⸣ : I gravé}
\nf{Traduction}{Ici a été placé Leo, qui vécut 26 ans ; déposé le 8e jour avant les ides d’octobre ; il était né sous le consulat de Taurus et de Florentius.}
\nf{Langue}{L’éditeur signale 2 graphies corrigées : TS gravé pour st ; I gravé pour e (confusion e/i).}
\nf{Chronologie}{Noms consulaires : Taurus (lu) et Florentius (lu). Dates du jour : déposition : VIII Idus Octobres = 8 octobre 361 (lu). Âge déclaré : 26 a. (lu).}
\nf{Bibliographie}{ICUR I 2979 ; ILCV 4368 ; Bencivenni-2018, 00038,1. Bibliographie complémentaire selon l’EDB : ILCV, 4368.}
\nf{Statut}{Lecture secondaire (EDCS) concordante avec l’EDB : lecture stabilisée entre deux bases, non collationnée sur l’ICVR ni sur le monument.}
\end{notice}

\begin{notice}{61}{361}{ICUR n.s. II, 6023 · EDCS-34200649 · EDB17463}
\label{not:61}
\nf{Datation}{361 — directe : nom consulaire lu}
\nf{Genre}{Épitaphe ; ensembles A, B, C ; 54 lettres transmises.}
\nf{Références}{ICUR II 6023 = ILCV 2569 = AE 1906, 83.}
\nf{Provenance}{Rattachement selon l’EDB : \emph{Coem. Commodillae}.}
\nf{Support}{Plaque de marbre. Inscription gravée. Signes décrits par l’EDB : chrisme. Fiche EDB : Carlo Carletti, première version 2015.}
\ntext{\lnum{1}Tertulla in pace\\
\lnum{2}quae vixit ann(os)\\
\lnum{3}VII menses VI\\
\lnum{4}d(ies) XVIII Tauro et Flor(entio)}
\nf{Témoins}{EDB (EDB17463) : texte concordant aux conventions de notation près (indice de similarité 0,937).}
\nf{Traduction}{Tertulla, en paix, qui vécut 7 ans 6 mois 18 jours ; sous (le consulat de) Taurus et Florentius.}
\nf{Langue}{3 abréviations développées par l’éditeur (dont ann(os), d(ies), Flor(entio)).}
\nf{Chronologie}{Noms consulaires : Taurus (lu) et Florentius (lu). Âge déclaré : 7 a. 6 m. 18 j. (lu).}
\nf{Commentaire}{La formule in pace est conservée.}
\nf{Bibliographie}{ICUR II 6023 ; ILCV 2569 ; AE 1906, 83.}
\nf{Statut}{Lecture secondaire (EDCS) concordante avec l’EDB : lecture stabilisée entre deux bases, non collationnée sur l’ICVR ni sur le monument.}
\end{notice}

\annee{362}

\begin{notice}{62}{362}{ICUR n.s. IV, 12526 · EDCS-18600145 · EDB41423}
\label{not:62}
\nf{Datation}{362 — directe : nom consulaire lu}
\nf{Genre}{Épitaphe ; ensembles A, B, C ; 55 lettres transmises ; 36 \% restituées.}
\nf{Références}{CIL VI 7429 = ILCV +4144b = ICUR IV 12526.}
\nf{Provenance}{Rattachement selon l’EDB : \emph{inter vias Appiam et Ardeatinam}.}
\nf{Support}{Plaque de marbre. Inscription gravée. Fiche EDB : Filippo Piazzolla, première version 2016.}
\ntext{\lnum{1}D(omino) n(ostro) ⸢C⸣l(audio) Mamertino et\\
\lnum{2}Fl(avio) Nev[itta co(n)s(ulibus)]\\
\lnum{3}Ac() Euda[emon(?) be]\\
\lnum{4}ne merenti}
\nf{Témoins}{Export EDCS de 2020 : \emph{D(omino) n(ostro) l(audio) Mamertino et / Fl(avio) Nev[itta co(n)s(ulibus)] / Ac() Euda[emon(?) be]/ne merenti}. EDB (EDB41423) : texte concordant aux conventions de notation près (indice de similarité 0,945).}
\nf{Apparat}{⸢C⸣ : K gravé}
\nf{Traduction}{Sous le consulat de notre seigneur (?) Claudius Mamertinus et de Flavius Nevitta […] ; Eudaemon, bien méritant.}
\nf{Langue}{L’éditeur signale une graphie corrigée : K gravé pour C.}
\nf{Chronologie}{Noms consulaires : Mamertinus (lu) et Nevitta (lu en partie).}
\nf{Bibliographie}{CIL VI 7429 ; ILCV +4144b ; ICUR IV 12526. Bibliographie complémentaire selon l’EDB : CIL VI, 7429 ; CLRE, 640 ; ILCV, 4144B adn.}
\nf{Statut}{Lecture secondaire (EDCS) concordante avec l’EDB : lecture stabilisée entre deux bases, non collationnée sur l’ICVR ni sur le monument.}
\end{notice}

\begin{notice}{63}{362}{ICUR n.s. I, 82 · EDCS-38600098 · EDB27455}
\label{not:63}
\nf{Datation}{362 — directe : nom consulaire lu}
\nf{Genre}{Épitaphe ; ensembles A, B, C ; 56 lettres transmises ; 2 \% restituées.}
\nf{Références}{ICUR I 82.}
\nf{Provenance}{Rattachement selon l’EDB : \emph{Incertae originis} (le volume I des ICVR réunit les inscriptions dont la provenance est perdue). Conservation selon l’EDB : \emph{s. Cecilia in Trastevere}.}
\nf{Support}{Plaque de marbre. 35 × 72 × 3,5 cm (h. × l. × ép.). Inscription gravée. Reproductions signalées par l’EDB : dessin, ICVR I (de Rossi) (droits réservés ; non reproduites). Fiche EDB : Filippo Piazzolla, première version 2013.}
\ntext{\lnum{1}] in pace qui vixit annis V[- - -]\\
\lnum{2}[- - -]XL dep(ositus) XV Kal(endas) Aug(ustas)\\
\lnum{3}[- - - M]amertino et Ne⸢v⸣i⸢t⸣\\
\lnum{4}ta conss(ulibus)}
\nf{Témoins}{Export EDCS de 2020 : \emph{] in pace qui vixit annis V[3] / [3]XL dep(ositus) XV Kal(endas) Aug(ustas) / [3 M]amertino et Nei/ta conss(ulibus)}. EDB (EDB27455 ; écarts à examiner) : \emph{------ / [---] in pace qui vixit annis V / [---? diebus] XL dep(ositus) XV kal(endas) aug(ustas) / [--- M]amertino et Nebi⌜t⌝/ta cons(ulibu)s / ((avis cum olea)) ((statera)) [texte tronqué dans la liste de résultats]}.}
\nf{Apparat}{⸢v⸣ : B gravé ; ⸢t⸣ : I gravé}
\nf{Traduction}{[…] en paix, qui vécut […] ans […] ; déposé le 15e jour avant les calendes d’août, sous le consulat de Mamertinus et de Nevitta.}
\nf{Langue}{L’éditeur signale 2 graphies corrigées : B gravé pour v (bêtacisme) ; I gravé pour t.}
\nf{Chronologie}{Noms consulaires : Mamertinus (lu en partie) et Nevitta (lu). Dates du jour : déposition : XV Kalendas Augustas = 18 juillet 362 (lu).}
\nf{Commentaire}{La formule in pace est conservée.}
\nf{Bibliographie}{ICUR I 82. Bibliographie complémentaire selon l’EDB : Ferrua 1979b, 10, n. 82 ; Bisconti 2000a, 254, n. XXIIIa2.1.}
\nf{Statut}{Lecture secondaire (EDCS) ; écarts avec l’EDB à arbitrer sur l’ICVR. L’EDCS signale une photographie (non consultée).}
\end{notice}

\begin{notice}{64}{362}{ICUR n.s. I, 2805 · EDCS-11700571 · EDB32728}
\label{not:64}
\nf{Datation}{362 — directe : nom consulaire lu}
\nf{Genre}{Épitaphe ; ensembles A, B, C ; 74 lettres transmises ; 43 \% restituées.}
\nf{Références}{ICUR I 2805 = ILMN-01, 00440.}
\nf{Provenance}{Rattachement selon l’EDB : \emph{Incertae originis} (le volume I des ICVR réunit les inscriptions dont la provenance est perdue). Conservation selon l’EDB : \emph{Napoli   ; Museo Archeologico Nazionale}.}
\nf{Support}{Plaque de marbre. 24 × 26 × 4,5 cm (h. × l. × ép.). Inscription gravée. Reproductions signalées par l’EDB : reproduction, Camodeca et alii 2000 (droits réservés ; non reproduites). Fiche EDB : Filippo Piazzolla, première version 2015.}
\ntext{\lnum{1}Mamert[ino et Nevit]\\
\lnum{2}ta co(n)ss(ulibus) XI[- - -]\\
\lnum{3}parentes fi[lio dulcissi]mo Diocleti[ano qui vixit]\\
\lnum{4}annis XVI[}
\nf{Témoins}{EDB (EDB32728 ; écarts à examiner) : \emph{Mamert[ino et Nevit]/ta co(n)s(ulibu)s XI[--- kalendas ---] / parentes fi[lio dulcissi]/mo Diocleti[ano qui vixit] / annis XVI[---] [texte tronqué dans la liste de résultats]}.}
\nf{Traduction}{Sous le consulat de Mamertinus et de Nevitta, 11 […] ; les parents à leur fils très doux Diocletianus, qui vécut 16 ans.}
\nf{Langue}{2 abréviations développées par l’éditeur (dont co(n)ss, (ulibus)).}
\nf{Chronologie}{Noms consulaires : Mamertinus (lu) et Nevitta (restitué). Âge déclaré : 16 a. (partiellement restitué).}
\nf{Bibliographie}{ICUR I 2805 ; ILMN-01, 00440. Bibliographie complémentaire selon l’EDB : Camodeca et alii 2000, 152, n. 440 con foto.}
\nf{Statut}{Lecture secondaire (EDCS) ; écarts avec l’EDB à arbitrer sur l’ICVR. L’EDCS signale une photographie (non consultée).}
\end{notice}

\begin{notice}{65}{362}{ICUR n.s. II, 4804 · EDCS-36101285}
\label{not:65}
\nf{Datation}{362 — directe : nom consulaire lu}
\nf{Genre}{Épitaphe ; ensembles A, B, C ; 46 lettres transmises ; 52 \% restituées.}
\nf{Références}{ICUR II 4804.}
\nf{Provenance}{Rome (indication de l’EDCS) ; aucun rattachement cimétérial transmis.}
\nf{Support}{Non documenté dans les sources consultées.}
\ntext{\lnum{1}Lacuriu[s - - -]\\
\lnum{2}in omni a[ctu(?) - - -]\\
\lnum{3}Mamertin[o et Nevitta cons(ulibus ?)}
\nf{Témoins}{Aucune fiche EDB appariée (numéro ICUR absent de la liste publique de l’EDB).}
\nf{Traduction}{Lacurius […] ; en toute action […] ; sous le consulat de Mamertinus et de Nevitta.}
\nf{Langue}{2 abréviations développées par l’éditeur (dont actu(?), cons(ulibus ?)).}
\nf{Chronologie}{Noms consulaires : Mamertinus (lu) et Nevitta (restitué).}
\nf{Bibliographie}{ICUR II 4804.}
\nf{Statut}{Lecture secondaire (EDCS) sans second témoin numérique. Plus de la moitié des lettres sont restituées.}
\end{notice}

\begin{notice}{66}{362}{ICUR n.s. III, 6500 · EDCS-34801300 · EDB19877}
\label{not:66}
\nf{Datation}{362 — directe : nom consulaire lu}
\nf{Genre}{Épitaphe ; ensembles A, B, C ; 46 lettres transmises.}
\nf{Références}{ICUR III 6500 = ILCV +2967.}
\nf{Provenance}{Rattachement selon l’EDB : \emph{Coem. Domitillae pars inferior}. Conservation selon l’EDB : \emph{Coem. Domitillae pars inferior}.}
\nf{Support}{Plaque de marbre. 29 × 61 × 3,7 cm (h. × l. × ép.). Inscription gravée. Reproductions signalées par l’EDB : photographie PCAS (droits réservés ; non reproduites). Fiche EDB : Carlo Carletti / Antonio Enrico Felle, première version 2015.}
\ntext{\lnum{1}Mame(r)tino e(t) (N)\\
\lnum{2}e⸢vi⸣t(t)a con(sulibus) Se(cu)\\
\lnum{3}ndian(a)es(!) IIX K(alendas) Dece(mbres) deposita in (p)ac(e)}
\nf{Témoins}{Export EDCS de 2020 : \emph{Mame(r)tino e(t) (N)/et(t)a con(sulibus) Se(cu)/ndian(a)es(!) IIX K(alendas) Dece(mbres) deposita in (p)ac(e)}. EDB (EDB19877 ; écarts à examiner) : \emph{Mame⟨r⟩tino e⟨t N⟩/ebeta con(sulibus) Se⟨cu⟩/ndianes IIX k(alendas) / dece(mbres) deposita in ⟨p⟩ac⟨e⟩ [texte tronqué dans la liste de résultats]}.}
\nf{Apparat}{⸢vi⸣ : BE gravé ; (!) : forme signalée par l’éditeur}
\nf{Traduction}{Sous le consulat de Mamertinus et de Nevitta, Secundiana, le 8e jour avant les calendes de décembre, déposée en paix.}
\nf{Langue}{L’éditeur signale une graphie corrigée : BE gravé pour vi. L’éditeur signale au moins une forme anormale (sic). 11 abréviations développées par l’éditeur (dont Mame(r)tino, e(t), (N), t(t)a, con(sulibus), Se(cu)).}
\nf{Chronologie}{Noms consulaires : Mamertinus (lu) et Nevitta (lu). Dates du jour : événement indéterminé : IIX Kalendas Decembres = 22 novembre 362 (lu).}
\nf{Commentaire}{La formule in pace est transmise sous forme abrégée, sans restitution de lacune ; son développement reste éditorial.}
\nf{Bibliographie}{ICUR III 6500 ; ILCV +2967.}
\nf{Statut}{Lecture secondaire (EDCS) ; écarts avec l’EDB à arbitrer sur l’ICVR. L’EDCS signale une photographie (non consultée).}
\end{notice}

\begin{notice}{67}{362}{ICUR n.s. III, 8143 · EDCS-35100219 · EDB24666}
\label{not:67}
\nf{Datation}{362 — directe : nom consulaire lu}
\nf{Genre}{Épitaphe ; ensembles A, B, C ; 101 lettres transmises ; 38 \% restituées.}
\nf{Références}{ICUR III 8143 = ILCV 1268 (em).}
\nf{Provenance}{Rattachement selon l’EDB : \emph{Basilica ss. Nerei et Achillei in coem. Domitillae}. Conservation selon l’EDB : \emph{Basilica ss. Nerei et Achillei in coem. Domitillae}.}
\nf{Support}{Plaque de marbre. 44 × 127 × 4 cm (h. × l. × ép.). Inscription gravée. Reproductions signalées par l’EDB : photographie PCAS (droits réservés ; non reproduites). Fiche EDB : Carlo Carletti / Antonio Enrico Felle, première version 2003.}
\ntext{\lnum{1}Iuli[us Innocen]tius lect[or]\\
\lnum{2}ec(c)les[iae sanctae c]at(h)olic(a)e\\
\lnum{3}vi\{c\}x[it pl(us) m(inus)] ann(os) LXXIII d(e)p(ositus) XIII K(alendas) Iun(ias)\\
\lnum{4}Mame[rtino] et Nevita conss(ulibus)}
\nf{Témoins}{EDB (EDB24666) : texte concordant aux conventions de notation près (indice de similarité 0,93).}
\nf{Apparat}{\{c\} : lettres superflues}
\nf{Traduction}{Iulius Innocentius, lecteur de la sainte Église catholique, vécut environ 73 ans ; déposé le 13e jour avant les calendes de juin, sous le consulat de Mamertinus et de Nevitta.}
\nf{Langue}{Des lettres gravées sont jugées superflues par l’éditeur (entre accolades). 11 abréviations développées par l’éditeur (dont ec(c)les, cat(h)olic, (a)e, pl(us), m(inus), ann(os)).}
\nf{Chronologie}{Noms consulaires : Mamertinus (lu en partie) et Nevitta (lu). Dates du jour : déposition : XIII Kalendas Iunias = 20 mai 362 (lu). Âge déclaré : env. 73 a. (partiellement restitué).}
\nf{Commentaire}{Vocabulaire religieux : Église (épithète) — sanctae (entièrement restitué) ; vocabulaire ecclésial — ecclesiae (partiellement restitué). Fonction ou métier mentionné : lector.}
\nf{Bibliographie}{ICUR III 8143 ; ILCV 1268 (em).}
\nf{Statut}{Lecture secondaire (EDCS) concordante avec l’EDB : lecture stabilisée entre deux bases, non collationnée sur l’ICVR ni sur le monument. L’EDCS signale une photographie (non consultée).}
\end{notice}

\begin{notice}{68}{362}{ICUR n.s. III, 8421 · EDCS-35802008 · EDB25696}
\label{not:68}
\nf{Datation}{362 — directe : nom consulaire lu}
\nf{Genre}{Épitaphe ; ensembles A, B, C ; 105 lettres transmises.}
\nf{Références}{ICUR III 8421 (p 408) = ILCV 2978b.}
\nf{Provenance}{Rattachement selon l’EDB : \emph{Incertum locum in coem. Domitillae}. Conservation selon l’EDB : \emph{Musei Vaticani: Lapidario Cristiano}.}
\nf{Support}{Plaque de marbre. Inscription gravée. Signes décrits par l’EDB : chrisme et chrisme avec alpha et oméga. Reproductions signalées par l’EDB : dessin, ICVR I (de Rossi) (droits réservés ; non reproduites). Fiche EDB : Carlo Carletti / Antonio Enrico Felle, première version 2014.}
\ntext{\lnum{1}Cl(audio) Mamertino et Fl(avio) Nevitta conss(ulibus)\\
\lnum{2}dulcissimo filio Petrio qui vixit ann(os) XII m(ensem) I d(ies) XV dep(ositus) prid(ie) Kal(endas) Sept(embres) in pace\\
\lnum{3}Eutyches pater fecit}
\nf{Témoins}{EDB (EDB25696) : texte concordant aux conventions de notation près (indice de similarité 0,992).}
\nf{Traduction}{Sous le consulat de Claudius Mamertinus et de Flavius Nevitta, à son fils très doux Petrius, qui vécut 12 ans 1 mois 15 jours, déposé la veille des calendes de septembre, en paix ; Eutychès, son père, a fait (ce monument).}
\nf{Langue}{10 abréviations développées par l’éditeur (dont Cl(audio), Fl(avio), conss(ulibus), ann(os), m(ensem), d(ies)).}
\nf{Chronologie}{Noms consulaires : Mamertinus (lu) et Nevitta (lu). Ne sont pas des consuls : Eutyches pater fecit (dédicant (père)). Dates du jour : déposition : pridie Kalendas Septembres = 31 août 362 (lu). Âge déclaré : 12 a. 1 m. 15 j. (lu).}
\nf{Commentaire}{La formule in pace est conservée. La photographie universitaire ILCV 2978B (Ohio State University) montre PETRIO, IN PACE, le collège MAMERTINO / NEVITTA et la séquence EVTYCHES PATER FECIT : \emph{Eutyches} est le père commémorant, non un consul. L'EDB décrit un chrisme.}
\nf{Bibliographie}{ICUR III 8421 (p 408) ; ILCV 2978b.}
\nf{Statut}{Lecture secondaire (EDCS) concordante avec l’EDB : lecture stabilisée entre deux bases, non collationnée sur l’ICVR ni sur le monument. L’EDCS signale une photographie (non consultée).}
\end{notice}

\begin{notice}{69}{362}{ICUR n.s. IV, 9557 · EDCS-36101086 · EDB20873}
\label{not:69}
\nf{Datation}{362 — directe : lecture partielle d’un nom consulaire}
\nf{Genre}{Épitaphe ; ensembles A, B, C ; 69 lettres transmises ; 32 \% restituées.}
\nf{Références}{ICUR IV 9557.}
\nf{Provenance}{Rattachement selon l’EDB : \emph{Coem. Callisti pars inferior}.}
\nf{Support}{Plaque de marbre. Inscription gravée. Reproductions signalées par l’EDB : photographie PCAS (droits réservés ; non reproduites). Fiche EDB : Antonio Enrico Felle, première version 2003.}
\ntext{\lnum{1}an]n⸢o⸣s\\
\lnum{2}[- - - dies] XXV depo\\
\lnum{3}[situs(?) - - -] Idus [S]eptembr\\
\lnum{4}[es Ma]mertino et Neve\\
\lnum{5}[tta co]nsulibu[s in] pace}
\nf{Témoins}{Export EDCS de 2020 : \emph{an]ns / [3 dies] XXV depo/[situs(?) 3] Idus [S]eptembr/[es Ma]mertino et Neve/[tta co]nsulibu[s in] pace}. EDB (EDB20873) : texte concordant aux conventions de notation près (indice de similarité 0,933).}
\nf{Apparat}{⸢o⸣ : V gravé}
\nf{Traduction}{[…] ans […] 25 jours ; déposé […] ides de septembre, sous le consulat de Mamertinus et de Nevitta, en paix.}
\nf{Langue}{L’éditeur signale une graphie corrigée : V gravé pour o (confusion o/u).}
\nf{Chronologie}{Noms consulaires : Mamertinus (lu en partie). Dates du jour : déposition : … Idus Septembres (partiellement restitué).}
\nf{Commentaire}{La formule in pace est attestée seulement par restitution.}
\nf{Bibliographie}{ICUR IV 9557.}
\nf{Statut}{Lecture secondaire (EDCS) concordante avec l’EDB : lecture stabilisée entre deux bases, non collationnée sur l’ICVR ni sur le monument. L’EDCS signale une photographie (non consultée).}
\end{notice}

\begin{notice}{70}{362}{ICUR n.s. IV, 9558 · EDCS-36101087 · EDB18727}
\label{not:70}
\nf{Datation}{362 — directe : initiales consulaires}
\nf{Genre}{Épitaphe ; ensembles A, B, C ; 126 lettres transmises.}
\nf{Références}{ICUR IV 9558 = ILCV 1563 = ICaRoma 85 = Carletti-2008, 00056.}
\nf{Provenance}{Rattachement selon l’EDB : \emph{Coem. Callisti pars inferior}.}
\nf{Support}{Plaque de marbre. Inscription gravée. Reproductions signalées par l’EDB : photographie PCAS (droits réservés ; non reproduites). Fiche EDB : Antonio Enrico Felle, première version 2014.}
\ntext{\lnum{1}Mirae bonitatis Secund(a)e\\
\lnum{2}quae vixit pura fide annos\\
\lnum{3}viginti pudica cessavit\\
\lnum{4}in pace \{id\} virgo fidelis\\
\lnum{5}bene merenti quiesc⸢i⸣t Id(ibus) Iul(iis)\\
\lnum{6}palumbo sine felle M(amertino) et N(evitta)}
\nf{Témoins}{Export EDCS de 2020 : \emph{Mirae bonitatis Secund(a)e / quae vixit pura fide annos / viginti pudica cessavit / in pace \{id\} virgo fidelis / bene merenti quiesct Id(ibus) Iul(iis) / palumbo sine felle M(amertino) et N(evitta)}. EDB (EDB18727) : texte concordant aux conventions de notation près (indice de similarité 0,99).}
\nf{Apparat}{⸢i⸣ : E gravé ; \{id\} : lettres superflues}
\nf{Traduction}{À Secunda, d’une bonté admirable, qui vécut d’une foi pure vingt ans ; chaste, elle cessa (de vivre) en paix, vierge fidèle ; bien méritante, elle repose ; aux ides de juillet ; colombe sans fiel ; sous (le consulat de) Mamertinus et Nevitta.}
\nf{Langue}{L’éditeur signale une graphie corrigée : E gravé pour i (confusion e/i). Des lettres gravées sont jugées superflues par l’éditeur (entre accolades). 5 abréviations développées par l’éditeur (dont Secund(a)e, Id(ibus), Iul(iis), M(amertino), N(evitta)).}
\nf{Chronologie}{Noms consulaires : Mamertinus (restitué) et Nevitta (restitué). Dates du jour : décès : Idibus Iuliis = 15 juillet 362 (lu). Âge déclaré : 20 a. (lu).}
\nf{Commentaire}{La formule in pace est conservée. Vocabulaire religieux : initiation et foi — fide (lu), fidelis (lu) ; repos et sommeil de la mort — quiescit (lu). Commentaire développé : dossier p.~\pageref{dos:70}.}
\nf{Bibliographie}{ICUR IV 9558 ; ILCV 1563 ; ICaRoma 85 ; Carletti-2008, 00056. Bibliographie complémentaire selon l’EDB : Carletti 2008, 174-175, n. 56 ; CLRE, 640.}
\nf{Statut}{Lecture secondaire (EDCS) concordante avec l’EDB : lecture stabilisée entre deux bases, non collationnée sur l’ICVR ni sur le monument. L’EDCS signale une photographie (non consultée).}
\end{notice}

\begin{notice}{71}{362}{ICUR n.s. IV, 11094 · EDCS-37200376 · EDB39553}
\label{not:71}
\nf{Datation}{362 — directe : lecture partielle d’un nom consulaire}
\nf{Genre}{Épitaphe ; ensembles A, B, C ; 47 lettres transmises ; 70 \% restituées.}
\nf{Références}{ICUR IV 11094.}
\nf{Provenance}{Rattachement selon l’EDB : \emph{Coem. Callisti area subdialis et incerti loci}. Conservation selon l’EDB : \emph{Coem. Callisti area subdialis et incerti loci}.}
\nf{Support}{Plaque de marbre. 12 × 13 × 4,5 cm (h. × l. × ép.). Inscription gravée. Fiche EDB : Marida Pierno, première version 2015.}
\ntext{\lnum{1}Mamertino e]t Nev[itta vv(iris) cc(larissimis) conss(ulibus)}
\nf{Témoins}{EDB (EDB39553 ; écarts à examiner) : \emph{------ / [--- Mamertino e]t Neṿ[itto ---] / ------}.}
\nf{Traduction}{Sous le consulat de Mamertinus et de Nevitta, clarissimes.}
\nf{Langue}{3 abréviations développées par l’éditeur (dont vv(iris), cc(larissimis), conss(ulibus)).}
\nf{Chronologie}{Noms consulaires : Mamertinus (restitué) et Nevitta (lu en partie).}
\nf{Bibliographie}{ICUR IV 11094.}
\nf{Statut}{Lecture secondaire (EDCS) ; écarts avec l’EDB à arbitrer sur l’ICVR. Plus de la moitié des lettres sont restituées.}
\end{notice}

\begin{notice}{72}{362}{ICUR n.s. IV, 11758 · EDCS-37400852 · EDB6730}
\label{not:72}
\nf{Datation}{362 — directe : nom consulaire lu}
\nf{Genre}{Épitaphe ; ensembles A, B, C ; 99 lettres transmises.}
\nf{Références}{ICUR IV 11758 = ILCV 3904 = RAC-1925-24,1 = RAC-1981-223 = AE 1904, 33 = AE 1904, 122.}
\nf{Provenance}{Rattachement selon l’EDB : \emph{Coem. anonymum ad viam Appiam}. Conservation selon l’EDB : \emph{Coem. anonymum ad viam Appiam}.}
\nf{Support}{Plaque de marbre. Inscription gravée. Reproductions signalées par l’EDB : photographie PCAS ×2 (droits réservés ; non reproduites). Fiche EDB : Maria Cristina Grisanzio, première version 2015.}
\ntext{\lnum{1}D(is) M(anibus) s(acrum)\\
\lnum{2}Fl(avia) Balbilla vixit ann(os) XXVIII\\
\lnum{3}menses VI dies XII requievit in pace\\
\lnum{4}Mamertino et Ne⸢v⸣i⸢tt⸣a co(n)ss(ulibus) pr(a)ef(ecto) urb(i)\\
\lnum{5}Maximo V Kal(endas) Feb(ruarias)}
\nf{Témoins}{Export EDCS de 2020 : \emph{D(is) M(anibus) s(acrum) / Fl(avia) Balbilla vixit ann(os) XXVIII / menses VI dies XII requievit in pace / Mamertino et Neia co(n)ss(ulibus) pr(a)ef(ecto) urb(i) / Maximo V Kal(endas) Feb(ruarias)}. EDB (EDB6730) : texte concordant aux conventions de notation près (indice de similarité 0,941).}
\nf{Apparat}{⸢v⸣ : B gravé ; ⸢tt⸣ : DD gravé}
\nf{Traduction}{Consacré aux dieux Mânes. Flavia Balbilla vécut 28 ans 6 mois 12 jours ; elle a reposé en paix sous le consulat de Mamertinus et de Nevitta, Maximus étant préfet de la Ville, le 5e jour avant les calendes de février.}
\nf{Langue}{L’éditeur signale 2 graphies corrigées : B gravé pour v (bêtacisme) ; DD gravé pour tt.}
\nf{Chronologie}{Noms consulaires : Mamertinus (lu) et Nevitta (lu). Dates du jour : événement indéterminé : V Kalendas Februarias = 28 janvier 362 (lu). Âge déclaré : 28 a. 6 m. 12 j. (lu).}
\nf{Commentaire}{La formule in pace est conservée. Fonction ou métier mentionné : praefecto. Commentaire développé : dossier p.~\pageref{dos:72}.}
\nf{Bibliographie}{ICUR IV 11758 ; ILCV 3904 ; RAC-1925-24,1 ; RAC-1981-223 ; AE 1904, 33 ; AE 1904, 122.}
\nf{Statut}{Lecture secondaire (EDCS) concordante avec l’EDB : lecture stabilisée entre deux bases, non collationnée sur l’ICVR ni sur le monument. L’EDCS signale une photographie (non consultée).}
\end{notice}

\annee{363}

\begin{notice}{73}{363}{ICUR n.s. I, 1426 · EDCS-35200401 · EDB31189}
\label{not:73}
\nf{Datation}{363 — directe : nom consulaire lu}
\nf{Genre}{Épitaphe ; ensembles A, B, C ; 456 lettres transmises ; 18 \% restituées.}
\nf{Références}{CLE 662 = ICUR I 1426 = ILCV 3359.}
\nf{Provenance}{Rattachement selon l’EDB : \emph{Incertae originis} (le volume I des ICVR réunit les inscriptions dont la provenance est perdue). Conservation selon l’EDB : \emph{Musei Vaticani: Lapidario Cristiano}.}
\nf{Support}{Plaque de marbre. 87 × 151 cm (h. × l.). Inscription gravée. Texte métrique. Signes décrits par l’EDB : alpha et oméga, chrisme, croix et croix gammée. Reproductions signalées par l’EDB : dessin, ICVR I (de Rossi) (droits réservés ; non reproduites). Fiche EDB : Filippo Piazzolla, première version 2015.}
\ntext{\lnum{1}[Haec aete]rna domus in qua nunc ipsa secura quiescis\\
\lnum{2}[at tuus in caelo] spiritus a carne recedens\\
\lnum{3}[vivit cum] sanctis pro meritis et opera tanta\\
\lnum{4}[quaeque deu]m metuisti semper quiescis secura p(alma) e(t) l(aurus)\\
\lnum{5}[mandasti corp]us terrae pignusque receptus\\
\lnum{6}[- - - spl]endori cum lumine claro\\
\lnum{7}[- - -]s quae te semper deo dicasti\\
\lnum{8}[- - -]tas sed pro factis ad alta vocaris\\
\lnum{9}[- - -]ca Pacatiano cum multa debebit\\
\lnum{10}[Pac]actianus Eleuteriae matri suae locum fecit\\
\lnum{11}[quae vixit an]n(os) pl(us) m(inus) LXXV mens(es) II d(ies) X dep(osita) IIII Kal(endas) Iun(ias) dormit\\
\lnum{12}[in pace cons(ulatu)] d(omini) n(ostri) Cl(audi) Iuliani Aug(usti) IIII et Fl(avi) Sallusti con[- - -]\\
\lnum{13}[H]esperius quem nutrit i(n)scripsit}
\nf{Témoins}{EDB (EDB31189 ; divergence forte) : \emph{[haec tibi aet]e̲rna domus in qua nunc ipsa secura quiescis ((palma)) / [corpore namque tu]u̲s spiritus a carne redecens ((sol)) / [est sociatu]s̲ sanctis pro meritis et opera tanta ((monogramma Christi)) / [quaeque de]m̲etuisti semper quiescis secura ((palma et laurus)) / [dedisti corp]us terrae pignus que receptus ((palma)) / [--- spl]e̲ndori cum lumine claro ((crux gammata )) / [---]s̲ quae te semper deo dica⟨vi⟩sti ((alpha inversa?)) / [---]t̲as sed pro factis ad alta vocaris ((omega)) / [---]c̲a Pacatiano cum multa debebit ((Y?)) / [Pac]ạtianus Eleuteriae matri suae locum fecit / [que vixit a]n̲ṇ(is) pl(us) m(inus) LXXV mens(ibus) II d(iebus) X dep(osita) IIII kal(endas) iun(ias) dormit / [in pace cons(ultum)] d̲(omini) n(ostri) Cl(audi) Iuliani Aug(usti) IIII et Fl(avi) Sallusti con(sule) / [---?] H̲esperius quem nutrit iscripsit [texte tronqué dans la liste de résultats]}. L'EDB (EDB31189) et l'EDCS restituent différemment le début des vers 2 et 3 : \emph{[at tuus in caelo] spiritus … [vivit cum] sanctis} (EDCS) ; \emph{[corpore namque tu]us spiritus … [est sociatu]s sanctis} (EDB). Les deux restitutions sont conjecturales et orientent différemment la lecture : séjour céleste de l'esprit d'un côté, association aux saints de l'autre. Seuls \emph{spiritus a carne recedens} et \emph{sanctis pro meritis} sont lus. L'EDB décrit aussi des palmes, un soleil, un chrisme, une croix gammée, un alpha et un oméga.}
\nf{Traduction}{(vers) Voici la demeure éternelle où tu reposes désormais en sécurité ; mais ton esprit, se retirant de la chair, vit au ciel avec les saints, selon tes mérites et tes si grandes œuvres ; toi qui as toujours craint Dieu, tu reposes en sécurité, (avec) la palme et le laurier ; tu as confié ton corps à la terre et (ton) gage reçu […] ; […] à la splendeur avec une lumière claire ; […] toi qui t’es toujours consacrée à Dieu […] ; […] mais pour tes actes tu es appelée vers les hauteurs […]. Pacatianus a fait cet emplacement pour Eleuteria, sa mère, qui vécut environ 75 ans 2 mois 10 jours ; déposée le 4e jour avant les calendes de juin ; elle dort en paix ; sous le consulat de notre seigneur Claudius Julien Auguste pour la 4e fois et de Flavius Sallustius. Hesperius, qu’elle a nourri, a gravé (l’inscription).}
\nf{Langue}{18 abréviations développées par l’éditeur (dont p(alma), e(t), l(aurus), ann(os), pl(us), m(inus)).}
\nf{Chronologie}{Noms consulaires : Iulianus (lu) et Sallustius (lu). Dates du jour : déposition : IIII Kalendas Iunias = 29 mai 363 (lu). Âge déclaré : env. 75 a. 2 m. 10 j. (partiellement restitué).}
\nf{Commentaire}{La formule in pace est attestée seulement par restitution. Vocabulaire religieux : désignation de Dieu — deum (partiellement restitué), deo (lu) ; âme humaine (non doctrinal) — spiritus (lu) ; saints (communion des saints) — sanctis (lu) ; repos et sommeil de la mort — dormit (lu), quiescis (lu) ; allusion interprétative (niveau 5) — caelo (entièrement restitué), lumine (lu). Désignation divine : retenue dans le profil strict. Commentaire développé : dossier p.~\pageref{dos:73}.}
\nf{Bibliographie}{CLE 662 ; ICUR I 1426 ; ILCV 3359. Bibliographie complémentaire selon l’EDB : Ferrua 1964, 108 ; Lega 2000, 102 ; ILCV, 3359 ; Ferrua 1979b, 36, n. 1426 ; Trout 2013, 20, nt. 59.}
\nf{Statut}{Lecture secondaire (EDCS) ; divergence forte avec l’EDB : le texte ne doit pas être exploité sans contrôle de l’édition imprimée. L’EDCS signale une photographie (non consultée).}
\end{notice}

\begin{notice}{74}{363}{ICUR n.s. I, 966 · EDCS-38700473 · EDB29599}
\label{not:74}
\nf{Datation}{363 — directe : lecture partielle d’un nom consulaire}
\nf{Genre}{Épitaphe ; ensembles A, B, C ; 69 lettres transmises ; 33 \% restituées.}
\nf{Références}{ICUR I 966 = ILCV 4653.}
\nf{Provenance}{Rattachement selon l’EDB : \emph{Incertae originis} (le volume I des ICVR réunit les inscriptions dont la provenance est perdue).}
\nf{Support}{Support non précisé par l’EDB. Fiche EDB : Filippo Piazzolla, première version 2014.}
\ntext{\lnum{1}e]loque(n)s qu[i vixit]\\
\lnum{2}[- - -]III d(ies) XV in pa[ce]\\
\lnum{3}[Iuliano A]ug(usto) IIII et Sall[ustio]\\
\lnum{4}[- - - pa]rentibus(?) M[- - -]\\
\lnum{5}[- - -] quib(us) MF[}
\nf{Témoins}{EDB (EDB29599) : texte concordant aux conventions de notation près (indice de similarité 0,944).}
\nf{Traduction}{[…] éloquent, qui vécut […] 15 jours, en paix ; sous (le consulat de) Julien Auguste pour la 4e fois et de Sallustius ; […] à ses parents […].}
\nf{Langue}{5 abréviations développées par l’éditeur (dont eloque(n)s, d(ies), Aug(usto), parentibus(?), quib(us)).}
\nf{Chronologie}{Noms consulaires : Iulianus (restitué) et Sallustius (lu en partie). Âge déclaré : 15 j. (partiellement restitué).}
\nf{Commentaire}{La formule in pace est attestée seulement par restitution.}
\nf{Bibliographie}{ICUR I 966 ; ILCV 4653. Bibliographie complémentaire selon l’EDB : ILCV, 4653.}
\nf{Statut}{Lecture secondaire (EDCS) concordante avec l’EDB : lecture stabilisée entre deux bases, non collationnée sur l’ICVR ni sur le monument.}
\end{notice}

\begin{notice}{75}{363}{ICUR n.s. I, 1425 · EDCS-39000411 · EDB31188}
\label{not:75}
\nf{Datation}{363 — directe : nom consulaire lu}
\nf{Genre}{Épitaphe ; ensembles A, B, C ; 108 lettres transmises ; 8 \% restituées.}
\nf{Références}{ICUR I 1425 = ILCV 2603.}
\nf{Provenance}{Rattachement selon l’EDB : \emph{Incertae originis} (le volume I des ICVR réunit les inscriptions dont la provenance est perdue). Conservation selon l’EDB : \emph{Musei Vaticani: Lapidario Cristiano}.}
\nf{Support}{Plaque de marbre. 20 × 70 cm (h. × l.). Inscription gravée. Reproductions signalées par l’EDB : dessin, ICVR I (de Rossi) et reproduction, Perret 1851 (droits réservés ; non reproduites). Fiche EDB : Filippo Piazzolla, première version 2015.}
\ntext{\lnum{1}Refrigerio bene merenti in\\
\lnum{2}pace qui vixit ann(os) plus minu[s]\\
\lnum{3}n(umero) XL depos\{s\}itus III Idus Maias\\
\lnum{4}d(omino) n(ostro) Iuliano Aug(usto) IIII et Sallustio co[ns(ulibus)]}
\nf{Témoins}{EDB (EDB31188) : texte concordant aux conventions de notation près (indice de similarité 0,958).}
\nf{Apparat}{\{s\} : lettres superflues}
\nf{Traduction}{À Refrigerius, bien méritant, en paix, qui vécut environ 40 ans ; déposé le 3e jour avant les ides de mai, sous le consulat de notre seigneur Julien Auguste pour la 4e fois et de Sallustius.}
\nf{Langue}{Des lettres gravées sont jugées superflues par l’éditeur (entre accolades). 6 abréviations développées par l’éditeur (dont ann(os), n(umero), d(omino), n(ostro), Aug(usto), cons(ulibus)).}
\nf{Chronologie}{Noms consulaires : Iulianus (lu) et Sallustius (lu). Dates du jour : déposition : III Idus Maias = 13 mai 363 (lu). Âge déclaré : env. 40 a. (partiellement restitué).}
\nf{Commentaire}{La formule in pace est conservée (répartie sur deux lignes).}
\nf{Bibliographie}{ICUR I 1425 ; ILCV 2603. Bibliographie complémentaire selon l’EDB : Lega 2000, 262 ; Ferrua 1979b, 36, n. 1425 ; ILCV, 2603 adn.}
\nf{Statut}{Lecture secondaire (EDCS) concordante avec l’EDB : lecture stabilisée entre deux bases, non collationnée sur l’ICVR ni sur le monument. L’EDCS signale une photographie (non consultée).}
\end{notice}

\begin{notice}{76}{363}{ICUR n.s. I, 2082 · EDCS-39300306 · EDB32044}
\label{not:76}
\nf{Datation}{363 — directe : nom consulaire lu}
\nf{Genre}{Épitaphe ; ensembles A, B, C ; 78 lettres transmises ; 22 \% restituées.}
\nf{Références}{ICUR I 2082.}
\nf{Provenance}{Rattachement selon l’EDB : \emph{Incertae originis} (le volume I des ICVR réunit les inscriptions dont la provenance est perdue). Conservation selon l’EDB : \emph{Musei Vaticani: Galleria Lapidaria}.}
\nf{Support}{Plaque de marbre. 26 × 67 cm (h. × l.). Inscription gravée. Reproductions signalées par l’EDB : dessin, ICVR I (de Rossi) (droits réservés ; non reproduites). Fiche EDB : Filippo Piazzolla, première version 2015.}
\ntext{\lnum{1}]e bene merenti ⸢i⸣n pace\\
\lnum{2}[- - - vixit ann(os) - - -]VII m(enses) II dep(ositus) IIII Non(as) Febr(uarias)\\
\lnum{3}[Iuliano] Aug(usto) IIII et \{et\} Sallustio conss(ulibus)}
\nf{Témoins}{Export EDCS de 2020 : \emph{]e bene merenti n pace / [3 vixit ann(os) 3]VII m(enses) II dep(ositus) IIII Non(as) Febr(uarias) / [Iuliano] Aug(usto) IIII et \{et\} Sallustio conss(ulibus)}. EDB (EDB32044 ; écarts à examiner) : \emph{------ / [---]e benemerenti en pace / [vixit annis] VII m(ensibus) II dep(osit---) IIII non(as) febr(uarias) / [ Iuliano] Ạug(usto) IIII et \{et\} Sallustio cons(ulibu)s [texte tronqué dans la liste de résultats]}.}
\nf{Apparat}{⸢i⸣ : E gravé ; \{et\} : lettres superflues}
\nf{Traduction}{[…] bien méritant(e), en paix […] ; vécut […]7 ans 2 mois ; déposé le 4e jour avant les nones de février, sous le consulat de Julien Auguste pour la 4e fois et de Sallustius.}
\nf{Langue}{L’éditeur signale une graphie corrigée : E gravé pour i (confusion e/i). Des lettres gravées sont jugées superflues par l’éditeur (entre accolades). 7 abréviations développées par l’éditeur (dont ann(os), m(enses), dep(ositus), Non(as), Febr(uarias), Aug(usto)).}
\nf{Chronologie}{Noms consulaires : Iulianus (restitué) et Sallustius (lu). Dates du jour : déposition : IIII Nonas Februarias = 2 février 363 (lu). Âge déclaré : 2 m. (partiellement restitué).}
\nf{Commentaire}{La formule in pace est conservée.}
\nf{Bibliographie}{ICUR I 2082. Bibliographie complémentaire selon l’EDB : ILCV, 2603 ; Lega 2000.}
\nf{Statut}{Lecture secondaire (EDCS) ; écarts avec l’EDB à arbitrer sur l’ICVR. L’EDCS signale une photographie (non consultée).}
\end{notice}

\begin{notice}{77}{363}{ICUR n.s. I, 2083 · EDCS-39300307 · EDB32045}
\label{not:77}
\nf{Datation}{363 — directe : nom consulaire lu}
\nf{Genre}{Épitaphe ; ensembles A, B, C ; 69 lettres transmises ; 1 \% restituées.}
\nf{Références}{ICUR I 2083.}
\nf{Provenance}{Rattachement selon l’EDB : \emph{Incertae originis} (le volume I des ICVR réunit les inscriptions dont la provenance est perdue). Conservation selon l’EDB : \emph{Musei Vaticani: Galleria Lapidaria}.}
\nf{Support}{Plaque de marbre. 21 × 38 cm (h. × l.). Inscription gravée. Reproductions signalées par l’EDB : dessin, ICVR I (de Rossi) (droits réservés ; non reproduites). Fiche EDB : Filippo Piazzolla, première version 2015.}
\ntext{\lnum{1}]\\
\lnum{2}e Cornificiae b(ene) m(erenti) in pac[e]\\
\lnum{3}marito ann(os) XIIII m(enses) XI dep(ositus)\\
\lnum{4}Iuliano Aug(usto) IIII et Sallustio co(n)ss(ulibus)}
\nf{Témoins}{EDB (EDB32045) : texte concordant aux conventions de notation près (indice de similarité 0,94).}
\nf{Traduction}{[…] à Cornificia, bien méritante, en paix ; (a vécu avec son) mari 14 ans 11 mois ; déposé(e) sous le consulat de Julien Auguste pour la 4e fois et de Sallustius.}
\nf{Langue}{8 abréviations développées par l’éditeur (dont b(ene), m(erenti), ann(os), m(enses), dep(ositus), Aug(usto)).}
\nf{Chronologie}{Noms consulaires : Iulianus (lu) et Sallustius (lu).}
\nf{Commentaire}{La formule in pace est attestée seulement par restitution.}
\nf{Bibliographie}{ICUR I 2083. Bibliographie complémentaire selon l’EDB : Lega 2000.}
\nf{Statut}{Lecture secondaire (EDCS) concordante avec l’EDB : lecture stabilisée entre deux bases, non collationnée sur l’ICVR ni sur le monument. L’EDCS signale une photographie (non consultée).}
\end{notice}

\begin{notice}{78}{363}{ICUR n.s. I, 2806 · EDCS-11700572 · EDB32729}
\label{not:78}
\nf{Datation}{363 — directe : nom consulaire lu}
\nf{Genre}{Épitaphe ; ensembles A, B, C ; 94 lettres transmises ; 26 \% restituées.}
\nf{Références}{ICUR I 2806 = ILCV 4338f = ILMN-01, 00441.}
\nf{Provenance}{Rattachement selon l’EDB : \emph{Incertae originis} (le volume I des ICVR réunit les inscriptions dont la provenance est perdue). Conservation selon l’EDB : \emph{Napoli   ; Museo Archeologico Nazionale}.}
\nf{Support}{Plaque de marbre. 27,5 × 160 × 4 cm (h. × l. × ép.). Inscription gravée. Reproductions signalées par l’EDB : reproduction, Camodeca et alii 2000 (droits réservés ; non reproduites). Fiche EDB : Filippo Piazzolla, première version 2015.}
\ntext{\lnum{1}Laborum autrix cas[tita]tis socia hic\\
\lnum{2}Erenes in pace [quies]cit\\
\lnum{3}d(e)p(osita) d(ie) III Kal(endas) Oct(obres) Iuliano [Aug(usto) III et Sal]lustio co(n)s(ulibus)}
\nf{Témoins}{EDB (EDB32729) : texte concordant aux conventions de notation près (indice de similarité 0,949).}
\nf{Traduction}{Compagne de ses travaux, associée de sa chasteté, ici Irène repose en paix ; déposée le 3e jour avant les calendes d’octobre, sous le consulat de Julien Auguste pour la 3e (sic) fois et de Sallustius.}
\nf{Langue}{8 abréviations développées par l’éditeur (dont d(e)p, (osita), d(ie), Kal(endas), Oct(obres), Aug(usto)).}
\nf{Chronologie}{Noms consulaires : Iulianus (lu) et Sallustius (lu en partie). Dates du jour : déposition : die III Kalendas Octobres = 29 septembre 363 (lu).}
\nf{Commentaire}{La formule in pace est conservée. Vocabulaire religieux : repos et sommeil de la mort — quiescit (partiellement restitué).}
\nf{Bibliographie}{ICUR I 2806 ; ILCV 4338f ; ILMN-01, 00441. Bibliographie complémentaire selon l’EDB : Camodeca et alii 2000, 152, n. 441 con foto ; ILCV, 4338F.}
\nf{Statut}{Lecture secondaire (EDCS) concordante avec l’EDB : lecture stabilisée entre deux bases, non collationnée sur l’ICVR ni sur le monument. L’EDCS signale une photographie (non consultée).}
\end{notice}

\begin{notice}{79}{363}{ICUR n.s. I, 2807 · EDCS-11700573 · EDB32730}
\label{not:79}
\nf{Datation}{363 — directe : nom consulaire lu}
\nf{Genre}{Épitaphe ; ensembles A, B, C ; 62 lettres transmises ; 14 \% restituées.}
\nf{Références}{ICUR I 2807 = ILMN-01, 00442.}
\nf{Provenance}{Rattachement selon l’EDB : \emph{Incertae originis} (le volume I des ICVR réunit les inscriptions dont la provenance est perdue). Conservation selon l’EDB : \emph{Napoli   ; Museo Archeologico Nazionale}.}
\nf{Support}{Plaque de marbre. 29 × 70 × 3,5 cm (h. × l. × ép.). Inscription gravée. Reproductions signalées par l’EDB : reproduction, Camodeca et alii 2000 (droits réservés ; non reproduites). Fiche EDB : Filippo Piazzolla, première version 2015.}
\ntext{\lnum{1}defu]ncta q(uae) vixit an\\
\lnum{2}[nos - - -]nta quattuor de\\
\lnum{3}[posita - - -] Cl(audio) Iuliano et Sallustio\\
\lnum{4}co(n)ss(ulibus)}
\nf{Témoins}{EDB (EDB32730 ; écarts à examiner) : \emph{[--- defu]ncta q(uae) vixit an/[nis ---gi]nta quattuor de/[posita ---] Cl(audio) Iuliano et Sallustio / co(n)s(ulibu)s [texte tronqué dans la liste de résultats]}.}
\nf{Traduction}{[…] défunte, qui vécut […] quatre ans ; déposée […] sous le consulat de Claudius Julien et de Sallustius.}
\nf{Langue}{4 abréviations développées par l’éditeur (dont q(uae), Cl(audio), co(n)ss, (ulibus)).}
\nf{Chronologie}{Noms consulaires : Iulianus (lu) et Sallustius (lu).}
\nf{Bibliographie}{ICUR I 2807 ; ILMN-01, 00442. Bibliographie complémentaire selon l’EDB : Camodeca et alii 2000, 152, n. 442 con foto.}
\nf{Statut}{Lecture secondaire (EDCS) ; écarts avec l’EDB à arbitrer sur l’ICVR. L’EDCS signale une photographie (non consultée).}
\end{notice}

\begin{notice}{80}{363}{ICUR n.s. I, 3167 · EDCS-40300314 · EDB33118}
\label{not:80}
\nf{Datation}{363 — directe : nom consulaire lu}
\nf{Genre}{Épitaphe ; ensembles A, B, C ; 74 lettres transmises.}
\nf{Références}{ICUR I 3167 = ILCV 3137 (em) = Louvre 852.}
\nf{Provenance}{Rattachement selon l’EDB : \emph{Incertae originis} (le volume I des ICVR réunit les inscriptions dont la provenance est perdue). Conservation selon l’EDB : \emph{Parigi (F)   ; Museo Louvre}.}
\nf{Support}{Sarcophage de marbre. Inscription gravée. Reproductions signalées par l’EDB : dessin, ICVR I (de Rossi) (droits réservés ; non reproduites). Fiche EDB : Filippo Piazzolla, première version 2015.}
\ntext{\lnum{1}Theodora requ(i)evit\\
\lnum{2}in pace die pridie Non(as)\\
\lnum{3}Septembr⸢e⸣s d(ominis) n(ostris)\\
\lnum{4}Iuliano Aug(usto) IIII et\\
\lnum{5}Sal(l)ustio cons(ulibus)}
\nf{Témoins}{Export EDCS de 2020 : \emph{Theodora requ(i)evit / in pace die pridie Non(as) / Septembrs d(ominis) n(ostris) / Iuliano Aug(usto) IIII et / Sal(l)ustio cons(ulibus)}. EDB (EDB33118) : texte concordant aux conventions de notation près (indice de similarité 0,989).}
\nf{Apparat}{⸢e⸣ : I gravé}
\nf{Traduction}{Theodora a reposé en paix, la veille des nones de septembre, sous le consulat de nos seigneurs Julien Auguste pour la 4e fois et Sallustius.}
\nf{Langue}{L’éditeur signale une graphie corrigée : I gravé pour e (confusion e/i).}
\nf{Chronologie}{Noms consulaires : Iulianus (lu) et Sallustius (lu). Dates du jour : décès : die pridie Nonas Septembres (lu).}
\nf{Commentaire}{La formule in pace est conservée. Vocabulaire religieux : anthroponyme (niveau 6) — theodora (anthroponyme).}
\nf{Bibliographie}{ICUR I 3167 ; ILCV 3137 (em) ; Louvre 852. Bibliographie complémentaire selon l’EDB : ILCV, 3137 ; Baratte et alii 1985, n. 202 ; Ducroux 1975, 223, n. 852 ; Christern-Briesenick et alii 2003, 203, pl. 104, 3, n. 433.}
\nf{Statut}{Lecture secondaire (EDCS) concordante avec l’EDB : lecture stabilisée entre deux bases, non collationnée sur l’ICVR ni sur le monument. L’EDCS signale une photographie (non consultée).}
\end{notice}

\begin{notice}{81}{363}{ICUR n.s. IV, 9559 · EDCS-36101088 · EDB19258}
\label{not:81}
\nf{Datation}{363 — directe : nom consulaire lu}
\nf{Genre}{Épitaphe ; ensembles A, B, C ; 109 lettres transmises.}
\nf{Références}{ICUR IV 9559 = ILCV 1327.}
\nf{Provenance}{Rattachement selon l’EDB : \emph{Coem. Callisti pars inferior}.}
\nf{Support}{Plaque de marbre. Inscription gravée. Reproductions signalées par l’EDB : photographie PCAS ×2 (droits réservés ; non reproduites). Fiche EDB : Antonio Enrico Felle, première version 2003.}
\ntext{\lnum{1}Felix qui\\
\lnum{2}vixit annis XXI m(en)ses(!) duo\\
\lnum{3}dep(ositus) die VI Idus Decembres Iuli\\
\lnum{4}ano IIII et Sallustio conss(ulibus) in pac(e)\\
\lnum{5}filio dulcissimo Calliste mater}
\nf{Témoins}{EDB (EDB19258) : texte concordant aux conventions de notation près (indice de similarité 0,992).}
\nf{Apparat}{(!) : forme signalée par l’éditeur}
\nf{Traduction}{Felix, qui vécut 21 ans 2 mois ; déposé le 6e jour avant les ides de décembre, sous le consulat de Julien pour la 4e fois et de Sallustius, en paix ; à son fils très doux, Callistè, sa mère.}
\nf{Langue}{L’éditeur signale au moins une forme anormale (sic). 4 abréviations développées par l’éditeur (dont m(en)ses, dep(ositus), conss(ulibus), pac(e)).}
\nf{Chronologie}{Noms consulaires : Iulianus (lu) et Sallustius (lu). Dates du jour : déposition : die VI Idus Decembres = 8 décembre 363 (lu). Âge déclaré : 21 a. 2 m. 6 j. (lu).}
\nf{Commentaire}{La formule in pace est transmise sous forme abrégée, sans restitution de lacune ; son développement reste éditorial.}
\nf{Bibliographie}{ICUR IV 9559 ; ILCV 1327.}
\nf{Statut}{Lecture secondaire (EDCS) concordante avec l’EDB : lecture stabilisée entre deux bases, non collationnée sur l’ICVR ni sur le monument. L’EDCS signale une photographie (non consultée).}
\end{notice}

\begin{notice}{82}{363}{ICUR n.s. IV, 11759 · EDCS-37400853 · EDB2091}
\label{not:82}
\nf{Datation}{363 — directe : lecture partielle d’un nom consulaire}
\nf{Genre}{Épitaphe ; ensembles A, B, C ; 48 lettres transmises ; 62 \% restituées.}
\nf{Références}{ICUR IV 11759.}
\nf{Provenance}{Rattachement selon l’EDB : \emph{Coem. anonymum ad viam Appiam}. Conservation selon l’EDB : \emph{Coem. anonymum ad viam Appiam}.}
\nf{Support}{Plaque de marbre. Inscription gravée. Fiche EDB : Maria Cristina Grisanzio, première version 2003.}
\ntext{\lnum{1}]MINIS[- - -]\\
\lnum{2}[- - - Iu]liano II[II et Sallustio cons(ulibus) - - -]\\
\lnum{3}[- - - e]t cum ma[rito}
\nf{Témoins}{EDB (EDB2091 ; écarts à examiner) : \emph{[---]MINIS[---] / [--- Iu]liano II[II et Sallustio] / [--- e]t cum ma[rito ---] [texte tronqué dans la liste de résultats]}.}
\nf{Traduction}{[…] sous le consulat de Julien pour la 4e fois et de Sallustius […] ; […] et avec son mari.}
\nf{Langue}{1 abréviation développée par l’éditeur (dont cons(ulibus)).}
\nf{Chronologie}{Noms consulaires : Iulianus (lu en partie) et Sallustius (restitué).}
\nf{Bibliographie}{ICUR IV 11759.}
\nf{Statut}{Lecture secondaire (EDCS) ; écarts avec l’EDB à arbitrer sur l’ICVR. Plus de la moitié des lettres sont restituées. L’EDCS signale une photographie (non consultée).}
\end{notice}

\annee{364}

\begin{notice}{83}{364}{EDCS-59400075}
\label{not:83}
\nf{Datation}{364 — directe : nom consulaire lu}
\nf{Genre}{Épitaphe ; ensembles A, B, C ; 107 lettres transmises ; 21 \% restituées.}
\nf{Références}{AE 2010, 207.}
\nf{Provenance}{Rome (indication de l’EDCS) ; aucun rattachement cimétérial transmis.}
\nf{Support}{Non documenté dans les sources consultées.}
\ntext{\lnum{1}Carissim(a)e ac dulcissim(a)e coni[ugi - - -]\\
\lnum{2}quae vix\{c\}(it) ann(os) p(lus) m(inus) XIIII et fecit [mecum]\\
\lnum{3}m(enses) VI d(ies) XXII dep(osita) in p(ace) III Idus No⸢v⸣(embres) [Ioviano Aug(usto)]\\
\lnum{4}et Varroniano (!)}
\nf{Témoins}{Export EDCS de 2020 : \emph{Carissim(a)e ac dulcissim(a)e coni[ugi 3] / quae vix\{c\}(it) ann(os) p(lus) m(inus) XIIII et fecit [mecum] / m(enses) VI d(ies) XXII dep(osita) in p(ace) III Idus No(embres) [Ioviano Aug(usto)] / et Varroniano (!)}. Aucune fiche EDB appariée (pas de numéro ICUR ; recherche par le texte infructueuse).}
\nf{Apparat}{⸢v⸣ : B gravé ; \{c\} : lettres superflues ; (!) : forme signalée par l’éditeur}
\nf{Traduction}{À son épouse très chère et très douce […] qui vécut environ 14 ans et a vécu avec moi 6 mois 22 jours ; déposée en paix le 3e jour avant les ides de novembre, sous (le consulat de) Jovien Auguste et Varronianus.}
\nf{Langue}{L’éditeur signale une graphie corrigée : B gravé pour v (bêtacisme). Des lettres gravées sont jugées superflues par l’éditeur (entre accolades). L’éditeur signale au moins une forme anormale (sic). 12 abréviations développées par l’éditeur (dont Carissim(a)e, dulcissim(a)e, (it), ann(os), p(lus), m(inus)).}
\nf{Chronologie}{Noms consulaires : Iovianus (restitué) et Varronianus (lu). Dates du jour : déposition : III Idus Novembres = 11 novembre 364 (lu). Âge déclaré : env. 14 a. (partiellement restitué).}
\nf{Commentaire}{La formule in pace est transmise sous forme abrégée, sans restitution de lacune ; son développement reste éditorial.}
\nf{Bibliographie}{AE 2010, 207.}
\nf{Statut}{Lecture secondaire (EDCS) sans second témoin numérique.}
\end{notice}

\begin{notice}{84}{364}{ICUR n.s. I, 2084 · EDCS-39300308 · EDB32048}
\label{not:84}
\nf{Datation}{364 — formule transmise par l’EDB (lue)}
\nf{Genre}{Épitaphe ; ensembles A, B, C ; 84 lettres transmises ; 34 \% restituées.}
\nf{Références}{ICUR I 2084.}
\nf{Provenance}{Rattachement selon l’EDB : \emph{Incertae originis} (le volume I des ICVR réunit les inscriptions dont la provenance est perdue).}
\nf{Support}{Support non précisé par l’EDB. Reproductions signalées par l’EDB : dessin, ICVR I (de Rossi) (droits réservés ; non reproduites). Fiche EDB : Filippo Piazzolla, première version 2015.}
\ntext{\lnum{1}[Bonif]atia filia Ausustin(a)e\\
\lnum{2}[dulcissim]ae qu(a)e vixit\\
\lnum{3}[annos - - -] mens⸢e⸣s IIII dies X\\
\lnum{4}[dep(osita)] VII [Ka]l(endas) Octobres\\
\lnum{5}[}
\nf{Témoins}{Export EDCS de 2020 : \emph{[Bonif]atia filia Ausustin(a)e / [dulcissim]ae qu(a)e vixit / [annos 3] menss IIII dies X / [dep(osita)] VII [Ka]l(endas) Octobres / [}. EDB (EDB32048 ; écarts à examiner) : \emph{[Bonif]atia filiae Asustine / [dulcissim?]ae que vixit / [annos ---] m̲e̲nsis IIII dies X / [deposita] V̲I̲I̲ [kal(endas)] octobres / [di]v̲o Ioviano et Fl(avio) Barronian`o c(on)s(ule)´ [texte tronqué dans la liste de résultats]}.}
\nf{Apparat}{⸢e⸣ : I gravé}
\nf{Traduction}{Bonifatia, fille d’Augustina, très douce, qui vécut […] ans 4 mois 10 jours ; déposée le 7e jour avant les calendes d’octobre.}
\nf{Langue}{L’éditeur signale une graphie corrigée : I gravé pour e (confusion e/i).}
\nf{Chronologie}{Le texte EDCS s'arrête avant la date consulaire. La fiche EDB (EDB32048) transmet \emph{[di]v̲o Ioviano et Fl(avio) Barronian`o c(on)s(ule)´} : les lettres soulignées ont été vues par un éditeur antérieur, et la formule nomme Jovien et Varronianus (364). L'année 364 est désormais fondée sur une formule consulaire lue, transmise par l'EDB. Dates du jour : déposition : VII Kalendas Octobres = 25 septembre 364 (partiellement restitué). Âge déclaré : 4 m. 10 j. (partiellement restitué).}
\nf{Bibliographie}{ICUR I 2084. Bibliographie complémentaire selon l’EDB : ILCV, 2941A adn.}
\nf{Statut}{Lecture secondaire (EDCS) ; écarts avec l’EDB à arbitrer sur l’ICVR. L’EDCS signale une photographie (non consultée).}
\end{notice}

\begin{notice}{85}{364}{ICUR n.s. I, 2770 · EDCS-39900319 · EDB32699}
\label{not:85}
\nf{Datation}{364 — directe : nom consulaire lu}
\nf{Genre}{Épitaphe ; ensembles A, B, C ; 136 lettres transmises ; 24 \% restituées.}
\nf{Références}{ICUR I 2770 = ILCV 3546.}
\nf{Provenance}{Rattachement selon l’EDB : \emph{Incertae originis} (le volume I des ICVR réunit les inscriptions dont la provenance est perdue). Conservation selon l’EDB : \emph{Fossombrone (PU)   ; Lapidario Diocesano}.}
\nf{Support}{Plaque de marbre. 50 × 60 cm (h. × l.). Inscription gravée. Reproductions signalées par l’EDB : reproduction (droits réservés ; non reproduites). Fiche EDB : Filippo Piazzolla, première version 2018.}
\ntext{\lnum{1}[In hoc t]umulo iacet Innocen\\
\lnum{2}[tia(?) - - -] omnia domino\\
\lnum{3}[- - - n]umquam super⸢b⸣ia\\
\lnum{4}[- - -]nio vel amicitiae ab OMI\\
\lnum{5}[- - - vixit a]⸢n⸣nos XXIII menses IIII die\\
\lnum{6}[s - - - depo]s(ita) III Kal(endas) April(es) Ioviano [et]\\
\lnum{7}[Varroniano}
\nf{Témoins}{Export EDCS de 2020 : \emph{[In hoc t]umulo iacet Innocen/[tia(?) 3] omnia domino / [3 n]umquam superia / [3]nio vel amicitiae ab OMI/[3 vixit a]nos XXIII menses IIII die/[s 3 depo]s(ita) III Kal(endas) April(es) Ioviano [et] / [Varroniano}. EDB (EDB32699 ; écarts à examiner) : \emph{[--- in hoc t]umulo iace̲t̲ in̲n̲o̲c̲e̲n̲/[---] omnia do̲m̲i̲̲n̲o̲ [---] / [--- n]umquam̲ s̲u̲p̲e̲r̲v̲i̲a̲ [---] / [---]ṇio vel amic̲i̲t̲ịae ab om̲i̲[---] / [--- vixit a]⌜n⌝nos XXIII menses III̲I̲ die [---] / [--- depo]s(it--) III kal(endas) ap̲ril(es) Ioviano [et] / [Varroniano consulibus] [texte tronqué dans la liste de résultats]}.}
\nf{Apparat}{⸢b⸣ : V gravé ; ⸢n⸣ : V gravé}
\nf{Traduction}{Dans ce tombeau repose Innocentia […] toutes choses au Seigneur […] jamais d’orgueil […] ou d’amitié […] ; elle vécut 23 ans 4 mois […] jours ; déposée le 3e jour avant les calendes d’avril, sous (le consulat de) Jovien et Varronianus.}
\nf{Langue}{L’éditeur signale 2 graphies corrigées : V gravé pour b (bêtacisme) ; V gravé pour n.}
\nf{Chronologie}{Noms consulaires : Iovianus (lu) et Varronianus (restitué). Dates du jour : déposition : III Kalendas Apriles = 30 mars 364 (lu). Âge déclaré : 23 a. 4 m. (partiellement restitué).}
\nf{Commentaire}{Vocabulaire religieux : désignation de Dieu — domino (lu). Désignation divine : retenue dans le profil strict.}
\nf{Bibliographie}{ICUR I 2770 ; ILCV 3546. Bibliographie complémentaire selon l’EDB : ILCV, 3546.}
\nf{Statut}{Lecture secondaire (EDCS) ; écarts avec l’EDB à arbitrer sur l’ICVR. L’EDCS signale une photographie (non consultée).}
\end{notice}

\begin{notice}{86}{364}{ICUR n.s. II, 5789 · EDCS-38302491 · EDB22178}
\label{not:86}
\nf{Datation}{364 — directe : lecture partielle d’un nom consulaire}
\nf{Genre}{Épitaphe ; ensembles A, B, C ; 39 lettres transmises ; 56 \% restituées.}
\nf{Références}{ICUR II 5789.}
\nf{Provenance}{Rattachement selon l’EDB : \emph{Area Lucinae}. Conservation selon l’EDB : \emph{s. Paolo f.l.m.}.}
\nf{Support}{Plaque de marbre. 29 × 33 cm (h. × l.). Inscription gravée. Fiche EDB : Anita Rocco, première version 2006.}
\ntext{\lnum{1}]AMRE[- - -]\\
\lnum{2}[- - -]iunus I[- - -]\\
\lnum{3}[- - - Io]viano et [Varroniano cons(ulibus)]}
\nf{Témoins}{EDB (EDB22178 ; écarts à examiner) : \emph{------ / [---]am re[---] / [---]iunus I[---] / [--- Iov]iano eṭ [Varroniano co(n)ss(ulibus)]}.}
\nf{Traduction}{[…] sous le consulat de Jovien et de Varronianus.}
\nf{Langue}{1 abréviation développée par l’éditeur (dont cons(ulibus)).}
\nf{Chronologie}{Noms consulaires : Iovianus (lu en partie) et Varronianus (restitué).}
\nf{Bibliographie}{ICUR II 5789.}
\nf{Statut}{Lecture secondaire (EDCS) ; écarts avec l’EDB à arbitrer sur l’ICVR. Plus de la moitié des lettres sont restituées. L’EDCS signale une photographie (non consultée).}
\end{notice}

\begin{notice}{87}{364}{ICUR n.s. II, 6024 · EDCS-34200650 · EDB17464}
\label{not:87}
\nf{Datation}{364 — directe : nom consulaire lu}
\nf{Genre}{Épitaphe ; ensembles A, B, C ; 31 lettres transmises ; 71 \% restituées.}
\nf{Références}{ICUR II 6024.}
\nf{Provenance}{Rattachement selon l’EDB : \emph{Coem. Commodillae}.}
\nf{Support}{Plaque de marbre. Inscription gravée. Fiche EDB : Carlo Carletti, première version 2003.}
\ntext{\lnum{1}DD(ominis) nn(ostris) Iovi(an)o [et Varroniano cons(ulibu)s]}
\nf{Témoins}{EDB (EDB17464 ; écarts à examiner) : \emph{[---] dd(ominis) nn(ostris) Iovio [et Varroniano cons(ulibu)s] [texte tronqué dans la liste de résultats]}.}
\nf{Traduction}{Sous le consulat de nos seigneurs Jovien et Varronianus.}
\nf{Langue}{4 abréviations développées par l’éditeur (dont DD(ominis), nn(ostris), Iovi(an)o, cons(ulibu)s).}
\nf{Chronologie}{Noms consulaires : Iovianus (lu) et Varronianus (restitué).}
\nf{Bibliographie}{ICUR II 6024.}
\nf{Statut}{Lecture secondaire (EDCS) ; écarts avec l’EDB à arbitrer sur l’ICVR. Plus de la moitié des lettres sont restituées.}
\end{notice}

\begin{notice}{88}{364}{ICUR n.s. III, 7381 · EDCS-34900709 · EDB24129}
\label{not:88}
\nf{Datation}{364 — directe : nom consulaire lu}
\nf{Genre}{Épitaphe ; ensembles A, B, C ; 69 lettres transmises ; 67 \% restituées.}
\nf{Références}{ICUR III 7381.}
\nf{Provenance}{Rattachement selon l’EDB : \emph{Coem. Domitillae pars superior}. Conservation selon l’EDB : \emph{Coem. Domitillae pars superior}.}
\nf{Support}{Plaque de marbre. 23 × 26 × 1,3 cm (h. × l. × ép.). Inscription gravée. Fiche EDB : Carlo Carletti / Antonio Enrico Felle, première version 2003.}
\ntext{\lnum{1}- - -] Iobian[i Augusti]\\
\lnum{2}[et Fl(avi) Varronian]i nob(ili) p[ueri - - -]\\
\lnum{3}[qui vixit ann(os)] LII (e)t men(ses) [- - -]\\
\lnum{4}[et die]s VIII}
\nf{Témoins}{EDB (EDB24129 ; écarts à examiner) : \emph{[cons(ulatu) d(omini) n(ostri)?] Iobian[i Aug(usti)] / [et Fl(avi) Varronian]i nob(ilissimi) p[ueri] / [qui vixit ann(os)] LII ⟨e⟩t men(ses) [---] / [et die]s VIII}.}
\nf{Traduction}{[…] (sous le consulat) de Jovien Auguste et de Flavius Varronianus, noble enfant […], qui vécut 52 ans et […] mois et 8 jours.}
\nf{Langue}{5 abréviations développées par l’éditeur (dont Fl(avi), nob(ili), ann(os), (e)t, men(ses)).}
\nf{Chronologie}{Noms consulaires : Varronianus (restitué) et Iovianus (lu). Âge déclaré : 52 a. 8 j. (partiellement restitué).}
\nf{Bibliographie}{ICUR III 7381. Bibliographie complémentaire selon l’EDB : CLRE, 641.}
\nf{Statut}{Lecture secondaire (EDCS) ; écarts avec l’EDB à arbitrer sur l’ICVR. Plus de la moitié des lettres sont restituées. L’EDCS signale une photographie (non consultée).}
\end{notice}

\begin{notice}{89}{364}{ICUR n.s. IV, 9560 · EDCS-36101089 · EDB19824}
\label{not:89}
\nf{Datation}{364 — directe : nom consulaire lu}
\nf{Genre}{Épitaphe ; ensembles A, B, C ; 68 lettres transmises ; 52 \% restituées.}
\nf{Références}{ICUR IV 9560.}
\nf{Provenance}{Rattachement selon l’EDB : \emph{Coem. Callisti pars inferior}.}
\nf{Support}{Plaque de marbre. Inscription gravée. Fiche EDB : Antonio Enrico Felle, première version 2003.}
\ntext{\lnum{1}D(omino) n(ostro) Io⸢v⸣ian[o Aug(usto) pri]\\
\lnum{2}mum et Ba[rroniano]\\
\lnum{3}[c]onss(ulibus) XI Kal(endas) M[artias d(e)p(ositio)]\\
\lnum{4}Laurenti [}
\nf{Témoins}{Export EDCS de 2020 : \emph{D(omino) n(ostro) Ioian[o Aug(usto) pri]/mum et Ba[rroniano] / [c]onss(ulibus) XI Kal(endas) M[artias d(e)p(ositio)] / Laurenti [}. EDB (EDB19824) : texte concordant aux conventions de notation près (indice de similarité 0,951).}
\nf{Apparat}{⸢v⸣ : B gravé}
\nf{Traduction}{Sous le consulat de notre seigneur Jovien Auguste pour la première fois et de Varronianus, le 11e jour avant les calendes de mars, déposition de Laurentius.}
\nf{Langue}{L’éditeur signale une graphie corrigée : B gravé pour v (bêtacisme).}
\nf{Chronologie}{Noms consulaires : Iovianus (lu) et Varronianus (restitué). Dates du jour : événement indéterminé : XI Kalendas Martias = 19 février 364 (partiellement restitué).}
\nf{Bibliographie}{ICUR IV 9560. Bibliographie complémentaire selon l’EDB : CLRE, 641.}
\nf{Statut}{Lecture secondaire (EDCS) concordante avec l’EDB : lecture stabilisée entre deux bases, non collationnée sur l’ICVR ni sur le monument. Plus de la moitié des lettres sont restituées. L’EDCS signale une photographie (non consultée).}
\end{notice}

\begin{notice}{90}{364}{ICUR n.s. IV, 11096 · EDCS-37200378 · EDB39555}
\label{not:90}
\nf{Datation}{364 — directe : nom consulaire lu}
\nf{Genre}{Épitaphe ; ensembles A, B, C ; 61 lettres transmises ; 12 \% restituées.}
\nf{Références}{ICUR IV 11096 = ILCV +3099.}
\nf{Provenance}{Rattachement selon l’EDB : \emph{Coem. Callisti area subdialis et incerti loci}. Conservation selon l’EDB : \emph{Coem. Callisti area subdialis et incerti loci}.}
\nf{Support}{Sarcophage de marbre. 18 × 28 × 6 cm (h. × l. × ép.). Inscription gravée. Reproductions signalées par l’EDB : photographie PCAS (droits réservés ; non reproduites). Fiche EDB : Marida Pierno, première version 2015.}
\ntext{\lnum{1}]\\
\lnum{2}qui vixit ann(os)\\
\lnum{3}XXXV m(enses) V dep(ositus)\\
\lnum{4}III Kal(endas) Iul(ias) Io⸢v⸣i\\
\lnum{5}[a]no et ⸢V⸣arronian[o]\\
\lnum{6}[ q]uescit in [pace]}
\nf{Témoins}{Export EDCS de 2020 : \emph{] / qui vixit ann(os) / XXXV m(enses) V dep(ositus) / III Kal(endas) Iul(ias) Ioi/[a]no et arronian[o] / [ q]uescit in [pace]}. EDB (EDB39555) : texte concordant aux conventions de notation près (indice de similarité 0,951).}
\nf{Apparat}{⸢v⸣ : B gravé ; ⸢V⸣ : B gravé}
\nf{Traduction}{[…] qui vécut 35 ans 5 mois ; déposé le 3e jour avant les calendes de juillet, sous (le consulat de) Jovien et Varronianus ; il repose en paix.}
\nf{Langue}{L’éditeur signale 2 graphies corrigées : B gravé pour v (bêtacisme) ; B gravé pour V (bêtacisme).}
\nf{Chronologie}{Noms consulaires : Iovianus (lu en partie) et Varronianus (lu). Dates du jour : déposition : III Kalendas Iulias = 29 juin 364 (lu). Âge déclaré : 35 a. 5 m. (lu).}
\nf{Commentaire}{La formule in pace est attestée seulement par restitution.}
\nf{Bibliographie}{ICUR IV 11096 ; ILCV +3099. Bibliographie complémentaire selon l’EDB : ILCV, 3099 adn.}
\nf{Statut}{Lecture secondaire (EDCS) concordante avec l’EDB : lecture stabilisée entre deux bases, non collationnée sur l’ICVR ni sur le monument. L’EDCS signale une photographie (non consultée).}
\end{notice}

\begin{notice}{91}{364}{ICUR n.s. IV, 11760 · EDCS-37400854 · EDB2092}
\label{not:91}
\nf{Datation}{364 — directe : nom consulaire lu}
\nf{Genre}{Épitaphe ; ensembles A, B, C ; 55 lettres transmises ; 16 \% restituées.}
\nf{Références}{ICUR IV 11760 = ILCV 4411 (em).}
\nf{Provenance}{Rattachement selon l’EDB : \emph{Coem. anonymum ad viam Appiam}. Conservation selon l’EDB : \emph{Coem. anonymum ad viam Appiam}.}
\nf{Support}{Plaque de marbre. Inscription gravée. Fiche EDB : Maria Cristina Grisanzio, première version 2003.}
\ntext{\lnum{1}[- - -]illa qu(a)e vix(it) ann(os) L et m(en)s(es) VIII d(ies) VIII\\
\lnum{2}[dep(osita) - - - K]al(endas) Sep(tembres) di⸢v⸣o Io⸢v⸣iano Aug(usto)}
\nf{Témoins}{Export EDCS de 2020 : \emph{[3]illa qu(a)e vix(it) ann(os) L et m(en)s(es) VIII d(ies) VIII / [dep(osita) 3 K]al(endas) Sep(tembres) dio Ioiano Aug(usto)}. EDB (EDB2092) : texte concordant aux conventions de notation près (indice de similarité 0,948).}
\nf{Apparat}{⸢v⸣ : B gravé ; ⸢v⸣ : B gravé}
\nf{Traduction}{[…]illa, qui vécut 50 ans et 8 mois 8 jours ; déposée […] calendes de septembre, (sous le consulat du) divin Jovien Auguste.}
\nf{Langue}{L’éditeur signale 2 graphies corrigées : B gravé pour v (bêtacisme) ×2.}
\nf{Chronologie}{Noms consulaires : Iovianus (lu). Dates du jour : déposition : … Kalendas Septembres (partiellement restitué). Âge déclaré : 50 a. 8 m. 8 j. (lu).}
\nf{Bibliographie}{ICUR IV 11760 ; ILCV 4411 (em).}
\nf{Statut}{Lecture secondaire (EDCS) concordante avec l’EDB : lecture stabilisée entre deux bases, non collationnée sur l’ICVR ni sur le monument.}
\end{notice}

\annee{365}

\begin{notice}{92}{365}{ICUR n.s. I, 1189 · EDCS-38700730 · EDB30135}
\label{not:92}
\nf{Datation}{365 — directe : nom consulaire lu}
\nf{Genre}{Épitaphe ; ensembles A, B, C ; 95 lettres transmises ; 25 \% restituées.}
\nf{Références}{ICUR I 1189.}
\nf{Provenance}{Rattachement selon l’EDB : \emph{Incertae originis} (le volume I des ICVR réunit les inscriptions dont la provenance est perdue).}
\nf{Support}{Support non précisé par l’EDB. Reproductions signalées par l’EDB : dessin, ICVR I (de Rossi) (droits réservés ; non reproduites). Fiche EDB : Filippo Piazzolla, première version 2014.}
\ntext{\lnum{1}[Bon(a)e memo]riae filio dulcissimo Laurentio\\
\lnum{2}[qui vixit an]n(os) XVI m(enses) VI d(ies) IIII decessit\\
\lnum{3}[- - - I]dus Aug(ustas) Valentiniano et\\
\lnum{4}[- - - Vale]nte}
\nf{Témoins}{EDB (EDB30135) : texte concordant aux conventions de notation près (indice de similarité 0,921).}
\nf{Traduction}{De bonne mémoire, à son fils très doux Laurentius, qui vécut 16 ans 6 mois 4 jours ; il mourut […] ides d’août, sous (le consulat de) Valentinien et Valens.}
\nf{Langue}{5 abréviations développées par l’éditeur (dont Bon(a)e, ann(os), m(enses), d(ies), Aug(ustas)).}
\nf{Chronologie}{Noms consulaires : Valentinianus (lu) et Valens (lu). Dates du jour : décès : … Idus Augustas (partiellement restitué). Âge déclaré : 16 a. 6 m. 4 j. (partiellement restitué).}
\nf{Bibliographie}{ICUR I 1189. Bibliographie complémentaire selon l’EDB : ILCV, 2807A adn.}
\nf{Statut}{Lecture secondaire (EDCS) concordante avec l’EDB : lecture stabilisée entre deux bases, non collationnée sur l’ICVR ni sur le monument. L’EDCS signale une photographie (non consultée).}
\end{notice}

\begin{notice}{93}{365}{ICUR n.s. I, 2085 · EDCS-39300309 · EDB32049}
\label{not:93}
\nf{Datation}{365 — directe : lecture partielle d’un nom consulaire}
\nf{Genre}{Épitaphe ; ensembles A, B, C ; 111 lettres transmises ; 32 \% restituées.}
\nf{Références}{ICUR I 2085 = ILCV +94a.}
\nf{Provenance}{Rattachement selon l’EDB : \emph{Incertae originis} (le volume I des ICVR réunit les inscriptions dont la provenance est perdue). Conservation selon l’EDB : \emph{Musei Vaticani: Galleria Lapidaria}.}
\nf{Support}{Plaque de marbre. 30 × 45 cm (h. × l.). Inscription gravée. Reproductions signalées par l’EDB : dessin, ICVR I (de Rossi) (droits réservés ; non reproduites). Fiche EDB : Filippo Piazzolla, première version 2013.}
\ntext{\lnum{1}[P]re(sbyteri ?) Urbi qui vix[it - - -]\\
\lnum{2}[- - -]rini mir(a)e sapie[ntiae - - -]\\
\lnum{3}[- - -]V obsecutus est ut [- - -]\\
\lnum{4}[i]nfelices parent[es - - -]\\
\lnum{5}[in] pace depositu[s - - - Valentiniano]\\
\lnum{6}[et V]alente con[s(ulibus)}
\nf{Témoins}{EDB (EDB32049) : texte concordant aux conventions de notation près (indice de similarité 0,936).}
\nf{Traduction}{[P]re(sbyteri ?) Urbi : titre et nom incertains ; ne pas traduire « prêtre Urbicus » comme une identification acquise. […] qui vécut […] ; […] d’une admirable sagesse […] ; […] obéit (?) […] ; parents malheureux […] ; [en] paix, déposé […].}
\nf{Langue}{3 abréviations développées par l’éditeur (dont Pre(sbyteri ?), mir(a)e, cons(ulibus)).}
\nf{Chronologie}{Noms consulaires : Valentinianus (restitué) et Valens (lu en partie).}
\nf{Commentaire}{In pace dépend d’une restitution. Pre(sbyteri ?) demeure un développement incertain ; le nom Urbi et sa construction ne sont pas établis. Aucune attestation certaine d’un prêtre « Urbicus » n’en est déduite.}
\nf{Bibliographie}{ICUR I 2085 ; ILCV +94a. Bibliographie complémentaire selon l’EDB : ILCV, 94A adn., 1 ; Ferrua 1981d, 9, ad n. 94A.}
\nf{Statut}{Lecture secondaire (EDCS) concordante avec l’EDB : lecture stabilisée entre deux bases, non collationnée sur l’ICVR ni sur le monument. L’EDCS signale une photographie (non consultée).}
\end{notice}

\begin{notice}{94}{365}{ICUR n.s. I, 2695 · EDCS-39900259 · EDB32628}
\label{not:94}
\nf{Datation}{365 — directe : nom consulaire lu}
\nf{Genre}{Épitaphe ; ensembles A, B, C ; 98 lettres transmises.}
\nf{Références}{ICUR I 2695 = ILCV 1327a (em) = Bagno 51.}
\nf{Provenance}{Rattachement selon l’EDB : \emph{Incertae originis} (le volume I des ICVR réunit les inscriptions dont la provenance est perdue). Conservation selon l’EDB : \emph{Civitella di Romagna (FC)   ; Cusercoli}.}
\nf{Support}{Plaque de marbre. 13,8 × 35,5 × 2,5 cm (h. × l. × ép.). Inscription gravée. Signes décrits par l’EDB : chrisme. Reproductions signalées par l’EDB : dessin, ICVR I (de Rossi) et reproduction, Geraci 1975 (droits réservés ; non reproduites). Fiche EDB : Filippo Piazzolla, première version 2015.}
\ntext{\lnum{1}Dominis nostris Valentiniano et\\
\lnum{2}Valente conss(ulibus) V Kal(endas) Iunias Gaianus\\
\lnum{3}in pace qui vixit ann⸢o⸣s VIIII mens⸢e⸣s\\
\lnum{4}III dies XXVI}
\nf{Témoins}{Export EDCS de 2020 : \emph{Dominis nostris Valentiniano et / Valente conss(ulibus) V Kal(endas) Iunias Gaianus / in pace qui vixit anns VIIII menss / III dies XXVI}. EDB (EDB32628) : texte concordant aux conventions de notation près (indice de similarité 0,939).}
\nf{Apparat}{⸢o⸣ : V gravé ; ⸢e⸣ : I gravé}
\nf{Traduction}{Sous le consulat de nos seigneurs Valentinien et Valens, le 5e jour avant les calendes de juin, Gaianus, en paix, qui vécut 9 ans 3 mois 26 jours.}
\nf{Langue}{L’éditeur signale 2 graphies corrigées : V gravé pour o (confusion o/u) ; I gravé pour e (confusion e/i).}
\nf{Chronologie}{Noms consulaires : Valentinianus (lu) et Valens (lu). Dates du jour : événement indéterminé : V Kalendas Iunias = 28 mai 365 (lu). Âge déclaré : 9 a. 3 m. 26 j. (lu).}
\nf{Commentaire}{La formule in pace est conservée.}
\nf{Bibliographie}{ICUR I 2695 ; ILCV 1327a (em) ; Bagno 51. Bibliographie complémentaire selon l’EDB : ILCV, 1327A ; Ferrua 1979b, 77 ; Geraci 1975, 99-101, n. 51 con foto.}
\nf{Statut}{Lecture secondaire (EDCS) concordante avec l’EDB : lecture stabilisée entre deux bases, non collationnée sur l’ICVR ni sur le monument. L’EDCS signale une photographie (non consultée).}
\end{notice}

\begin{notice}{95}{365}{ICUR n.s. II, 6025 · EDCS-34200651 · EDB17465}
\label{not:95}
\nf{Datation}{365 — directe : nom consulaire lu}
\nf{Genre}{Épitaphe ; ensembles A, B, C ; 75 lettres transmises ; 23 \% restituées.}
\nf{Références}{ICUR II 6025 = ILCV +2926.}
\nf{Provenance}{Rattachement selon l’EDB : \emph{Coem. Commodillae}.}
\nf{Support}{Plaque de marbre. Inscription gravée. Signes décrits par l’EDB : croix monogrammatique avec alpha et oméga. Fiche EDB : Carlo Carletti, première version 2003.}
\ntext{\lnum{1}Asterius innocens\\
\lnum{2}hic positus est qui vixit an(n)o\\
\lnum{3}m(ensibus) VIII d(iebus) X Kal(endis) Iul(iis) Val(entiniano) et [Val(ente) cons(ulibu)s]}
\nf{Témoins}{EDB (EDB17465) : texte concordant aux conventions de notation près (indice de similarité 0,949).}
\nf{Traduction}{Asterius, innocent, a été placé ici ; il vécut un an (?) 8 mois 10 jours ; aux calendes de juillet, sous le consulat de Valentinien et de Valens.}
\nf{Langue}{8 abréviations développées par l’éditeur (dont an(n)o, m(ensibus), d(iebus), Kal(endis), Iul(iis), Val(entiniano)).}
\nf{Chronologie}{Noms consulaires : Valentinianus (lu) et Valens (lu). Dates du jour : événement indéterminé : X Kalendis Iuliis = 22 juin 365 (lu). Âge déclaré : 8 m. 10 j. (lu).}
\nf{Bibliographie}{ICUR II 6025 ; ILCV +2926.}
\nf{Statut}{Lecture secondaire (EDCS) concordante avec l’EDB : lecture stabilisée entre deux bases, non collationnée sur l’ICVR ni sur le monument.}
\end{notice}

\begin{notice}{96}{365}{ICUR n.s. III, 7382 · EDCS-34900710 · EDB24130}
\label{not:96}
\nf{Datation}{365 — directe : nom consulaire lu}
\nf{Genre}{Épitaphe ; ensembles A, B, C ; 44 lettres transmises ; 55 \% restituées.}
\nf{Références}{ICUR III 7382.}
\nf{Provenance}{Rattachement selon l’EDB : \emph{Coem. Domitillae pars superior}. Conservation selon l’EDB : \emph{Coem. Domitillae pars inferior}.}
\nf{Support}{Plaque de marbre. 18 × 34 × 1,5 cm (h. × l. × ép.). Inscription gravée. Reproductions signalées par l’EDB : photographie PCAS (droits réservés ; non reproduites). Fiche EDB : Carlo Carletti / Antonio Enrico Felle, première version 2003.}
\ntext{\lnum{1}Fl(avi) Valentini[ano Aug(usto) et Val\\
\lnum{2}en]te Aug(usto) cons[s(ulibus) - - -]}
\nf{Témoins}{EDB (EDB24130 ; écarts à examiner) : \emph{Fl(avio) Valentini[ano Aug(usto) et Fl(avio)? Val]/[en]te Aug(usto) cons(ulibu)[s ---] / [--- in] pa[ce ---?]}.}
\nf{Traduction}{Sous le consulat de Flavius Valentinien Auguste et de Valens Auguste […].}
\nf{Langue}{4 abréviations développées par l’éditeur (dont Fl(avi), Aug(usto), conss(ulibus)).}
\nf{Chronologie}{Noms consulaires : Valentinianus (lu) et Valens (lu).}
\nf{Bibliographie}{ICUR III 7382. Bibliographie complémentaire selon l’EDB : CLRE, 642.}
\nf{Statut}{Lecture secondaire (EDCS) ; écarts avec l’EDB à arbitrer sur l’ICVR. Plus de la moitié des lettres sont restituées. L’EDCS signale une photographie (non consultée).}
\end{notice}

\begin{notice}{97}{365}{ICUR n.s. IV, 9561 · EDCS-36101090 · EDB18365}
\label{not:97}
\nf{Datation}{365 — directe : nom consulaire lu}
\nf{Genre}{Épitaphe ; ensembles A, B, C ; 107 lettres transmises.}
\nf{Références}{ICUR IV 9561.}
\nf{Provenance}{Rattachement selon l’EDB : \emph{Coem. Callisti pars inferior}.}
\nf{Support}{Plaque de marbre. Inscription gravée. Reproductions signalées par l’EDB : photographie PCAS (droits réservés ; non reproduites). Fiche EDB : Antonio Enrico Felle, première version 2003.}
\ntext{\lnum{1}DD(ominis) nn(ostris) ⸢V⸣alentiniano et ⸢V⸣alente Aug(ustis) conss(ulibus)\\
\lnum{2}Lea quae deposita est Kal(endis) Iun(iis) ann(orum) n(umero) XI\\
\lnum{3}Sozomene quae deposita est VI Idus Oct(obres) ann(orum) VIII m(ensium) VIII}
\nf{Témoins}{Export EDCS de 2020 : \emph{DD(ominis) nn(ostris) alentiniano et alente Aug(ustis) conss(ulibus) / Lea quae deposita est Kal(endis) Iun(iis) ann(orum) n(umero) XI / Sozomene quae deposita est VI Idus Oct(obres) ann(orum) VIII m(ensium) VIII}. EDB (EDB18365) : texte concordant aux conventions de notation près (indice de similarité 0,981).}
\nf{Apparat}{⸢V⸣ : B gravé ; ⸢V⸣ : B gravé}
\nf{Traduction}{Sous le consulat de nos seigneurs Valentinien et Valens Augustes : Lea, qui a été déposée aux calendes de juin, âgée de 11 ans ; Sozomenè, qui a été déposée le 6e jour avant les ides d’octobre, âgée de 8 ans 8 mois.}
\nf{Langue}{L’éditeur signale 2 graphies corrigées : B gravé pour V (bêtacisme) ×2.}
\nf{Chronologie}{Noms consulaires : Valentinianus (lu) et Valens (lu). Dates du jour : déposition : Kalendis Iuniis = 1 juin 365 (lu) ; déposition : VI Idus Octobres = 10 octobre 365 (lu). Âge déclaré : 8 a. (lu).}
\nf{Bibliographie}{ICUR IV 9561.}
\nf{Statut}{Lecture secondaire (EDCS) concordante avec l’EDB : lecture stabilisée entre deux bases, non collationnée sur l’ICVR ni sur le monument. L’EDCS signale une photographie (non consultée).}
\end{notice}

\begin{notice}{98}{461/482}{EDCS-16300190 · EDB20076 · EDH HD028015}
\label{not:98}
\nf{Datation}{461/482 — hors période}
\nf{Genre}{Épitaphe ; notice hors période : exclue des ensembles B et C ; 54 lettres transmises ; 56 \% restituées.}
\nf{Références}{AE 1914, 68.}
\nf{Provenance}{Rattachement selon l’EDB : \emph{Coem. Cyriacae ad s.Laurentium}. Conservation selon l’EDB : \emph{Musei Capitolini}.}
\nf{Support}{Plaque de marbre. 32 × 61,5 × 5 cm (h. × l. × ép.). Inscription gravée. Fiche EDB : Miriam Ramunni, première version 2015.}
\ntext{\lnum{1}Dep(ositus) (H)yppolytus(!) in p[ace - - - consulatu]\\
\lnum{2}Fl(avi) Severin[i et Fl(avi) Dagalaiphi]}
\nf{Témoins}{EDB (EDB20076 ; divergence forte) : \emph{dep(ositus) Yppolytus in p[ace ---] / Fl(avio) Severin[o v(iro) c(larissimo) cons(ule)]}.}
\nf{Apparat}{(!) : forme signalée par l’éditeur}
\nf{Traduction}{Déposé, Hippolytus, en paix […], sous le consulat de Flavius Severinus […].}
\nf{Langue}{L’éditeur signale au moins une forme anormale (sic). 4 abréviations développées par l’éditeur (dont Dep(ositus), (H)yppolytus, Fl(avi)).}
\nf{Chronologie}{Notice hors période. Le nom \emph{Severin[us]} est lu ; \emph{[Dagalaiphi]} est une restitution de l'EDCS qui fondait l'attribution à 366. L'EDB (EDB20076, d'après la CLRE, p. 680) et l'EDH (HD028015) lisent \emph{Fl(avio) Severin[o v(iro) c(larissimo) cons(ule)]} et datent le texte de 461 ou de 482, années de consulats de Severinus. La notice est conservée dans le catalogue pour la stabilité des numéros, mais exclue de tous les ensembles analytiques.}
\nf{Commentaire}{La formule in pace est attestée seulement par restitution.}
\nf{Bibliographie}{AE 1914, 68. Bibliographie complémentaire selon l’EDB : CLRE, 680.}
\nf{Statut}{Lecture secondaire (EDCS) ; divergence forte avec l’EDB : le texte ne doit pas être exploité sans contrôle de l’édition imprimée. Plus de la moitié des lettres sont restituées.}
\end{notice}

\annee{366}

\begin{notice}{99}{366}{ICUR n.s. III, 8723 · EDCS-22300012 · EDB20991}
\label{not:99}
\nf{Datation}{366 — directe : nom consulaire lu}
\nf{Genre}{Épitaphe ; ensembles A, B, C ; 76 lettres transmises.}
\nf{Références}{CIL X *1088,458 = ICUR III 8723 = ILCV 3252b = IMCCatania 360.}
\nf{Provenance}{Rattachement selon l’EDB : \emph{Tituli qui in coem. Callisti reperti traduntur}. Conservation selon l’EDB : \emph{Catania   ; Museo Civico}.}
\nf{Support}{Plaque de marbre. 28 × 30 × 3,8 cm (h. × l. × ép.). Inscription gravée. Signes décrits par l’EDB : chrisme. Reproductions signalées par l’EDB : dessin, ICVR I (de Rossi) (droits réservés ; non reproduites). Fiche EDB : Carlo Carletti, première version 2016.}
\ntext{\lnum{1}Bene merenti filiae\\
\lnum{2}Agapeni vir⸢g⸣ini in pace\\
\lnum{3}dep(osita) III Kal(endas) Dec(embres) Gratiano et\\
\lnum{4}Dagalaifo co(n)s(ulibus) v(ixit) a(nnos) XXI}
\nf{Témoins}{Export EDCS de 2020 : \emph{Bene merenti filiae / Agapeni virini in pace / dep(osita) III Kal(endas) Dec(embres) Gratiano et / Dagalaifo co(n)s(ulibus) v(ixit) a(nnos) XXI}. EDB (EDB20991) : texte concordant aux conventions de notation près (indice de similarité 0,953).}
\nf{Apparat}{⸢g⸣ : C gravé}
\nf{Traduction}{À leur fille bien méritante Agapè, vierge, en paix ; déposée le 3e jour avant les calendes de décembre, sous le consulat de Gratien et de Dagalaifus ; elle vécut 21 ans.}
\nf{Langue}{L’éditeur signale une graphie corrigée : C gravé pour g (c pour g).}
\nf{Chronologie}{Noms consulaires : Gratianus (lu) et Dagalaifus (lu). Dates du jour : déposition : III Kalendas Decembres = 29 novembre 366 (lu). Âge déclaré : 21 a. (lu).}
\nf{Commentaire}{La formule in pace est conservée.}
\nf{Bibliographie}{CIL X *1088,458 ; ICUR III 8723 ; ILCV 3252b ; IMCCatania 360. Bibliographie complémentaire selon l’EDB : Korhonen 2004, 247, n. 360 ; ILCV, 3252B ; CIL X, 1088\textbackslash{}*, 458.}
\nf{Statut}{Lecture secondaire (EDCS) concordante avec l’EDB : lecture stabilisée entre deux bases, non collationnée sur l’ICVR ni sur le monument. L’EDCS signale une photographie (non consultée).}
\end{notice}

\begin{notice}{100}{366}{ICUR n.s. I, 1429 · EDCS-39000413 · EDB31193}
\label{not:100}
\nf{Datation}{366 — directe : nom consulaire lu}
\nf{Genre}{Épitaphe ; ensembles A, B, C ; 109 lettres transmises ; 3 \% restituées.}
\nf{Références}{ICUR I 1429 = ILCV 2942.}
\nf{Provenance}{Rattachement selon l’EDB : \emph{Incertae originis} (le volume I des ICVR réunit les inscriptions dont la provenance est perdue). Conservation selon l’EDB : \emph{Musei Vaticani: Lapidario Cristiano}.}
\nf{Support}{Plaque de marbre. 22 × 108 cm (h. × l.). Inscription gravée. Signes décrits par l’EDB : chrisme. Reproductions signalées par l’EDB : reproduction, Perret 1851 (droits réservés ; non reproduites). Fiche EDB : Filippo Piazzolla, première version 2015.}
\ntext{\lnum{1}Ligitimus bene merent⸢i⸣ filio Asello\\
\lnum{2}dulcissimo qui v(i)xit annis VII et m(ensibus) X et d(iebus) XXI\\
\lnum{3}[dep]ositus d(ie) VIIII K(alendas) Oct(obres) Gratiano et Dagalaifo conss(ulibus)}
\nf{Témoins}{Export EDCS de 2020 : \emph{Ligitimus bene merent filio Asello / dulcissimo qui v(i)xit annis VII et m(ensibus) X et d(iebus) XXI / [dep]ositus d(ie) VIIII K(alendas) Oct(obres) Gratiano et Dagalaifo conss(ulibus)}. EDB (EDB31193) : texte concordant aux conventions de notation près (indice de similarité 0,936).}
\nf{Apparat}{⸢i⸣ : E gravé}
\nf{Traduction}{Legitimus à Asellus, son fils très doux et bien méritant, qui vécut 7 ans et 10 mois et 21 jours ; déposé le 9e jour avant les calendes d’octobre, sous le consulat de Gratien et de Dagalaifus.}
\nf{Langue}{L’éditeur signale une graphie corrigée : E gravé pour i (confusion e/i).}
\nf{Chronologie}{Noms consulaires : Gratianus (lu) et Dagalaifus (lu). Dates du jour : déposition : die VIIII Kalendas Octobres = 23 septembre 366 (lu). Âge déclaré : 7 a. 10 m. 21 j. (lu).}
\nf{Bibliographie}{ICUR I 1429 ; ILCV 2942. Bibliographie complémentaire selon l’EDB : Ferrua 1979b, 37, nn. 1429-1430 ; Ferrua 1964, 108 ; Lega 2000, 180 ; ILCV, 2942.}
\nf{Statut}{Lecture secondaire (EDCS) concordante avec l’EDB : lecture stabilisée entre deux bases, non collationnée sur l’ICVR ni sur le monument. L’EDCS signale une photographie (non consultée).}
\end{notice}

\begin{notice}{101}{366}{ICUR n.s. II, 4163 · EDCS-33900327 · EDB19222}
\label{not:101}
\nf{Datation}{366 — directe : lecture partielle d’un nom consulaire}
\nf{Genre}{Épitaphe ; ensembles A, B, C ; 40 lettres transmises ; 32 \% restituées.}
\nf{Références}{ICUR II 4163.}
\nf{Provenance}{Rattachement selon l’EDB : \emph{Basilica S.Petri apostoli}.}
\nf{Support}{Plaque de marbre. Inscription gravée. Fiche EDB : Antonio Enrico Felle, première version 2003.}
\ntext{\lnum{1}Gratiano et Daga]laifo [conss(ulibus) - - -]\\
\lnum{2}[- - - do]mini [- - -]\\
\lnum{3}[- - -]o dom[}
\nf{Témoins}{EDB (EDB19222) : texte concordant aux conventions de notation près (indice de similarité 0,955).}
\nf{Traduction}{Sous le consulat de Gratien et de Dagalaifus […] ; […] du Seigneur (?) […].}
\nf{Langue}{1 abréviation développée par l’éditeur (dont conss(ulibus)).}
\nf{Chronologie}{Noms consulaires : Gratianus (restitué) et Dagalaifus (lu en partie).}
\nf{Commentaire}{Vocabulaire religieux : Dominus (référent indéterminé) — domini (partiellement restitué).}
\nf{Bibliographie}{ICUR II 4163.}
\nf{Statut}{Lecture secondaire (EDCS) concordante avec l’EDB : lecture stabilisée entre deux bases, non collationnée sur l’ICVR ni sur le monument. L’EDCS signale une photographie (non consultée).}
\end{notice}

\begin{notice}{102}{366}{ICUR n.s. II, 4269 · EDCS-33900176 · EDB8229}
\label{not:102}
\nf{Datation}{366 — directe : nom consulaire lu}
\nf{Genre}{Épitaphe ; ensembles A, B, C ; 121 lettres transmises.}
\nf{Références}{ICUR II 4269 = ILCV 4606.}
\nf{Provenance}{Rattachement selon l’EDB : \emph{Coem. Octavillae ad s.Pancratium}.}
\nf{Support}{Support non précisé par l’EDB. Reproductions signalées par l’EDB : dessin, ICVR I (de Rossi) (droits réservés ; non reproduites). Fiche EDB : Antonio Enrico Felle, première version 2003.}
\ntext{\lnum{1}Cecili dominus Felicis post cons(ulatum) Valen\\
\lnum{2}tiniani et Valentis VII Kal(endas) Feb(ruarias) depos(itus)\\
\lnum{3}Apolinus qui vix(it) ann(os) XVIIII m(enses) VII d(ies) VI\\
\lnum{4}qu(i)escet in pace\\
\lnum{5}mater Leopardes}
\nf{Témoins}{EDB (EDB8229) : texte concordant aux conventions de notation près (indice de similarité 1).}
\nf{Traduction}{Caecilius, maître de Felix (?) ; après le consulat de Valentinien et Valens, le 7e jour avant les calendes de février, a été déposé Apollinus (?), qui vécut 19 ans 7 mois 6 jours ; il repose en paix ; sa mère Leopardis.}
\nf{Langue}{9 abréviations développées par l’éditeur (dont cons(ulatum), Kal(endas), Feb(ruarias), depos(itus), vix(it), ann(os)).}
\nf{Chronologie}{Noms consulaires : Valentinianus (lu). Datation par post-consulat. Dates du jour : événement indéterminé : VII Kalendas Februarias = 26 janvier 366 (lu). Âge déclaré : 19 a. 7 m. 6 j. (lu).}
\nf{Commentaire}{La formule in pace est conservée. Vocabulaire religieux : repos et sommeil de la mort — quiescet (lu).}
\nf{Bibliographie}{ICUR II 4269 ; ILCV 4606.}
\nf{Statut}{Lecture secondaire (EDCS) concordante avec l’EDB : lecture stabilisée entre deux bases, non collationnée sur l’ICVR ni sur le monument. L’EDCS signale une photographie (non consultée).}
\end{notice}

\begin{notice}{103}{366}{ICUR n.s. IV, 10853 · EDCS-36901720 · EDB39162}
\label{not:103}
\nf{Datation}{366 — directe : lecture partielle d’un nom consulaire}
\nf{Genre}{Épitaphe ; ensembles A, B, C ; 35 lettres transmises ; 63 \% restituées.}
\nf{Références}{ICUR IV 10853.}
\nf{Provenance}{Rattachement selon l’EDB : \emph{Coem. Callisti pars superior}. Conservation selon l’EDB : \emph{Coem. Callisti pars superior}.}
\nf{Support}{Plaque de marbre. 18 × 35 cm (h. × l.). Inscription gravée. Reproductions signalées par l’EDB : photographie PCAS (droits réservés ; non reproduites). Fiche EDB : Marida Pierno, première version 2015.}
\ntext{\lnum{1}]V[- - -]\\
\lnum{2}[- - -]I di(es) XX [- - -]\\
\lnum{3}[- - - Gratian]o et Daga[laifo co(n)s(ulibus)}
\nf{Témoins}{EDB (EDB39162 ; écarts à examiner) : \emph{[---]U[---] / [---]I di(ebus) XX[---] / [Gratian]o et Daga[laifo] [texte tronqué dans la liste de résultats]}.}
\nf{Traduction}{[…] 1 (?) […] 20 jours […] ; […] sous le consulat de Gratien et de Dagalaifus.}
\nf{Langue}{3 abréviations développées par l’éditeur (dont di(es), co(n)s, (ulibus)).}
\nf{Chronologie}{Noms consulaires : Gratianus (restitué) et Dagalaifus (lu en partie).}
\nf{Bibliographie}{ICUR IV 10853.}
\nf{Statut}{Lecture secondaire (EDCS) ; écarts avec l’EDB à arbitrer sur l’ICVR. Plus de la moitié des lettres sont restituées. L’EDCS signale une photographie (non consultée).}
\end{notice}

\begin{notice}{104}{366}{ICUR n.s. IV, 11097 · EDCS-37200379 · EDB39556}
\label{not:104}
\nf{Datation}{366 — directe : lecture partielle d’un nom consulaire}
\nf{Genre}{Épitaphe ; ensembles A, B, C ; 79 lettres transmises ; 38 \% restituées.}
\nf{Références}{ICUR IV 11097.}
\nf{Provenance}{Rattachement selon l’EDB : \emph{Coem. Callisti area subdialis et incerti loci}. Conservation selon l’EDB : \emph{Coem. Callisti area subdialis et incerti loci}.}
\nf{Support}{Plaque de marbre. 59 × 51 cm (h. × l.). Inscription gravée. Reproductions signalées par l’EDB : photographie PCAS ×2 (droits réservés ; non reproduites). Fiche EDB : Marida Pierno, première version 2015.}
\ntext{\lnum{1}[Gratiano et Dagal]aifo cons[s(ulibus) - - -]\\
\lnum{2}[- - -] cum uxore [- - -]\\
\lnum{3}[- - -] ann⸢o⸣s XLV[- - -]\\
\lnum{4}[- - -] Ma]rtias\\
\lnum{5}[- - -]r fecit\\
\lnum{6}[- - - bene me]renti in pace d[ei(?)]}
\nf{Témoins}{Export EDCS de 2020 : \emph{[Gratiano et Dagal]aifo cons[s(ulibus) 3] / [3] cum uxore [3] / [3] anns XLV[3] / [3] Ma]rtias / [3]r fecit / [3 bene me]renti in pace d[ei(?)]}. EDB (EDB39556 ; écarts à examiner) : \emph{[Gratiano et Dagal]aifo cons[ulibus] / [---] cum uxore [---] / [---] annu(s) XLV[---] / [kalendas ma]rtias / [---]r fecit / [beneme]renti in pace d[---]}.}
\nf{Apparat}{⸢o⸣ : V gravé}
\nf{Traduction}{Sous le consulat de Gratien et de Dagalaifus […] ; […] avec son épouse […] ; […] 45 ans […] ; […] mars ; […] a fait […] ; bien méritant, dans la paix de Dieu.}
\nf{Langue}{L’éditeur signale une graphie corrigée : V gravé pour o (confusion o/u).}
\nf{Chronologie}{Noms consulaires : Gratianus (restitué) et Dagalaifus (lu en partie).}
\nf{Commentaire}{La formule in pace est conservée. Vocabulaire religieux : désignation de Dieu — dei (partiellement restitué).}
\nf{Bibliographie}{ICUR IV 11097.}
\nf{Statut}{Lecture secondaire (EDCS) ; écarts avec l’EDB à arbitrer sur l’ICVR. L’EDCS signale une photographie (non consultée).}
\end{notice}

\begin{notice}{105}{366}{ICUR n.s. IV, 11763 · EDCS-35700942 · EDB1004}
\label{not:105}
\nf{Datation}{366 — directe : nom consulaire lu}
\nf{Genre}{Épitaphe ; ensembles A, B, C ; 52 lettres transmises.}
\nf{Références}{ICUR IV 11763 = ILCV 1126a (em) = Mossong-2022, 00127.}
\nf{Provenance}{Rattachement selon l’EDB : \emph{Coem. anonymum ad viam Appiam}. Conservation selon l’EDB : \emph{Coem. anonymum ad viam Appiam}.}
\nf{Support}{Paroi de enduit. Inscription peinte. Fiche EDB : Maria Cristina Grisanzio, première version 2003.}
\ntext{\lnum{1}Vincentius\\
\lnum{2}presb(yter)\\
\lnum{3}dep(ositus) XI Kal(endas) Iulias\\
\lnum{4}Gratiano et\\
\lnum{5}Dagalaifo co(n)s(ulibu)s}
\nf{Témoins}{EDB (EDB1004) : texte concordant aux conventions de notation près (indice de similarité 1).}
\nf{Traduction}{Vincentius, prêtre, déposé le 11e jour avant les calendes de juillet, sous le consulat de Gratien et de Dagalaifus.}
\nf{Langue}{5 abréviations développées par l’éditeur (dont presb(yter), dep(ositus), Kal(endas), co(n)s, (ulibu)s).}
\nf{Chronologie}{Noms consulaires : Gratianus (lu) et Dagalaifus (lu). Ne sont pas des consuls : Vincentius / presb(yter) (défunt, prêtre). Dates du jour : déposition : XI Kalendas Iulias = 21 juin 366 (lu).}
\nf{Commentaire}{Fonction ou métier mentionné : presbyter.}
\nf{Bibliographie}{ICUR IV 11763 ; ILCV 1126a (em) ; Mossong-2022, 00127.}
\nf{Statut}{Lecture secondaire (EDCS) concordante avec l’EDB : lecture stabilisée entre deux bases, non collationnée sur l’ICVR ni sur le monument.}
\end{notice}

\annee{367}

\begin{notice}{106}{367}{EDCS-16600352 · EDB13167 · EDH HD030246}
\label{not:106}
\nf{Statut}{Doublon établi : même défunt (Zosimus), même date (13 mai 367), même formule \emph{annorum plus minus}, même coupe inhabituelle des lignes (\emph{Ma/ias}, \emph{Lupicin/o}). L'EDH (HD030246) et l'EDB (EDB13167) rattachent le texte de l'AE 1905, 71 à ICUR II 6028. La notice est fusionnée dans le n° 120 ; sa leçon \emph{c(iu)esq(it!)}, contre \emph{recessit} dans l'ICVR et l'EDB, est donnée en apparat du n° 120.}
\nf{Références}{AE 1905, 71}
\ntext{\lnum{1}Zosimus c(iu)esq(it !) in pace\\
\lnum{2}annorum plus minus\\
\lnum{3}[- - -] III Idus Ma\\
\lnum{4}[ias] Lupicin\\
\lnum{5}o et Io⸢v⸣ino cons(ulibus)}
\nf{Traduction}{Zosimus repose en paix, âgé d’environ […] ans ; […] le 3e jour avant les ides de mai, sous le consulat de Lupicinus et de Iovinus.}
\nf{Voir}{la notice 120 (p.~\pageref{not:120}).}
\end{notice}

\begin{notice}{107}{367}{ICUR n.s. I, 2792 · EDCS-39900335 · EDB32714}
\label{not:107}
\nf{Datation}{367 — directe : nom consulaire lu}
\nf{Genre}{Épitaphe ; ensembles A, B, C ; 35 lettres transmises ; 51 \% restituées.}
\nf{Références}{CIL XIV *290,12 = ICUR I 2792.}
\nf{Provenance}{Rattachement selon l’EDB : \emph{Incertae originis} (le volume I des ICVR réunit les inscriptions dont la provenance est perdue).}
\nf{Support}{Support non précisé par l’EDB. Fiche EDB : Filippo Piazzolla, première version 2015.}
\ntext{\lnum{1}]OEI Antus[- - -]\\
\lnum{2}[- - - Lupicino et] Iovino con[ss(ulibus)}
\nf{Témoins}{EDB (EDB32714 ; écarts à examiner) : \emph{------ / [---]OEI Antus[---] / [--- Lupicino et] Iovino con[s(ulibu)s] [texte tronqué dans la liste de résultats]}.}
\nf{Traduction}{[…] sous le consulat de Lupicinus et de Iovinus.}
\nf{Langue}{1 abréviation développée par l’éditeur (dont conss(ulibus)).}
\nf{Chronologie}{Noms consulaires : Lupicinus (restitué) et Iovinus (lu).}
\nf{Bibliographie}{CIL XIV *290,12 ; ICUR I 2792. Bibliographie complémentaire selon l’EDB : CIL XIV, 290\textbackslash{}*.}
\nf{Statut}{Lecture secondaire (EDCS) ; écarts avec l’EDB à arbitrer sur l’ICVR. Plus de la moitié des lettres sont restituées.}
\end{notice}

\begin{notice}{108}{367}{ICUR n.s. I, 265 · EDCS-38600271 · EDB27783}
\label{not:108}
\nf{Datation}{367 — directe : nom consulaire lu}
\nf{Genre}{Épitaphe ; ensembles A, B, C ; 35 lettres transmises ; 34 \% restituées.}
\nf{Références}{ICUR I 265.}
\nf{Provenance}{Rattachement selon l’EDB : \emph{Incertae originis} (le volume I des ICVR réunit les inscriptions dont la provenance est perdue).}
\nf{Support}{Support non précisé par l’EDB. Fiche EDB : Filippo Piazzolla, première version 2014.}
\ntext{\lnum{1}[- - - St]ercoria se [viva - - -]\\
\lnum{2}[Iovino] et Lupicino cons(ulibus) [}
\nf{Témoins}{EDB (EDB27783) : texte concordant aux conventions de notation près (indice de similarité 0,955).}
\nf{Traduction}{[…] Stercoria, de son vivant […] ; sous le consulat de Iovinus et de Lupicinus.}
\nf{Langue}{1 abréviation développée par l’éditeur (dont cons(ulibus)).}
\nf{Chronologie}{Noms consulaires : Lupicinus (lu) et Iovinus (restitué).}
\nf{Bibliographie}{ICUR I 265.}
\nf{Statut}{Lecture secondaire (EDCS) concordante avec l’EDB : lecture stabilisée entre deux bases, non collationnée sur l’ICVR ni sur le monument.}
\end{notice}

\begin{notice}{109}{367}{ICUR n.s. I, 896 · EDCS-38700412 · EDB29441}
\label{not:109}
\nf{Datation}{367 — directe : nom consulaire lu}
\nf{Genre}{Épitaphe ; ensembles A, B, C ; 99 lettres transmises.}
\nf{Références}{ICUR I 896 = ILCV 2943 (em).}
\nf{Provenance}{Rattachement selon l’EDB : \emph{Incertae originis} (le volume I des ICVR réunit les inscriptions dont la provenance est perdue).}
\nf{Support}{Plaque de marbre. Inscription gravée. Signes décrits par l’EDB : chrisme et chrisme avec alpha et oméga. Reproductions signalées par l’EDB : dessin, ICVR I (de Rossi) (droits réservés ; non reproduites). Fiche EDB : Filippo Piazzolla, première version 2014.}
\ntext{\lnum{1}Bene merenti\\
\lnum{2}Valeriae quae\\
\lnum{3}vixit annos\\
\lnum{4}XXXVII deposita\\
\lnum{5}XIIII Kalendas\\
\lnum{6}Martias post con\\
\lnum{7}sulatu(m) Gratian\\
\lnum{8}i et Dagalaifi}
\nf{Témoins}{EDB (EDB29441) : texte concordant aux conventions de notation près (indice de similarité 0,979).}
\nf{Traduction}{À Valeria, bien méritante, qui vécut 37 ans ; déposée le 14e jour avant les calendes de mars, après le consulat de Gratien et de Dagalaifus.}
\nf{Langue}{1 abréviation développée par l’éditeur (dont sulatu(m)).}
\nf{Chronologie}{Noms consulaires : Gratianus (lu) et Dagalaifus (lu). Datation par post-consulat. Dates du jour : déposition : XIIII Kalendas Martias = 16 février 367 (lu). Âge déclaré : 37 a. (lu).}
\nf{Bibliographie}{ICUR I 896 ; ILCV 2943 (em). Bibliographie complémentaire selon l’EDB : ILCV, 2943.}
\nf{Statut}{Lecture secondaire (EDCS) concordante avec l’EDB : lecture stabilisée entre deux bases, non collationnée sur l’ICVR ni sur le monument. L’EDCS signale une photographie (non consultée).}
\end{notice}

\begin{notice}{110}{367}{ICUR n.s. I, 940 · EDCS-34700228 · EDB19681}
\label{not:110}
\nf{Datation}{367 — directe : nom consulaire lu}
\nf{Genre}{Épitaphe ; ensembles A, B, C ; 152 lettres transmises.}
\nf{Références}{ICUR I 940 = ILCV 2878.}
\nf{Provenance}{Rattachement selon l’EDB : \emph{Incertae originis} (le volume I des ICVR réunit les inscriptions dont la provenance est perdue).}
\nf{Support}{Plaque de marbre. Inscription gravée. Reproductions signalées par l’EDB : dessin, ICVR I (de Rossi) (droits réservés ; non reproduites). Fiche EDB : Antonio Enrico Felle, première version 2014.}
\ntext{\lnum{1}ἐν εἰρήνῃ ἀνεπάη Βεναινᾶτος ἀπ(ὸ) κο(μήτων)\\
\lnum{2}ζήσας ἔτη οεʹ ὃς κῖτε ἐν τῷδε τόπῳ ἐξελτὼν\\
\lnum{3}ἐκ τοῦ βίου τῇ πρ(ὸ) δʹ ἰδ(ῶ)ν μαρ(τίων) ὑπ(ατείᾳ) Λουππικί(νου)\\
\lnum{4}καὶ Ἰοβίνου Θεοφίλα ἐποίησα συμβίῳ\\
\lnum{5}Benenato Theofila fecit uxor eius}
\nf{Témoins}{Export EDCS de 2020 : \emph{// Benenato Theofila fecit uxor eius}. EDB (EDB19681) : texte concordant aux conventions de notation près (indice de similarité 0,966).}
\nf{Traduction}{(grec) En paix a reposé Benenatus, ancien comes, ayant vécu 75 ans ; il gît en ce lieu, sorti de la vie le 4e jour avant les ides de mars, sous le consulat de Lupicinus et Iovinus ; moi, Theophila, j’ai fait (ce monument) pour mon époux. // (latin) À Benenatus, Theophila son épouse a fait (ce monument).}
\nf{Langue}{Le texte est rédigé en grec, en tout ou en partie ; l’export de 2020 ne le transmettait pas. 7 abréviations développées par l’éditeur (dont (ὸ), (μήτων), (ῶ), (τίων), (ατείᾳ)).}
\nf{Chronologie}{Noms consulaires : Lupicinus (lu) et Iovinus (lu).}
\nf{Commentaire}{Vocabulaire religieux : anthroponyme (niveau 6) — theofila (anthroponyme). Commentaire développé : dossier p.~\pageref{dos:110}.}
\nf{Bibliographie}{ICUR I 940 ; ILCV 2878.}
\nf{Statut}{Lecture secondaire (EDCS) concordante avec l’EDB : lecture stabilisée entre deux bases, non collationnée sur l’ICVR ni sur le monument. L’EDCS signale une photographie (non consultée).}
\end{notice}

\begin{notice}{111}{367}{ICUR n.s. I, 1350 · EDCS-39000328 · EDB30676}
\label{not:111}
\nf{Datation}{367 — directe : nom consulaire lu}
\nf{Genre}{Épitaphe ; ensembles A, B, C ; 62 lettres transmises.}
\nf{Références}{ICUR I 1350 = ILCV 1296a (add).}
\nf{Provenance}{Rattachement selon l’EDB : \emph{Incertae originis} (le volume I des ICVR réunit les inscriptions dont la provenance est perdue). Conservation selon l’EDB : \emph{Musei Capitolini}.}
\nf{Support}{Plaque de marbre. 25 × 40 cm (h. × l.). Inscription gravée. Fiche EDB : Filippo Piazzolla, première version 2015.}
\ntext{\lnum{1}Aur(elio) Satullo puero ⸢q⸣u⸢i⸣\\
\lnum{2}vixit annos XX dep(ositus)\\
\lnum{3}VIIII Kal(endas) Mai(as) Lupici\\
\lnum{4}no Iovino cons(ulibus)}
\nf{Témoins}{Export EDCS de 2020 : \emph{Aur(elio) Satullo puero u / vixit annos XX dep(ositus) / VIIII Kal(endas) Mai(as) Lupici/no Iovino cons(ulibus)}. EDB (EDB30676) : texte concordant aux conventions de notation près (indice de similarité 0,932).}
\nf{Apparat}{⸢q⸣ : O gravé ; ⸢i⸣ : E gravé}
\nf{Traduction}{À Aurelius Satullus, jeune homme, qui vécut 20 ans ; déposé le 9e jour avant les calendes de mai, sous le consulat de Lupicinus et Iovinus.}
\nf{Langue}{L’éditeur signale 2 graphies corrigées : O gravé pour q (confusion de lettres) ; E gravé pour i (confusion e/i).}
\nf{Chronologie}{Noms consulaires : Lupicinus (lu) et Iovinus (lu). Dates du jour : déposition : VIIII Kalendas Maias = 23 avril 367 (lu). Âge déclaré : 20 a. (lu).}
\nf{Bibliographie}{ICUR I 1350 ; ILCV 1296a (add). Bibliographie complémentaire selon l’EDB : ILCV, 1296A.}
\nf{Statut}{Lecture secondaire (EDCS) concordante avec l’EDB : lecture stabilisée entre deux bases, non collationnée sur l’ICVR ni sur le monument.}
\end{notice}

\begin{notice}{112}{367}{ICUR n.s. I, 1401 · EDCS-34501139 · EDB30801}
\label{not:112}
\nf{Datation}{367 — directe : nom consulaire lu}
\nf{Genre}{Épitaphe ; ensembles A, B, C ; 37 lettres transmises ; 54 \% restituées.}
\nf{Références}{ICUR I 1401 = IGLFRPal 201.}
\nf{Provenance}{Rattachement selon l’EDB : \emph{Incertae originis} (le volume I des ICVR réunit les inscriptions dont la provenance est perdue). Conservation selon l’EDB : \emph{Lapidario Forense}.}
\nf{Support}{Plaque de marbre. 15 × 31 × 6 cm (h. × l. × ép.). Inscription gravée. Fiche EDB : Filippo Piazzolla, première version 2015.}
\ntext{\lnum{1}Quo]dvuldeu[s - - -]\\
\lnum{2}Io⸢v⸣ino e[t Lupicino cons(ulibus)]}
\nf{Témoins}{Export EDCS de 2020 : \emph{Quo]dvuldeu[s 3] / Ioino e[t Lupicino cons(ulibus)]}. EDB (EDB30801) : texte concordant aux conventions de notation près (indice de similarité 0,937).}
\nf{Apparat}{⸢v⸣ : B gravé}
\nf{Traduction}{Quodvultdeus […] ; sous le consulat de Iovinus et de Lupicinus.}
\nf{Langue}{L’éditeur signale une graphie corrigée : B gravé pour v (bêtacisme).}
\nf{Chronologie}{Noms consulaires : Lupicinus (restitué) et Iovinus (lu).}
\nf{Commentaire}{Vocabulaire religieux : anthroponyme (niveau 6) — quodvuldeus (anthroponyme).}
\nf{Bibliographie}{ICUR I 1401 ; IGLFRPal 201. Bibliographie complémentaire selon l’EDB : Panciera 1996, 370, n. 201, Tav. XLII, fig.1.}
\nf{Statut}{Lecture secondaire (EDCS) concordante avec l’EDB : lecture stabilisée entre deux bases, non collationnée sur l’ICVR ni sur le monument. Plus de la moitié des lettres sont restituées.}
\end{notice}

\begin{notice}{113}{367}{ICUR n.s. I, 1431 · EDCS-39000414 · EDB31195}
\label{not:113}
\nf{Datation}{367 — directe : lecture partielle d’un nom consulaire}
\nf{Genre}{Épitaphe ; ensembles A, B, C ; 39 lettres transmises ; 62 \% restituées.}
\nf{Références}{ICUR I 1431.}
\nf{Provenance}{Rattachement selon l’EDB : \emph{Incertae originis} (le volume I des ICVR réunit les inscriptions dont la provenance est perdue). Conservation selon l’EDB : \emph{Musei Vaticani: Lapidario Cristiano}.}
\nf{Support}{Plaque de marbre. 20 × 32 cm (h. × l.). Inscription gravée. Fiche EDB : Filippo Piazzolla, première version 2015.}
\ntext{\lnum{1}]ID[- - -]\\
\lnum{2}an(nos) V et [mens(es) - - - Lupici]\\
\lnum{3}no et Io⸢v⸣i[no cons(ulibus)}
\nf{Témoins}{Export EDCS de 2020 : \emph{]ID[3] / an(nos) V et [mens(es) 3 Lupici]/no et Ioi[no cons(ulibus)}. EDB (EDB31195) : même texte ; les écarts relevés dans la liste de résultats tiennent à la notation (relecture de la fiche détaillée). Relecture de la fiche détaillée de l’EDB : Même texte (lecture Lega 2000).}
\nf{Apparat}{⸢v⸣ : B gravé}
\nf{Traduction}{[…] 5 ans et […] mois ; sous le consulat de Lupicinus et de Iovinus.}
\nf{Langue}{L’éditeur signale une graphie corrigée : B gravé pour v (bêtacisme).}
\nf{Chronologie}{Noms consulaires : Lupicinus (restitué) et Iovinus (lu en partie).}
\nf{Bibliographie}{ICUR I 1431. Bibliographie complémentaire selon l’EDB : Lega 2000, 146.}
\nf{Statut}{Lecture secondaire (EDCS) concordante avec l’EDB : lecture stabilisée entre deux bases, non collationnée sur l’ICVR ni sur le monument. Plus de la moitié des lettres sont restituées. L’EDCS signale une photographie (non consultée).}
\end{notice}

\begin{notice}{114}{367}{ICUR n.s. I, 2086 · EDCS-39300310 · EDB32051}
\label{not:114}
\nf{Datation}{367 — directe : nom consulaire lu}
\nf{Genre}{Épitaphe ; ensembles A, B, C ; 99 lettres transmises ; 31 \% restituées.}
\nf{Références}{ICUR I 2086 = ILCV 4217a.}
\nf{Provenance}{Rattachement selon l’EDB : \emph{Incertae originis} (le volume I des ICVR réunit les inscriptions dont la provenance est perdue). Conservation selon l’EDB : \emph{Musei Vaticani: Galleria Lapidaria}.}
\nf{Support}{Plaque de marbre. 53 × 81 cm (h. × l.). Inscription gravée. Reproductions signalées par l’EDB : dessin, ICVR I (de Rossi) (droits réservés ; non reproduites). Fiche EDB : Filippo Piazzolla, première version 2015.}
\ntext{\lnum{1}Lupicino et Iovino cc[ons(ulibus) - - - dep(ositio)]\\
\lnum{2}Victori(a)es(!) q(uae) an(nos) XXV[- - - et cum]\\
\lnum{3}marito fecit an(nos) XIII[- - -]\\
\lnum{4}et pudicitia omnibus [nota(?) fuit]}
\nf{Témoins}{EDB (EDB32051 ; écarts à examiner) : \emph{Lupicino et Iovino cc̲[(on)s(ulibu)s --- dep(ositio)] / Victories q(ui) an(nis) XXV [--- et cum] / marito fecit an(nis) XIII [---] / et Pudicitia omnibus̲ [ ---] [texte tronqué dans la liste de résultats]}.}
\nf{Apparat}{(!) : forme signalée par l’éditeur}
\nf{Traduction}{Sous le consulat de Lupicinus et de Iovinus […] déposition de Victoria, qui (vécut) 25 ans […] et a vécu avec son mari 13 ans […] ; et sa chasteté fut connue de tous.}
\nf{Langue}{L’éditeur signale au moins une forme anormale (sic). 7 abréviations développées par l’éditeur (dont ccons(ulibus), dep(ositio), Victori(a)es, q(uae), an(nos)).}
\nf{Chronologie}{Noms consulaires : Lupicinus (lu) et Iovinus (lu).}
\nf{Bibliographie}{ICUR I 2086 ; ILCV 4217a. Bibliographie complémentaire selon l’EDB : Lega 2000 ; ILCV, 4217A.}
\nf{Statut}{Lecture secondaire (EDCS) ; écarts avec l’EDB à arbitrer sur l’ICVR. L’EDCS signale une photographie (non consultée).}
\end{notice}

\begin{notice}{115}{367}{ICUR n.s. I, 3175 · EDCS-40300321 · EDB33127}
\label{not:115}
\nf{Datation}{367 — directe : nom consulaire lu}
\nf{Genre}{Épitaphe ; ensembles A, B, C ; 210 lettres transmises ; 2 \% restituées.}
\nf{Références}{ICUR I 3175 = ILCV 4311.}
\nf{Provenance}{Rattachement selon l’EDB : \emph{Incertae originis} (le volume I des ICVR réunit les inscriptions dont la provenance est perdue).}
\nf{Support}{Support non précisé par l’EDB. Reproductions signalées par l’EDB : dessin, ICVR I (de Rossi) (droits réservés ; non reproduites). Fiche EDB : Filippo Piazzolla, première version 2015.}
\ntext{\lnum{1}Venerabili ac rari exempli feminae\\
\lnum{2}pud⸢i⸣cissim(a)e coniugi dulcissim(a)e\\
\lnum{3}Alexander fecit qu(a)e numquam\\
\lnum{4}mecum discordia qua fidelis\\
\lnum{5}habere potuit bene merenti\\
\lnum{6}in pace qu(a)e vixit annos XXVII\\
\lnum{7}et menses plus minus VIII deposita\\
\lnum{8}Non(is) Sept(embribus) conss(ulatu) Lupicini et Io[vini]}
\nf{Témoins}{Export EDCS de 2020 : \emph{Venerabili ac rari exempli feminae / pudcissim(a)e coniugi dulcissim(a)e / Alexander fecit qu(a)e numquam / mecum discordia qua fidelis / habere potuit bene merenti / in pace qu(a)e vixit annos XXVII / et menses plus minus VIII deposita / Non(is) Sept(embribus) conss(ulatu) Lupicini et Io[vini]}. EDB (EDB33127) : même texte ; les écarts relevés dans la liste de résultats tiennent à la notation (relecture de la fiche détaillée). Relecture de la fiche détaillée de l’EDB : Même texte ; non(is) / non(as).}
\nf{Apparat}{⸢i⸣ : U gravé}
\nf{Traduction}{À une femme vénérable et d’un exemple rare, épouse très chaste et très douce, Alexander a fait (ce monument) ; elle n’eut jamais de discorde avec moi, comme une femme fidèle peut l’avoir (?) ; bien méritante, en paix ; elle vécut 27 ans et environ 8 mois ; déposée aux nones de septembre, sous le consulat de Lupicinus et de Iovinus.}
\nf{Langue}{L’éditeur signale une graphie corrigée : U gravé pour i.}
\nf{Chronologie}{Noms consulaires : Lupicinus (lu) et Iovinus (restitué). Dates du jour : déposition : Nonis Septembribus = 5 septembre 367 (lu). Âge déclaré : env. 27 a. 8 m. (lu).}
\nf{Commentaire}{La formule in pace est conservée. Vocabulaire religieux : initiation et foi — fidelis (lu).}
\nf{Bibliographie}{ICUR I 3175 ; ILCV 4311. Bibliographie complémentaire selon l’EDB : ILCV, 4311.}
\nf{Statut}{Lecture secondaire (EDCS) concordante avec l’EDB : lecture stabilisée entre deux bases, non collationnée sur l’ICVR ni sur le monument. L’EDCS signale une photographie (non consultée).}
\end{notice}

\begin{notice}{116}{367}{ICUR n.s. II, 4805 · EDCS-36101286}
\label{not:116}
\nf{Datation}{367 — directe : nom consulaire lu}
\nf{Genre}{Épitaphe ; ensembles A, B, C ; 60 lettres transmises ; 52 \% restituées.}
\nf{Références}{ICUR II 4805 = ILCV +2807a.}
\nf{Provenance}{Rome (indication de l’EDCS) ; aucun rattachement cimétérial transmis.}
\nf{Support}{Non documenté dans les sources consultées.}
\ntext{\lnum{1}[I]ul(ius) Hilara[- - - qui vixit]\\
\lnum{2}[a]nn(os) XXVII de[positus - - -]\\
\lnum{3}Lupicino et Io[vino cons(ulibus)]}
\nf{Témoins}{Aucune fiche EDB appariée (numéro ICUR absent de la liste publique de l’EDB).}
\nf{Traduction}{Iulius Hilarus (?) […] qui vécut 27 ans ; déposé […] sous le consulat de Lupicinus et de Iovinus.}
\nf{Langue}{3 abréviations développées par l’éditeur (dont Iul(ius), ann(os), cons(ulibus)).}
\nf{Chronologie}{Noms consulaires : Lupicinus (lu) et Iovinus (restitué). Âge déclaré : 27 a. (partiellement restitué).}
\nf{Bibliographie}{ICUR II 4805 ; ILCV +2807a.}
\nf{Statut}{Lecture secondaire (EDCS) sans second témoin numérique. Plus de la moitié des lettres sont restituées.}
\end{notice}

\begin{notice}{117}{367}{ICUR n.s. II, 4807 · EDCS-36101288}
\label{not:117}
\nf{Datation}{367 — directe : nom consulaire lu}
\nf{Genre}{Épitaphe ; ensembles A, B, C ; 188 lettres transmises ; 21 \% restituées.}
\nf{Références}{ICUR II 4807 = ILCV 2609 (em).}
\nf{Provenance}{Rome (indication de l’EDCS) ; aucun rattachement cimétérial transmis.}
\nf{Support}{Non documenté dans les sources consultées.}
\ntext{\lnum{1}Bene merenti in pace Festo qui vixit i[- - -]\\
\lnum{2}annos XI et menses VIIII depositus XI K[al(endas ?) - - -]\\
\lnum{3}Iovino et Lupicino [cons(ulibus)]\\
\lnum{4}bene merenti compari Hilarino qui vixit a[nnos - - -]\\
\lnum{5}deposit(us) VI Kal(endas) Octobr(es) Valentiniano II et V[alente II cons(ulibus)]}
\nf{Témoins}{Aucune fiche EDB appariée (numéro ICUR absent de la liste publique de l’EDB).}
\nf{Traduction}{Au bien méritant, en paix, Festus, qui vécut 11 ans et 9 mois ; déposé le 11e jour avant les calendes de […], sous le consulat de Iovinus et de Lupicinus. À Hilarinus, compagnon bien méritant, qui vécut […] ans ; déposé le 6e jour avant les calendes d’octobre, sous le consulat de Valentinien pour la 2e fois et de Valens pour la 2e fois.}
\nf{Langue}{6 abréviations développées par l’éditeur (dont Kal(endas ?), cons(ulibus), deposit(us), Kal(endas), Octobr(es)).}
\nf{Chronologie}{Noms consulaires : Valentinianus (lu), Valens (lu), Lupicinus (lu) et Iovinus (lu). Dates du jour : déposition : VI Kalendas Octobres = 26 septembre 367 (lu). Âge déclaré : 11 a. 9 m. (lu).}
\nf{Commentaire}{La formule in pace est conservée.}
\nf{Bibliographie}{ICUR II 4807 ; ILCV 2609 (em).}
\nf{Statut}{Lecture secondaire (EDCS) sans second témoin numérique. L’EDCS signale une photographie (non consultée).}
\end{notice}

\begin{notice}{118}{367}{ICUR n.s. II, 6026 · EDCS-34200652 · EDB17466}
\label{not:118}
\nf{Datation}{367 — directe : nom consulaire lu}
\nf{Genre}{Épitaphe ; ensembles A, B, C ; 85 lettres transmises ; 12 \% restituées.}
\nf{Références}{ICUR II 6026 = ILCV 2824 (add) = Luciani-02, 00018.}
\nf{Provenance}{Rattachement selon l’EDB : \emph{Coem. Commodillae}. Conservation selon l’EDB : \emph{Coem. Commodillae}.}
\nf{Support}{Plaque de marbre. Inscription gravée. Signes décrits par l’EDB : chrisme. Fiche EDB : Carlo Carletti, première version 2003.}
\ntext{\lnum{1}In con[sulatu L]upicino(!) (e)t Io\\
\lnum{2}⸢v⸣ino(!) I[anu]aria qu(a)e ⸢v⸣ixit\\
\lnum{3}annos XXIIII diem I horas II s(emis) et recessit die Martis}
\nf{Témoins}{Export EDCS de 2020 : \emph{In con[sulatu L]upicino(!) (e)t Io/ino(!) I[anu]aria qu(a)e ixit / annos XXIIII diem I (h)(as) II s(emis) et recessit die Martis}. EDB (EDB17466 ; écarts à examiner) : \emph{((imago defunctae)) / / in con[sulatu L]upicino ⟨e⟩t Io/bino (!) I[anu]aria que bixit / annos XXIIII diem I ROII/S et recessit die martis / / ((avis)) ((monogramma Christi)) ((avis)) [texte tronqué dans la liste de résultats]}.}
\nf{Apparat}{⸢v⸣ : B gravé ; ⸢v⸣ : B gravé ; (!) : forme signalée par l’éditeur}
\nf{Traduction}{Sous le consulat de Lupicinus et Iovinus, Ianuaria, qui vécut 24 ans 1 jour 2 heures et demie, s’en est allée un mardi.}
\nf{Langue}{L’éditeur signale 2 graphies corrigées : B gravé pour v (bêtacisme) ×2. L’export de 2020 diffère de l’état 2022 (voir apparat). L’éditeur signale au moins une forme anormale (sic). 5 abréviations développées par l’éditeur (dont Lupicino(!), (e)t, ino(!), qu(a)e, s(emis)).}
\nf{Chronologie}{Noms consulaires : Lupicinus (lu en partie) et Iovinus (lu). Âge déclaré : 24 a. 2 h (lu).}
\nf{Commentaire}{Commentaire développé : dossier p.~\pageref{dos:118}.}
\nf{Bibliographie}{ICUR II 6026 ; ILCV 2824 (add) ; Luciani-02, 00018.}
\nf{Statut}{Lecture secondaire (EDCS) ; écarts avec l’EDB à arbitrer sur l’ICVR.}
\end{notice}

\begin{notice}{119}{367}{ICUR n.s. II, 6027 · EDCS-34200653 · EDB13166}
\label{not:119}
\nf{Datation}{367 — directe : nom consulaire lu}
\nf{Genre}{Épitaphe ; ensembles A, B, C ; 95 lettres transmises.}
\nf{Références}{ICUR II 6027 = ILCV 2976a = AE 1905, 72.}
\nf{Provenance}{Rattachement selon l’EDB : \emph{Coem. Commodillae}.}
\nf{Support}{Plaque de marbre. Inscription gravée. Fiche EDB : Carlo Carletti, première version 2015.}
\ntext{\lnum{1}Sosoro qui vixit\\
\lnum{2}(a)nnos XIIII dies\\
\lnum{3}VIIII decessit X Kal(endas)\\
\lnum{4}No⸢v⸣e⸢m⸣bres depo\\
\lnum{5}situs VIII K(a)l(endas) No⸢v⸣e⸢m⸣(bres)\\
\lnum{6}Lupicino et Iovino\\
\lnum{7}in pace}
\nf{Témoins}{Export EDCS de 2020 : \emph{Sosoro qui vixit / (a)nnos XIIII dies / VIIII decessit X Kal(endas) / Noebres depo/situs VIII K(a)l(endas) Noe(bres) / Lupicino et Iovino / in pace}. EDB (EDB13166) : texte concordant aux conventions de notation près (indice de similarité 0,925).}
\nf{Apparat}{⸢v⸣ : B gravé ; ⸢m⸣ : N gravé ; ⸢v⸣ : B gravé ; ⸢m⸣ : N gravé}
\nf{Traduction}{À Sosorus (?), qui vécut 14 ans 9 jours ; il mourut le 10e jour avant les calendes de novembre ; déposé le 8e jour avant les calendes de novembre, sous (le consulat de) Lupicinus et Iovinus, en paix.}
\nf{Langue}{L’éditeur signale 4 graphies corrigées : B gravé pour v (bêtacisme) ×2 ; N gravé pour m (n pour m) ×2.}
\nf{Chronologie}{Noms consulaires : Lupicinus (lu) et Iovinus (lu). Dates du jour : décès : X Kalendas Novembres = 23 octobre 367 (lu) ; déposition : VIII Kalendas Novembres = 25 octobre 367 (lu). Âge déclaré : 14 a. 9 j. (lu).}
\nf{Commentaire}{La formule in pace est conservée.}
\nf{Bibliographie}{ICUR II 6027 ; ILCV 2976a ; AE 1905, 72.}
\nf{Statut}{Lecture secondaire (EDCS) concordante avec l’EDB : lecture stabilisée entre deux bases, non collationnée sur l’ICVR ni sur le monument.}
\end{notice}

\begin{notice}{120}{367}{ICUR n.s. II, 6028 · EDCS-34200654 · EDB13167 · EDH HD030246}
\label{not:120}
\nf{Datation}{367 — directe : nom consulaire lu}
\nf{Genre}{Épitaphe ; ensembles A, B, C ; 69 lettres transmises.}
\nf{Références}{ICUR II 6028 = ILCV +3099.}
\nf{Provenance}{Rattachement selon l’EDB : \emph{Coem. Commodillae}.}
\nf{Support}{Plaque de marbre. Inscription gravée. Signes décrits par l’EDB : chrisme. Reproductions signalées par l’EDB : photographie PCAS ×2 (droits réservés ; non reproduites). Fiche EDB : Carlo Carletti, première version 2015.}
\ntext{\lnum{1}Zosimus recessit in pace\\
\lnum{2}annorum plus minus\\
\lnum{3}[- - -]III Idus Ma\\
\lnum{4}ias Lup\{p\}icin\\
\lnum{5}o et Io⸢v⸣ino cons(ulibus)}
\nf{Témoins}{Export EDCS de 2020 : \emph{Zosimus recessit in pace / annorum plus minus / [3]III Idus Ma/ias Lup\{p\}icin/o et Ioino cons(ulibus)}. EDB (EDB13167) : texte concordant aux conventions de notation près (indice de similarité 0,988). Le texte publié dans l'AE 1905, 71 (ancien n° 106 du catalogue) est le même monument : voir la notice 106. Leçon divergente du verbe : \emph{c(iu)esq(it!)} (AE) contre \emph{recessit} (ICVR, EDB). Seule l'ICVR II 6028, avec sa photographie, peut trancher. L'EDB (EDB13167) rattache la plaque à la catacombe de Commodilla et décrit un chrisme ; deux photographies de la Pontificia Commissione di Archeologia Sacra sont signalées. Absorbe la notice 106 (voir cette notice).}
\nf{Apparat}{⸢v⸣ : B gravé ; \{p\} : lettres superflues}
\nf{Traduction}{Zosimus s’en est allé en paix, âgé d’environ […] ans ; […] le 3e jour avant les ides de mai, sous le consulat de Lupicinus et de Iovinus.}
\nf{Langue}{L’éditeur signale une graphie corrigée : B gravé pour v (bêtacisme). Des lettres gravées sont jugées superflues par l’éditeur (entre accolades). 1 abréviation développée par l’éditeur (dont cons(ulibus)).}
\nf{Chronologie}{Noms consulaires : Lupicinus (lu) et Iovinus (lu). Dates du jour : décès : …III Idus Maias = 13 mai 367 (partiellement restitué).}
\nf{Commentaire}{La formule in pace est conservée. La variation d’édition ne peut être résolue par la seule similarité du texte.}
\nf{Bibliographie}{ICUR II 6028 ; ILCV +3099.}
\nf{Statut}{Lecture secondaire (EDCS) concordante avec l’EDB : lecture stabilisée entre deux bases, non collationnée sur l’ICVR ni sur le monument. L’EDCS signale une photographie (non consultée).}
\end{notice}

\begin{notice}{121}{367}{ICUR n.s. III, 6502 · EDCS-34801302 · EDB19879}
\label{not:121}
\nf{Datation}{367 — directe : lecture partielle d’un nom consulaire}
\nf{Genre}{Épitaphe ; ensembles A, B, C ; 55 lettres transmises ; 56 \% restituées.}
\nf{Références}{ICUR III 6502 = ILCV +2628.}
\nf{Provenance}{Rattachement selon l’EDB : \emph{Coem. Domitillae pars inferior}. Conservation selon l’EDB : \emph{Coem. Domitillae pars inferior}.}
\nf{Support}{Plaque de marbre. 29 × 26 × 1,5 cm (h. × l. × ép.). Inscription gravée. Reproductions signalées par l’EDB : photographie PCAS ×2 (droits réservés ; non reproduites). Fiche EDB : Carlo Carletti / Antonio Enrico Felle, première version 2003.}
\ntext{\lnum{1}f]ecit Bibi co[- - -]\\
\lnum{2}[- - - mere]nti i⸢n⸣ pace [- - -]\\
\lnum{3}[- - - feci]t Lup[icino et Iovino cons(ulibus)]}
\nf{Témoins}{Export EDCS de 2020 : \emph{f]ecit Bibi co[3] / [3 mere]nti i pace [3] / [3 feci]t Lup[icino et Iovino cons(ulibus)]}. EDB (EDB19879 ; écarts à examiner) : \emph{[--- f]ecit Bibi co[n --- / mere]nti im pace [---] / [--- feci]t Lupi[cino et Iovino cons(ulibus)] [texte tronqué dans la liste de résultats]}.}
\nf{Apparat}{⸢n⸣ : M gravé}
\nf{Traduction}{[…] a fait (pour) Bibus (?) […] ; […] bien méritant, en paix […] ; […] a fait, sous le consulat de Lupicinus et de Iovinus.}
\nf{Langue}{L’éditeur signale une graphie corrigée : M gravé pour n.}
\nf{Chronologie}{Noms consulaires : Lupicinus (lu en partie) et Iovinus (restitué).}
\nf{Commentaire}{La formule in pace est conservée.}
\nf{Bibliographie}{ICUR III 6502 ; ILCV +2628.}
\nf{Statut}{Lecture secondaire (EDCS) ; écarts avec l’EDB à arbitrer sur l’ICVR. Plus de la moitié des lettres sont restituées. L’EDCS signale une photographie (non consultée).}
\end{notice}

\begin{notice}{122}{367}{ICUR n.s. III, 8146 · EDCS-35700175 · EDB24669}
\label{not:122}
\nf{Datation}{367 — directe : nom consulaire lu}
\nf{Genre}{Épitaphe ; ensembles A, B, C ; 67 lettres transmises ; 40 \% restituées.}
\nf{Références}{ICUR III 8146.}
\nf{Provenance}{Rattachement selon l’EDB : \emph{Basilica ss. Nerei et Achillei in coem. Domitillae}. Conservation selon l’EDB : \emph{Basilica ss. Nerei et Achillei in coem. Domitillae}.}
\nf{Support}{Plaque de marbre. 29,5 × 28 × 3 cm (h. × l. × ép.). Inscription gravée. Reproductions signalées par l’EDB : photographie PCAS (droits réservés ; non reproduites). Fiche EDB : Carlo Carletti / Antonio Enrico Felle, première version 2003.}
\ntext{\lnum{1}] qu(a)e vixit anno[s - - -]\\
\lnum{2}[- - - de]pos(i)ta Nonis Iu[niis - - -]\\
\lnum{3}[- - - die Ma]rtis Lupicino et Io[vino conss(ulibus)}
\nf{Témoins}{EDB (EDB24669 ; écarts à examiner) : \emph{[---] que vixit anno[s ---] / [de]posta nonis iu[niis] / [die ma]rtis Lupicino et Io[vino cons(ulibu)s] [texte tronqué dans la liste de résultats]}.}
\nf{Traduction}{[…] qui vécut […] ans ; […] déposée aux nones de juin […] ; […] un mardi, sous le consulat de Lupicinus et de Iovinus.}
\nf{Langue}{3 abréviations développées par l’éditeur (dont qu(a)e, depos(i)ta, conss(ulibus)).}
\nf{Chronologie}{Noms consulaires : Lupicinus (lu) et Iovinus (restitué). Dates du jour : déposition : Nonis Iuniis = 5 juin 367 (partiellement restitué).}
\nf{Bibliographie}{ICUR III 8146. Bibliographie complémentaire selon l’EDB : IC I Supplementum, 1578.}
\nf{Statut}{Lecture secondaire (EDCS) ; écarts avec l’EDB à arbitrer sur l’ICVR. L’EDCS signale une photographie (non consultée).}
\end{notice}

\begin{notice}{123}{367}{ICUR n.s. III, 9344b · EDCS-35900809 · EDB24921}
\label{not:123}
\nf{Datation}{367 — conjecturale}
\nf{Genre}{Épitaphe ; ensembles A, B, C ; 19 lettres transmises ; 42 \% restituées.}
\nf{Références}{ICUR III 9344b.}
\nf{Provenance}{Rattachement selon l’EDB : \emph{Basilica ss. Nerei et Achillei in coem. Domitillae}.}
\nf{Support}{Plaque de marbre. Inscription gravée. Fiche EDB : Carlo Carletti / Antonio Enrico Felle, première version 2003.}
\ntext{\lnum{1}b]ene m[erenti - - -]\\
\lnum{2}[- - -]o et o⸢b⸣i[it - - -]}
\nf{Témoins}{Export EDCS de 2020 : \emph{b]ene m[erenti 3] / [3]o et oi[it 3]}. EDB (EDB24921 ; divergence forte) : \emph{[--- be]neme[renti ---] / [--- Lupicin?]o et ⟨I⟩ovi[no cons(ulibus?]}.}
\nf{Apparat}{⸢b⸣ : V gravé}
\nf{Traduction}{[…] bien méritant […] ; […] et mourut […].}
\nf{Langue}{L’éditeur signale une graphie corrigée : V gravé pour b (bêtacisme).}
\nf{Chronologie}{Le texte EDCS lit \emph{[---]o et o⸢b⸣i[it ---]} (« et mourut ») ; l'EDB (EDB24921 = ICVR IX 9344b) restitue au même endroit \emph{[--- Lupicin?]o et ⟨I⟩ovi[no cons(ulibus?)]} (367). Les deux lectures sont incompatibles. L'année 367 n'est qu'une hypothèse de l'EDB ; la datation est déclarée conjecturale.}
\nf{Bibliographie}{ICUR III 9344b.}
\nf{Statut}{Lecture secondaire (EDCS) ; divergence forte avec l’EDB : le texte ne doit pas être exploité sans contrôle de l’édition imprimée.}
\end{notice}

\begin{notice}{124}{367}{ICUR n.s. IV, 9563 · EDCS-36101092 · EDB18367}
\label{not:124}
\nf{Datation}{367 — directe : nom consulaire lu}
\nf{Genre}{Épitaphe ; ensembles A, B, C ; 106 lettres transmises.}
\nf{Références}{ICUR IV 9563 = ILCV 3253.}
\nf{Provenance}{Rattachement selon l’EDB : \emph{Coem. Callisti pars inferior}.}
\nf{Support}{Plaque de marbre. Inscription gravée. Signes décrits par l’EDB : chrisme. Reproductions signalées par l’EDB : photographie PCAS ×2 (droits réservés ; non reproduites). Fiche EDB : Antonio Enrico Felle, première version 2003.}
\ntext{\lnum{1}Iuliae bene merenti qu(a)e vixit ann(os)\\
\lnum{2}duo mens⸢e⸣s VI d(ies) XV dep(osita) XV Kal(endas) Iun(ias)\\
\lnum{3}Lupicino et Io⸢v⸣ino conss(ulibus) in p(ace)\\
\lnum{4}parentes Ursus et Seberane et sibi}
\nf{Témoins}{Export EDCS de 2020 : \emph{Iuliae bene merenti qu(a)e vixit ann(os) / duo menss VI d(ies) XV dep(osita) XV Kal(endas) Iun(ias) / Lupicino et Ioino conss(ulibus) in p(ace) / parentes Ursus et Seberane et sibi}. EDB (EDB18367) : texte concordant aux conventions de notation près (indice de similarité 0,981).}
\nf{Apparat}{⸢e⸣ : I gravé ; ⸢v⸣ : B gravé}
\nf{Traduction}{À Iulia, bien méritante, qui vécut 2 ans 6 mois 15 jours ; déposée le 15e jour avant les calendes de juin, sous le consulat de Lupicinus et de Iovinus, en paix ; ses parents Ursus et Severana, et pour eux-mêmes.}
\nf{Langue}{L’éditeur signale 2 graphies corrigées : I gravé pour e (confusion e/i) ; B gravé pour v (bêtacisme).}
\nf{Chronologie}{Noms consulaires : Lupicinus (lu) et Iovinus (lu). Dates du jour : déposition : XV Kalendas Iunias = 18 mai 367 (lu). Âge déclaré : 2 a. 6 m. 15 j. (lu).}
\nf{Commentaire}{La formule in pace est transmise sous forme abrégée, sans restitution de lacune ; son développement reste éditorial.}
\nf{Bibliographie}{ICUR IV 9563 ; ILCV 3253.}
\nf{Statut}{Lecture secondaire (EDCS) concordante avec l’EDB : lecture stabilisée entre deux bases, non collationnée sur l’ICVR ni sur le monument. L’EDCS signale une photographie (non consultée).}
\end{notice}

\annee{368}

\begin{notice}{125}{368}{ICUR n.s. I, 963 · EDCS-35200402 · EDB29591}
\label{not:125}
\nf{Datation}{368 — directe : nom consulaire lu}
\nf{Genre}{Épitaphe ; ensembles A, B, C ; 164 lettres transmises ; 21 \% restituées.}
\nf{Références}{CLE 664 = ILCV 2732.}
\nf{Provenance}{Rattachement selon l’EDB : \emph{Incertae originis} (le volume I des ICVR réunit les inscriptions dont la provenance est perdue).}
\nf{Support}{Support non précisé par l’EDB. Reproductions signalées par l’EDB : dessin, ICVR I (de Rossi) (droits réservés ; non reproduites). Fiche EDB : Filippo Piazzolla, première version 2014.}
\ntext{\lnum{1}Quie\{n\}sces in pace merens Leontia con[iunx]\\
\lnum{2}incompara⸢b⸣il⸢i⸣s qu(a)e vixit in s(a)eculo an[nos - - -]\\
\lnum{3}mens(es) VII d(ies) XXVIII Cresce(n)s maritus contra vo[tum posuit]\\
\lnum{4}d(eposita) IIII Non(as) Octrob(res) Valentiniano II et Va[lente II co(n)ss(ulibus)]}
\nf{Témoins}{Export EDCS de 2020 : \emph{Quie\{n\}sces in pace merens Leontia con[iux] / incomparails qu(a)e vixit in s(a)eculo an[nos 3] / mens(es) VII d(ies) XXVIII Cresce(n)s maritus contra vo[tum posuit] / d(eposita) IIII Non(as) Octrob(res) Valentiniano II et Va[lente II co(n)ss(ulibus)]}. EDB (EDB29591) : texte concordant aux conventions de notation près (indice de similarité 0,977).}
\nf{Apparat}{⸢b⸣ : V gravé ; ⸢i⸣ : E gravé ; \{n\} : lettres superflues}
\nf{Traduction}{Tu reposeras en paix, (toi qui le) mérites, Leontia, épouse incomparable, qui vécut dans le siècle […] ans 7 mois 28 jours ; Crescens, son mari, a posé (ce monument) contre son vœu ; déposée le 4e jour avant les nones d’octobre, sous le consulat de Valentinien pour la 2e fois et de Valens pour la 2e fois.}
\nf{Langue}{L’éditeur signale 2 graphies corrigées : V gravé pour b (bêtacisme) ; E gravé pour i (confusion e/i). L’export de 2020 diffère de l’état 2022 (voir apparat). Des lettres gravées sont jugées superflues par l’éditeur (entre accolades). 10 abréviations développées par l’éditeur (dont qu(a)e, s(a)eculo, mens(es), d(ies), Cresce(n)s, d(eposita)).}
\nf{Chronologie}{Noms consulaires : Valentinianus (lu) et Valens (lu). Ne sont pas des consuls : Leontia con[iunx] (défunte). Dates du jour : déposition : IIII Nonas Octrobres = 4 octobre 368 (lu). Âge déclaré : 7 m. 28 j. (partiellement restitué).}
\nf{Commentaire}{La formule in pace est conservée. Vocabulaire religieux : repos et sommeil de la mort — Quiesces (lu).}
\nf{Bibliographie}{CLE 664 ; ILCV 2732. Bibliographie complémentaire selon l’EDB : ILCV, 2732.}
\nf{Statut}{Lecture secondaire (EDCS) concordante avec l’EDB : lecture stabilisée entre deux bases, non collationnée sur l’ICVR ni sur le monument. L’EDCS signale une photographie (non consultée).}
\end{notice}

\begin{notice}{126}{368}{ICUR n.s. I, 479 · EDCS-38700887 · EDB28428}
\label{not:126}
\nf{Datation}{368 — directe : nom consulaire lu}
\nf{Genre}{Épitaphe ; ensembles A, B, C ; 271 lettres transmises.}
\nf{Références}{ICUR I 479 = ILCV 3650 = AE 1892, 31.}
\nf{Provenance}{Rattachement selon l’EDB : \emph{Incertae originis} (le volume I des ICVR réunit les inscriptions dont la provenance est perdue). Conservation selon l’EDB : \emph{s. Maria in Cappella}.}
\nf{Support}{Plaque de marbre. 60 × 29 cm (h. × l.). Inscription gravée. Fiche EDB : Filippo Piazzolla, première version 2014.}
\ntext{\lnum{1}M(a)erentes parentes fecerunt quod\\
\lnum{2}ipse debuit facere infelices parentes\\
\lnum{3}dolo debere et hunc tumulum fecerunt\\
\lnum{4}qui natus est Fl(avio) Sergio et Fl(avio) Nigriniano conss(ulibus)\\
\lnum{5}III Idus Aug(ustas) die Saturnis depositus in hac\\
\lnum{6}domo aeterna dd(ominis) nn(ostris) Valentiniano et Valente\\
\lnum{7}Aug(ustis) iterum conss(ulibus) VII Idus Dece⸢m⸣b(res) die Saturnis\\
\lnum{8}qui vixit ann(os) XVIIII m(enses) III d(ies) XXVI\\
\lnum{9}in pace}
\nf{Témoins}{Export EDCS de 2020 : \emph{M(a)erentes parentes fecerunt quod / ipse debuit facere infelices parentes / dolo debere et hunc tumulum fecerunt / qui natus est Fl(avio) Sergio et Fl(avio) Nigriniano conss(ulibus) / III Idus Aug(ustas) die Saturnis depositus in hac / domo aeterna dd(ominis) nn(ostris) Valentiniano et Valente / Aug(ustis) iterum conss(ulibus) VII Idus Deceb(res) die Saturnis / qui vixit ann(os) XVIIII m(enses) I…}. EDB (EDB28428) : texte concordant aux conventions de notation près (indice de similarité 0,997).}
\nf{Apparat}{⸢m⸣ : N gravé}
\nf{Traduction}{Ses parents affligés ont fait ce qu’il aurait dû faire lui-même ; malheureux parents, (contraints) par la douleur au devoir, ils ont fait aussi ce tombeau ; il naquit sous le consulat de Flavius Sergius et de Flavius Nigrinianus, le 3e jour avant les ides d’août, un samedi ; il a été déposé dans cette demeure éternelle sous le consulat de nos seigneurs Valentinien et Valens Augustes pour la 2e fois, le 7e jour avant les ides de décembre, un samedi ; il vécut 19 ans 3 mois 26 jours ; en paix.}
\nf{Langue}{L’éditeur signale une graphie corrigée : N gravé pour m (n pour m).}
\nf{Chronologie}{Noms consulaires : Valentinianus (lu), Valens (lu), Sergius (lu) et Nigrinianus (lu). Dates du jour : événement indéterminé : III Idus Augustas = 11 août 368 (lu) ; événement indéterminé : VII Idus Decembres = 7 décembre 368 (lu). Âge déclaré : 19 a. 3 m. 26 j. (lu).}
\nf{Commentaire}{La formule in pace est conservée. Commentaire développé : dossier p.~\pageref{dos:126}.}
\nf{Bibliographie}{ICUR I 479 ; ILCV 3650 ; AE 1892, 31. Bibliographie complémentaire selon l’EDB : ILCV, 3650.}
\nf{Statut}{Lecture secondaire (EDCS) concordante avec l’EDB : lecture stabilisée entre deux bases, non collationnée sur l’ICVR ni sur le monument.}
\end{notice}

\begin{notice}{127}{368}{ICUR n.s. I, 725 · EDCS-38701126 · EDB28952}
\label{not:127}
\nf{Datation}{368 — directe : nom consulaire lu}
\nf{Genre}{Épitaphe ; ensembles A, B, C ; 118 lettres transmises ; 19 \% restituées.}
\nf{Références}{ICUR I 725.}
\nf{Provenance}{Rattachement selon l’EDB : \emph{Incertae originis} (le volume I des ICVR réunit les inscriptions dont la provenance est perdue).}
\nf{Support}{Support non précisé par l’EDB. Reproductions signalées par l’EDB : dessin, ICVR I (de Rossi) (droits réservés ; non reproduites). Fiche EDB : Filippo Piazzolla, première version 2014.}
\ntext{\lnum{1}Dulcissim(a)e filiae Victori[nae]\\
\lnum{2}bene merenti in pace qu(a)e vix[it]\\
\lnum{3}ann(os) p(lus) m(inus) XXV d(ies) XXVI dep(osita) Kal(endis) Feb(ruariis) [cons(ulatu)]\\
\lnum{4}dd(ominorum) nn(ostrorum) Valentiniani et [Valentis]\\
\lnum{5}Augg(ustorum) iterum}
\nf{Témoins}{EDB (EDB28952) : texte concordant aux conventions de notation près (indice de similarité 0,949).}
\nf{Traduction}{À sa fille très douce Victorina, bien méritante, en paix, qui vécut environ 25 ans 26 jours ; déposée aux calendes de février, sous le consulat de nos seigneurs Valentinien et Valens Augustes pour la 2e fois.}
\nf{Langue}{13 abréviations développées par l’éditeur (dont Dulcissim(a)e, qu(a)e, ann(os), p(lus), m(inus), d(ies)).}
\nf{Chronologie}{Noms consulaires : Valentinianus (lu). Dates du jour : déposition : Kalendis Februariis = 1 février 368 (lu). Âge déclaré : env. 25 a. 26 j. (partiellement restitué).}
\nf{Commentaire}{La formule in pace est conservée.}
\nf{Bibliographie}{ICUR I 725. Bibliographie complémentaire selon l’EDB : ILCV, 2603 adn.}
\nf{Statut}{Lecture secondaire (EDCS) concordante avec l’EDB : lecture stabilisée entre deux bases, non collationnée sur l’ICVR ni sur le monument. L’EDCS signale une photographie (non consultée).}
\end{notice}

\begin{notice}{128}{368}{ICUR n.s. I, 2808 · EDCS-11700574 · EDB32731}
\label{not:128}
\nf{Datation}{368 — directe : nom consulaire lu}
\nf{Genre}{Épitaphe ; ensembles A, B, C ; 101 lettres transmises.}
\nf{Références}{ICUR I 2808 = ILCV 4393a = ILMN-01, 00443.}
\nf{Provenance}{Rattachement selon l’EDB : \emph{Incertae originis} (le volume I des ICVR réunit les inscriptions dont la provenance est perdue). Conservation selon l’EDB : \emph{Napoli   ; Museo Archeologico Nazionale}.}
\nf{Support}{Plaque de marbre. 15,5 × 95 × 3 cm (h. × l. × ép.). Inscription gravée. Signes décrits par l’EDB : chrisme. Reproductions signalées par l’EDB : dessin, ICVR I (de Rossi) et reproduction, Camodeca et alii 2000 (droits réservés ; non reproduites). Fiche EDB : Filippo Piazzolla, première version 2015.}
\ntext{\lnum{1}⸢V⸣alentiniano et ⸢V⸣alenti iterum Balerius(!)\\
\lnum{2}qui vixit annos VI m(enses) XI d(ies) XVI depositus\\
\lnum{3}est ⸢XV⸣ ⸢K⸣al(endas) Dec(embres) d(ie) Martis bene\\
\lnum{4}merenti in pace}
\nf{Témoins}{Export EDCS de 2020 : \emph{alentiniano et alenti iterum Balerius(!) / qui vixit annos VI m(enses) XI d(ies) XVI depositus / est al(endas) Dec(embres) d(ie) Martis bene / merenti in pace}. EDB (EDB32731) : texte concordant aux conventions de notation près (indice de similarité 0,927).}
\nf{Apparat}{⸢V⸣ : B gravé ; ⸢V⸣ : B gravé ; ⸢XV⸣ : VX gravé ; ⸢K⸣ : C gravé ; (!) : forme signalée par l’éditeur}
\nf{Traduction}{Sous (le consulat de) Valentinien et Valens pour la 2e fois, Valerius, qui vécut 6 ans 11 mois 16 jours ; il a été déposé le 15e jour avant les calendes de décembre, un mardi ; bien méritant, en paix.}
\nf{Langue}{L’éditeur signale 4 graphies corrigées : B gravé pour V (bêtacisme) ×2 ; VX gravé pour XV ; C gravé pour K. L’éditeur signale au moins une forme anormale (sic). 6 abréviations développées par l’éditeur (dont Balerius(!), m(enses), d(ies), al(endas), Dec(embres), d(ie)).}
\nf{Chronologie}{Noms consulaires : Valentinianus (lu) et Valens (lu). Dates du jour : déposition : XV Kalendas Decembres = 17 novembre 368 (lu). Âge déclaré : 6 a. 11 m. 16 j. (lu).}
\nf{Commentaire}{La formule in pace est conservée.}
\nf{Bibliographie}{ICUR I 2808 ; ILCV 4393a ; ILMN-01, 00443. Bibliographie complémentaire selon l’EDB : ILCV, 4393A ; Camodeca et alii 2000, 152-153, n. 443 con foto ; Ferrua 1979b, 80.}
\nf{Statut}{Lecture secondaire (EDCS) concordante avec l’EDB : lecture stabilisée entre deux bases, non collationnée sur l’ICVR ni sur le monument. L’EDCS signale une photographie (non consultée).}
\end{notice}

\begin{notice}{129}{368}{ICUR n.s. II, 4500 · EDCS-34100563 · EDB20916}
\label{not:129}
\nf{Datation}{368 — directe : consulat restitué}
\nf{Genre}{Épitaphe ; ensembles A, B, C ; 57 lettres transmises ; 88 \% restituées.}
\nf{Références}{ICUR II 4500.}
\nf{Provenance}{Rattachement selon l’EDB : \emph{Coem. Pontiani}.}
\nf{Support}{Plaque de marbre. Inscription gravée. Fiche EDB : Antonio Enrico Felle, première version 2003.}
\ntext{\lnum{1}M]ai[as(?) - - - dd(ominis) nn(ostris) Valentiniano II]\\
\lnum{2}[et Valen]te II [conss(ulibus)]}
\nf{Témoins}{EDB (EDB20916) : texte concordant aux conventions de notation près (indice de similarité 0,98).}
\nf{Traduction}{[…] mai […] sous le consulat de nos seigneurs Valentinien pour la 2e fois et Valens pour la 2e fois.}
\nf{Langue}{4 abréviations développées par l’éditeur (dont Maias(?), dd(ominis), nn(ostris), conss(ulibus)).}
\nf{Chronologie}{Noms consulaires : Valentinianus (restitué) et Valens (restitué).}
\nf{Bibliographie}{ICUR II 4500.}
\nf{Statut}{Lecture secondaire (EDCS) concordante avec l’EDB : lecture stabilisée entre deux bases, non collationnée sur l’ICVR ni sur le monument. Plus de la moitié des lettres sont restituées.}
\end{notice}

\begin{notice}{130}{368}{ICUR n.s. II, 5931 · EDCS-38100789 · EDB21154}
\label{not:130}
\nf{Datation}{368 — directe : nom consulaire lu}
\nf{Genre}{Épitaphe ; ensembles A, B, C ; 114 lettres transmises ; 44 \% restituées.}
\nf{Références}{ICUR II 5931 = ILCV 3797b (em).}
\nf{Provenance}{Rattachement selon l’EDB : \emph{Museum paulianum}. Conservation selon l’EDB : \emph{Basilica s.Pauli apostoli}.}
\nf{Support}{Plaque de marbre. 82 × 77 cm (h. × l.). Inscription gravée. Reproductions signalées par l’EDB : dessin, ICVR I (de Rossi) (droits réservés ; non reproduites). Fiche EDB : Maria Cristina Grisanzio, première version 2006.}
\ntext{\lnum{1}[- - -] in pace qui vi⸢x⸣it\\
\lnum{2}[annos - - - et sibi feci]t locu(m) cum c⸢o⸣mpare\\
\lnum{3}[sua - - -]t Maura\\
\lnum{4}bisomu(m)\\
\lnum{5}[- - - dep(ositus) - - - I]ulias dd(ominis) nn(ostris) Valentiniano II\\
\lnum{6}[et Valente II conss(ulibus)]}
\nf{Témoins}{Export EDCS de 2020 : \emph{[3] in pace qui viit / [annos 3 et sibi feci]t locu(m) cum cmpare / [sua 3]t Maura / bisomu(m) / [3 dep(ositus) 3 I]ulias dd(ominis) nn(ostris) Valentiniano II / [et Valente II conss(ulibus)]}. EDB (EDB21154) : texte concordant aux conventions de notation près (indice de similarité 0,937).}
\nf{Apparat}{⸢x⸣ : S gravé ; ⸢o⸣ : V gravé}
\nf{Traduction}{[…] en paix, qui vécut […] ans ; et il s’est fait un emplacement avec sa compagne […] Maura ; tombe double ; […] déposé […] juillet, sous le consulat de nos seigneurs Valentinien pour la 2e fois et Valens pour la 2e fois.}
\nf{Langue}{L’éditeur signale 2 graphies corrigées : S gravé pour x (s pour x) ; V gravé pour o (confusion o/u).}
\nf{Chronologie}{Noms consulaires : Valentinianus (lu) et Valens (lu).}
\nf{Commentaire}{La formule in pace est conservée.}
\nf{Bibliographie}{ICUR II 5931 ; ILCV 3797b (em).}
\nf{Statut}{Lecture secondaire (EDCS) concordante avec l’EDB : lecture stabilisée entre deux bases, non collationnée sur l’ICVR ni sur le monument. L’EDCS signale une photographie (non consultée).}
\end{notice}

\begin{notice}{131}{368}{ICUR n.s. III, 8147 · EDCS-35700176 · EDB24670}
\label{not:131}
\nf{Datation}{368 — directe : nom consulaire lu}
\nf{Genre}{Épitaphe ; ensembles A, B, C ; 85 lettres transmises ; 13 \% restituées.}
\nf{Références}{ICUR III 8147 = ILCV 4399a.}
\nf{Provenance}{Rattachement selon l’EDB : \emph{Basilica ss. Nerei et Achillei in coem. Domitillae}. Conservation selon l’EDB : \emph{Basilica ss. Nerei et Achillei in coem. Domitillae}.}
\nf{Support}{Plaque de marbre. 46 × 86 × 5,5 cm (h. × l. × ép.). Inscription gravée. Reproductions signalées par l’EDB : photographie PCAS (droits réservés ; non reproduites). Fiche EDB : Carlo Carletti / Antonio Enrico Felle, première version 2003.}
\ntext{\lnum{1}Gelasio bene merenti qui vix\{s\}it [- - - depo]\\
\lnum{2}situs prid(i)e Nonas Martias die Ve[neris]\\
\lnum{3}Valen[t]iniano [e]t Fl(avio) Valen⸢te⸣}
\nf{Témoins}{Export EDCS de 2020 : \emph{Gelasio bene merenti qui vix\{s\}it [3 depo]/situs prid(i)e Nonas Martias die Ve[neris] / Valen[t]iniano [e]t Fl(avio) Valen}. EDB (EDB24670 ; écarts à examiner) : \emph{Gelasino benemerenti qui vixsit [--- depo]/situs prid⟨i⟩e nonas martias die Ve[neris] / Valen⌜t⌝iniano e̲t Fl(avius) Valens [---?] [texte tronqué dans la liste de résultats]}.}
\nf{Apparat}{⸢te⸣ : S gravé ; \{s\} : lettres superflues}
\nf{Traduction}{À Gelasius, bien méritant, qui vécut […] ; déposé la veille des nones de mars, un vendredi, sous (le consulat de) Valentinien et Flavius Valens.}
\nf{Langue}{L’éditeur signale une graphie corrigée : S gravé pour te. Des lettres gravées sont jugées superflues par l’éditeur (entre accolades). 2 abréviations développées par l’éditeur (dont prid(i)e, Fl(avio)).}
\nf{Chronologie}{Noms consulaires : Valentinianus (lu en partie) et Valens (lu). Dates du jour : déposition : pridie Nonas Martias = 6 mars 368 (lu).}
\nf{Bibliographie}{ICUR III 8147 ; ILCV 4399a.}
\nf{Statut}{Lecture secondaire (EDCS) ; écarts avec l’EDB à arbitrer sur l’ICVR. L’EDCS signale une photographie (non consultée).}
\end{notice}

\annee{369}

\begin{notice}{132}{369}{EDCS-29900387}
\label{not:132}
\nf{Datation}{369 — directe : nom consulaire lu}
\nf{Genre}{Épitaphe ; ensembles A, B, C ; 111 lettres transmises.}
\nf{Références}{AE 2003, 216.}
\nf{Provenance}{Rome (indication de l’EDCS) ; aucun rattachement cimétérial transmis.}
\nf{Support}{Non documenté dans les sources consultées.}
\ntext{\lnum{1}Criscen⸢t⸣in(a)e quae vi⸢x⸣it\\
\lnum{2}annis n(umero) XXX m(ensibus) VI d(iebus) VII be\\
\lnum{3}ne merenti in pace frater\\
\lnum{4}eius Priscus fecit\\
\lnum{5}d(eposita) V Non(as) Oct(obres) Valentiniano n(obilissimo) p(uero) et Victore conss(ulibus)}
\nf{Témoins}{Export EDCS de 2020 : \emph{Criscenin(a)e quae viit / annis n(umero) XXX m(ensibus) VI d(iebus) VII be/ne merenti in pace frater / eius Priscus fecit / d(eposita) V Non(as) Oct(obres) Valentiniano n(obilissimo) p(uero) et Victore conss(ulibus)}. Aucune fiche EDB appariée (pas de numéro ICUR ; recherche par le texte infructueuse).}
\nf{Apparat}{⸢t⸣ : I gravé ; ⸢x⸣ : S gravé}
\nf{Traduction}{À Crescentina, qui vécut 30 ans 6 mois 7 jours ; bien méritante, en paix ; son frère Priscus a fait (ce monument) ; déposée le 5e jour avant les nones d’octobre, sous le consulat de Valentinien, très noble enfant, et de Victor.}
\nf{Langue}{L’éditeur signale 2 graphies corrigées : I gravé pour t ; S gravé pour x (s pour x).}
\nf{Chronologie}{Noms consulaires : Valentinianus (lu) et Victor (lu). Dates du jour : déposition : V Nonas Octobres = 3 octobre 369 (lu). Âge déclaré : 30 a. 6 m. 7 j. (lu).}
\nf{Commentaire}{La formule in pace est conservée.}
\nf{Bibliographie}{AE 2003, 216.}
\nf{Statut}{Lecture secondaire (EDCS) sans second témoin numérique.}
\end{notice}

\begin{notice}{133}{369}{ICUR n.s. II, 4165 · EDCS-33900077 · EDB19224}
\label{not:133}
\nf{Datation}{369 — directe : nom consulaire lu}
\nf{Genre}{Épitaphe ; ensembles A, B, C ; 272 lettres transmises.}
\nf{Références}{ICUR II 4165 = ILCV 1603 (add, em).}
\nf{Provenance}{Rattachement selon l’EDB : \emph{Basilica S.Petri apostoli}.}
\nf{Support}{Support non précisé par l’EDB. Reproductions signalées par l’EDB : dessin, ICVR I (de Rossi) (droits réservés ; non reproduites). Fiche EDB : Antonio Enrico Felle, première version 2003.}
\ntext{\lnum{1}Felix sanctae fidei vocitus iit in pace cuius tantus amor et caritas retenetur ab amicis in aevo\\
\lnum{2}qui cum esset fuit solacius misericors omnibus notus\\
\lnum{3}Agrippina fecit dulcissimo suo marito\\
\lnum{4}cum qu⸢o⸣ vixit sine l(a)esione animi annos III et m(enses) X\\
\lnum{5}fuit in saeculum quo(a)d vixit annos XXXII dep(ositus) XIII Kal(endas) Sept(embres) Valentiniano\\
\lnum{6}n(obilissimo) p(uero) et Victori conss(ulibus)}
\nf{Témoins}{Export EDCS de 2020 : \emph{Felix sanctae fidei vocitus iit in pace cuius tantus amor et caritas retenetur ab amicis in aevo / qui cum esset fuit solacius misericors omnibus notus / Agrippina fecit dulcissimo suo marito / cum qu vixit sine l(a)esione animi annos III et m(enses) X / fuit in saeculum quo(a)d vixit annos XXXII dep(ositus) XIII Kal(endas) Sept(embres) Valentiniano / n(obilissimo) p(uero) et Victori conss(ulibus)}. EDB (EDB19224 ; écarts à examiner) : \emph{Felix sanctae fidei vocitus iit in pace / cuius tantus amor et caritas retenetur ab amicis in aevo / qui cum esset fuit solacius misericors omnibus notus / Agrippina fecit dulcissimo suo marito /cum quem vixit sine lesione animi annos III det m(enses) X / fuit in saeculum quod vixit annos XXXII dep(ositus) XIII kal(endas) sept(embres) Valentiniano n(obilissimo) p(uero) et Victori conss(ulibus)}.}
\nf{Apparat}{⸢o⸣ : EM gravé}
\nf{Traduction}{Felix, appelé à la sainte foi, s’en est allé en paix ; un si grand amour et une si grande affection sont gardés pour lui par ses amis pour toujours ; lui qui, de son vivant, fut un réconfort, miséricordieux, connu de tous. Agrippina a fait (ce monument) pour son très doux mari, avec qui elle vécut sans blessure d’âme 3 ans et 10 mois ; il fut dans le siècle tant qu’il vécut 32 ans ; déposé le 13e jour avant les calendes de septembre, sous le consulat de Valentinien, très noble enfant, et de Victor.}
\nf{Langue}{L’éditeur signale une graphie corrigée : EM gravé pour o.}
\nf{Chronologie}{Noms consulaires : Valentinianus (lu). Dates du jour : déposition : XIII Kalendas Septembres = 20 août 369 (lu). Âge déclaré : 32 a. (lu).}
\nf{Commentaire}{La formule in pace est conservée. Vocabulaire religieux : épithète du défunt (non doctrinal) — sanctae (lu) ; initiation et foi — fidei (lu). Commentaire développé : dossier p.~\pageref{dos:133}.}
\nf{Bibliographie}{ICUR II 4165 ; ILCV 1603 (add, em).}
\nf{Statut}{Lecture secondaire (EDCS) ; écarts avec l’EDB à arbitrer sur l’ICVR. L’EDCS signale une photographie (non consultée).}
\end{notice}

\begin{notice}{134}{369}{ICUR n.s. III, 8148 · EDCS-35700177 · EDB24671}
\label{not:134}
\nf{Datation}{369 — directe : nom consulaire lu}
\nf{Genre}{Épitaphe ; ensembles A, B, C ; 172 lettres transmises ; 24 \% restituées.}
\nf{Références}{ICUR III 8148 = ILCV 1128 (add, em) = AnalEpi p 101 = AE 1912, 261 = Mossong-2022, 00129.}
\nf{Provenance}{Rattachement selon l’EDB : \emph{Basilica ss. Nerei et Achillei in coem. Domitillae}. Conservation selon l’EDB : \emph{Basilica ss. Nerei et Achillei in coem. Domitillae}.}
\nf{Support}{Plaque de marbre. 105 × 190 × 15 cm (h. × l. × ép.). Inscription gravée. Reproductions signalées par l’EDB : photographie PCAS (droits réservés ; non reproduites). Fiche EDB : Carlo Carletti / Antonio Enrico Felle, première version 2015.}
\ntext{\lnum{1}Dormitio Dulciti presb(yteri)\\
\lnum{2}Kal(endis) Feb(ruariis) et Greg[ori] diac(oni)\\
\lnum{3}XIII Kal(endas) Mai(as) et Ca[t]afro\\
\lnum{4}ni(a)es(!) s[o]roris eo[r(um)] V Id(us) Mar(tias)\\
\lnum{5}d(e)[p(ositio) F]aust[i]nae IIII [K(a)l(endas)] Nov(embres)\\
\lnum{6}Valen[tinia]no n(obilissimo) p(uero) et Victore [conss(ulibus)]\\
\lnum{7}dormitio Agapenis IIII N(onas) Oc(tobres) in p(ace)}
\nf{Témoins}{EDB (EDB24671 ; divergence forte) : \emph{⟨a:⟩ dormitio Dulciti presb(yteri) / kal(endis) feb(ruariis) et Greg[ori] diac(oni) / XIII kal(endas) mai(as) et Ca[l]ifro/nies s[o]roris eo[rum] V id(us) mar(tias) / / ⟨b:⟩ d(e)[p(ositio) F]aust[i]nae IIII [kal(endas)] nov(embres) / Valen[tinia]no n(obilissimo) p(uero) et Victore [cons(ulibu)s] / / ⟨c:⟩ dormitio Agapenis n(onis) oc(tobris) in p(ace) [texte tronqué dans la liste de résultats]}. L'EDB (EDB24671) décrit une grande dalle de marbre (105 × 190 × 15 cm) portant plusieurs textes (a, b, c), remployée en morceaux dans un mur (Marucchi 1914, cité par l'EDB) ; une photographie PCAS est signalée. Le nom de la sœur est lu \emph{Ca[t]afronies} par l'EDCS, \emph{Ca[l]ifronies} par l'EDB. L'EDB date l'ensemble « 350–399 ; 369 » : l'année 369 n'est portée que par la partie b (déposition de Faustina, sous Valentinien \emph{nobilissimus puer} et Victor). Les dates des autres défunts ne sont pas nécessairement de la même année.}
\nf{Apparat}{(!) : forme signalée par l’éditeur}
\nf{Traduction}{Endormissement de Dulcitius, prêtre, aux calendes de février, et de Gregorius, diacre, le 13e jour avant les calendes de mai, et de Catafronia, leur sœur, le 5e jour avant les ides de mars ; déposition de Faustina le 4e jour avant les calendes de novembre, sous le consulat de Valentinien, très noble enfant, et de Victor ; endormissement d’Agapè le 4e jour avant les nones d’octobre, en paix.}
\nf{Langue}{L’éditeur signale au moins une forme anormale (sic). 21 abréviations développées par l’éditeur (dont presb(yteri), Kal(endis), Feb(ruariis), diac(oni), Kal(endas), Mai(as)).}
\nf{Chronologie}{Noms consulaires : Valentinianus (lu en partie) et Victor (lu). Dates du jour : événement indéterminé : Kalendis Februariis = 1 février 369 (lu) ; événement indéterminé : XIII Kalendas Maias = 19 avril 369 (lu) ; événement indéterminé : V Idus Martias = 11 mars 369 (lu) ; déposition : IIII Kalendas Novembres = 29 octobre 369 (partiellement restitué) ; événement indéterminé : IIII Nonas Octobres = 4 octobre 369 (lu).}
\nf{Commentaire}{La formule in pace est transmise sous forme abrégée, sans restitution de lacune ; son développement reste éditorial. Vocabulaire religieux : repos et sommeil de la mort — Dormitio (lu), dormitio (lu). Fonction ou métier mentionné : presbyteri, diaconi. Commentaire développé : dossier p.~\pageref{dos:134}.}
\nf{Bibliographie}{ICUR III 8148 ; ILCV 1128 (add, em) ; AnalEpi p 101 ; AE 1912, 261 ; Mossong-2022, 00129.}
\nf{Statut}{Lecture secondaire (EDCS) ; divergence forte avec l’EDB : le texte ne doit pas être exploité sans contrôle de l’édition imprimée. L’EDCS signale une photographie (non consultée).}
\end{notice}

\begin{notice}{135}{369}{ICUR n.s. IV, 11101 · EDCS-37200383 · EDB39560}
\label{not:135}
\nf{Datation}{369 — directe : nom consulaire lu}
\nf{Genre}{Épitaphe ; ensembles A, B, C ; 89 lettres transmises ; 4 \% restituées.}
\nf{Références}{ICUR IV 11101 = ILCV 4213a.}
\nf{Provenance}{Rattachement selon l’EDB : \emph{Coem. Callisti area subdialis et incerti loci}. Conservation selon l’EDB : \emph{Coem. Callisti area subdialis et incerti loci}.}
\nf{Support}{La fiche EDB ne figure pas dans l’état 2019 des données EpiDoc ; support non documenté.}
\ntext{\lnum{1}]NECI[- - -]\\
\lnum{2}[- - -] si merita feli[x - - -]\\
\lnum{3}[- - -] hoc fecisset Fe[- - -]\\
\lnum{4}[- - -]o suo Leopardo\\
\lnum{5}[- - - X]XXVII fecit mecum a(n)n(os) XV[- - -]\\
\lnum{6}[- - - Va]lentiniano ⸢d⸣(omino) n(ostro) et Victore cons(ulibus) [}
\nf{Témoins}{Export EDCS de 2020 : \emph{]NECI[3] / [3] si merita feli[x 3] / [3] hoc fecisset Fe[3] / [3]o suo Leopardo / [3 X]XXVII fecit mecum a(n)n(os) XV[3] / [3 Va]lentiniano (omino) n(ostro) et Victore cons(ulibus) [}. EDB (EDB39560) : texte concordant aux conventions de notation près (indice de similarité 0,956).}
\nf{Apparat}{⸢d⸣ : P gravé}
\nf{Traduction}{[…] si heureux par ses mérites […] ; […] avait fait […] ; […] à son […] Leopardus ; […] 37 (ans), il a vécu avec moi 15 ans […] ; sous le consulat de notre seigneur Valentinien et de Victor.}
\nf{Langue}{L’éditeur signale une graphie corrigée : P gravé pour d.}
\nf{Chronologie}{Noms consulaires : Valentinianus (lu en partie) et Victor (lu).}
\nf{Bibliographie}{ICUR IV 11101 ; ILCV 4213a.}
\nf{Statut}{Lecture secondaire (EDCS) concordante avec l’EDB : lecture stabilisée entre deux bases, non collationnée sur l’ICVR ni sur le monument. L’EDCS signale une photographie (non consultée).}
\end{notice}

\annee{370}

\begin{notice}{136}{370}{ICUR n.s. I, 3179 · EDCS-40300325 · EDB33131}
\label{not:136}
\nf{Datation}{370 — directe : nom consulaire lu}
\nf{Genre}{Épitaphe ; ensembles A, B, C ; 74 lettres transmises.}
\nf{Références}{CIL XIII *381 = ICUR I 3179.}
\nf{Provenance}{Rattachement selon l’EDB : \emph{Incertae originis} (le volume I des ICVR réunit les inscriptions dont la provenance est perdue).}
\nf{Support}{Support non précisé par l’EDB. Signes décrits par l’EDB : chrisme et chrisme avec alpha et oméga. Fiche EDB : Filippo Piazzolla, première version 2015.}
\ntext{\lnum{1}Vixit Secundi\\
\lnum{2}nus ann(os) XXX et\\
\lnum{3}def(unctus) die VI Kal(endas)\\
\lnum{4}Apri(les) dd(ominis) nn(ostris) Vale\\
\lnum{5}ntiniano et Va\\
\lnum{6}len⸢t⸣e Augg(ustis) ter\\
\lnum{7}cons(ulibus)}
\nf{Témoins}{Export EDCS de 2020 : \emph{Vixit Secundi/nus ann(os) XXX et / def(unctus) die VI Kal(endas) / Apri(les) dd(ominis) nn(ostris) Vale/ntiniano et Va/lene Augg(ustis) ter / cons(ulibus)}. EDB (EDB33131) : texte concordant aux conventions de notation près (indice de similarité 0,945).}
\nf{Apparat}{⸢t⸣ : I gravé}
\nf{Traduction}{Secundinus vécut 30 ans et mourut le 6e jour avant les calendes d’avril, sous le consulat de nos seigneurs Valentinien et Valens Augustes pour la 3e fois.}
\nf{Langue}{L’éditeur signale une graphie corrigée : I gravé pour t.}
\nf{Chronologie}{Noms consulaires : Valentinianus (lu) et Valens (lu). Dates du jour : décès : die VI Kalendas Apriles = 27 mars 370 (lu). Âge déclaré : 30 a. 6 j. (lu).}
\nf{Bibliographie}{CIL XIII *381 ; ICUR I 3179. Bibliographie complémentaire selon l’EDB : CIL XIII, 40\textbackslash{}*, 381 ; ILCV, 2795B adn.}
\nf{Statut}{Lecture secondaire (EDCS) concordante avec l’EDB : lecture stabilisée entre deux bases, non collationnée sur l’ICVR ni sur le monument.}
\end{notice}

\begin{notice}{137}{370}{ICUR n.s. I, 442 · EDCS-38600420 · EDB28273}
\label{not:137}
\nf{Datation}{370 — directe : consulat restitué}
\nf{Genre}{Mixte (épitaphe et acte d’achat) ; ensembles A, B, C ; 99 lettres transmises ; 46 \% restituées.}
\nf{Références}{ICUR I 442.}
\nf{Provenance}{Rattachement selon l’EDB : \emph{Incertae originis} (le volume I des ICVR réunit les inscriptions dont la provenance est perdue).}
\nf{Support}{Plaque de marbre. Inscription gravée. Fiche EDB : Filippo Piazzolla, première version 2014.}
\ntext{\lnum{1}[- - -] et Artemia\\
\lnum{2}[locum se vivi si]bi cum suis\\
\lnum{3}[comparaverunt] d(e)p(osita) Artemia XIIII Kal(endas) Sept(embres)\\
\lnum{4}[Valentin(iano) III et Val]e(n)t(is) III Augg(ustis) conss(ulibus)}
\nf{Témoins}{EDB (EDB28273) : texte concordant aux conventions de notation près (indice de similarité 0,967).}
\nf{Traduction}{[…] et Artemia se sont acquis de leur vivant un emplacement pour eux et les leurs ; Artemia a été déposée le 14e jour avant les calendes de septembre, sous le consulat de Valentinien pour la 3e fois et de Valens pour la 3e fois, Augustes.}
\nf{Langue}{9 abréviations développées par l’éditeur (dont d(e)p, (osita), Kal(endas), Sept(embres), Valentin(iano), Vale(n)t).}
\nf{Chronologie}{Noms consulaires : Valentinianus (restitué). Dates du jour : déposition : XIIII Kalendas Septembres = 19 août 370 (lu).}
\nf{Bibliographie}{ICUR I 442. Bibliographie complémentaire selon l’EDB : ILCV, 3734 A adn.}
\nf{Statut}{Lecture secondaire (EDCS) concordante avec l’EDB : lecture stabilisée entre deux bases, non collationnée sur l’ICVR ni sur le monument.}
\end{notice}

\begin{notice}{138}{370}{ICUR n.s. I, 1351 · EDCS-39000329 · EDB30683}
\label{not:138}
\nf{Datation}{370 — directe : consulat restitué}
\nf{Genre}{Épitaphe ; ensembles A, B, C ; 68 lettres transmises ; 62 \% restituées.}
\nf{Références}{ICUR I 1351.}
\nf{Provenance}{Rattachement selon l’EDB : \emph{Incertae originis} (le volume I des ICVR réunit les inscriptions dont la provenance est perdue). Conservation selon l’EDB : \emph{Collegio Teutonico}.}
\nf{Support}{Plaque de marbre. 20 × 21 cm (h. × l.). Inscription gravée. Fiche EDB : Filippo Piazzolla, première version 2013.}
\ntext{\lnum{1}]IIA[- - -]\\
\lnum{2}[- - - Io]vina de[p(osita) - - -]\\
\lnum{3}[- - - Valentiniano Aug(usto) III et Va]l(ente) Aug(usto) III co[ns(ulibus) - - -]\\
\lnum{4}[- - -] Constant[- - -]}
\nf{Témoins}{EDB (EDB30683 ; écarts à examiner) : \emph{------ / [---]++ạ[---] / [--- Io]vina de[posita ---] / [dd(ominis) nn(ostris) Valentiniano et Va]lente Aug(ustis) III co[ns(ulibu)s ---] / [---] C̣onstant[---] / ------ [texte tronqué dans la liste de résultats]}.}
\nf{Traduction}{[…] Iovina, déposée […] sous le consulat de Valentinien Auguste pour la 3e fois et de Valens Auguste pour la 3e fois […] ; Constant[…].}
\nf{Langue}{5 abréviations développées par l’éditeur (dont dep(osita), Aug(usto), Val(ente), cons(ulibus)).}
\nf{Chronologie}{Noms consulaires : Valentinianus (restitué), Valens (restitué), Constantius (lu) et Constans (lu).}
\nf{Bibliographie}{ICUR I 1351. Bibliographie complémentaire selon l’EDB : CLRE, 645.}
\nf{Statut}{Lecture secondaire (EDCS) ; écarts avec l’EDB à arbitrer sur l’ICVR. Plus de la moitié des lettres sont restituées.}
\end{notice}

\begin{notice}{139}{370}{ICUR n.s. I, 3180 · EDCS-40300326 · EDB33132}
\label{not:139}
\nf{Datation}{370 — directe : nom consulaire lu}
\nf{Genre}{Épitaphe ; ensembles A, B, C ; 104 lettres transmises ; 38 \% restituées.}
\nf{Références}{ICUR I 3180.}
\nf{Provenance}{Rattachement selon l’EDB : \emph{Incertae originis} (le volume I des ICVR réunit les inscriptions dont la provenance est perdue).}
\nf{Support}{Support non précisé par l’EDB. Signes décrits par l’EDB : chrisme. Fiche EDB : Filippo Piazzolla, première version 2015.}
\ntext{\lnum{1}Bene mer(enti) Mart[yrius anima]\\
\lnum{2}dulcis qui vix[it ann(os) - - -]\\
\lnum{3}d(efunctus) in pace d(epositus) Non(is) Se[ptembribus do]\\
\lnum{4}minis nostris V[alen]\\
\lnum{5}tin(iano) et Val(ente) III cons[s(ulibus)}
\nf{Témoins}{EDB (EDB33132 ; écarts à examiner) : \emph{((monogramma Christi)) / benemer(enti) Mart[yrius anima?] / dulcis qui vix[it annis ---] / d(epositus) in pace d(ie) non(as) se[ptembres do]/minis nostris V[alen]/tin(iano) et Va⌜l⌝(ente) ⌜A⌝A(u)⌜gg⌝(ustis) III cons(ulibu)[s ---] [texte tronqué dans la liste de résultats]}.}
\nf{Traduction}{Au bien méritant Martyrius, âme douce, qui vécut […] ans ; mort en paix, déposé aux nones de septembre, sous (le consulat de) nos seigneurs Valentinien et Valens pour la 3e fois.}
\nf{Langue}{8 abréviations développées par l’éditeur (dont mer(enti), ann(os), d(efunctus), d(epositus), Non(is), tin(iano)).}
\nf{Chronologie}{Noms consulaires : Valentinianus (lu en partie) et Valens (lu). Dates du jour : déposition : Nonis Septembribus = 5 septembre 370 (partiellement restitué).}
\nf{Commentaire}{La formule in pace est conservée.}
\nf{Bibliographie}{ICUR I 3180.}
\nf{Statut}{Lecture secondaire (EDCS) ; écarts avec l’EDB à arbitrer sur l’ICVR.}
\end{notice}

\begin{notice}{140}{370}{ICUR n.s. II, 4501 · EDCS-34100564 · EDB20917}
\label{not:140}
\nf{Datation}{370 — directe : nom consulaire lu}
\nf{Genre}{Épitaphe ; ensembles A, B, C ; 100 lettres transmises.}
\nf{Références}{ICUR II 4501 = ILCV 2603a.}
\nf{Provenance}{Rattachement selon l’EDB : \emph{Coem. Pontiani}.}
\nf{Support}{Support non précisé par l’EDB. Signes décrits par l’EDB : chrisme. Fiche EDB : Antonio Enrico Felle, première version 2003.}
\ntext{\lnum{1}Bonifati(a)e be\\
\lnum{2}ne merenti\\
\lnum{3}in pace Chr(isti)\\
\lnum{4}qu(a)e vixit ann\\
\lnum{5}is XXVIII deposi\\
\lnum{6}ta in pace XII Kal(endas)\\
\lnum{7}Sept(embres) Valentiniano\\
\lnum{8}et Valente III cons(ulibus)}
\nf{Témoins}{EDB (EDB20917) : texte concordant aux conventions de notation près (indice de similarité 0,992).}
\nf{Traduction}{À Bonifatia, bien méritante, dans la paix du Christ, qui vécut 28 ans ; déposée en paix le 12e jour avant les calendes de septembre, sous le consulat de Valentinien et Valens pour la 3e fois.}
\nf{Langue}{6 abréviations développées par l’éditeur (dont Bonifati(a)e, Chr(isti), qu(a)e, Kal(endas), Sept(embres), cons(ulibus)).}
\nf{Chronologie}{Noms consulaires : Valentinianus (lu) et Valens (lu). Dates du jour : déposition : XII Kalendas Septembres = 21 août 370 (lu). Âge déclaré : 28 a. (lu).}
\nf{Commentaire}{La formule in pace est conservée. Vocabulaire religieux : désignation du Christ — Christi (abréviation développée). L'EDCS et l'EDB transcrivent \emph{in pace Chr(isti)}, mais la fiche EpiDoc de l'EDB (EDB20917) décrit un chrisme parmi les données matérielles. Il est probable que \emph{Chr(isti)} soit ici le développement éditorial d'un monogramme employé comme abréviation, comme au n° 38, où l'EDB le transcrit au contraire comme signe. Le texte est maintenu dans le profil lexical large, mais exclu du profil strict, qui ne retient que les désignations écrites en lettres. Désignation divine : retenue dans le profil restrictif. Commentaire développé : dossier p.~\pageref{dos:140}.}
\nf{Bibliographie}{ICUR II 4501 ; ILCV 2603a.}
\nf{Statut}{Lecture secondaire (EDCS) concordante avec l’EDB : lecture stabilisée entre deux bases, non collationnée sur l’ICVR ni sur le monument. L’EDCS signale une photographie (non consultée).}
\end{notice}

\begin{notice}{141}{370}{ICUR n.s. II, 6030 · EDCS-34200656 · EDB9633}
\label{not:141}
\nf{Datation}{370 — directe : nom consulaire lu ; datation contestée}
\nf{Genre}{Acte d’achat ou de réservation ; ensembles A, B ; 48 lettres transmises.}
\nf{Références}{ICUR II 6030 = ILCV 3808 (em).}
\nf{Provenance}{Rattachement selon l’EDB : \emph{Coem. Commodillae}.}
\nf{Support}{Plaque de marbre. Inscription gravée. Fiche EDB : Carlo Carletti, première version 2003.}
\ntext{\lnum{1}Concordius se\\
\lnum{2}vivo locum bisomum\\
\lnum{3}emit c(onsule) Valentiniano III}
\nf{Témoins}{EDB (EDB9633) : texte concordant aux conventions de notation près (indice de similarité 0,941).}
\nf{Traduction}{Concordius a acheté de son vivant un emplacement double, sous le consulat de Valentinien pour la 3e fois.}
\nf{Langue}{1 abréviation développée par l’éditeur (dont c(onsule)).}
\nf{Chronologie}{Noms consulaires : Valentinianus (lu). Valentinien III : Valentinien Ier (370) ou Valentinien II (387) ; la CLRE (p. 644, citée par l’EDB) note « Year doubtful: 387 possible ».}
\nf{Bibliographie}{ICUR II 6030 ; ILCV 3808 (em). Bibliographie complémentaire selon l’EDB : ILCV, 3808 ; CLRE, 644.}
\nf{Statut}{Lecture secondaire (EDCS) concordante avec l’EDB : lecture stabilisée entre deux bases, non collationnée sur l’ICVR ni sur le monument.}
\end{notice}

\begin{notice}{142}{370}{ICUR n.s. IV, 12421 · EDCS-37600945 · EDB41696}
\label{not:142}
\nf{Datation}{370 — directe : nom consulaire lu}
\nf{Genre}{Épitaphe ; ensembles A, B, C ; 85 lettres transmises ; 4 \% restituées.}
\nf{Références}{ICUR IV 12421 = ILCV 1297 (em).}
\nf{Provenance}{Rattachement selon l’EDB : \emph{inter vias Appiam et Ardeatinam}. Conservation selon l’EDB : \emph{Musei Vaticani: Lapidario Cristiano}.}
\nf{Support}{La fiche EDB ne figure pas dans l’état 2019 des données EpiDoc ; support non documenté.}
\ntext{\lnum{1}Puer Rufin[us]\\
\lnum{2}qui vixit anni[s]\\
\lnum{3}XVIIII d(ies) LVIIII d(epositus) prid(ie)\\
\lnum{4}Idus Oct(obres) dd(ominis) nn(ostris) ⸢V⸣a\\
\lnum{5}lentiniano et Valente\\
\lnum{6}Aug(ustis) III conss(ulibus)}
\nf{Témoins}{Export EDCS de 2020 : \emph{Puer Rufin[us] / qui vixit anni[s] / XVIIII d(ies) LVIIII d(epositus) prid(ie) / Idus Oct(obres) dd(ominis) nn(ostris) a/lentiniano et Valente / Aug(ustis) III conss(ulibus)}. EDB (EDB41696) : texte concordant aux conventions de notation près (indice de similarité 0,991).}
\nf{Apparat}{⸢V⸣ : Y gravé}
\nf{Traduction}{Le garçon Rufinus, qui vécut 19 (?) ans, 59 jours ; déposé la veille des ides d’octobre, sous le consulat de nos seigneurs Valentinien et Valens Augustes pour la 3e fois.}
\nf{Langue}{L’éditeur signale une graphie corrigée : Y gravé pour V.}
\nf{Chronologie}{Noms consulaires : Valentinianus (lu) et Valens (lu). Ne sont pas des consuls : Puer Rufin[us] (défunt). Dates du jour : déposition : Idus Octobres = 15 octobre 370 (lu). Âge déclaré : 19 a. 59 j. (partiellement restitué).}
\nf{Commentaire}{L’EDB transcrit les signes suivants : chrisme.}
\nf{Bibliographie}{ICUR IV 12421 ; ILCV 1297 (em). Bibliographie complémentaire selon l’EDB : ILCV, 1297 ; Vilella 2019b, 147-148.}
\nf{Statut}{Lecture secondaire (EDCS) concordante avec l’EDB : lecture stabilisée entre deux bases, non collationnée sur l’ICVR ni sur le monument.}
\end{notice}

\annee{371}

\begin{notice}{143}{371}{ICUR n.s. I, 1432 · EDCS-39000415 · EDB31196}
\label{not:143}
\nf{Datation}{371 — directe : nom consulaire lu}
\nf{Genre}{Épitaphe ; ensembles A, B, C ; 101 lettres transmises.}
\nf{Références}{ICUR I 1432.}
\nf{Provenance}{Rattachement selon l’EDB : \emph{Incertae originis} (le volume I des ICVR réunit les inscriptions dont la provenance est perdue). Conservation selon l’EDB : \emph{Musei Vaticani: Lapidario Cristiano}.}
\nf{Support}{Plaque de marbre. 24 × 35 cm (h. × l.). Inscription gravée. Signes décrits par l’EDB : chrisme et chrisme avec alpha et oméga. Fiche EDB : Filippo Piazzolla, première version 2015.}
\ntext{\lnum{1}]\\
\lnum{2}sia [- - -]\\
\lnum{3}dulcis(s)ima fili(a)e bene me\\
\lnum{4}re(n)ti qu(a)e vix(i)t an(nos) VI m(enses) VI\\
\lnum{5}d(ies) VIII deposita in pa\\
\lnum{6}ce pri(di)e Kal(endas) Maias\\
\lnum{7}depos(ita) cons⸢u⸣\\
\lnum{8}latu Gratiani\\
\lnum{9}et Probi}
\nf{Témoins}{Export EDCS de 2020 : \emph{]/sia [3] / dulcis(s)ima fili(a)e bene me/re(n)ti qu(a)e vix(i)t an(nos) VI m(enses) VI / d(ies) VIII deposita in pa/ce pri(di)e Kal(endas) Maias / depos(ita) cons/latu Gratiani / et Probi}. EDB (EDB31196 ; écarts à examiner) : \emph{------ / [-c.1-2-]sia [---] / [---] dulcisima filie beneme/re⟨n⟩ti que vix⟨i⟩t an(nis) VI m̅(ensibus) VI / d̅(iebus) VIII deposita in pa/ce pri(di)e kal(endas) maias / ((monogramma Christi cum alpha et omega)) depos(uit) conso/latu Gratiani / et Probi [texte tronqué dans la liste de résultats]}.}
\nf{Apparat}{⸢u⸣ : O gravé}
\nf{Traduction}{[…] à sa fille très douce et bien méritante, qui vécut 6 ans 6 mois 8 jours ; déposée en paix la veille des calendes de mai, sous le consulat de Gratien et de Probus.}
\nf{Langue}{L’éditeur signale une graphie corrigée : O gravé pour u (confusion o/u).}
\nf{Chronologie}{Noms consulaires : Gratianus (lu) et Probus (lu). Dates du jour : déposition : pridie Kalendas Maias = 30 avril 371 (lu). Âge déclaré : 6 a. 6 m. 8 j. (lu).}
\nf{Commentaire}{La formule in pace est conservée.}
\nf{Bibliographie}{ICUR I 1432. Bibliographie complémentaire selon l’EDB : ILCV, 2968 adn. ; Lega 2000, 104 ; Ferrua 1979b, 37, n. 1432 ; Ferrua 1964, 108.}
\nf{Statut}{Lecture secondaire (EDCS) ; écarts avec l’EDB à arbitrer sur l’ICVR.}
\end{notice}

\begin{notice}{144}{371}{ICUR n.s. I, 1936 · EDCS-35300743 · EDB35584}
\label{not:144}
\nf{Datation}{371 — directe : nom consulaire lu}
\nf{Genre}{Épitaphe ; ensembles A, B, C ; 100 lettres transmises.}
\nf{Références}{ICUR I 1936 = ICUR VIII 23061 = ILCV 2874a.}
\nf{Provenance}{Rattachement selon l’EDB : \emph{Inscriptiones priscillianae}. Conservation selon l’EDB : \emph{Museo Archeologico Nazionale}.}
\nf{Support}{Plaque de marbre. 26 × 43 × 3,5 cm (h. × l. × ép.). Inscription gravée. Fiche EDB : Filippo Piazzolla, première version 2015.}
\ntext{\lnum{1}Maxim(a)e bene me\\
\lnum{2}renti qu(a)e vixit ann(os)\\
\lnum{3}p(lus) m(inus) LXII i⸢v⸣it in pace VII Id\\
\lnum{4}us Mart(ias) d(omino) n(ostro) Gratian\\
\lnum{5}o et Petronio Probo cons(ulibus)\\
\lnum{6}maritus fecit sibi [}
\nf{Témoins}{Export EDCS de 2020 : \emph{Maxim(a)e bene me/renti qu(a)e vixit ann(os) / p(lus) m(inus) LXII iit in pace VII Id/us Mart(ias) d(omino) n(ostro) Gratian/o et Petronio Probo cons(ulibus) / maritus fecit sibi [}. EDB (EDB35584) : texte concordant aux conventions de notation près (indice de similarité 0,947).}
\nf{Apparat}{⸢v⸣ : B gravé}
\nf{Traduction}{À Maxima, bien méritante, qui vécut environ 62 ans ; elle s’en alla en paix le 7e jour avant les ides de mars, sous le consulat de notre seigneur Gratien et de Petronius Probus ; son mari a fait (ce monument) pour lui aussi.}
\nf{Langue}{L’éditeur signale une graphie corrigée : B gravé pour v (bêtacisme).}
\nf{Chronologie}{Noms consulaires : Gratianus (lu) et Probus (lu). Dates du jour : décès : VII Idus Martias = 9 mars 371 (lu). Âge déclaré : env. 62 a. (lu).}
\nf{Commentaire}{La formule in pace est conservée.}
\nf{Bibliographie}{ICUR I 1936 ; ICUR VIII 23061 ; ILCV 2874a. Bibliographie complémentaire selon l’EDB : ILCV, 2874A ; ICVR VIII, 23061.}
\nf{Statut}{Lecture secondaire (EDCS) concordante avec l’EDB : lecture stabilisée entre deux bases, non collationnée sur l’ICVR ni sur le monument.}
\end{notice}

\begin{notice}{145}{371}{ICUR n.s. I, 2088 · EDCS-39300312 · EDB32053}
\label{not:145}
\nf{Datation}{371 — directe : nom consulaire lu}
\nf{Genre}{Épitaphe ; ensembles A, B, C ; 95 lettres transmises.}
\nf{Références}{ICUR I 2088 = ILCV 1478a.}
\nf{Provenance}{Rattachement selon l’EDB : \emph{Incertae originis} (le volume I des ICVR réunit les inscriptions dont la provenance est perdue). Conservation selon l’EDB : \emph{Musei Vaticani: Galleria Lapidaria}.}
\nf{Support}{Plaque de marbre. 21 × 33 cm (h. × l.). Inscription gravée. Reproductions signalées par l’EDB : reproduction, Perret 1851 (droits réservés ; non reproduites). Fiche EDB : Filippo Piazzolla, première version 2015.}
\ntext{\lnum{1}Romano neof⸢y⸣to\\
\lnum{2}bene merenti qui vi\\
\lnum{3}xit annos VIIII d(ies) XV\\
\lnum{4}requiescit in pace d(omino) n(ostro)\\
\lnum{5}Fl(avio) Gratiano Aug(usto) II et\\
\lnum{6}Petronio Probo c(on)s(ulibus)}
\nf{Témoins}{Export EDCS de 2020 : \emph{Romano neofto / bene merenti qui vi/xit annos VIIII d(ies) XV / requiescit in pace d(omino) n(ostro) / Fl(avio) Gratiano Aug(usto) II et / Petronio Probo c(on)s(ulibus)}. EDB (EDB32053) : texte concordant aux conventions de notation près (indice de similarité 0,951).}
\nf{Apparat}{⸢y⸣ : I gravé}
\nf{Traduction}{À Romanus, néophyte, bien méritant, qui vécut 9 ans 15 jours ; il repose en paix ; sous le consulat de notre seigneur Flavius Gratien Auguste pour la 2e fois et de Petronius Probus.}
\nf{Langue}{L’éditeur signale une graphie corrigée : I gravé pour y (graphie de y).}
\nf{Chronologie}{Noms consulaires : Gratianus (lu) et Probus (lu). Âge déclaré : 9 a. 15 j. (lu).}
\nf{Commentaire}{La formule in pace est conservée. Vocabulaire religieux : initiation et foi — neofyto (lu) ; repos et sommeil de la mort — requiescit (lu).}
\nf{Bibliographie}{ICUR I 2088 ; ILCV 1478a. Bibliographie complémentaire selon l’EDB : ILCV, 1478A ; Lega 2000.}
\nf{Statut}{Lecture secondaire (EDCS) concordante avec l’EDB : lecture stabilisée entre deux bases, non collationnée sur l’ICVR ni sur le monument. L’EDCS signale une photographie (non consultée).}
\end{notice}

\begin{notice}{146}{371}{ICUR n.s. I, 2809 · EDCS-11700575 · EDB32732}
\label{not:146}
\nf{Datation}{371 — directe : nom consulaire lu}
\nf{Genre}{Épitaphe ; ensembles A, B, C ; 108 lettres transmises.}
\nf{Références}{ICUR I 2809 (p 491) = ILCV 4351 = ILMN-01, 00444.}
\nf{Provenance}{Rattachement selon l’EDB : \emph{Incertae originis} (le volume I des ICVR réunit les inscriptions dont la provenance est perdue). Conservation selon l’EDB : \emph{Napoli   ; Museo Archeologico Nazionale}.}
\nf{Support}{Plaque de marbre. 29,5 × 36,5 × 5 cm (h. × l. × ép.). Inscription gravée. Signes décrits par l’EDB : chrisme et chrisme avec alpha et oméga. Reproductions signalées par l’EDB : reproduction, Camodeca et alii 2000 (droits réservés ; non reproduites). Fiche EDB : Filippo Piazzolla, première version 2015.}
\ntext{\lnum{1}Mire sapientiae\\
\lnum{2}Augendo qui vixit\\
\lnum{3}ann(os) plus min(us) LXXII\\
\lnum{4}cum uxore fecit\\
\lnum{5}ann(os) XXX depositus\\
\lnum{6}XVI Kal(endas) Octob(res) d(omino) n(ostro) Gra\\
\lnum{7}tiano Aug(usto) II et Probo con(sulibus)}
\nf{Témoins}{EDB (EDB32732) : texte concordant aux conventions de notation près (indice de similarité 0,961).}
\nf{Traduction}{D’une admirable sagesse, à Augendus, qui vécut environ 72 ans, qui a vécu avec son épouse 30 ans ; déposé le 16e jour avant les calendes d’octobre, sous le consulat de notre seigneur Gratien Auguste pour la 2e fois et de Probus.}
\nf{Langue}{9 abréviations développées par l’éditeur (dont ann(os), min(us), Kal(endas), Octob(res), d(omino)).}
\nf{Chronologie}{Noms consulaires : Gratianus (lu) et Probus (lu). Dates du jour : déposition : XVI Kalendas Octobres = 16 septembre 371 (lu). Âge déclaré : env. 72 a. (lu).}
\nf{Bibliographie}{ICUR I 2809 (p 491) ; ILCV 4351 ; ILMN-01, 00444. Bibliographie complémentaire selon l’EDB : ILCV, 4351 ; Camodeca et alii 2000, 153, n. 444 con foto.}
\nf{Statut}{Lecture secondaire (EDCS) concordante avec l’EDB : lecture stabilisée entre deux bases, non collationnée sur l’ICVR ni sur le monument. L’EDCS signale une photographie (non consultée).}
\end{notice}

\begin{notice}{147}{371}{ICUR n.s. I, 3181 · EDCS-40300327 · EDB33133}
\label{not:147}
\nf{Datation}{371 — directe : nom consulaire lu}
\nf{Genre}{Épitaphe ; ensembles A, B, C ; 104 lettres transmises.}
\nf{Références}{ICUR I 3181 = ILCV 4258.}
\nf{Provenance}{Rattachement selon l’EDB : \emph{Incertae originis} (le volume I des ICVR réunit les inscriptions dont la provenance est perdue).}
\nf{Support}{Support non précisé par l’EDB. Signes décrits par l’EDB : chrisme. Fiche EDB : Filippo Piazzolla, première version 2015.}
\ntext{\lnum{1}Gratiano et Probo co(n)s(ulibus)\\
\lnum{2}tert(ium) Non(as) Feb(ruarias) ego Cons\\
\lnum{3}tantius compari meae\\
\lnum{4}Agrige qu(a)e vixit\\
\lnum{5}ann(os) LX et mecum\\
\lnum{6}fecit ann(os) XL m(enses)\\
\lnum{7}VIII d(ies) VIII\\
\lnum{8}XIIII}
\nf{Témoins}{EDB (EDB33133) : texte concordant aux conventions de notation près (indice de similarité 0,956).}
\nf{Traduction}{Sous le consulat de Gratien et de Probus, le 3e jour avant les nones de février, moi Constantius, à ma compagne Agrigia (?), qui vécut 60 ans et a vécu avec moi 40 ans 8 mois 8 jours ; 14 [chiffre isolé].}
\nf{Langue}{10 abréviations développées par l’éditeur (dont co(n)s, (ulibus), tert(ium), Non(as), Feb(ruarias), qu(a)e).}
\nf{Chronologie}{Noms consulaires : Gratianus (lu) et Probus (lu). Ne sont pas des consuls : ego Cons/tantius compari meae (dédicant (époux)). Dates du jour : événement indéterminé : tertium Nonas Februarias = 3 février 371 (lu). Âge déclaré : 60 a. 8 m. 8 j. (lu).}
\nf{Bibliographie}{ICUR I 3181 ; ILCV 4258. Bibliographie complémentaire selon l’EDB : ILCV, 4258 ; Ferrua 1979b, 91.}
\nf{Statut}{Lecture secondaire (EDCS) concordante avec l’EDB : lecture stabilisée entre deux bases, non collationnée sur l’ICVR ni sur le monument.}
\end{notice}

\begin{notice}{148}{371}{ICUR n.s. II, 4502 · EDCS-34100565 · EDB20918}
\label{not:148}
\nf{Datation}{371 — directe : lecture partielle d’un nom consulaire}
\nf{Genre}{Épitaphe ; ensembles A, B, C ; 62 lettres transmises ; 60 \% restituées.}
\nf{Références}{ICUR II 4502.}
\nf{Provenance}{Rattachement selon l’EDB : \emph{Coem. Pontiani}.}
\nf{Support}{Plaque de marbre. Inscription gravée. Reproductions signalées par l’EDB : photographie PCAS (droits réservés ; non reproduites). Fiche EDB : Antonio Enrico Felle, première version 2003.}
\ntext{\lnum{1}]V[- - -]\\
\lnum{2}[- - -]ntia\\
\lnum{3}[- - -]irm[- - -]\\
\lnum{4}[- - - hun]c locum [- - -]\\
\lnum{5}[- - - mecu]m fecit [- - -]\\
\lnum{6}[Gratiano Aug(usto) II et P]robo c[onss(ulibus) - - -]}
\nf{Témoins}{EDB (EDB20918) : texte concordant aux conventions de notation près (indice de similarité 0,938).}
\nf{Traduction}{[…] cet emplacement […] a vécu avec moi […] ; sous le consulat de Gratien Auguste pour la 2e fois et de Probus […].}
\nf{Langue}{2 abréviations développées par l’éditeur (dont Aug(usto), conss(ulibus)).}
\nf{Chronologie}{Noms consulaires : Gratianus (restitué) et Probus (lu en partie).}
\nf{Bibliographie}{ICUR II 4502.}
\nf{Statut}{Lecture secondaire (EDCS) concordante avec l’EDB : lecture stabilisée entre deux bases, non collationnée sur l’ICVR ni sur le monument. Plus de la moitié des lettres sont restituées. L’EDCS signale une photographie (non consultée).}
\end{notice}

\begin{notice}{149}{371}{ICUR n.s. II, 5712 · EDCS-38302410}
\label{not:149}
\nf{Datation}{371 — directe : nom consulaire lu}
\nf{Genre}{Épitaphe ; ensembles A, B, C ; 65 lettres transmises ; 28 \% restituées.}
\nf{Références}{ICUR II 5712.}
\nf{Provenance}{Rome (indication de l’EDCS) ; aucun rattachement cimétérial transmis.}
\nf{Support}{Non documenté dans les sources consultées.}
\ntext{\lnum{1}[- - -] Fl(avius) NINIRIT[- - -]\\
\lnum{2}[- - - q]ui vixit anno[s - - -]\\
\lnum{3}[- - - d]epositus d[ie - - -]\\
\lnum{4}[- - - De]c(embres) Gr(a)t(iano) Aug(usto) II\\
\lnum{5}[et Petronio P]robo conss(ulibus)}
\nf{Témoins}{Aucune fiche EDB appariée (numéro ICUR absent de la liste publique de l’EDB).}
\nf{Traduction}{[…] Flavius […] qui vécut […] ans ; […] déposé le […] décembre, sous le consulat de Gratien Auguste pour la 2e fois et de Petronius Probus.}
\nf{Langue}{6 abréviations développées par l’éditeur (dont Fl(avius), Dec(embres), Gr(a)t, (iano), Aug(usto), conss(ulibus)).}
\nf{Chronologie}{Noms consulaires : Gratianus (lu) et Probus (lu en partie).}
\nf{Bibliographie}{ICUR II 5712.}
\nf{Statut}{Lecture secondaire (EDCS) sans second témoin numérique. L’EDCS signale une photographie (non consultée).}
\end{notice}

\begin{notice}{150}{371}{ICUR n.s. II, 5933 · EDCS-38100790 · EDB21458}
\label{not:150}
\nf{Datation}{371 — directe : nom consulaire lu}
\nf{Genre}{Épitaphe ; ensembles A, B, C ; 31 lettres transmises ; 10 \% restituées.}
\nf{Références}{ICUR II 5933 (p 373) = ILCV 3813a.}
\nf{Provenance}{Rattachement selon l’EDB : \emph{Museum paulianum}. Conservation selon l’EDB : \emph{Basilica s.Pauli apostoli}.}
\nf{Support}{Plaque de marbre. 87 × 34 cm (h. × l.). Inscription gravée. Fiche EDB : Maria Cristina Grisanzio, première version 2006.}
\ntext{\lnum{1}dd(ominis) nn(ostris)] Grat(iano) et Pro(bo) i⸢n⸣ p(ace)\\
\lnum{2}[- - - bis]omu(m) locum sum(ptu) p(roprio)}
\nf{Témoins}{Export EDCS de 2020 : \emph{dd(ominis) nn(ostris)] Grat(iano) et Pro(bo) i p(ace) / [3 bis]omu(m) locum sum(ptu) p(roprio)}. EDB (EDB21458) : texte concordant aux conventions de notation près (indice de similarité 0,977).}
\nf{Apparat}{⸢n⸣ : M gravé}
\nf{Traduction}{Sous (le consulat de) nos seigneurs Gratien et Probus, en paix […] ; un emplacement double, à ses propres frais.}
\nf{Langue}{L’éditeur signale une graphie corrigée : M gravé pour n.}
\nf{Chronologie}{Noms consulaires : Gratianus (lu) et Probus (lu).}
\nf{Commentaire}{La formule in pace est transmise sous forme abrégée, sans restitution de lacune ; son développement reste éditorial.}
\nf{Bibliographie}{ICUR II 5933 (p 373) ; ILCV 3813a.}
\nf{Statut}{Lecture secondaire (EDCS) concordante avec l’EDB : lecture stabilisée entre deux bases, non collationnée sur l’ICVR ni sur le monument. L’EDCS signale une photographie (non consultée).}
\end{notice}

\begin{notice}{151}{371}{ICUR n.s. II, 6031 · EDCS-34200657 · EDB17432}
\label{not:151}
\nf{Datation}{371 — directe : nom consulaire lu}
\nf{Genre}{Épitaphe ; ensembles A, B, C ; 102 lettres transmises.}
\nf{Références}{ICUR II 6031 = ILCV 3253a (em).}
\nf{Provenance}{Rattachement selon l’EDB : \emph{Coem. Commodillae}. Conservation selon l’EDB : \emph{Musei Vaticani: Lapidario Cristiano}.}
\nf{Support}{Plaque de marbre. Inscription gravée. Signes décrits par l’EDB : chrisme. Reproductions signalées par l’EDB : reproduction, Perret 1851 (droits réservés ; non reproduites). Fiche EDB : Carlo Carletti, première version 2003.}
\ntext{\lnum{1}Hic iacet Muscula quae et Galatea\\
\lnum{2}quae vix(it) ann(is) duob(us) mens(ibus) duob(us) et d(iebus) XVII\\
\lnum{3}dep(osita) XV Kal(endas) Aug(ustas) Gratiano Aug(usto) II ⸢e⸣t\\
\lnum{4}Probo cons(ulibu)s in pace Chr(isti)}
\nf{Témoins}{Export EDCS de 2020 : \emph{Hic iacet Muscula quae et Galata / quae vix(it) ann(is) duob(us) mens(ibus) duob(us) et d(iebus) XVII / dep(osita) XV Kal(endas) Aug(ustas) Gratiano Aug(usto) II t / Probo cons(ulibu)s in pace Chr(isti)}. EDB (EDB17432) : texte concordant aux conventions de notation près (indice de similarité 1).}
\nf{Apparat}{⸢e⸣ : I gravé}
\nf{Traduction}{Ici gît Muscula, dite aussi Galatea, qui vécut 2 ans 2 mois et 17 jours ; déposée le 15e jour avant les calendes d’août, sous le consulat de Gratien Auguste pour la 2e fois et de Probus, dans la paix du Christ.}
\nf{Langue}{L’éditeur signale une graphie corrigée : I gravé pour e (confusion e/i). L’export de 2020 diffère de l’état 2022 (voir apparat). 12 abréviations développées par l’éditeur (dont vix(it), ann(is), duob(us), mens(ibus), d(iebus)).}
\nf{Chronologie}{Noms consulaires : Gratianus (lu) et Probus (lu). Dates du jour : déposition : XV Kalendas Augustas = 18 juillet 371 (lu). Âge déclaré : 2 a. 2 m. 17 j. (lu).}
\nf{Commentaire}{La formule in pace est conservée. Vocabulaire religieux : désignation du Christ — Christi (abréviation développée). Même situation qu'au n° 140 : \emph{in pace Chr(isti)} est transcrit par les deux bases, et l'EDB (EDB17432) décrit un chrisme sur le support. La formule clôt le texte, après la date consulaire. Maintenu dans le profil large, exclu du profil strict. Désignation divine : retenue dans le profil restrictif. Commentaire développé : dossier p.~\pageref{dos:151}.}
\nf{Bibliographie}{ICUR II 6031 ; ILCV 3253a (em).}
\nf{Statut}{Lecture secondaire (EDCS) concordante avec l’EDB : lecture stabilisée entre deux bases, non collationnée sur l’ICVR ni sur le monument. L’EDCS signale une photographie (non consultée).}
\end{notice}

\begin{notice}{152}{371}{ICUR n.s. IV, 11102 · EDCS-37200384 · EDB39561}
\label{not:152}
\nf{Datation}{371 — directe : nom consulaire lu}
\nf{Genre}{Épitaphe ; ensembles A, B, C ; 177 lettres transmises ; 28 \% restituées.}
\nf{Références}{ICUR IV 11102 = ILCV 1614 (em).}
\nf{Provenance}{Rattachement selon l’EDB : \emph{Coem. Callisti area subdialis et incerti loci}. Conservation selon l’EDB : \emph{Coem. Callisti area subdialis et incerti loci}.}
\nf{Support}{Plaque de marbre. Inscription gravée. Fiche EDB : Marida Pierno, première version 2015.}
\ntext{\lnum{1}Marcena [- - - d]ulcis qui tam [- - -]\\
\lnum{2}[r]edemptu[m - - -]co te aeterna tellu[re - - -]\\
\lnum{3}[- - -]PQ redemptus vulner[e - - -]\\
\lnum{4}[- - -]ritus ipse haec enim [- - -]\\
\lnum{5}[- - -]imus t[e pa]rentibus [- - -]\\
\lnum{6}[- - -]CO[- - - d]ulcis pace [- - -]\\
\lnum{7}[- - - de]p(ositus) d(ie) XVI Kal(endas) Se[pt(embres) - - -]\\
\lnum{8}[- - -] m(enses) II b(ene) m(erenti) paren[tes - - -]\\
\lnum{9}[- - - Gratiano Aug(usto) I]I et Probo [c]o[nss(ulibus)]}
\nf{Témoins}{EDB (EDB39561 ; divergence forte) : \emph{Marce na[--- d]ulcis qui tam [---] / [r]ẹḍemptu[--- e]go te aeterna tellu[re ---] / [---]+++[--- cor]p⌜o⌝re demptus vulner[e acerbo?] / [--- ir]ritus ipse haec enim / [---]imus t[e pa]rentibus / [---]CO[--- d]ulcis pace / --- dp(ositus) d(ie) XVI kal(endas) se[p(tembres)] / [---] m(ensibus) II b(ene)m(erenti) parent[es] / [Gratiano Augusto I]I et Probo [v(iro) c(larissimo) c]o[nsulibus]}. La dalle est perdue ; elle est connue par de Rossi (1864) et par le dessin d'ICVR I. L'EDB (EDB39561) lit \emph{Marce na[---]} (EDCS : \emph{Marcena}) et \emph{[cor]p⟨o⟩re demptus} (EDCS : \emph{PQ redemptus}). La lecture \emph{redemptus}, qui fondait une allusion à la rédemption, est donc contestée par la seconde base : l'occurrence est retirée du vocabulaire de la rédemption et signalée comme lecture incertaine.}
\nf{Traduction}{Marcena (?) […] doux, qui si […] ; racheté […] toi, en la terre éternelle […] ; […] racheté par la blessure […] ; […] lui-même, en effet, ces choses […] ; […] toi à tes parents […] ; […] doux, (en) paix […] ; déposé le 16e jour avant les calendes de septembre […] ; […] 2 mois ; bien méritant, ses parents […] ; sous le consulat de Gratien Auguste pour la 2e fois et de Probus.}
\nf{Langue}{9 abréviations développées par l’éditeur (dont dep(ositus), d(ie), Kal(endas), Sept(embres), m(enses), b(ene)).}
\nf{Chronologie}{Noms consulaires : Gratianus (restitué) et Probus (lu). Dates du jour : déposition : die XVI Kalendas Septembres = 17 août 371 (partiellement restitué).}
\nf{Commentaire}{Le catalogue antérieur relevait une allusion à la rédemption (\emph{redemptum}, \emph{redemptus}) ; la lecture \emph{redemptus} est contestée par l’EDB (voir les témoins).}
\nf{Bibliographie}{ICUR IV 11102 ; ILCV 1614 (em). Bibliographie complémentaire selon l’EDB : IC I, 239, n. 563 (right part) ; de Rossi 1864, 325 (left part) ; IC I Supplementum, 1614 (left part) ; ILCV, 1614.}
\nf{Statut}{Lecture secondaire (EDCS) ; divergence forte avec l’EDB : le texte ne doit pas être exploité sans contrôle de l’édition imprimée.}
\end{notice}

\begin{notice}{153}{371}{ICUR n.s. IV, 11103 · EDCS-37200385 · EDB39562}
\label{not:153}
\nf{Datation}{371 — directe : nom consulaire lu}
\nf{Genre}{Épitaphe ; ensembles A, B, C ; 44 lettres transmises ; 61 \% restituées.}
\nf{Références}{ICUR IV 11103.}
\nf{Provenance}{Rattachement selon l’EDB : \emph{Coem. Callisti area subdialis et incerti loci}. Conservation selon l’EDB : \emph{Coem. Callisti area subdialis et incerti loci}.}
\nf{Support}{Plaque de marbre. 10,5 × 21,5 × 2,5 cm (h. × l. × ép.). Inscription gravée. Fiche EDB : Marida Pierno, première version 2015.}
\ntext{\lnum{1}]\\
\lnum{2}IIDILIDCE[- - - Gratiano Aug(usto) II]\\
\lnum{3}e⸢t⸣ Probo c[onss(ulibus)}
\nf{Témoins}{Export EDCS de 2020 : \emph{] / IIDILIDCE[3 Gratiano Aug(usto) II] / e Probo c[onss(ulibus)}. EDB (EDB39562 ; divergence forte) : \emph{------ / +++++++CE+[---] / et Probo c[onsulibus]}.}
\nf{Apparat}{⸢t⸣ : I gravé}
\nf{Traduction}{[…] sous le consulat de Gratien Auguste pour la 2e fois et de Probus.}
\nf{Langue}{L’éditeur signale une graphie corrigée : I gravé pour t.}
\nf{Chronologie}{Noms consulaires : Gratianus (restitué) et Probus (lu). L'EDB (EDB39562) ne restitue pas le nom de Gratien : \emph{+++++++CE+[---] / et Probo c[onsulibus]}. L'année 371 repose sur le seul nom de Probus, lu.}
\nf{Bibliographie}{ICUR IV 11103.}
\nf{Statut}{Lecture secondaire (EDCS) ; divergence forte avec l’EDB : le texte ne doit pas être exploité sans contrôle de l’édition imprimée. Plus de la moitié des lettres sont restituées.}
\end{notice}

\begin{notice}{154}{371}{ICUR n.s. IV, 11765 · EDCS-35700944 · EDB2099}
\label{not:154}
\nf{Datation}{371 — directe : nom consulaire lu}
\nf{Genre}{Épitaphe ; ensembles A, B, C ; 31 lettres transmises.}
\nf{Références}{ICUR IV 11765.}
\nf{Provenance}{Rattachement selon l’EDB : \emph{Coem. anonymum ad viam Appiam}. Conservation selon l’EDB : \emph{Coem. anonymum ad viam Appiam}.}
\nf{Support}{Enduit de mortier de mortier. Inscription tracée dans le mortier frais. Fiche EDB : Maria Cristina Grisanzio, première version 2003.}
\ntext{\lnum{1}D(epositio) Petri IIII Kal(endas) Oct(obres) Gratiano et Probo}
\nf{Témoins}{EDB (EDB2099) : texte concordant aux conventions de notation près (indice de similarité 1).}
\nf{Traduction}{Déposition de Petrus le 4e jour avant les calendes d’octobre, sous (le consulat de) Gratien et Probus.}
\nf{Langue}{3 abréviations développées par l’éditeur (dont D(epositio), Kal(endas), Oct(obres)).}
\nf{Chronologie}{Noms consulaires : Gratianus (lu) et Probus (lu). Dates du jour : déposition : IIII Kalendas Octobres = 28 septembre 371 (lu).}
\nf{Bibliographie}{ICUR IV 11765.}
\nf{Statut}{Lecture secondaire (EDCS) concordante avec l’EDB : lecture stabilisée entre deux bases, non collationnée sur l’ICVR ni sur le monument.}
\end{notice}

\annee{372}

\begin{notice}{155}{372}{ICUR n.s. I, 83 · EDCS-38600099 · EDB27456}
\label{not:155}
\nf{Datation}{372 — directe : nom consulaire lu}
\nf{Genre}{Épitaphe ; ensembles A, B, C ; 118 lettres transmises.}
\nf{Références}{ICUR I 83 = ILCV 4670.}
\nf{Provenance}{Rattachement selon l’EDB : \emph{Incertae originis} (le volume I des ICVR réunit les inscriptions dont la provenance est perdue). Conservation selon l’EDB : \emph{s. Cecilia in Trastevere}.}
\nf{Support}{Plaque de marbre. 85 × 197 × 3 cm (h. × l. × ép.). Inscription gravée. Fiche EDB : Filippo Piazzolla, première version 2013.}
\ntext{\lnum{1}Praestantissimo adque amatori totius domus\\
\lnum{2}Romuliano bene merenti in pace\\
\lnum{3}depositus IIII Nonas Decembres\\
\lnum{4}Modesto et Arantheo consulibus}
\nf{Témoins}{EDB (EDB27456) : texte concordant aux conventions de notation près (indice de similarité 0,964).}
\nf{Traduction}{Au très excellent et ami de toute la maison, Romulianus, bien méritant, en paix ; déposé le 4e jour avant les nones de décembre, sous le consulat de Modestus et d’Arintheus.}
\nf{Langue}{Aucune particularité graphique n’est signalée par l’éditeur dans l’état 2022 de la base.}
\nf{Chronologie}{Noms consulaires : Modestus (lu). Ne sont pas des consuls : Romuliano bene merenti (défunt (le second consul lu est Arantheo = Arintheus)). Dates du jour : déposition : IIII Nonas Decembres = 2 décembre 372 (lu).}
\nf{Commentaire}{La formule in pace est conservée.}
\nf{Bibliographie}{ICUR I 83 ; ILCV 4670. Bibliographie complémentaire selon l’EDB : ILCV, 4670.}
\nf{Statut}{Lecture secondaire (EDCS) concordante avec l’EDB : lecture stabilisée entre deux bases, non collationnée sur l’ICVR ni sur le monument.}
\end{notice}

\begin{notice}{156}{372}{ICUR n.s. I, 1433 · EDCS-39000416 · EDB31197}
\label{not:156}
\nf{Datation}{372 — directe : nom consulaire lu}
\nf{Genre}{Épitaphe ; ensembles A, B, C ; 136 lettres transmises.}
\nf{Références}{ICUR I 1433.}
\nf{Provenance}{Rattachement selon l’EDB : \emph{Incertae originis} (le volume I des ICVR réunit les inscriptions dont la provenance est perdue). Conservation selon l’EDB : \emph{Musei Vaticani: Lapidario Cristiano}.}
\nf{Support}{Plaque de marbre. 17 × 27 cm (h. × l.). Inscription gravée. Reproductions signalées par l’EDB : reproduction, Perret 1851 (droits réservés ; non reproduites). Fiche EDB : Filippo Piazzolla, première version 2015.}
\ntext{\lnum{1}Depositus est puer\\
\lnum{2}nomine Exsuperanti\\
\lnum{3}us qui vixit anno⸢s⸣ pl\{l\}(us)\\
\lnum{4}minus IIII et menses VII dies V\\
\lnum{5}depositus VII Idus Maias sub\\
\lnum{6}consulatu Modest⸢i⸣ et Fl(avi) Arin\\
\lnum{7}the(i) cc(on)ss(ulum) in pace}
\nf{Témoins}{Export EDCS de 2020 : \emph{Depositus est puer / nomine Exsuperanti/us qui vixit anno pl\{l\}(us) / minus IIII et menses VII dies V / depositus VII Idus Maias sub / consulatu Modest et Fl(avi) Arin/the(i) cc(on)ss(ulum) in pace}. EDB (EDB31197) : texte concordant aux conventions de notation près (indice de similarité 0,959).}
\nf{Apparat}{⸢s⸣ : RO gravé ; ⸢i⸣ : O gravé ; \{l\} : lettres superflues}
\nf{Traduction}{A été déposé le garçon nommé Exsuperantius, qui vécut environ 4 ans et 7 mois 5 jours ; déposé le 7e jour avant les ides de mai, sous le consulat de Modestus et de Flavius Arintheus, en paix.}
\nf{Langue}{L’éditeur signale 2 graphies corrigées : RO gravé pour s ; O gravé pour i. Des lettres gravées sont jugées superflues par l’éditeur (entre accolades). 5 abréviations développées par l’éditeur (dont (us), Fl(avi), the(i), cc(on)ss, (ulum)).}
\nf{Chronologie}{Noms consulaires : Modestus (lu) et Arintheus (lu). Dates du jour : déposition : VII Idus Maias = 9 mai 372 (lu). Âge déclaré : env. 4 a. 7 m. 5 j. (lu).}
\nf{Commentaire}{La formule in pace est conservée.}
\nf{Bibliographie}{ICUR I 1433. Bibliographie complémentaire selon l’EDB : Lega 2000, 90 ; ILCV, 2976B adn.}
\nf{Statut}{Lecture secondaire (EDCS) concordante avec l’EDB : lecture stabilisée entre deux bases, non collationnée sur l’ICVR ni sur le monument. L’EDCS signale une photographie (non consultée).}
\end{notice}

\begin{notice}{157}{372}{ICUR n.s. I, 1434 · EDCS-39000417 · EDB31198}
\label{not:157}
\nf{Datation}{372 — directe : nom consulaire lu}
\nf{Genre}{Épitaphe ; ensembles A, B, C ; 91 lettres transmises.}
\nf{Références}{ICUR I 1434 = ILCV 2690 (em).}
\nf{Provenance}{Rattachement selon l’EDB : \emph{Incertae originis} (le volume I des ICVR réunit les inscriptions dont la provenance est perdue). Conservation selon l’EDB : \emph{Musei Vaticani: Lapidario Cristiano}.}
\nf{Support}{Plaque de marbre. 28 × 54 cm (h. × l.). Inscription gravée. Fiche EDB : Filippo Piazzolla, première version 2015.}
\ntext{\lnum{1}Fl(avio) Domitio et Fl(avio) Ar⸢i⸣ntheo co(n)ss(ulibus)\\
\lnum{2}XV Kal(endas) Sep(tem)b(res) deposita Floren\\
\lnum{3}tia qu(a)e vixit annos VIII\\
\lnum{4}mens\{s\}⸢e⸣s IIII dies XII\\
\lnum{5}in pace}
\nf{Témoins}{Export EDCS de 2020 : \emph{Fl(avio) Domitio et Fl(avio) Arntheo co(n)ss(ulibus) / XV Kal(endas) Sep(tem)b(res) deposita Floren/tia qu(a)e vixit annos VIII / mens\{s\}s IIII dies XII / in pace}. EDB (EDB31198) : texte concordant aux conventions de notation près (indice de similarité 0,922).}
\nf{Apparat}{⸢i⸣ : V gravé ; ⸢e⸣ : I gravé ; \{s\} : lettres superflues}
\nf{Traduction}{Sous le consulat de Flavius Domitius (Modestus) et de Flavius Arintheus, le 15e jour avant les calendes de septembre, déposée Florentia, qui vécut 8 ans 4 mois 12 jours, en paix.}
\nf{Langue}{L’éditeur signale 2 graphies corrigées : V gravé pour i ; I gravé pour e (confusion e/i). Des lettres gravées sont jugées superflues par l’éditeur (entre accolades). 8 abréviations développées par l’éditeur (dont Fl(avio), co(n)ss, (ulibus), Kal(endas), Sep(tem)b).}
\nf{Chronologie}{Noms consulaires : Arintheus (lu). Dates du jour : événement indéterminé : XV Kalendas Septembres = 18 août 372 (lu). Âge déclaré : 8 a. 4 m. 12 j. (lu).}
\nf{Commentaire}{La formule in pace est conservée.}
\nf{Bibliographie}{ICUR I 1434 ; ILCV 2690 (em). Bibliographie complémentaire selon l’EDB : ILCV, 2690 ; Lega 2000, 135.}
\nf{Statut}{Lecture secondaire (EDCS) concordante avec l’EDB : lecture stabilisée entre deux bases, non collationnée sur l’ICVR ni sur le monument.}
\end{notice}

\begin{notice}{158}{372}{ICUR n.s. I, 2090 · EDCS-39300313 · EDB32056}
\label{not:158}
\nf{Datation}{372 — directe : lecture partielle d’un nom consulaire}
\nf{Genre}{Épitaphe ; ensembles A, B, C ; 49 lettres transmises ; 61 \% restituées.}
\nf{Références}{ICUR I 2090.}
\nf{Provenance}{Rattachement selon l’EDB : \emph{Incertae originis} (le volume I des ICVR réunit les inscriptions dont la provenance est perdue). Conservation selon l’EDB : \emph{Musei Vaticani: Galleria Lapidaria}.}
\nf{Support}{Plaque de marbre. 16 × 18 cm (h. × l.). Inscription gravée. Fiche EDB : Filippo Piazzolla, première version 2015.}
\ntext{\lnum{1}]s in pa[ce - - -]\\
\lnum{2}[vixit annos] XLV d(e)p(ositus) V[- - -]\\
\lnum{3}[consulatu M]odesti et [Arinthei}
\nf{Témoins}{EDB (EDB32056) : texte concordant aux conventions de notation près (indice de similarité 0,968).}
\nf{Traduction}{[…] en paix […] ; vécut 45 ans ; déposé le 5e […] ; sous le consulat de Modestus et d’Arintheus.}
\nf{Langue}{2 abréviations développées par l’éditeur (dont d(e)p, (ositus)).}
\nf{Chronologie}{Noms consulaires : Modestus (lu en partie) et Arintheus (restitué). Âge déclaré : 45 a. (partiellement restitué).}
\nf{Commentaire}{La formule in pace est attestée seulement par restitution.}
\nf{Bibliographie}{ICUR I 2090. Bibliographie complémentaire selon l’EDB : Lega 2000.}
\nf{Statut}{Lecture secondaire (EDCS) concordante avec l’EDB : lecture stabilisée entre deux bases, non collationnée sur l’ICVR ni sur le monument. Plus de la moitié des lettres sont restituées.}
\end{notice}

\begin{notice}{159}{372}{ICUR n.s. I, 3183 · EDCS-40300329 · EDB33135}
\label{not:159}
\nf{Datation}{372 — directe : lecture partielle d’un nom consulaire}
\nf{Genre}{Épitaphe ; ensembles A, B, C ; 37 lettres transmises ; 54 \% restituées.}
\nf{Références}{ICUR I 3183.}
\nf{Provenance}{Rattachement selon l’EDB : \emph{Incertae originis} (le volume I des ICVR réunit les inscriptions dont la provenance est perdue). Conservation selon l’EDB : \emph{Antiquarium del Celio}.}
\nf{Support}{Plaque de marbre. 16 × 21,5 × 5 cm (h. × l. × ép.). Inscription gravée. Fiche EDB : Filippo Piazzolla, première version 2015.}
\ntext{\lnum{1}Her]culiu[s - - -]\\
\lnum{2}[- - -]EN [- - -]\\
\lnum{3}[- - - Modesto et Ari]nteo con[s(ulibus)}
\nf{Témoins}{EDB (EDB33135 ; écarts à examiner) : \emph{------ / [--- Her]culiu[s ---] / [---]en [---] / [--- Modesto et Ari]ṇteo con[s(ulibu)s ---] / ------ [texte tronqué dans la liste de résultats]}.}
\nf{Traduction}{Herculius […] ; sous le consulat de Modestus et d’Arintheus.}
\nf{Langue}{1 abréviation développée par l’éditeur (dont cons(ulibus)).}
\nf{Chronologie}{Noms consulaires : Modestus (restitué) et Arintheus (lu en partie).}
\nf{Bibliographie}{ICUR I 3183. Bibliographie complémentaire selon l’EDB : Ferrua 1979b, 91.}
\nf{Statut}{Lecture secondaire (EDCS) ; écarts avec l’EDB à arbitrer sur l’ICVR. Plus de la moitié des lettres sont restituées.}
\end{notice}

\begin{notice}{160}{372}{ICUR n.s. II, 4748 · EDCS-34201004 · EDB13470}
\label{not:160}
\nf{Datation}{372 — directe : nom consulaire lu}
\nf{Genre}{Épitaphe ; ensembles A, B, C ; 51 lettres transmises ; 8 \% restituées.}
\nf{Références}{ICUR II 4748.}
\nf{Provenance}{Rattachement selon l’EDB : \emph{Coem. Generosae}.}
\nf{Support}{Enduit de mortier de mortier. Inscription tracée dans le mortier frais. Signes décrits par l’EDB : chrisme avec alpha et oméga. Fiche EDB : Antonio Enrico Felle, première version 2003.}
\ntext{\lnum{1}[P]aula defuncta est die\\
\lnum{2}[- - - Sep]tembrium cons(ulibu)s Modesto et Arenteo}
\nf{Témoins}{EDB (EDB13470) : texte concordant aux conventions de notation près (indice de similarité 0,984).}
\nf{Traduction}{Paula est morte le […] (de) septembre, sous le consulat de Modestus et d’Arintheus.}
\nf{Langue}{1 abréviation développée par l’éditeur (dont cons(ulibu)s).}
\nf{Chronologie}{Noms consulaires : Modestus (lu).}
\nf{Bibliographie}{ICUR II 4748.}
\nf{Statut}{Lecture secondaire (EDCS) concordante avec l’EDB : lecture stabilisée entre deux bases, non collationnée sur l’ICVR ni sur le monument.}
\end{notice}

\begin{notice}{161}{372}{ICUR n.s. II, 4808 · EDCS-36101289}
\label{not:161}
\nf{Datation}{372 — directe : consulat restitué}
\nf{Genre}{Épitaphe ; ensembles A, B, C ; 57 lettres transmises ; 58 \% restituées.}
\nf{Références}{ICUR II 4808.}
\nf{Provenance}{Rome (indication de l’EDCS) ; aucun rattachement cimétérial transmis.}
\nf{Support}{Non documenté dans les sources consultées.}
\ntext{\lnum{1}]bendi\\
\lnum{2}[- - - vixit a]nn⸢o⸣s\\
\lnum{3}[- - - K]al(endas) Septe[mb(res) Arinth]io et Mo\\
\lnum{4}[desto cons(ulibus) - - -]us}
\nf{Témoins}{Export EDCS de 2020 : \emph{]bendi / [3 vixit a]nns / [3 K]al(endas) Septe[mb(res) Arinth]io et Mo/[desto cons(ulibus) 3]us}. Aucune fiche EDB appariée (numéro ICUR absent de la liste publique de l’EDB).}
\nf{Apparat}{⸢o⸣ : V gravé}
\nf{Traduction}{[…] vécut […] ans ; […] calendes de septembre, sous le consulat d’Arintheus et de Modestus […].}
\nf{Langue}{L’éditeur signale une graphie corrigée : V gravé pour o (confusion o/u).}
\nf{Chronologie}{Noms consulaires : Modestus (restitué) et Arintheus (restitué). Dates du jour : événement indéterminé : … Kalendas Septembres (partiellement restitué).}
\nf{Bibliographie}{ICUR II 4808.}
\nf{Statut}{Lecture secondaire (EDCS) sans second témoin numérique. Plus de la moitié des lettres sont restituées. L’EDCS signale une photographie (non consultée).}
\end{notice}

\begin{notice}{162}{372}{ICUR n.s. III, 6503 · EDCS-34801303 · EDB22061}
\label{not:162}
\nf{Datation}{372 — directe : lecture partielle d’un nom consulaire}
\nf{Genre}{Épitaphe ; ensembles A, B, C ; 82 lettres transmises ; 51 \% restituées.}
\nf{Références}{ICUR III 6503.}
\nf{Provenance}{Rattachement selon l’EDB : \emph{Coem. Domitillae pars inferior}. Conservation selon l’EDB : \emph{Coem. Domitillae pars inferior}.}
\nf{Support}{Plaque de marbre. 52 × 72 × 4 cm (h. × l. × ép.). Inscription gravée. Reproductions signalées par l’EDB : photographie PCAS (droits réservés ; non reproduites). Fiche EDB : Carlo Carletti / Antonio Enrico Felle, première version 2003.}
\ntext{\lnum{1}Lu[- - - be]ne merenti qui [vixit an]n(os) LIII[- - -]\\
\lnum{2}do[- - - Mod]esto et Fl(avio) Aronteo [vv(iris) cc(larissimis) cons(ulibu)s i]n pac[e]}
\nf{Témoins}{EDB (EDB22061 ; écarts à examiner) : \emph{Lu[--- be]nemerenti qui [vixit an]n(os) LIII[---] / do[--- Mod]esto et Fl(avio) Aronteo [cons(ulibu)s i]n pac[e] [texte tronqué dans la liste de résultats]}.}
\nf{Traduction}{[…] bien méritant, qui vécut 53 ans […] ; […] sous le consulat de Modestus et de Flavius Arintheus, clarissimes ; en paix.}
\nf{Langue}{5 abréviations développées par l’éditeur (dont ann(os), Fl(avio), vv(iris), cc(larissimis), cons(ulibu)s).}
\nf{Chronologie}{Noms consulaires : Modestus (lu en partie).}
\nf{Commentaire}{La formule in pace est attestée seulement par restitution.}
\nf{Bibliographie}{ICUR III 6503.}
\nf{Statut}{Lecture secondaire (EDCS) ; écarts avec l’EDB à arbitrer sur l’ICVR. Plus de la moitié des lettres sont restituées. L’EDCS signale une photographie (non consultée).}
\end{notice}

\begin{notice}{163}{372}{ICUR n.s. IV, 9565 · EDCS-36101094 · EDB18369}
\label{not:163}
\nf{Datation}{372 — directe : lecture partielle d’un nom consulaire}
\nf{Genre}{Épitaphe ; ensembles A, B, C ; 36 lettres transmises ; 39 \% restituées.}
\nf{Références}{ICUR IV 9565.}
\nf{Provenance}{Rattachement selon l’EDB : \emph{Coem. Callisti pars inferior}.}
\nf{Support}{Plaque de marbre. Inscription gravée. Fiche EDB : Antonio Enrico Felle, première version 2003.}
\ntext{\lnum{1}] qui vixit\\
\lnum{2}[- - - K]al(endas)\\
\lnum{3}[- - - Modesto et Ar]onteo\\
\lnum{4}[- - -]um\\
\lnum{5}[- - - he]redum}
\nf{Témoins}{EDB (EDB18369) : texte concordant aux conventions de notation près (indice de similarité 0,939).}
\nf{Traduction}{[…] qui vécut […] ; […] calendes […] ; (sous le consulat de) Modestus et Arintheus ; […] des héritiers.}
\nf{Langue}{1 abréviation développée par l’éditeur (dont Kal(endas)).}
\nf{Chronologie}{Noms consulaires : Modestus (restitué).}
\nf{Bibliographie}{ICUR IV 9565.}
\nf{Statut}{Lecture secondaire (EDCS) concordante avec l’EDB : lecture stabilisée entre deux bases, non collationnée sur l’ICVR ni sur le monument.}
\end{notice}

\begin{notice}{164}{372}{ICUR n.s. IV, 11104 · EDCS-37200386 · EDB39563}
\label{not:164}
\nf{Datation}{372 — directe : lecture partielle d’un nom consulaire}
\nf{Genre}{Épitaphe ; ensembles A, B, C ; 48 lettres transmises ; 44 \% restituées.}
\nf{Références}{ICUR IV 11104 = ILCV 4461a.}
\nf{Provenance}{Rattachement selon l’EDB : \emph{Coem. Callisti area subdialis et incerti loci}. Conservation selon l’EDB : \emph{Coem. Callisti area subdialis et incerti loci}.}
\nf{Support}{Plaque de marbre. 84 × 60 cm (h. × l.). Inscription gravée. Reproductions signalées par l’EDB : photographie PCAS (droits réservés ; non reproduites). Fiche EDB : Marida Pierno, première version 2015.}
\ntext{\lnum{1}civ]is Alicensi\\
\lnum{2}[- - -]III d(ies) XVI\\
\lnum{3}[- - - Mo]desto e⸢t⸣\\
\lnum{4}[Arintheo conss(ulibus)}
\nf{Témoins}{Export EDCS de 2020 : \emph{civ]is Alicensi / [3]III d(ies) XVI / [3 Mo]desto e / [Arintheo conss(ulibus)}. EDB (EDB39563 ; écarts à examiner) : \emph{[--- civ]i salicensi (?) / [----]III d(iebus) XVI / [Mo]desto e⌜t⌝ / [Arintheo] ((a̲v̲i̲s̲))}.}
\nf{Apparat}{⸢t⸣ : I gravé}
\nf{Traduction}{[…] citoyen d’Alicensis (?) […] ; […] 3 (?) […] 16 jours ; sous le consulat de Modestus et d’Arintheus.}
\nf{Langue}{L’éditeur signale une graphie corrigée : I gravé pour t.}
\nf{Chronologie}{Noms consulaires : Modestus (lu en partie) et Arintheus (restitué).}
\nf{Bibliographie}{ICUR IV 11104 ; ILCV 4461a. Bibliographie complémentaire selon l’EDB : ILCV, 4461A.}
\nf{Statut}{Lecture secondaire (EDCS) ; écarts avec l’EDB à arbitrer sur l’ICVR. L’EDCS signale une photographie (non consultée).}
\end{notice}

\begin{notice}{165}{372}{ICUR n.s. IV, 11613t · EDCS-37300523 · EDB40555}
\label{not:165}
\nf{Datation}{372 — directe : lecture partielle d’un nom consulaire ; datation contestée}
\nf{Genre}{Épitaphe ; ensembles A, B, C ; 34 lettres transmises ; 68 \% restituées.}
\nf{Références}{ICUR IV 11613t.}
\nf{Provenance}{Rattachement selon l’EDB : \emph{Coem. Callisti area subdialis et incerti loci}. Conservation selon l’EDB : \emph{Coem. Callisti area subdialis et incerti loci}.}
\nf{Support}{Plaque de marbre. 16 × 26 cm (h. × l.). Inscription gravée. Fiche EDB : Marida Pierno, première version 2015.}
\ntext{\lnum{1}p]ace VII [- - -]\\
\lnum{2}[- - - Mo]dest[o et Arintheo cons(ulibus ?)}
\nf{Témoins}{EDB (EDB40555 ; divergence forte) : \emph{[--- in p]ace VII / [---]dest[---] / ------}.}
\nf{Traduction}{[…] paix, 7 […] ; sous le consulat de Modestus et d’Arintheus.}
\nf{Langue}{1 abréviation développée par l’éditeur (dont cons(ulibus ?)).}
\nf{Chronologie}{Noms consulaires : Modestus (lu en partie) et Arintheus (restitué). Seule la séquence \emph{dest} est lue. L'année 372 suppose la restitution \emph{[Mo]dest[o et Arintheo]} proposée par Ferrua ; l'EDB (EDB40555) date le texte « 300–399 ». Datation conjecturale.}
\nf{Bibliographie}{ICUR IV 11613t.}
\nf{Statut}{Lecture secondaire (EDCS) ; divergence forte avec l’EDB : le texte ne doit pas être exploité sans contrôle de l’édition imprimée. Plus de la moitié des lettres sont restituées.}
\end{notice}

\annee{373}

\begin{notice}{166}{373}{ICUR n.s. I, 1254 · EDCS-39101964 · EDB30476}
\label{not:166}
\nf{Datation}{373 — directe : nom consulaire lu}
\nf{Genre}{Épitaphe ; ensembles A, B, C ; 81 lettres transmises.}
\nf{Références}{ICUR I 1254 = ILCV 2604.}
\nf{Provenance}{Rattachement selon l’EDB : \emph{Incertae originis} (le volume I des ICVR réunit les inscriptions dont la provenance est perdue). Conservation selon l’EDB : \emph{Musei Capitolini}.}
\nf{Support}{Plaque de marbre. 31 × 42 cm (h. × l.). Inscription gravée. Reproductions signalées par l’EDB : reproduction, Luni et alii 1986 (droits réservés ; non reproduites). Fiche EDB : Filippo Piazzolla, première version 2015.}
\ntext{\lnum{1}Amantio bene mere\\
\lnum{2}nti in pace q(u)i vix\\
\lnum{3}\{x\}it an(nos) XVIII m(enses) II d(ies) VII\\
\lnum{4}II d(epositus) X Kal(endas) Sept(embres) Valentinia\\
\lnum{5}no et Valenti\\
\lnum{6}IIII}
\nf{Témoins}{EDB (EDB30476) : texte concordant aux conventions de notation près (indice de similarité 0,936).}
\nf{Apparat}{\{x\} : lettres superflues}
\nf{Traduction}{À Amantius, bien méritant, en paix, qui vécut 18 ans 2 mois 9 jours ; déposé le 10e jour avant les calendes de septembre, sous (le consulat de) Valentinien et Valens pour la 4e fois.}
\nf{Langue}{Des lettres gravées sont jugées superflues par l’éditeur (entre accolades). 7 abréviations développées par l’éditeur (dont q(u)i, an(nos), m(enses), d(ies), d(epositus), Kal(endas)).}
\nf{Chronologie}{Noms consulaires : Valentinianus (lu) et Valens (lu). Ne sont pas des consuls : Amantio bene mere/nti (défunt). Dates du jour : déposition : X Kalendas Septembres = 23 août 373 (lu). Âge déclaré : 18 a. 2 m. 9 j. (lu).}
\nf{Commentaire}{La formule in pace est conservée.}
\nf{Bibliographie}{ICUR I 1254 ; ILCV 2604. Bibliographie complémentaire selon l’EDB : ILCV, 2604 ; Molisani 1973, 110 ; Luni et alii 1986, 178-179 tab. XX-1.}
\nf{Statut}{Lecture secondaire (EDCS) concordante avec l’EDB : lecture stabilisée entre deux bases, non collationnée sur l’ICVR ni sur le monument. L’EDCS signale une photographie (non consultée).}
\end{notice}

\begin{notice}{167}{373}{ICUR n.s. I, 2091 · EDCS-39300314 · EDB32057}
\label{not:167}
\nf{Datation}{373 — directe : lecture partielle d’un nom consulaire}
\nf{Genre}{Épitaphe ; ensembles A, B, C ; 80 lettres transmises ; 41 \% restituées.}
\nf{Références}{ICUR I 2091,1.}
\nf{Provenance}{Rattachement selon l’EDB : \emph{Incertae originis} (le volume I des ICVR réunit les inscriptions dont la provenance est perdue). Conservation selon l’EDB : \emph{Musei Vaticani: Galleria Lapidaria}.}
\nf{Support}{Plaque de marbre. 34 × 40 cm (h. × l.). Inscription gravée. Fiche EDB : Filippo Piazzolla, première version 2015.}
\ntext{\lnum{1}Niceforidi uxor[i dulci quae vix(it)]\\
\lnum{2}ann(os) p(lus) m(inus) XXV dep(osita) [- - -]\\
\lnum{3}Septemb(res) Valen[tiniano et Valente]\\
\lnum{4}[Au]g(usto) IIII conss(ulibus) [}
\nf{Témoins}{EDB (EDB32057) : texte concordant aux conventions de notation près (indice de similarité 0,96).}
\nf{Traduction}{À Nicephoris, épouse douce, qui vécut environ 25 ans ; déposée […] septembre, sous le consulat de Valentinien et de Valens Auguste(s) pour la 4e fois.}
\nf{Langue}{8 abréviations développées par l’éditeur (dont vix(it), ann(os), p(lus), m(inus), dep(osita), Septemb(res)).}
\nf{Chronologie}{Noms consulaires : Valentinianus (lu en partie) et Valens (lu en partie). Âge déclaré : env. 25 a. (partiellement restitué).}
\nf{Bibliographie}{ICUR I 2091,1. Bibliographie complémentaire selon l’EDB : Lega 2000.}
\nf{Statut}{Lecture secondaire (EDCS) concordante avec l’EDB : lecture stabilisée entre deux bases, non collationnée sur l’ICVR ni sur le monument.}
\end{notice}

\begin{notice}{168}{373}{ICUR n.s. I, 2092 · EDCS-34700234 · EDB11297}
\label{not:168}
\nf{Datation}{373 — directe : nom consulaire lu}
\nf{Genre}{Épitaphe ; ensembles A, B, C ; 183 lettres transmises ; 37 \% restituées.}
\nf{Références}{ICUR I 2092 = ILCV 4458.}
\nf{Provenance}{Rattachement selon l’EDB : \emph{Incertae originis} (le volume I des ICVR réunit les inscriptions dont la provenance est perdue). Conservation selon l’EDB : \emph{Musei Vaticani: Depositi}.}
\nf{Support}{Plaque de marbre. 53 × 100 cm (h. × l.). Inscription gravée. Signes décrits par l’EDB : croix monogrammatique. Fiche EDB : Antonio Enrico Felle, première version 2003.}
\ntext{\lnum{1}[Locus - - - u]bi posita es[t - - -]\\
\lnum{2}[Valente Aug(usto)] IIII et Valentinia[no Aug(usto) IIII conss(ulibus)]\\
\lnum{3}[- - -] in pace qu(a)e vixit annis XXV[- - -]\\
\lnum{4}[τόπος - - -]νου ζῶντος ἐν ᾧ κατάκιτε\\
\lnum{5}πεδ[ὶν Λέων(?) κ]αὶ εὐγενεστάτη Κ[ω]νσταντινου\\
\lnum{6}πολιτί[σσα Βι]ταλία ἐτῶ[ν εἴκο]σι ἑπτά}
\nf{Témoins}{Export EDCS de 2020 : \emph{[Locus 3 u]bi posita es[t 3] / [Valente Aug(usto)] IIII et Valentinia[no Aug(usto) IIII conss(ulibus)] / [3] in pace qu(a)e vixit annis XXV[3] //}. EDB (EDB11297) : texte concordant aux conventions de notation près (indice de similarité 0,982).}
\nf{Traduction}{Emplacement […] où a été placée […] ; sous le consulat de Valens Auguste pour la 4e fois et de Valentinien Auguste pour la 4e fois ; […] en paix, qui vécut 25 ans […] // (grec) Emplacement de […]nus, acquis de son vivant, dans lequel reposent l’enfant Léon et la très noble Constantinopolitaine Vitalia, âgée de vingt-sept ans.}
\nf{Langue}{Le texte est rédigé en grec, en tout ou en partie ; l’export de 2020 ne le transmettait pas. 4 abréviations développées par l’éditeur (dont Aug(usto), conss(ulibus), qu(a)e).}
\nf{Chronologie}{Noms consulaires : Valentinianus (lu) et Valens (lu).}
\nf{Commentaire}{La formule in pace est conservée.}
\nf{Bibliographie}{ICUR I 2092 ; ILCV 4458.}
\nf{Statut}{Lecture secondaire (EDCS) concordante avec l’EDB : lecture stabilisée entre deux bases, non collationnée sur l’ICVR ni sur le monument.}
\end{notice}

\begin{notice}{169}{373}{ICUR n.s. III, 8724 · EDCS-35900968 · EDB20992}
\label{not:169}
\nf{Datation}{373 — directe : nom consulaire lu}
\nf{Genre}{Épitaphe ; ensembles A, B, C ; 70 lettres transmises.}
\nf{Références}{ICUR III 8724 = ILCV 4392.}
\nf{Provenance}{Rattachement selon l’EDB : \emph{Tituli qui in coem. Callisti reperti traduntur}.}
\nf{Support}{Plaque de marbre. Inscription gravée. Fiche EDB : Carlo Carletti, première version 2003.}
\ntext{\lnum{1}Depossio Ursi qui vixit an(nos) LV\\
\lnum{2}d(ies) XXVII dep(ositus) XV Kal(endas) April(es) d(ie) Lun(ae)\\
\lnum{3}cons(ulatu) Valentin(iani)\\
\lnum{4}et Vale(nti) XII DI()}
\nf{Témoins}{EDB (EDB20992) : texte concordant aux conventions de notation près (indice de similarité 1).}
\nf{Traduction}{Déposition d’Ursus, qui vécut 55 ans 27 jours ; déposé le 15e jour avant les calendes d’avril, un lundi, sous le consulat de Valentinien et de Valens ; 12 [lettres finales obscures : DI].}
\nf{Langue}{11 abréviations développées par l’éditeur (dont an(nos), d(ies), dep(ositus), Kal(endas), April(es), d(ie)).}
\nf{Chronologie}{Noms consulaires : Valentinianus (lu) et Valens (lu). Dates du jour : déposition : XV Kalendas Apriles = 18 mars 373 (lu). Âge déclaré : 55 a. 27 j. (lu).}
\nf{Bibliographie}{ICUR III 8724 ; ILCV 4392. Bibliographie complémentaire selon l’EDB : ILCV, 4392.}
\nf{Statut}{Lecture secondaire (EDCS) concordante avec l’EDB : lecture stabilisée entre deux bases, non collationnée sur l’ICVR ni sur le monument.}
\end{notice}

\begin{notice}{170}{373}{ICUR n.s. III, 9344c · EDCS-35900810 · EDB24922}
\label{not:170}
\nf{Datation}{373 — directe : nom consulaire lu}
\nf{Genre}{Épitaphe ; ensembles A, B, C ; 64 lettres transmises ; 45 \% restituées.}
\nf{Références}{ICUR III 9344c.}
\nf{Provenance}{Rattachement selon l’EDB : \emph{Basilica ss. Nerei et Achillei in coem. Domitillae}.}
\nf{Support}{Plaque de marbre. Inscription gravée. Fiche EDB : Carlo Carletti / Antonio Enrico Felle, première version 2003.}
\ntext{\lnum{1}] co[- - -]\\
\lnum{2}[- - -] vixit ann[os - - -]\\
\lnum{3}[- - -] Idus Iul(ias) dd(ominis) n[n(ostris) Valen]\\
\lnum{4}[tiniano e]t Valente IIII con[s(ulibus) - - -]}
\nf{Témoins}{EDB (EDB24922 ; écarts à examiner) : \emph{[---] con[iugi ---] / [---] vixit ann[is ---] / [---] idus iul(ias) d(ominis) n(ostris) / [Valentiniano e]t Valente IIII con[sulibus]}.}
\nf{Traduction}{[…] vécut […] ans ; […] ides de juillet, sous le consulat de nos seigneurs Valentinien et Valens pour la 4e fois […].}
\nf{Langue}{4 abréviations développées par l’éditeur (dont Iul(ias), dd(ominis), nn(ostris), cons(ulibus)).}
\nf{Chronologie}{Noms consulaires : Valentinianus (restitué) et Valens (lu). Dates du jour : événement indéterminé : … Idus Iulias (partiellement restitué).}
\nf{Bibliographie}{ICUR III 9344c.}
\nf{Statut}{Lecture secondaire (EDCS) ; écarts avec l’EDB à arbitrer sur l’ICVR.}
\end{notice}

\begin{notice}{171}{373}{ICUR n.s. IV, 11766 · EDCS-35700945 · EDB2100}
\label{not:171}
\nf{Datation}{373 — directe : nom consulaire lu}
\nf{Genre}{Épitaphe ; ensembles A, B, C ; 66 lettres transmises ; 12 \% restituées.}
\nf{Références}{ICUR IV 11766.}
\nf{Provenance}{Rattachement selon l’EDB : \emph{Coem. anonymum ad viam Appiam}. Conservation selon l’EDB : \emph{Coem. anonymum ad viam Appiam}.}
\nf{Support}{Enduit de mortier de mortier. Inscription tracée dans le mortier frais. Fiche EDB : Maria Cristina Grisanzio, première version 2003.}
\ntext{\lnum{1}Depositus Ex[tric]atu[s - - - die] IIII Kal(endas) Decembre(s) Valentiniano et Valente IIII co(n)ss(ulibus)}
\nf{Témoins}{EDB (EDB2100) : texte concordant aux conventions de notation près (indice de similarité 0,928).}
\nf{Traduction}{Déposé, Extricatus […], le 4e jour avant les calendes de décembre, sous le consulat de Valentinien et de Valens pour la 4e fois.}
\nf{Langue}{4 abréviations développées par l’éditeur (dont Kal(endas), Decembre(s), co(n)ss, (ulibus)).}
\nf{Chronologie}{Noms consulaires : Valentinianus (lu) et Valens (lu). Dates du jour : déposition : … die IIII Kalendas Decembres = 28 novembre 373 (partiellement restitué).}
\nf{Bibliographie}{ICUR IV 11766.}
\nf{Statut}{Lecture secondaire (EDCS) concordante avec l’EDB : lecture stabilisée entre deux bases, non collationnée sur l’ICVR ni sur le monument.}
\end{notice}

\annee{374}

\begin{notice}{172}{374}{ICUR n.s. I, 1937 · EDCS-39200381 · EDB31855}
\label{not:172}
\nf{Datation}{374 — directe : nom consulaire lu}
\nf{Genre}{Épitaphe ; ensembles A, B, C ; 100 lettres transmises.}
\nf{Références}{ICUR I 1937 = ILCV 2007 (add).}
\nf{Provenance}{Rattachement selon l’EDB : \emph{Incertae originis} (le volume I des ICVR réunit les inscriptions dont la provenance est perdue). Conservation selon l’EDB : \emph{Museo Nazionale Romano}.}
\nf{Support}{Plaque de marbre. 18 × 20 × 0,1 cm (h. × l. × ép.). Inscription gravée. Reproductions signalées par l’EDB : reproduction, Perret 1851 (droits réservés ; non reproduites). Fiche EDB : Filippo Piazzolla, première version 2015.}
\ntext{\lnum{1}Bene merenti\\
\lnum{2}in pace Libera\\
\lnum{3}qu(a)e ⸢v⸣ixit an(nos) VIII\\
\lnum{4}qu(a)e ⸢v⸣ixit an(nos) VIII\\
\lnum{5}neof⸢y⸣ta dep(osita) die\\
\lnum{6}III Nonas Maias con(sulibus)\\
\lnum{7}Gratiano III et Equitio}
\nf{Témoins}{Export EDCS de 2020 : \emph{Bene merenti / in pace Libera / qu(a)e ixit an(nos) VIII / qu(a)e ixit an(nos) VIII / neofta dep(osita) die / III Nonas Maias con(sulibus) / Gratiano III et Equitio}. EDB (EDB31855 ; écarts à examiner) : \emph{benemerenti / in pace Libera / que bixit an(nis) VIII / neofita dep(osita) die / III nonas maias con(sulibus) / Gratiano III et Equitio [texte tronqué dans la liste de résultats]}.}
\nf{Apparat}{⸢v⸣ : B gravé ; ⸢v⸣ : B gravé ; ⸢y⸣ : I gravé}
\nf{Traduction}{À Libera, bien méritante, en paix, qui vécut 8 ans [ligne répétée sur la pierre], néophyte ; déposée le 3e jour avant les nones de mai, sous le consulat de Gratien pour la 3e fois et d’Equitius.}
\nf{Langue}{L’éditeur signale 3 graphies corrigées : B gravé pour v (bêtacisme) ×2 ; I gravé pour y (graphie de y).}
\nf{Chronologie}{Noms consulaires : Gratianus (lu) et Equitius (lu). Dates du jour : déposition : III Nonas Maias = 5 mai 374 (lu). Âge déclaré : 8 a. 3 j. (lu).}
\nf{Commentaire}{La formule in pace est conservée. Vocabulaire religieux : initiation et foi — neofyta (lu).}
\nf{Bibliographie}{ICUR I 1937 ; ILCV 2007 (add). Bibliographie complémentaire selon l’EDB : ILCV, 2007.}
\nf{Statut}{Lecture secondaire (EDCS) ; écarts avec l’EDB à arbitrer sur l’ICVR. L’EDCS signale une photographie (non consultée).}
\end{notice}

\begin{notice}{173}{374}{ICUR n.s. I, 1938 · EDCS-39200382 · EDB31856}
\label{not:173}
\nf{Datation}{374 — directe : lecture partielle d’un nom consulaire}
\nf{Genre}{Épitaphe ; ensembles A, B, C ; 47 lettres transmises ; 68 \% restituées.}
\nf{Références}{ICUR I 1938.}
\nf{Provenance}{Rattachement selon l’EDB : \emph{Incertae originis} (le volume I des ICVR réunit les inscriptions dont la provenance est perdue). Conservation selon l’EDB : \emph{Museo Nazionale Romano}.}
\nf{Support}{Plaque de marbre. 23 × 18 × 4 cm (h. × l. × ép.). Inscription gravée. Fiche EDB : Filippo Piazzolla, première version 2015.}
\ntext{\lnum{1}]IE\\
\lnum{2}[- - -] duo\\
\lnum{3}[- - -]L depo\\
\lnum{4}[- - - cons(ulatu) Gratiani Aug(usti) III et Equ]iti v(iri) c(larissimi)}
\nf{Témoins}{EDB (EDB31856 ; écarts à examiner) : \emph{------ / [---]ịẹ [---] / [---] duo [---] / [---] L depo/[(sit---) cons(ulat---) Gratiani Aug(usti) III et Equ]iti v(iri) c(larissimi) [texte tronqué dans la liste de résultats]}.}
\nf{Traduction}{[…] deux […] ; […] déposé(e) […] sous le consulat de Gratien Auguste pour la 3e fois et d’Equitius, clarissime.}
\nf{Langue}{4 abréviations développées par l’éditeur (dont cons(ulatu), Aug(usti), v(iri), c(larissimi)).}
\nf{Chronologie}{Noms consulaires : Gratianus (restitué) et Equitius (lu en partie).}
\nf{Bibliographie}{ICUR I 1938.}
\nf{Statut}{Lecture secondaire (EDCS) ; écarts avec l’EDB à arbitrer sur l’ICVR. Plus de la moitié des lettres sont restituées.}
\end{notice}

\begin{notice}{174}{374}{ICUR n.s. IV, 9566 · EDCS-35700797 · EDB18370}
\label{not:174}
\nf{Datation}{374 — directe : nom consulaire lu}
\nf{Genre}{Épitaphe ; ensembles A, B, C ; 24 lettres transmises ; 33 \% restituées.}
\nf{Références}{ICUR IV 9566.}
\nf{Provenance}{Rattachement selon l’EDB : \emph{Coem. Callisti pars inferior}.}
\nf{Support}{Enduit de mortier de mortier. Inscription tracée dans le mortier frais. Fiche EDB : Antonio Enrico Felle, première version 2003.}
\ntext{\lnum{1}D(omino) n(ostro) Gr(a)tiano Aug(usto) III e[t Equitio - - -]}
\nf{Témoins}{EDB (EDB18370) : texte concordant aux conventions de notation près (indice de similarité 0,954).}
\nf{Traduction}{Sous (le consulat de) notre seigneur Gratien Auguste pour la 3e fois et d’Equitius […].}
\nf{Langue}{4 abréviations développées par l’éditeur (dont D(omino), n(ostro), Gr(a)tiano, Aug(usto)).}
\nf{Chronologie}{Noms consulaires : Gratianus (lu) et Equitius (restitué).}
\nf{Bibliographie}{ICUR IV 9566.}
\nf{Statut}{Lecture secondaire (EDCS) concordante avec l’EDB : lecture stabilisée entre deux bases, non collationnée sur l’ICVR ni sur le monument.}
\end{notice}

\begin{notice}{175}{374}{ICUR n.s. IV, 9567 · EDCS-35700798 · EDB18371}
\label{not:175}
\nf{Datation}{374 — directe : nom consulaire lu}
\nf{Genre}{Épitaphe ; ensembles A, B, C ; 19 lettres transmises.}
\nf{Références}{ICUR IV 9567.}
\nf{Provenance}{Rattachement selon l’EDB : \emph{Coem. Callisti pars inferior}.}
\nf{Support}{Enduit de mortier de mortier. Inscription tracée dans le mortier frais. Signes décrits par l’EDB : chrisme avec alpha et oméga. Fiche EDB : Antonio Enrico Felle, première version 2003.}
\ntext{\lnum{1}Deposita est\\
\lnum{2}G\\
\lnum{3}rat\\
\lnum{4}ia\\
\lnum{5}n\\
\lnum{6}o [- - -]}
\nf{Témoins}{EDB (EDB18371) : texte concordant aux conventions de notation près (indice de similarité 0,926).}
\nf{Traduction}{A été déposée […] sous (le consulat de) Gratien […].}
\nf{Langue}{Aucune particularité graphique n’est signalée par l’éditeur dans l’état 2022 de la base.}
\nf{Chronologie}{Noms consulaires : Gratianus (lu).}
\nf{Bibliographie}{ICUR IV 9567.}
\nf{Statut}{Lecture secondaire (EDCS) concordante avec l’EDB : lecture stabilisée entre deux bases, non collationnée sur l’ICVR ni sur le monument.}
\end{notice}

\begin{notice}{176}{374}{ICUR n.s. IV, 11105 · EDCS-37200387 · EDB39564}
\label{not:176}
\nf{Datation}{374 — directe : lecture partielle d’un nom consulaire}
\nf{Genre}{Épitaphe ; ensembles A, B, C ; 25 lettres transmises ; 68 \% restituées.}
\nf{Références}{ICUR IV 11105.}
\nf{Provenance}{Rattachement selon l’EDB : \emph{Coem. Callisti area subdialis et incerti loci}. Conservation selon l’EDB : \emph{Coem. Callisti area subdialis et incerti loci}.}
\nf{Support}{Sarcophage de marbre. 10 × 36 cm (h. × l.). Inscription gravée. Fiche EDB : Marida Pierno, première version 2015.}
\ntext{\lnum{1}]\\
\lnum{2}Grat[iano Aug(usto)]\\
\lnum{3}III c[ons(ule)}
\nf{Témoins}{EDB (EDB39564 ; écarts à examiner) : \emph{((pars genii)) / ------/ Grat[iano Augusto] / III [consule]}.}
\nf{Traduction}{[…] sous le consulat de Gratien Auguste pour la 3e fois.}
\nf{Langue}{2 abréviations développées par l’éditeur (dont Aug(usto), cons(ule)).}
\nf{Chronologie}{Noms consulaires : Gratianus (lu en partie).}
\nf{Bibliographie}{ICUR IV 11105.}
\nf{Statut}{Lecture secondaire (EDCS) ; écarts avec l’EDB à arbitrer sur l’ICVR. Plus de la moitié des lettres sont restituées.}
\end{notice}

\begin{notice}{177}{374}{ICUR n.s. IV, 11106 · EDCS-37200388 · EDB39565}
\label{not:177}
\nf{Datation}{374 — directe : lecture partielle d’un nom consulaire}
\nf{Genre}{Épitaphe ; ensembles A, B, C ; 47 lettres transmises ; 68 \% restituées.}
\nf{Références}{ICUR IV 11106.}
\nf{Provenance}{Rattachement selon l’EDB : \emph{Coem. Callisti area subdialis et incerti loci}. Conservation selon l’EDB : \emph{Coem. Callisti area subdialis et incerti loci}.}
\nf{Support}{Plaque de marbre. 37 × 18 cm (h. × l.). Inscription gravée. Fiche EDB : Marida Pierno, première version 2015.}
\ntext{\lnum{1}] qui [- - -]\\
\lnum{2}[- - - K]al(endas) Oct(obres) [Gratiano]\\
\lnum{3}[Aug(usto) I]II et Equ[itio conss(ulibus)]}
\nf{Témoins}{EDB (EDB39565 ; écarts à examiner) : \emph{[---] qui [---] / [k]al(endas) oct(obres) [Gratiano] / [Augusto I]II et Equ[itio consulibus]}.}
\nf{Traduction}{[…] qui […] ; […] calendes d’octobre, sous le consulat de Gratien Auguste pour la 3e fois et d’Equitius.}
\nf{Langue}{4 abréviations développées par l’éditeur (dont Kal(endas), Oct(obres), Aug(usto), conss(ulibus)).}
\nf{Chronologie}{Noms consulaires : Gratianus (restitué) et Equitius (lu en partie). Dates du jour : événement indéterminé : … Kalendas Octobres (partiellement restitué).}
\nf{Bibliographie}{ICUR IV 11106.}
\nf{Statut}{Lecture secondaire (EDCS) ; écarts avec l’EDB à arbitrer sur l’ICVR. Plus de la moitié des lettres sont restituées.}
\end{notice}

\begin{notice}{178}{374}{ICUR n.s. IV, 12243 · EDCS-37600678 · EDB8839}
\label{not:178}
\nf{Datation}{374 — directe : lecture partielle d’un nom consulaire}
\nf{Genre}{Épitaphe ; ensembles A, B, C ; 62 lettres transmises ; 50 \% restituées.}
\nf{Références}{ICUR IV 12243.}
\nf{Provenance}{Rattachement selon l’EDB : \emph{Coem. anonymum ad viam Ardeatinam}. Conservation selon l’EDB : \emph{Coem. anonymum ad viam Ardeatinam}.}
\nf{Support}{Plaque de marbre. 30 × 23 × 2,4 cm (h. × l. × ép.). Inscription gravée. Fiche EDB : Marida Pierno, première version 2003.}
\ntext{\lnum{1}plus min]us XX[- - -]\\
\lnum{2}[- - - c]ollaud[ant - - -]\\
\lnum{3}[- - - Gratiano III et E]quitio co[ns(ulibus) - - -]\\
\lnum{4}[- - - be]ne mere[nti}
\nf{Témoins}{EDB (EDB8839 ; écarts à examiner) : \emph{------ / [--- plus min]us XX [--- ]/ [--- c]ollaud[ant] / [--- Gratiano III et E]quitio co[ns(ulibu)s)] / [--- be]nemere[nti]}.}
\nf{Traduction}{[…] environ 20 […] ; […] louent ensemble […] ; […] sous le consulat de Gratien pour la 3e fois et d’Equitius […] ; […] bien méritant.}
\nf{Langue}{1 abréviation développée par l’éditeur (dont cons(ulibus)).}
\nf{Chronologie}{Noms consulaires : Gratianus (restitué) et Equitius (lu en partie).}
\nf{Bibliographie}{ICUR IV 12243.}
\nf{Statut}{Lecture secondaire (EDCS) ; écarts avec l’EDB à arbitrer sur l’ICVR.}
\end{notice}

\annee{375}

\begin{notice}{179}{375}{ICUR n.s. V, 13109 · EDCS-19900430 · EDB5945}
\label{not:179}
\nf{Datation}{375 — directe : nom consulaire lu}
\nf{Genre}{Épitaphe ; ensembles A, B, C ; 113 lettres transmises ; 9 \% restituées.}
\nf{Références}{CIL VI 32045a (p 4806) = ICUR V 13109 (p 414) = ILCV 95 = AE 2009, 145.}
\nf{Provenance}{Rattachement selon l’EDB : \emph{Coem. subterraneum ad Catacumbas}. Conservation selon l’EDB : \emph{Coem. subterraneum ad Catacumbas}.}
\nf{Support}{Sarcophage de marbre. Inscription gravée. Signes décrits par l’EDB : chrisme et chrisme avec alpha et oméga. Reproductions signalées par l’EDB : photographie PCAS (droits réservés ; non reproduites). Fiche EDB : Carlo Carletti, première version 2014.}
\ntext{\lnum{1}Simplicius v(ir) c(larissimus) et Didyme c(larissima) f(emina) et Innocentius [dormi]entes in pace et\\
\lnum{2}Roscia Calc(h)edonia c(larissima) f(emina) d(eposita) die V Non(as) Mai(as) post con(sulatum) d(omini) n(ostri) Gra[tiani] Aug(usti) III et Equiti}
\nf{Témoins}{EDB (EDB5945) : texte concordant aux conventions de notation près (indice de similarité 0,96).}
\nf{Traduction}{Simplicius, clarissime, et Didymè, clarissime, et Innocentius, endormis en paix ; et Roscia Calchedonia, clarissime, déposée le 5e jour avant les nones de mai, après le consulat de notre seigneur Gratien Auguste pour la 3e fois et d’Equitius.}
\nf{Langue}{14 abréviations développées par l’éditeur (dont v(ir), c(larissimus), c(larissima), f(emina), Calc(h)edonia).}
\nf{Chronologie}{Noms consulaires : Gratianus (lu en partie) et Equitius (lu). Datation par post-consulat. Dates du jour : déposition : die V Nonas Maias = 3 mai 375 (lu).}
\nf{Commentaire}{La formule in pace est conservée. Vocabulaire religieux : repos et sommeil de la mort — dormientes (partiellement restitué). Rang sénatorial d’une personne commémorée (simplicius vir clarissimus et d ; ir clarissimus et didyme clarissima femina et i).}
\nf{Bibliographie}{CIL VI 32045a (p 4806) ; ICUR V 13109 (p 414) ; ILCV 95 ; AE 2009, 145. Bibliographie complémentaire selon l’EDB : CIL VI, 32045a ; ILCV, 95.}
\nf{Statut}{Lecture secondaire (EDCS) concordante avec l’EDB : lecture stabilisée entre deux bases, non collationnée sur l’ICVR ni sur le monument. L’EDCS signale une photographie (non consultée).}
\end{notice}

\begin{notice}{180}{375}{ICUR n.s. I, 1847 · EDCS-39200296 · EDB31761}
\label{not:180}
\nf{Datation}{375 — formule transmise par l’EDB (partiellement restituée)}
\nf{Genre}{Épitaphe ; ensembles A, B, C ; 24 lettres transmises ; 8 \% restituées.}
\nf{Références}{ICUR I 1847,02.}
\nf{Provenance}{Rattachement selon l’EDB : \emph{Incertae originis} (le volume I des ICVR réunit les inscriptions dont la provenance est perdue). Conservation selon l’EDB : \emph{Musei Vaticani: Lapidario Cristiano}.}
\nf{Support}{Plaque de marbre. 15 × 17 cm (h. × l.). Inscription gravée. Fiche EDB : Filippo Piazzolla, première version 2015.}
\ntext{\lnum{1}]NSACP[- - -]\\
\lnum{2}[- - -] Dec(em)b(res) post [- - -]\\
\lnum{3}[- - -]tiae qu[ae - - -]\\
\lnum{4}[- - -]IIX[}
\nf{Témoins}{EDB (EDB31761 ; divergence forte) : \emph{------ / [---]++SAC+O+++[---] / [---] dec(em)b(r--) post c[onsulatum ---] / [--- et] Fl(avii) Aequ[itii ---] / [---] vix(it) [---] / ------}.}
\nf{Traduction}{[…] décembre, après […] ; […] qui […] ; […] 8 (?).}
\nf{Langue}{2 abréviations développées par l’éditeur (dont Dec(em)b, (res)).}
\nf{Chronologie}{Le texte EDCS ne conserve aucun nom consulaire. La fiche EDB (EDB31761) transmet \emph{[---] dec(em)b(r--) post c[onsulatum ---] / [--- et] Fl(avii) Aequ[itii ---]} : la fin du nom d'Equitius est lue, dans une formule de post-consulat qui convient à 375 (après Gratien III et Equitius).}
\nf{Bibliographie}{ICUR I 1847,02. Bibliographie complémentaire selon l’EDB : Ferrua 1979b, 53 ; Lega 2000.}
\nf{Statut}{Lecture secondaire (EDCS) ; divergence forte avec l’EDB : le texte ne doit pas être exploité sans contrôle de l’édition imprimée.}
\end{notice}

\begin{notice}{181}{375}{ICUR n.s. I, 2093 · EDCS-39300315 · EDB32058}
\label{not:181}
\nf{Datation}{375 — directe : nom consulaire lu}
\nf{Genre}{Épitaphe ; ensembles A, B, C ; 145 lettres transmises ; 30 \% restituées.}
\nf{Références}{ICUR I 2093 = ILCV 4369 (add).}
\nf{Provenance}{Rattachement selon l’EDB : \emph{Incertae originis} (le volume I des ICVR réunit les inscriptions dont la provenance est perdue). Conservation selon l’EDB : \emph{Musei Vaticani: Galleria Lapidaria}.}
\nf{Support}{Plaque de marbre. 50 × 70 cm (h. × l.). Inscription gravée. Fiche EDB : Filippo Piazzolla, première version 2015.}
\ntext{\lnum{1}Innox\\
\lnum{2}[- - -]antia qu(a)e vix\{s\}it [- - -]\\
\lnum{3}[- - -] d(e)pos(ita ?) Gratiano III e[t Equitio cons(ulibus) - - -]\\
\lnum{4}[- - - Calli]ste(?) soror qu(a)e vix\{s\}[it - - -]\\
\lnum{5}[- - -]III consules ips[i - - -]\\
\lnum{6}[- - -] XVIII K(a)l(endas) [- - -]\\
\lnum{7}[- - - cons]ules s(upra) s(cripti) d(epositus ?) Callistratu[s qui]\\
\lnum{8}[vixit ann(os)] XVII me(nses) VI die(s) XXI[}
\nf{Témoins}{EDB (EDB32058) : texte concordant aux conventions de notation près (indice de similarité 0,941).}
\nf{Apparat}{\{s\} : lettres superflues ; \{s\} : lettres superflues}
\nf{Traduction}{Innocent (?) […] ; […]antia, qui vécut […] ; déposée (sous le consulat de) Gratien pour la 3e fois et Equitius […] ; […] Callistè, sa sœur, qui vécut […] ; […] les mêmes consuls […] ; […] le 18e jour avant les calendes […] ; […] sous les consuls susdits, a été déposé Callistratus, qui vécut 17 ans 6 mois 21 jours.}
\nf{Langue}{Des lettres gravées sont jugées superflues par l’éditeur (entre accolades). 14 abréviations développées par l’éditeur (dont qu(a)e, d(e)pos, (ita ?), cons(ulibus), Calliste(?)).}
\nf{Chronologie}{Noms consulaires : Gratianus (lu) et Equitius (restitué). Âge déclaré : 17 a. 6 m. 21 j. (partiellement restitué).}
\nf{Bibliographie}{ICUR I 2093 ; ILCV 4369 (add). Bibliographie complémentaire selon l’EDB : ILCV, 4369 ; Lega 2000.}
\nf{Statut}{Lecture secondaire (EDCS) concordante avec l’EDB : lecture stabilisée entre deux bases, non collationnée sur l’ICVR ni sur le monument.}
\end{notice}

\begin{notice}{182}{375}{ICUR n.s. I, 3184 · EDCS-40300330 · EDB33136}
\label{not:182}
\nf{Datation}{375 — directe : lecture partielle d’un nom consulaire}
\nf{Genre}{Épitaphe ; ensembles A, B, C ; 127 lettres transmises ; 17 \% restituées.}
\nf{Références}{ICUR I 3184 = ILCV 3305.}
\nf{Provenance}{Rattachement selon l’EDB : \emph{Incertae originis} (le volume I des ICVR réunit les inscriptions dont la provenance est perdue).}
\nf{Support}{Support non précisé par l’EDB. Fiche EDB : Filippo Piazzolla, première version 2015.}
\ntext{\lnum{1}f]ati munus [- - -]\\
\lnum{2}[- - -]um corpori MENBI[- - -]\\
\lnum{3}[- - -]tes reliqui⸢t⸣ supe[- - -]\\
\lnum{4}[- - -] quae germine vivi[- - -]\\
\lnum{5}[- - - vixi]t annos XXIII men[ses - - -]\\
\lnum{6}[- - - Id]us Maias pos(t) consula[tum - - -]\\
\lnum{7}[- - - Gr]atiano Aug(usto) III et Fl(avio) \{A\}E[quitio - - -]\\
\lnum{8}[- - - in] pace}
\nf{Témoins}{Export EDCS de 2020 : \emph{f]ati munus [3] / [3]um corpori MENBI[3] / [3]tes reliqui supe[3] / [3] quae germine vivi[3] / [3 vixi]t annos XXIII men[ses 3] / [3 Id]us Maias pos(t) consula[tum 3] / [3 Gr]atiano Aug(usto) III et Fl(avio) \{A\}E[quitio 3] / [3 in] pace}. EDB (EDB33136 ; écarts à examiner) : \emph{[--- f]ati munus [---] / [---]ụm corpori menbṛ[---] / [---]tes reliquid supẹ[---] / [---] quae germine vivi[---] / [--- vixi]ṭ annos XXIII meṇ[ses --- ] / [--- dep(osit--) ---id]us maias pos(t) consula[tum ---] / [--- Gr]atiano Aug(usto) III et Fl(avio) Ae[quitio ---] / [in] pace [texte tronqué dans la liste de résultats]}.}
\nf{Apparat}{⸢t⸣ : D gravé ; \{A\} : lettres superflues}
\nf{Traduction}{[…] le présent du destin […] au corps […] ; […] laissa […] ; […] qui, de la semence vivante (?) […] ; vécut 23 ans […] mois ; […] ides de mai, après le consulat de Gratien Auguste pour la 3e fois et de Flavius Equitius ; […] en paix.}
\nf{Langue}{L’éditeur signale une graphie corrigée : D gravé pour t. Des lettres gravées sont jugées superflues par l’éditeur (entre accolades). 3 abréviations développées par l’éditeur (dont pos(t), Aug(usto), Fl(avio)).}
\nf{Chronologie}{Noms consulaires : Gratianus (lu en partie) et Equitius (restitué). Datation par post-consulat. Dates du jour : événement indéterminé : … Idus Maias (partiellement restitué). Âge déclaré : 23 a. (partiellement restitué).}
\nf{Commentaire}{La formule in pace est attestée seulement par restitution.}
\nf{Bibliographie}{ICUR I 3184 ; ILCV 3305. Bibliographie complémentaire selon l’EDB : ILCV, 3305.}
\nf{Statut}{Lecture secondaire (EDCS) ; écarts avec l’EDB à arbitrer sur l’ICVR.}
\end{notice}

\begin{notice}{183}{375}{ICUR n.s. II, 6009 · EDCS-37700196 · EDB21312}
\label{not:183}
\nf{Datation}{375 — directe : nom consulaire lu}
\nf{Genre}{Épitaphe ; ensembles A, B, C ; 124 lettres transmises.}
\nf{Références}{ICUR II 6009 = ILCV 775 (em).}
\nf{Provenance}{Rattachement selon l’EDB : \emph{Coem. s.Zenonis}. Conservation selon l’EDB : \emph{Collegio Teutonico}.}
\nf{Support}{Plaque de marbre. 38 × 31 cm (h. × l.). Inscription gravée. Signes décrits par l’EDB : chrisme. Fiche EDB : Maria Cristina Grisanzio, première version 2006.}
\ntext{\lnum{1}Depossio Petroni(a)es(!) VI Kal(endas)\\
\lnum{2}Dec(embres) post consulatum\\
\lnum{3}Gratiani Aug(usti) III et Fl(avi) Equiti\\
\lnum{4}qu(a)e vixit virgo annos p(lus) m(inus)\\
\lnum{5}XVIIII in pace est de regio\\
\lnum{6}ne VIIII a lacu cunic(u)li}
\nf{Témoins}{EDB (EDB21312) : texte concordant aux conventions de notation près (indice de similarité 0,978).}
\nf{Apparat}{(!) : forme signalée par l’éditeur}
\nf{Traduction}{Déposition de Petronia le 6e jour avant les calendes de décembre, après le consulat de Gratien Auguste pour la 3e fois et de Flavius Equitius ; elle vécut vierge environ 19 ans ; elle est en paix ; (elle était) de la IXe région, du (quartier du) lac du Cuniculus (?).}
\nf{Langue}{L’éditeur signale au moins une forme anormale (sic). 9 abréviations développées par l’éditeur (dont Petroni(a)es, Kal(endas), Dec(embres), Aug(usti), Fl(avi), qu(a)e).}
\nf{Chronologie}{Noms consulaires : Gratianus (lu) et Equitius (lu). Datation par post-consulat. Dates du jour : déposition : VI Kalendas Decembres = 26 novembre 375 (lu). Âge déclaré : env. 19 a. (lu).}
\nf{Commentaire}{La formule in pace est conservée. Commentaire développé : dossier p.~\pageref{dos:183}.}
\nf{Bibliographie}{ICUR II 6009 ; ILCV 775 (em). Bibliographie complémentaire selon l’EDB : ILCV, 775.}
\nf{Statut}{Lecture secondaire (EDCS) concordante avec l’EDB : lecture stabilisée entre deux bases, non collationnée sur l’ICVR ni sur le monument. L’EDCS signale une photographie (non consultée).}
\end{notice}

\begin{notice}{184}{375}{ICUR n.s. II, 6032 · EDCS-34200658 · EDB17433}
\label{not:184}
\nf{Datation}{375 — directe : nom consulaire lu}
\nf{Genre}{Épitaphe ; ensembles A, B, C ; 147 lettres transmises.}
\nf{Références}{ICUR II 6032 = ILCV 4226 (em).}
\nf{Provenance}{Rattachement selon l’EDB : \emph{Coem. Commodillae}. Conservation selon l’EDB : \emph{Musei Vaticani: Lapidario Cristiano}.}
\nf{Support}{Plaque de marbre. Inscription gravée. Reproductions signalées par l’EDB : reproduction, Perret 1851 (droits réservés ; non reproduites). Fiche EDB : Carlo Carletti, première version 2003.}
\ntext{\lnum{1}Petroniae dignae co(n)iugi qu(a)e vixit annis\\
\lnum{2}XXXI et fecit cum co⸢m⸣pare suo m(enses) X d(ies ?) V\\
\lnum{3}Kal(endis) No⸢v⸣(embribus) pos(t) cons\{s\}(ulatu) Gratiani ter et Equiti\\
\lnum{4}Ursus maritus sibi et innocenti co\\
\lnum{5}mpari fecit ⸢q⸣(ui)ces⸢c⸣uet in pace}
\nf{Témoins}{Export EDCS de 2020 : \emph{Petroniae dignae co(n)iugi qu(a)e vixit annis / XXXI et fecit cum copare suo m(enses) X d(ies ?) V / Kal(endis) No(embribus) pos(t) cons\{s\}(ulatu) Gratiani ter et Equiti / Ursus maritus sibi et innocenti co/mpari fecit (ui)cesuet in pace}. EDB (EDB17433) : texte concordant aux conventions de notation près (indice de similarité 0,978).}
\nf{Apparat}{⸢m⸣ : N gravé ; ⸢v⸣ : B gravé ; ⸢q⸣ : C gravé ; ⸢c⸣ : Q gravé ; \{s\} : lettres superflues}
\nf{Traduction}{À Petronia, épouse digne, qui vécut 31 ans et a vécu avec son compagnon 10 mois 5 jours ; aux calendes de novembre, après le consulat de Gratien pour la 3e fois et d’Equitius ; Ursus, son mari, a fait (ce monument) pour lui et pour sa compagne innocente ; qu’elle repose en paix.}
\nf{Langue}{L’éditeur signale 4 graphies corrigées : N gravé pour m (n pour m) ; B gravé pour v (bêtacisme) ; C gravé pour q ; Q gravé pour c. Des lettres gravées sont jugées superflues par l’éditeur (entre accolades). 9 abréviations développées par l’éditeur (dont co(n)iugi, qu(a)e, m(enses), d(ies ?), Kal(endis), (embribus)).}
\nf{Chronologie}{Noms consulaires : Gratianus (lu) et Equitius (lu). Datation par post-consulat. Dates du jour : achat : Kalendis Novembribus = 1 novembre 375 (lu). Âge déclaré : 31 a. (lu).}
\nf{Commentaire}{La formule in pace est conservée.}
\nf{Bibliographie}{ICUR II 6032 ; ILCV 4226 (em).}
\nf{Statut}{Lecture secondaire (EDCS) concordante avec l’EDB : lecture stabilisée entre deux bases, non collationnée sur l’ICVR ni sur le monument. L’EDCS signale une photographie (non consultée).}
\end{notice}

\begin{notice}{185}{375}{ICUR n.s. III, 8725 · EDCS-35900969 · EDB20993}
\label{not:185}
\nf{Datation}{375 — directe : nom consulaire lu}
\nf{Genre}{Épitaphe ; ensembles A, B, C ; 128 lettres transmises.}
\nf{Références}{ICUR III 8725 = ILCV 4258a.}
\nf{Provenance}{Rattachement selon l’EDB : \emph{Tituli qui in coem. Callisti reperti traduntur}.}
\nf{Support}{Plaque de marbre. Inscription gravée. Signes décrits par l’EDB : chrisme avec alpha et oméga. Fiche EDB : Carlo Carletti, première version 2003.}
\ntext{\lnum{1}Apro qui vixit annos XLVIII menses IIII dies\\
\lnum{2}XVI depositus in pace VI Kalendas Martias\\
\lnum{3}post cons(ulatum) Gratiani et Equitii Limfiria marito\\
\lnum{4}fecit mecum annos XX}
\nf{Témoins}{EDB (EDB20993) : texte concordant aux conventions de notation près (indice de similarité 1).}
\nf{Traduction}{À Aper, qui vécut 48 ans 4 mois 16 jours ; déposé en paix le 6e jour avant les calendes de mars, après le consulat de Gratien et d’Equitius ; Lymphyria (?) a fait (ce monument) pour son mari ; il a vécu avec moi 20 ans.}
\nf{Langue}{1 abréviation développée par l’éditeur (dont cons(ulatum)).}
\nf{Chronologie}{Noms consulaires : Gratianus (lu) et Equitius (lu). Datation par post-consulat. Dates du jour : déposition : VI Kalendas Martias = 24 février 375 (lu). Âge déclaré : 48 a. 4 m. 16 j. (lu).}
\nf{Commentaire}{La formule in pace est conservée.}
\nf{Bibliographie}{ICUR III 8725 ; ILCV 4258a.}
\nf{Statut}{Lecture secondaire (EDCS) concordante avec l’EDB : lecture stabilisée entre deux bases, non collationnée sur l’ICVR ni sur le monument.}
\end{notice}

\begin{notice}{186}{375}{ICUR n.s. III, 8726 · EDCS-35900970 · EDB20994}
\label{not:186}
\nf{Datation}{375 — directe : nom consulaire lu}
\nf{Genre}{Épitaphe ; ensembles A, B, C ; 91 lettres transmises ; 6 \% restituées.}
\nf{Références}{ICUR III 8726.}
\nf{Provenance}{Rattachement selon l’EDB : \emph{Tituli qui in coem. Callisti reperti traduntur}.}
\nf{Support}{Plaque de marbre. Inscription gravée. Signes décrits par l’EDB : chrisme. Fiche EDB : Carlo Carletti, première version 2014.}
\ntext{\lnum{1}Dulcissimo filio Vita⸢l⸣[ini]\\
\lnum{2}[q]ui vixi⸢t⸣ annos XXII d(ies) XXVII\\
\lnum{3}[i]n p(ace) XII Kal(endas) Decemb(res) pos⸢t⸣ cons(ulatum)\\
\lnum{4}d(ominorum) n(ostrorum) Gratiani Aug(usti) et Aequiti}
\nf{Témoins}{Export EDCS de 2020 : \emph{Dulcissimo filio Vita[ini] / [q]ui vixi annos XXII d(ies) XXVII / [i]n p(ace) XII Kal(endas) Decemb(res) pos cons(ulatum) / d(ominorum) n(ostrorum) Gratiani Aug(usti) et Aequiti}. EDB (EDB20994) : texte concordant aux conventions de notation près (indice de similarité 0,983).}
\nf{Apparat}{⸢l⸣ : I gravé ; ⸢t⸣ : I gravé ; ⸢t⸣ : I gravé}
\nf{Traduction}{À son fils très doux Vitalinus, qui vécut 22 ans 27 jours ; en paix ; le 12e jour avant les calendes de décembre, après le consulat de nos seigneurs Gratien Auguste et Equitius.}
\nf{Langue}{L’éditeur signale 3 graphies corrigées : I gravé pour l ; I gravé pour t ×2.}
\nf{Chronologie}{Noms consulaires : Gratianus (lu) et Equitius (lu). Datation par post-consulat. Dates du jour : événement indéterminé : XII Kalendas Decembres = 20 novembre 375 (lu). Âge déclaré : 22 a. 27 j. (partiellement restitué).}
\nf{Commentaire}{La formule in pace est attestée seulement par restitution.}
\nf{Bibliographie}{ICUR III 8726.}
\nf{Statut}{Lecture secondaire (EDCS) concordante avec l’EDB : lecture stabilisée entre deux bases, non collationnée sur l’ICVR ni sur le monument.}
\end{notice}

\begin{notice}{187}{375}{ICUR n.s. IV, 11767 · EDCS-37400856 · EDB876}
\label{not:187}
\nf{Datation}{375 — directe : lecture partielle d’un nom consulaire}
\nf{Genre}{Épitaphe ; ensembles A, B, C ; 66 lettres transmises ; 65 \% restituées.}
\nf{Références}{ICUR IV 11767.}
\nf{Provenance}{Rattachement selon l’EDB : \emph{Coem. anonymum ad viam Appiam}. Conservation selon l’EDB : \emph{Coem. anonymum ad viam Appiam}.}
\nf{Support}{Plaque de marbre. Inscription gravée. Reproductions signalées par l’EDB : photographie PCAS (droits réservés ; non reproduites). Fiche EDB : Maria Cristina Grisanzio, première version 2003.}
\ntext{\lnum{1}depos]itus\\
\lnum{2}[- - - Kalendas Ian]uarias\\
\lnum{3}[in pace post con]s(ulatum) Gra\\
\lnum{4}[tiani Aug(usti) III et Eq]uiti}
\nf{Témoins}{EDB (EDB876 ; écarts à examiner) : \emph{------(?) / [---]+++ / [--- depos]itus / [--- ian]uarias / [--- con]s(ulatu) Gra/[tiani Aug(usti) III et Eq]uiti}.}
\nf{Traduction}{Déposé […] calendes de janvier, en paix, après le consulat de Gratien Auguste pour la 3e fois et d’Equitius.}
\nf{Langue}{2 abréviations développées par l’éditeur (dont cons(ulatum), Aug(usti)).}
\nf{Chronologie}{Noms consulaires : Gratianus (lu en partie) et Equitius (lu en partie). Datation par post-consulat. Dates du jour : déposition : … Kalendas Ianuarias (restitué).}
\nf{Commentaire}{La formule in pace est attestée seulement par restitution.}
\nf{Bibliographie}{ICUR IV 11767. Bibliographie complémentaire selon l’EDB : CLRE, 646.}
\nf{Statut}{Lecture secondaire (EDCS) ; écarts avec l’EDB à arbitrer sur l’ICVR. Plus de la moitié des lettres sont restituées. L’EDCS signale une photographie (non consultée).}
\end{notice}

\begin{notice}{188}{375}{ICUR n.s. IV, 12531 · EDCS-37801594 · EDB42433}
\label{not:188}
\nf{Datation}{375 — directe : nom consulaire lu}
\nf{Genre}{Épitaphe ; ensembles A, B, C ; 96 lettres transmises ; 46 \% restituées.}
\nf{Références}{ICUR IV 12531.}
\nf{Provenance}{Rattachement selon l’EDB : \emph{inter vias Appiam et Ardeatinam}. Conservation selon l’EDB : \emph{Musei Vaticani: Lapidario Cristiano}.}
\nf{Support}{La fiche EDB ne figure pas dans l’état 2019 des données EpiDoc ; support non documenté.}
\ntext{\lnum{1}[- - - P]raetextatu[s]\\
\lnum{2}[- - -] vix(it) ann(os) ann(os) LXXVII\\
\lnum{3}[m(enses) - - -] d(ies) XI dep(ositus) IIII Id(us) Ian(uarias)\\
\lnum{4}[pos]t cons(ulatum) Gratian[i]\\
\lnum{5}[Aug(usti) III et Equiti v(iri) c(larissimi)]}
\nf{Témoins}{EDB (EDB42433) : texte concordant aux conventions de notation près (indice de similarité 0,927).}
\nf{Traduction}{[…] Praetextatus […] vécut 77 ans […] mois 11 jours ; déposé le 4e jour avant les ides de janvier, après le consulat de Gratien Auguste pour la 3e fois et d’Equitius, clarissime.}
\nf{Langue}{12 abréviations développées par l’éditeur (dont vix(it), ann(os), m(enses), d(ies), dep(ositus)).}
\nf{Chronologie}{Noms consulaires : Gratianus (lu) et Equitius (restitué). Datation par post-consulat. Dates du jour : déposition : IIII Idus Ianuarias = 10 janvier 375 (lu). Âge déclaré : 77 a. 11 j. (partiellement restitué).}
\nf{Bibliographie}{ICUR IV 12531. Bibliographie complémentaire selon l’EDB : Dresken-Weiland 2003, 221.}
\nf{Statut}{Lecture secondaire (EDCS) concordante avec l’EDB : lecture stabilisée entre deux bases, non collationnée sur l’ICVR ni sur le monument.}
\end{notice}

\annee{376}

\begin{notice}{189}{376}{ICUR n.s. VIII, 20794 · EDCS-24100416 · EDB15528}
\label{not:189}
\nf{Datation}{376 — directe : nom consulaire lu}
\nf{Genre}{Épitaphe ; ensembles A, B, C ; 117 lettres transmises ; 56 \% restituées.}
\nf{Références}{CIL VI 33902 = ICUR VIII 20794 = ILCV 657 (em).}
\nf{Provenance}{Rattachement selon l’EDB : \emph{Coem. s.Agnetis}.}
\nf{Support}{Plaque de marbre. Inscription gravée. Fiche EDB : Carolina Ventura, première version 2003.}
\ntext{\lnum{1}]R T[- - -]\\
\lnum{2}[- - -]rentis vel[- - -]\\
\lnum{3}[- - -]untinq(ue) [- - -]\\
\lnum{4}[qua]dratari(i) AM[- - -]\\
\lnum{5}[- - -] qui vixit [ann(os) - - -]\\
\lnum{6}[menses] sex d(ies) XII d(onum) p(osuit) [- - -]\\
\lnum{7}[Valente A]ug(usto) Valenti[niano Aug(usto) co(n)s(ulibus)\\
\lnum{8}[par]entis et [a]mici valete [}
\nf{Témoins}{EDB (EDB15528 ; écarts à examiner) : \emph{((ascia)) ((falx)) / [---]r ((scalprum)) t[------] / [---]rentis vel [---] / [---]unt in q[ua]/dratariam [---] / [---] qui vixit [---] / [---] sex d(ies) XII d(e)p(ositus) [---] / [--- A]ug(usto) V et Valenti[niano consulibus] / [par]entis et âmici valete [---]}.}
\nf{Traduction}{[…] ou […] ; […] des tailleurs de pierre (quadratarii) […] ; […] qui vécut […] ans 6 mois 12 jours ; a posé (ce) don […] ; sous le consulat de Valens Auguste et de Valentinien Auguste ; parents et amis, adieu.}
\nf{Langue}{10 abréviations développées par l’éditeur (dont untinq(ue), quadratari(i), ann(os), d(ies), d(onum), p(osuit)).}
\nf{Chronologie}{Noms consulaires : Valentinianus (lu en partie) et Valens (lu). Âge déclaré : 6 m. 12 j. (partiellement restitué).}
\nf{Commentaire}{Fonction ou métier mentionné : quadratarii.}
\nf{Bibliographie}{CIL VI 33902 ; ICUR VIII 20794 ; ILCV 657 (em).}
\nf{Statut}{Lecture secondaire (EDCS) ; écarts avec l’EDB à arbitrer sur l’ICVR. Plus de la moitié des lettres sont restituées.}
\end{notice}

\begin{notice}{190}{376}{ICUR n.s. I, 3186 · EDCS-40300332 · EDB33138}
\label{not:190}
\nf{Datation}{376 — directe : lecture partielle d’un nom consulaire}
\nf{Genre}{Épitaphe ; ensembles A, B, C ; 111 lettres transmises ; 75 \% restituées.}
\nf{Références}{ICUR I 3186.}
\nf{Provenance}{Rattachement selon l’EDB : \emph{Incertae originis} (le volume I des ICVR réunit les inscriptions dont la provenance est perdue).}
\nf{Support}{Plaque de marbre. Fiche EDB : Filippo Piazzolla, première version 2015.}
\ntext{\lnum{1}[Vale]ntinian[o bene merenti in pace]\\
\lnum{2}[qui vi]xit annos [- - -]\\
\lnum{3}[- - -]I Kal(endas) N[o]vem[bres dd(ominis) nn(ostris) Valenti]\\
\lnum{4}[et Valen]tin[iano - - -] Au[gg(ustis) cons(ulibus)]}
\nf{Témoins}{EDB (EDB33138) : texte concordant aux conventions de notation près (indice de similarité 0,937).}
\nf{Traduction}{À Valentinianus, bien méritant, en paix, qui vécut […] ans ; […] calendes de novembre, sous le consulat de nos seigneurs Valens et Valentinien […] Augustes.}
\nf{Langue}{5 abréviations développées par l’éditeur (dont Kal(endas), dd(ominis), nn(ostris), Augg(ustis), cons(ulibus)).}
\nf{Chronologie}{Noms consulaires : Valentinianus (lu en partie) et Valens (lu en partie). Dates du jour : événement indéterminé : …I Kalendas Novembres = 1 novembre 376 (partiellement restitué).}
\nf{Commentaire}{La formule in pace est attestée seulement par restitution.}
\nf{Bibliographie}{ICUR I 3186.}
\nf{Statut}{Lecture secondaire (EDCS) concordante avec l’EDB : lecture stabilisée entre deux bases, non collationnée sur l’ICVR ni sur le monument. Plus de la moitié des lettres sont restituées.}
\end{notice}

\begin{notice}{191}{376}{ICUR n.s. IV, 9568 · EDCS-36101095 · EDB18372}
\label{not:191}
\nf{Datation}{376 — directe : nom consulaire lu}
\nf{Genre}{Épitaphe ; ensembles A, B, C ; 76 lettres transmises.}
\nf{Références}{ICUR IV 9568 = ILCV 1328 (em).}
\nf{Provenance}{Rattachement selon l’EDB : \emph{Coem. Callisti pars inferior}.}
\nf{Support}{Plaque de marbre. Inscription gravée. Signes décrits par l’EDB : chrisme. Fiche EDB : Antonio Enrico Felle, première version 2003.}
\ntext{\lnum{1}Eusebius qui vixit\\
\lnum{2}ann⸢o⸣s XXXVIII d(e)p(ositus) V Kal(endas)\\
\lnum{3}Oct(obres) in pace Valent⸢e⸣\\
\lnum{4}V et Valentini\\
\lnum{5}ano Iuniori cons(ulibus)}
\nf{Témoins}{Export EDCS de 2020 : \emph{Eusebius qui vixit / anns XXXVIII d(e)p(ositus) V Kal(endas) / Oct(obres) in pace Valent / V et Valentini/ano Iuniori cons(ulibus)}. EDB (EDB18372) : texte concordant aux conventions de notation près (indice de similarité 0,989).}
\nf{Apparat}{⸢o⸣ : V gravé ; ⸢e⸣ : I gravé}
\nf{Traduction}{Eusebius, qui vécut 38 ans ; déposé le 5e jour avant les calendes d’octobre, en paix, sous le consulat de Valens pour la 5e fois et de Valentinien le Jeune.}
\nf{Langue}{L’éditeur signale 2 graphies corrigées : V gravé pour o (confusion o/u) ; I gravé pour e (confusion e/i).}
\nf{Chronologie}{Noms consulaires : Valentinianus (lu) et Valens (lu). Ne sont pas des consuls : Eusebius qui vixit (défunt). Dates du jour : déposition : V Kalendas Octobres = 27 septembre 376 (lu). Âge déclaré : 38 a. (lu).}
\nf{Commentaire}{La formule in pace est conservée.}
\nf{Bibliographie}{ICUR IV 9568 ; ILCV 1328 (em).}
\nf{Statut}{Lecture secondaire (EDCS) concordante avec l’EDB : lecture stabilisée entre deux bases, non collationnée sur l’ICVR ni sur le monument.}
\end{notice}

\begin{notice}{192}{376}{ICUR n.s. IV, 9569 · EDCS-35700799 · EDB17958}
\label{not:192}
\nf{Datation}{376 — directe : lecture partielle d’un nom consulaire}
\nf{Genre}{Épitaphe ; ensembles A, B, C ; 52 lettres transmises ; 79 \% restituées.}
\nf{Références}{ICUR IV 9569 (p 533).}
\nf{Provenance}{Rattachement selon l’EDB : \emph{Coem. Callisti pars inferior}.}
\nf{Support}{Support non précisé par l’EDB. Fiche EDB : Antonio Enrico Felle, première version 2003.}
\ntext{\lnum{1}]s[- - -]\\
\lnum{2}[- - - dd(ominis) nn(ostris)] Val(e)n\\
\lnum{3}[te V et Valentini]ano\\
\lnum{4}Aug(ustis) [conss(ulibus)]}
\nf{Témoins}{EDB (EDB17958) : texte concordant aux conventions de notation près (indice de similarité 0,967).}
\nf{Traduction}{[…] sous le consulat de nos seigneurs Valens pour la 5e fois et Valentinien, Augustes.}
\nf{Langue}{5 abréviations développées par l’éditeur (dont dd(ominis), nn(ostris), Val(e)n, Aug(ustis), conss(ulibus)).}
\nf{Chronologie}{Noms consulaires : Valentinianus (lu en partie) et Valens (lu en partie).}
\nf{Bibliographie}{ICUR IV 9569 (p 533).}
\nf{Statut}{Lecture secondaire (EDCS) concordante avec l’EDB : lecture stabilisée entre deux bases, non collationnée sur l’ICVR ni sur le monument. Plus de la moitié des lettres sont restituées.}
\end{notice}

\begin{notice}{193}{376}{ICUR n.s. IV, 11769 · EDCS-37400858 · EDB878}
\label{not:193}
\nf{Datation}{376 — directe : nom consulaire lu}
\nf{Genre}{Épitaphe ; ensembles A, B, C ; 187 lettres transmises.}
\nf{Références}{ICUR IV 11769 = ILCV 4219.}
\nf{Provenance}{Rattachement selon l’EDB : \emph{Coem. anonymum ad viam Appiam}. Conservation selon l’EDB : \emph{Coem. anonymum ad viam Appiam}.}
\nf{Support}{Plaque de marbre. Inscription gravée. Signes décrits par l’EDB : chrisme et chrisme avec alpha et oméga. Fiche EDB : Maria Cristina Grisanzio, première version 2003.}
\ntext{\lnum{1}Decessit nomine Calleroe die tertiu(m) Idus Decembres\\
\lnum{2}consul⸢ibus⸣ dominis nostris Val⸢e⸣nte quinqu⸢e⸣(!) et Valentini\\
\lnum{3}ano iuniore quae vixit ann⸢o⸣s plus minus triginta\\
\lnum{4}et fecit cum marito ann⸢o⸣s plus minus quin\\
\lnum{5}decim et dies duodecim}
\nf{Témoins}{Export EDCS de 2020 : \emph{Decessit nomine Calleroe die tertiu(m) Idus Decembres / consul dominis nostris Valnte quinqu(!) et Valentini/ano iuniore quae vixit anns plus minus triginta / et fecit cum marito anns plus minus quin/decim et dies duodecim}. EDB (EDB878) : même texte ; les écarts relevés dans la liste de résultats tiennent à la notation (relecture de la fiche détaillée). Relecture de la fiche détaillée de l’EDB : Même texte ; consul<ibus=ATV> (EDCS) / consulatu (EDB). Chrisme avec alpha et oméga.}
\nf{Apparat}{⸢ibus⸣ : ATV gravé ; ⸢e⸣ : I gravé ; ⸢e⸣ : I gravé ; ⸢o⸣ : V gravé ; ⸢o⸣ : V gravé ; (!) : forme signalée par l’éditeur}
\nf{Traduction}{Est morte, nommée Calleroè, le 3e jour avant les ides de décembre, sous les consuls nos seigneurs Valens pour la 5e fois et Valentinien le Jeune ; elle vécut environ trente ans et a vécu avec son mari environ quinze ans et douze jours.}
\nf{Langue}{L’éditeur signale 5 graphies corrigées : ATV gravé pour ibus ; I gravé pour e (confusion e/i) ×2 ; V gravé pour o (confusion o/u) ×2. L’éditeur signale au moins une forme anormale (sic). 1 abréviation développée par l’éditeur (dont tertiu(m)).}
\nf{Chronologie}{Noms consulaires : Valentinianus (lu) et Valens (lu). Dates du jour : décès : die tertium Idus Decembres = 11 décembre 376 (lu). Âge déclaré : env. 30 a. 12 j. (lu).}
\nf{Bibliographie}{ICUR IV 11769 ; ILCV 4219.}
\nf{Statut}{Lecture secondaire (EDCS) concordante avec l’EDB : lecture stabilisée entre deux bases, non collationnée sur l’ICVR ni sur le monument.}
\end{notice}

\begin{notice}{194}{376}{ICUR n.s. IV, 12244 · EDCS-37600679 · EDB5083}
\label{not:194}
\nf{Datation}{376 — directe : nom consulaire lu}
\nf{Genre}{Épitaphe ; ensembles A, B, C ; 75 lettres transmises ; 23 \% restituées.}
\nf{Références}{ICUR IV 12244 = ILCV 4146.}
\nf{Provenance}{Rattachement selon l’EDB : \emph{Coem. anonymum ad viam Ardeatinam}.}
\nf{Support}{Plaque de marbre. Inscription gravée. Fiche EDB : Marida Pierno, première version 2003.}
\ntext{\lnum{1}Leo se vi⸢v⸣us fecit sibi et [- - -]\\
\lnum{2}Marcellinae quae vi[xit - - -]\\
\lnum{3}dec(essit) XVII Kal(endas) Dec(embres) Valente [- - - et Valentiniano - - -]}
\nf{Témoins}{Export EDCS de 2020 : \emph{Leo se vius fecit sibi et [3] / Marcellinae quae vi[xit 3] / dec(essit) XVII Kal(endas) Dec(embres) Valente [3 et Valentiniano 3]}. EDB (EDB5083) : texte concordant aux conventions de notation près (indice de similarité 0,938).}
\nf{Apparat}{⸢v⸣ : B gravé}
\nf{Traduction}{Leo, de son vivant, a fait (ce monument) pour lui-même et […] (pour) Marcellina, qui vécut […] ; elle mourut le 17e jour avant les calendes de décembre, sous (le consulat de) Valens […] et Valentinien […].}
\nf{Langue}{L’éditeur signale une graphie corrigée : B gravé pour v (bêtacisme).}
\nf{Chronologie}{Noms consulaires : Valentinianus (restitué) et Valens (lu). Ne sont pas des consuls : Marcellinae quae vi[xit (défunte). Dates du jour : décès : XVII Kalendas Decembres = 15 novembre 376 (lu).}
\nf{Bibliographie}{ICUR IV 12244 ; ILCV 4146.}
\nf{Statut}{Lecture secondaire (EDCS) concordante avec l’EDB : lecture stabilisée entre deux bases, non collationnée sur l’ICVR ni sur le monument.}
\end{notice}

\annee{377}

\begin{notice}{195}{377}{EDCS-09400838 · EDB33615}
\label{not:195}
\nf{Datation}{377 — directe : nom consulaire lu}
\nf{Genre}{Épitaphe ; ensembles A, B, C ; 32 lettres transmises ; 25 \% restituées.}
\nf{Références}{AE 1974, 85.}
\nf{Provenance}{Rattachement selon l’EDB : \emph{Coem. Calepodii}. Conservation selon l’EDB : \emph{Coem. Calepodii}.}
\nf{Support}{Plaque de marbre. 46,5 × 81 × 4 cm (h. × l. × ép.). Inscription gravée. Reproductions signalées par l’EDB : photographie PCAS (droits réservés ; non reproduites). Fiche EDB : Carlo Carletti, première version 2015.}
\ntext{\lnum{1}Grat(iano) et Merob(aude) co[ns(ulibus)]\\
\lnum{2}dep(ositio) Anastasi}
\nf{Témoins}{EDB (EDB33615) : texte concordant aux conventions de notation près (indice de similarité 0,97).}
\nf{Traduction}{Sous le consulat de Gratien et de Merobaudes, déposition d’Anastasius.}
\nf{Langue}{4 abréviations développées par l’éditeur (dont Grat(iano), Merob(aude), cons(ulibus), dep(ositio)).}
\nf{Chronologie}{Noms consulaires : Gratianus (lu) et Merobaudes (lu).}
\nf{Bibliographie}{AE 1974, 85. Bibliographie complémentaire selon l’EDB : Nestori 1972, 206, nota 111, fig. 94 ; AE, 1974, 85 ; Mazzoleni et alii 1999, 605, 611, 615 ; Solin 2003, 1283.}
\nf{Statut}{Lecture secondaire (EDCS) concordante avec l’EDB : lecture stabilisée entre deux bases, non collationnée sur l’ICVR ni sur le monument. L’EDCS signale une photographie (non consultée).}
\end{notice}

\begin{notice}{196}{377}{EDCS-02700284}
\label{not:196}
\nf{Datation}{377 — directe : nom consulaire lu}
\nf{Genre}{Épitaphe ; ensembles A, B, C ; 104 lettres transmises ; 48 \% restituées.}
\nf{Références}{AE 1994, 263.}
\nf{Provenance}{Rome (indication de l’EDCS) ; aucun rattachement cimétérial transmis.}
\nf{Support}{Non documenté dans les sources consultées.}
\ntext{\lnum{1}u]nius bonitatis\\
\lnum{2}[- - - dulcis]simo filio\\
\lnum{3}[qui vixit ann(os) - - -] m(enses) II d(ies) XII dep(ositus) in pace\\
\lnum{4}[- - - d(omino) n(ostro) Gratiano Aug(usto) IIII] et Merobaude cons(ulibus)}
\nf{Témoins}{Aucune fiche EDB appariée (pas de numéro ICUR ; recherche par le texte infructueuse).}
\nf{Traduction}{[…] d’une bonté unique […] à son fils très doux, qui vécut […] ans 2 mois 12 jours ; déposé en paix […] sous le consulat de notre seigneur Gratien Auguste pour la 4e fois et de Merobaudes.}
\nf{Langue}{8 abréviations développées par l’éditeur (dont ann(os), m(enses), d(ies), dep(ositus), d(omino), n(ostro)).}
\nf{Chronologie}{Noms consulaires : Gratianus (restitué) et Merobaudes (lu). Âge déclaré : 2 m. 12 j. (partiellement restitué).}
\nf{Commentaire}{La formule in pace est conservée.}
\nf{Bibliographie}{AE 1994, 263.}
\nf{Statut}{Lecture secondaire (EDCS) sans second témoin numérique.}
\end{notice}

\begin{notice}{197}{377}{ICUR n.s. I, 84 · EDCS-38600100 · EDB27457}
\label{not:197}
\nf{Datation}{377 — directe : nom consulaire lu}
\nf{Genre}{Épitaphe ; ensembles A, B, C ; 138 lettres transmises ; 29 \% restituées.}
\nf{Références}{ICUR I 84 (p 490).}
\nf{Provenance}{Rattachement selon l’EDB : \emph{Incertae originis} (le volume I des ICVR réunit les inscriptions dont la provenance est perdue). Conservation selon l’EDB : \emph{s. Cecilia in Trastevere}.}
\nf{Support}{Plaque de marbre. 44 × 31 × 4 cm (h. × l. × ép.). Inscription gravée. Fiche EDB : Filippo Piazzolla, première version 2013.}
\ntext{\lnum{1}Marcus e[t - - -]\\
\lnum{2}fecerunt [sibi se vivi]\\
\lnum{3}[- - -] Petrus vixit [- - - dec(essit)]\\
\lnum{4}X Kal(endas) Apr(iles) dd(ominis) nn(ostris)\\
\lnum{5}Gratiano III et Mer[obaude conss(ulibus)]\\
\lnum{6}Anastasa dec(essit) III[- - -]\\
\lnum{7}qu(a)e vix(it) ann(os) IIII d(ies) XX[- - -]\\
\lnum{8}Marcus qui vixit a[nn(os)}
\nf{Témoins}{EDB (EDB27457) : texte concordant aux conventions de notation près (indice de similarité 0,955).}
\nf{Traduction}{Marcus et […] se sont fait (ce tombeau) de leur vivant ; […] Petrus vécut […] ; il mourut le 10e jour avant les calendes d’avril, sous le consulat de nos seigneurs Gratien pour la 3e (sic) fois et Merobaudes ; Anastasia mourut le 3 […], qui vécut 4 ans 20 jours ; Marcus, qui vécut […] ans.}
\nf{Langue}{12 abréviations développées par l’éditeur (dont dec(essit), Kal(endas), Apr(iles), dd(ominis), nn(ostris), conss(ulibus)).}
\nf{Chronologie}{Noms consulaires : Gratianus (lu) et Merobaudes (lu en partie). Dates du jour : décès : X Kalendas Apriles = 23 mars 377 (lu). Âge déclaré : 4 a. (lu).}
\nf{Bibliographie}{ICUR I 84 (p 490). Bibliographie complémentaire selon l’EDB : Ferrua 1979b, 11, n. 84.}
\nf{Statut}{Lecture secondaire (EDCS) concordante avec l’EDB : lecture stabilisée entre deux bases, non collationnée sur l’ICVR ni sur le monument.}
\end{notice}

\begin{notice}{198}{377}{ICUR n.s. I, 1435 · EDCS-39000418 · EDB31199}
\label{not:198}
\nf{Datation}{377 — directe : lecture partielle d’un nom consulaire}
\nf{Genre}{Épitaphe ; ensembles A, B, C ; 92 lettres transmises ; 34 \% restituées.}
\nf{Références}{ICUR I 1435.}
\nf{Provenance}{Rattachement selon l’EDB : \emph{Incertae originis} (le volume I des ICVR réunit les inscriptions dont la provenance est perdue). Conservation selon l’EDB : \emph{Musei Vaticani: Lapidario Cristiano}.}
\nf{Support}{Plaque de marbre. 41 × 45 cm (h. × l.). Inscription gravée. Fiche EDB : Filippo Piazzolla, première version 2015.}
\ntext{\lnum{1}vi]xit an(nos) XXXVIII men[ses - - -]\\
\lnum{2}[- - -]I bene merenti in p[ace - - -]\\
\lnum{3}[- - - d(omino) n(ostro) Gratia]no Aug(usto) IIII et M⸢e⸣ro[baude]\\
\lnum{4}[co]nss(ulibus) ⸢XV⸣ Kal(endas) Iulias}
\nf{Témoins}{Export EDCS de 2020 : \emph{vi]xit an(nos) XXXVIII men[ses 3] / [3]I bene merenti in p[ace 3] / [3 d(omino) n(ostro) Gratia]no Aug(usto) IIII et Mro[baude] / [co]nss(ulibus) Kal(endas) Iulias}. EDB (EDB31199) : texte concordant aux conventions de notation près (indice de similarité 0,931).}
\nf{Apparat}{⸢e⸣ : Y gravé ; ⸢XV⸣ : VX gravé}
\nf{Traduction}{[…] vécut 38 ans […] mois ; […] bien méritant(e), en paix […] ; sous le consulat de notre seigneur Gratien Auguste pour la 4e fois et de Merobaudes, le 15e jour avant les calendes de juillet.}
\nf{Langue}{L’éditeur signale 2 graphies corrigées : Y gravé pour e ; VX gravé pour XV.}
\nf{Chronologie}{Noms consulaires : Gratianus (restitué) et Merobaudes (lu en partie). Dates du jour : événement indéterminé : XV Kalendas Iulias = 17 juin 377 (lu). Âge déclaré : 38 a. (partiellement restitué).}
\nf{Commentaire}{La formule in pace est attestée seulement par restitution.}
\nf{Bibliographie}{ICUR I 1435. Bibliographie complémentaire selon l’EDB : Lega 2000, 35.}
\nf{Statut}{Lecture secondaire (EDCS) concordante avec l’EDB : lecture stabilisée entre deux bases, non collationnée sur l’ICVR ni sur le monument. L’EDCS signale une photographie (non consultée).}
\end{notice}

\begin{notice}{199}{377}{ICUR n.s. I, 2094 · EDCS-39300316 · EDB32059}
\label{not:199}
\nf{Datation}{377 — directe : lecture partielle d’un nom consulaire}
\nf{Genre}{Épitaphe ; ensembles A, B, C ; 43 lettres transmises ; 40 \% restituées.}
\nf{Références}{ICUR I 2094.}
\nf{Provenance}{Rattachement selon l’EDB : \emph{Incertae originis} (le volume I des ICVR réunit les inscriptions dont la provenance est perdue). Conservation selon l’EDB : \emph{Musei Vaticani: Galleria Lapidaria}.}
\nf{Support}{Plaque de marbre. 20 × 32 cm (h. × l.). Inscription gravée. Fiche EDB : Filippo Piazzolla, première version 2015.}
\ntext{\lnum{1}[- - -]XXX deposita\\
\lnum{2}[- - - Gratiano III et Mero]baude consulibus}
\nf{Témoins}{EDB (EDB32059) : texte concordant aux conventions de notation près (indice de similarité 0,971).}
\nf{Traduction}{[…]30 (ans), déposée […] sous le consulat de Gratien pour la 3e (sic) fois et de Merobaudes.}
\nf{Chronologie}{Noms consulaires : Gratianus (restitué) et Merobaudes (lu en partie).}
\nf{Bibliographie}{ICUR I 2094. Bibliographie complémentaire selon l’EDB : Lega 2000.}
\nf{Statut}{Lecture secondaire (EDCS) concordante avec l’EDB : lecture stabilisée entre deux bases, non collationnée sur l’ICVR ni sur le monument.}
\end{notice}

\begin{notice}{200}{377}{ICUR n.s. I, 3188 · EDCS-40300334 · EDB33139}
\label{not:200}
\nf{Datation}{377 — directe : nom consulaire lu}
\nf{Genre}{Épitaphe ; ensembles A, B, C ; 89 lettres transmises.}
\nf{Références}{ICUR I 3188 = ILCV 4289.}
\nf{Provenance}{Rattachement selon l’EDB : \emph{Incertae originis} (le volume I des ICVR réunit les inscriptions dont la provenance est perdue).}
\nf{Support}{Support non précisé par l’EDB. Fiche EDB : Filippo Piazzolla, première version 2015.}
\ntext{\lnum{1}Gratiano IIII et Merobaude c(onsulibus)\\
\lnum{2}VII Id(us) Feb(ruarias) deposita est puella\\
\lnum{3}Sabina in pace qu(a)e vixit\\
\lnum{4}ann(os) XV in coniu(gio) fec(it) ann(os) II}
\nf{Témoins}{EDB (EDB33139) : texte concordant aux conventions de notation près (indice de similarité 0,961).}
\nf{Traduction}{Sous le consulat de Gratien pour la 4e fois et de Merobaudes, le 7e jour avant les ides de février, a été déposée la jeune fille Sabina, en paix, qui vécut 15 ans ; en mariage elle passa 2 ans.}
\nf{Langue}{8 abréviations développées par l’éditeur (dont c(onsulibus), Id(us), Feb(ruarias), qu(a)e, ann(os), coniu(gio)).}
\nf{Chronologie}{Noms consulaires : Gratianus (lu) et Merobaudes (lu). Dates du jour : événement indéterminé : VII Idus Februarias = 7 février 377 (lu). Âge déclaré : 15 a. (lu).}
\nf{Commentaire}{La formule in pace est conservée.}
\nf{Bibliographie}{ICUR I 3188 ; ILCV 4289. Bibliographie complémentaire selon l’EDB : ILCV, 4289.}
\nf{Statut}{Lecture secondaire (EDCS) concordante avec l’EDB : lecture stabilisée entre deux bases, non collationnée sur l’ICVR ni sur le monument.}
\end{notice}

\begin{notice}{201}{377}{ICUR n.s. I, 3189 · EDCS-40300335 · EDB33140}
\label{not:201}
\nf{Datation}{377 — directe : nom consulaire lu}
\nf{Genre}{Épitaphe ; ensembles A, B, C ; 79 lettres transmises ; 28 \% restituées.}
\nf{Références}{ICUR I 3189 = ILCV 1548.}
\nf{Provenance}{Rattachement selon l’EDB : \emph{Incertae originis} (le volume I des ICVR réunit les inscriptions dont la provenance est perdue).}
\nf{Support}{Support non précisé par l’EDB. Fiche EDB : Filippo Piazzolla, première version 2015.}
\ntext{\lnum{1}[- - -]C solum nostrum [- - -]\\
\lnum{2}[- - -]L C(h)ristianae uncta [- - -]\\
\lnum{3}[- - - qu]ae vixit ann(os) p(lus) m(inus) XXV[- - -]\\
\lnum{4}[Gratia]no Au⸢g⸣(usto) IIII M(er)o⸢ba⸣[aude cons(ulibus) - - -]}
\nf{Témoins}{Export EDCS de 2020 : \emph{[3]C solum nostrum [3] / [3]L C(h)ristianae uncta [3] / [3 qu]ae vixit ann(os) p(lus) m(inus) XXV[3] / [Gratia]no Au(usto) IIII M(er)o[aude cons(ulibus) 3]}. EDB (EDB33140 ; écarts à examiner) : \emph{[---]c solum nostrum [---] / [---]l cristianae uncta [---] / [--- qu]ae vixit ann(os) p(lus) m(inus) XXV [---] / [--- Gratia]no Au⌜g⌝(usto) IIII M⟨er⟩o⌜ba⌝[ude ---] [texte tronqué dans la liste de résultats]}.}
\nf{Apparat}{⸢g⸣ : C gravé ; ⸢ba⸣ : PR gravé}
\nf{Traduction}{[…] notre sol […] ; […] ointe chrétienne (?) […] ; […] qui vécut environ 25 ans […] ; sous le consulat de Gratien Auguste pour la 4e fois et de Merobaudes […].}
\nf{Langue}{L’éditeur signale 2 graphies corrigées : C gravé pour g (c pour g) ; PR gravé pour ba.}
\nf{Chronologie}{Noms consulaires : Gratianus (restitué) et Merobaudes (lu).}
\nf{Commentaire}{Vocabulaire religieux : identité chrétienne (christianus) — Christianae (lu).}
\nf{Bibliographie}{ICUR I 3189 ; ILCV 1548. Bibliographie complémentaire selon l’EDB : ILCV, 1548.}
\nf{Statut}{Lecture secondaire (EDCS) ; écarts avec l’EDB à arbitrer sur l’ICVR.}
\end{notice}

\begin{notice}{202}{377}{ICUR n.s. II, 4815 · EDCS-36101297}
\label{not:202}
\nf{Datation}{377 — directe : nom consulaire lu}
\nf{Genre}{Épitaphe ; ensembles A, B, C ; 124 lettres transmises.}
\nf{Références}{ICUR II 4815 = ILCV 1269 (em) = ICaRoma 130 = Carletti-2008, 00058.}
\nf{Provenance}{Rome (indication de l’EDCS) ; aucun rattachement cimétérial transmis.}
\nf{Support}{Non documenté dans les sources consultées.}
\ntext{\lnum{1}Cinnamius Opas lector tituli Fasciol(a)e amicus pauperum\\
\lnum{2}qui vixit ann(os) XLVI mens(es) VII d(ies) VIIII deposit(us) in pace X Kal(endas) Mart(ias)\\
\lnum{3}Gratiano IIII et Merobaude conss(ulibus)}
\nf{Témoins}{Aucune fiche EDB appariée (numéro ICUR absent de la liste publique de l’EDB).}
\nf{Traduction}{Cinnamius Opas, lecteur du titre de Fasciola, ami des pauvres, qui vécut 46 ans 7 mois 9 jours ; déposé en paix le 10e jour avant les calendes de mars, sous le consulat de Gratien pour la 4e fois et de Merobaudes.}
\nf{Langue}{8 abréviations développées par l’éditeur (dont Fasciol(a)e, ann(os), mens(es), d(ies), deposit(us), Kal(endas)).}
\nf{Chronologie}{Noms consulaires : Gratianus (lu) et Merobaudes (lu). Dates du jour : déposition : X Kalendas Martias = 20 février 377 (lu). Âge déclaré : 46 a. 7 m. 9 j. (lu).}
\nf{Commentaire}{La formule in pace est conservée. Fonction ou métier mentionné : lector. Commentaire développé : dossier p.~\pageref{dos:202}.}
\nf{Bibliographie}{ICUR II 4815 ; ILCV 1269 (em) ; ICaRoma 130 ; Carletti-2008, 00058.}
\nf{Statut}{Lecture secondaire (EDCS) sans second témoin numérique. L’EDCS signale une photographie (non consultée).}
\end{notice}

\begin{notice}{203}{377}{ICUR n.s. II, 5790 · EDCS-38302492 · EDB22177}
\label{not:203}
\nf{Datation}{377 — directe : consulat restitué}
\nf{Genre}{Épitaphe ; ensembles A, B, C ; 135 lettres transmises ; 42 \% restituées.}
\nf{Références}{ICUR II 5790 = ILCV +1002.}
\nf{Provenance}{Rattachement selon l’EDB : \emph{Area Lucinae}. Conservation selon l’EDB : \emph{Museo Archeologico Nazionale}.}
\nf{Support}{Plaque de marbre. 43 × 64 × 4 cm (h. × l. × ép.). Inscription gravée. Fiche EDB : Anita Rocco, première version 2013.}
\ntext{\lnum{1}[- - - Re]verendo patri IV[- - -]\\
\lnum{2}[- - -] qui vixit annos LXVII[- - -]\\
\lnum{3}[- - - depos]itu[s in] pace die III Non(as) No\\
\lnum{4}[vemb(res) d(omino) n(ostro) Gratiano Aug(usto) IIII et Merobau]de v(iro) c(larissimo) cons(ulibus)\\
\lnum{5}[- - -] Leoparda\\
\lnum{6}[- - -]ius quiesc(it) in p(ace)}
\nf{Témoins}{EDB (EDB22177) : texte concordant aux conventions de notation près (indice de similarité 0,962).}
\nf{Traduction}{[…] au révérend père […] qui vécut 67 ans […] ; déposé en paix le 3e jour avant les nones de novembre, sous le consulat de notre seigneur Gratien Auguste pour la 4e fois et de Merobaudes, clarissime ; […] Leoparda ; […] repose en paix.}
\nf{Langue}{10 abréviations développées par l’éditeur (dont Non(as), vemb(res), d(omino), n(ostro), Aug(usto), v(iro)).}
\nf{Chronologie}{Noms consulaires : Gratianus (restitué) et Merobaudes (restitué). Dates du jour : déposition : die III Nonas Novembres = 3 novembre 377 (partiellement restitué). Âge déclaré : 3 j. (partiellement restitué).}
\nf{Commentaire}{La formule in pace est transmise sous forme abrégée, sans restitution de lacune ; son développement reste éditorial. Vocabulaire religieux : repos et sommeil de la mort — quiescit (lu).}
\nf{Bibliographie}{ICUR II 5790 ; ILCV +1002.}
\nf{Statut}{Lecture secondaire (EDCS) concordante avec l’EDB : lecture stabilisée entre deux bases, non collationnée sur l’ICVR ni sur le monument.}
\end{notice}

\begin{notice}{204}{377}{ICUR n.s. III, 8423 · EDCS-35802010 · EDB25698}
\label{not:204}
\nf{Datation}{377 — directe : nom consulaire lu}
\nf{Genre}{Épitaphe ; ensembles A, B, C ; 123 lettres transmises ; 7 \% restituées.}
\nf{Références}{ICUR III 8423 = ILCV 3598 (em).}
\nf{Provenance}{Rattachement selon l’EDB : \emph{Incertum locum in coem. Domitillae}. Conservation selon l’EDB : \emph{coem. Domitillae, aula supra coemeterium}.}
\nf{Support}{Plaque de marbre. Inscription gravée. Reproductions signalées par l’EDB : photographie PCAS ×2 (droits réservés ; non reproduites). Fiche EDB : Carlo Carletti / Antonio Enrico Felle, première version 2013.}
\ntext{\lnum{1}Fortunius se viv⸢us⸣ fecit\\
\lnum{2}d(omino) n(ostro) Gratiano \{ano\} Aug(usto) IIII et Merobaude cons[s(ulibus)]\\
\lnum{3}deposita est Caelia rari exempli femina quod vides h[ic]\\
\lnum{4}quiescit\\
\lnum{5}Porfurius}
\nf{Témoins}{Export EDCS de 2020 : \emph{Fortunius se viv fecit / d(omino) n(ostro) Gratiano \{ano\} Aug(usto) IIII et Merobaude cons[s(ulibus)] / deposita est Caelia rari exempli femina quod vides h[ic] / quiescit / Porfurius}. EDB (EDB25698 ; écarts à examiner) : \emph{⟨a:⟩ d(omino) n(ostro) Gratiano \{ano\} Aug(usto) IIII et Merobaude cons(ulibu)[s] / deposita est Caelia rari exempli femina quod vides h[ic] / quiescit / / ⟨b:⟩ Fortunius se vivom fecit / /⟨c:⟩ Porfurius}.}
\nf{Apparat}{⸢us⸣ : OM gravé ; \{ano\} : lettres superflues}
\nf{Traduction}{Fortunius, de son vivant, a fait (ce monument) ; sous le consulat de notre seigneur Gratien Auguste pour la 4e fois et de Merobaudes a été déposée Caelia, femme d’un rare exemple ; ce que tu vois ici, elle y repose ; Porphyrius.}
\nf{Langue}{L’éditeur signale une graphie corrigée : OM gravé pour us. Des lettres gravées sont jugées superflues par l’éditeur (entre accolades). 4 abréviations développées par l’éditeur (dont d(omino), n(ostro), Aug(usto), conss(ulibus)).}
\nf{Chronologie}{Noms consulaires : Gratianus (lu) et Merobaudes (lu).}
\nf{Commentaire}{Vocabulaire religieux : repos et sommeil de la mort — quiescit (lu) ; allusion interprétative (niveau 5) — Caelia (lu).}
\nf{Bibliographie}{ICUR III 8423 ; ILCV 3598 (em). Bibliographie complémentaire selon l’EDB : ILCV, 3598.}
\nf{Statut}{Lecture secondaire (EDCS) ; écarts avec l’EDB à arbitrer sur l’ICVR. L’EDCS signale une photographie (non consultée).}
\end{notice}

\begin{notice}{205}{377}{ICUR n.s. IV, 11107 · EDCS-37200389 · EDB39569}
\label{not:205}
\nf{Datation}{377 — directe : nom consulaire lu}
\nf{Genre}{Épitaphe ; ensembles A, B, C ; 70 lettres transmises.}
\nf{Références}{ICUR IV 11107 = ILCV 2969 (em).}
\nf{Provenance}{Rattachement selon l’EDB : \emph{extra portam Capenam}. Conservation selon l’EDB : \emph{Musei Vaticani: Galleria Lapidaria}.}
\nf{Support}{Plaque de marbre. 25 × 58 cm (h. × l.). Inscription gravée. Fiche EDB : Marida Pierno, première version 2015.}
\ntext{\lnum{1}D(omino) n(ostro) Gratiano Aug(usto) IIII et\\
\lnum{2}M⸢e⸣r⸢o⸣ba(u)de conss(ulibus) XVII Kal(endas) Ma()\\
\lnum{3}depositus in pace Faustinus\\
\lnum{4}m(ensium) VIII}
\nf{Témoins}{Export EDCS de 2020 : \emph{D(omino) n(ostro) Gratiano Aug(usto) IIII et / Mrba(u)de conss(ulibus) XVII Kal(endas) Ma() / depositus in pace Faustinus / m(ensium) VIII}. EDB (EDB39569) : texte concordant aux conventions de notation près (indice de similarité 0,955).}
\nf{Apparat}{⸢e⸣ : A gravé ; ⸢o⸣ : A gravé}
\nf{Traduction}{Sous le consulat de notre seigneur Gratien Auguste pour la 4e fois et de Merobaudes, le 17e jour avant les calendes de mai (ou mars), déposé en paix Faustinus, de 8 mois.}
\nf{Langue}{L’éditeur signale 2 graphies corrigées : A gravé pour e ; A gravé pour o.}
\nf{Chronologie}{Noms consulaires : Gratianus (lu) et Merobaudes (lu).}
\nf{Commentaire}{La formule in pace est conservée.}
\nf{Bibliographie}{ICUR IV 11107 ; ILCV 2969 (em). Bibliographie complémentaire selon l’EDB : Lega 2000.}
\nf{Statut}{Lecture secondaire (EDCS) concordante avec l’EDB : lecture stabilisée entre deux bases, non collationnée sur l’ICVR ni sur le monument.}
\end{notice}

\annee{370}

\begin{notice}{206}{370}{ICUR n.s. IV, 11770 · EDCS-35700946 · EDB880}
\label{not:206}
\nf{Datation}{370 — directe : nom consulaire lu}
\nf{Genre}{Épitaphe ; ensembles A, B, C ; 53 lettres transmises.}
\nf{Références}{ICUR IV 11770 (p 535) = ILCV 3996c.}
\nf{Provenance}{Rattachement selon l’EDB : \emph{Coem. anonymum ad viam Appiam}. Conservation selon l’EDB : \emph{Coem. anonymum ad viam Appiam}.}
\nf{Support}{Paroi de enduit. Inscription peinte. Fiche EDB : Maria Cristina Grisanzio, première version 2013.}
\ntext{\lnum{1}Elpidius\\
\lnum{2}d(epositus) XVIII Kal(endas) Oct(obres)\\
\lnum{3}Valentiniano III\\
\lnum{4}et Valente Augg(ustis)\\
\lnum{5}conss(ulibus)}
\nf{Témoins}{EDB (EDB880) : texte concordant aux conventions de notation près (indice de similarité 1).}
\nf{Traduction}{Elpidius, déposé le 18e jour avant les calendes d’octobre, sous le consulat de Valentinien pour la 3e fois et de Valens, Augustes.}
\nf{Langue}{5 abréviations développées par l’éditeur (dont d(epositus), Kal(endas), Oct(obres), Augg(ustis), conss(ulibus)).}
\nf{Chronologie}{Noms consulaires : Valentinianus (lu) et Valens (lu). La formule \emph{Valentiniano III et Valente Augg.} ne convient qu'à 370 (Valentinien Ier III et Valens III). L'EDB (EDB880) date de 377 : il s'agit vraisemblablement d'une erreur de saisie de la base, l'inscription étant peinte. L'année 370 est maintenue. Dates du jour : déposition : XVIII Kalendas Octobres = 14 septembre 370 (lu).}
\nf{Bibliographie}{ICUR IV 11770 (p 535) ; ILCV 3996c.}
\nf{Statut}{Lecture secondaire (EDCS) concordante avec l’EDB : lecture stabilisée entre deux bases, non collationnée sur l’ICVR ni sur le monument.}
\end{notice}

\annee{378}

\begin{notice}{207}{378}{ICUR n.s. IV, 9570 · EDCS-36101096 · EDB17959}
\label{not:207}
\nf{Datation}{378 — directe : nom consulaire lu}
\nf{Genre}{Épitaphe ; ensembles A, B, C ; 93 lettres transmises ; 13 \% restituées.}
\nf{Références}{ICUR IV 9570 = ILCV 4146a.}
\nf{Provenance}{Rattachement selon l’EDB : \emph{Coem. Callisti pars inferior}. Conservation selon l’EDB : \emph{Musei Vaticani: Lapidario Cristiano}.}
\nf{Support}{Plaque de marbre. Inscription gravée. Fiche EDB : Antonio Enrico Felle, première version 2003.}
\ntext{\lnum{1}Titiana qu(a)e vixit annos [- - -]\\
\lnum{2}Agape se ⸢v⸣i⸢v⸣a fecit\\
\lnum{3}dd(ominis) nn(ostris) Fl(aviis) Valente VI et Valentiniano A(u)g(ustis) I[I conss(ulibus)]\\
\lnum{4}d(e)p(osita) Agape in pace}
\nf{Témoins}{Export EDCS de 2020 : \emph{Titiana qu(a)e vixit annos [3] / Agape se ia fecit / dd(ominis) nn(ostris) Fl(aviis) Valente VI et Valentiniano A(u)g(ustis) I[I conss(ulibus)] / d(e)p(osita) Agape in pace}. EDB (EDB17959) : texte concordant aux conventions de notation près (indice de similarité 0,963).}
\nf{Apparat}{⸢v⸣ : B gravé ; ⸢v⸣ : B gravé}
\nf{Traduction}{Titiana, qui vécut […] ans ; Agapè, de son vivant, a fait (ce monument) ; sous le consulat de nos seigneurs les Flavii Valens pour la 6e fois et Valentinien pour la 2e fois, Augustes, Agapè a été déposée en paix.}
\nf{Langue}{L’éditeur signale 2 graphies corrigées : B gravé pour v (bêtacisme) ×2.}
\nf{Chronologie}{Noms consulaires : Valentinianus (lu) et Valens (lu).}
\nf{Commentaire}{La formule in pace est conservée.}
\nf{Bibliographie}{ICUR IV 9570 ; ILCV 4146a.}
\nf{Statut}{Lecture secondaire (EDCS) concordante avec l’EDB : lecture stabilisée entre deux bases, non collationnée sur l’ICVR ni sur le monument. L’EDCS signale une photographie (non consultée).}
\end{notice}

\begin{notice}{208}{378}{ICUR n.s. IV, 11771 · EDCS-37400859 · EDB1342}
\label{not:208}
\nf{Datation}{378 — directe : nom consulaire lu}
\nf{Genre}{Épitaphe ; ensembles A, B, C ; 110 lettres transmises ; 23 \% restituées.}
\nf{Références}{ICUR IV 11771 = ILCV +2945a.}
\nf{Provenance}{Rattachement selon l’EDB : \emph{Coem. anonymum ad viam Appiam}. Conservation selon l’EDB : \emph{Coem. anonymum ad viam Appiam}.}
\nf{Support}{Plaque de marbre. Inscription gravée. Reproductions signalées par l’EDB : photographie PCAS (droits réservés ; non reproduites). Fiche EDB : Maria Cristina Grisanzio, première version 2003.}
\ntext{\lnum{1}Germa[na qua]e ⸢v⸣(i)xit mense[s - - -]\\
\lnum{2}deposita d[ie] pridie Nonas Septe\\
\lnum{3}mbres dd(ominis) nn(ostris) Valente VI et\\
\lnum{4}[Flavio] Valentiniano iterum\\
\lnum{5}[Augustis con]sulibus}
\nf{Témoins}{Export EDCS de 2020 : \emph{Germa[na qua]e (i)xit mense[s 3] / deposita d[ie] pridie Nonas Septe/mbres dd(ominis) nn(ostris) Valente VI et / [Flavio] Valentiniano iterum / [Augustis con]sulibus}. EDB (EDB1342) : texte concordant aux conventions de notation près (indice de similarité 0,985).}
\nf{Apparat}{⸢v⸣ : B gravé}
\nf{Traduction}{Germana, qui vécut […] mois ; déposée la veille des nones de septembre, sous le consulat de nos seigneurs Valens pour la 6e fois et Flavius Valentinien pour la 2e fois, Augustes.}
\nf{Langue}{L’éditeur signale une graphie corrigée : B gravé pour v (bêtacisme).}
\nf{Chronologie}{Noms consulaires : Valentinianus (lu) et Valens (lu). Dates du jour : déposition : die pridie Nonas Septembres (partiellement restitué).}
\nf{Bibliographie}{ICUR IV 11771 ; ILCV +2945a. Bibliographie complémentaire selon l’EDB : CLRE, 648.}
\nf{Statut}{Lecture secondaire (EDCS) concordante avec l’EDB : lecture stabilisée entre deux bases, non collationnée sur l’ICVR ni sur le monument. L’EDCS signale une photographie (non consultée).}
\end{notice}

\begin{notice}{209}{378}{ICUR n.s. IV, 11772 · EDCS-37400860 · EDB5268}
\label{not:209}
\nf{Datation}{378 — directe : nom consulaire lu}
\nf{Genre}{Acte d’achat ou de réservation ; ensembles A, B ; 65 lettres transmises.}
\nf{Références}{ICUR IV 11772 = ILCV 3754a (em).}
\nf{Provenance}{Rattachement selon l’EDB : \emph{Coem. anonymum ad viam Appiam}. Conservation selon l’EDB : \emph{Coem. anonymum ad viam Appiam}.}
\nf{Support}{Plaque de marbre. Inscription gravée. Signes décrits par l’EDB : chrisme et chrisme avec alpha et oméga. Reproductions signalées par l’EDB : photographie PCAS ×2 (droits réservés ; non reproduites). Fiche EDB : Maria Cristina Grisanzio, première version 2003.}
\ntext{\lnum{1}Hic locus\\
\lnum{2}est Victo\\
\lnum{3}rini quem\\
\lnum{4}emit a Mar\\
\lnum{5}tinian⸢o⸣\\
\lnum{6}fossore\\
\lnum{7}Valente VI\\
\lnum{8}et Valentin(iano !)}
\nf{Témoins}{Export EDCS de 2020 : \emph{Hic locus / est Victo/rini quem / emit a Mar/tinian / fossore / Valente VI / et Valentin(iano !)}. EDB (EDB5268 ; écarts à examiner) : \emph{((monogramma Christi cum alpha et omega)) / hic locus / est Victo/rini quem / emit / a Mar/tinianu / fossore ((folium)) / Valente VI / et Valentin[iano bis]}.}
\nf{Apparat}{⸢o⸣ : V gravé}
\nf{Traduction}{Cet emplacement est celui de Victorinus, qui l’a acheté au fossoyeur Martinianus, sous (le consulat de) Valens pour la 6e fois et Valentinien.}
\nf{Langue}{L’éditeur signale une graphie corrigée : V gravé pour o (confusion o/u).}
\nf{Chronologie}{Noms consulaires : Valentinianus (lu) et Valens (lu).}
\nf{Commentaire}{Fonction ou métier mentionné : fossore. Commentaire développé : dossier p.~\pageref{dos:209}.}
\nf{Bibliographie}{ICUR IV 11772 ; ILCV 3754a (em).}
\nf{Statut}{Lecture secondaire (EDCS) ; écarts avec l’EDB à arbitrer sur l’ICVR. L’EDCS signale une photographie (non consultée).}
\end{notice}

\annee{379}

\begin{notice}{210}{379}{ICUR n.s. I, 3142 · EDCS-39900598 · EDB33094}
\label{not:210}
\nf{Datation}{379 — directe : nom consulaire lu}
\nf{Genre}{Épitaphe ; ensembles A, B, C ; 86 lettres transmises.}
\nf{Références}{CIL V *429 = ICUR I 3142 = ILCV 1479.}
\nf{Provenance}{Rattachement selon l’EDB : \emph{Incertae originis} (le volume I des ICVR réunit les inscriptions dont la provenance est perdue). Conservation selon l’EDB : \emph{Verona   ; Museo Lapidario Maffeiano}.}
\nf{Support}{Plaque de marbre. 17 × 54 × 1,8 cm (h. × l. × ép.). Inscription gravée. Fiche EDB : Filippo Piazzolla, première version 2015.}
\ntext{\lnum{1}Quiminati neofyto bene me\\
\lnum{2}renti in pace qui vixit ann(os) p(lus) m(inus) V\\
\lnum{3}depositus VIII Kal(endas) Dec(embres) Ausonio\\
\lnum{4}et Olybrio co(n)ss(ulibus)}
\nf{Témoins}{EDB (EDB33094) : texte concordant aux conventions de notation près (indice de similarité 0,995).}
\nf{Traduction}{À Quiminatus (?), néophyte, bien méritant, en paix, qui vécut environ 5 ans ; déposé le 8e jour avant les calendes de décembre, sous le consulat d’Ausone et d’Olybrius.}
\nf{Langue}{7 abréviations développées par l’éditeur (dont ann(os), p(lus), m(inus), Kal(endas), Dec(embres), co(n)ss).}
\nf{Chronologie}{Noms consulaires : Ausonius (lu) et Olybrius (lu). Dates du jour : déposition : VIII Kalendas Decembres = 24 novembre 379 (lu). Âge déclaré : env. 5 a. (lu).}
\nf{Commentaire}{La formule in pace est conservée. Vocabulaire religieux : initiation et foi — neofyto (lu).}
\nf{Bibliographie}{CIL V *429 ; ICUR I 3142 ; ILCV 1479. Bibliographie complémentaire selon l’EDB : CIL X, 429\textbackslash{}* ; Ferrua 1979b, 89 ; ILCV, 1479.}
\nf{Statut}{Lecture secondaire (EDCS) concordante avec l’EDB : lecture stabilisée entre deux bases, non collationnée sur l’ICVR ni sur le monument.}
\end{notice}

\begin{notice}{211}{379}{ICUR n.s. I, 3191 · EDCS-40300336 · EDB33143}
\label{not:211}
\nf{Datation}{379 — directe : nom consulaire lu}
\nf{Genre}{Épitaphe ; ensembles A, B, C ; 133 lettres transmises ; 4 \% restituées.}
\nf{Références}{ICUR I 3191 = ILCV 2772.}
\nf{Provenance}{Rattachement selon l’EDB : \emph{Incertae originis} (le volume I des ICVR réunit les inscriptions dont la provenance est perdue).}
\nf{Support}{Support non précisé par l’EDB. Fiche EDB : Filippo Piazzolla, première version 2015.}
\ntext{\lnum{1}In pace\\
\lnum{2}sancto filio ⸢G⸣erontio inno\\
\lnum{3}centi qui vixit anno uno\\
\lnum{4}[m]ensibus octo et di(e)bus octo abso\\
\lnum{5}[lutus] de corpore Idibus Iul(iis) depos(ita) XVII K(alendas) Aug(ustas)\\
\lnum{6}Ausonio et Olybrio conss(ulibus)}
\nf{Témoins}{Export EDCS de 2020 : \emph{In pace / sancto filio erontio inno/centi qui vixit anno uno / [m]ensibus octo et di(e)bus octo abso/[lutus] de corpore Idibus Iul(iis) depos(ita) XVII K(alendas) Aug(ustas) / Ausonio et Olybrio conss(ulibus)}. EDB (EDB33143) : texte concordant aux conventions de notation près (indice de similarité 0,967).}
\nf{Apparat}{⸢G⸣ : C gravé}
\nf{Traduction}{En paix. À leur saint fils Gerontius, innocent, qui vécut un an huit mois et huit jours ; délivré de son corps aux ides de juillet ; déposé le 17e jour avant les calendes d’août, sous le consulat d’Ausone et d’Olybrius.}
\nf{Langue}{L’éditeur signale une graphie corrigée : C gravé pour G.}
\nf{Chronologie}{Noms consulaires : Ausonius (lu) et Olybrius (lu). Dates du jour : décès : Idibus Iuliis = 15 juillet 379 (lu) ; déposition : XVII Kalendas Augustas = 16 juillet 379 (lu). Âge déclaré : 1 a. 8 m. 8 j. (partiellement restitué).}
\nf{Commentaire}{La formule in pace est conservée. Vocabulaire religieux : épithète du défunt (non doctrinal) — sancto (lu).}
\nf{Bibliographie}{ICUR I 3191 ; ILCV 2772. Bibliographie complémentaire selon l’EDB : ILCV, 2772.}
\nf{Statut}{Lecture secondaire (EDCS) concordante avec l’EDB : lecture stabilisée entre deux bases, non collationnée sur l’ICVR ni sur le monument.}
\end{notice}

\begin{notice}{212}{379}{ICUR n.s. II, 4948 · EDCS-37700147}
\label{not:212}
\nf{Datation}{379 — directe : nom consulaire lu}
\nf{Genre}{Épitaphe ; ensembles A, B, C ; 15 lettres transmises.}
\nf{Références}{ICUR II 4948.}
\nf{Provenance}{Rome (indication de l’EDCS) ; aucun rattachement cimétérial transmis.}
\nf{Support}{Non documenté dans les sources consultées.}
\ntext{\lnum{1}]io et Olybrio cons(ulibus)}
\nf{Témoins}{Aucune fiche EDB appariée (numéro ICUR absent de la liste publique de l’EDB).}
\nf{Traduction}{[…] sous le consulat d’Ausone et d’Olybrius.}
\nf{Langue}{1 abréviation développée par l’éditeur (dont cons(ulibus)).}
\nf{Chronologie}{Noms consulaires : Olybrius (lu).}
\nf{Bibliographie}{ICUR II 4948.}
\nf{Statut}{Lecture secondaire (EDCS) sans second témoin numérique.}
\end{notice}

\annee{380}

\begin{notice}{213}{380}{ICUR n.s. II, 4166 · EDCS-19900369 · EDB19166}
\label{not:213}
\nf{Datation}{380 — directe : nom consulaire lu}
\nf{Genre}{Épitaphe ; ensembles A, B, C ; 85 lettres transmises ; 20 \% restituées.}
\nf{Références}{CIL VI 31972 (p 4798) = ICUR II 4166 = ILCV +139.}
\nf{Provenance}{Rattachement selon l’EDB : \emph{Basilica S.Petri apostoli}.}
\nf{Support}{Support non précisé par l’EDB. Signes décrits par l’EDB : chrisme et croix monogrammatique. Fiche EDB : Antonio Enrico Felle, première version 2003.}
\ntext{\lnum{1}H(ic) q(uiescit)\\
\lnum{2}Eustolius v(ir) c(larissimus) cum filia The[odora(?)]\\
\lnum{3}deposita X Kal(endas) Aug(ustas) depo[sit - - -]\\
\lnum{4}dd(ominis) nn(ostris) Gratiano V(?)\\
\lnum{5}et [Theodosio] Augg(ustis) co(n)ss(ulibus)\\
\lnum{6}i(n) p(ace)}
\nf{Témoins}{EDB (EDB19166) : texte concordant aux conventions de notation près (indice de similarité 0,951).}
\nf{Traduction}{Ici repose Eustolius, clarissime, avec sa fille Theodora ; (elle a été) déposée le 10e jour avant les calendes d’août ; déposé(e) […] ; sous le consulat de nos seigneurs Gratien pour la 5e fois et Théodose, Augustes ; en paix.}
\nf{Langue}{15 abréviations développées par l’éditeur (dont H(ic), q(uiescit), v(ir), c(larissimus), Theodora(?), Kal(endas)).}
\nf{Chronologie}{Noms consulaires : Gratianus (lu) et Theodosius (restitué). Dates du jour : déposition : X Kalendas Augustas = 23 juillet 380 (lu).}
\nf{Commentaire}{La formule in pace est transmise sous forme abrégée, sans restitution de lacune ; son développement reste éditorial. Vocabulaire religieux : repos et sommeil de la mort — quiescit (abréviation développée) ; anthroponyme (niveau 6) — theodora (anthroponyme). Rang sénatorial d’une personne commémorée (hic quiescit eustolius vir clarissimus cum ).}
\nf{Bibliographie}{CIL VI 31972 (p 4798) ; ICUR II 4166 ; ILCV +139. Bibliographie complémentaire selon l’EDB : CIL VI, 31972.}
\nf{Statut}{Lecture secondaire (EDCS) concordante avec l’EDB : lecture stabilisée entre deux bases, non collationnée sur l’ICVR ni sur le monument. L’EDCS signale une photographie (non consultée).}
\end{notice}

\begin{notice}{214}{380}{ICUR n.s. I, 3194 · EDCS-35200283 · EDB33145}
\label{not:214}
\nf{Datation}{380 — formule transmise par l’EDB (lue)}
\nf{Genre}{Épitaphe ; ensembles A, B, C ; 129 lettres transmises.}
\nf{Références}{CLE 1619 = ICUR I 3194 = ILCV 1464.}
\nf{Provenance}{Rattachement selon l’EDB : \emph{Incertae originis} (le volume I des ICVR réunit les inscriptions dont la provenance est perdue). Conservation selon l’EDB : \emph{villa Albani}.}
\nf{Support}{Plaque de marbre. 69 × 128 cm (h. × l.). Inscription gravée. Reproductions signalées par l’EDB : reproduction, Perret 1851 (droits réservés ; non reproduites). Fiche EDB : Filippo Piazzolla, première version 2016.}
\ntext{\lnum{1}Hic quiescit ancilla dei qu(a)e de\\
\lnum{2}sua omnia possedit domum ista\\
\lnum{3}qu⸢a⸣m amic(a)e deflen(t) solaciumq(ue) requirunt\\
\lnum{4}pro hunc unum ora subolem qu⸢a⸣m superis\\
\lnum{5}titem re(li)quisti}
\nf{Témoins}{Export EDCS de 2020 : \emph{Hic quiescit ancilla dei qu(a)e de / sua omnia possedit domum ista / qum amic(a)e deflen(t) solaciumq(ue) requirunt / pro hunc unum ora subolem qum superis/titem re(li)quisti}. EDB (EDB33145 ; divergence forte) : \emph{hic quescit ancilla dei ⌜q⌝ue de / sua omnia possedit domum ista / quem amice defien solaciumq(ue) requirun̂t / pro hunc unum ora subolem quem superis / titem requisti eterna requiem felicita/⟨ti⟩s cuasa man(e)bis IIIIX k⌜a⌝lendas octobris / Cucurbitinus et Abu⌜nd⌝antius hic ⌜s⌝imul quiescit / dd(ominis) nn(ostris) Gratiano V et Teodosio aaugg(ustis) ((aves coronam tenentes)) [texte tronqué dans la liste de résultats]}. Traduction des lignes transmises par l'EDB seule : « … tu demeureras dans le repos éternel de la félicité (?) ; le … jour avant les calendes d'octobre. Cucurbitinus et Abundantius reposent ici ensemble. Sous le consulat de nos seigneurs Gratien, pour la cinquième fois, et Théodose, Augustes. »}
\nf{Apparat}{⸢a⸣ : E gravé ; ⸢a⸣ : E gravé}
\nf{Traduction}{Ici repose une servante de Dieu qui, de ses biens, posséda toute cette maison, que ses amies pleurent et dont elles cherchent le réconfort ; prie pour ce seul (enfant), la descendance que tu as laissée survivante.}
\nf{Langue}{L’éditeur signale 2 graphies corrigées : E gravé pour a (e pour a) ×2.}
\nf{Chronologie}{Le texte EDCS s'arrête à \emph{re(li)quisti}. La fiche EDB (EDB33145) transmet trois lignes supplémentaires : \emph{eterna requiem felicita⟨ti⟩s cuasa man(e)bis IIIIX k⌜a⌝lendas octobris / Cucurbitinus et Abu⌜nd⌝antius hic ⌜s⌝imul quiescit / dd(ominis) nn(ostris) Gratiano V et Teodosio aaugg(ustis)}. L'année 380 est donc fondée sur une formule consulaire lue. Le quantième \emph{IIIIX} est anomal. Deux autres défunts, Cucurbitinus et Abundantius, reposent dans la même tombe ; l'EDB décrit des oiseaux tenant une couronne. La plaque (69 × 128 cm) est conservée à la villa Albani.}
\nf{Commentaire}{Vocabulaire religieux : désignation de Dieu — dei (lu), ancilla dei (lu) ; intercession — ora (lu) ; repos et sommeil de la mort — quiescit (lu). Désignation divine : retenue dans le profil strict. Commentaire développé : dossier p.~\pageref{dos:214}.}
\nf{Bibliographie}{CLE 1619 ; ICUR I 3194 ; ILCV 1464. Bibliographie complémentaire selon l’EDB : ILCV, 1464.}
\nf{Statut}{Lecture secondaire (EDCS) ; divergence forte avec l’EDB : le texte ne doit pas être exploité sans contrôle de l’édition imprimée. L’EDCS signale une photographie (non consultée).}
\end{notice}

\begin{notice}{215}{380}{ICUR n.s. I, 1438 · EDCS-39000421 · EDB31203}
\label{not:215}
\nf{Datation}{380 — directe : nom consulaire lu}
\nf{Genre}{Épitaphe ; ensembles A, B, C ; 69 lettres transmises ; 1 \% restituées.}
\nf{Références}{ICUR I 1438.}
\nf{Provenance}{Rattachement selon l’EDB : \emph{Incertae originis} (le volume I des ICVR réunit les inscriptions dont la provenance est perdue). Conservation selon l’EDB : \emph{Musei Vaticani: Lapidario Cristiano}.}
\nf{Support}{Plaque de marbre. 28 × 85 cm (h. × l.). Inscription gravée. Fiche EDB : Filippo Piazzolla, première version 2015.}
\ntext{\lnum{1}i]nnocentissimo qui ⸢v⸣ixit annis\\
\lnum{2}[- - - I]dus Augustas Gratiano V et T⸢h⸣eo\\
\lnum{3}dosio consulibus}
\nf{Témoins}{Export EDCS de 2020 : \emph{i]nnocentissimo qui ixit annis / [3 I]dus Augustas Gratiano V et Teo/dosio consulibus}. EDB (EDB31203) : texte concordant aux conventions de notation près (indice de similarité 0,964).}
\nf{Apparat}{⸢v⸣ : B gravé ; ⸢h⸣ : A gravé}
\nf{Traduction}{[…] au très innocent, qui vécut […] ans ; […] ides d’août, sous le consulat de Gratien pour la 5e fois et de Théodose.}
\nf{Langue}{L’éditeur signale 2 graphies corrigées : B gravé pour v (bêtacisme) ; A gravé pour h.}
\nf{Chronologie}{Noms consulaires : Gratianus (lu) et Theodosius (lu). Dates du jour : événement indéterminé : … Idus Augustas (partiellement restitué).}
\nf{Bibliographie}{ICUR I 1438. Bibliographie complémentaire selon l’EDB : Ferrua 1964, 108 ; Lega 2000, 167 ; Ferrua 1979b, 37, n. 1438.}
\nf{Statut}{Lecture secondaire (EDCS) concordante avec l’EDB : lecture stabilisée entre deux bases, non collationnée sur l’ICVR ni sur le monument. L’EDCS signale une photographie (non consultée).}
\end{notice}

\begin{notice}{216}{380}{ICUR n.s. I, 3192 · EDCS-38801427 · EDB17710}
\label{not:216}
\nf{Datation}{380 — directe : lecture partielle d’un nom consulaire}
\nf{Genre}{Épitaphe ; ensembles A, B, C ; 45 lettres transmises ; 47 \% restituées.}
\nf{Références}{ICUR I 3192 = ICUR VII 17475.}
\nf{Provenance}{Rattachement selon l’EDB : \emph{Coem. Cyriacae ad s.Laurentium}.}
\nf{Support}{Support non précisé par l’EDB. Fiche EDB : Miriam Ramunni, première version 2014.}
\ntext{\lnum{1}[Gr]atiano V [et Theodosio]\\
\lnum{2}[- - - fecit lo]um cum ipsa sibi e[t]\\
\lnum{3}dum [}
\nf{Témoins}{EDB (EDB17710 ; écarts à examiner) : \emph{------ / [---Gr]atiano V [et Theodosio---] / [---]um cum ipsa sibi e[t ---] / [---]+ dum [---]}.}
\nf{Traduction}{Sous (le consulat de) Gratien pour la 5e fois et Théodose […] a fait un emplacement avec elle-même, pour lui et […].}
\nf{Chronologie}{Noms consulaires : Gratianus (lu en partie) et Theodosius (restitué).}
\nf{Bibliographie}{ICUR I 3192 ; ICUR VII 17475. Bibliographie complémentaire selon l’EDB : ICVR I, 3192.}
\nf{Statut}{Lecture secondaire (EDCS) ; écarts avec l’EDB à arbitrer sur l’ICVR.}
\end{notice}

\begin{notice}{217}{380}{ICUR n.s. II, 6033 · EDCS-34200659 · EDB17434}
\label{not:217}
\nf{Datation}{380 — directe : nom consulaire lu}
\nf{Genre}{Acte d’achat ou de réservation ; ensembles A, B ; 106 lettres transmises ; 34 \% restituées.}
\nf{Références}{ICUR II 6033 = ILCV 3754b (em).}
\nf{Provenance}{Rattachement selon l’EDB : \emph{Coem. Commodillae}.}
\nf{Support}{Plaque de marbre. Inscription gravée. Reproductions signalées par l’EDB : photographie PCAS (droits réservés ; non reproduites). Fiche EDB : Carlo Carletti, première version 2015.}
\ntext{\lnum{1}Gratiano V et Th(e)odosio [conss(ulibus)]\\
\lnum{2}ego Fl(avius) Victor me vi⸢v⸣o s[cripsi]\\
\lnum{3}co⸢m⸣parasse cum Cons[- - -]\\
\lnum{4}[coniuge me]a locum\\
\lnum{5}[bisomum a La]ur() fossorem}
\nf{Témoins}{Export EDCS de 2020 : \emph{Gratiano V et Th(e)odosio [conss(ulibus)] / ego Fl(avius) Victor me vio s[cripsi] / coparasse cum Cons[3] / [coniuge me]a locum / [bisomum a La]ur() fossorem}. EDB (EDB17434) : texte concordant aux conventions de notation près (indice de similarité 0,984).}
\nf{Apparat}{⸢v⸣ : B gravé ; ⸢m⸣ : N gravé}
\nf{Traduction}{Sous le consulat de Gratien pour la 5e fois et de Théodose, moi, Flavius Victor, de mon vivant, j’ai écrit avoir acquis avec Cons[…], mon épouse, un emplacement double auprès du fossoyeur Laur(entius ?).}
\nf{Langue}{L’éditeur signale 2 graphies corrigées : B gravé pour v (bêtacisme) ; N gravé pour m (n pour m).}
\nf{Chronologie}{Noms consulaires : Gratianus (lu) et Theodosius (lu).}
\nf{Commentaire}{Fonction ou métier mentionné : fossorem.}
\nf{Bibliographie}{ICUR II 6033 ; ILCV 3754b (em).}
\nf{Statut}{Lecture secondaire (EDCS) concordante avec l’EDB : lecture stabilisée entre deux bases, non collationnée sur l’ICVR ni sur le monument. L’EDCS signale une photographie (non consultée).}
\end{notice}

\begin{notice}{218}{380}{ICUR n.s. IV, 11774 · EDCS-37400862 · EDB5272}
\label{not:218}
\nf{Datation}{380 — directe : nom consulaire lu}
\nf{Genre}{Épitaphe ; ensembles A, B, C ; 85 lettres transmises ; 39 \% restituées.}
\nf{Références}{ICUR IV 11774 = ILCV 4248a (em).}
\nf{Provenance}{Rattachement selon l’EDB : \emph{Coem. anonymum ad viam Appiam}. Conservation selon l’EDB : \emph{Coem. anonymum ad viam Appiam}.}
\nf{Support}{Plaque de marbre. Inscription gravée. Fiche EDB : Maria Cristina Grisanzio, première version 2013.}
\ntext{\lnum{1}] amoris Pauli[- - -]\\
\lnum{2}[- - -]am cast(a)e fuit vit[ae - - -]\\
\lnum{3}[- - - cu]m virginio suo ann[os - - -]\\
\lnum{4}[- - - dd(ominis) n]n(ostris) Grat(iano) Aug(usto) V\{I\} et T[heodosio cons(ulibus)}
\nf{Témoins}{EDB (EDB5272 ; écarts à examiner) : \emph{[------] / amoris Pauli / [---]am caste fuit vit[ae] / [cu]m virginio suo ann[is] dd(ominis) n Grat(iano) Aug(usto) V⌜I⌝ et T[heodosio ---]}.}
\nf{Apparat}{\{I\} : lettres superflues}
\nf{Traduction}{[…] de l’amour de Paulus (?) […] ; […] elle fut d’une vie chaste […] ; […] avec son ‘vierge’ (époux) […] ans ; sous le consulat de nos seigneurs Gratien Auguste pour la 5e fois et Théodose.}
\nf{Langue}{Des lettres gravées sont jugées superflues par l’éditeur (entre accolades). 6 abréviations développées par l’éditeur (dont cast(a)e, dd(ominis), nn(ostris), Grat(iano), Aug(usto), cons(ulibus)).}
\nf{Chronologie}{Noms consulaires : Gratianus (lu) et Theodosius (restitué). Ne sont pas des consuls : amoris Pauli[ (personne nommée dans l’éloge (non consul en 380)).}
\nf{Bibliographie}{ICUR IV 11774 ; ILCV 4248a (em).}
\nf{Statut}{Lecture secondaire (EDCS) ; écarts avec l’EDB à arbitrer sur l’ICVR. L’EDCS signale une photographie (non consultée).}
\end{notice}

\annee{381}

\begin{notice}{219}{381}{ICUR n.s. I, 1209 · EDCS-38700748 · EDB30156}
\label{not:219}
\nf{Datation}{381 — directe : consulat restitué}
\nf{Genre}{Épitaphe ; ensembles A, B, C ; 36 lettres transmises ; 50 \% restituées.}
\nf{Références}{ICUR I 1209.}
\nf{Provenance}{Rattachement selon l’EDB : \emph{Incertae originis} (le volume I des ICVR réunit les inscriptions dont la provenance est perdue). Conservation selon l’EDB : \emph{Musei Capitolini}.}
\nf{Support}{Plaque de marbre. 21 × 30 cm (h. × l.). Inscription gravée. Fiche EDB : Filippo Piazzolla, première version 2014.}
\ntext{\lnum{1}]LI[- - -]\\
\lnum{2}[- - - qua]e vixit cu[m marito suo]\\
\lnum{3}[- - -]IIII ⸢F⸣l(avio) Sy[agrio}
\nf{Témoins}{Export EDCS de 2020 : \emph{]LI[3] / [3 qua]e vixit cu[m marito suo] / [3]IIII l(avio) Sy[agrio}. EDB (EDB30156) : texte concordant aux conventions de notation près (indice de similarité 0,947).}
\nf{Apparat}{⸢F⸣ : E gravé}
\nf{Traduction}{[…] qui vécut avec son mari […] 4 […] ; (sous le consulat de) Flavius Syagrius.}
\nf{Langue}{L’éditeur signale une graphie corrigée : E gravé pour F.}
\nf{Chronologie}{Noms consulaires : Syagrius (restitué).}
\nf{Bibliographie}{ICUR I 1209. Bibliographie complémentaire selon l’EDB : ICVR I, 3196 ; Ferrua 1979b, 29, n. 1209.}
\nf{Statut}{Lecture secondaire (EDCS) concordante avec l’EDB : lecture stabilisée entre deux bases, non collationnée sur l’ICVR ni sur le monument.}
\end{notice}

\begin{notice}{220}{381}{ICUR n.s. I, 1352 · EDCS-39000330 · EDB30687}
\label{not:220}
\nf{Datation}{381 — directe : lecture partielle d’un nom consulaire ; datation contestée}
\nf{Genre}{Épitaphe ; ensembles A, B, C ; 64 lettres transmises ; 83 \% restituées.}
\nf{Références}{ICUR I 1352 (p 490).}
\nf{Provenance}{Rattachement selon l’EDB : \emph{Incertae originis} (le volume I des ICVR réunit les inscriptions dont la provenance est perdue). Conservation selon l’EDB : \emph{Collegio Teutonico}.}
\nf{Support}{Plaque de marbre. 28 × 13 cm (h. × l.). Inscription gravée. Fiche EDB : Filippo Piazzolla, première version 2015.}
\ntext{\lnum{1}[- - -]ECIM[- - -]\\
\lnum{2}[- - -]NOM[- - -]\\
\lnum{3}[- - - Fl(avio) Sygario et Fl(avio) Euch]erio [vv(iris) cc(larissimis) cons(ulibus)}
\nf{Témoins}{EDB (EDB30687 ; divergence forte) : \emph{[---]ecim[---] / [---]nom[---] / [--- et Fl(avio Euch?]erio [vv(iris) cc(larissimis) cons(ulibu)s] / [---?] ((avis cum olea )) [---?] [texte tronqué dans la liste de résultats]}.}
\nf{Traduction}{[…] sous le consulat de Flavius Syagrius et de Flavius Eucherius, clarissimes.}
\nf{Langue}{5 abréviations développées par l’éditeur (dont Fl(avio), vv(iris), cc(larissimis), cons(ulibus)).}
\nf{Chronologie}{Noms consulaires : Eucherius (lu en partie). Seules les lettres \emph{erio} sont lues. L'EDCS restitue \emph{[Fl(avio) Sygario et Fl(avio) Euch]erio} sans réserve ; l'EDB (EDB30687) restitue \emph{[--- et Fl(avio) Euch?]erio} et signale d'autres lectures possibles (\emph{Val]erio}, \emph{Gal]erio}, \emph{Neot]erio}, \emph{Ast]erio}). L'année 381 est conservée comme année de classement ; 390 (Neoterius) reste possible. Datation contestée.}
\nf{Bibliographie}{ICUR I 1352 (p 490).}
\nf{Statut}{Lecture secondaire (EDCS) ; divergence forte avec l’EDB : le texte ne doit pas être exploité sans contrôle de l’édition imprimée. Plus de la moitié des lettres sont restituées.}
\end{notice}

\begin{notice}{221}{381}{ICUR n.s. I, 2095 · EDCS-39300317 · EDB32075}
\label{not:221}
\nf{Datation}{381 — directe : nom consulaire lu}
\nf{Genre}{Épitaphe ; ensembles A, B, C ; 144 lettres transmises ; 38 \% restituées.}
\nf{Références}{ICUR I 2095.}
\nf{Provenance}{Rattachement selon l’EDB : \emph{Incertae originis} (le volume I des ICVR réunit les inscriptions dont la provenance est perdue). Conservation selon l’EDB : \emph{Musei Vaticani: Galleria Lapidaria}.}
\nf{Support}{Plaque de marbre. 18 × 22 cm (h. × l.). Inscription gravée. Fiche EDB : Filippo Piazzolla, première version 2015.}
\ntext{\lnum{1}[Mirae bonitati]s et innocentiae Fl(aviae) N[- - -]ae Iuliae innocen[tissimae]\\
\lnum{2}[filiae(?) quae vixit] annos V et menses IIII de[posi]ta in pace die IIII[- - -]\\
\lnum{3}[- - - Sya]⸢g⸣rio et Eucherio vv(iris) c[c(larissimis) c]ons(ulibus) s(upra) s(criptis)}
\nf{Témoins}{Export EDCS de 2020 : \emph{[Mirae bonitati]s et innocentiae Fl(aviae) N[3]ae Iuliae innocen[tissimae] / [filiae(?) quae vixit] annos V et menses IIII de[posi]ta in pace die IIII[3] / [3 Sya]rio et Eucherio vv(iris) c[c(larissimis) c]ons(ulibus) s(upra) s(criptis)}. EDB (EDB32075) : texte concordant aux conventions de notation près (indice de similarité 0,961).}
\nf{Apparat}{⸢g⸣ : C gravé}
\nf{Traduction}{D’une bonté et d’une innocence admirables, à Flavia N[…]a Iulia, fille très innocente, qui vécut 5 ans et 4 mois ; déposée en paix le 4e jour […] ; […] sous le consulat susdit de Syagrius et d’Eucherius, clarissimes.}
\nf{Langue}{L’éditeur signale une graphie corrigée : C gravé pour g (c pour g).}
\nf{Chronologie}{Noms consulaires : Syagrius (lu en partie) et Eucherius (lu). Âge déclaré : 5 a. 4 m. (partiellement restitué).}
\nf{Commentaire}{La formule in pace est conservée.}
\nf{Bibliographie}{ICUR I 2095. Bibliographie complémentaire selon l’EDB : ILCV, 4332 adn. ; Lega 2000.}
\nf{Statut}{Lecture secondaire (EDCS) concordante avec l’EDB : lecture stabilisée entre deux bases, non collationnée sur l’ICVR ni sur le monument.}
\end{notice}

\begin{notice}{222}{381}{ICUR n.s. I, 3196 · EDCS-40300339 · EDB33147}
\label{not:222}
\nf{Datation}{381 — directe : consulat restitué}
\nf{Genre}{Épitaphe ; ensembles A, B, C ; 30 lettres transmises ; 57 \% restituées.}
\nf{Références}{ICUR I 3196.}
\nf{Provenance}{Rattachement selon l’EDB : \emph{Incertae originis} (le volume I des ICVR réunit les inscriptions dont la provenance est perdue). Conservation selon l’EDB : \emph{Musei Capitolini}.}
\nf{Support}{Plaque de marbre. 21 × 30 cm (h. × l.). Inscription gravée. Fiche EDB : Filippo Piazzolla, première version 2014.}
\ntext{\lnum{1}qua]e vix(it) pl[us minus ann(os) - - -]\\
\lnum{2}[- - -] Fl(avio) Sy[agrio}
\nf{Témoins}{EDB (EDB33147 ; divergence forte) : \emph{------? / [---]E+[---] / [--- qua]e vixit cu[m marito suo] / [---] IIII ⌜F⌝l(avio) Sy[agrio ---] [texte tronqué dans la liste de résultats]}. Relecture de la fiche détaillée de l’EDB : Même texte pour la date (Fl(avio) Sy[agrio). L'EDB suit ICVR I 1209 et non I 1396 (Ferrua 1979b).}
\nf{Traduction}{[…] qui vécut environ […] ans ; (sous le consulat de) Flavius Syagrius.}
\nf{Langue}{3 abréviations développées par l’éditeur (dont vix(it), ann(os), Fl(avio)).}
\nf{Chronologie}{Noms consulaires : Syagrius (restitué).}
\nf{Bibliographie}{ICUR I 3196. Bibliographie complémentaire selon l’EDB : ICVR I, 1209 ; Ferrua 1979b, 92, n. 3196.}
\nf{Statut}{Lecture secondaire (EDCS) ; divergence forte avec l’EDB : le texte ne doit pas être exploité sans contrôle de l’édition imprimée. Plus de la moitié des lettres sont restituées.}
\end{notice}

\begin{notice}{223}{381}{ICUR n.s. II, 6034 · EDCS-34200660 · EDB17435}
\label{not:223}
\nf{Datation}{381 — directe : nom consulaire lu}
\nf{Genre}{Épitaphe ; ensembles A, B, C ; 112 lettres transmises.}
\nf{Références}{ICUR II 6034 = ILCV 2851.}
\nf{Provenance}{Rattachement selon l’EDB : \emph{Coem. Commodillae}. Conservation selon l’EDB : \emph{Musei Vaticani: Lapidario Cristiano}.}
\nf{Support}{Plaque de marbre. Inscription gravée. Fiche EDB : Carlo Carletti, première version 2015.}
\ntext{\lnum{1}Paterniae bene merenti\\
\lnum{2}in pace qu(a)e vixit annis\\
\lnum{3}plus m(inus) XL et fecit cu\\
\lnum{4}m marito suo an(nos) XX\\
\lnum{5}discessit pri(die) Kal(endas)\\
\lnum{6}Octob(res) Fl(avio) S⸢y⸣a\\
\lnum{7}grio et Eu⸢ch⸣e\\
\lnum{8}rio cons(ulibus)}
\nf{Témoins}{Export EDCS de 2020 : \emph{Paterniae bene merenti / in pace qu(a)e vixit annis / plus m(inus) XL et fecit cu/m marito suo an(nos) XX / discessit pri(die) Kal(endas) / Octob(res) Fl(avio) Sa/grio et Eue/rio cons(ulibus)}. EDB (EDB17435) : texte concordant aux conventions de notation près (indice de similarité 0,971).}
\nf{Apparat}{⸢y⸣ : V gravé ; ⸢ch⸣ : T gravé}
\nf{Traduction}{À Paternia, bien méritante, en paix, qui vécut environ 40 ans et a vécu avec son mari 20 ans ; elle s’en alla la veille des calendes d’octobre, sous le consulat de Flavius Syagrius et d’Eucherius.}
\nf{Langue}{L’éditeur signale 2 graphies corrigées : V gravé pour y (graphie de y) ; T gravé pour ch.}
\nf{Chronologie}{Noms consulaires : Syagrius (lu) et Eucherius (lu). Dates du jour : décès : pridie Kalendas Octobres = 30 septembre 381 (lu). Âge déclaré : env. 40 a. (lu).}
\nf{Commentaire}{La formule in pace est conservée.}
\nf{Bibliographie}{ICUR II 6034 ; ILCV 2851.}
\nf{Statut}{Lecture secondaire (EDCS) concordante avec l’EDB : lecture stabilisée entre deux bases, non collationnée sur l’ICVR ni sur le monument. L’EDCS signale une photographie (non consultée).}
\end{notice}

\begin{notice}{224}{381}{ICUR n.s. IV, 9571 · EDCS-36101097 · EDB17960}
\label{not:224}
\nf{Datation}{381 — directe : lecture partielle d’un nom consulaire}
\nf{Genre}{Épitaphe ; ensembles A, B, C ; 32 lettres transmises ; 62 \% restituées.}
\nf{Références}{ICUR IV 9571.}
\nf{Provenance}{Rattachement selon l’EDB : \emph{Coem. Callisti pars inferior}.}
\nf{Support}{Plaque de marbre. Inscription gravée. Fiche EDB : Antonio Enrico Felle, première version 2003.}
\ntext{\lnum{1}]II die(s) XI\\
\lnum{2}[- - - cons(ulatu) Syagri et Euc]herii}
\nf{Témoins}{EDB (EDB17960) : texte concordant aux conventions de notation près (indice de similarité 0,958).}
\nf{Traduction}{[…] 2 (mois ?) 11 jours ; […] sous le consulat de Syagrius et d’Eucherius.}
\nf{Langue}{2 abréviations développées par l’éditeur (dont die(s), cons(ulatu)).}
\nf{Chronologie}{Noms consulaires : Syagrius (restitué) et Eucherius (lu en partie).}
\nf{Bibliographie}{ICUR IV 9571.}
\nf{Statut}{Lecture secondaire (EDCS) concordante avec l’EDB : lecture stabilisée entre deux bases, non collationnée sur l’ICVR ni sur le monument. Plus de la moitié des lettres sont restituées. L’EDCS signale une photographie (non consultée).}
\end{notice}

\begin{notice}{225}{381}{ICUR n.s. IV, 9572 · EDCS-36101098 · EDB17961}
\label{not:225}
\nf{Datation}{381 — directe : nom consulaire lu}
\nf{Genre}{Épitaphe ; ensembles A, B, C ; 49 lettres transmises.}
\nf{Références}{ICUR IV 9572.}
\nf{Provenance}{Rattachement selon l’EDB : \emph{Coem. Callisti pars inferior}.}
\nf{Support}{Plaque de marbre. Inscription gravée. Signes décrits par l’EDB : chrisme avec alpha et oméga. Reproductions signalées par l’EDB : photographie PCAS (droits réservés ; non reproduites). Fiche EDB : Antonio Enrico Felle, première version 2003.}
\ntext{\lnum{1}q]uae vixit ann⸢o⸣s p(lus) m(inus) XL d(e)p(osita) d(ie) IV [- - -]\\
\lnum{2}[- - -]s Syagrio et Euc(h)erio vv(iris) cc(larissimis) conss(ulibus)}
\nf{Témoins}{Export EDCS de 2020 : \emph{q]uae vixit anns p(lus) m(inus) XL d(e)p(osita) d(ie) IV [3] / [3]s Syagrio et Euc(h)erio vv(iris) cc(larissimis) conss(ulibus)}. EDB (EDB17961) : texte concordant aux conventions de notation près (indice de similarité 0,985).}
\nf{Apparat}{⸢o⸣ : V gravé}
\nf{Traduction}{[…] qui vécut environ 40 ans ; déposée le 4e jour […] ; sous le consulat de Syagrius et d’Eucherius, clarissimes.}
\nf{Langue}{L’éditeur signale une graphie corrigée : V gravé pour o (confusion o/u).}
\nf{Chronologie}{Noms consulaires : Syagrius (lu) et Eucherius (lu). Âge déclaré : env. 40 a. 4 j. (lu).}
\nf{Bibliographie}{ICUR IV 9572.}
\nf{Statut}{Lecture secondaire (EDCS) concordante avec l’EDB : lecture stabilisée entre deux bases, non collationnée sur l’ICVR ni sur le monument. L’EDCS signale une photographie (non consultée).}
\end{notice}

\begin{notice}{226}{381}{ICUR n.s. IV, 11109 · EDCS-37200391 · EDB39571}
\label{not:226}
\nf{Datation}{381 — directe : nom consulaire lu}
\nf{Genre}{Épitaphe ; ensembles A, B, C ; 66 lettres transmises ; 9 \% restituées.}
\nf{Références}{ICUR IV 11109.}
\nf{Provenance}{Rattachement selon l’EDB : \emph{Coem. Callisti pars inferior}. Conservation selon l’EDB : \emph{Coem. Callisti area subdialis et incerti loci}.}
\nf{Support}{Plaque de marbre. 35 × 70 cm (h. × l.). Inscription gravée. Reproductions signalées par l’EDB : photographie PCAS ×2 (droits réservés ; non reproduites). Fiche EDB : Marida Pierno, première version 2015.}
\ntext{\lnum{1}Dulcissima[e - - -]\\
\lnum{2}qu(a)e vixit ann(os) p(lus) m(inus) [- - -]\\
\lnum{3}d(e)p(osita) XVII K(alendas) Sept(embres) FF(laviis) Syagrio [et Eu]\\
\lnum{4}caerio co(n)ss(ulibus) b(ene) m(erenti) in pac\{a\}[e}
\nf{Témoins}{EDB (EDB39571) : texte concordant aux conventions de notation près (indice de similarité 0,966).}
\nf{Apparat}{\{a\} : lettres superflues}
\nf{Traduction}{À la très douce […] qui vécut environ […] ans ; déposée le 17e jour avant les calendes de septembre, sous le consulat des Flavii Syagrius et Eucherius ; bien méritante, en paix.}
\nf{Langue}{Des lettres gravées sont jugées superflues par l’éditeur (entre accolades). 13 abréviations développées par l’éditeur (dont qu(a)e, ann(os), p(lus), m(inus), d(e)p, (osita)).}
\nf{Chronologie}{Noms consulaires : Syagrius (lu). Dates du jour : déposition : XVII Kalendas Septembres = 16 août 381 (lu).}
\nf{Commentaire}{La formule in pace est attestée seulement par restitution.}
\nf{Bibliographie}{ICUR IV 11109.}
\nf{Statut}{Lecture secondaire (EDCS) concordante avec l’EDB : lecture stabilisée entre deux bases, non collationnée sur l’ICVR ni sur le monument. L’EDCS signale une photographie (non consultée).}
\end{notice}

\begin{notice}{227}{381}{ICUR n.s. IV, 12259e · EDCS-35700732 · EDB4654}
\label{not:227}
\nf{Datation}{381 — directe : lecture partielle d’un nom consulaire}
\nf{Genre}{Épitaphe ; ensembles A, B, C ; 7 lettres transmises ; 29 \% restituées.}
\nf{Références}{ICUR IV 12259e.}
\nf{Provenance}{Rattachement selon l’EDB : \emph{Coem. anonymum ad viam Ardeatinam}. Conservation selon l’EDB : \emph{Coem. anonymum ad viam Ardeatinam}.}
\nf{Support}{Plaque de marbre. Inscription gravée. Fiche EDB : Marida Pierno, première version 2003.}
\ntext{\lnum{1}]t Euce[ri(?)}
\nf{Témoins}{EDB (EDB4654) : texte concordant aux conventions de notation près (indice de similarité 1).}
\nf{Traduction}{[…] Eucherius.}
\nf{Langue}{1 abréviation développée par l’éditeur (dont Euceri(?)).}
\nf{Chronologie}{Noms consulaires : Eucherius (lu en partie).}
\nf{Bibliographie}{ICUR IV 12259e.}
\nf{Statut}{Lecture secondaire (EDCS) concordante avec l’EDB : lecture stabilisée entre deux bases, non collationnée sur l’ICVR ni sur le monument.}
\end{notice}

\annee{382}

\begin{notice}{228}{382}{ICUR n.s. II, 6036 · EDCS-16600356 · EDB3284}
\label{not:228}
\nf{Datation}{382 — directe : lecture partielle d’un nom consulaire}
\nf{Genre}{Épitaphe ; ensembles A, B, C ; 39 lettres transmises ; 46 \% restituées.}
\nf{Références}{AE 1905, 75 = ICUR II 6036.}
\nf{Provenance}{Rattachement selon l’EDB : \emph{Coem. Commodillae}.}
\nf{Support}{Plaque de marbre. Inscription gravée. Fiche EDB : Carlo Carletti, première version 2003.}
\ntext{\lnum{1}vixit an]nos [- - -]\\
\lnum{2}[- - -] Cl(audio) Anto[nio et]\\
\lnum{3}[Fl(avio) Syagrio] conss(ulibus)}
\nf{Témoins}{EDB (EDB3284) : texte concordant aux conventions de notation près (indice de similarité 1).}
\nf{Traduction}{[…] vécut […] ans ; […] sous le consulat de Claudius Antonius et de Flavius Syagrius.}
\nf{Langue}{3 abréviations développées par l’éditeur (dont Cl(audio), Fl(avio), conss(ulibus)).}
\nf{Chronologie}{Noms consulaires : Syagrius (restitué) et Antonius (lu en partie).}
\nf{Bibliographie}{AE 1905, 75 ; ICUR II 6036.}
\nf{Statut}{Lecture secondaire (EDCS) concordante avec l’EDB : lecture stabilisée entre deux bases, non collationnée sur l’ICVR ni sur le monument.}
\end{notice}

\begin{notice}{229}{382}{ICUR n.s. I, 2096 · EDCS-19400473 · EDB32081}
\label{not:229}
\nf{Datation}{382 — directe : nom consulaire lu}
\nf{Genre}{Épitaphe ; ensembles A, B, C ; 42 lettres transmises ; 43 \% restituées.}
\nf{Références}{CIL VI 9787 = ICUR I 2096 = ILCV +669 = ICVaticano p 334 = Pittori 8 = AE 1997, 166.}
\nf{Provenance}{Rattachement selon l’EDB : \emph{Incertae originis} (le volume I des ICVR réunit les inscriptions dont la provenance est perdue). Conservation selon l’EDB : \emph{Musei Vaticani: Galleria Lapidaria}.}
\nf{Support}{Plaque de marbre. 19,3 × 32,5 cm (h. × l.). Inscription gravée. Fiche EDB : Filippo Piazzolla, première version 2015.}
\ntext{\lnum{1}Aur(elius) Felix pict[or]\\
\lnum{2}Cl(audio) Antonio Fl(avio) S[yagrio cons(ulibus)]}
\nf{Témoins}{EDB (EDB32081) : texte concordant aux conventions de notation près (indice de similarité 0,957).}
\nf{Traduction}{Aurelius Felix, peintre ; sous le consulat de Claudius Antonius et de Flavius Syagrius.}
\nf{Langue}{4 abréviations développées par l’éditeur (dont Aur(elius), Cl(audio), Fl(avio), cons(ulibus)).}
\nf{Chronologie}{Noms consulaires : Syagrius (restitué) et Antonius (lu).}
\nf{Commentaire}{Fonction ou métier mentionné : pictor.}
\nf{Bibliographie}{CIL VI 9787 ; ICUR I 2096 ; ILCV +669 ; ICVaticano p 334 ; Pittori 8 ; AE 1997, 166. Bibliographie complémentaire selon l’EDB : Di Stefano Manzella 1997, 334-335 n. 3.11.14 con foto (C. Lega) ; CIL VI, 9787 ; AE, 1997, 166 ; CLRE, 650 ; Lega 2000 ; Bisconti 2000a, 247 n.XIXb2.1 fig.83 ; ILCV, 669.}
\nf{Statut}{Lecture secondaire (EDCS) concordante avec l’EDB : lecture stabilisée entre deux bases, non collationnée sur l’ICVR ni sur le monument. L’EDCS signale une photographie (non consultée).}
\end{notice}

\begin{notice}{230}{382}{ICUR n.s. VI, 17248 · EDCS-32900465 · EDB12805}
\label{not:230}
\nf{Datation}{382 — directe : nom consulaire lu}
\nf{Genre}{Épitaphe ; ensembles A, B, C ; 74 lettres transmises ; 54 \% restituées.}
\nf{Références}{Collatia p 639 = ICUR VI 17248 = AnalEpi p 69 = AE 1939, 152 = AE 1949, 163.}
\nf{Provenance}{Rattachement selon l’EDB : \emph{Coem. s.Zotici}.}
\nf{Support}{Plaque de marbre. Inscription gravée. Fiche EDB : Carlo Carletti, première version 2015.}
\ntext{\lnum{1}Abudantia qu[ae vixit annis - - -]\\
\lnum{2}menses VII de[p(osita) - - -]\\
\lnum{3}Fl(avio) Antonio et F[l(avio) Syagrio cons(ulibus)]}
\nf{Témoins}{EDB (EDB12805 ; écarts à examiner) : \emph{Abudantia qu[ae vixit annis ---] / menses VII de[p(ositus) ---] /Fl(avio) Antonio et F[l(avio) Syagrio cons(ulibus)] [texte tronqué dans la liste de résultats]}.}
\nf{Traduction}{Abundantia, qui vécut […] ans 7 mois ; déposée […] ; sous le consulat de Flavius (sic) Antonius et de Flavius Syagrius.}
\nf{Langue}{4 abréviations développées par l’éditeur (dont dep(osita), Fl(avio), cons(ulibus)).}
\nf{Chronologie}{Noms consulaires : Syagrius (restitué) et Antonius (lu). Âge déclaré : 7 m. (partiellement restitué).}
\nf{Bibliographie}{Collatia p 639 ; ICUR VI 17248 ; AnalEpi p 69 ; AE 1939, 152 ; AE 1949, 163.}
\nf{Statut}{Lecture secondaire (EDCS) ; écarts avec l’EDB à arbitrer sur l’ICVR. Plus de la moitié des lettres sont restituées.}
\end{notice}

\begin{notice}{231}{382}{ICUR n.s. I, 1439 · EDCS-39000422 · EDB31204}
\label{not:231}
\nf{Datation}{382 — directe : nom consulaire lu}
\nf{Genre}{Épitaphe ; ensembles A, B, C ; 199 lettres transmises.}
\nf{Références}{ICUR I 1439 = ILCV 4700.}
\nf{Provenance}{Rattachement selon l’EDB : \emph{Incertae originis} (le volume I des ICVR réunit les inscriptions dont la provenance est perdue). Conservation selon l’EDB : \emph{Musei Vaticani: Lapidario Cristiano}.}
\nf{Support}{Plaque de marbre. 30 × 58 cm (h. × l.). Inscription gravée. Fiche EDB : Filippo Piazzolla, première version 2015.}
\ntext{\lnum{1}Innocens recessit innocentium misericors fuit\\
\lnum{2}quis non doluit etati tuae piasq(ue) lacrimas fudit\\
\lnum{3}in te spes futura ex(s)pectabatur \{i per te\}\\
\lnum{4}per te ⸢g⸣loria perennis Celerine fili\\
\lnum{5}fidelis quiescis in pace ⸢q⸣ui ⸢v⸣ixi(t) ann(um) I\\
\lnum{6}m(enses) IIII d(epositus) VI Kal(endas) Oct(obres) Ant⸢o⸣n(i)o e⸢t⸣ Syagrio}
\nf{Témoins}{Export EDCS de 2020 : \emph{Innocens recessit innocentium misericors fuit / quis non doluit etati tuae piasq(ue) lacrimas fudit / in te spes futura ex(s)pectabatur \{i per te\} / per te loria perennis Celerine fili / fidelis quiescis in pace ui ixi(t) ann(um) I / m(enses) IIII d(epositus) VI Kal(endas) Oct(obres) Antn(i)o e Syagrio}. EDB (EDB31204 ; divergence forte) : \emph{innocens recessit innocen̂tium mi`s´eric`o´rs fui`t´ / quis non(as) doluit ⟨a⟩etati tuae piasq(ue) lacrima`s´ fud`it´ / in te spes futura expectabatur 〚i〛 per te / per te ⌜g⌝loria perennis Celerine fili / fidelis quiescis in pace ⌜q⌝ui ⟨v⟩i\{i\}xit ann(is) I / m(ensibus) IIII d(iebus) VI kal(endas) oct(obres) Antino e⌜t⌝ Syagrio}. Relecture de la fiche détaillée de l’EDB : Antino (EDB, d'après Lega) / Ant<o=I>n(i)o (EDCS) ; d(iebus) / d(epositus).}
\nf{Apparat}{⸢g⸣ : C gravé ; ⸢q⸣ : O gravé ; ⸢v⸣ : I gravé ; ⸢o⸣ : I gravé ; ⸢t⸣ : I gravé ; \{i per te\} : lettres superflues}
\nf{Traduction}{Innocent, il s’en est allé, (enfant) innocent et miséricordieux ; qui n’a pas souffert de ton âge (si tendre) et n’a pas versé de pieuses larmes ? En toi était attendue l’espérance de l’avenir, par toi une gloire durable, Celerinus, mon fils ; fidèle, tu reposes en paix ; il vécut 1 an 4 mois ; déposé le 6e jour avant les calendes d’octobre, sous (le consulat d’)Antonius et Syagrius.}
\nf{Langue}{L’éditeur signale 5 graphies corrigées : C gravé pour g (c pour g) ; O gravé pour q (confusion de lettres) ; I gravé pour v ; I gravé pour o ; I gravé pour t. Des lettres gravées sont jugées superflues par l’éditeur (entre accolades). 9 abréviations développées par l’éditeur (dont piasq(ue), ex(s)pectabatur, ixi(t), ann(um), m(enses), d(epositus)).}
\nf{Chronologie}{Noms consulaires : Syagrius (lu) et Antonius (lu). Dates du jour : déposition : VI Kalendas Octobres = 26 septembre 382 (lu). Âge déclaré : 1 a. 4 m. (lu).}
\nf{Commentaire}{La formule in pace est conservée. Vocabulaire religieux : initiation et foi — fidelis (lu) ; repos et sommeil de la mort — quiescis (lu). Commentaire développé : dossier p.~\pageref{dos:231}.}
\nf{Bibliographie}{ICUR I 1439 ; ILCV 4700. Bibliographie complémentaire selon l’EDB : ILCV, 4700 ; Lega 2000, 166.}
\nf{Statut}{Lecture secondaire (EDCS) ; divergence forte avec l’EDB : le texte ne doit pas être exploité sans contrôle de l’édition imprimée. L’EDCS signale une photographie (non consultée).}
\end{notice}

\begin{notice}{232}{382}{ICUR n.s. I, 3127 · EDCS-39900589 · EDB33081}
\label{not:232}
\nf{Datation}{382 — directe : lecture partielle d’un nom consulaire}
\nf{Genre}{Épitaphe ; ensembles A, B, C ; 150 lettres transmises ; 42 \% restituées.}
\nf{Références}{ICUR I 3127 = ILCV 2148.}
\nf{Provenance}{Rattachement selon l’EDB : \emph{Incertae originis} (le volume I des ICVR réunit les inscriptions dont la provenance est perdue).}
\nf{Support}{Support non précisé par l’EDB. Fiche EDB : Filippo Piazzolla, première version 2015.}
\ntext{\lnum{1}[- - -]na in dom(ino ?) cultrix\\
\lnum{2}[- - -]P nutrivit\\
\lnum{3}[- - - s]uis\\
\lnum{4}[- - - amatrix paup]e\{ro\}rum\\
\lnum{5}[quae pro tanta] merita accepit\\
\lnum{6}[sepulcrum intra l]imina san(c)torum\\
\lnum{7}[- - -]t accepit\\
\lnum{8}[quod multi cupiu]n(t) et rari accipiun(t)\\
\lnum{9}[- - - Antonio et S]⸢y⸣agrio c(on)s(ulibus)}
\nf{Témoins}{Export EDCS de 2020 : \emph{[3]na in dom(ino ?) cultrix / [3]P nutrivit / [3 s]uis /[3 amatrix paup]e\{ro\}rum / [quae pro tanta] merita accepit / [sepulcrum intra l]imina san(c)torum / [3]t accepit / [quod multi cupiu]n(t) et rari accipiun(t) / [3 Antonio et S]agrio c(on)s(ulibus)}. EDB (EDB33081) : texte concordant aux conventions de notation près (indice de similarité 0,945).}
\nf{Apparat}{⸢y⸣ : U gravé ; \{ro\} : lettres superflues}
\nf{Traduction}{[…]na, in dom(ino ?) cultrix : « dans le Seigneur » est une interprétation douteuse ; […] a nourri […] ; […] les siens ; [amie des pauv]res ; [pour de tels] mérites, reçut [une sépulture à l’intérieur des] seuils des saints ; […] reçut [ce que beaucoup désire]nt et que peu obtiennent. Traduction lacunaire, non texte continu assuré.}
\nf{Langue}{L’éditeur signale une graphie corrigée : U gravé pour y. Des lettres gravées sont jugées superflues par l’éditeur (entre accolades). 6 abréviations développées par l’éditeur (dont dom(ino ?), san(c)torum, cupiun(t), accipiun(t), c(on)s, (ulibus)).}
\nf{Chronologie}{Noms consulaires : Syagrius (lu en partie) et Antonius (restitué).}
\nf{Commentaire}{Dom(ino ?) : développement douteux, exclu du scénario restrictif. San(c)torum : abrégé, non entièrement écrit. L'abréviation \emph{dom(ino ?)} est développée avec un point d'interrogation par l'éditeur. Le texte est retenu dans le profil large, exclu du profil restrictif. Désignation divine : retenue dans le profil large seulement. Commentaire développé : dossier p.~\pageref{dos:232}.}
\nf{Bibliographie}{ICUR I 3127 ; ILCV 2148. Bibliographie complémentaire selon l’EDB : Ferrua 1979b, 88 ; ILCV, 2148.}
\nf{Statut}{Lecture secondaire (EDCS) concordante avec l’EDB : lecture stabilisée entre deux bases, non collationnée sur l’ICVR ni sur le monument.}
\end{notice}

\begin{notice}{233}{382}{ICUR n.s. I, 3197 · EDCS-40300340 · EDB33148}
\label{not:233}
\nf{Datation}{382 — directe : nom consulaire lu}
\nf{Genre}{Épitaphe ; ensembles A, B, C ; 58 lettres transmises.}
\nf{Références}{ICUR I 3197 = ILCV 1479a.}
\nf{Provenance}{Rattachement selon l’EDB : \emph{Incertae originis} (le volume I des ICVR réunit les inscriptions dont la provenance est perdue).}
\nf{Support}{Support non précisé par l’EDB. Fiche EDB : Filippo Piazzolla, première version 2015.}
\ntext{\lnum{1}Victori ne\\
\lnum{2}of⸢y⸣to in pace\\
\lnum{3}filio Lauren\\
\lnum{4}tia d(e)p(osita) XVIIII K(alendas) Sep(tembres)\\
\lnum{5}Ant(onio) et Syagrio}
\nf{Témoins}{Export EDCS de 2020 : \emph{Victori ne/ofto in pace / filio Lauren/tia d(e)p(osita) XVIIII K(alendas) Sep(tembres) / Ant(onio) et Syagrio}. EDB (EDB33148) : texte concordant aux conventions de notation près (indice de similarité 0,928).}
\nf{Apparat}{⸢y⸣ : I gravé}
\nf{Traduction}{À Victor, néophyte, en paix, à son fils, Laurentia ; déposé le 19e jour avant les calendes de septembre, sous (le consulat d’)Antonius et Syagrius.}
\nf{Langue}{L’éditeur signale une graphie corrigée : I gravé pour y (graphie de y).}
\nf{Chronologie}{Noms consulaires : Syagrius (lu) et Antonius (lu). Dates du jour : déposition : XVIIII Kalendas Septembres = 14 août 382 (lu).}
\nf{Commentaire}{La formule in pace est conservée. Vocabulaire religieux : initiation et foi — neofyto (lu).}
\nf{Bibliographie}{ICUR I 3197 ; ILCV 1479a. Bibliographie complémentaire selon l’EDB : ILCV, 1479A.}
\nf{Statut}{Lecture secondaire (EDCS) concordante avec l’EDB : lecture stabilisée entre deux bases, non collationnée sur l’ICVR ni sur le monument.}
\end{notice}

\begin{notice}{234}{382}{ICUR n.s. II, 4167 · EDCS-33900078 · EDB16941}
\label{not:234}
\nf{Datation}{382 — directe : nom consulaire lu}
\nf{Genre}{Épitaphe ; ensembles A, B, C ; 73 lettres transmises ; 71 \% restituées.}
\nf{Références}{ICUR II 4167.}
\nf{Provenance}{Rattachement selon l’EDB : \emph{Basilica S.Petri apostoli}.}
\nf{Support}{Plaque de marbre. Inscription gravée. Reproductions signalées par l’EDB : reproduction, Perret 1851 (droits réservés ; non reproduites). Fiche EDB : Antonio Enrico Felle, première version 2003.}
\ntext{\lnum{1}]o qui ⸢v⸣[ixit - - -]\\
\lnum{2}[- - -] in pace [depositus - - -]\\
\lnum{3}[- - - S]ept(embres) [- - -]\\
\lnum{4}[- - - Antonio e]t Syagri[o vv(iris) cc(larissimis) conss(ulibus)}
\nf{Témoins}{Export EDCS de 2020 : \emph{]o qui [ixit 3] / [3] in pace [depositus 3] / [3 S]ept(embres) [3] / [3 Antonio e]t Syagri[o vv(iris) cc(larissimis) conss(ulibus)}. EDB (EDB16941 ; écarts à examiner) : \emph{[---] / 1[---]o1 qui b[ixit ---] / [---] in pace [depositus ---] /[--- s]ept(embres) [---] / [--- Antonio e]t Syagri[o vv(iris) cc(larissimis) conss(ulibus) ---?]}.}
\nf{Apparat}{⸢v⸣ : B gravé}
\nf{Traduction}{[…] qui vécut […] ; […] déposé en paix […] ; […] septembre […] ; sous le consulat d’Antonius et de Syagrius, clarissimes.}
\nf{Langue}{L’éditeur signale une graphie corrigée : B gravé pour v (bêtacisme).}
\nf{Chronologie}{Noms consulaires : Syagrius (lu) et Antonius (restitué).}
\nf{Commentaire}{La formule in pace est conservée.}
\nf{Bibliographie}{ICUR II 4167.}
\nf{Statut}{Lecture secondaire (EDCS) ; écarts avec l’EDB à arbitrer sur l’ICVR. Plus de la moitié des lettres sont restituées. L’EDCS signale une photographie (non consultée).}
\end{notice}

\begin{notice}{235}{382}{ICUR n.s. II, 5791 · EDCS-38302493 · EDB22056}
\label{not:235}
\nf{Datation}{382 — directe : nom consulaire lu}
\nf{Genre}{Épitaphe ; ensembles A, B, C ; 87 lettres transmises.}
\nf{Références}{ICUR II 5791 = ILCV 4146c.}
\nf{Provenance}{Rattachement selon l’EDB : \emph{Area Lucinae}. Conservation selon l’EDB : \emph{Basilica s.Pauli apostoli}.}
\nf{Support}{Plaque de marbre. 200 × 74 cm (h. × l.). Inscription gravée. Fiche EDB : Anita Rocco, première version 2006.}
\ntext{\lnum{1}Serbulus se ⸢v⸣i⸢v⸣u(s) fecit\\
\lnum{2}prid(ie) Non(as) Decemb(res) d(e)p(ositus) Pascasius qui vixit ann(os) IIII m(enses) VIII d(ies) IIII\\
\lnum{3}Antonio et Syacrio cc(on)s(ulibus)}
\nf{Témoins}{Export EDCS de 2020 : \emph{Serbulus se iu(s) fecit / prid(ie) Non(as) Decemb(res) d(e)p(ositus) Pascasius qui vixit ann(os) IIII m(enses) VIII d(ies) IIII / Antonio et Syacrio cc(on)s(ulibus)}. EDB (EDB22056) : texte concordant aux conventions de notation près (indice de similarité 0,938).}
\nf{Apparat}{⸢v⸣ : B gravé ; ⸢v⸣ : B gravé}
\nf{Traduction}{Servulus a fait (ce monument) de son vivant ; la veille des nones de décembre a été déposé Paschasius, qui vécut 4 ans 8 mois 4 jours ; sous le consulat d’Antonius et de Syagrius.}
\nf{Langue}{L’éditeur signale 2 graphies corrigées : B gravé pour v (bêtacisme) ×2.}
\nf{Chronologie}{Noms consulaires : Antonius (lu). Dates du jour : événement indéterminé : pridie Nonas Decembres = 4 décembre 382 (lu). Âge déclaré : 4 a. 8 m. 4 j. (lu).}
\nf{Bibliographie}{ICUR II 5791 ; ILCV 4146c.}
\nf{Statut}{Lecture secondaire (EDCS) concordante avec l’EDB : lecture stabilisée entre deux bases, non collationnée sur l’ICVR ni sur le monument. L’EDCS signale une photographie (non consultée).}
\end{notice}

\begin{notice}{236}{382}{ICUR n.s. II, 5995 · EDCS-38100823 · EDB20813}
\label{not:236}
\nf{Datation}{382 — directe : nom consulaire lu}
\nf{Genre}{Épitaphe ; ensembles A, B, C ; 124 lettres transmises ; 68 \% restituées.}
\nf{Références}{ICUR II 5995.}
\nf{Provenance}{Rattachement selon l’EDB : \emph{Coem. s.Theclae}. Conservation selon l’EDB : \emph{Collegio Teutonico}.}
\nf{Support}{Plaque de marbre. 35 × 37 cm (h. × l.). Inscription gravée. Fiche EDB : Maria Cristina Grisanzio, première version 2006.}
\ntext{\lnum{1}[- - - vixit] ann(os) XII men[s(e)s - - -]\\
\lnum{2}[- - - Cl(audio) Antonio et] Fl(avio) Syagrio v[v(iris) cc(larissimis) cons(ulibu)s]\\
\lnum{3}[- - - vixit a]nn(os) XI mens(e)s I[- - -]\\
\lnum{4}[- - - FFll(avis) Saturnino et] Merobaude c[ons(ulibus)}
\nf{Témoins}{EDB (EDB20813) : texte concordant aux conventions de notation près (indice de similarité 0,939).}
\nf{Traduction}{[…] vécut 12 ans […] mois ; […] sous le consulat de Claudius Antonius et de Flavius Syagrius, clarissimes ; […] vécut 11 ans 1 mois […] ; […] sous le consulat des Flavii Saturninus et Merobaudes.}
\nf{Langue}{11 abréviations développées par l’éditeur (dont ann(os), mens(e)s, Cl(audio), Fl(avio), vv(iris), cc(larissimis)).}
\nf{Chronologie}{Noms consulaires : Merobaudes (lu), Syagrius (lu), Antonius (restitué) et Saturninus (restitué). Âges déclarés : 12 a. (partiellement restitué) ; 11 a. (partiellement restitué).}
\nf{Bibliographie}{ICUR II 5995.}
\nf{Statut}{Lecture secondaire (EDCS) concordante avec l’EDB : lecture stabilisée entre deux bases, non collationnée sur l’ICVR ni sur le monument. Plus de la moitié des lettres sont restituées.}
\end{notice}

\begin{notice}{237}{382}{ICUR n.s. III, 6504 · EDCS-34801304 · EDB22062}
\label{not:237}
\nf{Datation}{382 — directe : lecture partielle d’un nom consulaire}
\nf{Genre}{Épitaphe ; ensembles A, B, C ; 102 lettres transmises ; 18 \% restituées.}
\nf{Références}{ICUR III 6504 = ILCV +4333a.}
\nf{Provenance}{Rattachement selon l’EDB : \emph{Coem. Domitillae pars inferior}. Conservation selon l’EDB : \emph{Coem. Domitillae pars inferior}.}
\nf{Support}{Plaque de marbre. 59 × 150 × 5 cm (h. × l. × ép.). Inscription gravée. Reproductions signalées par l’EDB : photographie PCAS (droits réservés ; non reproduites). Fiche EDB : Carlo Carletti / Antonio Enrico Felle, première version 2003.}
\ntext{\lnum{1}[Depositi]o Laurenti XV Kalendas\\
\lnum{2}[- - - qui] vixit annos sexaginta sex\\
\lnum{3}[- - - bene] merenti in pace totius\\
\lnum{4}innocentiae [A]ntonio et [Sy]agrio}
\nf{Témoins}{EDB (EDB22062) : texte concordant aux conventions de notation près (indice de similarité 0,939).}
\nf{Traduction}{Déposition de Laurentius le 15e jour avant les calendes […], qui vécut soixante-six ans ; bien méritant, en paix, d’une entière innocence ; sous (le consulat d’)Antonius et Syagrius.}
\nf{Langue}{Aucune particularité graphique n’est signalée par l’éditeur dans l’état 2022 de la base.}
\nf{Chronologie}{Noms consulaires : Syagrius (lu en partie) et Antonius (lu en partie). Âge déclaré : 66 a. (partiellement restitué).}
\nf{Commentaire}{La formule in pace est conservée.}
\nf{Bibliographie}{ICUR III 6504 ; ILCV +4333a.}
\nf{Statut}{Lecture secondaire (EDCS) concordante avec l’EDB : lecture stabilisée entre deux bases, non collationnée sur l’ICVR ni sur le monument. L’EDCS signale une photographie (non consultée).}
\end{notice}

\begin{notice}{238}{382}{ICUR n.s. III, 8150 · EDCS-35700179 · EDB24673}
\label{not:238}
\nf{Datation}{382 — directe : lecture partielle d’un nom consulaire}
\nf{Genre}{Épitaphe ; ensembles A, B, C ; 38 lettres transmises ; 74 \% restituées.}
\nf{Références}{ICUR III 8150.}
\nf{Provenance}{Rattachement selon l’EDB : \emph{Basilica ss. Nerei et Achillei in coem. Domitillae}. Conservation selon l’EDB : \emph{Basilica ss. Nerei et Achillei in coem. Domitillae}.}
\nf{Support}{Plaque de marbre. 9 × 14 × 4 cm (h. × l. × ép.). Inscription gravée. Fiche EDB : Carlo Carletti / Antonio Enrico Felle, première version 2003.}
\ntext{\lnum{1}ann(os)] XV m[ens(es) - - -]\\
\lnum{2}[ Antonio e]t S⸢y⸣a[grio conss(ulibus)}
\nf{Témoins}{Export EDCS de 2020 : \emph{ann(os)] XV m[ens(es) 3] / [ Antonio e]t Sa[grio conss(ulibus)}. EDB (EDB24673 ; écarts à examiner) : \emph{------ / [--- annos] XV m[enses ---] / [Antonio e]t Sia[grio cons(ulibu)s] / ------? [texte tronqué dans la liste de résultats]}.}
\nf{Apparat}{⸢y⸣ : I gravé}
\nf{Traduction}{[…] 15 ans […] mois ; sous le consulat d’Antonius et de Syagrius.}
\nf{Langue}{L’éditeur signale une graphie corrigée : I gravé pour y (graphie de y).}
\nf{Chronologie}{Noms consulaires : Syagrius (lu en partie) et Antonius (restitué).}
\nf{Bibliographie}{ICUR III 8150.}
\nf{Statut}{Lecture secondaire (EDCS) ; écarts avec l’EDB à arbitrer sur l’ICVR. Plus de la moitié des lettres sont restituées. L’EDCS signale une photographie (non consultée).}
\end{notice}

\begin{notice}{239}{382}{ICUR n.s. III, 8151 · EDCS-35700180 · EDB24674}
\label{not:239}
\nf{Datation}{382 — directe : lecture partielle d’un nom consulaire}
\nf{Genre}{Épitaphe ; ensembles A, B, C ; 34 lettres transmises ; 59 \% restituées.}
\nf{Références}{ICUR III 8151.}
\nf{Provenance}{Rattachement selon l’EDB : \emph{Basilica ss. Nerei et Achillei in coem. Domitillae}.}
\nf{Support}{Plaque de marbre. 15 × 28 cm (h. × l.). Inscription gravée. Fiche EDB : Carlo Carletti / Antonio Enrico Felle, première version 2003.}
\ntext{\lnum{1}]vit(?) [- - -]\\
\lnum{2}[- - - Anton]io et S⸢y⸣a[grio conss(ulibus) - - -]\\
\lnum{3}[- - -]tici(?) [}
\nf{Témoins}{Export EDCS de 2020 : \emph{]vit(?) [3] / [3 Anton]io et Sa[grio conss(ulibus) 3] / [3]tici(?) [}. EDB (EDB24674 ; écarts à examiner) : \emph{[---]vit [---] / [--- Anton]io et Sua[grio cons(ulibu)s] / [---]tici[---] [texte tronqué dans la liste de résultats]}.}
\nf{Apparat}{⸢y⸣ : V gravé}
\nf{Traduction}{[…] sous le consulat d’Antonius et de Syagrius […].}
\nf{Langue}{L’éditeur signale une graphie corrigée : V gravé pour y (graphie de y).}
\nf{Chronologie}{Noms consulaires : Syagrius (lu en partie) et Antonius (restitué).}
\nf{Bibliographie}{ICUR III 8151.}
\nf{Statut}{Lecture secondaire (EDCS) ; écarts avec l’EDB à arbitrer sur l’ICVR. Plus de la moitié des lettres sont restituées. L’EDCS signale une photographie (non consultée).}
\end{notice}

\begin{notice}{240}{382}{ICUR n.s. IV, 12259a · EDCS-35700730 · EDB4650}
\label{not:240}
\nf{Datation}{382 — directe : lecture partielle d’un nom consulaire}
\nf{Genre}{Épitaphe ; ensembles A, B, C ; 29 lettres transmises ; 31 \% restituées.}
\nf{Références}{ICUR IV 12259a.}
\nf{Provenance}{Rattachement selon l’EDB : \emph{Coem. anonymum ad viam Ardeatinam}. Conservation selon l’EDB : \emph{Coem. anonymum ad viam Ardeatinam}.}
\nf{Support}{Plaque de marbre. Inscription gravée. Fiche EDB : Marida Pierno, première version 2003.}
\ntext{\lnum{1}o] et Sya[grio - - -]\\
\lnum{2}[- - - depo]sitio\\
\lnum{3}[\\
\lnum{4}n]isi [- - -]\\
\lnum{5}[- - -] in pace [}
\nf{Témoins}{EDB (EDB4650 ; écarts à examiner) : \emph{[-- o] et Sya[grio ---] / [---]sitio / [------]/ [-- n]isi [---]/ [---] in pace}.}
\nf{Traduction}{[…] (Antoni)us et Syagrius […] ; déposition […] // si ce n’est […] ; […] en paix.}
\nf{Chronologie}{Noms consulaires : Syagrius (lu en partie).}
\nf{Commentaire}{La formule in pace est conservée.}
\nf{Bibliographie}{ICUR IV 12259a.}
\nf{Statut}{Lecture secondaire (EDCS) ; écarts avec l’EDB à arbitrer sur l’ICVR.}
\end{notice}

\annee{383}

\begin{notice}{241}{383}{EDCS-13700250 · EDB42969}
\label{not:241}
\nf{Datation}{383 — directe : lecture partielle d’un nom consulaire}
\nf{Genre}{Épitaphe ; ensembles A, B, C ; 192 lettres transmises ; 22 \% restituées.}
\nf{Références}{AE 1999, 307.}
\nf{Provenance}{Rattachement selon l’EDB : \emph{Incertae originis} (le volume I des ICVR réunit les inscriptions dont la provenance est perdue). Conservation selon l’EDB : \emph{s. Giovanni in Laterano}.}
\nf{Support}{La fiche EDB ne figure pas dans l’état 2019 des données EpiDoc ; support non documenté.}
\ntext{\lnum{1}]itur ostend[- - -]\\
\lnum{2}[- - -] somno replet[- - -]\\
\lnum{3}[- - - h]ortabatur ad[- - -]\\
\lnum{4}[- - -] obsequiis sem[- - -]\\
\lnum{5}[- - -]it pater et Uran[ia(?) - - -]\\
\lnum{6}[- - -]or teneros inn[- - -]\\
\lnum{7}[- - -] domum quem [- - -]\\
\lnum{8}[- - -]uas quabus(!) aet[- - -]\\
\lnum{9}[- - -] desolavit et [- - -]\\
\lnum{10}[- - -] in pace perenn[- - -]\\
\lnum{11}[- - -]rus scripsit ma[- - -]\\
\lnum{12}[- - -] deposita in p[ace - - -]\\
\lnum{13}[Fl(avio) Mer]obaude iterum [et Fl(avio) Saturnino conss(ulibus)]}
\nf{Témoins}{EDB (EDB42969) : texte concordant aux conventions de notation près (indice de similarité 0,98).}
\nf{Apparat}{(!) : forme signalée par l’éditeur}
\nf{Traduction}{(vers, fragmentaire) […] est montré […] ; […] rempli de sommeil […] ; […] il exhortait à […] ; […] par des services […] ; […] son père et Urania […] ; […] tendres […] ; […] la maison que […] ; […] a désolé et […] ; […] dans la paix perpétuelle […] ; […] a écrit […] ; déposée en paix […] ; sous le consulat de Flavius Merobaudes pour la 2e fois et de Flavius Saturninus.}
\nf{Langue}{L’éditeur signale au moins une forme anormale (sic). 5 abréviations développées par l’éditeur (dont Urania(?), quabus(!), Fl(avio), conss(ulibus)).}
\nf{Chronologie}{Noms consulaires : Merobaudes (lu en partie) et Saturninus (restitué).}
\nf{Commentaire}{La formule in pace est conservée.}
\nf{Bibliographie}{AE 1999, 307. Bibliographie complémentaire selon l’EDB : AE, 1999, 307 ; Liverani 1998, 93, n. 448 (G. Spinola).}
\nf{Statut}{Lecture secondaire (EDCS) concordante avec l’EDB : lecture stabilisée entre deux bases, non collationnée sur l’ICVR ni sur le monument.}
\end{notice}

\begin{notice}{242}{383}{ICUR n.s. I, 504 · EDCS-38700907 · EDB28484}
\label{not:242}
\nf{Datation}{383 — directe : lecture partielle d’un nom consulaire}
\nf{Genre}{Épitaphe ; ensembles A, B, C ; 47 lettres transmises ; 51 \% restituées.}
\nf{Références}{ICUR I 504.}
\nf{Provenance}{Rattachement selon l’EDB : \emph{Incertae originis} (le volume I des ICVR réunit les inscriptions dont la provenance est perdue).}
\nf{Support}{Support non précisé par l’EDB. Fiche EDB : Filippo Piazzolla, première version 2014.}
\ntext{\lnum{1}c]oniugii\\
\lnum{2}[- - -]XI dep(ositus ?) pr(idie) Id(us)\\
\lnum{3}[- - - cons(ulatu) Merobaudis et Sat]urnini}
\nf{Témoins}{EDB (EDB28484 ; écarts à examiner) : \emph{[ --- c]oniugii / [---]XI dep(osit---) pr(idie) id(us) / [--- post? consulatu Merobaudis et Sat]urnini [texte tronqué dans la liste de résultats]}.}
\nf{Traduction}{[…] du mariage […] ; […] 11 […] ; déposé la veille des ides […] ; […] sous le consulat de Merobaudes et de Saturninus.}
\nf{Langue}{4 abréviations développées par l’éditeur (dont dep(ositus ?), pr(idie), Id(us), cons(ulatu)).}
\nf{Chronologie}{Noms consulaires : Merobaudes (restitué) et Saturninus (lu en partie).}
\nf{Bibliographie}{ICUR I 504.}
\nf{Statut}{Lecture secondaire (EDCS) ; écarts avec l’EDB à arbitrer sur l’ICVR. Plus de la moitié des lettres sont restituées.}
\end{notice}

\begin{notice}{243}{383}{ICUR n.s. I, 1169 · EDCS-38700710 · EDB30112}
\label{not:243}
\nf{Datation}{383 — directe : lecture partielle d’un nom consulaire}
\nf{Genre}{Épitaphe ; ensembles A, B, C ; 37 lettres transmises.}
\nf{Références}{ICUR I 1169.}
\nf{Provenance}{Rattachement selon l’EDB : \emph{Incertae originis} (le volume I des ICVR réunit les inscriptions dont la provenance est perdue).}
\nf{Support}{Support non précisé par l’EDB. Fiche EDB : Filippo Piazzolla, première version 2014.}
\ntext{\lnum{1}Merobaude et Satur]nino vv(iris) cc(larissimis) conss(ulibus) [- - -]\\
\lnum{2}[- - -] deposita [}
\nf{Témoins}{EDB (EDB30112) : texte concordant aux conventions de notation près (indice de similarité 1).}
\nf{Traduction}{Sous le consulat de Merobaudes et de Saturninus, clarissimes […] ; […] déposée.}
\nf{Langue}{3 abréviations développées par l’éditeur (dont vv(iris), cc(larissimis), conss(ulibus)).}
\nf{Chronologie}{Noms consulaires : Merobaudes (restitué) et Saturninus (lu en partie).}
\nf{Bibliographie}{ICUR I 1169.}
\nf{Statut}{Lecture secondaire (EDCS) concordante avec l’EDB : lecture stabilisée entre deux bases, non collationnée sur l’ICVR ni sur le monument.}
\end{notice}

\begin{notice}{244}{383}{ICUR n.s. I, 1440 · EDCS-34501146 · EDB31205}
\label{not:244}
\nf{Datation}{383 — directe : nom consulaire lu}
\nf{Genre}{Épitaphe ; ensembles A, B, C ; 394 lettres transmises.}
\nf{Références}{ICUR I 1440 = ILCV 3446 (add) = CLE 670 = CLENuovo p 108 = ICVaticano p 284 = Damaso p 110 = Veleia-2012-31 = Trout-01, 00051 = AE 1997, 166.}
\nf{Provenance}{Rattachement selon l’EDB : \emph{Incertae originis} (le volume I des ICVR réunit les inscriptions dont la provenance est perdue). Conservation selon l’EDB : \emph{Musei Vaticani: Lapidario Cristiano}.}
\nf{Support}{Plaque de marbre. 56 × 198 cm (h. × l.). Inscription gravée. Texte métrique. Fiche EDB : Filippo Piazzolla, première version 2015.}
\ntext{\lnum{1}Quid loquar aut sileam prohibet dolor ipse fateri\\
\lnum{2}hic tumulus lacrimas retinet cognosce parentu(m)\\
\lnum{3}Proiectae fuerat Primo quae iuncta marito\\
\lnum{4}pulc(h)ra decore suo solo contenta pudore\\
\lnum{5}heu dilecta satis miserae genetricis amore\\
\lnum{6}accipe quid multis thalami post foedera prima\\
\lnum{7}erepta ex oculis Flori genitoris abiit\\
\lnum{8}aetheriam cupiens caeli conscendere lucem\\
\lnum{9}haes Damasus pr(a)estat cunctis solacia fletus\\
\lnum{10}vixit ann(os) XVI « m(enses) X » dies XXV dep(osita) III Kal(endas) Ian(uarias) Fl(avio) Merobaude\\
\lnum{11}et Fl(avio) Saturnin(o) cons(ulibus)}
\nf{Témoins}{EDB (EDB31205 ; écarts à examiner) : \emph{quid? loquar aut sileam? prohibet dolor ipse fateri / hic tumulus lacrimas retinet cognosce parentu̅(m) / Proiectae fuerat primo quae iuncta marito / pulcra decore suo solo contenta pudore / heu dilecta satis miserae genitricis amore / accipe quid multis thalami post foedera prima / erepta ex oculis flori genitoris abiit / aetheriam cupiens caeli conscendere lucem / haec Damasus prestat cunctis solacia fletus / vixit ann(os) XVI 《m(enses) X》 dies XXV dep(osita) III kal(endas) ian(uarias) Fl(avio) Merobaude et Fl(avio) Saturnin⟨o⟩ cons(ulibu)s}.}
\nf{Traduction}{(vers, épigramme de Damase) Que dire ? ou me taire ? La douleur même m’interdit de parler. Ce tombeau retient les larmes des parents — sache-le — de Proiecta, qui avait été unie à Primus, son mari ; belle de sa grâce, contente de sa seule pudeur. Hélas, si chérie de l’amour d’une mère malheureuse ! Apprends ce que (elle eut) de plus que beaucoup : après les premiers liens du mariage, arrachée aux yeux de Florus, son père, elle s’en alla, désireuse de monter à la lumière éthérée du ciel. Damase offre à tous cette consolation des pleurs. Elle vécut 16 ans 10 mois 25 jours ; déposée le 3e jour avant les calendes de janvier, sous le consulat de Flavius Merobaudes et de Flavius Saturninus.}
\nf{Langue}{12 abréviations développées par l’éditeur (dont parentu(m), pulc(h)ra, pr(a)estat, ann(os), m(enses), dep(osita)).}
\nf{Chronologie}{Noms consulaires : Merobaudes (lu) et Saturninus (lu). Dates du jour : déposition : III Kalendas Ianuarias = 30 décembre 383 (lu). Âge déclaré : 16 a. 10 m. 25 j. (lu).}
\nf{Commentaire}{Vocabulaire religieux : allusion interprétative (niveau 5) — aetheriam (lu), caeli (lu). Commentaire développé : dossier p.~\pageref{dos:244}.}
\nf{Bibliographie}{ICUR I 1440 ; ILCV 3446 (add) ; CLE 670 ; CLENuovo p 108 ; ICVaticano p 284 ; Damaso p 110 ; Veleia-2012-31 ; Trout-01, 00051 ; AE 1997, 166. Bibliographie complémentaire selon l’EDB : ILCV, 3446 ; Ferrua 1942, 201-205, n.51 con foto ; Ferrua 1964, 108 ; Ferrua 1979b, 37, n. 1441 ; Di Stefano Manzella 1997, 284-285, 3.6.5, con foto (G. De Felice) ; Lega 2001 ; Trout 2015, 180-182, n. 51 ; Mazzoleni 2015, 165 ; Sághy 2015, 48-49.}
\nf{Statut}{Lecture secondaire (EDCS) ; écarts avec l’EDB à arbitrer sur l’ICVR. L’EDCS signale une photographie (non consultée).}
\end{notice}

\begin{notice}{245}{383}{ICUR n.s. I, 2097 · EDCS-39300318 · EDB32082}
\label{not:245}
\nf{Datation}{383 — directe : nom consulaire lu}
\nf{Genre}{Épitaphe ; ensembles A, B, C ; 61 lettres transmises ; 39 \% restituées.}
\nf{Références}{ICUR I 2097.}
\nf{Provenance}{Rattachement selon l’EDB : \emph{Incertae originis} (le volume I des ICVR réunit les inscriptions dont la provenance est perdue). Conservation selon l’EDB : \emph{Musei Vaticani: Galleria Lapidaria}.}
\nf{Support}{Plaque de marbre. 55 × 68 cm (h. × l.). Inscription gravée. Fiche EDB : Filippo Piazzolla, première version 2015.}
\ntext{\lnum{1}[Merobaude et] Saturnino\\
\lnum{2}[conss(ulibus) - - -] Septembres\\
\lnum{3}[- - - Ex]uperantius qui vixi⸢t⸣\\
\lnum{4}[}
\nf{Témoins}{Export EDCS de 2020 : \emph{[Merobaude et] Saturnino / [conss(ulibus) 3] Septembres / [3 Ex]uperantius qui vixi / [}. EDB (EDB32082) : texte concordant aux conventions de notation près (indice de similarité 0,969).}
\nf{Apparat}{⸢t⸣ : I gravé}
\nf{Traduction}{Sous le consulat de Merobaudes et de Saturninus […] septembre […] Exuperantius, qui vécut […].}
\nf{Langue}{L’éditeur signale une graphie corrigée : I gravé pour t.}
\nf{Chronologie}{Noms consulaires : Merobaudes (restitué) et Saturninus (lu).}
\nf{Bibliographie}{ICUR I 2097. Bibliographie complémentaire selon l’EDB : Lega 2000 ; Ferrua 1979b, 63.}
\nf{Statut}{Lecture secondaire (EDCS) concordante avec l’EDB : lecture stabilisée entre deux bases, non collationnée sur l’ICVR ni sur le monument.}
\end{notice}

\begin{notice}{246}{383}{ICUR n.s. I, 2098 · EDCS-30600273 · EDB35608}
\label{not:246}
\nf{Datation}{383 — directe : nom consulaire lu}
\nf{Genre}{Épitaphe ; ensembles A, B, C ; 62 lettres transmises ; 37 \% restituées.}
\nf{Références}{ICUR I 2098 (p 491) = ICUR X 26644 (p 244).}
\nf{Provenance}{Rattachement selon l’EDB : \emph{Area subdialis coem. Pamphili et inscriptiones incerto loco repertae}. Conservation selon l’EDB : \emph{Musei Vaticani: Galleria Lapidaria}.}
\nf{Support}{Plaque de marbre. 39 × 72 cm (h. × l.). Inscription gravée. Fiche EDB : Filippo Piazzolla, première version 2015.}
\ntext{\lnum{1}[Merobaude II et] Saturnino conss(ulibus)\\
\lnum{2}[deposita - - -]donia annorum XI\\
\lnum{3}[in] pace Pelagia}
\nf{Témoins}{EDB (EDB35608) : texte concordant aux conventions de notation près (indice de similarité 0,961).}
\nf{Traduction}{Sous le consulat de Merobaudes pour la 2e fois et de Saturninus, déposée […]donia, âgée de 11 ans, en paix ; Pelagia.}
\nf{Langue}{1 abréviation développée par l’éditeur (dont conss(ulibus)).}
\nf{Chronologie}{Noms consulaires : Merobaudes (restitué) et Saturninus (lu). Âge déclaré : 11 a. (partiellement restitué).}
\nf{Commentaire}{La formule in pace est attestée seulement par restitution.}
\nf{Bibliographie}{ICUR I 2098 (p 491) ; ICUR X 26644 (p 244). Bibliographie complémentaire selon l’EDB : Lega 2000, 268 ; ICVR X, 26644.}
\nf{Statut}{Lecture secondaire (EDCS) concordante avec l’EDB : lecture stabilisée entre deux bases, non collationnée sur l’ICVR ni sur le monument.}
\end{notice}

\begin{notice}{247}{383}{ICUR n.s. I, 2944 · EDCS-39900394 · EDB32852}
\label{not:247}
\nf{Datation}{383 — directe : nom consulaire lu}
\nf{Genre}{Épitaphe ; ensembles A, B, C ; 87 lettres transmises ; 24 \% restituées.}
\nf{Références}{ICUR I 2944 = ILCV 4144b.}
\nf{Provenance}{Rattachement selon l’EDB : \emph{Incertae originis} (le volume I des ICVR réunit les inscriptions dont la provenance est perdue).}
\nf{Support}{Support non précisé par l’EDB. Signes décrits par l’EDB : chrisme. Fiche EDB : Filippo Piazzolla, première version 2015.}
\ntext{\lnum{1}Ioviano karissimo\\
\lnum{2}fecit Lampadius et Soteres fra\\
\lnum{3}tres pientissimi Merobaude\\
\lnum{4}[et Saturnino cons(ulibus)}
\nf{Témoins}{EDB (EDB32852) : texte concordant aux conventions de notation près (indice de similarité 0,939).}
\nf{Traduction}{À Iovianus très cher, Lampadius et Soterès, frères très pieux, ont fait (ce monument) ; sous le consulat de Merobaudes et de Saturninus.}
\nf{Langue}{1 abréviation développée par l’éditeur (dont cons(ulibus)).}
\nf{Chronologie}{Noms consulaires : Merobaudes (lu) et Saturninus (restitué). Ne sont pas des consuls : Ioviano karissimo (défunt).}
\nf{Bibliographie}{ICUR I 2944 ; ILCV 4144b. Bibliographie complémentaire selon l’EDB : ILCV, 4144B.}
\nf{Statut}{Lecture secondaire (EDCS) concordante avec l’EDB : lecture stabilisée entre deux bases, non collationnée sur l’ICVR ni sur le monument.}
\end{notice}

\begin{notice}{248}{383}{ICUR n.s. I, 3043 · EDCS-39900484 · EDB32975}
\label{not:248}
\nf{Datation}{383 — directe : lecture partielle d’un nom consulaire}
\nf{Genre}{Épitaphe ; ensembles A, B, C ; 30 lettres transmises ; 63 \% restituées.}
\nf{Références}{ICUR I 3043.}
\nf{Provenance}{Rattachement selon l’EDB : \emph{Incertae originis} (le volume I des ICVR réunit les inscriptions dont la provenance est perdue). Conservation selon l’EDB : \emph{Tarquinia (Roma)   ; s. Maria di Castello}.}
\nf{Support}{Plaque de marbre. 4 × 20 cm (h. × l.). Inscription gravée. Fiche EDB : Filippo Piazzolla, première version 2015.}
\ntext{\lnum{1}Merob]aude et [Saturnino cons(ulibus)}
\nf{Témoins}{EDB (EDB32975 ; écarts à examiner) : \emph{[--- Merob]aude et [Saturnino cons(ulibu)s] / ------ [texte tronqué dans la liste de résultats]}.}
\nf{Traduction}{Sous le consulat de Merobaudes et de Saturninus.}
\nf{Langue}{1 abréviation développée par l’éditeur (dont cons(ulibus)).}
\nf{Chronologie}{Noms consulaires : Merobaudes (lu en partie) et Saturninus (restitué).}
\nf{Bibliographie}{ICUR I 3043.}
\nf{Statut}{Lecture secondaire (EDCS) ; écarts avec l’EDB à arbitrer sur l’ICVR. Plus de la moitié des lettres sont restituées.}
\end{notice}

\begin{notice}{249}{383}{ICUR n.s. I, 3198 · EDCS-33500350 · EDB35683}
\label{not:249}
\nf{Datation}{383 — directe : nom consulaire lu}
\nf{Genre}{Épitaphe ; ensembles A, B, C ; 63 lettres transmises ; 48 \% restituées.}
\nf{Références}{ICUR I 3198 = ICUR VIII 20717 = ILCV 3390a.}
\nf{Provenance}{Rattachement selon l’EDB : \emph{Coem. Nicomedis}. Conservation selon l’EDB : \emph{Musei Capitolini}.}
\nf{Support}{Plaque de marbre. 20 × 23,5 × 8,5 cm (h. × l. × ép.). Inscription gravée. Fiche EDB : Filippo Piazzolla, première version 2015.}
\ntext{\lnum{1}]\\
\lnum{2}tus est qui ac[cepti sanc]\\
\lnum{3}tum spiritum [XVI Kal(endas)]\\
\lnum{4}Mai(as) Merobaud[e II con(sule)]}
\nf{Témoins}{EDB (EDB35683) : texte concordant aux conventions de notation près (indice de similarité 0,934).}
\nf{Traduction}{[…] qui aurait reçu l’[Esprit saint] (restitution et syntaxe incertaines) ; [16e jour avant les calendes] de mai ; Merobaud[e II…]. Latin ac[cepti…] maintenu, sans correction en accepit.}
\nf{Langue}{3 abréviations développées par l’éditeur (dont Kal(endas), Mai(as), con(sule)).}
\nf{Chronologie}{Noms consulaires : Merobaudes (lu). Dates du jour : événement indéterminé : XVI Kalendas Maias = 16 avril 383 (restitué).}
\nf{Commentaire}{Spiritum est transmis ; sanctum est partiellement restitué. L’identification de l’Esprit saint dépend de cette restitution. Syntaxe ac[cepti…] non arbitrée ; exclusion du scénario restrictif. Le mot \emph{spiritum} est lu ; une partie de \emph{sanctum} est restituée. L'identification de l'Esprit saint dépend de cette restitution : le texte est retenu dans le profil large, exclu du profil restrictif. Désignation divine : retenue dans le profil large seulement. Commentaire développé : dossier p.~\pageref{dos:249}.}
\nf{Bibliographie}{ICUR I 3198 ; ICUR VIII 20717 ; ILCV 3390a. Bibliographie complémentaire selon l’EDB : ICVR I, 3198.}
\nf{Statut}{Lecture secondaire (EDCS) concordante avec l’EDB : lecture stabilisée entre deux bases, non collationnée sur l’ICVR ni sur le monument.}
\end{notice}

\begin{notice}{250}{383}{ICUR n.s. II, 6038 · EDCS-16600357 · EDB8750}
\label{not:250}
\nf{Datation}{383 — directe : nom consulaire lu}
\nf{Genre}{Épitaphe ; ensembles A, B, C ; 96 lettres transmises.}
\nf{Références}{ICUR II 6038 = ILCV 2868 = AE 1905, 76.}
\nf{Provenance}{Rattachement selon l’EDB : \emph{Coem. Commodillae}.}
\nf{Support}{Plaque de marbre. Inscription gravée. Fiche EDB : Carlo Carletti, première version 2003.}
\ntext{\lnum{1}Sepulcrum Filonetis secessit in pace III Idus\\
\lnum{2}Septemb(res) annorum plus minus XXVIII Merobau\\
\lnum{3}de bis et Saturnino conss(ulibus)}
\nf{Témoins}{EDB (EDB8750) : texte concordant aux conventions de notation près (indice de similarité 0,996).}
\nf{Traduction}{Sépulture de Philonès ; il/elle s’en est allé(e) en paix le 3e jour avant les ides de septembre, âgé(e) d’environ 28 ans, sous le consulat de Merobaudes pour la 2e fois et de Saturninus.}
\nf{Langue}{2 abréviations développées par l’éditeur (dont Septemb(res), conss(ulibus)).}
\nf{Chronologie}{Noms consulaires : Merobaudes (lu) et Saturninus (lu). Dates du jour : événement indéterminé : III Idus Septembres = 11 septembre 383 (lu).}
\nf{Commentaire}{La formule in pace est conservée.}
\nf{Bibliographie}{ICUR II 6038 ; ILCV 2868 ; AE 1905, 76.}
\nf{Statut}{Lecture secondaire (EDCS) concordante avec l’EDB : lecture stabilisée entre deux bases, non collationnée sur l’ICVR ni sur le monument.}
\end{notice}

\begin{notice}{251}{383}{ICUR n.s. II, 6039 · EDCS-34200665 · EDB8751}
\label{not:251}
\nf{Datation}{383 — directe : nom consulaire lu}
\nf{Genre}{Épitaphe ; ensembles A, B, C ; 67 lettres transmises ; 25 \% restituées.}
\nf{Références}{ICUR II 6039 = AE 1905, 77.}
\nf{Provenance}{Rattachement selon l’EDB : \emph{Coem. Commodillae}.}
\nf{Support}{Plaque de marbre. Inscription gravée. Reproductions signalées par l’EDB : photographie PCAS (droits réservés ; non reproduites). Fiche EDB : Carlo Carletti, première version 2003.}
\ntext{\lnum{1}qu]i vixit annos X\\
\lnum{2}[mens(es)] II dep(ositus ?) di\{a\}e VI Kal(endas)\\
\lnum{3}[- - -] Fl(avio) Merobaude II et Sat[urni]\\
\lnum{4}no con[s(ulibu)s]}
\nf{Témoins}{EDB (EDB8751) : texte concordant aux conventions de notation près (indice de similarité 0,968).}
\nf{Apparat}{\{a\} : lettres superflues}
\nf{Traduction}{[…] qui vécut 10 ans 2 mois ; déposé le 6e jour avant les calendes […] ; sous le consulat de Flavius Merobaudes pour la 2e fois et de Saturninus.}
\nf{Langue}{Des lettres gravées sont jugées superflues par l’éditeur (entre accolades). 5 abréviations développées par l’éditeur (dont mens(es), dep(ositus ?), Kal(endas), Fl(avio), cons(ulibu)s).}
\nf{Chronologie}{Noms consulaires : Merobaudes (lu) et Saturninus (lu en partie). Âge déclaré : 10 a. 2 m. 6 j. (partiellement restitué).}
\nf{Bibliographie}{ICUR II 6039 ; AE 1905, 77.}
\nf{Statut}{Lecture secondaire (EDCS) concordante avec l’EDB : lecture stabilisée entre deux bases, non collationnée sur l’ICVR ni sur le monument. L’EDCS signale une photographie (non consultée).}
\end{notice}

\begin{notice}{252}{383}{ICUR n.s. III, 7383 · EDCS-34900711 · EDB24131}
\label{not:252}
\nf{Datation}{383 — directe : lecture partielle d’un nom consulaire}
\nf{Genre}{Épitaphe ; ensembles A, B, C ; 68 lettres transmises ; 66 \% restituées.}
\nf{Références}{ICUR III 7383.}
\nf{Provenance}{Rattachement selon l’EDB : \emph{Coem. Domitillae pars superior}. Conservation selon l’EDB : \emph{Coem. Domitillae pars superior}.}
\nf{Support}{Plaque de marbre. 22 × 25 × 2 cm (h. × l. × ép.). Inscription gravée. Fiche EDB : Carlo Carletti / Antonio Enrico Felle, première version 2015.}
\ntext{\lnum{1}vix]it anno[s - - -]\\
\lnum{2}[- - - Se]pt(embres) conss(ulibus) [FFll(avvis)]\\
\lnum{3}[- - - Merobaude II et S]aturnin[o vv(iris) cc(larissimis)]}
\nf{Témoins}{EDB (EDB24131) : texte concordant aux conventions de notation près (indice de similarité 0,955).}
\nf{Traduction}{[…] vécut […] ans ; […] septembre, sous le consulat des Flavii […] Merobaudes pour la 2e fois et Saturninus, clarissimes.}
\nf{Langue}{5 abréviations développées par l’éditeur (dont Sept(embres), conss(ulibus), FFll(avvis), vv(iris), cc(larissimis)).}
\nf{Chronologie}{Noms consulaires : Merobaudes (restitué) et Saturninus (lu en partie).}
\nf{Bibliographie}{ICUR III 7383.}
\nf{Statut}{Lecture secondaire (EDCS) concordante avec l’EDB : lecture stabilisée entre deux bases, non collationnée sur l’ICVR ni sur le monument. Plus de la moitié des lettres sont restituées.}
\end{notice}

\begin{notice}{253}{383}{ICUR n.s. III, 8152 · EDCS-35700181 · EDB24675}
\label{not:253}
\nf{Datation}{383 — directe : lecture partielle d’un nom consulaire}
\nf{Genre}{Épitaphe ; ensembles A, B, C ; 22 lettres transmises.}
\nf{Références}{ICUR III 8152.}
\nf{Provenance}{Rattachement selon l’EDB : \emph{Basilica ss. Nerei et Achillei in coem. Domitillae}. Conservation selon l’EDB : \emph{Basilica ss. Nerei et Achillei in coem. Domitillae}.}
\nf{Support}{Plaque de marbre. 37 × 33 × 3,5 cm (h. × l. × ép.). Inscription gravée. Fiche EDB : Carlo Carletti / Antonio Enrico Felle, première version 2003.}
\ntext{\lnum{1}Merobaude II et S]aturnino}
\nf{Témoins}{EDB (EDB24675) : texte concordant aux conventions de notation près (indice de similarité 1).}
\nf{Traduction}{(Sous le consulat de) Merobaudes pour la 2e fois et Saturninus.}
\nf{Langue}{Aucune particularité graphique n’est signalée par l’éditeur dans l’état 2022 de la base.}
\nf{Chronologie}{Noms consulaires : Merobaudes (restitué) et Saturninus (lu en partie).}
\nf{Bibliographie}{ICUR III 8152.}
\nf{Statut}{Lecture secondaire (EDCS) concordante avec l’EDB : lecture stabilisée entre deux bases, non collationnée sur l’ICVR ni sur le monument. L’EDCS signale une photographie (non consultée).}
\end{notice}

\begin{notice}{254}{383}{ICUR n.s. III, 8425 · EDCS-35802012 · EDB25700}
\label{not:254}
\nf{Datation}{383 — directe : lecture partielle d’un nom consulaire}
\nf{Genre}{Épitaphe ; ensembles A, B, C ; 59 lettres transmises ; 68 \% restituées.}
\nf{Références}{ICUR III 8425 (p 408).}
\nf{Provenance}{Rattachement selon l’EDB : \emph{Incertum locum in coem. Domitillae}.}
\nf{Support}{Plaque de marbre. Inscription gravée. Reproductions signalées par l’EDB : photographie PCAS ×2 (droits réservés ; non reproduites). Fiche EDB : Carlo Carletti / Antonio Enrico Felle, première version 2003.}
\ntext{\lnum{1}] Primia C[- - -]\\
\lnum{2}[- - - fecit c]um marit[o ann(os) - - -]\\
\lnum{3}Merob[aude II et Saturnino conss(ulibus)]}
\nf{Témoins}{EDB (EDB25700) : texte concordant aux conventions de notation près (indice de similarité 0,931).}
\nf{Traduction}{[…] Primia […] a vécu avec son mari […] ans ; sous le consulat de Merobaudes pour la 2e fois et de Saturninus.}
\nf{Langue}{2 abréviations développées par l’éditeur (dont ann(os), conss(ulibus)).}
\nf{Chronologie}{Noms consulaires : Merobaudes (lu en partie) et Saturninus (restitué).}
\nf{Bibliographie}{ICUR III 8425 (p 408).}
\nf{Statut}{Lecture secondaire (EDCS) concordante avec l’EDB : lecture stabilisée entre deux bases, non collationnée sur l’ICVR ni sur le monument. Plus de la moitié des lettres sont restituées. L’EDCS signale une photographie (non consultée).}
\end{notice}

\begin{notice}{255}{383}{ICUR n.s. III, 8728 · EDCS-35900972 · EDB20583}
\label{not:255}
\nf{Datation}{383 — directe : nom consulaire lu}
\nf{Genre}{Épitaphe ; ensembles A, B, C ; 70 lettres transmises.}
\nf{Références}{ICUR III 8728 = ILCV 3057a (add).}
\nf{Provenance}{Rattachement selon l’EDB : \emph{Tituli qui in coem. Callisti reperti traduntur}.}
\nf{Support}{Plaque de marbre. Inscription gravée. Fiche EDB : Carlo Carletti, première version 2003.}
\ntext{\lnum{1}Hic iacet Ticira vixit annis V mensib(us) III\\
\lnum{2}dep(osita) III Kal(endas) Mai(as) Merubaude et\\
\lnum{3}Saturnino cons(ulibus)}
\nf{Témoins}{EDB (EDB20583) : texte concordant aux conventions de notation près (indice de similarité 1).}
\nf{Traduction}{Ici gît Ticira ; elle vécut 5 ans 3 mois ; déposée le 3e jour avant les calendes de mai, sous le consulat de Merobaudes et de Saturninus.}
\nf{Langue}{5 abréviations développées par l’éditeur (dont mensib(us), dep(osita), Kal(endas), Mai(as), cons(ulibus)).}
\nf{Chronologie}{Noms consulaires : Saturninus (lu). Dates du jour : déposition : III Kalendas Maias = 29 avril 383 (lu). Âge déclaré : 5 a. 3 m. (lu).}
\nf{Bibliographie}{ICUR III 8728 ; ILCV 3057a (add).}
\nf{Statut}{Lecture secondaire (EDCS) concordante avec l’EDB : lecture stabilisée entre deux bases, non collationnée sur l’ICVR ni sur le monument.}
\end{notice}

\begin{notice}{256}{383}{ICUR n.s. IV, 11111 · EDCS-37200393 · EDB39573}
\label{not:256}
\nf{Datation}{383 — directe : consulat restitué}
\nf{Genre}{Épitaphe ; ensembles A, B, C ; 42 lettres transmises ; 71 \% restituées.}
\nf{Références}{ICUR IV 11111.}
\nf{Provenance}{Rattachement selon l’EDB : \emph{Coem. Callisti area subdialis et incerti loci}. Conservation selon l’EDB : \emph{Coem. Callisti area subdialis et incerti loci}.}
\nf{Support}{Plaque de marbre. 25 × 18 × 4,5 cm (h. × l. × ép.). Inscription gravée. Fiche EDB : Marida Pierno, première version 2015.}
\ntext{\lnum{1}]RE[- - -]\\
\lnum{2}[- - - men]ses X[- - -]\\
\lnum{3}[- - - Merobaud]e II et S[aturnino conss(ulibus)]}
\nf{Témoins}{EDB (EDB39573 ; écarts à examiner) : \emph{------ / [---]RE[---]/[--- men]ses X[---]/ [Merobaud]e II et S[aturnino consulibus]}.}
\nf{Traduction}{[…] 10 mois […] ; sous le consulat de Merobaudes pour la 2e fois et de Saturninus.}
\nf{Langue}{1 abréviation développée par l’éditeur (dont conss(ulibus)).}
\nf{Chronologie}{Noms consulaires : Merobaudes (restitué) et Saturninus (restitué).}
\nf{Bibliographie}{ICUR IV 11111.}
\nf{Statut}{Lecture secondaire (EDCS) ; écarts avec l’EDB à arbitrer sur l’ICVR. Plus de la moitié des lettres sont restituées.}
\end{notice}

\begin{notice}{257}{383}{ICUR n.s. IV, 12245 · EDCS-37600680 · EDB5084}
\label{not:257}
\nf{Datation}{383 — directe : nom consulaire lu}
\nf{Genre}{Épitaphe ; ensembles A, B, C ; 116 lettres transmises ; 51 \% restituées.}
\nf{Références}{ICUR IV 12245.}
\nf{Provenance}{Rattachement selon l’EDB : \emph{Coem. anonymum ad viam Ardeatinam}. Conservation selon l’EDB : \emph{Coem. anonymum ad viam Ardeatinam}.}
\nf{Support}{Plaque de marbre. Inscription gravée. Reproductions signalées par l’EDB : photographie PCAS (droits réservés ; non reproduites). Fiche EDB : Antonio Enrico Felle, première version 2003.}
\ntext{\lnum{1}[Hic iacet in somno] pacis dig[na gra]taq(ue) p[er]\\
\lnum{2}[omnia - - - Lau]rentia q(uae) v[ixit a]nn(os) p(lus) m(inus)\\
\lnum{3}[- - - et cum ma]rito fecit an[n(os) - - -] m(ensem) I d(ies) XX\\
\lnum{4}[dep(osita) - - -] Merob(aude) II et Satu[rnino c]onss(ulibus)}
\nf{Témoins}{EDB (EDB5084) : texte concordant aux conventions de notation près (indice de similarité 0,974).}
\nf{Traduction}{Ici gît dans le sommeil de la paix, digne et agréable en toutes choses, […] Laurentia, qui vécut environ […] ans, et a vécu avec son mari […] ans 1 mois 20 jours ; déposée […] sous le consulat de Merobaudes pour la 2e fois et de Saturninus.}
\nf{Langue}{11 abréviations développées par l’éditeur (dont grataq(ue), q(uae), ann(os), p(lus), m(inus)).}
\nf{Chronologie}{Noms consulaires : Merobaudes (lu) et Saturninus (lu en partie). Âge déclaré : env. 1 m. 20 j. (partiellement restitué).}
\nf{Bibliographie}{ICUR IV 12245.}
\nf{Statut}{Lecture secondaire (EDCS) concordante avec l’EDB : lecture stabilisée entre deux bases, non collationnée sur l’ICVR ni sur le monument. Plus de la moitié des lettres sont restituées. L’EDCS signale une photographie (non consultée).}
\end{notice}

\begin{notice}{258}{383}{ICUR n.s. IV, 12246 · EDCS-37600681 · EDB5345}
\label{not:258}
\nf{Datation}{383 — directe : nom consulaire lu}
\nf{Genre}{Épitaphe ; ensembles A, B, C ; 91 lettres transmises ; 19 \% restituées.}
\nf{Références}{ICUR IV 12246.}
\nf{Provenance}{Rattachement selon l’EDB : \emph{Coem. anonymum ad viam Ardeatinam}. Conservation selon l’EDB : \emph{Coem. anonymum ad viam Ardeatinam}.}
\nf{Support}{Plaque de marbre. Inscription gravée. Texte métrique. Reproductions signalées par l’EDB : photographie PCAS (droits réservés ; non reproduites). Fiche EDB : Marida Pierno, première version 2003.}
\ntext{\lnum{1}requiescu]nt corpora casta\\
\lnum{2}[- - - in p]ace pridie Kal(endas) Sept(embres)\\
\lnum{3}[- - - Merobaude II] et Saturnino conss(ulibus)\\
\lnum{4}[- - - Vi]ctoria deposita in pac[e]}
\nf{Témoins}{EDB (EDB5345) : texte concordant aux conventions de notation près (indice de similarité 0,961).}
\nf{Traduction}{[…] reposent des corps chastes […] ; […] en paix, la veille des calendes de septembre […] ; sous le consulat de Merobaudes pour la 2e fois et de Saturninus ; […] Victoria déposée en paix.}
\nf{Langue}{3 abréviations développées par l’éditeur (dont Kal(endas), Sept(embres), conss(ulibus)).}
\nf{Chronologie}{Noms consulaires : Merobaudes (restitué) et Saturninus (lu). Dates du jour : décès : pridie Kalendas Septembres = 31 août 383 (lu).}
\nf{Commentaire}{La formule in pace est attestée seulement par restitution. Vocabulaire religieux : repos et sommeil de la mort — requiescunt (lu).}
\nf{Bibliographie}{ICUR IV 12246.}
\nf{Statut}{Lecture secondaire (EDCS) concordante avec l’EDB : lecture stabilisée entre deux bases, non collationnée sur l’ICVR ni sur le monument. L’EDCS signale une photographie (non consultée).}
\end{notice}

\begin{notice}{259}{383}{ICUR n.s. IV, 12534 · EDCS-37801595 · EDB42430}
\label{not:259}
\nf{Datation}{383 — directe : nom consulaire lu}
\nf{Genre}{Épitaphe ; ensembles A, B, C ; 89 lettres transmises ; 37 \% restituées.}
\nf{Références}{ICUR IV 12534.}
\nf{Provenance}{Rattachement selon l’EDB : \emph{inter vias Appiam et Ardeatinam}. Conservation selon l’EDB : \emph{Collegio Teutonico}.}
\nf{Support}{La fiche EDB ne figure pas dans l’état 2019 des données EpiDoc ; support non documenté.}
\ntext{\lnum{1}No]nas M(a)ias Merobaude\\
\lnum{2}[et Saturnino conss(ulibus) vi]xit annis XXX[- - -]\\
\lnum{3}[- - - ann]is XII et mensi(bu)s VI et\\
\lnum{4}[diebus - - -] sine filiis}
\nf{Témoins}{EDB (EDB42430) : texte concordant aux conventions de notation près (indice de similarité 0,957).}
\nf{Traduction}{[…] nones de mai, sous le consulat de Merobaudes et de Saturninus ; elle vécut 30 […] ans ; […] 12 ans et 6 mois et […] jours ; sans enfants.}
\nf{Langue}{3 abréviations développées par l’éditeur (dont M(a)ias, conss(ulibus), mensi(bu)s).}
\nf{Chronologie}{Noms consulaires : Merobaudes (lu) et Saturninus (restitué). Dates du jour : événement indéterminé : Nonas Maias = 7 mai 383 (lu). Âge déclaré : 12 a. 6 m. (partiellement restitué).}
\nf{Bibliographie}{ICUR IV 12534.}
\nf{Statut}{Lecture secondaire (EDCS) concordante avec l’EDB : lecture stabilisée entre deux bases, non collationnée sur l’ICVR ni sur le monument.}
\end{notice}

\annee{384}

\begin{notice}{260}{384}{ICUR n.s. VII, 19963 · EDCS-03900106 · EDB21504}
\label{not:260}
\nf{Datation}{384 — directe : nom consulaire lu}
\nf{Genre}{Épitaphe ; ensembles A, B, C ; 115 lettres transmises.}
\nf{Références}{Anagni 155 = ICUR VII 19963 = ILCV 2819a (em).}
\nf{Provenance}{Rattachement selon l’EDB : \emph{Coem. s.Hippolyti}. Conservation selon l’EDB : \emph{Musei Capitolini}.}
\nf{Support}{Plaque de marbre. 26 × 150 × 3 cm (h. × l. × ép.). Inscription gravée. Signes décrits par l’EDB : chrisme et chrisme avec alpha et oméga. Reproductions signalées par l’EDB : reproduction, Perret 1851 et reproduction, Solin et alii 1996 (droits réservés ; non reproduites). Fiche EDB : Anita Rocco, première version 2015.}
\ntext{\lnum{1}D(ominis) n(ostris) Claearco et Ricomede v(iris) c(larissimis)\\
\lnum{2}consulibus bene merenti Olibioni qui vixit\\
\lnum{3}ann⸢o⸣s XV me(n)s⸢e⸣s VI dies XX decessit\\
\lnum{4}die XII Kalendas Octobres in pace}
\nf{Témoins}{Export EDCS de 2020 : \emph{D(ominis) n(ostris) Claearco et Ricomede v(iris) c(larissimis) / consulibus bene merenti Olibioni qui vixit / anns XV me(n)ss VI dies XX decessit / die XII Kalendas Octobres in pace}. EDB (EDB21504) : texte concordant aux conventions de notation près (indice de similarité 0,958).}
\nf{Apparat}{⸢o⸣ : U gravé ; ⸢e⸣ : I gravé}
\nf{Traduction}{Sous le consulat de nos seigneurs (sic) Clearchus et Richomeres, clarissimes : au bien méritant Olibio, qui vécut 15 ans 6 mois 20 jours ; il mourut le 12e jour avant les calendes d’octobre, en paix.}
\nf{Langue}{L’éditeur signale 2 graphies corrigées : U gravé pour o ; I gravé pour e (confusion e/i).}
\nf{Chronologie}{Noms consulaires : Richomeres (lu). Dates du jour : décès : die XII Kalendas Octobres = 20 septembre 384 (lu). Âge déclaré : 15 a. 6 m. 20 j. (lu).}
\nf{Commentaire}{La formule in pace est conservée.}
\nf{Bibliographie}{Anagni 155 ; ICUR VII 19963 ; ILCV 2819a (em). Bibliographie complémentaire selon l’EDB : Solin et alii 1996, 155, con foto.}
\nf{Statut}{Lecture secondaire (EDCS) concordante avec l’EDB : lecture stabilisée entre deux bases, non collationnée sur l’ICVR ni sur le monument. L’EDCS signale une photographie (non consultée).}
\end{notice}

\begin{notice}{261}{384}{ICUR n.s. I, 3201 · EDCS-40300342 · EDB33151}
\label{not:261}
\nf{Datation}{384 — directe : lecture partielle d’un nom consulaire}
\nf{Genre}{Épitaphe ; ensembles A, B, C ; 97 lettres transmises ; 19 \% restituées.}
\nf{Références}{CIL V *672,67 = ICUR I 3201.}
\nf{Provenance}{Rattachement selon l’EDB : \emph{Incertae originis} (le volume I des ICVR réunit les inscriptions dont la provenance est perdue). Conservation selon l’EDB : \emph{Milano   ; Biblioteca Ambrosiana}.}
\nf{Support}{Plaque de marbre. 30 × 60 cm (h. × l.). Inscription gravée. Fiche EDB : Filippo Piazzolla, première version 2015.}
\ntext{\lnum{1}bene mer]⸢e⸣(n)ti in pace Ioannti(!) qui ⸢v⸣i\\
\lnum{2}[xit annos - - -]VIII et m(enses) X et di(e)s depositus\\
\lnum{3}[- - -]as Maias pos(t) conss(ulatum) Me\\
\lnum{4}[robaude et S]aturnino vv(iris) cc(larissimis)}
\nf{Témoins}{Export EDCS de 2020 : \emph{bene mer](n)ti in pace Ioannti(!) qui i/[xit annos 3]VIII et m(enses) X et di(e)s depositus / [3]as Maias pos(t) conss(ulatum) Me/[robaude et S]aturnino vv(iris) cc(larissimis)}. EDB (EDB33151) : texte concordant aux conventions de notation près (indice de similarité 0,92).}
\nf{Apparat}{⸢e⸣ : I gravé ; ⸢v⸣ : B gravé ; (!) : forme signalée par l’éditeur}
\nf{Traduction}{Au bien méritant, en paix, Ioannes (?), qui vécut […]8 ans et 10 mois et […] jours ; déposé […] mai, après le consulat de Merobaudes et de Saturninus, clarissimes.}
\nf{Langue}{L’éditeur signale 2 graphies corrigées : I gravé pour e (confusion e/i) ; B gravé pour v (bêtacisme). L’éditeur signale au moins une forme anormale (sic). 8 abréviations développées par l’éditeur (dont (n)ti, Ioannti(!), m(enses), di(e)s, pos(t), conss(ulatum)).}
\nf{Chronologie}{Noms consulaires : Merobaudes (restitué) et Saturninus (lu en partie). Datation par post-consulat. Âge déclaré : 10 m. (partiellement restitué).}
\nf{Commentaire}{La formule in pace est conservée.}
\nf{Bibliographie}{CIL V *672,67 ; ICUR I 3201. Bibliographie complémentaire selon l’EDB : ILCV, 2609A adn.}
\nf{Statut}{Lecture secondaire (EDCS) concordante avec l’EDB : lecture stabilisée entre deux bases, non collationnée sur l’ICVR ni sur le monument.}
\end{notice}

\begin{notice}{262}{384}{ICUR n.s. II, 4168 · EDCS-19900458 · EDB16942}
\label{not:262}
\nf{Datation}{384 — directe : nom consulaire lu}
\nf{Genre}{Épitaphe ; ensembles A, B, C ; 95 lettres transmises ; 35 \% restituées.}
\nf{Références}{CIL VI 32078 (p 4810) = ICUR II 4168 = ILCV +172 = ILCV 4620 (em).}
\nf{Provenance}{Rattachement selon l’EDB : \emph{Basilica S.Petri apostoli}. Conservation selon l’EDB : \emph{Basilica S.Petri apostoli}.}
\nf{Support}{Plaque de marbre. Inscription gravée. Fiche EDB : Antonio Enrico Felle, première version 2003.}
\ntext{\lnum{1}castissima]e(?) et in omn[ibus - - -]\\
\lnum{2}[- - -] admirabili co[niugi - - -]\\
\lnum{3}[- - -]e c(larissimae) m(emoriae) f(eminae) Valerius [- - -]\\
\lnum{4}[- - -] fecit V Idus Iul[ias conss(ulibus)]\\
\lnum{5}[Ricomere et] Clearc(h)o}
\nf{Témoins}{EDB (EDB16942) : texte concordant aux conventions de notation près (indice de similarité 0,978).}
\nf{Traduction}{À son épouse très chaste et admirable en toutes choses […] femme de clarissime mémoire, Valerius […] a fait (ce monument) le 5e jour avant les ides de juillet, sous le consulat de Richomeres et de Clearchus.}
\nf{Langue}{6 abréviations développées par l’éditeur (dont castissimae(?), c(larissimae), m(emoriae), f(eminae), conss(ulibus), Clearc(h)o).}
\nf{Chronologie}{Noms consulaires : Richomeres (restitué) et Clearchus (lu). Dates du jour : événement indéterminé : V Idus Iulias = 11 juillet 384 (partiellement restitué).}
\nf{Commentaire}{Rang sénatorial d’une personne commémorée ( admirabili coniugi … …e clarissimae memoriae femi).}
\nf{Bibliographie}{CIL VI 32078 (p 4810) ; ICUR II 4168 ; ILCV +172 ; ILCV 4620 (em). Bibliographie complémentaire selon l’EDB : CIL VI, 32078.}
\nf{Statut}{Lecture secondaire (EDCS) concordante avec l’EDB : lecture stabilisée entre deux bases, non collationnée sur l’ICVR ni sur le monument. L’EDCS signale une photographie (non consultée).}
\end{notice}

\begin{notice}{263}{384}{ICUR n.s. I, 62 · EDCS-38600081 · EDB27435}
\label{not:263}
\nf{Datation}{384 — directe : lecture partielle d’un nom consulaire}
\nf{Genre}{Épitaphe ; ensembles A, B, C ; 44 lettres transmises ; 36 \% restituées.}
\nf{Références}{ICUR I 62.}
\nf{Provenance}{Rattachement selon l’EDB : \emph{Incertae originis} (le volume I des ICVR réunit les inscriptions dont la provenance est perdue).}
\nf{Support}{Support non précisé par l’EDB. Fiche EDB : Filippo Piazzolla, première version 2013.}
\ntext{\lnum{1}Ricomere] et Clea[rcho conss(ulibus) - - -]\\
\lnum{2}[- - -] qui vixit [- - -]\\
\lnum{3}[- - - I]dus Sep(tembres) [}
\nf{Témoins}{EDB (EDB27435 ; écarts à examiner) : \emph{[--- Ricomere] et Clea[rco cons(ulibu)s ---] / [---] qui vixit [---] / [--- i]dus sep[tembres ---]}.}
\nf{Traduction}{Sous le consulat de Richomeres et de Clearchus […] ; […] qui vécut […] ; […] ides de septembre.}
\nf{Langue}{2 abréviations développées par l’éditeur (dont conss(ulibus), Sep(tembres)).}
\nf{Chronologie}{Noms consulaires : Richomeres (restitué) et Clearchus (lu en partie). Dates du jour : événement indéterminé : … Idus Septembres (partiellement restitué).}
\nf{Bibliographie}{ICUR I 62.}
\nf{Statut}{Lecture secondaire (EDCS) ; écarts avec l’EDB à arbitrer sur l’ICVR.}
\end{notice}

\begin{notice}{264}{384}{ICUR n.s. I, 177 · EDCS-38600196 · EDB27580}
\label{not:264}
\nf{Datation}{384 — directe : lecture partielle d’un nom consulaire}
\nf{Genre}{Épitaphe ; ensembles A, B, C ; 50 lettres transmises ; 58 \% restituées.}
\nf{Références}{ICUR I 177.}
\nf{Provenance}{Rattachement selon l’EDB : \emph{Incertae originis} (le volume I des ICVR réunit les inscriptions dont la provenance est perdue).}
\nf{Support}{Support non précisé par l’EDB. Fiche EDB : Filippo Piazzolla, première version 2013.}
\ntext{\lnum{1}vixit] annos LX de[p(ositus ?) - - -]\\
\lnum{2}[- - - Ricomere] et Clear[cho conss(ulibus)}
\nf{Témoins}{EDB (EDB27580 ; écarts à examiner) : \emph{[ --- vixit] annos LX de[p(osit---) ---] / [--- Ricomere] et Clear[cho conss(ulibus)] [texte tronqué dans la liste de résultats]}.}
\nf{Traduction}{[…] vécut 60 ans ; déposé […] ; sous le consulat de Richomeres et de Clearchus.}
\nf{Langue}{2 abréviations développées par l’éditeur (dont dep(ositus ?), conss(ulibus)).}
\nf{Chronologie}{Noms consulaires : Richomeres (restitué) et Clearchus (lu en partie). Âge déclaré : 60 a. (partiellement restitué).}
\nf{Bibliographie}{ICUR I 177.}
\nf{Statut}{Lecture secondaire (EDCS) ; écarts avec l’EDB à arbitrer sur l’ICVR. Plus de la moitié des lettres sont restituées.}
\end{notice}

\begin{notice}{265}{384}{ICUR n.s. I, 331 · EDCS-38600325 · EDB27988}
\label{not:265}
\nf{Datation}{384 — directe : nom consulaire lu}
\nf{Genre}{Épitaphe ; ensembles A, B, C ; 117 lettres transmises ; 13 \% restituées.}
\nf{Références}{ICUR I 331 = ILCV 4248.}
\nf{Provenance}{Rattachement selon l’EDB : \emph{Incertae originis} (le volume I des ICVR réunit les inscriptions dont la provenance est perdue).}
\nf{Support}{Support non précisé par l’EDB. Fiche EDB : Filippo Piazzolla, première version 2014.}
\ntext{\lnum{1}[- - -]o s(u)ae Eudoxiae ca(ri)ssimae femina[e - - -]\\
\lnum{2}[- - -]XVI et fecit cum virginio suo an[nos - - -]\\
\lnum{3}[- - - Rico]mere et Clearcho consulibus V Idus No[v(embres) - - -]\\
\lnum{4}[- - -] depositus opera eorum}
\nf{Témoins}{EDB (EDB27988) : texte concordant aux conventions de notation près (indice de similarité 0,921).}
\nf{Traduction}{[…] à sa […] Eudoxia, femme très chère […] ; […]16 (ans), et elle a vécu avec son ‘vierge’ (époux) […] ans ; […] sous le consulat de Richomeres et de Clearchus, le 5e jour avant les ides de novembre […] ; […] déposé ; ouvrage de ceux-ci (?).}
\nf{Langue}{3 abréviations développées par l’éditeur (dont s(u)ae, ca(ri)ssimae, Nov(embres)).}
\nf{Chronologie}{Noms consulaires : Richomeres (lu en partie) et Clearchus (lu). Dates du jour : événement indéterminé : V Idus Novembres = 9 novembre 384 (partiellement restitué).}
\nf{Bibliographie}{ICUR I 331 ; ILCV 4248. Bibliographie complémentaire selon l’EDB : ILCV, 4248.}
\nf{Statut}{Lecture secondaire (EDCS) concordante avec l’EDB : lecture stabilisée entre deux bases, non collationnée sur l’ICVR ni sur le monument.}
\end{notice}

\begin{notice}{266}{384}{ICUR n.s. I, 481 · EDCS-38700889 · EDB28430}
\label{not:266}
\nf{Datation}{384 — directe : nom consulaire lu}
\nf{Genre}{Épitaphe ; ensembles A, B, C ; 52 lettres transmises ; 40 \% restituées.}
\nf{Références}{ICUR I 481 = ILCV 4763.}
\nf{Provenance}{Rattachement selon l’EDB : \emph{Incertae originis} (le volume I des ICVR réunit les inscriptions dont la provenance est perdue). Conservation selon l’EDB : \emph{s. Maria in Cosmedin}.}
\nf{Support}{Plaque de marbre. 45 × 30 cm (h. × l.). Inscription gravée. Fiche EDB : Filippo Piazzolla, première version 2014.}
\ntext{\lnum{1}Floridus [- - -]\\
\lnum{2}floruit ac [- - -]\\
\lnum{3}-----4 pr(a)eces[sit - - - Ricomere] 5 et C⸢l⸣(e)arc(h)o [cons(ulibus)]}
\nf{Témoins}{Export EDCS de 2020 : \emph{Floridus [3] / floruit ac [3] / [6] / pr(a)eces[sit 3 Ricomere] / et C(e)arc(h)o [cons(ulibus)]}. EDB (EDB28430) : texte concordant aux conventions de notation près (indice de similarité 0,956).}
\nf{Apparat}{⸢l⸣ : I gravé}
\nf{Traduction}{Floridus […] fleurit et […] ; […] précéda […] ; sous le consulat de Richomeres et de Clearchus.}
\nf{Langue}{L’éditeur signale une graphie corrigée : I gravé pour l.}
\nf{Chronologie}{Noms consulaires : Richomeres (restitué) et Clearchus (lu).}
\nf{Bibliographie}{ICUR I 481 ; ILCV 4763.}
\nf{Statut}{Lecture secondaire (EDCS) concordante avec l’EDB : lecture stabilisée entre deux bases, non collationnée sur l’ICVR ni sur le monument.}
\end{notice}

\begin{notice}{267}{384}{ICUR n.s. I, 1210 · EDCS-38700749 · EDB30157}
\label{not:267}
\nf{Datation}{384 — indirecte}
\nf{Genre}{Épitaphe ; ensembles A, B, C ; 26 lettres transmises ; 15 \% restituées.}
\nf{Références}{ICUR I 1210.}
\nf{Provenance}{Rattachement selon l’EDB : \emph{Incertae originis} (le volume I des ICVR réunit les inscriptions dont la provenance est perdue). Conservation selon l’EDB : \emph{Antiquarium del Celio}.}
\nf{Support}{Plaque de marbre. 21 × 23 cm (h. × l.). Inscription gravée. Fiche EDB : Filippo Piazzolla, première version 2014.}
\ntext{\lnum{1}]AEIDE RO[- - -]\\
\lnum{2}[- - - anno]s XLVII NIC[- - -]\\
\lnum{3}[- - -]CVARCV[}
\nf{Témoins}{EDB (EDB30157) : texte concordant aux conventions de notation près (indice de similarité 0,939).}
\nf{Traduction}{[…] 47 ans […] [texte inintelligible].}
\nf{Langue}{Aucune particularité graphique n’est signalée par l’éditeur dans l’état 2022 de la base.}
\nf{Bibliographie}{ICUR I 1210.}
\nf{Statut}{Lecture secondaire (EDCS) concordante avec l’EDB : lecture stabilisée entre deux bases, non collationnée sur l’ICVR ni sur le monument.}
\end{notice}

\begin{notice}{268}{384}{ICUR n.s. I, 2099 · EDCS-39300319 · EDB32083}
\label{not:268}
\nf{Datation}{384 — formule transmise par l’EDB (lue)}
\nf{Genre}{Épitaphe ; ensembles A, B, C ; 251 lettres transmises ; 16 \% restituées.}
\nf{Références}{ICUR I 2099 = ILCV 2782 (add).}
\nf{Provenance}{Rattachement selon l’EDB : \emph{Incertae originis} (le volume I des ICVR réunit les inscriptions dont la provenance est perdue). Conservation selon l’EDB : \emph{Musei Vaticani: Galleria Lapidaria}.}
\nf{Support}{Plaque de marbre. 64 × 98 cm (h. × l.). Inscription gravée. Fiche EDB : Filippo Piazzolla, première version 2015.}
\ntext{\lnum{1}Hoc tumulum patris filius Benedict[us]\\
\lnum{2}fieri volvit causa amoris paternae re[cordatio]\\
\lnum{3}nis Lucifero homini innocentiss[imo amico(?)]\\
\lnum{4}amicorum qui dum in rebus humanis [fuit]\\
\lnum{5}omnibus paratissimus et dulcis adqu[e]\\
\lnum{6}fuisse videbatur vix enim semper sin[e culpa(?) tem]\\
\lnum{7}poris sui annis LXV et m(ensibus) VI d(iebus) XII decess[it in pace(?)]}
\nf{Témoins}{EDB (EDB32083) : texte concordant aux conventions de notation près (indice de similarité 0,937).}
\nf{Traduction}{Ce tombeau de son père, son fils Benedictus a voulu qu’il fût fait par amour et souvenir paternel, pour Lucifer, homme très innocent, ami des amis, qui, tant qu’il fut dans les choses humaines, se montra très disponible et doux envers tous ; car il vécut presque toujours sans faute de son temps ; à 65 ans 6 mois 12 jours il mourut en paix.}
\nf{Langue}{5 abréviations développées par l’éditeur (dont amico(?), culpa(?), m(ensibus), d(iebus), pace(?)).}
\nf{Chronologie}{La fiche EDB (EDB32083) transmet la dernière ligne, absente du texte EDCS : \emph{VII k(alendas) iul(ias) Ricomere et Clearcho cons(ulibu)s}. La date est donc le 25 juin 384, fondée sur une formule consulaire lue. L'EDB décrit une palme et un oiseau ; la plaque (64 × 98 cm) est conservée aux Musées du Vatican (Galleria Lapidaria). Âge déclaré : 65 a. 6 m. 12 j. (partiellement restitué).}
\nf{Commentaire}{La formule in pace est attestée seulement par restitution. Commentaire développé : dossier p.~\pageref{dos:268}.}
\nf{Bibliographie}{ICUR I 2099 ; ILCV 2782 (add). Bibliographie complémentaire selon l’EDB : ILCV, 2782 ; Lega 2000.}
\nf{Statut}{Lecture secondaire (EDCS) concordante avec l’EDB : lecture stabilisée entre deux bases, non collationnée sur l’ICVR ni sur le monument.}
\end{notice}

\begin{notice}{269}{384}{ICUR n.s. I, 3200 · EDCS-40300341 · EDB33150}
\label{not:269}
\nf{Datation}{384 — directe : nom consulaire lu}
\nf{Genre}{Épitaphe ; ensembles A, B, C ; 126 lettres transmises.}
\nf{Références}{ICUR I 3200 = ILCV 1270 (add).}
\nf{Provenance}{Rattachement selon l’EDB : \emph{Incertae originis} (le volume I des ICVR réunit les inscriptions dont la provenance est perdue).}
\nf{Support}{Support non précisé par l’EDB. Fiche EDB : Filippo Piazzolla, première version 2015.}
\ntext{\lnum{1}Mirae innocentiae adq(ue) eximiae\\
\lnum{2}bonitatis hic requiesicit Leopardus\\
\lnum{3}lector de Pudentiana qui vixit\\
\lnum{4}ann(os) XXIIII def(unctus) VIII Kal(endas) Dec(embres)\\
\lnum{5}Ricome⸢r⸣e et Clearc(h)o con(sulibus)}
\nf{Témoins}{Export EDCS de 2020 : \emph{Mirae innocentiae adq(ue) eximiae / bonitatis hic requiesicit Leopardus / lector de Pudentiana qui vixit / ann(os) XXIIII def(unctus) VIII Kal(endas) Dec(embres) / Ricomee et Clearc(h)o con(sulibus)}. EDB (EDB33150) : texte concordant aux conventions de notation près (indice de similarité 0,961).}
\nf{Apparat}{⸢r⸣ : D gravé}
\nf{Traduction}{D’une admirable innocence et d’une bonté exceptionnelle, ici repose Leopardus, lecteur de (titre de) Pudentiana, qui vécut 24 ans ; mort le 8e jour avant les calendes de décembre, sous le consulat de Richomeres et de Clearchus.}
\nf{Langue}{L’éditeur signale une graphie corrigée : D gravé pour r.}
\nf{Chronologie}{Noms consulaires : Richomeres (lu) et Clearchus (lu). Dates du jour : décès : VIII Kalendas Decembres = 24 novembre 384 (lu). Âge déclaré : 24 a. (lu).}
\nf{Commentaire}{Fonction ou métier mentionné : lector.}
\nf{Bibliographie}{ICUR I 3200 ; ILCV 1270 (add). Bibliographie complémentaire selon l’EDB : ILCV, 1270 ; Petrignani 1934, 5 ; Angelelli 2010, 3-4.}
\nf{Statut}{Lecture secondaire (EDCS) concordante avec l’EDB : lecture stabilisée entre deux bases, non collationnée sur l’ICVR ni sur le monument.}
\end{notice}

\begin{notice}{270}{384}{ICUR n.s. II, 4480 · EDCS-34100543 · EDB19138}
\label{not:270}
\nf{Datation}{384 — directe : nom consulaire lu}
\nf{Genre}{Épitaphe ; ensembles A, B, C ; 29 lettres transmises ; 72 \% restituées.}
\nf{Références}{ICUR II 4480.}
\nf{Provenance}{Rattachement selon l’EDB : \emph{Coem. ad duos Felices}.}
\nf{Support}{Enduit de mortier de mortier. Inscription tracée dans le mortier frais. Fiche EDB : Antonio Enrico Felle, première version 2013.}
\ntext{\lnum{1}] Ricomede [et Clearcho conss(ulibus)]}
\nf{Témoins}{EDB (EDB19138) : texte concordant aux conventions de notation près (indice de similarité 1).}
\nf{Traduction}{[…] sous le consulat de Richomeres et de Clearchus.}
\nf{Langue}{1 abréviation développée par l’éditeur (dont conss(ulibus)).}
\nf{Chronologie}{Noms consulaires : Richomeres (lu) et Clearchus (restitué).}
\nf{Bibliographie}{ICUR II 4480.}
\nf{Statut}{Lecture secondaire (EDCS) concordante avec l’EDB : lecture stabilisée entre deux bases, non collationnée sur l’ICVR ni sur le monument. Plus de la moitié des lettres sont restituées.}
\end{notice}

\begin{notice}{271}{384}{ICUR n.s. II, 4750 · EDCS-34201006 · EDB13472}
\label{not:271}
\nf{Datation}{384 — directe : nom consulaire lu}
\nf{Genre}{Épitaphe ; ensembles A, B, C ; 82 lettres transmises.}
\nf{Références}{ICUR II 4750 = ILCV 2879 (em).}
\nf{Provenance}{Rattachement selon l’EDB : \emph{Coem. Generosae}. Conservation selon l’EDB : \emph{coem. Generosae}.}
\nf{Support}{Plaque de marbre. Inscription gravée. Fiche EDB : Antonio Enrico Felle, première version 2013.}
\ntext{\lnum{1}Vincenti(a)e bene merente in pace\\
\lnum{2}qu(a)e vixit annos p(lus) m(inus) LXXX\\
\lnum{3}exivit d(ie) III Idus Dec(em)b(res) Re(c)em(e)de\\
\lnum{4}et Clearc(h)o cons(ulibus)}
\nf{Témoins}{EDB (EDB13472) : texte concordant aux conventions de notation près (indice de similarité 0,985).}
\nf{Traduction}{À Vincentia, bien méritante, en paix, qui vécut environ 80 ans ; elle s’en alla le 3e jour avant les ides de décembre, sous le consulat de Richomeres et de Clearchus.}
\nf{Langue}{11 abréviations développées par l’éditeur (dont Vincenti(a)e, qu(a)e, p(lus), m(inus), d(ie), Dec(em)b).}
\nf{Chronologie}{Noms consulaires : Clearchus (lu). Dates du jour : décès : die III Idus Decembres = 11 décembre 384 (lu). Âge déclaré : env. 80 a. 3 j. (lu).}
\nf{Commentaire}{La formule in pace est conservée.}
\nf{Bibliographie}{ICUR II 4750 ; ILCV 2879 (em).}
\nf{Statut}{Lecture secondaire (EDCS) concordante avec l’EDB : lecture stabilisée entre deux bases, non collationnée sur l’ICVR ni sur le monument.}
\end{notice}

\begin{notice}{272}{384}{ICUR n.s. II, 6040 · EDCS-34200666 · EDB8752}
\label{not:272}
\nf{Datation}{384 — directe : nom consulaire lu}
\nf{Genre}{Épitaphe ; ensembles A, B, C ; 69 lettres transmises ; 33 \% restituées.}
\nf{Références}{ICUR II 6040 = AE 1905, 78.}
\nf{Provenance}{Rattachement selon l’EDB : \emph{Coem. Commodillae}.}
\nf{Support}{Plaque de marbre. Inscription gravée. Reproductions signalées par l’EDB : photographie PCAS ×2 (droits réservés ; non reproduites). Fiche EDB : Carlo Carletti, première version 2013.}
\ntext{\lnum{1}deposita Le[a in pa]\\
\lnum{2}ce Ricomedi e[t Clea]\\
\lnum{3}rc(h)o cons(ulibu)s X[- - - Kal(endas)]\\
\lnum{4}No⸢v⸣(embres) qu(a)e vi⸢x⸣it [ann(os)]\\
\lnum{5}XXXII}
\nf{Témoins}{Export EDCS de 2020 : \emph{deposita Le[a in pa]/ce Ricomedi e[t Clea]/rc(h)o cons(ulibu)s X[3 Kal(endas)] / No(embres) qu(a)e viit [ann(os)] / XXXII}. EDB (EDB8752) : texte concordant aux conventions de notation près (indice de similarité 0,963).}
\nf{Apparat}{⸢v⸣ : B gravé ; ⸢x⸣ : S gravé}
\nf{Traduction}{Déposée, Lea, en paix, sous le consulat de Richomeres et de Clearchus, le 1x jour avant les calendes de novembre ; elle vécut 32 ans.}
\nf{Langue}{L’éditeur signale 2 graphies corrigées : B gravé pour v (bêtacisme) ; S gravé pour x (s pour x).}
\nf{Chronologie}{Noms consulaires : Richomeres (lu) et Clearchus (lu en partie). Dates du jour : événement indéterminé : … Kalendas Novembres (restitué). Âge déclaré : 32 a. (partiellement restitué).}
\nf{Commentaire}{La formule in pace est attestée seulement par restitution.}
\nf{Bibliographie}{ICUR II 6040 ; AE 1905, 78.}
\nf{Statut}{Lecture secondaire (EDCS) concordante avec l’EDB : lecture stabilisée entre deux bases, non collationnée sur l’ICVR ni sur le monument. L’EDCS signale une photographie (non consultée).}
\end{notice}

\begin{notice}{273}{384}{ICUR n.s. III, 8426 · EDCS-35802013 · EDB25701}
\label{not:273}
\nf{Datation}{384 — directe : lecture partielle d’un nom consulaire}
\nf{Genre}{Épitaphe ; ensembles A, B, C ; 46 lettres transmises ; 20 \% restituées.}
\nf{Références}{ICUR III 8426.}
\nf{Provenance}{Rattachement selon l’EDB : \emph{Incertum locum in coem. Domitillae}. Conservation selon l’EDB : \emph{coem. Domitillae, aula supra coemeterium}.}
\nf{Support}{Plaque de marbre. Inscription gravée. Signes décrits par l’EDB : chrisme et chrisme avec alpha et oméga. Reproductions signalées par l’EDB : photographie PCAS ×2 (droits réservés ; non reproduites). Fiche EDB : Carlo Carletti / Antonio Enrico Felle, première version 2013.}
\ntext{\lnum{1}Ricomere et Clea]rco c(on)s(ulibus)\\
\lnum{2}[dep(ositus ?) - - -]uda VIII Kal(endas) Augustas}
\nf{Témoins}{EDB (EDB25701 ; écarts à examiner) : \emph{[---?] ((monogramma Christi cum alpha et omega)) / [--- Ricomere et Cle]arco cc〚o〛(n)s(ulibu)s / [---]uda VIII kal(endas) agustas}.}
\nf{Traduction}{Sous le consulat de Richomeres et de Clearchus, déposé […]uda le 8e jour avant les calendes d’août.}
\nf{Langue}{4 abréviations développées par l’éditeur (dont c(on)s, (ulibus), dep(ositus ?), Kal(endas)).}
\nf{Chronologie}{Noms consulaires : Richomeres (restitué). Dates du jour : déposition : VIII Kalendas Augustas = 25 juillet 384 (lu).}
\nf{Bibliographie}{ICUR III 8426.}
\nf{Statut}{Lecture secondaire (EDCS) ; écarts avec l’EDB à arbitrer sur l’ICVR. L’EDCS signale une photographie (non consultée).}
\end{notice}

\begin{notice}{274}{384}{ICUR n.s. IV, 11112 · EDCS-37200394 · EDB39574}
\label{not:274}
\nf{Datation}{384 — directe : lecture partielle d’un nom consulaire}
\nf{Genre}{Épitaphe ; ensembles A, B, C ; 29 lettres transmises ; 41 \% restituées.}
\nf{Références}{ICUR IV 11112.}
\nf{Provenance}{Rattachement selon l’EDB : \emph{Coem. Callisti area subdialis et incerti loci}. Conservation selon l’EDB : \emph{Coem. Callisti area subdialis et incerti loci}.}
\nf{Support}{Plaque de marbre. 24,5 × 23,5 × 3,5 cm (h. × l. × ép.). Inscription gravée. Fiche EDB : Marida Pierno, première version 2015.}
\ntext{\lnum{1}Ricomede et Cle]arch[o conss(ulibus)]}
\nf{Témoins}{EDB (EDB39574 ; écarts à examiner) : \emph{------ / [---Ricomede et Cle]arch[o]}.}
\nf{Traduction}{Sous le consulat de Richomeres et de Clearchus.}
\nf{Langue}{1 abréviation développée par l’éditeur (dont conss(ulibus)).}
\nf{Chronologie}{Noms consulaires : Richomeres (restitué) et Clearchus (lu en partie).}
\nf{Bibliographie}{ICUR IV 11112.}
\nf{Statut}{Lecture secondaire (EDCS) ; écarts avec l’EDB à arbitrer sur l’ICVR.}
\end{notice}

\annee{385}

\begin{notice}{275}{385}{ICUR n.s. I, 1441 · EDCS-34501147 · EDB31206}
\label{not:275}
\nf{Datation}{385 — directe : nom consulaire lu}
\nf{Genre}{Épitaphe ; ensembles A, B, C ; 161 lettres transmises.}
\nf{Références}{ICUR I 1441 = ICVaticano p 350 = ILCV 4460 = AE 1997, 166.}
\nf{Provenance}{Rattachement selon l’EDB : \emph{n.d.}. Conservation selon l’EDB : \emph{Musei Vaticani: Lapidario Cristiano}.}
\nf{Support}{Plaque de marbre. 41 × 52 cm (h. × l.). Inscription gravée. Signes décrits par l’EDB : chrisme et chrisme avec alpha et oméga. Reproductions signalées par l’EDB : reproduction, Perret 1851 (droits réservés ; non reproduites). Fiche EDB : Filippo Piazzolla, première version 2015.}
\ntext{\lnum{1}Civem(!) Armeniacum Cappadocem\\
\lnum{2}n⸢o⸣mine Quirillus pius omni\\
\lnum{3}bus amicatus vixit annos\\
\lnum{4}p(lus) m(inus) LXX praecessit ad pacem\\
\lnum{5}VI Id(us) Mart(ia)\\
\lnum{6}s\\
\lnum{7}post conss(ulatum ) Rico\\
\lnum{8}mere(s) et Clearc(h)⸢i⸣ die Lunae\\
\lnum{9}bisomu(m) cum Ype« r »echiam(!) co⸢m⸣parem(!)}
\nf{Témoins}{Export EDCS de 2020 : \emph{Civem(!) Armeniacum Cappadocem / nmine Quirillus pius omni/bus amicatus vixit annos / p(lus) m(inus) LXX praecessit ad pacem / VI Id(us) Mart(ia)//s // post conss(ulatum ) Rico/mere(s) et Clearc(h) die Lunae / bisomu(m) cum Ype« r »echiam(!) coparem(!)}. EDB (EDB31206) : texte concordant aux conventions de notation près (indice de similarité 0,962).}
\nf{Apparat}{⸢o⸣ : V gravé ; ⸢i⸣ : O gravé ; ⸢m⸣ : N gravé ; (!) : forme signalée par l’éditeur}
\nf{Traduction}{Citoyen arménien de Cappadoce, nommé Quirillus (Cyrillus), pieux, ami de tous ; il vécut environ 70 ans ; il partit en avant vers la paix le 6e jour avant les ides de mars, après le consulat de Richomeres et de Clearchus, un lundi ; tombe double avec sa compagne Hyperechia.}
\nf{Langue}{L’éditeur signale 3 graphies corrigées : V gravé pour o (confusion o/u) ; O gravé pour i ; N gravé pour m (n pour m). L’éditeur signale au moins une forme anormale (sic). 11 abréviations développées par l’éditeur (dont Civem(!), p(lus), m(inus), Id(us), Mart(ia), conss(ulatum )).}
\nf{Chronologie}{Noms consulaires : Richomeres (lu) et Clearchus (lu). Datation par post-consulat. Dates du jour : événement indéterminé : VI Idus Martia = 10 mars 385 (lu). Âge déclaré : env. 70 a. (lu).}
\nf{Commentaire}{Commentaire développé : dossier p.~\pageref{dos:275}.}
\nf{Bibliographie}{ICUR I 1441 ; ICVaticano p 350 ; ILCV 4460 ; AE 1997, 166. Bibliographie complémentaire selon l’EDB : Di Stefano Manzella 1997, 350-351, 3.12.25, con foto (C. Ricci) ; Ferrua 1979b, 37, n. 1441 ; ILCV, 4460 ; Lega 2000, 65.}
\nf{Statut}{Lecture secondaire (EDCS) concordante avec l’EDB : lecture stabilisée entre deux bases, non collationnée sur l’ICVR ni sur le monument. L’EDCS signale une photographie (non consultée).}
\end{notice}

\begin{notice}{276}{385}{ICUR n.s. I, 3202 · EDCS-40300343 · EDB33152}
\label{not:276}
\nf{Datation}{385 — directe : nom consulaire lu}
\nf{Genre}{Épitaphe ; ensembles A, B, C ; 108 lettres transmises.}
\nf{Références}{ICUR I 3202 = ILCV 1480 = RIRIOe 95.}
\nf{Provenance}{Rattachement selon l’EDB : \emph{Incertae originis} (le volume I des ICVR réunit les inscriptions dont la provenance est perdue). Conservation selon l’EDB : \emph{Wien   ; Kunsthistorisches Museum}.}
\nf{Support}{Plaque de marbre. 30 × 27 × 2 cm (h. × l. × ép.). Inscription gravée. Signes décrits par l’EDB : chrisme. Fiche EDB : Filippo Piazzolla, première version 2015.}
\ntext{\lnum{1}Fl(avio) Arcadio et Baud\\
\lnum{2}one conss(ulibus) X Kal(endas) Iu\\
\lnum{3}lias defunctus Leo\\
\lnum{4}ntius neof⸢y⸣tus qui\\
\lnum{5}vixit ann(os) plus minu\\
\lnum{6}s XXVIII me(n)s(es) V d(ies) XV\\
\lnum{7}⸢b⸣ene merenti in pa\\
\lnum{8}ce}
\nf{Témoins}{Export EDCS de 2020 : \emph{Fl(avio) Arcadio et Baud/one conss(ulibus) X Kal(endas) Iu/lias defunctus Leo/ntius neoftus qui / vixit ann(os) plus minu/s XXVIII me(n)s(es) V d(ies) XV / ene merenti in pa/ce}. EDB (EDB33152) : texte concordant aux conventions de notation près (indice de similarité 0,931).}
\nf{Apparat}{⸢y⸣ : I gravé ; ⸢b⸣ : V gravé}
\nf{Traduction}{Sous le consulat de Flavius Arcadius et de Bauto, le 10e jour avant les calendes de juillet, est mort Leontius, néophyte, qui vécut environ 28 ans 5 mois 15 jours ; bien méritant, en paix.}
\nf{Langue}{L’éditeur signale 2 graphies corrigées : I gravé pour y (graphie de y) ; V gravé pour b (bêtacisme).}
\nf{Chronologie}{Noms consulaires : Arcadius (lu) et Bauto (lu). Ne sont pas des consuls : defunctus Leo/ntius neof<y=I>tus (défunt, néophyte). Dates du jour : événement indéterminé : X Kalendas Iulias = 22 juin 385 (lu). Âge déclaré : env. 28 a. 5 m. 15 j. (lu).}
\nf{Commentaire}{La formule in pace est conservée. Vocabulaire religieux : initiation et foi — neofytus (lu).}
\nf{Bibliographie}{ICUR I 3202 ; ILCV 1480 ; RIRIOe 95. Bibliographie complémentaire selon l’EDB : ILCV, 1480.}
\nf{Statut}{Lecture secondaire (EDCS) concordante avec l’EDB : lecture stabilisée entre deux bases, non collationnée sur l’ICVR ni sur le monument.}
\end{notice}

\begin{notice}{277}{385}{ICUR n.s. II, 4481 · EDCS-34100544 · EDB19139}
\label{not:277}
\nf{Datation}{385 — directe : nom consulaire lu}
\nf{Genre}{Épitaphe ; ensembles A, B, C ; 35 lettres transmises.}
\nf{Références}{ICUR II 4481.}
\nf{Provenance}{Rattachement selon l’EDB : \emph{Coem. ad duos Felices}.}
\nf{Support}{Enduit de mortier de mortier. Inscription tracée dans le mortier frais. Fiche EDB : Antonio Enrico Felle, première version 2013.}
\ntext{\lnum{1}]X Kalendas Ian(uarias) [- - -]CIIE\\
\lnum{2}Ar\\
\lnum{3}ca\\
\lnum{4}di\\
\lnum{5}[- - -]elio consulatu}
\nf{Témoins}{EDB (EDB19139 ; écarts à examiner) : \emph{[dep(osit---)] X kalendas ian(uarias) [in pa]c⟨a⟩e / Ar/ca/di /[---] Elio consulatu}.}
\nf{Traduction}{Le 10e jour avant les calendes de janvier […] ; (sous le consulat d’)Arcadius […].}
\nf{Langue}{1 abréviation développée par l’éditeur (dont Ian(uarias)).}
\nf{Chronologie}{Noms consulaires : Arcadius (lu). Dates du jour : événement indéterminé : X Kalendas Ianuarias = 23 décembre 385 (lu).}
\nf{Bibliographie}{ICUR II 4481.}
\nf{Statut}{Lecture secondaire (EDCS) ; écarts avec l’EDB à arbitrer sur l’ICVR.}
\end{notice}

\annee{383}

\begin{notice}{278}{383}{ICUR n.s. II, 6037 · EDCS-34200663 · EDB3285}
\label{not:278}
\nf{Datation}{383 — directe : lecture partielle d’un nom consulaire}
\nf{Genre}{Épitaphe ; ensembles A, B, C ; 81 lettres transmises ; 26 \% restituées.}
\nf{Références}{ICUR II 6037.}
\nf{Provenance}{Rattachement selon l’EDB : \emph{Coem. Commodillae}.}
\nf{Support}{Plaque de marbre. Inscription gravée. Reproductions signalées par l’EDB : photographie PCAS ×2 (droits réservés ; non reproduites). Fiche EDB : Carlo Carletti, première version 2013.}
\ntext{\lnum{1}[- - -]tus fidelis qui vixit\\
\lnum{2}[anni]s X mens(ibus) VI diebus XXIIII\\
\lnum{3}[depositus - - -] Aug(ustas) [Merob]aude II et Satur\\
\lnum{4}[nin]o cons(ulibu)s}
\nf{Témoins}{EDB (EDB3285) : texte concordant aux conventions de notation près (indice de similarité 0,971).}
\nf{Traduction}{[…], fidèle, qui vécut 10 ans 6 mois 24 jours ; déposé […] août, sous le consulat de Merobaudes pour la 2e fois et de Saturninus.}
\nf{Langue}{3 abréviations développées par l’éditeur (dont mens(ibus), Aug(ustas), cons(ulibu)s).}
\nf{Chronologie}{Noms consulaires : Merobaudes (lu en partie) et Saturninus (lu en partie). \emph{Merobaude II et Saturnino} désigne 383. L'EDB (EDB3285) date de 385, alors que la fiche voisine EDB8750 (même consulat) est datée de 383 : incohérence interne probable de l'EDB. L'année 383 est maintenue. Âge déclaré : 10 a. 6 m. 24 j. (partiellement restitué).}
\nf{Commentaire}{Vocabulaire religieux : initiation et foi — fidelis (lu).}
\nf{Bibliographie}{ICUR II 6037.}
\nf{Statut}{Lecture secondaire (EDCS) concordante avec l’EDB : lecture stabilisée entre deux bases, non collationnée sur l’ICVR ni sur le monument. L’EDCS signale une photographie (non consultée).}
\end{notice}

\annee{385}

\begin{notice}{279}{385}{ICUR n.s. II, 6041 · EDCS-34200667 · EDB8753}
\label{not:279}
\nf{Datation}{385 — directe : nom consulaire lu}
\nf{Genre}{Épitaphe ; ensembles A, B, C ; 48 lettres transmises.}
\nf{Références}{ICUR II 6041 = ILCV +2969a.}
\nf{Provenance}{Rattachement selon l’EDB : \emph{Coem. Commodillae}.}
\nf{Support}{Plaque de marbre. Inscription gravée. Fiche EDB : Carlo Carletti, première version 2013.}
\ntext{\lnum{1}Apronian(a)e\\
\lnum{2}deposita in p(a)c(e)\\
\lnum{3}XV ⸢K⸣al(endas) I⸢a⸣n(uarias)\\
\lnum{4}cons(ulatu) d(o)m(ini) n(ostri) Arc(adi)\\
\lnum{5}et Baudo(nis) v(iri) c(larissimi)}
\nf{Témoins}{Export EDCS de 2020 : \emph{Apronian(a)e / deposita in p(a)c(e) / XV al(endas) In(uarias) / cons(ulatu) d(o)m(ini) n(ostri) Arc(adi) / et Baudo(nis) v(iri) c(larissimi)}. EDB (EDB8753) : texte concordant aux conventions de notation près (indice de similarité 0,968).}
\nf{Apparat}{⸢K⸣ : C gravé ; ⸢a⸣ : E gravé}
\nf{Traduction}{À Aproniana, déposée en paix le 15e jour avant les calendes de janvier, sous le consulat de notre seigneur Arcadius et de Bauto, clarissime.}
\nf{Langue}{L’éditeur signale 2 graphies corrigées : C gravé pour K ; E gravé pour a (e pour a).}
\nf{Chronologie}{Noms consulaires : Arcadius (lu) et Bauto (lu). Dates du jour : déposition : XV Kalendas Ianuarias = 18 décembre 385 (lu).}
\nf{Commentaire}{La formule in pace est transmise sous forme abrégée, sans restitution de lacune ; son développement reste éditorial.}
\nf{Bibliographie}{ICUR II 6041 ; ILCV +2969a.}
\nf{Statut}{Lecture secondaire (EDCS) concordante avec l’EDB : lecture stabilisée entre deux bases, non collationnée sur l’ICVR ni sur le monument.}
\end{notice}

\begin{notice}{280}{385}{ICUR n.s. IV, 9573 · EDCS-36101099 · EDB17962}
\label{not:280}
\nf{Datation}{385 — directe : nom consulaire lu}
\nf{Genre}{Épitaphe ; ensembles A, B, C ; 61 lettres transmises ; 48 \% restituées.}
\nf{Références}{ICUR IV 9573 (p 533).}
\nf{Provenance}{Rattachement selon l’EDB : \emph{Coem. Callisti pars inferior}.}
\nf{Support}{Plaque de marbre. Inscription gravée. Fiche EDB : Antonio Enrico Felle, première version 2013.}
\ntext{\lnum{1}bene mer]enti Astu[to(?) qui vixit]\\
\lnum{2}[annos] p(lus) m(inus) LX depo[situs - - -]\\
\lnum{3}[- - -] Fl(avio) Arcadio [et Bautone}
\nf{Témoins}{EDB (EDB17962 ; écarts à examiner) : \emph{[--- benemer]enti Astu[rio qui vixit] / [annos] p(lus) m(inus) LX depo[situs ---] / [---] Fl(avio) Arcadio [et Bautone ---] [texte tronqué dans la liste de résultats]}.}
\nf{Traduction}{Au bien méritant Astutus, qui vécut environ 60 ans ; déposé […] ; (sous le consulat de) Flavius Arcadius et Bauto.}
\nf{Langue}{4 abréviations développées par l’éditeur (dont Astuto(?), p(lus), m(inus), Fl(avio)).}
\nf{Chronologie}{Noms consulaires : Arcadius (lu) et Bauto (restitué). Âge déclaré : env. 60 a. (partiellement restitué).}
\nf{Bibliographie}{ICUR IV 9573 (p 533).}
\nf{Statut}{Lecture secondaire (EDCS) ; écarts avec l’EDB à arbitrer sur l’ICVR.}
\end{notice}

\begin{notice}{281}{385}{ICUR n.s. IV, 12247 · EDCS-37600682 · EDB5346}
\label{not:281}
\nf{Datation}{385 — directe : nom consulaire lu}
\nf{Genre}{Acte d’achat ou de réservation ; ensembles A, B ; 91 lettres transmises ; 50 \% restituées.}
\nf{Références}{ICUR IV 12247.}
\nf{Provenance}{Rattachement selon l’EDB : \emph{Coem. anonymum ad viam Ardeatinam}. Conservation selon l’EDB : \emph{Coem. anonymum ad viam Ardeatinam}.}
\nf{Support}{Plaque de marbre. Inscription gravée. Reproductions signalées par l’EDB : photographie PCAS (droits réservés ; non reproduites). Fiche EDB : Marida Pierno, première version 2013.}
\ntext{\lnum{1}Locus Rufin[ae quae]\\
\lnum{2}se ⸢v⸣i⸢v⸣a sibi c⸢o⸣[mpa]\\
\lnum{3}ra⸢v⸣it bisomu(m) [d(omino) n(ostro) Arcad]\\
\lnum{4}io et Baudone v(iro) [c(larissimo) cons(ulibus)]}
\nf{Témoins}{Export EDCS de 2020 : \emph{Locus Rufin[ae quae] / se ia sibi c[mpa]/rait bisomu(m) [d(omino) n(ostro) Arcad]/io et Baudone v(iro) [c(larissimo) cons(ulibus)]}. EDB (EDB5346) : texte concordant aux conventions de notation près (indice de similarité 0,947).}
\nf{Apparat}{⸢v⸣ : B gravé ; ⸢v⸣ : B gravé ; ⸢o⸣ : V gravé ; ⸢v⸣ : B gravé}
\nf{Traduction}{Emplacement de Rufina, qui, de son vivant, s’est acquis une tombe double, sous le consulat de notre seigneur Arcadius et de Bauto, clarissime.}
\nf{Langue}{L’éditeur signale 4 graphies corrigées : B gravé pour v (bêtacisme) ×3 ; V gravé pour o (confusion o/u).}
\nf{Chronologie}{Noms consulaires : Arcadius (restitué) et Bauto (lu). Ne sont pas des consuls : Locus Rufin[ae (propriétaire de la tombe).}
\nf{Bibliographie}{ICUR IV 12247.}
\nf{Statut}{Lecture secondaire (EDCS) concordante avec l’EDB : lecture stabilisée entre deux bases, non collationnée sur l’ICVR ni sur le monument. Plus de la moitié des lettres sont restituées. L’EDCS signale une photographie (non consultée).}
\end{notice}

\annee{386}

\begin{notice}{282}{386}{EDCS-16600245 · EDB11278 · EDH HD021839}
\label{not:282}
\nf{Statut}{Doublon établi : même âge au jour près (25 ans 2 mois 27 jours), même date (3 juillet 386), même coupe des lignes 1 à 3. L'EDH (HD021839) et l'EDB (EDB11278) rattachent le texte de l'AE 1906, 86 à ICUR II 6043. La notice est fusionnée dans le n° 288. L'EDH lit le nom \emph{An]iciae} là où l'ICVR et l'EDB lisent \emph{Euticiae} : la variante est donnée en apparat du n° 288.}
\nf{Références}{AE 1906, 86}
\ntext{\lnum{1}]\\
\lnum{2}iciae co(n)iugi quae vixit\\
\lnum{3}ann⸢o⸣s XXV m(enses) II d(ies) XXVII deposita\\
\lnum{4}V Nonas Iul(ias) in pace consulatu\\
\lnum{5}Honori\{o\} n(obilissimi) p(ueri) et Euodi\{o\}}
\nf{Traduction}{[…] à son épouse […]icia, qui vécut 25 ans 2 mois 27 jours ; déposée le 5e jour avant les nones de juillet, en paix, sous le consulat d’Honorius, très noble enfant, et d’Euodius.}
\nf{Voir}{la notice 288 (p.~\pageref{not:288}).}
\end{notice}

\begin{notice}{283}{386}{ICUR n.s. I, 1257 · EDCS-39101967 · EDB30485}
\label{not:283}
\nf{Datation}{386 — directe : nom consulaire lu}
\nf{Genre}{Épitaphe ; ensembles A, B, C ; 61 lettres transmises ; 56 \% restituées.}
\nf{Références}{ICUR I 1257.}
\nf{Provenance}{Rattachement selon l’EDB : \emph{Incertae originis} (le volume I des ICVR réunit les inscriptions dont la provenance est perdue). Conservation selon l’EDB : \emph{Musei Capitolini}.}
\nf{Support}{Plaque de marbre. 16 × 23 cm (h. × l.). Inscription gravée. Fiche EDB : Filippo Piazzolla, première version 2014.}
\ntext{\lnum{1}]C[- - -]\\
\lnum{2}[in pa]ce qui vix[it ann(os) - - -]\\
\lnum{3}[d]ep(ositus) Honor[io n(obilissimo) p(uero) et]\\
\lnum{4}[Eu]odio v(iro) c(larissimo) conss(ulibus) [- - -]}
\nf{Témoins}{EDB (EDB30485 ; écarts à examiner) : \emph{------ / [--- in pa]c̣e qui vix[it annos ---] / [--- d]ẹp(osit---) Honor[io n(obilissimo) puer(o) et] / [Ev]odio v(iro) c(larissimo) cons(ulibu)s [---?] [texte tronqué dans la liste de résultats]}.}
\nf{Traduction}{[…] en paix, qui vécut […] ans ; déposé sous le consulat d’Honorius, très noble enfant, et d’Euodius, clarissime […].}
\nf{Langue}{7 abréviations développées par l’éditeur (dont ann(os), dep(ositus), n(obilissimo), p(uero), v(iro), c(larissimo)).}
\nf{Chronologie}{Noms consulaires : Honorius (lu) et Euodius (lu en partie).}
\nf{Commentaire}{La formule in pace est attestée seulement par restitution.}
\nf{Bibliographie}{ICUR I 1257. Bibliographie complémentaire selon l’EDB : Ferrua 1979b, 31, n. 1257.}
\nf{Statut}{Lecture secondaire (EDCS) ; écarts avec l’EDB à arbitrer sur l’ICVR. Plus de la moitié des lettres sont restituées.}
\end{notice}

\begin{notice}{284}{386}{ICUR n.s. I, 1353 · EDCS-39000331 · EDB30696}
\label{not:284}
\nf{Datation}{386 — directe : lecture partielle d’un nom consulaire}
\nf{Genre}{Épitaphe ; ensembles A, B, C ; 105 lettres transmises ; 58 \% restituées.}
\nf{Références}{ICUR I 1353.}
\nf{Provenance}{Rattachement selon l’EDB : \emph{Incertae originis} (le volume I des ICVR réunit les inscriptions dont la provenance est perdue). Conservation selon l’EDB : \emph{Collegio Teutonico}.}
\nf{Support}{Plaque de marbre. 31 × 20 cm (h. × l.). Inscription gravée. Fiche EDB : Filippo Piazzolla, première version 2015.}
\ntext{\lnum{1}i]n pace IIII [- - -]\\
\lnum{2}[- - - L]aurenti [- - -]\\
\lnum{3}[- - - Fl(avio) Honorio n(obilissimo) p(uero) et Eu]\{b\}odio vv(iris) [cc(larissimis) cons(ulibus) - - -]\\
\lnum{4}[- - -]ae bene me[renti - - -]\\
\lnum{5}[- - -] Urbicae in pac\{a\}e [- - -]}
\nf{Témoins}{EDB (EDB30696 ; écarts à examiner) : \emph{------ / [dep(osit---) i]n pace IIII [---] / [--- L]aurenti [---] / [Fl(avio) Honorio n(obilissimo) p(uero) et Eu]bodio vv(iris) [cc(larissimis) cons(ulibu)s] / [---]ae beneme[renti ---] / [---] Urbicae in pacae [texte tronqué dans la liste de résultats]}.}
\nf{Apparat}{\{b\} : lettres superflues ; \{a\} : lettres superflues}
\nf{Traduction}{[…] en paix, 4 […] ; […] Laurentius […] ; […] sous le consulat de Flavius Honorius, très noble enfant, et d’Euodius, clarissime […] ; […] bien méritante […] ; […] à Urbica, en paix […].}
\nf{Langue}{Des lettres gravées sont jugées superflues par l’éditeur (entre accolades). 6 abréviations développées par l’éditeur (dont Fl(avio), n(obilissimo), p(uero), vv(iris), cc(larissimis), cons(ulibus)).}
\nf{Chronologie}{Noms consulaires : Honorius (restitué) et Euodius (lu en partie).}
\nf{Commentaire}{La formule in pace est conservée.}
\nf{Bibliographie}{ICUR I 1353.}
\nf{Statut}{Lecture secondaire (EDCS) ; écarts avec l’EDB à arbitrer sur l’ICVR. Plus de la moitié des lettres sont restituées.}
\end{notice}

\begin{notice}{285}{386}{ICUR n.s. I, 1443 · EDCS-39000424 · EDB31210}
\label{not:285}
\nf{Datation}{386 — directe : lecture partielle d’un nom consulaire}
\nf{Genre}{Épitaphe ; ensembles A, B, C ; 74 lettres transmises ; 81 \% restituées.}
\nf{Références}{ICUR I 1443.}
\nf{Provenance}{Rattachement selon l’EDB : \emph{Incertae originis} (le volume I des ICVR réunit les inscriptions dont la provenance est perdue). Conservation selon l’EDB : \emph{Musei Vaticani: Lapidario Cristiano}.}
\nf{Support}{Plaque de marbre. 25 × 16 cm (h. × l.). Inscription gravée. Fiche EDB : Filippo Piazzolla, première version 2015.}
\ntext{\lnum{1}a]nni[s - - -]\\
\lnum{2}[- - - co]ntra v[otum(?) - - -]\\
\lnum{3}[- - - Honorio n(obilissimo) p(uero) et Fl(avio) Eu]odio v(iro) [c(larissimo) cons(ulibus)]}
\nf{Témoins}{EDB (EDB31210) : texte concordant aux conventions de notation près (indice de similarité 0,927).}
\nf{Traduction}{[…] ans […] ; […] contre son vœu […] ; sous le consulat d’Honorius, très noble enfant, et de Flavius Euodius, clarissime.}
\nf{Langue}{7 abréviations développées par l’éditeur (dont votum(?), n(obilissimo), p(uero), Fl(avio), v(iro), c(larissimo)).}
\nf{Chronologie}{Noms consulaires : Honorius (restitué) et Euodius (lu en partie).}
\nf{Bibliographie}{ICUR I 1443. Bibliographie complémentaire selon l’EDB : Lega 2000, 28.}
\nf{Statut}{Lecture secondaire (EDCS) concordante avec l’EDB : lecture stabilisée entre deux bases, non collationnée sur l’ICVR ni sur le monument. Plus de la moitié des lettres sont restituées. L’EDCS signale une photographie (non consultée).}
\end{notice}

\begin{notice}{286}{386}{ICUR n.s. I, 2101 · EDCS-30600274 · EDB35609}
\label{not:286}
\nf{Datation}{386 — directe : nom consulaire lu}
\nf{Genre}{Épitaphe ; ensembles A, B, C ; 128 lettres transmises ; 59 \% restituées.}
\nf{Références}{ICUR I 2101 (p 491) = ICUR X 26645 (p 244) = ILCV +4249a.}
\nf{Provenance}{Rattachement selon l’EDB : \emph{Area subdialis coem. Pamphili et inscriptiones incerto loco repertae}. Conservation selon l’EDB : \emph{Musei Vaticani: Galleria Lapidaria}.}
\nf{Support}{Plaque de marbre. 39 × 58 cm (h. × l.). Inscription gravée. Fiche EDB : Filippo Piazzolla, première version 2015.}
\ntext{\lnum{1}Sabiniane qu[i vixit annos]\\
\lnum{2}LV depositus [- - -]\\
\lnum{3}fecit cum vi[rginia sua annos - - - consu]\\
\lnum{4}latu Honori[o n(obilissimo) p(uero) et Fl(avio) Euo]\\
\lnum{5}dio viro clar[issimo conss(ulibus)]}
\nf{Témoins}{EDB (EDB35609) : texte concordant aux conventions de notation près (indice de similarité 0,981).}
\nf{Traduction}{À Sabinianus, qui vécut 55 ans ; déposé […] ; il a vécu avec sa ‘vierge’ (épouse) […] ans ; sous le consulat d’Honorius, très noble enfant, et de Flavius Euodius, clarissime.}
\nf{Langue}{4 abréviations développées par l’éditeur (dont n(obilissimo), p(uero), Fl(avio), conss(ulibus)).}
\nf{Chronologie}{Noms consulaires : Honorius (lu) et Euodius (lu en partie). Âge déclaré : 55 a. (partiellement restitué).}
\nf{Bibliographie}{ICUR I 2101 (p 491) ; ICUR X 26645 (p 244) ; ILCV +4249a. Bibliographie complémentaire selon l’EDB : ICVR X, 26645 ; Lega 2000, 266.}
\nf{Statut}{Lecture secondaire (EDCS) concordante avec l’EDB : lecture stabilisée entre deux bases, non collationnée sur l’ICVR ni sur le monument. Plus de la moitié des lettres sont restituées.}
\end{notice}

\begin{notice}{287}{386}{ICUR n.s. II, 6042 · EDCS-34200668 · EDB11277}
\label{not:287}
\nf{Datation}{386 — directe : nom consulaire lu}
\nf{Genre}{Épitaphe ; ensembles A, B, C ; 132 lettres transmises ; 49 \% restituées.}
\nf{Références}{ICUR II 6042 = ILCV 4379 = AE 1905, 79.}
\nf{Provenance}{Rattachement selon l’EDB : \emph{Coem. Commodillae}.}
\nf{Support}{Plaque de marbre. Inscription gravée. Reproductions signalées par l’EDB : photographie PCAS (droits réservés ; non reproduites). Fiche EDB : Carlo Carletti, première version 2003.}
\ntext{\lnum{1}]s Honorio\\
\lnum{2}[n(obilissimo) p(uero) et Euodio] v(iro) c(larissimo) cons(ulibu)s\\
\lnum{3}[X Kal(endas) Sept(embres) di]e So⸢l⸣is\\
\lnum{4}[hora noctis - - - Lu]na XII signo\\
\lnum{5}[Capricor]nus\\
\lnum{6}[- - -]II ⸢d⸣(i)es XXX\\
\lnum{7}[- - -]ti s⸢e⸣ptimu(m)\\
\lnum{8}[- - -] ⸢i⸣a⸢c⸣et in pace\\
\lnum{9}[- - -]terus}
\nf{Témoins}{Export EDCS de 2020 : \emph{]s Honorio / [n(obilissimo) p(uero) et Euodio] v(iro) c(larissimo) cons(ulibu)s / [X Kal(endas) Sept(embres) di]e Sois / [hora noctis 3 Lu]na XII signo / [Capricor]nus / [3]II (i)es XXX / [3]ti sptimu(m) / [3] aet in pace / [3]terus}. EDB (EDB11277) : texte concordant aux conventions de notation près (indice de similarité 0,975).}
\nf{Apparat}{⸢l⸣ : I gravé ; ⸢d⸣ : O gravé ; ⸢e⸣ : I gravé ; ⸢i⸣ : B gravé ; ⸢c⸣ : Q gravé}
\nf{Traduction}{[…] sous le consulat d’Honorius, très noble enfant, et d’Euodius, clarissime, le 10e jour avant les calendes de septembre, un dimanche, à l’heure […] de la nuit, au 12e jour de la lune, sous le signe du Capricorne ; […] 30 jours ; […] le septième ; […] gît en paix ; […].}
\nf{Langue}{L’éditeur signale 5 graphies corrigées : I gravé pour l ; O gravé pour d ; I gravé pour e (confusion e/i) ; B gravé pour i ; Q gravé pour c.}
\nf{Chronologie}{Noms consulaires : Honorius (lu) et Euodius (restitué). Dates du jour : événement indéterminé : X Kalendas Septembres = 23 août 386 (restitué).}
\nf{Commentaire}{La formule in pace est conservée. Commentaire développé : dossier p.~\pageref{dos:287}.}
\nf{Bibliographie}{ICUR II 6042 ; ILCV 4379 ; AE 1905, 79.}
\nf{Statut}{Lecture secondaire (EDCS) concordante avec l’EDB : lecture stabilisée entre deux bases, non collationnée sur l’ICVR ni sur le monument. L’EDCS signale une photographie (non consultée).}
\end{notice}

\begin{notice}{288}{386}{ICUR n.s. II, 6043 · EDCS-34200669 · EDB11278 · EDH HD021839}
\label{not:288}
\nf{Datation}{386 — directe : nom consulaire lu}
\nf{Genre}{Épitaphe ; ensembles A, B, C ; 100 lettres transmises.}
\nf{Références}{ICUR II 6043 = ILCV 2978c.}
\nf{Provenance}{Rattachement selon l’EDB : \emph{Coem. Commodillae}.}
\nf{Support}{Plaque de marbre. Inscription gravée. Fiche EDB : Carlo Carletti, première version 2015.}
\ntext{\lnum{1}Euticiae coniugi quae vixit\\
\lnum{2}ann⸢o⸣s XXV m(enses) II d(ies) XXVII deposita\\
\lnum{3}V Nonas Iul(ias) in pace consulatum\\
\lnum{4}Fl(avio) Honorio n(obilissimo) p(uero) et Euodio v(iro) c(larissimo) cons(ulibu)s}
\nf{Témoins}{Export EDCS de 2020 : \emph{Euticiae coniugi quae vixit / anns XXV m(enses) II d(ies) XXVII deposita / V Nonas Iul(ias) in pace consulatum / Fl(avio) Honorio n(obilissimo) p(uero) et Euodio v(iro) c(larissimo) cons(ulibu)s}. EDB (EDB11278) : texte concordant aux conventions de notation près (indice de similarité 0,992). Le texte publié dans l'AE 1906, 86 (ancien n° 282) est le même monument, sous une forme plus fragmentaire : voir la notice 282. Pour le nom, l'EDH lit \emph{An]iciae}, l'ICVR et l'EDB \emph{Euticiae}. Absorbe la notice 282 (voir cette notice).}
\nf{Apparat}{⸢o⸣ : V gravé}
\nf{Traduction}{À Euticia, son épouse, qui vécut 25 ans 2 mois 27 jours ; déposée le 5e jour avant les nones de juillet, en paix, sous le consulat de Flavius Honorius, très noble enfant, et d’Euodius, clarissime.}
\nf{Langue}{L’éditeur signale une graphie corrigée : V gravé pour o (confusion o/u).}
\nf{Chronologie}{Noms consulaires : Honorius (lu) et Euodius (lu). Dates du jour : déposition : V Nonas Iulias = 3 juillet 386 (lu). Âge déclaré : 25 a. 2 m. 27 j. (lu).}
\nf{Commentaire}{La formule in pace est conservée.}
\nf{Bibliographie}{ICUR II 6043 ; ILCV 2978c.}
\nf{Statut}{Lecture secondaire (EDCS) concordante avec l’EDB : lecture stabilisée entre deux bases, non collationnée sur l’ICVR ni sur le monument.}
\end{notice}

\begin{notice}{289}{386}{ICUR n.s. III, 8154 · EDCS-35700183 · EDB24676}
\label{not:289}
\nf{Datation}{386 — directe : nom consulaire lu}
\nf{Genre}{Épitaphe ; ensembles A, B, C ; 59 lettres transmises ; 46 \% restituées.}
\nf{Références}{ICUR III 8154 = ILCV +1119 = Mossong-2022, 00131.}
\nf{Provenance}{Rattachement selon l’EDB : \emph{Basilica ss. Nerei et Achillei in coem. Domitillae}. Conservation selon l’EDB : \emph{Basilica ss. Nerei et Achillei in coem. Domitillae}.}
\nf{Support}{Plaque de marbre. 50 × 48 cm (h. × l.). Inscription gravée. Reproductions signalées par l’EDB : photographie PCAS (droits réservés ; non reproduites). Fiche EDB : Carlo Carletti / Antonio Enrico Felle, première version 2003.}
\ntext{\lnum{1}]IISTE[- - -]\\
\lnum{2}[- - - s]acerdot[- - -]\\
\lnum{3}[- - -]onis presb[yteri - - -]\\
\lnum{4}[- - - Fl(avio) H]onorio et Eu⸢o⸣[dio conss(ulibus)}
\nf{Témoins}{Export EDCS de 2020 : \emph{]IISTE[3] / [3 s]acerdot[3] / [3]onis presb[yteri 3] / [3 Fl(avio) H]onorio et Eu[dio conss(ulibus)}. EDB (EDB24676) : texte concordant aux conventions de notation près (indice de similarité 0,938).}
\nf{Apparat}{⸢o⸣ : V gravé}
\nf{Traduction}{[…] prêtre (sacerdos) […] ; […] du prêtre (presbyter) […] ; sous le consulat de Flavius Honorius et d’Euodius.}
\nf{Langue}{L’éditeur signale une graphie corrigée : V gravé pour o (confusion o/u).}
\nf{Chronologie}{Noms consulaires : Honorius (lu) et Euodius (lu en partie).}
\nf{Commentaire}{Fonction ou métier mentionné : sacerdot, presbyteri.}
\nf{Bibliographie}{ICUR III 8154 ; ILCV +1119 ; Mossong-2022, 00131. Bibliographie complémentaire selon l’EDB : IC I Supplementum, 1751 ; ILCV, 1119 adn.}
\nf{Statut}{Lecture secondaire (EDCS) concordante avec l’EDB : lecture stabilisée entre deux bases, non collationnée sur l’ICVR ni sur le monument. L’EDCS signale une photographie (non consultée).}
\end{notice}

\begin{notice}{290}{386}{ICUR n.s. IV, 11114 · EDCS-37200396 · EDB39576}
\label{not:290}
\nf{Datation}{386 — directe : consulat restitué}
\nf{Genre}{Épitaphe ; ensembles A, B, C ; 77 lettres transmises ; 74 \% restituées.}
\nf{Références}{ICUR IV 11114.}
\nf{Provenance}{Rattachement selon l’EDB : \emph{Coem. Callisti area subdialis et incerti loci}. Conservation selon l’EDB : \emph{Musei Vaticani: Lapidario Cristiano}.}
\nf{Support}{Plaque de marbre. 31 × 22 cm (h. × l.). Inscription gravée. Fiche EDB : Marida Pierno, première version 2016.}
\ntext{\lnum{1}]ina de[posita - - -]\\
\lnum{2}[- - -]e Fl(avio) Ho[norio n(obilissimo) p(uero) et Fl(avio) Euodio conss(ulibus) - - -]\\
\lnum{3}[- - - ann(os)] XXIIII dies [}
\nf{Témoins}{EDB (EDB39576 ; divergence forte) : \emph{[---]ina de[posita ---] / [---]e Fl(avio) Ho[norio ---] / [---] XXIIII dies [---]}.}
\nf{Traduction}{[…]ina, déposée […] ; sous le consulat de Flavius Honorius, très noble enfant, et de Flavius Euodius […] ; […] 24 ans […] jours.}
\nf{Langue}{6 abréviations développées par l’éditeur (dont Fl(avio), n(obilissimo), p(uero), conss(ulibus), ann(os)).}
\nf{Chronologie}{Noms consulaires : Honorius (restitué) et Euodius (restitué). L'EDB (EDB39576) s'arrête à \emph{Ho[norio ---]}. Ferrua, cité par l'EDB, note que seul le premier consulat d'Honorius (386) le nomme \emph{Flavius} : l'année 386 est argumentée par \emph{Fl(avio) Ho-}, mais le reste de la formule est restitué.}
\nf{Bibliographie}{ICUR IV 11114. Bibliographie complémentaire selon l’EDB : Lega 2000 ; ICVR IV, 11171 ; ICVR IV, 11137.}
\nf{Statut}{Lecture secondaire (EDCS) ; divergence forte avec l’EDB : le texte ne doit pas être exploité sans contrôle de l’édition imprimée. Plus de la moitié des lettres sont restituées.}
\end{notice}

\annee{387}

\begin{notice}{291}{387}{ICUR n.s. I, 1258 · EDCS-39101968 · EDB30488}
\label{not:291}
\nf{Datation}{387 — directe : nom consulaire lu}
\nf{Genre}{Épitaphe ; ensembles A, B, C ; 134 lettres transmises ; 46 \% restituées.}
\nf{Références}{ICUR I 1258.}
\nf{Provenance}{Rattachement selon l’EDB : \emph{Incertae originis} (le volume I des ICVR réunit les inscriptions dont la provenance est perdue). Conservation selon l’EDB : \emph{Musei Capitolini}.}
\nf{Support}{Plaque de marbre. 32 × 16 cm (h. × l.). Inscription gravée. Fiche EDB : Filippo Piazzolla, première version 2014.}
\ntext{\lnum{1}bene mere]nti in pace q[uae vix(it) ann(os) - - - quae]\\
\lnum{2}[fecit cum ma]rit⸢o⸣ ann[os - - -] dies X[- - - recessit]\\
\lnum{3}[in pace - - - K]al(endas) Ian(uarias) et deposita est XIII [Kal(endas) - - -]\\
\lnum{4}[Valentinia]no III et Eutropio conss(ulibus) [}
\nf{Témoins}{Export EDCS de 2020 : \emph{bene mere]nti in pace q[uae vix(it) ann(os) 3 quae] / [fecit cum ma]rit ann[os 3] dies X[3 recessit] / [in pace 3 K]al(endas) Ian(uarias) et deposita est XIII [Kal(endas) 3] / [Valentinia]no III et Eutropio conss(ulibus) [}. EDB (EDB30488 ; écarts à examiner) : \emph{[--- benemere]nti in pace q[ue vixit annos --- quae] / [fecit cum ma]ritum ann[os ---] ḍịẹṣ X [recessit] / [in pace --- k]al(endas) ian(uarias) et deposita est XIII [kal(endas) easdem] / [Valentinia]no III et Eutropio cons(ulibu)s [---?] [texte tronqué dans la liste de résultats]}.}
\nf{Apparat}{⸢o⸣ : VM gravé}
\nf{Traduction}{Bien méritante, en paix, qui vécut […] ans, qui a vécu avec son mari […] ans […] jours ; elle s’en alla en paix […] calendes de janvier et a été déposée le 13e jour avant les calendes […], sous le consulat de Valentinien pour la 3e fois et d’Eutropius.}
\nf{Langue}{L’éditeur signale une graphie corrigée : VM gravé pour o.}
\nf{Chronologie}{Noms consulaires : Valentinianus (restitué) et Eutropius (lu). Dates du jour : décès : … Kalendas Ianuarias (partiellement restitué).}
\nf{Commentaire}{La formule in pace est conservée.}
\nf{Bibliographie}{ICUR I 1258. Bibliographie complémentaire selon l’EDB : ILCV, 4219 adn. ; Molisani 1973, 110.}
\nf{Statut}{Lecture secondaire (EDCS) ; écarts avec l’EDB à arbitrer sur l’ICVR.}
\end{notice}

\begin{notice}{292}{387}{ICUR n.s. I, 1444 · EDCS-39000425 · EDB31208}
\label{not:292}
\nf{Datation}{387 — directe : nom consulaire lu}
\nf{Genre}{Épitaphe ; ensembles A, B, C ; 56 lettres transmises.}
\nf{Références}{ICUR I 1444.}
\nf{Provenance}{Rattachement selon l’EDB : \emph{Incertae originis} (le volume I des ICVR réunit les inscriptions dont la provenance est perdue). Conservation selon l’EDB : \emph{Musei Vaticani: Lapidario Cristiano}.}
\nf{Support}{Plaque de marbre. 21 × 84 cm (h. × l.). Inscription gravée. Fiche EDB : Filippo Piazzolla, première version 2018.}
\ntext{\lnum{1}]\\
\lnum{2}qu(a)e vixit an(nos) II m(enses) III d(ies) VII dep(osita)\\
\lnum{3}XV Kal(endas) Octob(res) Valenti(ni)ano III et Eutropio}
\nf{Témoins}{EDB (EDB31208) : texte concordant aux conventions de notation près (indice de similarité 0,986).}
\nf{Traduction}{[…] qui vécut 2 ans 3 mois 7 jours ; déposée le 15e jour avant les calendes d’octobre, sous (le consulat de) Valentinien pour la 3e fois et Eutropius.}
\nf{Langue}{8 abréviations développées par l’éditeur (dont qu(a)e, an(nos), m(enses), d(ies), dep(osita), Kal(endas)).}
\nf{Chronologie}{Noms consulaires : Valentinianus (lu) et Eutropius (lu). Dates du jour : déposition : XV Kalendas Octobres = 17 septembre 387 (lu). Âge déclaré : 2 a. 3 m. 7 j. (lu).}
\nf{Bibliographie}{ICUR I 1444. Bibliographie complémentaire selon l’EDB : Ferrua 1975c, 144 fig. 8 ; Ferrua 1979b, 37, n. 1444 ; Lega 2000, 251.}
\nf{Statut}{Lecture secondaire (EDCS) concordante avec l’EDB : lecture stabilisée entre deux bases, non collationnée sur l’ICVR ni sur le monument. L’EDCS signale une photographie (non consultée).}
\end{notice}

\begin{notice}{293}{387}{ICUR n.s. I, 3204 · EDCS-40300344 · EDB33154}
\label{not:293}
\nf{Datation}{387 — directe : nom consulaire lu}
\nf{Genre}{Épitaphe ; ensembles A, B, C ; 47 lettres transmises ; 68 \% restituées.}
\nf{Références}{ICUR I 3204.}
\nf{Provenance}{Rattachement selon l’EDB : \emph{Incertae originis} (le volume I des ICVR réunit les inscriptions dont la provenance est perdue).}
\nf{Support}{Support non précisé par l’EDB. Fiche EDB : Filippo Piazzolla, première version 2015.}
\ntext{\lnum{1}DD(ominis) nn(ostris) [Valentiniano Aug(usto) III]\\
\lnum{2}et Eutropio [cons(ulibus) - - -]\\
\lnum{3}V[}
\nf{Témoins}{EDB (EDB33154 ; écarts à examiner) : \emph{dd(ominis) nn(ostris) [Valentiniano Aug(usto) III] / et Eutropio [cons(ulibu)s ---] / V[---] [texte tronqué dans la liste de résultats]}.}
\nf{Traduction}{Sous le consulat de nos seigneurs Valentinien Auguste pour la 3e fois et Eutropius […] 5.}
\nf{Langue}{4 abréviations développées par l’éditeur (dont DD(ominis), nn(ostris), Aug(usto), cons(ulibus)).}
\nf{Chronologie}{Noms consulaires : Valentinianus (restitué) et Eutropius (lu).}
\nf{Bibliographie}{ICUR I 3204.}
\nf{Statut}{Lecture secondaire (EDCS) ; écarts avec l’EDB à arbitrer sur l’ICVR. Plus de la moitié des lettres sont restituées.}
\end{notice}

\begin{notice}{294}{387}{ICUR n.s. III, 8155 · EDCS-35700184 · EDB24677}
\label{not:294}
\nf{Datation}{387 — directe : lecture partielle d’un nom consulaire}
\nf{Genre}{Épitaphe ; ensembles A, B, C ; 52 lettres transmises ; 44 \% restituées.}
\nf{Références}{ICUR III 8155.}
\nf{Provenance}{Rattachement selon l’EDB : \emph{Basilica ss. Nerei et Achillei in coem. Domitillae}. Conservation selon l’EDB : \emph{Basilica ss. Nerei et Achillei in coem. Domitillae}.}
\nf{Support}{Plaque de marbre. 10 × 36 × 2 cm (h. × l. × ép.). Inscription gravée. Fiche EDB : Carlo Carletti / Antonio Enrico Felle, première version 2003.}
\ntext{\lnum{1}plus minu]s anno[s - - -]\\
\lnum{2}[menses - - - Va]lentiniano\\
\lnum{3}[Aug(usto) III et Eu]tropio}
\nf{Témoins}{EDB (EDB24677) : texte concordant aux conventions de notation près (indice de similarité 0,952).}
\nf{Traduction}{[…] environ […] ans […] mois ; (sous le consulat de) Valentinien Auguste pour la 3e fois et Eutropius.}
\nf{Langue}{1 abréviation développée par l’éditeur (dont Aug(usto)).}
\nf{Chronologie}{Noms consulaires : Valentinianus (lu en partie) et Eutropius (lu en partie).}
\nf{Bibliographie}{ICUR III 8155.}
\nf{Statut}{Lecture secondaire (EDCS) concordante avec l’EDB : lecture stabilisée entre deux bases, non collationnée sur l’ICVR ni sur le monument. L’EDCS signale une photographie (non consultée).}
\end{notice}

\annee{388}

\begin{notice}{295}{388}{ICUR n.s. VII, 17494 · EDCS-03900101 · EDB21373}
\label{not:295}
\nf{Datation}{388 — directe : nom consulaire lu}
\nf{Genre}{Épitaphe ; ensembles A, B, C ; 91 lettres transmises.}
\nf{Références}{Anagni 150 = ICUR VII 17494 = ILCV 3306 (em).}
\nf{Provenance}{Rattachement selon l’EDB : \emph{Coem. Cyriacae ad s.Laurentium}. Conservation selon l’EDB : \emph{Anagni (FR)   ; monastero delle suore Cistercensi della Carità}.}
\nf{Support}{Plaque de marbre. 13 × 38 cm (h. × l.). Inscription gravée. Reproductions signalées par l’EDB : reproduction, Perret 1851 et reproduction, Solin et alii 1996 (droits réservés ; non reproduites). Fiche EDB : Anita Rocco, première version 2015.}
\ntext{\lnum{1}D(omino) n(ostro) Magno Maximo Aug(usto) II conss(ule)\\
\lnum{2}III Idus Maias fatum fecit Leo et\\
\lnum{3}depositus pridie Idus Maias bene\\
\lnum{4}merenti in pace}
\nf{Témoins}{EDB (EDB21373) : texte concordant aux conventions de notation près (indice de similarité 0,996).}
\nf{Traduction}{Sous le consulat de notre seigneur Magnus Maximus Auguste pour la 2e fois, le 3e jour avant les ides de mai, Leo accomplit son destin et fut déposé la veille des ides de mai ; bien méritant, en paix.}
\nf{Langue}{4 abréviations développées par l’éditeur (dont D(omino), n(ostro), Aug(usto), conss(ule)).}
\nf{Chronologie}{Noms consulaires : Magnus Maximus (lu). Dates du jour : événement indéterminé : III Idus Maias = 13 mai 388 (lu) ; déposition : pridie Idus Maias = 14 mai 388 (lu).}
\nf{Commentaire}{La formule in pace est conservée.}
\nf{Bibliographie}{Anagni 150 ; ICUR VII 17494 ; ILCV 3306 (em). Bibliographie complémentaire selon l’EDB : Solin et alii 1996, 94, n. 150 con foto.}
\nf{Statut}{Lecture secondaire (EDCS) concordante avec l’EDB : lecture stabilisée entre deux bases, non collationnée sur l’ICVR ni sur le monument. L’EDCS signale une photographie (non consultée).}
\end{notice}

\begin{notice}{296}{388}{ICUR n.s. VII, 17495 · EDCS-19301140 · EDB21374}
\label{not:296}
\nf{Datation}{388 — directe : nom consulaire lu}
\nf{Genre}{Épitaphe ; ensembles A, B, C ; 83 lettres transmises.}
\nf{Références}{CIL VI 9597 (p 3470) = ICUR VII 17495 = ILCV 610 = ICVaticano p 310 = Gummerus-01, 00098 = AE 1997, 166.}
\nf{Provenance}{Rattachement selon l’EDB : \emph{Coem. Cyriacae ad s.Laurentium}. Conservation selon l’EDB : \emph{Musei Vaticani: Lapidario Cristiano}.}
\nf{Support}{Plaque de marbre. 22 × 22 × 2,5 cm (h. × l. × ép.). Inscription gravée. Reproductions signalées par l’EDB : reproduction, Perret 1851 (droits réservés ; non reproduites). Fiche EDB : Miriam Ramunni, première version 2015.}
\ntext{\lnum{1}Rapetiga me\\
\lnum{2}dicus civis\\
\lnum{3}Hispanus qui\\
\lnum{4}vix⸢it⸣ in p(ace) ann(os) p(lus) m(inus)\\
\lnum{5}XXV hoc pater Ni\\
\lnum{6}ce⸢t⸣ius fecit d(omino) n(ostro) Ma(gno)\\
\lnum{7}Maximo Aug(usto) II}
\nf{Témoins}{Export EDCS de 2020 : \emph{Rapetiga me/dicus civis / Hispanus qui / vix in p(ace) ann(os) p(lus) m(inus) / XXV hoc pater Ni/ceius fecit d(omino) n(ostro) Ma(gno) / Maximo Aug(usto) II}. EDB (EDB21374) : texte concordant aux conventions de notation près (indice de similarité 0,986).}
\nf{Apparat}{⸢it⸣ : TI gravé ; ⸢t⸣ : I gravé}
\nf{Traduction}{Rapetiga (?), médecin, citoyen hispanique, qui vécut en paix environ 25 ans ; son père Nicetius a fait ceci, sous (le consulat de) notre seigneur Magnus Maximus Auguste pour la 2e fois.}
\nf{Langue}{L’éditeur signale 2 graphies corrigées : TI gravé pour it ; I gravé pour t.}
\nf{Chronologie}{Noms consulaires : Magnus Maximus (lu). Âge déclaré : env. 25 a. (lu).}
\nf{Commentaire}{La formule in pace est transmise sous forme abrégée, sans restitution de lacune ; son développement reste éditorial. Fonction ou métier mentionné : medicus.}
\nf{Bibliographie}{CIL VI 9597 (p 3470) ; ICUR VII 17495 ; ILCV 610 ; ICVaticano p 310 ; Gummerus-01, 00098 ; AE 1997, 166. Bibliographie complémentaire selon l’EDB : Di Stefano Manzella 1997, 318 n. 3.11.2 con foto (C. Ricci) ; Lega 2000 ; Bisconti 2000a, 236 n.XVc1.5 fig.82 ; CIL VI, 9597.}
\nf{Statut}{Lecture secondaire (EDCS) concordante avec l’EDB : lecture stabilisée entre deux bases, non collationnée sur l’ICVR ni sur le monument. L’EDCS signale une photographie (non consultée).}
\end{notice}

\begin{notice}{297}{388}{ICUR n.s. I, 3206 · EDCS-33500351 · EDB35685}
\label{not:297}
\nf{Datation}{388 — directe : nom consulaire lu}
\nf{Genre}{Épitaphe ; ensembles A, B, C ; 126 lettres transmises.}
\nf{Références}{ICUR I 3206 = ICUR VIII 20718 = ILCV 4219a.}
\nf{Provenance}{Rattachement selon l’EDB : \emph{Coem. Nicomedis}.}
\nf{Support}{Support non précisé par l’EDB. Fiche EDB : Filippo Piazzolla, première version 2015.}
\ntext{\lnum{1}Dulcissimae filiae Adeodatae\\
\lnum{2}qu(a)e vixit annos XVIII\\
\lnum{3}et m(en)s(es) V\\
\lnum{4}deposita XVI Kal(endas) F(e)br(uarias)\\
\lnum{5}cons(ule)\\
\lnum{6}Magno Maximo Augusto\\
\lnum{7}pater Celestinus fecit\\
\lnum{8}et cum marito f(ecit) mens(es) XI}
\nf{Témoins}{EDB (EDB35685) : texte concordant aux conventions de notation près (indice de similarité 0,99).}
\nf{Traduction}{À sa fille très douce Adeodata, qui vécut 18 ans et 5 mois ; déposée le 16e jour avant les calendes de février, sous le consulat de Magnus Maximus Auguste ; son père Caelestinus a fait (ce monument) ; elle a vécu avec son mari 11 mois.}
\nf{Langue}{9 abréviations développées par l’éditeur (dont qu(a)e, m(en)s, (es), Kal(endas), F(e)br, (uarias)).}
\nf{Chronologie}{Noms consulaires : Magnus Maximus (lu). Dates du jour : déposition : XVI Kalendas Februarias = 17 janvier 388 (lu). Âge déclaré : 18 a. 5 m. (lu).}
\nf{Commentaire}{Vocabulaire religieux : anthroponyme (niveau 6) — adeodatae (anthroponyme).}
\nf{Bibliographie}{ICUR I 3206 ; ICUR VIII 20718 ; ILCV 4219a. Bibliographie complémentaire selon l’EDB : ICVR I, 3206.}
\nf{Statut}{Lecture secondaire (EDCS) concordante avec l’EDB : lecture stabilisée entre deux bases, non collationnée sur l’ICVR ni sur le monument.}
\end{notice}

\begin{notice}{298}{388}{ICUR n.s. II, 5996 · EDCS-38100824 · EDB21298}
\label{not:298}
\nf{Datation}{388 — directe : nom consulaire lu}
\nf{Genre}{Épitaphe ; ensembles A, B, C ; 138 lettres transmises.}
\nf{Références}{ICUR II 5996 = ILCV 4623 (em).}
\nf{Provenance}{Rattachement selon l’EDB : \emph{Coem. s.Theclae}. Conservation selon l’EDB : \emph{Musei Vaticani: Lapidario Cristiano}.}
\nf{Support}{Plaque de marbre. 70 × 40 cm (h. × l.). Inscription gravée. Reproductions signalées par l’EDB : reproduction, Perret 1851 (droits réservés ; non reproduites). Fiche EDB : Maria Cristina Grisanzio, première version 2013.}
\ntext{\lnum{1}Hic posita est anima dulcis\\
\lnum{2}innocens sapiens et pulc(h)ra nomine\\
\lnum{3}Quiriace qu(a)e vixit annos III m(enses) III d(ies) VIII\\
\lnum{4}d(e)p(osita) in pace IIII Id(us) Ian(uarias) cons(ulibu)s d(omino) n(ostro) T(h)eodosio Aug(usto) II\\
\lnum{5}et Merobaude v(iro) c(larissimo) III}
\nf{Témoins}{EDB (EDB21298) : texte concordant aux conventions de notation près (indice de similarité 0,977).}
\nf{Traduction}{Ici a été placée une âme douce, innocente, sage et belle, nommée Quiriacè, qui vécut 3 ans 3 mois 8 jours ; déposée en paix le 4e jour avant les ides de janvier, sous le consulat de notre seigneur Théodose Auguste pour la 2e fois et de Merobaudes, clarissime, pour la 3e fois.}
\nf{Langue}{15 abréviations développées par l’éditeur (dont pulc(h)ra, qu(a)e, m(enses), d(ies), d(e)p, (osita)).}
\nf{Chronologie}{Noms consulaires : Theodosius (lu) et Merobaudes (lu). Dates du jour : déposition : IIII Idus Ianuarias = 10 janvier 388 (lu). Âge déclaré : 3 a. 3 m. 8 j. (lu).}
\nf{Commentaire}{La formule in pace est conservée.}
\nf{Bibliographie}{ICUR II 5996 ; ILCV 4623 (em). Bibliographie complémentaire selon l’EDB : ILCV, 4623.}
\nf{Statut}{Lecture secondaire (EDCS) concordante avec l’EDB : lecture stabilisée entre deux bases, non collationnée sur l’ICVR ni sur le monument. L’EDCS signale une photographie (non consultée).}
\end{notice}

\begin{notice}{299}{388}{ICUR n.s. II, 6044 · EDCS-34200670 · EDB11279}
\label{not:299}
\nf{Datation}{388 — directe : nom consulaire lu}
\nf{Genre}{Épitaphe ; ensembles A, B, C ; 70 lettres transmises.}
\nf{Références}{ICUR II 6044 = ILCV 2925a.}
\nf{Provenance}{Rattachement selon l’EDB : \emph{Coem. Commodillae}.}
\nf{Support}{Plaque de marbre. Inscription gravée. Reproductions signalées par l’EDB : photographie PCAS (droits réservés ; non reproduites). Fiche EDB : Carlo Carletti, première version 2013.}
\ntext{\lnum{1}Hic est posita puella\\
\lnum{2}Perseberatia qu(a)e vixit\\
\lnum{3}an(num) I m(enses) IIII d(ies) XI conss(ulatu) T(h)e⸢o⸣dos(i)\\
\lnum{4}M\{o\}ero⸢b⸣a⸢u⸣dis}
\nf{Témoins}{Export EDCS de 2020 : \emph{Hic est posita puella / Perseberatia qu(a)e vixit / an(num) I m(enses) IIII d(ies) XI conss(ulatu) T(h)edos(i) / M\{o\}eroadis}. EDB (EDB11279) : texte concordant aux conventions de notation près (indice de similarité 0,959).}
\nf{Apparat}{⸢o⸣ : V gravé ; ⸢b⸣ : P gravé ; ⸢u⸣ : O gravé ; \{o\} : lettres superflues}
\nf{Traduction}{Ici a été placée la fillette Perseverantia, qui vécut 1 an 4 mois 11 jours, sous le consulat de Théodose et de Merobaudes.}
\nf{Langue}{L’éditeur signale 3 graphies corrigées : V gravé pour o (confusion o/u) ; P gravé pour b ; O gravé pour u (confusion o/u). Des lettres gravées sont jugées superflues par l’éditeur (entre accolades). 7 abréviations développées par l’éditeur (dont qu(a)e, an(num), m(enses), d(ies), conss(ulatu), T(h)e).}
\nf{Chronologie}{Noms consulaires : Theodosius (lu) et Merobaudes (lu). Âge déclaré : 1 a. 4 m. 11 j. (lu).}
\nf{Bibliographie}{ICUR II 6044 ; ILCV 2925a.}
\nf{Statut}{Lecture secondaire (EDCS) concordante avec l’EDB : lecture stabilisée entre deux bases, non collationnée sur l’ICVR ni sur le monument. L’EDCS signale une photographie (non consultée).}
\end{notice}

\begin{notice}{300}{388}{ICUR n.s. II, 6045 · EDCS-34200671 · EDB11280}
\label{not:300}
\nf{Datation}{388 — directe : nom consulaire lu}
\nf{Genre}{Épitaphe ; ensembles A, B, C ; 90 lettres transmises.}
\nf{Références}{ICUR II 6045 = ILCV 2795c (add, em) = ILCV 4146d (add, em) = AE 1905, 80.}
\nf{Provenance}{Rattachement selon l’EDB : \emph{Coem. Commodillae}.}
\nf{Support}{Plaque de marbre. Inscription gravée. Signes décrits par l’EDB : croix monogrammatique avec alpha et oméga. Fiche EDB : Carlo Carletti, première version 2013.}
\ntext{\lnum{1}Mag(no) Maximo Aug(usto)\\
\lnum{2}II cons(ule) Dionisus\\
\lnum{3}et Rufina se vivos(!) fece\\
\lnum{4}runt et defuncta est\\
\lnum{5}VIIII Kal(endas) Aug(ustas) bene\\
\lnum{6}merenti in pace}
\nf{Témoins}{EDB (EDB11280) : texte concordant aux conventions de notation près (indice de similarité 1).}
\nf{Apparat}{(!) : forme signalée par l’éditeur}
\nf{Traduction}{Sous le consulat de Magnus Maximus Auguste pour la 2e fois, Dionysus et Rufina se sont fait (ce tombeau) de leur vivant ; et elle est morte le 9e jour avant les calendes d’août ; bien méritante, en paix.}
\nf{Langue}{L’éditeur signale au moins une forme anormale (sic). 6 abréviations développées par l’éditeur (dont Mag(no), Aug(usto), cons(ule), vivos(!), Kal(endas), Aug(ustas)).}
\nf{Chronologie}{Noms consulaires : Magnus Maximus (lu) et Rufinus (lu). Dates du jour : décès : VIIII Kalendas Augustas = 24 juillet 388 (lu).}
\nf{Commentaire}{La formule in pace est conservée.}
\nf{Bibliographie}{ICUR II 6045 ; ILCV 2795c (add, em) ; ILCV 4146d (add, em) ; AE 1905, 80.}
\nf{Statut}{Lecture secondaire (EDCS) concordante avec l’EDB : lecture stabilisée entre deux bases, non collationnée sur l’ICVR ni sur le monument.}
\end{notice}

\begin{notice}{301}{388}{ICUR n.s. II, 6046 · EDCS-34200672 · EDB11281}
\label{not:301}
\nf{Datation}{388 — directe : nom consulaire lu}
\nf{Genre}{Épitaphe ; ensembles A, B, C ; 62 lettres transmises.}
\nf{Références}{ICUR II 6046.}
\nf{Provenance}{Rattachement selon l’EDB : \emph{Coem. Commodillae}.}
\nf{Support}{Plaque de marbre. Inscription gravée. Fiche EDB : Carlo Carletti, première version 2013.}
\ntext{\lnum{1}Felicitas qu(a)e ⸢v⸣ix(it)\\
\lnum{2}an(nos) III m(enses) III et d(ies) XIII\\
\lnum{3}d(e)p(osita) IIIX K(a)l(endas) Feb(ruarias) post c(onsulatum)\\
\lnum{4}Valentiniano(!) III}
\nf{Témoins}{Export EDCS de 2020 : \emph{Felicitas qu(a)e ix(it) / an(nos) III m(enses) III et d(ies) XIII / d(e)p(osita) IIIX K(a)l(endas) Feb(ruarias) post c(onsulatum) / Valentiniano(!) III}. EDB (EDB11281) : texte concordant aux conventions de notation près (indice de similarité 0,975).}
\nf{Apparat}{⸢v⸣ : B gravé ; (!) : forme signalée par l’éditeur}
\nf{Traduction}{Felicitas, qui vécut 3 ans 3 mois et 13 jours ; déposée le 8e (?) jour avant les calendes de février, après le consulat de Valentinien pour la 3e fois.}
\nf{Langue}{L’éditeur signale une graphie corrigée : B gravé pour v (bêtacisme). L’éditeur signale au moins une forme anormale (sic). 12 abréviations développées par l’éditeur (dont qu(a)e, ix(it), an(nos), m(enses), d(ies), d(e)p).}
\nf{Chronologie}{Noms consulaires : Valentinianus (lu). Datation par post-consulat. Dates du jour : déposition : IIIX Kalendas Februarias = 22 janvier 388 (lu). Âge déclaré : 3 a. 3 m. 13 j. (lu).}
\nf{Bibliographie}{ICUR II 6046.}
\nf{Statut}{Lecture secondaire (EDCS) concordante avec l’EDB : lecture stabilisée entre deux bases, non collationnée sur l’ICVR ni sur le monument.}
\end{notice}

\begin{notice}{302}{388}{ICUR n.s. IV, 9575 · EDCS-36101100 · EDB17188}
\label{not:302}
\nf{Datation}{388 — directe : lecture partielle d’un nom consulaire}
\nf{Genre}{Épitaphe ; ensembles A, B, C ; 60 lettres transmises ; 45 \% restituées.}
\nf{Références}{ICUR IV 9575.}
\nf{Provenance}{Rattachement selon l’EDB : \emph{Coem. Callisti pars inferior}.}
\nf{Support}{Plaque de marbre. Inscription gravée. Fiche EDB : Antonio Enrico Felle, première version 2013.}
\ntext{\lnum{1}Infanti Eu[- - - in]\\
\lnum{2}pace qui vi⸢x⸣i[t - - -]\\
\lnum{3}depositus I[- - - Magno Ma]xim[o Augusto II consule]}
\nf{Témoins}{Export EDCS de 2020 : \emph{Infanti Eu[3 in] / pace qui vii[t 3] / depositus I[3 Magno Ma]xim[o Augusto II consule]}. EDB (EDB17188 ; écarts à examiner) : \emph{infanti Eu[--- in] / pace qui visi[t ---] / depositus I[--- Magno Ma]xim[o Augusto II consule] [texte tronqué dans la liste de résultats]}.}
\nf{Apparat}{⸢x⸣ : S gravé}
\nf{Traduction}{À l’enfant Eu[…], en paix, qui vécut […] ; déposé […] sous le consulat de Magnus Maximus Auguste pour la 2e fois.}
\nf{Langue}{L’éditeur signale une graphie corrigée : S gravé pour x (s pour x).}
\nf{Chronologie}{Noms consulaires : Magnus Maximus (lu en partie).}
\nf{Commentaire}{La formule in pace est attestée seulement par restitution (répartie sur deux lignes).}
\nf{Bibliographie}{ICUR IV 9575.}
\nf{Statut}{Lecture secondaire (EDCS) ; écarts avec l’EDB à arbitrer sur l’ICVR. L’EDCS signale une photographie (non consultée).}
\end{notice}

\begin{notice}{303}{388}{ICUR n.s. IV, 12250 · EDCS-37600684 · EDB5349}
\label{not:303}
\nf{Datation}{388 — directe : nom consulaire lu}
\nf{Genre}{Épitaphe ; ensembles A, B, C ; 88 lettres transmises.}
\nf{Références}{ICUR IV 12250.}
\nf{Provenance}{Rattachement selon l’EDB : \emph{Coem. anonymum ad viam Ardeatinam}. Conservation selon l’EDB : \emph{Coem. anonymum ad viam Ardeatinam}.}
\nf{Support}{Plaque de marbre. Inscription gravée. Reproductions signalées par l’EDB : photographie PCAS (droits réservés ; non reproduites). Fiche EDB : Marida Pierno, première version 2013.}
\ntext{\lnum{1}Hic Laurentia iacet qu(a)e vixit\\
\lnum{2}annum unum menses V\\
\lnum{3}cons(ulibus) d(omino) n(ostro) Fl(avio) T(h)eodosio II\\
\lnum{4}et M⸢e⸣robaude II d(eposita) IIII N(onas) Mart(ias)\\
\lnum{5}in pace}
\nf{Témoins}{Export EDCS de 2020 : \emph{Hic Laurentia iacet qu(a)e vixit / annum unum menses V / cons(ulibus) d(omino) n(ostro) Fl(avio) T(h)eodosio II / et Mrobaude II d(eposita) IIII N(onas) Mart(ias) / in pace}. EDB (EDB5349) : texte concordant aux conventions de notation près (indice de similarité 0,973).}
\nf{Apparat}{⸢e⸣ : Y gravé}
\nf{Traduction}{Ici gît Laurentia, qui vécut un an 5 mois ; sous le consulat de notre seigneur Flavius Théodose pour la 2e fois et de Merobaudes pour la 2e (sic) fois ; déposée le 4e jour avant les nones de mars, en paix.}
\nf{Langue}{L’éditeur signale une graphie corrigée : Y gravé pour e.}
\nf{Chronologie}{Noms consulaires : Theodosius (lu) et Merobaudes (lu). Dates du jour : déposition : IIII Nonas Martias = 4 mars 388 (lu). Âge déclaré : 1 a. 5 m. (lu).}
\nf{Commentaire}{La formule in pace est conservée.}
\nf{Bibliographie}{ICUR IV 12250.}
\nf{Statut}{Lecture secondaire (EDCS) concordante avec l’EDB : lecture stabilisée entre deux bases, non collationnée sur l’ICVR ni sur le monument. L’EDCS signale une photographie (non consultée).}
\end{notice}

\annee{389}

\begin{notice}{304}{389}{EDCS-08900756 · EDB27361}
\label{not:304}
\nf{Datation}{389 — directe : nom consulaire lu}
\nf{Genre}{Épitaphe ; ensembles A, B, C ; 105 lettres transmises.}
\nf{Références}{AE 1979, 43.}
\nf{Provenance}{Rattachement selon l’EDB : \emph{Incertae originis} (le volume I des ICVR réunit les inscriptions dont la provenance est perdue). Conservation selon l’EDB : \emph{Pontificio Istituto Biblico}.}
\nf{Support}{Plaque de marbre. 17 × 66 × 3,5 cm (h. × l. × ép.). Inscription gravée. Signes décrits par l’EDB : chrisme et chrisme avec alpha et oméga. Fiche EDB : Anita Rocco, première version 2013.}
\ntext{\lnum{1}Timasio et Promoto cons(ulibus) bene merenti Alexa\\
\lnum{2}ndri(a)e coniu⸢g⸣i qu(a)e vi⸢x⸣it annos plus minus\\
\lnum{3}XXV deposita in pace XVIII Kal(endas) Septembres}
\nf{Témoins}{Export EDCS de 2020 : \emph{Timasio et Promoto cons(ulibus) bene merenti Alexa/ndri(a)e coniui qu(a)e viit annos plus minus / XXV deposita in pace XVIII Kal(endas) Septembres}. EDB (EDB27361) : texte concordant aux conventions de notation près (indice de similarité 0,927).}
\nf{Apparat}{⸢g⸣ : C gravé ; ⸢x⸣ : S gravé}
\nf{Traduction}{Sous le consulat de Timasius et de Promotus, à Alexandria, son épouse bien méritante, qui vécut environ 25 ans ; déposée en paix le 18e jour avant les calendes de septembre.}
\nf{Langue}{L’éditeur signale 2 graphies corrigées : C gravé pour g (c pour g) ; S gravé pour x (s pour x).}
\nf{Chronologie}{Noms consulaires : Timasius (lu) et Promotus (lu). Dates du jour : déposition : XVIII Kalendas Septembres = 15 août 389 (lu). Âge déclaré : env. 25 a. (lu).}
\nf{Commentaire}{La formule in pace est conservée.}
\nf{Bibliographie}{AE 1979, 43. Bibliographie complémentaire selon l’EDB : Ferrua 1979a, 31-32, n. 21, Tav. IV, 3 ; AE, 1979, 43.}
\nf{Statut}{Lecture secondaire (EDCS) concordante avec l’EDB : lecture stabilisée entre deux bases, non collationnée sur l’ICVR ni sur le monument.}
\end{notice}

\begin{notice}{305}{389}{EDCS-08300019 · EDB34442}
\label{not:305}
\nf{Datation}{389 — directe : lecture partielle d’un nom consulaire}
\nf{Genre}{Épitaphe ; ensembles A, B, C ; 58 lettres transmises ; 19 \% restituées.}
\nf{Références}{AE 1985, 60.}
\nf{Provenance}{Rattachement selon l’EDB : \emph{Ipogeo della Torretta}. Conservation selon l’EDB : \emph{Ipogeo della Torretta}.}
\nf{Support}{La fiche EDB ne figure pas dans l’état 2019 des données EpiDoc ; support non documenté.}
\ntext{\lnum{1}consulatu Timasi et Pro]moti vv(irorum) cc(larissimorum)\\
\lnum{2}[- - -]ne veram [- - -]\\
\lnum{3}[- - - vixi]t dies XX X [horas ?]\\
\lnum{4}[in] pace}
\nf{Témoins}{EDB (EDB34442) : texte concordant aux conventions de notation près (indice de similarité 0,987).}
\nf{Traduction}{Sous le consulat de Timasius et de Promotus, clarissimes […] ; […] vraie […] ; […] vécut 20 jours 10 heures ; en paix.}
\nf{Langue}{2 abréviations développées par l’éditeur (dont vv(irorum), cc(larissimorum)).}
\nf{Chronologie}{Noms consulaires : Timasius (restitué) et Promotus (lu en partie). Âge déclaré : 20 j. (partiellement restitué).}
\nf{Commentaire}{La formule in pace est attestée seulement par restitution.}
\nf{Bibliographie}{AE 1985, 60. Bibliographie complémentaire selon l’EDB : Fasola 1985, 20 ; AE, 1985, 60.}
\nf{Statut}{Lecture secondaire (EDCS) concordante avec l’EDB : lecture stabilisée entre deux bases, non collationnée sur l’ICVR ni sur le monument.}
\end{notice}

\begin{notice}{306}{389}{ICUR n.s. V, 13355 · EDCS-01000473 · EDB27045}
\label{not:306}
\nf{Datation}{389 — directe : nom consulaire lu}
\nf{Genre}{Épitaphe ; ensembles A, B, C ; 664 lettres transmises ; 4 \% restituées.}
\nf{Références}{CIL VI 41342 = ICUR V 13355 = AE 2009, 145.}
\nf{Provenance}{Rattachement selon l’EDB : \emph{Coem. subdiale ad Catacumbas}. Conservation selon l’EDB : \emph{Coem. subdiale ad Catacumbas}.}
\nf{Support}{Sarcophage de marbre. 72 × 212 cm (h. × l.). Inscription gravée. Reproductions signalées par l’EDB : photographie PCAS ×2 (droits réservés ; non reproduites). Fiche EDB : Anita Rocco, première version 2010.}
\ntext{\lnum{1}Maximilla poten[s servare f(o)edera C]⸢hr⸣isti propositumque fidem mors abstulit omnia vitae sed melius vivit vitam C(h)risti qui sperat (a)eternam si animo recolat fidem datam omn[ibus] unam tunc poterit mors ipsa mori cum tempore toto vita perennis eri⸢t⸣ veniam dabit omnibus unam sed meritis [ve]⸢g⸣etans su[per h]anc Lucceia fecisti nam m⸢e⸣mor ipsa tui cum reddis vota sepulcr\{h\}i ostend⸢i⸣s mentem religiosa vita repleta(m) mirae fide(i) et innocentiae optimique propositi Maximilla virgo ancilla dei civis Pannonia qu(a)e vixit annis n(umero) L huic sepulturae condita est a \{L\} Lucceia c(larissima) f(emina) filia Viventi c(larissimae) r(ecordationis) v(iri) ex praefecto pr(a)etorio et urbis aeternae prop⸢te⸣r amicitiam sibi coniun(c)tam qu(a)e diem clusit una cum matre sua pariter Nunita qu(a)e fuit matrona diaconi\{s\} et vixit annis n(umero) LXX die IIII Non(as) Octobr(es) FFll(aviis) Timasio et Promoto vv(viris) cc(larissimis) conss(ulibus)}
\nf{Témoins}{Export EDCS de 2020 : \emph{Maximilla poten[s servare f(o)edera C]isti propositumque fidem mors abstulit omnia vitae sed melius vivit vitam C(h)risti qui sperat (a)eternam si animo recolat fidem / datam omn[ibus] unam tunc poterit mors ipsa mori cum tempore toto vita perennis eri veniam dabit omnibus unam sed meritis / [ve]etans su[per h]anc Lucceia fecisti nam mmor ipsa tui cum reddis vota sepulcr\{h\}i ostends mentem religio…}. EDB (EDB27045 ; divergence forte) : \emph{Maximilla poten[s servare foedera C]rhisti propositumque fidem mors abstulit omnia vitae sed melius vivit vitam Cristi qui sperat eternam si animo recolat fidem / datam omn[ibus] unam tunc poterit mors ipsa mori cum tempore toto vita perennis eri⌜t⌝ veniam dabit omnibus unam sed meritis / [ve]getans su[per h]anc Lucceia fecisti nam memor ipsa tui cum reddis vota sepulcrhi ostendis mentem religiosa vita repleta⟨m⟩. Mirae fide⟨i⟩}. La ligne \emph{et innocentiae optimique propositi Maximilla virgo ancilla dei civis Pannonia qu(a)e} est transmise par l'EDCS, mais absente du texte de la liste EDB (EDB27045), qui passe de \emph{Mirae fide⟨i⟩} à \emph{vixit annis}. Comme cette ligne contient \emph{virgo} et \emph{ancilla dei}, son existence doit être contrôlée sur l'ICVR V 13355 avant toute exploitation. Le support est un sarcophage de marbre (72 × 212 cm) de la Platonia de Saint-Sébastien ; deux photographies PCAS sont signalées.}
\nf{Apparat}{⸢hr⸣ : RH gravé ; ⸢t⸣ : I gravé ; ⸢g⸣ : C gravé ; ⸢e⸣ : I gravé ; ⸢i⸣ : E gravé ; ⸢te⸣ : I gravé ; \{h\} : lettres superflues ; \{L\} : lettres superflues ; \{s\} : lettres superflues}
\nf{Traduction}{(vers) Maximilla, capable de garder l’alliance du Christ et le propos (de virginité) : la mort a emporté la foi (?) et tout ce qui est de la vie ; mais il vit mieux la vie du Christ, celui qui espère l’éternelle, s’il médite en son âme la foi unique donnée à tous ; alors la mort elle-même pourra mourir, quand avec le temps entier la vie sera perpétuelle ; il accordera à tous un pardon unique, mais selon les mérites ; fortifiant au-dessus de celle-ci, Lucceia, tu as fait (ce monument) ; car, te souvenant d’elle quand tu rends tes vœux au sépulcre, tu montres un esprit rempli d’une vie religieuse. — D’une foi admirable, d’une innocence et du meilleur propos, Maximilla, vierge, servante de Dieu, citoyenne de Pannonie, qui vécut 50 ans, a été ensevelie dans cette sépulture par Lucceia, femme clarissime, fille de Viventius, homme de clarissime mémoire, ancien préfet du prétoire et de la Ville éternelle, en raison de l’amitié qui les unissait ; elle termina ses jours avec sa mère Nunita, qui fut épouse d’un diacre et vécut 70 ans, le 4e jour avant les nones d’octobre, sous le consulat des Flavii Timasius et Promotus, clarissimes.}
\nf{Langue}{L’éditeur signale 6 graphies corrigées : RH gravé pour hr ; I gravé pour t ; C gravé pour g (c pour g) ; I gravé pour e (confusion e/i) ; E gravé pour i (confusion e/i) ; I gravé pour te. Des lettres gravées sont jugées superflues par l’éditeur (entre accolades). 23 abréviations développées par l’éditeur (dont f(o)edera, C(h)risti, (a)eternam, repleta(m), fide(i), qu(a)e).}
\nf{Chronologie}{Noms consulaires : Timasius (lu) et Promotus (lu). Dates du jour : événement indéterminé : die IIII Nonas Octobres = 4 octobre 389 (lu). Âges déclarés : 50 a. (lu) ; 70 a. 4 j. (lu).}
\nf{Commentaire}{Vocabulaire religieux : désignation de Dieu — dei (lu), ancilla dei (lu) ; désignation du Christ — Christi (partiellement restitué), Christi (lu) ; initiation et foi — fidem (lu), fidei (lu). Rang sénatorial d’une personne commémorée (ae condita est a lucceia clarissima femina fili ; clarissima femina filia viventi clarissimae reco). Fonction ou métier mentionné : praefecto, diaconi. Désignation divine : retenue dans le profil strict. Commentaire développé : dossier p.~\pageref{dos:306}.}
\nf{Bibliographie}{CIL VI 41342 ; ICUR V 13355 ; AE 2009, 145. Bibliographie complémentaire selon l’EDB : Mazzoleni 2015, 167.}
\nf{Statut}{Lecture secondaire (EDCS) ; divergence forte avec l’EDB : le texte ne doit pas être exploité sans contrôle de l’édition imprimée. L’EDCS signale une photographie (non consultée).}
\end{notice}

\begin{notice}{307}{389}{EDCS-74700288 · EDB31230}
\label{not:307}
\nf{Datation}{389 — directe : nom consulaire lu}
\nf{Genre}{Épitaphe ; ensembles A, B, C ; 19 lettres transmises.}
\nf{Références}{Ferrua-1979, 01444.}
\nf{Provenance}{Rattachement selon l’EDB : \emph{Incertae originis} (le volume I des ICVR réunit les inscriptions dont la provenance est perdue). Conservation selon l’EDB : \emph{Musei Vaticani: Lapidario Cristiano}.}
\nf{Support}{Plaque de marbre. 21 × 84 cm (h. × l.). Inscription gravée. Fiche EDB : Filippo Piazzolla, première version 2018.}
\ntext{\lnum{1}OIT [- - -]\\
\lnum{2}Promoto vv(iris) cc(larissimis) conss(ulibus)}
\nf{Témoins}{EDB (EDB31230) : texte concordant aux conventions de notation près (indice de similarité 0,936).}
\nf{Traduction}{[…] (sous le consulat de Timasius et) Promotus, clarissimes.}
\nf{Langue}{3 abréviations développées par l’éditeur (dont vv(iris), cc(larissimis), conss(ulibus)).}
\nf{Chronologie}{Noms consulaires : Promotus (lu).}
\nf{Bibliographie}{Ferrua-1979, 01444. Bibliographie complémentaire selon l’EDB : ICVR I, 1444 ; Ferrua 1979b, 37, n. 1444 ; Ferrua 1975c, 144-145 fig. 8.}
\nf{Statut}{Lecture secondaire (EDCS) concordante avec l’EDB : lecture stabilisée entre deux bases, non collationnée sur l’ICVR ni sur le monument. L’EDCS signale une photographie (non consultée).}
\end{notice}

\begin{notice}{308}{389}{ICUR n.s. I, 2771 · EDCS-39900320 · EDB32700}
\label{not:308}
\nf{Datation}{389 — directe : nom consulaire lu}
\nf{Genre}{Épitaphe ; ensembles A, B, C ; 90 lettres transmises.}
\nf{Références}{ICUR I 2771 = ILCV 1481.}
\nf{Provenance}{Rattachement selon l’EDB : \emph{Incertae originis} (le volume I des ICVR réunit les inscriptions dont la provenance est perdue). Conservation selon l’EDB : \emph{Fossombrone (PU)   ; Lapidario Diocesano}.}
\nf{Support}{Plaque de marbre. 28 × 79 cm (h. × l.). Inscription gravée. Reproductions signalées par l’EDB : reproduction (droits réservés ; non reproduites). Fiche EDB : Filippo Piazzolla, première version 2018.}
\ntext{\lnum{1}Aristo puer innox qui\\
\lnum{2}vixi⸢t⸣ mens\{s\}es oc⸢t⸣o neo\\
\lnum{3}f⸢y⸣tus reces(s)it p⸢r⸣idie Nonas\\
\lnum{4}Iunias Timasio et Promoto vv(iris)\\
\lnum{5}cc(larissimis) conss(ulibus)}
\nf{Témoins}{Export EDCS de 2020 : \emph{Aristo puer innox qui / vixi mens\{s\}es oco neo/ftus reces(s)it pidie Nonas / Iunias Timasio et Promoto vv(iris) / cc(larissimis) conss(ulibus)}. EDB (EDB32700) : texte concordant aux conventions de notation près (indice de similarité 0,951).}
\nf{Apparat}{⸢t⸣ : I gravé ; ⸢t⸣ : I gravé ; ⸢y⸣ : I gravé ; ⸢r⸣ : B gravé ; \{s\} : lettres superflues}
\nf{Traduction}{Aristus, enfant innocent, qui vécut huit mois, néophyte ; il s’en alla la veille des nones de juin, sous le consulat de Timasius et de Promotus, clarissimes.}
\nf{Langue}{L’éditeur signale 4 graphies corrigées : I gravé pour t ×2 ; I gravé pour y (graphie de y) ; B gravé pour r. Des lettres gravées sont jugées superflues par l’éditeur (entre accolades). 4 abréviations développées par l’éditeur (dont reces(s)it, vv(iris), cc(larissimis), conss(ulibus)).}
\nf{Chronologie}{Noms consulaires : Timasius (lu) et Promotus (lu). Dates du jour : décès : pridie Nonas Iunias = 4 juin 389 (lu). Âge déclaré : 8 m. (lu).}
\nf{Commentaire}{Vocabulaire religieux : initiation et foi — neofytus (lu).}
\nf{Bibliographie}{ICUR I 2771 ; ILCV 1481. Bibliographie complémentaire selon l’EDB : ILCV, 1481.}
\nf{Statut}{Lecture secondaire (EDCS) concordante avec l’EDB : lecture stabilisée entre deux bases, non collationnée sur l’ICVR ni sur le monument. L’EDCS signale une photographie (non consultée).}
\end{notice}

\begin{notice}{309}{389}{ICUR n.s. I, 3208 · EDCS-38100791 · EDB21459}
\label{not:309}
\nf{Datation}{389 — directe : nom consulaire lu}
\nf{Genre}{Épitaphe ; ensembles A, B, C ; 96 lettres transmises ; 23 \% restituées.}
\nf{Références}{ICUR I 3208 = ICUR II 5934.}
\nf{Provenance}{Rattachement selon l’EDB : \emph{Museum paulianum}. Conservation selon l’EDB : \emph{Basilica s.Pauli apostoli}.}
\nf{Support}{Plaque de marbre. 30 × 23 cm (h. × l.). Inscription gravée. Fiche EDB : Maria Cristina Grisanzio, première version 2014.}
\ntext{\lnum{1}[- - - d]isces[sit]\\
\lnum{2}[de ha]c luce X Kal(endas)\\
\lnum{3}[Iu]n(ias) Timatio\\
\lnum{4}[et] Promoto co⸢n⸣\\
\lnum{5}sulibus ann⸢is⸣\\
\lnum{6}[- - -]XI m(ensibus) III d(iebus) XXVI\\
\lnum{7}[domus] aeterna Eusebi\\
\lnum{8}[et fil]i(a)e eius}
\nf{Témoins}{Export EDCS de 2020 : \emph{[3 d]isces[sit] / [de ha]c luce X Kal(endas) / [Iu]n(ias) Timatio / [et] Promoto co/sulibus ann / [3]XI m(ensibus) III d(iebus) XXVI / [domus] aeterna Eusebi / [et fil]i(a)e eius}. EDB (EDB21459) : texte concordant aux conventions de notation près (indice de similarité 0,992).}
\nf{Apparat}{⸢n⸣ : V gravé ; ⸢is⸣ : SI gravé}
\nf{Traduction}{[…] il/elle quitta cette lumière le 10e jour avant les calendes de juin, sous le consulat de Timasius et de Promotus, à […]11 ans 3 mois 26 jours ; demeure éternelle d’Eusebius et de sa fille.}
\nf{Langue}{L’éditeur signale 2 graphies corrigées : V gravé pour n ; SI gravé pour is.}
\nf{Chronologie}{Noms consulaires : Promotus (lu). Ne sont pas des consuls : [domus] aeterna Eusebi (défunt ; le consul lu est Timatio (Timasius), absent de la liste v2.0). Dates du jour : décès : X Kalendas Iunias = 23 mai 389 (partiellement restitué).}
\nf{Commentaire}{Vocabulaire religieux : allusion interprétative (niveau 5) — luce (lu).}
\nf{Bibliographie}{ICUR I 3208 ; ICUR II 5934. Bibliographie complémentaire selon l’EDB : ICVR I, 3208.}
\nf{Statut}{Lecture secondaire (EDCS) concordante avec l’EDB : lecture stabilisée entre deux bases, non collationnée sur l’ICVR ni sur le monument. L’EDCS signale une photographie (non consultée).}
\end{notice}

\begin{notice}{310}{389}{ICUR n.s. II, 4823 · EDCS-36101305 · EDB43052}
\label{not:310}
\nf{Datation}{389 — directe : nom consulaire lu}
\nf{Genre}{Épitaphe ; ensembles A, B, C ; 130 lettres transmises ; 14 \% restituées.}
\nf{Références}{ICUR II 4823 = ILCV 1130 = ICaRoma 124 = Mossong-2022, 00132.}
\nf{Provenance}{Rome (indication de l’EDCS) ; aucun rattachement cimétérial transmis. Conservation selon l’EDB : \emph{s. Paolo f.l.m.}.}
\nf{Support}{La fiche EDB ne figure pas dans l’état 2019 des données EpiDoc ; support non documenté.}
\ntext{\lnum{1}Gaudentius presb(yter) sibi\\
\lnum{2}et coniugi suae Severae castae hac(!) sanc[tissimae]\\
\lnum{3}feminae quae vixit ann(os) XLII m(enses) III d(ies) X\\
\lnum{4}dep(osita) IIII Non(as) April(es) Timasio et Promoto [cons(ulibus)]}
\nf{Témoins}{EDB (EDB43052) : texte concordant aux conventions de notation près (indice de similarité 0,968).}
\nf{Apparat}{(!) : forme signalée par l’éditeur}
\nf{Traduction}{Gaudentius, prêtre, pour lui-même et pour son épouse Severa, chaste et très sainte femme, qui vécut 42 ans 3 mois 10 jours ; déposée le 4e jour avant les nones d’avril, sous le consulat de Timasius et de Promotus.}
\nf{Langue}{L’éditeur signale au moins une forme anormale (sic). 9 abréviations développées par l’éditeur (dont presb(yter), hac(!), ann(os), m(enses), d(ies), dep(osita)).}
\nf{Chronologie}{Noms consulaires : Timasius (lu) et Promotus (lu). Dates du jour : déposition : IIII Nonas Apriles = 2 avril 389 (lu). Âge déclaré : 42 a. 3 m. 10 j. (lu).}
\nf{Commentaire}{Vocabulaire religieux : épithète du défunt (non doctrinal) — sanctissimae (partiellement restitué). Fonction ou métier mentionné : presbyter.}
\nf{Bibliographie}{ICUR II 4823 ; ILCV 1130 ; ICaRoma 124 ; Mossong-2022, 00132. Bibliographie complémentaire selon l’EDB : ILCV, 1130.}
\nf{Statut}{Lecture secondaire (EDCS) concordante avec l’EDB : lecture stabilisée entre deux bases, non collationnée sur l’ICVR ni sur le monument. L’EDCS signale une photographie (non consultée).}
\end{notice}

\annee{388}

\begin{notice}{311}{388}{ICUR n.s. III, 6506 · EDCS-34801306 · EDB22064}
\label{not:311}
\nf{Datation}{388 — directe : nom consulaire lu ; datation contestée}
\nf{Genre}{Épitaphe ; ensembles A, B, C ; 53 lettres transmises.}
\nf{Références}{ICUR III 6506 = ILCV 3287a (em).}
\nf{Provenance}{Rattachement selon l’EDB : \emph{Coem. Domitillae pars inferior}. Conservation selon l’EDB : \emph{Coem. Domitillae pars inferior}.}
\nf{Support}{Enduit de mortier de mortier. Inscription tracée dans le mortier frais. Signes décrits par l’EDB : chrisme. Fiche EDB : Carlo Carletti / Antonio Enrico Felle, première version 2013.}
\ntext{\lnum{1}Depositus Nisus pridie N(onas) deposita in Chr(isto) VI idus Feb(ruarias) Maximo cons(ule)}
\nf{Témoins}{EDB (EDB22064) : texte concordant aux conventions de notation près (indice de similarité 0,934).}
\nf{Traduction}{Déposé, Nisus, la veille des nones ; déposée dans le Christ le 6e jour avant les ides de février, sous le consulat de Maximus.}
\nf{Langue}{4 abréviations développées par l’éditeur (dont N(onas), Chr(isto), Feb(ruarias), cons(ule)).}
\nf{Chronologie}{Noms consulaires : Magnus Maximus (lu). \emph{Maximo cons(ule)} désigne le consulat de Magnus Maximus (388), en usage à Rome au début de cette année ; la date du 8 février convient. L'EDB (EDB22064) date de 389 et note que sa transcription diffère de l'ICVR. Datation contestée (388/389). Dates du jour : déposition : VI idus Februarias = 8 février 388 (lu).}
\nf{Commentaire}{Vocabulaire religieux : désignation du Christ — Christo (abréviation développée). \emph{in Chr(isto)} est transcrit par les deux bases, mais l'EDB décrit un chrisme sur le support : comme aux n° 140 et 151, l'abréviation développée rend probablement un monogramme. Maintenu dans le profil large, exclu du profil strict. Désignation divine : retenue dans le profil restrictif. Commentaire développé : dossier p.~\pageref{dos:311}.}
\nf{Bibliographie}{ICUR III 6506 ; ILCV 3287a (em).}
\nf{Statut}{Lecture secondaire (EDCS) concordante avec l’EDB : lecture stabilisée entre deux bases, non collationnée sur l’ICVR ni sur le monument.}
\end{notice}

\annee{389}

\begin{notice}{312}{389}{ICUR n.s. III, 6507 · EDCS-34801307 · EDB22065}
\label{not:312}
\nf{Datation}{389 — directe : lecture partielle d’un nom consulaire}
\nf{Genre}{Épitaphe ; ensembles A, B, C ; 54 lettres transmises ; 61 \% restituées.}
\nf{Références}{ICUR III 6507.}
\nf{Provenance}{Rattachement selon l’EDB : \emph{Coem. Domitillae pars inferior}. Conservation selon l’EDB : \emph{Coem. Domitillae pars inferior}.}
\nf{Support}{Plaque de marbre. 14 × 75 × 3,3 cm (h. × l. × ép.). Inscription gravée. Fiche EDB : Carlo Carletti / Antonio Enrico Felle, première version 2003.}
\ntext{\lnum{1}Lau[ren]tia quae vixi[t ann(os) - - -]\\
\lnum{2}[- - -]bres Tim[asio et Promoto conss(ulibus)]}
\nf{Témoins}{EDB (EDB22065 ; écarts à examiner) : \emph{Lau[ren]tia quae vixi[t ---] / [---]bres Tim[asio et Promoto cons(ulibu)s] [texte tronqué dans la liste de résultats]}.}
\nf{Traduction}{Laurentia, qui vécut […] ans ; […] (mois en -bres), sous le consulat de Timasius et de Promotus.}
\nf{Langue}{2 abréviations développées par l’éditeur (dont ann(os), conss(ulibus)).}
\nf{Chronologie}{Noms consulaires : Timasius (lu en partie) et Promotus (restitué).}
\nf{Bibliographie}{ICUR III 6507.}
\nf{Statut}{Lecture secondaire (EDCS) ; écarts avec l’EDB à arbitrer sur l’ICVR. Plus de la moitié des lettres sont restituées. L’EDCS signale une photographie (non consultée).}
\end{notice}

\begin{notice}{313}{389}{ICUR n.s. IV, 11117a · EDCS-37200399 · EDB39579}
\label{not:313}
\nf{Datation}{389 — directe : nom consulaire lu}
\nf{Genre}{Épitaphe ; ensembles A, B, C ; 47 lettres transmises ; 11 \% restituées.}
\nf{Références}{ICUR IV 11117a.}
\nf{Provenance}{Rattachement selon l’EDB : \emph{Coem. Callisti area subdialis et incerti loci}. Conservation selon l’EDB : \emph{Coem. Callisti area subdialis et incerti loci}.}
\nf{Support}{Plaque de marbre. 22 × 105 × 3,6 cm (h. × l. × ép.). Inscription gravée. Reproductions signalées par l’EDB : photographie PCAS ×3 (droits réservés ; non reproduites). Fiche EDB : Marida Pierno, première version 2015.}
\ntext{\lnum{1}Ti]masio et Promoto conss(ulibus) deposit[us(?) est - - -]\\
\lnum{2}XV Kal(endas) No⸢v⸣embres}
\nf{Témoins}{Export EDCS de 2020 : \emph{Ti]masio et Promoto conss(ulibus) deposit[us(?) est 3] / XV Kal(endas) Noembres}. EDB (EDB39579 ; divergence forte) : \emph{⟨: a⟩ [--- Ti]masio et Promoto con(sulibu)s deposit[us est ---] / XV kal(endas) nobembres /⟨: b⟩ et fra[ter eius ---] / d̅(omi)n̅(o) V[alentiniano augusto IIII et Neoterio consulibus]}. L'EDB (EDB39579 = ICVR IV 11117) publie une fiche unique, en deux sections marquées a et b : une plaque de marbre (22 × 105 × 3,6 cm) portant deux dépositions successives, celle d'un défunt en 389 (a, ici) et celle de son frère (b, n° 314). Il s'agit d'un monument et de deux commémorations, non de deux fragments d'un même texte.}
\nf{Apparat}{⸢v⸣ : B gravé}
\nf{Traduction}{Sous le consulat de Timasius et de Promotus a été déposé […] le 15e jour avant les calendes de novembre.}
\nf{Langue}{L’éditeur signale une graphie corrigée : B gravé pour v (bêtacisme).}
\nf{Chronologie}{Noms consulaires : Timasius (lu en partie) et Promotus (lu). Dates du jour : déposition : XV Kalendas Novembres = 18 octobre 389 (lu).}
\nf{Bibliographie}{ICUR IV 11117a.}
\nf{Statut}{Lecture secondaire (EDCS) ; divergence forte avec l’EDB : le texte ne doit pas être exploité sans contrôle de l’édition imprimée. L’EDCS signale une photographie (non consultée).}
\end{notice}

\annee{390}

\begin{notice}{314}{390}{ICUR n.s. IV, 11117b · EDCS-37200400 · EDB39579}
\label{not:314}
\nf{Datation}{390 — directe : consulat restitué ; datation contestée}
\nf{Genre}{Épitaphe ; ensembles A, B, C ; 58 lettres transmises ; 86 \% restituées.}
\nf{Références}{ICUR IV 11117b.}
\nf{Provenance}{Rattachement selon l’EDB : \emph{Coem. Callisti area subdialis et incerti loci}. Conservation selon l’EDB : \emph{Coem. Callisti area subdialis et incerti loci}.}
\nf{Support}{Plaque de marbre. 22 × 105 × 3,6 cm (h. × l. × ép.). Inscription gravée. Reproductions signalées par l’EDB : photographie PCAS ×3 (droits réservés ; non reproduites). Fiche EDB : Marida Pierno, première version 2015.}
\ntext{\lnum{1}Et fra[ter eius - - -]\\
\lnum{2}d(omino) n(ostro) V[alentiniano Aug(usto) IIII et Neoterio conss(ulibus)]}
\nf{Témoins}{EDB (EDB39579 ; divergence forte) : \emph{⟨: a⟩ [--- Ti]masio et Promoto con(sulibu)s deposit[us est ---] / XV kal(endas) nobembres /⟨: b⟩ et fra[ter eius ---] / d̅(omi)n̅(o) V[alentiniano augusto IIII et Neoterio consulibus]}.}
\nf{Traduction}{Et son frère […] ; sous le consulat de notre seigneur Valentinien Auguste pour la 4e fois et de Neoterius.}
\nf{Langue}{4 abréviations développées par l’éditeur (dont d(omino), n(ostro), Aug(usto), conss(ulibus)).}
\nf{Chronologie}{Noms consulaires : Valentinianus (restitué) et Neoterius (restitué). Seul \emph{d(omino) n(ostro) V[…]} est lu. Le collège \emph{V[alentiniano Aug(usto) IIII et Neoterio]} (390) est une restitution partagée par l'EDCS et l'EDB. La note de l'EDB fait observer que le frère est mort avant le défunt de la section a, ce qui s'accorde mal avec une date de 390 : la datation de cette section est contestée.}
\nf{Bibliographie}{ICUR IV 11117b.}
\nf{Statut}{Lecture secondaire (EDCS) ; divergence forte avec l’EDB : le texte ne doit pas être exploité sans contrôle de l’édition imprimée. Plus de la moitié des lettres sont restituées. L’EDCS signale une photographie (non consultée).}
\end{notice}

\annee{389}

\begin{notice}{315}{389}{ICUR n.s. IV, 11118 · EDCS-37200401 · EDB39580}
\label{not:315}
\nf{Datation}{389 — directe : nom consulaire lu}
\nf{Genre}{Épitaphe ; ensembles A, B, C ; 44 lettres transmises ; 34 \% restituées.}
\nf{Références}{ICUR IV 11118.}
\nf{Provenance}{Rattachement selon l’EDB : \emph{Coem. Callisti area subdialis et incerti loci}. Conservation selon l’EDB : \emph{Musei Vaticani: Lapidario Cristiano}.}
\nf{Support}{Plaque de marbre. 46 × 37 cm (h. × l.). Inscription gravée. Reproductions signalées par l’EDB : photographie PCAS (droits réservés ; non reproduites). Fiche EDB : Marida Pierno, première version 2016.}
\ntext{\lnum{1}a]nnis de\\
\lnum{2}[cem - - - dep]osita VI Kal(endas)\\
\lnum{3}[- - - Timasio et] Promoto conss(ulibus)}
\nf{Témoins}{EDB (EDB39580) : texte concordant aux conventions de notation près (indice de similarité 0,965).}
\nf{Traduction}{[…] dix ans […] ; déposée le 6e jour avant les calendes […] ; sous le consulat de Timasius et de Promotus.}
\nf{Langue}{2 abréviations développées par l’éditeur (dont Kal(endas), conss(ulibus)).}
\nf{Chronologie}{Noms consulaires : Timasius (restitué) et Promotus (lu).}
\nf{Bibliographie}{ICUR IV 11118. Bibliographie complémentaire selon l’EDB : ICVR IV, 11578 ; Lega 2000.}
\nf{Statut}{Lecture secondaire (EDCS) concordante avec l’EDB : lecture stabilisée entre deux bases, non collationnée sur l’ICVR ni sur le monument. L’EDCS signale une photographie (non consultée).}
\end{notice}

\begin{notice}{316}{389}{ICUR n.s. IV, 11119 · EDCS-37200402 · EDB39581}
\label{not:316}
\nf{Datation}{389 — directe : lecture partielle d’un nom consulaire}
\nf{Genre}{Épitaphe ; ensembles A, B, C ; 38 lettres transmises ; 71 \% restituées.}
\nf{Références}{ICUR IV 11119.}
\nf{Provenance}{Rattachement selon l’EDB : \emph{Coem. Callisti area subdialis et incerti loci}. Conservation selon l’EDB : \emph{Coem. Callisti area subdialis et incerti loci}.}
\nf{Support}{Plaque de marbre. 16 × 15 cm (h. × l.). Inscription gravée. Reproductions signalées par l’EDB : photographie PCAS (droits réservés ; non reproduites). Fiche EDB : Marida Pierno, première version 2015.}
\ntext{\lnum{1}] qu[- - -]\\
\lnum{2}[- - - dies] V depo[- - -]\\
\lnum{3}[- - - Timasio et] Prom[oto conss(ulibus)}
\nf{Témoins}{EDB (EDB39581 ; écarts à examiner) : \emph{[---]QU[---]/ [--- dies] V depo[sit---] / [--- Timasio et] Prom[oto]}.}
\nf{Traduction}{[…] 5 jours ; déposé(e) […] sous le consulat de Timasius et de Promotus.}
\nf{Langue}{1 abréviation développée par l’éditeur (dont conss(ulibus)).}
\nf{Chronologie}{Noms consulaires : Timasius (restitué) et Promotus (lu en partie).}
\nf{Bibliographie}{ICUR IV 11119.}
\nf{Statut}{Lecture secondaire (EDCS) ; écarts avec l’EDB à arbitrer sur l’ICVR. Plus de la moitié des lettres sont restituées. L’EDCS signale une photographie (non consultée).}
\end{notice}

\begin{notice}{317}{389}{ICUR n.s. IV, 11120 · EDCS-37200403 · EDB39582}
\label{not:317}
\nf{Datation}{389 — directe : nom consulaire lu}
\nf{Genre}{Épitaphe ; ensembles A, B, C ; 174 lettres transmises ; 48 \% restituées.}
\nf{Références}{ICUR IV 11120.}
\nf{Provenance}{Rattachement selon l’EDB : \emph{Coem. Callisti area subdialis et incerti loci}. Conservation selon l’EDB : \emph{Coem. Callisti area subdialis et incerti loci}.}
\nf{Support}{Plaque de marbre. 67 × 89 cm (h. × l.). Inscription gravée. Reproductions signalées par l’EDB : photographie PCAS ×2 (droits réservés ; non reproduites). Fiche EDB : Marida Pierno, première version 2015.}
\ntext{\lnum{1}[Dep(ositio) - - -]ies Id(ib)us Dec(embribus) conss(ulibus) Timasio [et]\\
\lnum{2}[Promoto vv(iris) cc(larissimis)] qu(i ?) vixit ann(os) XXXV et m(enses) [- - -]\\
\lnum{3}[d(omino) n(ostro) Valentin]iano IIII et (Ne)ute\{i\}rio\\
\lnum{4}[- - - qui] vixit ann(os) p(lus) m(inus) XIII deposit[us - - -]\\
\lnum{5}[Valentiniano II]II et Neuterio ben[e merenti}
\nf{Témoins}{EDB (EDB39582 ; écarts à examiner) : \emph{⟨: a⟩ [---]ies idus dec(embres) cons(ulibu)s Timasio [et] / [Promoto v(iris) c(larissimis)] que vixit ann(is) XXXV et m[ensibus ---] / ⟨: b⟩ [d(omino) n(ostro) Valentin]iano IIII et Uteirio / [--- qui] vixit ann(is) p(lus) m(inus) XIII deposit[us] / [Valentiniano II]II et Neuterio ben(emerenti) / D}.}
\nf{Apparat}{\{i\} : lettres superflues}
\nf{Traduction}{Déposition de […] aux ides de décembre, sous le consulat de Timasius et de Promotus, clarissimes ; qui vécut 35 ans et […] mois ; sous (le consulat de) notre seigneur Valentinien pour la 4e fois et Neoterius […] qui vécut environ 13 ans ; déposé […] sous (le consulat de) Valentinien pour la 4e fois et Neoterius ; bien méritant.}
\nf{Langue}{Des lettres gravées sont jugées superflues par l’éditeur (entre accolades). 15 abréviations développées par l’éditeur (dont Dep(ositio), Id(ib)us, Dec(embribus), conss(ulibus), vv(iris), cc(larissimis)).}
\nf{Chronologie}{Noms consulaires : Valentinianus (lu en partie), Timasius (lu), Promotus (restitué) et Neoterius (lu). Dates du jour : déposition : Idibus Decembribus = 13 décembre 389 (lu). Âges déclarés : env. 35 a. (partiellement restitué) ; env. 13 a. (partiellement restitué).}
\nf{Bibliographie}{ICUR IV 11120.}
\nf{Statut}{Lecture secondaire (EDCS) ; écarts avec l’EDB à arbitrer sur l’ICVR. L’EDCS signale une photographie (non consultée).}
\end{notice}

\begin{notice}{318}{389}{ICUR n.s. IV, 12251 · EDCS-37600685 · EDB6536}
\label{not:318}
\nf{Datation}{389 — directe : lecture partielle d’un nom consulaire}
\nf{Genre}{Épitaphe ; ensembles A, B, C ; 60 lettres transmises ; 32 \% restituées.}
\nf{Références}{ICUR IV 12251.}
\nf{Provenance}{Rattachement selon l’EDB : \emph{Coem. anonymum ad viam Ardeatinam}. Conservation selon l’EDB : \emph{Coem. anonymum ad viam Ardeatinam}.}
\nf{Support}{Plaque de marbre. Inscription gravée. Reproductions signalées par l’EDB : photographie PCAS (droits réservés ; non reproduites). Fiche EDB : Marida Pierno, première version 2013.}
\ntext{\lnum{1}Hic sita ⸢L⸣ucida [quae vix(it) a]nn(os) VIIII\\
\lnum{2}dep(osita) XII Kal(endas) Apr(iles) Fl(avio) T[imasio] et Prom[oto]}
\nf{Témoins}{Export EDCS de 2020 : \emph{Hic sita ucida [quae vix(it) a]nn(os) VIIII / dep(osita) XII Kal(endas) Apr(iles) Fl(avio) T[imasio] et Prom[oto]}. EDB (EDB6536) : texte concordant aux conventions de notation près (indice de similarité 0,98).}
\nf{Apparat}{⸢L⸣ : I gravé}
\nf{Traduction}{Ici gît Lucida, qui vécut 9 ans ; déposée le 12e jour avant les calendes d’avril, sous (le consulat de) Flavius Timasius et Promotus.}
\nf{Langue}{L’éditeur signale une graphie corrigée : I gravé pour L.}
\nf{Chronologie}{Noms consulaires : Timasius (restitué) et Promotus (lu en partie). Dates du jour : déposition : XII Kalendas Apriles = 21 mars 389 (lu). Âge déclaré : 9 a. (partiellement restitué).}
\nf{Bibliographie}{ICUR IV 12251.}
\nf{Statut}{Lecture secondaire (EDCS) concordante avec l’EDB : lecture stabilisée entre deux bases, non collationnée sur l’ICVR ni sur le monument. L’EDCS signale une photographie (non consultée).}
\end{notice}

\annee{390}

\begin{notice}{319}{390}{ICUR n.s. I, 1445 · EDCS-39000426 · EDB31209}
\label{not:319}
\nf{Datation}{390 — directe : nom consulaire lu}
\nf{Genre}{Épitaphe ; ensembles A, B, C ; 64 lettres transmises.}
\nf{Références}{ICUR I 1445 = AnalEpi p 377.}
\nf{Provenance}{Rattachement selon l’EDB : \emph{Incertae originis} (le volume I des ICVR réunit les inscriptions dont la provenance est perdue). Conservation selon l’EDB : \emph{Musei Vaticani: Lapidario Cristiano}.}
\nf{Support}{Plaque de marbre. 23 × 84 cm (h. × l.). Inscription gravée. Fiche EDB : Filippo Piazzolla, première version 2015.}
\ntext{\lnum{1}C⸢e⸣r\{r\}eal(is) bene me(re)nti coniugi Quriaceti\\
\lnum{2}feci⸢t⸣ d(epositio) XII Kal(endas) Aug(ustas) d(omino) n(ostro) Fl(avio) Val(entiniano) Aug(usto) et Neot(erio) v(iro) ⸢c⸣(larissimo)}
\nf{Témoins}{Export EDCS de 2020 : \emph{Cr\{r\}eal(is) bene me(re)nti coniugi Quriaceti / feci d(epositio) XII Kal(endas) Aug(ustas) d(omino) n(ostro) Fl(avio) Val(entiniano) Aug(usto) et Neot(erio) v(iro) (larissimo)}. EDB (EDB31209) : texte concordant aux conventions de notation près (indice de similarité 0,925).}
\nf{Apparat}{⸢e⸣ : C gravé ; ⸢t⸣ : D gravé ; ⸢c⸣ : E gravé ; \{r\} : lettres superflues}
\nf{Traduction}{Cerealis à Quiriacè (?), son épouse bien méritante, a fait (ce monument) ; déposition le 12e jour avant les calendes d’août, sous (le consulat de) notre seigneur Flavius Valentinien Auguste et Neoterius, clarissime.}
\nf{Langue}{L’éditeur signale 3 graphies corrigées : C gravé pour e ; D gravé pour t ; E gravé pour c. Des lettres gravées sont jugées superflues par l’éditeur (entre accolades). 13 abréviations développées par l’éditeur (dont eal(is), me(re)nti, d(epositio), Kal(endas), Aug(ustas), d(omino)).}
\nf{Chronologie}{Noms consulaires : Valentinianus (lu) et Neoterius (lu). Ne sont pas des consuls : C<e=C>r\{r\}eal(is) bene me(re)nti coniugi (dédicant (époux)). Dates du jour : déposition : XII Kalendas Augustas = 21 juillet 390 (lu).}
\nf{Bibliographie}{ICUR I 1445 ; AnalEpi p 377. Bibliographie complémentaire selon l’EDB : Lega 2000, 52 ; ILCV, 2946 adn.}
\nf{Statut}{Lecture secondaire (EDCS) concordante avec l’EDB : lecture stabilisée entre deux bases, non collationnée sur l’ICVR ni sur le monument. L’EDCS signale une photographie (non consultée).}
\end{notice}

\begin{notice}{320}{390}{ICUR n.s. I, 2102 · EDCS-39300321 · EDB32090}
\label{not:320}
\nf{Datation}{390 — directe : nom consulaire lu}
\nf{Genre}{Épitaphe ; ensembles A, B, C ; 105 lettres transmises.}
\nf{Références}{ICUR I 2102 = ILCV 3501.}
\nf{Provenance}{Rattachement selon l’EDB : \emph{Incertae originis} (le volume I des ICVR réunit les inscriptions dont la provenance est perdue). Conservation selon l’EDB : \emph{Musei Vaticani: Galleria Lapidaria}.}
\nf{Support}{Plaque de marbre. 52 × 98 cm (h. × l.). Inscription gravée. Fiche EDB : Filippo Piazzolla, première version 2015.}
\ntext{\lnum{1}Maximimasi(a)e\\
\lnum{2}viva c⸢um⸣\\
\lnum{3}co(n)iuge suo Hila\\
\lnum{4}ro fecit locum\\
\lnum{5}bene merenti\\
\lnum{6}in pace\{m\}\\
\lnum{7}d(e)positus XV Kal(endas)\\
\lnum{8}Aug(ustas)\\
\lnum{9}Consulatu\{m\} d(omini) n(ostri)\\
\lnum{10}Valen(t)i(ni)ani et\\
\lnum{11}Ne⸢o⸣teri}
\nf{Témoins}{Export EDCS de 2020 : \emph{Maximimasi(a)e / viva c / co(n)iuge suo Hila/ro fecit locum / bene merenti / in pace\{m\} / d(e)positus XV Kal(endas) / Aug(ustas) // Consulatu\{m\} d(omini) n(ostri) / Valen(t)i(ni)ani et / Neteri}. EDB (EDB32090) : texte concordant aux conventions de notation près (indice de similarité 0,935).}
\nf{Apparat}{⸢um⸣ : ON gravé ; ⸢o⸣ : V gravé ; \{m\} : lettres superflues ; \{m\} : lettres superflues}
\nf{Traduction}{Maximasia (?), de son vivant, avec son époux Hilarus, a fait un emplacement ; bien méritant, en paix ; déposé le 15e jour avant les calendes d’août // sous le consulat de notre seigneur Valentinien et de Neoterius.}
\nf{Langue}{L’éditeur signale 2 graphies corrigées : ON gravé pour um ; V gravé pour o (confusion o/u). Des lettres gravées sont jugées superflues par l’éditeur (entre accolades). 9 abréviations développées par l’éditeur (dont Maximimasi(a)e, co(n)iuge, d(e)positus, Kal(endas), Aug(ustas), d(omini)).}
\nf{Chronologie}{Noms consulaires : Valentinianus (lu) et Neoterius (lu). Dates du jour : déposition : XV Kalendas Augustas = 18 juillet 390 (lu).}
\nf{Commentaire}{La formule in pace est conservée.}
\nf{Bibliographie}{ICUR I 2102 ; ILCV 3501. Bibliographie complémentaire selon l’EDB : Lega 2000 ; ILCV, 3501.}
\nf{Statut}{Lecture secondaire (EDCS) concordante avec l’EDB : lecture stabilisée entre deux bases, non collationnée sur l’ICVR ni sur le monument.}
\end{notice}

\begin{notice}{321}{390}{ICUR n.s. II, 4778c · EDCS-20100025 · EDB42643}
\label{not:321}
\nf{Datation}{390 — directe : nom consulaire lu}
\nf{Genre}{Non funéraire (inscription architecturale) ; ensembles A, B ; 230 lettres transmises ; 36 \% restituées.}
\nf{Références}{ICUR II 4778 = ILCV 1857 (em) = AE 1959, 64 = AE 2000, 187 = AE 2003, 202 = AE 2009, 140 = Mossong-2022, 00027.}
\nf{Provenance}{Rattachement selon l’EDB : \emph{Basilica s.Pauli apostoli}. Conservation selon l’EDB : \emph{Basilica s.Pauli apostoli}.}
\nf{Support}{La fiche EDB ne figure pas dans l’état 2019 des données EpiDoc ; support non documenté.}
\ntext{\lnum{1}Columna Paul(i) a[postol(i)] natale X[III Kal(endas) Dec(embres) consulatu d(omini) n(ostri)] Valentin[i]ani Aug(usti) IIII et Neoteri v(iri) c(larissimi) administrante Fl(avio) Filippo vir[o clarissimo - - - curato]re Fl(avio) Anastasio [v(iro) c(larissimo) t]rib(uno) praetoria[no]\\
\lnum{2}Siricius episcopus tota mente devotus\\
\lnum{3}Episcopus devotus}
\nf{Témoins}{Export EDCS de 2020 : \emph{Columna Paul(i) a[postol(i)] natale X[III Kal(endas) Dec(embres) consulatu d(omini) n(ostri)] Valentin[i]ani Aug(usti) IIII et Neoteri v(iri) c(larissimi) administrante Fl(avio) Filippo vir[o clarissimo 3 curato]re Fl(avio) Anastasio [v(iro) c(larissimo) t]rib(uno) praetoria[no]}. EDB (EDB42643 ; divergence forte) : \emph{⟨: a below the capital⟩ / Siricius episcopus ((monogramma Christi cum alpha et omega)) tota mente devotus / / ⟨: b on the middle of the column⟩ / e̲p̲i̲s̲c̲o̲p̲u̲s̲ d̲e̲v̲o̲t̲u̲s̲ / / ⟨: c⟩ / columna Paul(i) a[postol(i)] natale X[III kal(endas) dec(embres) consulatu d(omini) n(ostri)] Valentin[i]ani Aug(usti) IIII et Neoteri v(iri) c(larissimi) administrante Fl(avio) Filippo v̲i̲r̲[o clarissimo --- curato]re Fl(avio) Anastasio [v(iro) c(larissimo) t]rib(uno) praetoria[no] [texte tronqué dans la liste de résultats]}. L'EDB (EDB42643) distingue trois textes sur la colonne de Saint-Paul : a, sous le chapiteau (\emph{Siricius episcopus ☧ tota mente devotus}), b, au milieu du fût (\emph{episcopus devotus}), c, la dédicace datée. Lecture de Filippi (2000) ; AE 2000, 187 ; EDR074191. L'EDB décrit un chrisme avec alpha et oméga après \emph{episcopus}.}
\nf{Traduction}{[Inscription non funéraire] Colonne de l’apôtre Paul ; jour natal (de dédicace) le 13e jour avant les calendes de décembre, sous le consulat de notre seigneur Valentinien Auguste pour la 4e fois et de Neoterius, clarissime ; Flavius Philippus, clarissime, […] administrant (les travaux) ; curateur Flavius Anastasius, clarissime, tribun prétorien // Siricius, évêque, dévoué de tout cœur // L’évêque dévoué.}
\nf{Langue}{L’export de 2020 diffère de l’état 2022 (voir apparat). 14 abréviations développées par l’éditeur (dont Paul(i), apostol(i), Kal(endas), Dec(embres), d(omini), n(ostri)).}
\nf{Chronologie}{Noms consulaires : Valentinianus (lu) et Neoterius (lu). Ne sont pas des consuls : administrante Fl(avio) Filippo vir[o clarissimo (fonctionnaire nommé dans l’inscription (non consul)); Columna Paul(i) a[postol(i)] (apôtre Paul). Dates du jour : événement indéterminé : XIII Kalendas Decembres = 19 novembre 390 (restitué).}
\nf{Commentaire}{Rang sénatorial d’une personne commémorée (uratore flavio anastasio viro clarissimo trib). Fonction ou métier mentionné : episcopus, Episcopus. L’EDB transcrit les signes suivants : chrisme. Commentaire développé : dossier p.~\pageref{dos:321}.}
\nf{Bibliographie}{ICUR II 4778 ; ILCV 1857 (em) ; AE 1959, 64 ; AE 2000, 187 ; AE 2003, 202 ; AE 2009, 140 ; Mossong-2022, 00027. Bibliographie complémentaire selon l’EDB : ILCV, 1857 ; Chastagnol 1957 ; AE, 1959, 64 ; Filippi 2000 ; AE, 2000, 187.}
\nf{Statut}{Lecture secondaire (EDCS) ; divergence forte avec l’EDB : le texte ne doit pas être exploité sans contrôle de l’édition imprimée.}
\end{notice}

\begin{notice}{322}{390}{ICUR n.s. II, 6047 · EDCS-34200673 · EDB11282}
\label{not:322}
\nf{Datation}{390 — directe : nom consulaire lu}
\nf{Genre}{Épitaphe ; ensembles A, B, C ; 68 lettres transmises ; 44 \% restituées.}
\nf{Références}{ICUR II 6047.}
\nf{Provenance}{Rattachement selon l’EDB : \emph{Coem. Commodillae}.}
\nf{Support}{Plaque de marbre. Inscription gravée. Reproductions signalées par l’EDB : photographie PCAS (droits réservés ; non reproduites). Fiche EDB : Carlo Carletti, première version 2003.}
\ntext{\lnum{1}]mum\\
\lnum{2}[in] pace q(ui) vixit\\
\lnum{3}[an]n(os) XXV d(ies) XV\\
\lnum{4}[- - -] Mai(as) Val(entiniano) IIII\\
\lnum{5}[e]t Ne⸢o⸣teri\\
\lnum{6}[o v(iro) c(larissimo) cons(ulibu)s]}
\nf{Témoins}{Export EDCS de 2020 : \emph{]mum / [in] pace q(ui) vixit / [an]n(os) XXV d(ies) XV / [3] Mai(as) Val(entiniano) IIII / [e]t Neteri/[o v(iro) c(larissimo) cons(ulibu)s]}. EDB (EDB11282) : texte concordant aux conventions de notation près (indice de similarité 0,939).}
\nf{Apparat}{⸢o⸣ : V gravé}
\nf{Traduction}{[…] en paix, qui vécut 25 ans 15 jours ; […] mai, sous le consulat de Valentinien pour la 4e fois et de Neoterius, clarissime.}
\nf{Langue}{L’éditeur signale une graphie corrigée : V gravé pour o (confusion o/u).}
\nf{Chronologie}{Noms consulaires : Valentinianus (lu) et Neoterius (lu). Âge déclaré : 25 a. 15 j. (partiellement restitué).}
\nf{Commentaire}{La formule in pace est attestée seulement par restitution.}
\nf{Bibliographie}{ICUR II 6047.}
\nf{Statut}{Lecture secondaire (EDCS) concordante avec l’EDB : lecture stabilisée entre deux bases, non collationnée sur l’ICVR ni sur le monument. L’EDCS signale une photographie (non consultée).}
\end{notice}

\begin{notice}{323}{390}{ICUR n.s. II, 6048 · EDCS-34200674 · EDB14639}
\label{not:323}
\nf{Datation}{390 — directe : nom consulaire lu}
\nf{Genre}{Épitaphe ; ensembles A, B, C ; 73 lettres transmises ; 3 \% restituées.}
\nf{Références}{ICUR II 6048 = ILCV 2977 = AE 1906, 85.}
\nf{Provenance}{Rattachement selon l’EDB : \emph{Coem. Commodillae}.}
\nf{Support}{Plaque de marbre. Inscription gravée. Fiche EDB : Carlo Carletti, première version 2015.}
\ntext{\lnum{1}D(omino) n(ostro) Valentiniano Aug(usto) III[I]\\
\lnum{2}et Ne⸢o⸣teri[o] v(iro) c(larissimo) cons(ulibus)\\
\lnum{3}deposita Benerosa\\
\lnum{4}die Kal(endis) Octobri(bu)s\\
\lnum{5}in pace}
\nf{Témoins}{Export EDCS de 2020 : \emph{D(omino) n(ostro) Valentiniano Aug(usto) III[I] / et Neteri[o] v(iro) c(larissimo) cons(ulibus) / deposita Benerosa / die Kal(endis) Octobri(bu)s / in pace}. EDB (EDB14639) : texte concordant aux conventions de notation près (indice de similarité 0,989).}
\nf{Apparat}{⸢o⸣ : V gravé}
\nf{Traduction}{Sous le consulat de notre seigneur Valentinien Auguste pour la 4e fois et de Neoterius, clarissime, déposée Venerosa, aux calendes d’octobre, en paix.}
\nf{Langue}{L’éditeur signale une graphie corrigée : V gravé pour o (confusion o/u).}
\nf{Chronologie}{Noms consulaires : Valentinianus (lu) et Neoterius (lu). Dates du jour : déposition : Kalendis Octobribus = 1 octobre 390 (lu).}
\nf{Commentaire}{La formule in pace est conservée.}
\nf{Bibliographie}{ICUR II 6048 ; ILCV 2977 ; AE 1906, 85.}
\nf{Statut}{Lecture secondaire (EDCS) concordante avec l’EDB : lecture stabilisée entre deux bases, non collationnée sur l’ICVR ni sur le monument.}
\end{notice}

\begin{notice}{324}{390}{ICUR n.s. II, 6049 · EDCS-34200675 · EDB10125}
\label{not:324}
\nf{Datation}{390 — directe : nom consulaire lu}
\nf{Genre}{Épitaphe ; ensembles A, B, C ; 164 lettres transmises.}
\nf{Références}{ICUR II 6049 = ILCV 4215 = ICaRoma 26 = Carletti-2008, 00095 = AE 1905, 81.}
\nf{Provenance}{Rattachement selon l’EDB : \emph{Coem. Commodillae}.}
\nf{Support}{Plaque de marbre. Inscription gravée. Signes décrits par l’EDB : chrisme avec alpha et oméga. Fiche EDB : Carlo Carletti, première version 2003.}
\ntext{\lnum{1}Bene merenti coniugi Hilaritati q(u)ae vixit\\
\lnum{2}ann(os) XXV sponsata ann(os) XI fuit sponsa ann(os) VII\\
\lnum{3}vixit cum marito ann(os) VII menses VIII deposita die XVII Kal(endas)\\
\lnum{4}Septemb(res) Fl(avio) Valentiniano Aug(usto) IIII et Neoterio v(iro) c(larissimo) conss(ulibus)}
\nf{Témoins}{EDB (EDB10125) : texte concordant aux conventions de notation près (indice de similarité 0,925).}
\nf{Traduction}{À Hilaritas, épouse bien méritante, qui vécut 25 ans ; fiancée à 11 ans, elle fut fiancée 7 ans ; elle vécut avec son mari 7 ans 8 mois ; déposée le 17e jour avant les calendes de septembre, sous le consulat de Flavius Valentinien Auguste pour la 4e fois et de Neoterius, clarissime.}
\nf{Langue}{12 abréviations développées par l’éditeur (dont q(u)ae, ann(os), Kal(endas)).}
\nf{Chronologie}{Noms consulaires : Valentinianus (lu) et Neoterius (lu). Dates du jour : déposition : die XVII Kalendas Septembres = 16 août 390 (lu). Âge déclaré : 25 a. (lu).}
\nf{Commentaire}{Commentaire développé : dossier p.~\pageref{dos:324}.}
\nf{Bibliographie}{ICUR II 6049 ; ILCV 4215 ; ICaRoma 26 ; Carletti-2008, 00095 ; AE 1905, 81. Bibliographie complémentaire selon l’EDB : Carletti 2008, 203, n. 95.}
\nf{Statut}{Lecture secondaire (EDCS) concordante avec l’EDB : lecture stabilisée entre deux bases, non collationnée sur l’ICVR ni sur le monument.}
\end{notice}

\begin{notice}{325}{390}{ICUR n.s. III, 8156 · EDCS-35700185 · EDB24678}
\label{not:325}
\nf{Datation}{390 — directe : lecture partielle d’un nom consulaire}
\nf{Genre}{Épitaphe ; ensembles A, B, C ; 62 lettres transmises ; 50 \% restituées.}
\nf{Références}{ICUR III 8156.}
\nf{Provenance}{Rattachement selon l’EDB : \emph{Basilica ss. Nerei et Achillei in coem. Domitillae}. Conservation selon l’EDB : \emph{Basilica ss. Nerei et Achillei in coem. Domitillae}.}
\nf{Support}{Brique de maçonnerie de briques. 23 × 27 cm (h. × l.). Inscription graffitée. Reproductions signalées par l’EDB : photographie PCAS ×2 (droits réservés ; non reproduites). Fiche EDB : Carlo Carletti / Antonio Enrico Felle, première version 2003.}
\ntext{\lnum{1}] depositus\\
\lnum{2}[- - - in] pace cons(ulatu)\\
\lnum{3}[Valentiniani] Aug(usti) IIII et Neute\\
\lnum{4}[rii v(iri) c(larissimi)]}
\nf{Témoins}{EDB (EDB24678) : texte concordant aux conventions de notation près (indice de similarité 0,984).}
\nf{Traduction}{Déposé […] en paix, sous le consulat de Valentinien Auguste pour la 4e fois et de Neoterius, clarissime.}
\nf{Langue}{4 abréviations développées par l’éditeur (dont cons(ulatu), Aug(usti), v(iri), c(larissimi)).}
\nf{Chronologie}{Noms consulaires : Valentinianus (restitué) et Neoterius (lu en partie).}
\nf{Commentaire}{La formule in pace est attestée seulement par restitution.}
\nf{Bibliographie}{ICUR III 8156. Bibliographie complémentaire selon l’EDB : IC I Supplementum, 1816.}
\nf{Statut}{Lecture secondaire (EDCS) concordante avec l’EDB : lecture stabilisée entre deux bases, non collationnée sur l’ICVR ni sur le monument. L’EDCS signale une photographie (non consultée).}
\end{notice}

\begin{notice}{326}{390}{ICUR n.s. III, 9344a · EDCS-35900808 · EDB24920}
\label{not:326}
\nf{Datation}{390 — conjecturale}
\nf{Genre}{Épitaphe ; ensembles A, B, C ; 4 lettres transmises.}
\nf{Références}{ICUR III 9344a.}
\nf{Provenance}{Rattachement selon l’EDB : \emph{Basilica ss. Nerei et Achillei in coem. Domitillae}.}
\nf{Support}{Plaque de marbre. Inscription gravée. Fiche EDB : Carlo Carletti / Antonio Enrico Felle, première version 2003.}
\ntext{\lnum{1}]eute[}
\nf{Témoins}{EDB (EDB24920 ; divergence forte) : \emph{[---] / [Valentiniano Aug(usto) IIII et N]eute[rio cons(ulibus)]}.}
\nf{Traduction}{[fragment : …eute…]}
\nf{Langue}{Aucune particularité graphique n’est signalée par l’éditeur dans l’état 2022 de la base.}
\nf{Chronologie}{Seule la séquence \emph{eute} est lue. L'EDB (EDB24920) restitue \emph{[Valentiniano Aug(usto) IIII et N]eute[rio cons(ulibus)]} (390). Cette restitution est plausible, mais elle repose sur quatre lettres. Datation conjecturale.}
\nf{Bibliographie}{ICUR III 9344a.}
\nf{Statut}{Lecture secondaire (EDCS) ; divergence forte avec l’EDB : le texte ne doit pas être exploité sans contrôle de l’édition imprimée.}
\end{notice}

\begin{notice}{327}{390}{ICUR n.s. IV, 9576 · EDCS-36101101 · EDB17189}
\label{not:327}
\nf{Datation}{390 — directe : lecture partielle d’un nom consulaire}
\nf{Genre}{Épitaphe ; ensembles A, B, C ; 54 lettres transmises ; 68 \% restituées.}
\nf{Références}{ICUR IV 9576.}
\nf{Provenance}{Rattachement selon l’EDB : \emph{Coem. Callisti pars inferior}.}
\nf{Support}{Plaque de marbre. Inscription gravée. Fiche EDB : Antonio Enrico Felle, première version 2003.}
\ntext{\lnum{1}di]es XV [- - -]\\
\lnum{2}[- - -]i domi[nis nostris Valentiniano Aug(usto) IIII]\\
\lnum{3}[et Ne]⸢o⸣terio [}
\nf{Témoins}{Export EDCS de 2020 : \emph{di]es XV [3] / [3]i domi[nis nostris Valentiniano Aug(usto) IIII] / [et Ne]terio [}. EDB (EDB17189 ; écarts à examiner) : \emph{[---] / [--- di]es XV [---] / [---]i domi[nis nostris Valentiniano Aug(usto) IIII] / [et Ne]uterio [---] [texte tronqué dans la liste de résultats]}.}
\nf{Apparat}{⸢o⸣ : V gravé}
\nf{Traduction}{[…] 15 jours […] ; sous (le consulat de) nos seigneurs (sic) Valentinien Auguste pour la 4e fois et Neoterius.}
\nf{Langue}{L’éditeur signale une graphie corrigée : V gravé pour o (confusion o/u).}
\nf{Chronologie}{Noms consulaires : Valentinianus (restitué) et Neoterius (lu en partie).}
\nf{Bibliographie}{ICUR IV 9576.}
\nf{Statut}{Lecture secondaire (EDCS) ; écarts avec l’EDB à arbitrer sur l’ICVR. Plus de la moitié des lettres sont restituées. L’EDCS signale une photographie (non consultée).}
\end{notice}

\annee{381}

\begin{notice}{328}{381}{ICUR n.s. IV, 9578 · EDCS-36101103 · EDB17191}
\label{not:328}
\nf{Datation}{381 — directe : lecture partielle d’un nom consulaire ; datation contestée}
\nf{Genre}{Épitaphe ; ensembles A, B, C ; 50 lettres transmises ; 72 \% restituées.}
\nf{Références}{ICUR IV 9578.}
\nf{Provenance}{Rattachement selon l’EDB : \emph{Coem. Callisti pars inferior}.}
\nf{Support}{Plaque de marbre. Inscription gravée. Fiche EDB : Antonio Enrico Felle, première version 2003.}
\ntext{\lnum{1}] depo[situs - - -]\\
\lnum{2}[- - - qu]i vixi[t - - -]\\
\lnum{3}[- - - Fl(avio) Syagrio et Eu]terio [conss(ulibus)]}
\nf{Témoins}{EDB (EDB17191 ; écarts à examiner) : \emph{[---] depo[situs ---] / [--- qu]i vixi[t ---] / [--- Fl(avio) Syagrio et Eu]terio [conss(ulibus)] [texte tronqué dans la liste de résultats]}.}
\nf{Traduction}{Déposé […] qui vécut […] ; sous le consulat de Flavius Syagrius et d’Euterius (sic).}
\nf{Langue}{2 abréviations développées par l’éditeur (dont Fl(avio), conss(ulibus)).}
\nf{Chronologie}{Noms consulaires : Syagrius (restitué). Seul \emph{…terio} est lu. La formule \emph{[Fl(avio) Syagrio et Eu]terio} est une restitution commune aux deux bases ; l'EDB (EDB17191) date pourtant le texte de 390, ce qui suppose \emph{[Valentiniano IIII et Neo]terio}. L'année dépend entièrement de la restitution : 381 (Eucherius) ou 390 (Neoterius). L'année 381 est conservée comme année de classement ; la datation est contestée.}
\nf{Bibliographie}{ICUR IV 9578.}
\nf{Statut}{Lecture secondaire (EDCS) ; écarts avec l’EDB à arbitrer sur l’ICVR. Plus de la moitié des lettres sont restituées.}
\end{notice}

\annee{391}

\begin{notice}{329}{391}{EDCS-07400001 · EDB10126 · EDH HD009698}
\label{not:329}
\nf{Statut}{Doublon établi : l'EDB (EDB10126 = ICVR II 6050) donne dans sa bibliographie l'AE 1987, 61, unique référence de cette notice ; l'EDH (HD009698) confirme le rattachement. La notice est fusionnée dans le n° 336. L'âge est lu \emph{L[1]V} par l'AE, \emph{LXV} par l'EDCS pour l'ICVR, \emph{L[X]V} par l'EDB : le chiffre X est restitué, et l'âge de 65 ans est donc partiellement restitué.}
\nf{Références}{AE 1987, 61}
\ntext{\lnum{1}[P]etili\{li\}ano bene merenti in pace qui vix\{s\}i\\
\lnum{2}anns n(umero) L[.]V d(epositus) in pace\{m\} V Idus Maias Tatiano et Symmac(h)o cons(ulibus)}
\nf{Traduction}{À Petilianus, bien méritant, en paix, qui vécut 54 (?) ans ; déposé en paix le 5e jour avant les ides de mai, sous le consulat de Tatianus et de Symmaque.}
\nf{Voir}{la notice 336 (p.~\pageref{not:336}).}
\end{notice}

\begin{notice}{330}{391}{ICUR n.s. VII, 19974 · EDCS-19900403 · EDB20295}
\label{not:330}
\nf{Datation}{391 — directe : nom consulaire lu}
\nf{Genre}{Épitaphe ; ensembles A, B, C ; 117 lettres transmises ; 33 \% restituées.}
\nf{Références}{CIL VI 32018 (p 4804) = CIL VI 39440b (p 4068) = ICUR VII 19974 = ILCV 197 (add).}
\nf{Provenance}{Rattachement selon l’EDB : \emph{Coem. s.Hippolyti}. Conservation selon l’EDB : \emph{Coem. s.Hippolyti}.}
\nf{Support}{Plaque de marbre. 20,5 × 34 × 2 cm (h. × l. × ép.). Inscription gravée. Reproductions signalées par l’EDB : photographie PCAS ×3 (droits réservés ; non reproduites). Fiche EDB : Domenico Schiraldi, première version 2005.}
\ntext{\lnum{1}[- - -] Egnatia Susanna\\
\lnum{2}Abita\\
\lnum{3}c(larissima) p(uella)\\
\lnum{4}[quae vi]xit annis quinque\\
\lnum{5}[mense] uno dies quinque\\
\lnum{6}[dep(osita) - - -] Kal(endas) Iunias\\
\lnum{7}[cons(ulibus) Fl(avio) Tati]ano et Aur(elio) Symmac(h)o\\
\lnum{8}p(ax) t(ibi) c(um) s(anctis)}
\nf{Témoins}{EDB (EDB20295 ; écarts à examiner) : \emph{[---] Egnatia Susanna c(larissima) p(uella) / Abita / [quae vi]xit annis quinque / [mense] uno dies quinque / [dep(osita) ---] kal(endas) iunias p(ax) t(ibi) c(um) s(anctis) / [Fl(avio) Tati]ano et Aur[elio] Symacho (consulibus)}.}
\nf{Traduction}{[…] Egnatia Susanna // Habita (?) // fillette clarissime, qui vécut cinq ans un mois cinq jours ; déposée […] calendes de juin, sous le consulat de Flavius Tatianus et d’Aurelius Symmachus // paix à toi avec les saints.}
\nf{Langue}{12 abréviations développées par l’éditeur (dont c(larissima), p(uella), dep(osita), Kal(endas), cons(ulibus), Fl(avio)).}
\nf{Chronologie}{Noms consulaires : Tatianus (lu en partie) et Symmachus (lu). Dates du jour : déposition : … Kalendas Iunias (partiellement restitué). Âge déclaré : 5 a. 1 m. 5 j. (partiellement restitué).}
\nf{Commentaire}{Vocabulaire religieux : saints (communion des saints) — sanctis (abréviation développée). Rang sénatorial d’une personne commémorée (… egnatia susanna abita clarissima puella quae).}
\nf{Bibliographie}{CIL VI 32018 (p 4804) ; CIL VI 39440b (p 4068) ; ICUR VII 19974 ; ILCV 197 (add). Bibliographie complémentaire selon l’EDB : ILCV, 197 ; CIL VI, 32018 ; IC I Supplementum, 1824.}
\nf{Statut}{Lecture secondaire (EDCS) ; écarts avec l’EDB à arbitrer sur l’ICVR. L’EDCS signale une photographie (non consultée).}
\end{notice}

\begin{notice}{331}{391}{ICUR n.s. I, 2103 · EDCS-39300322 · EDB32097}
\label{not:331}
\nf{Datation}{391 — directe : lecture partielle d’un nom consulaire}
\nf{Genre}{Épitaphe ; ensembles A, B, C ; 56 lettres transmises ; 46 \% restituées.}
\nf{Références}{ICUR I 2103.}
\nf{Provenance}{Rattachement selon l’EDB : \emph{Incertae originis} (le volume I des ICVR réunit les inscriptions dont la provenance est perdue). Conservation selon l’EDB : \emph{Musei Vaticani: Galleria Lapidaria}.}
\nf{Support}{Plaque de marbre. 33 × 23 cm (h. × l.). Inscription gravée. Fiche EDB : Filippo Piazzolla, première version 2015.}
\ntext{\lnum{1}]AL[- - -]\\
\lnum{2}[- - -] Kal(endas) Oct(obres) [- - -]\\
\lnum{3}[- - - Ta]tiano et Sy[mmacho cons(ulibus) - - -]\\
\lnum{4}[- - -] Criscent[- - -]\\
\lnum{5}[- - - I]ul(ias) dd(ominis) n[n(ostris)}
\nf{Témoins}{EDB (EDB32097 ; écarts à examiner) : \emph{------ / [---]+E[---] / [---]kal(end--) oct(obr--) [---] / [--- Ta]tiano et Su[mmacho cons(ulibu)s] / [--- et] Criscent[---] / [--- i]ul(ias?) dd(ominis?) n[n(ostris ?) ---]}.}
\nf{Traduction}{[…] calendes d’octobre […] ; sous le consulat de Tatianus et de Symmaque […] ; Crescent[…] ; […] juillet, (sous) nos seigneurs […].}
\nf{Langue}{6 abréviations développées par l’éditeur (dont Kal(endas), Oct(obres), cons(ulibus), Iul(ias), dd(ominis), nn(ostris)).}
\nf{Chronologie}{Noms consulaires : Tatianus (lu en partie) et Symmachus (restitué). Dates du jour : événement indéterminé : … Kalendas Octobres (partiellement restitué).}
\nf{Bibliographie}{ICUR I 2103. Bibliographie complémentaire selon l’EDB : Lega 2000.}
\nf{Statut}{Lecture secondaire (EDCS) ; écarts avec l’EDB à arbitrer sur l’ICVR.}
\end{notice}

\begin{notice}{332}{391}{ICUR n.s. I, 3209 · EDCS-40300347 · EDB33157}
\label{not:332}
\nf{Datation}{391 — directe : nom consulaire lu}
\nf{Genre}{Épitaphe ; ensembles A, B, C ; 62 lettres transmises ; 47 \% restituées.}
\nf{Références}{ICUR I 3209.}
\nf{Provenance}{Rattachement selon l’EDB : \emph{Incertae originis} (le volume I des ICVR réunit les inscriptions dont la provenance est perdue).}
\nf{Support}{Support non précisé par l’EDB. Fiche EDB : Filippo Piazzolla, première version 2015.}
\ntext{\lnum{1}]ISIV bene [merenti - - -]\\
\lnum{2}[- - - N]o⸢v⸣(embres ?) cons(ulibus) Tatian[o et Symmacho - - -]\\
\lnum{3}[- - - vix(it) annos - - -]XIII menses III[}
\nf{Témoins}{Export EDCS de 2020 : \emph{]ISIV bene [merenti 3] / [3 N]o(embres ?) cons(ulibus) Tatian[o et Symmacho 3] / [3 vix(it) annos 3]XIII menses III[}. EDB (EDB33157 ; écarts à examiner) : \emph{------ / [---]ịṣịụ bene[merenti ---] / [--- n]ọb(ilissimis?) cons(ulibus) Tatiaṇ[o et Symmacho ---] / [vixit annos ---]XIII menses III [---] / ------ [texte tronqué dans la liste de résultats]}.}
\nf{Apparat}{⸢v⸣ : B gravé}
\nf{Traduction}{[…] bien méritant […] ; […] novembre, sous le consulat de Tatianus et de Symmaque […] ; vécut […]13 ans 3 mois.}
\nf{Langue}{L’éditeur signale une graphie corrigée : B gravé pour v (bêtacisme).}
\nf{Chronologie}{Noms consulaires : Tatianus (lu) et Symmachus (restitué). Âge déclaré : 3 m. (partiellement restitué).}
\nf{Bibliographie}{ICUR I 3209.}
\nf{Statut}{Lecture secondaire (EDCS) ; écarts avec l’EDB à arbitrer sur l’ICVR.}
\end{notice}

\begin{notice}{333}{391}{ICUR n.s. II, 4169 · EDCS-33900079 · EDB16943}
\label{not:333}
\nf{Datation}{391 — directe : nom consulaire lu}
\nf{Genre}{Épitaphe ; ensembles A, B, C ; 103 lettres transmises ; 16 \% restituées.}
\nf{Références}{ICUR II 4169 = ILCV +569.}
\nf{Provenance}{Rattachement selon l’EDB : \emph{Basilica S.Petri apostoli}.}
\nf{Support}{Plaque de marbre. Inscription gravée. Reproductions signalées par l’EDB : reproduction, Perret 1851 (droits réservés ; non reproduites). Fiche EDB : Antonio Enrico Felle, première version 2003.}
\ntext{\lnum{1}Hic Leo requiescet omniu[m - - -]\\
\lnum{2}bonorum consiliator q(ui) an[nos vixit - - -]\\
\lnum{3}Tatiano et Symmacho conss(ulibus) [- - -]\\
\lnum{4}et eius A⸢l⸣exand[ra coniux ?] a s⸢t⸣a⸢t⸣uis i[}
\nf{Témoins}{Export EDCS de 2020 : \emph{Hic Leo requiescet omniu[m 3] / bonorum consiliator q(ui) an[nos vixit 3] / Tatiano et Symmacho conss(ulibus) [3] / et eius Aexand[ra coniux ?] a sauis i[}. EDB (EDB16943) : texte concordant aux conventions de notation près (indice de similarité 0,984).}
\nf{Apparat}{⸢l⸣ : I gravé ; ⸢t⸣ : I gravé ; ⸢t⸣ : I gravé}
\nf{Traduction}{Ici repose Leo, conseiller de tous les gens de bien, qui vécut […] ans ; sous le consulat de Tatianus et de Symmaque […] ; et son épouse Alexandra, (du quartier) des statues (?) […].}
\nf{Langue}{L’éditeur signale 3 graphies corrigées : I gravé pour l ; I gravé pour t ×2.}
\nf{Chronologie}{Noms consulaires : Tatianus (lu) et Symmachus (lu).}
\nf{Commentaire}{Vocabulaire religieux : repos et sommeil de la mort — requiescet (lu).}
\nf{Bibliographie}{ICUR II 4169 ; ILCV +569.}
\nf{Statut}{Lecture secondaire (EDCS) concordante avec l’EDB : lecture stabilisée entre deux bases, non collationnée sur l’ICVR ni sur le monument. L’EDCS signale une photographie (non consultée).}
\end{notice}

\begin{notice}{334}{391}{ICUR n.s. II, 4751 · EDCS-34201007 · EDB13473}
\label{not:334}
\nf{Datation}{391 — directe : lecture partielle d’un nom consulaire}
\nf{Genre}{Épitaphe ; ensembles A, B, C ; 46 lettres transmises ; 26 \% restituées.}
\nf{Références}{ICUR II 4751.}
\nf{Provenance}{Rattachement selon l’EDB : \emph{Coem. Generosae}. Conservation selon l’EDB : \emph{coem. Generosae}.}
\nf{Support}{Plaque de marbre. Inscription gravée. Reproductions signalées par l’EDB : photographie PCAS (droits réservés ; non reproduites). Fiche EDB : Antonio Enrico Felle, première version 2003.}
\ntext{\lnum{1}se] v(i)vo sibi et [- - -]\\
\lnum{2}[- - -] q(ui) vix(it) ann(os) LV[- - -]\\
\lnum{3}[Tati]ano et Sy[mm]aco cons[ulibus]}
\nf{Témoins}{EDB (EDB13473) : texte concordant aux conventions de notation près (indice de similarité 0,992).}
\nf{Traduction}{De son vivant, pour lui-même et […] ; […] qui vécut 55 […] ans ; sous le consulat de Tatianus et de Symmaque.}
\nf{Langue}{4 abréviations développées par l’éditeur (dont v(i)vo, q(ui), vix(it), ann(os)).}
\nf{Chronologie}{Noms consulaires : Tatianus (lu en partie).}
\nf{Bibliographie}{ICUR II 4751.}
\nf{Statut}{Lecture secondaire (EDCS) concordante avec l’EDB : lecture stabilisée entre deux bases, non collationnée sur l’ICVR ni sur le monument. L’EDCS signale une photographie (non consultée).}
\end{notice}

\begin{notice}{335}{391}{ICUR n.s. II, 5935 · EDCS-38100792 · EDB21460}
\label{not:335}
\nf{Datation}{391 — directe : nom consulaire lu}
\nf{Genre}{Épitaphe ; ensembles A, B, C ; 75 lettres transmises ; 69 \% restituées.}
\nf{Références}{ICUR II 5935.}
\nf{Provenance}{Rattachement selon l’EDB : \emph{Museum paulianum}.}
\nf{Support}{Support non précisé par l’EDB. Fiche EDB : Maria Cristina Grisanzio, première version 2006.}
\ntext{\lnum{1}[- - -] Sil[- - - vixit annos - - -]\\
\lnum{2}[mense]s VII dies V c[onsulatu]\\
\lnum{3}[Symmach]i et Tatiani [vv(irorum) cc(larissimorum)]}
\nf{Témoins}{EDB (EDB21460) : texte concordant aux conventions de notation près (indice de similarité 0,948).}
\nf{Traduction}{[…] vécut […] ans 7 mois 5 jours ; sous le consulat de Symmaque et de Tatianus, clarissimes.}
\nf{Langue}{2 abréviations développées par l’éditeur (dont vv(irorum), cc(larissimorum)).}
\nf{Chronologie}{Noms consulaires : Tatianus (lu) et Symmachus (restitué). Âge déclaré : 7 m. 5 j. (partiellement restitué).}
\nf{Bibliographie}{ICUR II 5935.}
\nf{Statut}{Lecture secondaire (EDCS) concordante avec l’EDB : lecture stabilisée entre deux bases, non collationnée sur l’ICVR ni sur le monument. Plus de la moitié des lettres sont restituées.}
\end{notice}

\begin{notice}{336}{391}{ICUR n.s. II, 6050 · EDCS-34200676 · EDB10126 · EDH HD009698}
\label{not:336}
\nf{Datation}{391 — directe : nom consulaire lu}
\nf{Genre}{Épitaphe ; ensembles A, B, C ; 81 lettres transmises.}
\nf{Références}{AE 1987, 61 = ICUR II 6050 = ILCV 2971 (em).}
\nf{Provenance}{Rattachement selon l’EDB : \emph{Coem. Commodillae}. Conservation selon l’EDB : \emph{Ostia Antica   ; Lapidario}.}
\nf{Support}{Plaque de marbre. Inscription gravée. Fiche EDB : Carlo Carletti, première version 2003.}
\ntext{\lnum{1}Petiliano bene merenti in pace qui vix\{s\}i⸢t⸣\\
\lnum{2}ann⸢o⸣s n(umero) LXV d(epositus) in pace\{m\}\\
\lnum{3}V Idus Maias Tatiano et Symmac(h)o cons(ulibu)s}
\nf{Témoins}{Export EDCS de 2020 : \emph{Petiliano bene merenti in pace qui vix\{s\}i / anns n(umero) LXV d(epositus) in pace\{m\} / V Idus Maias Tatiano et Symmac(h)o cons(ulibu)s}. EDB (EDB10126) : texte concordant aux conventions de notation près (indice de similarité 0,946). Le texte publié dans l'AE 1987, 61 (ancien n° 329) est le même monument. L'EDB (EDB10126) signale deux références que le catalogue antérieur ignorait : AE 1991, 361 et Marinucci 1991, p. 104–105, n° 26. La plaque, en deux fragments, est conservée au lapidaire d'Ostie (\emph{Ostia Antica, Lapidario}) ; l'âge est lu \emph{L[X]V} par l'EDB, avec un X restitué. Absorbe la notice 329 (voir cette notice).}
\nf{Apparat}{⸢t⸣ : S gravé ; ⸢o⸣ : V gravé ; \{s\} : lettres superflues ; \{m\} : lettres superflues}
\nf{Traduction}{À Petilianus, bien méritant, en paix, qui vécut 65 ans ; déposé en paix le 5e jour avant les ides de mai, sous le consulat de Tatianus et de Symmaque.}
\nf{Langue}{L’éditeur signale 2 graphies corrigées : S gravé pour t ; V gravé pour o (confusion o/u). Des lettres gravées sont jugées superflues par l’éditeur (entre accolades). 4 abréviations développées par l’éditeur (dont n(umero), d(epositus), Symmac(h)o, cons(ulibu)s).}
\nf{Chronologie}{Noms consulaires : Tatianus (lu) et Symmachus (lu). Dates du jour : déposition : V Idus Maias = 11 mai 391 (lu). Âge déclaré : 65 a. (lu).}
\nf{Commentaire}{La formule in pace est conservée. La valeur de soixante-cinq ans demeure celle de cet état transmis, sans choix prétendument fondé sur une collation.}
\nf{Bibliographie}{AE 1987, 61 ; ICUR II 6050 ; ILCV 2971 (em). Bibliographie complémentaire selon l’EDB : AE, 1987, 61 ; Paci 1985, 542-544 ; Marinucci 1991, 104-105, n. 26 ; AE, 1991, 361.}
\nf{Statut}{Lecture secondaire (EDCS) concordante avec l’EDB : lecture stabilisée entre deux bases, non collationnée sur l’ICVR ni sur le monument. L’EDCS signale une photographie (non consultée).}
\end{notice}

\begin{notice}{337}{391}{ICUR n.s. II, 6051 · EDCS-34200677 · EDB14891}
\label{not:337}
\nf{Datation}{391 — directe : nom consulaire lu}
\nf{Genre}{Épitaphe ; ensembles A, B, C ; 109 lettres transmises ; 62 \% restituées.}
\nf{Références}{ICUR II 6051.}
\nf{Provenance}{Rattachement selon l’EDB : \emph{Coem. Commodillae}. Conservation selon l’EDB : \emph{Coem. Commodillae}.}
\nf{Support}{Plaque de marbre. Inscription gravée. Reproductions signalées par l’EDB : photographie PCAS (droits réservés ; non reproduites). Fiche EDB : Carlo Carletti, première version 2003.}
\ntext{\lnum{1}F(lavio) Tatian[o et Aur(elio) Sym]\\
\lnum{2}mac(h)o v[v(iris) cc(larissimis) cons(ulibu)s]\\
\lnum{3}[K]al(endis) Feb(ruariis) d⸢e⸣p[osita est]\\
\lnum{4}[- - -]dafini qu(a)e vi[xit]\\
\lnum{5}[p(lus) m(inus)] ann(os) XLI m(enses) V d(ies) II [in p(ace ?)]}
\nf{Témoins}{Export EDCS de 2020 : \emph{F(lavio) Tatian[o et Aur(elio) Sym]/mac(h)o v[v(iris) cc(larissimis) cons(ulibu)s] / [K]al(endis) Feb(ruariis) dp[osita est] / [3]dafini qu(a)e vi[xit] / [p(lus) m(inus)] ann(os) XLI m(enses) V d(ies) II [in p(ace ?)]}. EDB (EDB14891) : texte concordant aux conventions de notation près (indice de similarité 0,967).}
\nf{Apparat}{⸢e⸣ : I gravé}
\nf{Traduction}{Sous le consulat de Flavius Tatianus et d’Aurelius Symmachus, clarissimes, aux calendes de février, a été déposée […]dafinè (?), qui vécut environ 41 ans 5 mois 2 jours ; en paix.}
\nf{Langue}{L’éditeur signale une graphie corrigée : I gravé pour e (confusion e/i).}
\nf{Chronologie}{Noms consulaires : Tatianus (lu) et Symmachus (lu en partie). Dates du jour : événement indéterminé : Kalendis Februariis = 1 février 391 (partiellement restitué). Âge déclaré : env. 41 a. 5 m. 2 j. (partiellement restitué).}
\nf{Commentaire}{La formule in pace est attestée seulement par restitution.}
\nf{Bibliographie}{ICUR II 6051.}
\nf{Statut}{Lecture secondaire (EDCS) concordante avec l’EDB : lecture stabilisée entre deux bases, non collationnée sur l’ICVR ni sur le monument. Plus de la moitié des lettres sont restituées. L’EDCS signale une photographie (non consultée).}
\end{notice}

\begin{notice}{338}{391}{ICUR n.s. III, 8157 · EDCS-35700186 · EDB24679}
\label{not:338}
\nf{Datation}{391 — directe : nom consulaire lu}
\nf{Genre}{Épitaphe ; ensembles A, B, C ; 41 lettres transmises ; 46 \% restituées.}
\nf{Références}{ICUR III 8157.}
\nf{Provenance}{Rattachement selon l’EDB : \emph{Basilica ss. Nerei et Achillei in coem. Domitillae}. Conservation selon l’EDB : \emph{Basilica ss. Nerei et Achillei in coem. Domitillae}.}
\nf{Support}{Plaque de marbre. 29 × 34 × 3 cm (h. × l. × ép.). Inscription gravée. Fiche EDB : Carlo Carletti / Antonio Enrico Felle, première version 2003.}
\ntext{\lnum{1}Tatiano] et S⸢y⸣mma[cho conss(ulibus) - - -]\\
\lnum{2}[- - - q]ui vixit a[nn(os)}
\nf{Témoins}{Export EDCS de 2020 : \emph{Tatiano] et Smma[cho conss(ulibus) 3] / [3 q]ui vixit a[nn(os)}. EDB (EDB24679) : texte concordant aux conventions de notation près (indice de similarité 0,932).}
\nf{Apparat}{⸢y⸣ : V gravé}
\nf{Traduction}{Sous le consulat de Tatianus et de Symmaque […] qui vécut […] ans.}
\nf{Langue}{L’éditeur signale une graphie corrigée : V gravé pour y (graphie de y).}
\nf{Chronologie}{Noms consulaires : Tatianus (restitué) et Symmachus (lu).}
\nf{Bibliographie}{ICUR III 8157.}
\nf{Statut}{Lecture secondaire (EDCS) concordante avec l’EDB : lecture stabilisée entre deux bases, non collationnée sur l’ICVR ni sur le monument. L’EDCS signale une photographie (non consultée).}
\end{notice}

\begin{notice}{339}{391}{ICUR n.s. III, 8158 · EDCS-35700187 · EDB24680}
\label{not:339}
\nf{Datation}{391 — directe : nom consulaire lu}
\nf{Genre}{Épitaphe ; ensembles A, B, C ; 85 lettres transmises ; 44 \% restituées.}
\nf{Références}{ICUR III 8158 = ILCV +2132.}
\nf{Provenance}{Rattachement selon l’EDB : \emph{Basilica ss. Nerei et Achillei in coem. Domitillae}. Conservation selon l’EDB : \emph{Basilica ss. Nerei et Achillei in coem. Domitillae}.}
\nf{Support}{Plaque de marbre. 20 × 52 × 2,3 cm (h. × l. × ép.). Inscription gravée. Reproductions signalées par l’EDB : photographie PCAS (droits réservés ; non reproduites). Fiche EDB : Carlo Carletti / Antonio Enrico Felle, première version 2003.}
\ntext{\lnum{1}Accill[- - -]\\
\lnum{2}Laurentia Emerita Sam[batiae(?) - - -]\\
\lnum{3}quae vixit an[n(os) - - -]\\
\lnum{4}A(ulo) Tat(iano) et Q(uinto) S⸢y⸣m(m)a(cho) v[v(iris) cc(larissimis) co(n)ss(ulibus)}
\nf{Témoins}{Export EDCS de 2020 : \emph{Accill[3] / Laurentia Emerita Sam[batiae(?) 3] / quae vixit an[n(os) 3] / A(ulo) Tat(iano) et Q(uinto) Sm(m)a(cho) v[v(iris) cc(larissimis) co(n)ss(ulibus)}. EDB (EDB24680 ; écarts à examiner) : \emph{------ / Accilẹ[um? ---] / Emerita Sam[batiae? ---] / ((monogramma: Laurentia)) quae vixit an[n(os) ---] / A(ulo) Tat(iano) et Q(uinto) Suma(cho) v[v(iris) cc(larissimis)]}.}
\nf{Apparat}{⸢y⸣ : V gravé}
\nf{Traduction}{[…] Laurentia Emerita (fille ?) de Sambatia […] qui vécut […] ans ; sous le consulat d’Aulus (sic) Tatianus et de Quintus Symmachus, clarissimes.}
\nf{Langue}{L’éditeur signale une graphie corrigée : V gravé pour y (graphie de y).}
\nf{Chronologie}{Noms consulaires : Tatianus (lu) et Symmachus (lu).}
\nf{Bibliographie}{ICUR III 8158 ; ILCV +2132. Bibliographie complémentaire selon l’EDB : IC I Supplementum, 1828 ; ILCV, 2132 adn.}
\nf{Statut}{Lecture secondaire (EDCS) ; écarts avec l’EDB à arbitrer sur l’ICVR. L’EDCS signale une photographie (non consultée).}
\end{notice}

\begin{notice}{340}{391}{ICUR n.s. III, 8428 · EDCS-35802015 · EDB25702}
\label{not:340}
\nf{Datation}{391 — directe : nom consulaire lu}
\nf{Genre}{Acte d’achat ou de réservation ; ensembles A, B ; 80 lettres transmises ; 1 \% restituées.}
\nf{Références}{ICUR III 8428 = ILCV 3821 (em).}
\nf{Provenance}{Rattachement selon l’EDB : \emph{Incertum locum in coem. Domitillae}. Conservation selon l’EDB : \emph{coem. Domitillae, aula supra coemeterium}.}
\nf{Support}{Plaque de marbre. Inscription gravée. Reproductions signalées par l’EDB : photographie PCAS (droits réservés ; non reproduites). Fiche EDB : Carlo Carletti / Antonio Enrico Felle, première version 2003.}
\ntext{\lnum{1}Fl(avio) Tatiano et Quinto\\
\lnum{2}Aur(elio) S⸢y⸣mmaco v⸢i⸣ris\\
\lnum{3}cl(a)rissimis ego Zita\\
\lnum{4}locum quadri\{c\}\\
\lnum{5}somu(m) in b(a)silic(a)\\
\lnum{6}al⸢b⸣a emi[t]}
\nf{Témoins}{Export EDCS de 2020 : \emph{Fl(avio) Tatiano et Quinto / Aur(elio) Smmaco vris / cl(a)rissimis ego Zita / locum quadri\{c\}/somu(m) in b(a)silic(a) / ala emi[t]}. EDB (EDB25702) : texte concordant aux conventions de notation près (indice de similarité 0,954).}
\nf{Apparat}{⸢y⸣ : V gravé ; ⸢i⸣ : C gravé ; ⸢b⸣ : V gravé ; \{c\} : lettres superflues}
\nf{Traduction}{Sous (le consulat de) Flavius Tatianus et Quintus Aurelius Symmachus, clarissimes, moi, Zita, j’ai acheté un emplacement quadruple dans la basilique blanche (?).}
\nf{Langue}{L’éditeur signale 3 graphies corrigées : V gravé pour y (graphie de y) ; C gravé pour i ; V gravé pour b (bêtacisme). Des lettres gravées sont jugées superflues par l’éditeur (entre accolades). 6 abréviations développées par l’éditeur (dont Fl(avio), Aur(elio), cl(a)rissimis, somu(m), b(a)silic, (a)).}
\nf{Chronologie}{Noms consulaires : Tatianus (lu).}
\nf{Bibliographie}{ICUR III 8428 ; ILCV 3821 (em). Bibliographie complémentaire selon l’EDB : ILCV, 3821.}
\nf{Statut}{Lecture secondaire (EDCS) concordante avec l’EDB : lecture stabilisée entre deux bases, non collationnée sur l’ICVR ni sur le monument. L’EDCS signale une photographie (non consultée).}
\end{notice}

\begin{notice}{341}{391}{ICUR n.s. IV, 11121 · EDCS-37200404 · EDB39583}
\label{not:341}
\nf{Datation}{391 — directe : lecture partielle d’un nom consulaire}
\nf{Genre}{Épitaphe ; ensembles A, B, C ; 38 lettres transmises ; 63 \% restituées.}
\nf{Références}{ICUR IV 11121.}
\nf{Provenance}{Rattachement selon l’EDB : \emph{Coem. Callisti area subdialis et incerti loci}. Conservation selon l’EDB : \emph{Coem. Callisti area subdialis et incerti loci}.}
\nf{Support}{Plaque de marbre. 14 × 17 cm (h. × l.). Inscription gravée. Fiche EDB : Marida Pierno, première version 2015.}
\ntext{\lnum{1}in p]ace qua[e - - -]\\
\lnum{2}[- - - Tatiano e]t Summ[acho conss(ulibus)]}
\nf{Témoins}{EDB (EDB39583 ; écarts à examiner) : \emph{[--- in p]ace qua[e ---] / [--- e]t Summ[acho]}.}
\nf{Traduction}{[…] en paix, qui […] ; sous le consulat de Tatianus et de Symmaque.}
\nf{Langue}{1 abréviation développée par l’éditeur (dont conss(ulibus)).}
\nf{Chronologie}{Noms consulaires : Tatianus (restitué).}
\nf{Commentaire}{La formule in pace est conservée.}
\nf{Bibliographie}{ICUR IV 11121.}
\nf{Statut}{Lecture secondaire (EDCS) ; écarts avec l’EDB à arbitrer sur l’ICVR. Plus de la moitié des lettres sont restituées.}
\end{notice}

\begin{notice}{342}{391}{ICUR n.s. IV, 11122 · EDCS-37200405 · EDB39584}
\label{not:342}
\nf{Datation}{391 — directe : lecture partielle d’un nom consulaire}
\nf{Genre}{Épitaphe ; ensembles A, B, C ; 33 lettres transmises ; 46 \% restituées.}
\nf{Références}{ICUR IV 11122 (p 534) = ICUR IV 11667u.}
\nf{Provenance}{Rattachement selon l’EDB : \emph{Coem. Callisti area subdialis et incerti loci}. Conservation selon l’EDB : \emph{Coem. Callisti area subdialis et incerti loci}.}
\nf{Support}{Plaque de marbre. Inscription gravée. Fiche EDB : Marida Pierno, première version 2015.}
\ntext{\lnum{1}Tatiano et S]ymm[acho conss(ulibus) - - -]\\
\lnum{2}[- - -]se N[- - -]\\
\lnum{3}[- - -]EV[}
\nf{Témoins}{EDB (EDB39584) : texte concordant aux conventions de notation près (indice de similarité 0,972).}
\nf{Traduction}{Sous le consulat de Tatianus et de Symmaque […].}
\nf{Langue}{1 abréviation développée par l’éditeur (dont conss(ulibus)).}
\nf{Chronologie}{Noms consulaires : Tatianus (restitué) et Symmachus (lu en partie).}
\nf{Bibliographie}{ICUR IV 11122 (p 534) ; ICUR IV 11667u.}
\nf{Statut}{Lecture secondaire (EDCS) concordante avec l’EDB : lecture stabilisée entre deux bases, non collationnée sur l’ICVR ni sur le monument.}
\end{notice}

\annee{392}

\begin{notice}{343}{392}{ICUR n.s. I, 986 · EDCS-38700493 · EDB29619}
\label{not:343}
\nf{Datation}{392 — directe : nom consulaire lu}
\nf{Genre}{Épitaphe ; ensembles A, B, C ; 85 lettres transmises ; 29 \% restituées.}
\nf{Références}{ICUR I 986.}
\nf{Provenance}{Rattachement selon l’EDB : \emph{Incertae originis} (le volume I des ICVR réunit les inscriptions dont la provenance est perdue). Conservation selon l’EDB : \emph{ss. Quattro Coronati}.}
\nf{Support}{Plaque de marbre. 70 × 110 × 4,5 cm (h. × l. × ép.). Inscription gravée. Fiche EDB : Filippo Piazzolla, première version 2014.}
\ntext{\lnum{1}Sextiliu[s - - -]orum [- - -]\\
\lnum{2}q(ui) v(ixit) ann(os) XLVI d[ep(ositus)] pr(idie) I[d]us Iu[lias - - -]\\
\lnum{3}d(omino) n(ostro) Arcadio II et Ruf[i]no v(iro) c(larissimo) conss(ulibus)\\
\lnum{4}ben(e) m[erenti] in [pace]}
\nf{Témoins}{EDB (EDB29619) : texte concordant aux conventions de notation près (indice de similarité 0,938).}
\nf{Traduction}{Sextilius […] qui vécut 46 ans ; déposé la veille des ides de juillet […] ; sous le consulat de notre seigneur Arcadius pour la 2e fois et de Rufinus, clarissime ; bien méritant, en paix.}
\nf{Langue}{11 abréviations développées par l’éditeur (dont q(ui), v(ixit), ann(os), dep(ositus), pr(idie), d(omino)).}
\nf{Chronologie}{Noms consulaires : Arcadius (lu) et Rufinus (lu en partie). Dates du jour : déposition : pridie Idus Iulias = 14 juillet 392 (partiellement restitué). Âge déclaré : 46 a. (partiellement restitué).}
\nf{Commentaire}{La formule in pace est attestée seulement par restitution.}
\nf{Bibliographie}{ICUR I 986. Bibliographie complémentaire selon l’EDB : ILCV, 2611 adn.}
\nf{Statut}{Lecture secondaire (EDCS) concordante avec l’EDB : lecture stabilisée entre deux bases, non collationnée sur l’ICVR ni sur le monument.}
\end{notice}

\begin{notice}{344}{392}{ICUR n.s. I, 1355 · EDCS-39000332 · EDB30697}
\label{not:344}
\nf{Datation}{392 — directe : nom consulaire lu}
\nf{Genre}{Épitaphe ; ensembles A, B, C ; 70 lettres transmises ; 23 \% restituées.}
\nf{Références}{ICUR I 1355.}
\nf{Provenance}{Rattachement selon l’EDB : \emph{Incertae originis} (le volume I des ICVR réunit les inscriptions dont la provenance est perdue). Conservation selon l’EDB : \emph{Collegio Teutonico}.}
\nf{Support}{Plaque de marbre. 24 × 24 cm (h. × l.). Inscription gravée. Signes décrits par l’EDB : croix monogrammatique et croix monogrammatique avec alpha et oméga. Fiche EDB : Filippo Piazzolla, première version 2015.}
\ntext{\lnum{1}ben]e merenti qu[- - -]\\
\lnum{2}[- - -]II et dies VII[- - -]\\
\lnum{3}[- - - i]n pace d(ie) VIIII Id[us - - -]\\
\lnum{4}[- - - Pall]adio II et Rufino v(iris) c(larissimis) c[ons(ulibus)}
\nf{Témoins}{EDB (EDB30697 ; écarts à examiner) : \emph{[bene]merenti qụ[i vixit annos ---] / [menses ---]II et dies VII[---] / ((piscis)) ((monogramma Christi in crucem translatum cum alpha et omega in corona taeniata)) / [--- i]n pace d(eposit---) VIIII id[us ---] / [Arc]adio II et Rufino v(iro) c(larissimo) c[ons(ulibu)s ---] [texte tronqué dans la liste de résultats]}.}
\nf{Traduction}{Bien méritant […] ; […] 2 (?) et 7 jours […] ; […] en paix, le 9e jour avant les ides […] ; sous le consulat de Palladius (sic) pour la 2e fois et de Rufinus, clarissimes.}
\nf{Langue}{4 abréviations développées par l’éditeur (dont d(ie), v(iris), c(larissimis), cons(ulibus)).}
\nf{Chronologie}{Noms consulaires : Rufinus (lu).}
\nf{Commentaire}{La formule in pace est attestée seulement par restitution.}
\nf{Bibliographie}{ICUR I 1355. Bibliographie complémentaire selon l’EDB : Ferrua 1986c, 46, n. 1355.}
\nf{Statut}{Lecture secondaire (EDCS) ; écarts avec l’EDB à arbitrer sur l’ICVR.}
\end{notice}

\begin{notice}{345}{392}{ICUR n.s. I, 2104 · EDCS-39300323 · EDB32098}
\label{not:345}
\nf{Datation}{392 — directe : nom consulaire lu}
\nf{Genre}{Épitaphe ; ensembles A, B, C ; 37 lettres transmises ; 57 \% restituées.}
\nf{Références}{ICUR I 2104.}
\nf{Provenance}{Rattachement selon l’EDB : \emph{Incertae originis} (le volume I des ICVR réunit les inscriptions dont la provenance est perdue). Conservation selon l’EDB : \emph{Musei Vaticani: Galleria Lapidaria}.}
\nf{Support}{Plaque de marbre. 18 × 24 cm (h. × l.). Inscription gravée. Fiche EDB : Filippo Piazzolla, première version 2015.}
\ntext{\lnum{1}]EI[- - -]\\
\lnum{2}[- - -]XII depos[- - -]\\
\lnum{3}[- - - Arcadio II et] Rufino [cons(ulibus)}
\nf{Témoins}{EDB (EDB32098) : texte concordant aux conventions de notation près (indice de similarité 0,938).}
\nf{Traduction}{[…] 12 […] déposé(e) […] ; sous le consulat d’Arcadius pour la 2e fois et de Rufinus.}
\nf{Langue}{1 abréviation développée par l’éditeur (dont cons(ulibus)).}
\nf{Chronologie}{Noms consulaires : Arcadius (restitué) et Rufinus (lu).}
\nf{Bibliographie}{ICUR I 2104. Bibliographie complémentaire selon l’EDB : Lega 2000.}
\nf{Statut}{Lecture secondaire (EDCS) concordante avec l’EDB : lecture stabilisée entre deux bases, non collationnée sur l’ICVR ni sur le monument. Plus de la moitié des lettres sont restituées.}
\end{notice}

\begin{notice}{346}{392}{ICUR n.s. I, 2980 · EDCS-39900424 · EDB32879}
\label{not:346}
\nf{Datation}{392 — directe : nom consulaire lu}
\nf{Genre}{Mixte (épitaphe et acte d’achat) ; ensembles A, B, C ; 114 lettres transmises ; 24 \% restituées.}
\nf{Références}{ICUR I 2980 = ILCV 3727a.}
\nf{Provenance}{Rattachement selon l’EDB : \emph{Incertae originis} (le volume I des ICVR réunit les inscriptions dont la provenance est perdue). Conservation selon l’EDB : \emph{Ravenna   ; Museo Nazionale}.}
\nf{Support}{Plaque de marbre. 31 × 61 × 3,5 cm (h. × l. × ép.). Inscription gravée. Fiche EDB : Filippo Piazzolla, première version 2015.}
\ntext{\lnum{1}[C]⸢o⸣mparaverun[t sibi locum]\\
\lnum{2}se vivi Iustinus et Feli[citas]\\
\lnum{3}dep(osita) Pitzinnina in pace R[- - -]\\
\lnum{4}qu(a)e vixit ann(os) XXVI et [- - -]\\
\lnum{5}Arcadio Aug(usto) II et Rufin[o cons(ulibus)}
\nf{Témoins}{Export EDCS de 2020 : \emph{[C]mparaverun[t sibi locum] / se vivi Iustinus et Feli[citas] / dep(osita) Pitzinnina in pace R[3] / qu(a)e vixit ann(os) XXVI et [3] / Arcadio Aug(usto) II et Rufin[o cons(ulibus)}. EDB (EDB32879) : texte concordant aux conventions de notation près (indice de similarité 0,982).}
\nf{Apparat}{⸢o⸣ : V gravé}
\nf{Traduction}{Iustinus et Felicitas se sont acquis de leur vivant un emplacement ; déposée Pitzinnina, en paix […], qui vécut 26 ans et […] ; sous le consulat d’Arcadius Auguste pour la 2e fois et de Rufinus.}
\nf{Langue}{L’éditeur signale une graphie corrigée : V gravé pour o (confusion o/u).}
\nf{Chronologie}{Noms consulaires : Arcadius (lu) et Rufinus (lu). Âge déclaré : 26 a. (lu).}
\nf{Commentaire}{La formule in pace est conservée.}
\nf{Bibliographie}{ICUR I 2980 ; ILCV 3727a. Bibliographie complémentaire selon l’EDB : ILCV, 3727A ; Ferrua 1979b, 85.}
\nf{Statut}{Lecture secondaire (EDCS) concordante avec l’EDB : lecture stabilisée entre deux bases, non collationnée sur l’ICVR ni sur le monument.}
\end{notice}

\begin{notice}{347}{392}{ICUR n.s. II, 4271a · EDCS-33900177 · EDB11967}
\label{not:347}
\nf{Datation}{392 — directe : nom consulaire lu}
\nf{Genre}{Épitaphe ; ensembles A, B, C ; 155 lettres transmises.}
\nf{Références}{ICUR II 4271a = ILCV 3502.}
\nf{Provenance}{Rattachement selon l’EDB : \emph{Coem. Octavillae ad s.Pancratium}.}
\nf{Support}{Plaque de marbre. Inscription gravée. Fiche EDB : Antonio Enrico Felle, première version 2003.}
\ntext{\lnum{1}D(omino) n(ostro) Arcadio Aug(usto) II et Rufino con\\
\lnum{2}ss(ulibus) Id(ib)us Septembri(bu)s depositus puer\\
\lnum{3}Heraclius in pace qui ⸢v⸣ixit ann⸢o⸣s VII\\
\lnum{4}mens⸢e⸣s V dies XXVII Celerinus et Lea se\\
\lnum{5}vivis sibi et filio dulcissimo locu(m) f⸢e⸣cerunt}
\nf{Témoins}{Export EDCS de 2020 : \emph{D(omino) n(ostro) Arcadio Aug(usto) II et Rufino con/ss(ulibus) Id(ib)us Septembri(bu)s depositus puer / Heraclius in pace qui ixit anns VII / menss V dies XXVII Celerinus et Lea se / vivis sibi et filio dulcissimo locu(m) fcerunt}. EDB (EDB11967) : texte concordant aux conventions de notation près (indice de similarité 0,979).}
\nf{Apparat}{⸢v⸣ : B gravé ; ⸢o⸣ : V gravé ; ⸢e⸣ : I gravé ; ⸢e⸣ : I gravé}
\nf{Traduction}{Sous le consulat de notre seigneur Arcadius Auguste pour la 2e fois et de Rufinus, aux ides de septembre, a été déposé le garçon Heraclius, en paix, qui vécut 7 ans 5 mois 27 jours ; Celerinus et Lea, de leur vivant, ont fait un emplacement pour eux et pour leur fils très doux.}
\nf{Langue}{L’éditeur signale 4 graphies corrigées : B gravé pour v (bêtacisme) ; V gravé pour o (confusion o/u) ; I gravé pour e (confusion e/i) ×2.}
\nf{Chronologie}{Noms consulaires : Arcadius (lu) et Rufinus (lu). Dates du jour : événement indéterminé : Idibus Septembribus = 13 septembre 392 (lu). Âge déclaré : 7 a. 5 m. 27 j. (lu).}
\nf{Commentaire}{La formule in pace est conservée.}
\nf{Bibliographie}{ICUR II 4271a ; ILCV 3502.}
\nf{Statut}{Lecture secondaire (EDCS) concordante avec l’EDB : lecture stabilisée entre deux bases, non collationnée sur l’ICVR ni sur le monument.}
\end{notice}

\begin{notice}{348}{392}{ICUR n.s. II, 4271b · EDCS-25401950 · EDB11272}
\label{not:348}
\nf{Datation}{392 — indirecte}
\nf{Genre}{Épitaphe ; ensembles A, B, C ; 98 lettres transmises.}
\nf{Références}{ICUR II 4271b.}
\nf{Provenance}{Rattachement selon l’EDB : \emph{Coem. Octavillae ad s.Pancratium}.}
\nf{Support}{Plaque de marbre. Inscription gravée. Signes décrits par l’EDB : croix monogrammatique. Fiche EDB : Antonio Enrico Felle, première version 2003.}
\ntext{\lnum{1}ἐνθάδε κατάκι\{τι\}τε Ἀλέξα(ν)δρο(ς) ζήσας\\
\lnum{2}ἔτη εἴκοσι τελευτᾷ μηνὶ νοενβρί\\
\lnum{3}ου υ\{ε\}ἱὸς Ἀ(μ)βρωσίου (χ)ωρίῳ Μικρᾶς Κ\\
\lnum{4}ώμης ζέσας μ(ε)τὰ τῶν ἁ(γ)ίων}
\nf{Témoins}{Export EDCS de 2020 : \emph{texte grec absent}. EDB (EDB11272) : texte concordant aux conventions de notation près (indice de similarité 0,967).}
\nf{Apparat}{\{τι\} : lettres superflues ; \{ε\} : lettres superflues}
\nf{Traduction}{(grec) Ici repose Alexandros, ayant vécu vingt ans ; il meurt au mois de novembre ; fils d’Ambrosios, du village de Mikra Kômè ; vis avec les saints.}
\nf{Langue}{Le texte est rédigé en grec, en tout ou en partie ; l’export de 2020 ne le transmettait pas. Des lettres gravées sont jugées superflues par l’éditeur (entre accolades). 6 abréviations développées par l’éditeur (dont (ν), (ς), (μ), (χ), (ε), (γ)).}
\nf{Commentaire}{Vocabulaire religieux : saints (communion des saints) — ἁγίων (lu).}
\nf{Bibliographie}{ICUR II 4271b.}
\nf{Statut}{Lecture secondaire (EDCS) concordante avec l’EDB : lecture stabilisée entre deux bases, non collationnée sur l’ICVR ni sur le monument.}
\end{notice}

\begin{notice}{349}{392}{ICUR n.s. II, 6052 · EDCS-34200678 · EDB14892}
\label{not:349}
\nf{Datation}{392 — indirecte}
\nf{Genre}{Épitaphe ; ensembles A, B, C ; 44 lettres transmises.}
\nf{Références}{ICUR II 6052 = ILCV 4381.}
\nf{Provenance}{Rattachement selon l’EDB : \emph{Coem. Commodillae}.}
\nf{Support}{Plaque de marbre. Inscription gravée. Signes décrits par l’EDB : chrisme avec alpha et oméga. Fiche EDB : Carlo Carletti, première version 2003.}
\ntext{\lnum{1}XV Kal(endas) Oct(obres)\\
\lnum{2}die\{s\} Veneris\\
\lnum{3}Luna XV\\
\lnum{4}vix(it) Leoparda\\
\lnum{5}ann(os) d⸢u⸣o(bus) m(enses) XI}
\nf{Témoins}{Export EDCS de 2020 : \emph{XV Kal(endas) Oct(obres) / die\{s\} Veneris / Luna XV / vix(it) Leoparda / ann(os) do(bus) m(enses) XI}. EDB (EDB14892) : texte concordant aux conventions de notation près (indice de similarité 0,973).}
\nf{Apparat}{⸢u⸣ : Y gravé ; \{s\} : lettres superflues}
\nf{Traduction}{Le 15e jour avant les calendes d’octobre, un vendredi, 15e jour de la lune ; Leoparda vécut 2 ans 11 mois.}
\nf{Langue}{L’éditeur signale une graphie corrigée : Y gravé pour u. Des lettres gravées sont jugées superflues par l’éditeur (entre accolades). 6 abréviations développées par l’éditeur (dont Kal(endas), Oct(obres), vix(it), ann(os), o(bus), m(enses)).}
\nf{Chronologie}{Dates du jour : événement indéterminé : XV Kalendas Octobres = 17 septembre 392 (lu). Âge déclaré : 2 a. 11 m. (lu).}
\nf{Bibliographie}{ICUR II 6052 ; ILCV 4381.}
\nf{Statut}{Lecture secondaire (EDCS) concordante avec l’EDB : lecture stabilisée entre deux bases, non collationnée sur l’ICVR ni sur le monument.}
\end{notice}

\begin{notice}{350}{392}{ICUR n.s. III, 6508 · EDCS-34801308 · EDB22066}
\label{not:350}
\nf{Datation}{392 — directe : nom consulaire lu}
\nf{Genre}{Épitaphe ; ensembles A, B, C ; 100 lettres transmises ; 29 \% restituées.}
\nf{Références}{ICUR III 6508.}
\nf{Provenance}{Rattachement selon l’EDB : \emph{Coem. Domitillae pars inferior}. Conservation selon l’EDB : \emph{Coem. Domitillae pars inferior}.}
\nf{Support}{Plaque de marbre. 51 × 87 × 2 cm (h. × l. × ép.). Inscription gravée. Reproductions signalées par l’EDB : photographie PCAS ×2 (droits réservés ; non reproduites). Fiche EDB : Carlo Carletti / Antonio Enrico Felle, première version 2003.}
\ntext{\lnum{1}Post cons(ulatum) Fl(avi) Tatiano(!) et Quin\\
\lnum{2}ti Aureli Symmac(h)i vv(irorum) cc(larissimorum)\\
\lnum{3}Iust[i]nus recessi[t in pace]\\
\lnum{4}qu[i vixi]t anno[s - - -]\\
\lnum{5}de[positu]s V Idus [Ianuarias]}
\nf{Témoins}{EDB (EDB22066) : texte concordant aux conventions de notation près (indice de similarité 0,959).}
\nf{Apparat}{(!) : forme signalée par l’éditeur}
\nf{Traduction}{Après le consulat de Flavius Tatianus et de Quintus Aurelius Symmachus, clarissimes, Iustinus s’en est allé en paix, qui vécut […] ans ; déposé le 5e jour avant les ides de janvier.}
\nf{Langue}{L’éditeur signale au moins une forme anormale (sic). 6 abréviations développées par l’éditeur (dont cons(ulatum), Fl(avi), Tatiano(!), Symmac(h)i, vv(irorum), cc(larissimorum)).}
\nf{Chronologie}{Noms consulaires : Tatianus (lu) et Symmachus (lu). Datation par post-consulat. Dates du jour : déposition : V Idus Ianuarias = 9 janvier 392 (partiellement restitué).}
\nf{Commentaire}{La formule in pace est attestée seulement par restitution.}
\nf{Bibliographie}{ICUR III 6508. Bibliographie complémentaire selon l’EDB : CLRE, 656.}
\nf{Statut}{Lecture secondaire (EDCS) concordante avec l’EDB : lecture stabilisée entre deux bases, non collationnée sur l’ICVR ni sur le monument. L’EDCS signale une photographie (non consultée).}
\end{notice}

\annee{393}

\begin{notice}{351}{393}{EDCS-07000089 · EDB41240}
\label{not:351}
\nf{Datation}{393 — directe : nom consulaire lu}
\nf{Genre}{Épitaphe ; ensembles A, B, C ; 91 lettres transmises.}
\nf{Références}{AE 1988, 111.}
\nf{Provenance}{Rattachement selon l’EDB : \emph{Coem. s.Theclae}. Conservation selon l’EDB : \emph{Coem. s.Theclae}.}
\nf{Support}{Plaque de marbre. 28 × 46 × 1,4 cm (h. × l. × ép.). Inscription gravée. Reproductions signalées par l’EDB : photographie PCAS (droits réservés ; non reproduites). Fiche EDB : Filippo Piazzolla, première version 2015.}
\ntext{\lnum{1}Feliciae bene merenti in pace\\
\lnum{2}q(uae) vixit an(nos) p(lus) m(inus) XXXI depo\\
\lnum{3}sita d(ie) XVI Kal(endas) Sep(tembres) dd(ominis) nn(ostris)\\
\lnum{4}Theodosio III et Eugenio Augg(ustis)\\
\lnum{5}conss(ulibus)}
\nf{Témoins}{EDB (EDB41240) : texte concordant aux conventions de notation près (indice de similarité 0,996).}
\nf{Traduction}{À Felicia, bien méritante, en paix, qui vécut environ 31 ans ; déposée le 16e jour avant les calendes de septembre, sous le consulat de nos seigneurs Théodose pour la 3e fois et Eugène, Augustes.}
\nf{Langue}{11 abréviations développées par l’éditeur (dont q(uae), an(nos), p(lus), m(inus), d(ie), Kal(endas)).}
\nf{Chronologie}{Noms consulaires : Theodosius (lu) et Eugenius (lu). Dates du jour : déposition : die XVI Kalendas Septembres = 17 août 393 (lu). Âge déclaré : env. 31 a. 16 j. (lu).}
\nf{Commentaire}{La formule in pace est conservée.}
\nf{Bibliographie}{AE 1988, 111. Bibliographie complémentaire selon l’EDB : Fasola 1964, 43-44 fig.19 ; AE, 1988, 111.}
\nf{Statut}{Lecture secondaire (EDCS) concordante avec l’EDB : lecture stabilisée entre deux bases, non collationnée sur l’ICVR ni sur le monument.}
\end{notice}

\begin{notice}{352}{393}{ICUR n.s. I, 727 · EDCS-35100467 · EDB28954}
\label{not:352}
\nf{Datation}{393 — directe : nom consulaire lu}
\nf{Genre}{Épitaphe ; ensembles A, B, C ; 236 lettres transmises ; 6 \% restituées.}
\nf{Références}{ICUR I 727 = ILCV 1599 = CLE 1346.}
\nf{Provenance}{Rattachement selon l’EDB : \emph{Incertae originis} (le volume I des ICVR réunit les inscriptions dont la provenance est perdue).}
\nf{Support}{Support non précisé par l’EDB. Fiche EDB : Filippo Piazzolla, première version 2014.}
\ntext{\lnum{1}Hic funus crudele situm primaq(u)e iu⸢v⸣enta\\
\lnum{2}ereptus iu⸢v⸣enis limina mortis adiit\\
\lnum{3}Eut⸢y⸣chius sapiens sapiens pius adq(ue) beni⸢g⸣nus\\
\lnum{4}[qui Ch]ristum credens pr(a)emia lucis (h)abet\\
\lnum{5}[Euty]chius iacet hic decus uni ipse suorum\\
\lnum{6}[dep]ositus quartu(m) Nonas Oct(obres)\\
\lnum{7}[con]s(ul)i(bus) Theodosi(o) Aug(usto) III et Eugen(io)\\
\lnum{8}[- - -]so ann(os) XXVIIII m(enses) III}
\nf{Témoins}{Export EDCS de 2020 : \emph{Hic funus crudele situm primaq(u)e iubenta / ereptus iuenis limina mortis adiit / Eutichius sapiens sapiens pius adq(ue) beninus / [qui Ch]ristum credens pr(a)emia lucis (h)abet / [Euty]chius iacet hic decus uni ipse suorum / [dep]ositus quartu(m) Nonas Oct(obres) / [con]s(ul)i(bus) Theodosi(o) Aug(usto) III et Eugen(io) / [3]so ann(os) XXVIIII m(enses) III}. EDB (EDB28954) : même texte ; les écarts relevés dans la liste de résultats tiennent à la notation (relecture de la fiche détaillée).}
\nf{Apparat}{⸢v⸣ : B gravé ; ⸢v⸣ : B gravé ; ⸢y⸣ : I gravé ; ⸢g⸣ : C gravé}
\nf{Traduction}{(vers) Ici a été placé un cruel trépas : à la première jeunesse, le jeune homme, arraché, franchit le seuil de la mort ; Eutychius, sage, sage, pieux et bienveillant, qui, croyant au Christ, possède les récompenses de la lumière ; Eutychius gît ici, lui-même l’honneur unique des siens ; déposé le 4e jour avant les nones d’octobre, sous le consulat de Théodose Auguste pour la 3e fois et d’Eugène ; […] 29 ans 3 mois.}
\nf{Langue}{L’éditeur signale 4 graphies corrigées : B gravé pour v (bêtacisme) ×2 ; I gravé pour y (graphie de y) ; C gravé pour g (c pour g). L’export de 2020 diffère de l’état 2022 (voir apparat). 13 abréviations développées par l’éditeur (dont primaq(u)e, adq(ue), pr(a)emia, (h)abet, quartu(m), Oct(obres)).}
\nf{Chronologie}{Noms consulaires : Theodosius (lu) et Eugenius (lu). Ne sont pas des consuls : Eut<y=I>chius sapiens (défunt). Dates du jour : déposition : quartum Nonas Octobres = 4 octobre 393 (lu).}
\nf{Commentaire}{Vocabulaire religieux : désignation du Christ — Christum (partiellement restitué) ; allusion interprétative (niveau 5) — lucis (lu). \emph{[qui Ch]ristum credens} : les deux premières lettres du nom et le relatif sont restitués ; la restitution est contrainte par la terminaison \emph{-ristum} et par le mètre. La pierre est perdue (EDB28954). Le texte figure parmi les désignations du Christ avec restitution. Commentaire développé : dossier p.~\pageref{dos:352}.}
\nf{Bibliographie}{ICUR I 727 ; ILCV 1599 ; CLE 1346. Bibliographie complémentaire selon l’EDB : Ferrua 1979b, 22, n. 727 ; Ferrua 1986a, 44, n. 727.}
\nf{Statut}{Lecture secondaire (EDCS) concordante avec l’EDB : lecture stabilisée entre deux bases, non collationnée sur l’ICVR ni sur le monument.}
\end{notice}

\begin{notice}{353}{393}{ICUR n.s. I, 1448 · EDCS-30300391 · EDB35388}
\label{not:353}
\nf{Datation}{393 — directe : nom consulaire lu}
\nf{Genre}{Épitaphe ; ensembles A, B, C ; 103 lettres transmises.}
\nf{Références}{ICUR I 1448 = ICUR X 27319 = ILCV 2633 (add).}
\nf{Provenance}{Rattachement selon l’EDB : \emph{Coem. et basilica s.Valentini}. Conservation selon l’EDB : \emph{Musei Vaticani: Lapidario Cristiano}.}
\nf{Support}{Sarcophage de marbre. 47 × 139 × 4 cm (h. × l. × ép.). Inscription gravée. Fiche EDB : Filippo Piazzolla, première version 2015.}
\ntext{\lnum{1}Tigriti bene merit⸢e⸣ in pace\\
\lnum{2}qu(a)e vi⸢x⸣it annos LXXX\\
\lnum{3}men(ses) II deposita VIII Kal(endas)\\
\lnum{4}Ian(uarias) dd(ominis) nn(ostris)\\
\lnum{5}T(h)e⸢o⸣dosio III et\\
\lnum{6}Eu⸢g⸣enio\\
\lnum{7}filius fec⸢i⸣t matri}
\nf{Témoins}{Export EDCS de 2020 : \emph{Tigriti bene merit in pace / qu(a)e viit annos LXXX / men(ses) II deposita VIII Kal(endas) / Ian(uarias) dd(ominis) nn(ostris) / T(h)edosio III et / Euenio / filius fect matri}. EDB (EDB35388) : texte concordant aux conventions de notation près (indice de similarité 0,96).}
\nf{Apparat}{⸢e⸣ : I gravé ; ⸢x⸣ : CS gravé ; ⸢o⸣ : V gravé ; ⸢g⸣ : C gravé ; ⸢i⸣ : E gravé}
\nf{Traduction}{À Tigris, bien méritante, en paix, qui vécut 80 ans 2 mois ; déposée le 8e jour avant les calendes de janvier, sous (le consulat de) nos seigneurs Théodose pour la 3e fois et Eugène ; son fils a fait (ce monument) pour sa mère.}
\nf{Langue}{L’éditeur signale 5 graphies corrigées : I gravé pour e (confusion e/i) ; CS gravé pour x ; V gravé pour o (confusion o/u) ; C gravé pour g (c pour g) ; E gravé pour i (confusion e/i).}
\nf{Chronologie}{Noms consulaires : Theodosius (lu) et Eugenius (lu). Dates du jour : déposition : VIII Kalendas Ianuarias = 25 décembre 393 (lu). Âge déclaré : 80 a. 2 m. (lu).}
\nf{Commentaire}{La formule in pace est conservée.}
\nf{Bibliographie}{ICUR I 1448 ; ICUR X 27319 ; ILCV 2633 (add). Bibliographie complémentaire selon l’EDB : ICVR X, 27319 ; Ferrua 1979b, 38, n. 1448.}
\nf{Statut}{Lecture secondaire (EDCS) concordante avec l’EDB : lecture stabilisée entre deux bases, non collationnée sur l’ICVR ni sur le monument. L’EDCS signale une photographie (non consultée).}
\end{notice}

\begin{notice}{354}{393}{ICUR n.s. I, 1940 · EDCS-39200384 · EDB31858}
\label{not:354}
\nf{Datation}{393 — directe : nom consulaire lu}
\nf{Genre}{Épitaphe ; ensembles A, B, C ; 56 lettres transmises ; 29 \% restituées.}
\nf{Références}{ICUR I 1940 (p 490).}
\nf{Provenance}{Rattachement selon l’EDB : \emph{Incertae originis} (le volume I des ICVR réunit les inscriptions dont la provenance est perdue). Conservation selon l’EDB : \emph{Museo Nazionale Romano}.}
\nf{Support}{Plaque de marbre. 41 × 70 × 3,6 cm (h. × l. × ép.). Inscription gravée. Fiche EDB : Filippo Piazzolla, première version 2015.}
\ntext{\lnum{1}[- - -]e qui vixit annis XXVIII m[ensibus - - -]\\
\lnum{2}[- - - Theo]dosio Au⸢g⸣(usto) et Eugenio Au[g(usto)}
\nf{Témoins}{Export EDCS de 2020 : \emph{[3]e qui vixit annis XXVIII m[ensibus 3] / [3 Theo]dosio Au(usto) et Eugenio Au[g(usto)}. EDB (EDB31858 ; écarts à examiner) : \emph{[---]e qui vixit annis XXVIII m[ensibus ---] / [--- Theo]dosio Au⌜g⌝(usto) III et Eugenio Au[g(usto) ---] [texte tronqué dans la liste de résultats]}.}
\nf{Apparat}{⸢g⸣ : C gravé}
\nf{Traduction}{[…] qui vécut 28 ans […] mois ; (sous le consulat de) Théodose Auguste et Eugène Auguste.}
\nf{Langue}{L’éditeur signale une graphie corrigée : C gravé pour g (c pour g).}
\nf{Chronologie}{Noms consulaires : Theodosius (lu en partie) et Eugenius (lu). Âge déclaré : 28 a. (partiellement restitué).}
\nf{Bibliographie}{ICUR I 1940 (p 490).}
\nf{Statut}{Lecture secondaire (EDCS) ; écarts avec l’EDB à arbitrer sur l’ICVR.}
\end{notice}

\begin{notice}{355}{393}{ICUR n.s. III, 8159 · EDCS-35700188 · EDB24681}
\label{not:355}
\nf{Datation}{393 — directe : nom consulaire lu}
\nf{Genre}{Épitaphe ; ensembles A, B, C ; 92 lettres transmises ; 26 \% restituées.}
\nf{Références}{ICUR III 8159 = ILCV +4642.}
\nf{Provenance}{Rattachement selon l’EDB : \emph{Basilica ss. Nerei et Achillei in coem. Domitillae}. Conservation selon l’EDB : \emph{Basilica ss. Nerei et Achillei in coem. Domitillae}.}
\nf{Support}{Plaque de marbre. 47 × 81 × 3,2 cm (h. × l. × ép.). Inscription gravée. Reproductions signalées par l’EDB : photographie PCAS (droits réservés ; non reproduites). Fiche EDB : Carlo Carletti / Antonio Enrico Felle, première version 2003.}
\ntext{\lnum{1}d]ulcis ca[ra - - -]\\
\lnum{2}[- - -]entia q(uae) vix(it) [ann(os) - - -]\\
\lnum{3}XLIIII d(e)p(osita) II[I] Id(us) Oct(obres) co(n)ss(ulatu) Fl(avi)\\
\lnum{4}Theodosi [Aug(usti)] III et Fl(avi) Eugeni\\
\lnum{5}Aug(usti) [in pac]e\\
\lnum{6}b(i)scan[dens(?)]}
\nf{Témoins}{EDB (EDB24681) : texte concordant aux conventions de notation près (indice de similarité 0,949).}
\nf{Traduction}{Douce, chère […]entia, qui vécut 44 ans ; déposée le 3e jour avant les ides d’octobre, sous le consulat de Flavius Théodose Auguste pour la 3e fois et de Flavius Eugène Auguste, en paix ; [biscandens : sens obscur].}
\nf{Langue}{15 abréviations développées par l’éditeur (dont q(uae), vix(it), ann(os), d(e)p, (osita), Id(us)).}
\nf{Chronologie}{Noms consulaires : Theodosius (lu) et Eugenius (lu). Dates du jour : déposition : III Idus Octobres = 13 octobre 393 (partiellement restitué).}
\nf{Commentaire}{La formule in pace est attestée seulement par restitution.}
\nf{Bibliographie}{ICUR III 8159 ; ILCV +4642. Bibliographie complémentaire selon l’EDB : IC I Supplementum, 1846 ; ILCV, 4642 adn.}
\nf{Statut}{Lecture secondaire (EDCS) concordante avec l’EDB : lecture stabilisée entre deux bases, non collationnée sur l’ICVR ni sur le monument. L’EDCS signale une photographie (non consultée).}
\end{notice}

\begin{notice}{356}{393}{ICUR n.s. III, 8160 · EDCS-35700189 · EDB24682}
\label{not:356}
\nf{Datation}{393 — directe : lecture partielle d’un nom consulaire}
\nf{Genre}{Épitaphe ; ensembles A, B, C ; 20 lettres transmises ; 30 \% restituées.}
\nf{Références}{ICUR III 8160.}
\nf{Provenance}{Rattachement selon l’EDB : \emph{Basilica ss. Nerei et Achillei in coem. Domitillae}. Conservation selon l’EDB : \emph{Basilica ss. Nerei et Achillei in coem. Domitillae}.}
\nf{Support}{Plaque de marbre. 18 × 17 × 2,7 cm (h. × l. × ép.). Inscription gravée. Fiche EDB : Carlo Carletti / Antonio Enrico Felle, première version 2003.}
\ntext{\lnum{1}] et Euge[nio - - -]\\
\lnum{2}[- - - qua]e vixit [- - -]\\
\lnum{3}[- - -]LI[}
\nf{Témoins}{EDB (EDB24682) : texte concordant aux conventions de notation près (indice de similarité 0,929).}
\nf{Traduction}{[…] et Eugène […] ; […] qui vécut […] ; […] 51 (?).}
\nf{Chronologie}{Noms consulaires : Eugenius (lu en partie).}
\nf{Bibliographie}{ICUR III 8160.}
\nf{Statut}{Lecture secondaire (EDCS) concordante avec l’EDB : lecture stabilisée entre deux bases, non collationnée sur l’ICVR ni sur le monument. L’EDCS signale une photographie (non consultée).}
\end{notice}

\begin{notice}{357}{393}{ICUR n.s. III, 8429 · EDCS-43700747 · EDB25703}
\label{not:357}
\nf{Datation}{393 — directe : nom consulaire lu}
\nf{Genre}{Épitaphe ; ensembles A, B, C ; 86 lettres transmises ; 22 \% restituées.}
\nf{Références}{ICUR III 8429 = ILCV 4556.}
\nf{Provenance}{Rattachement selon l’EDB : \emph{Incertum locum in coem. Domitillae}. Conservation selon l’EDB : \emph{coem. Domitillae, aula supra coemeterium}.}
\nf{Support}{Plaque de marbre. Inscription gravée. Signes décrits par l’EDB : croix monogrammatique. Reproductions signalées par l’EDB : photographie PCAS ×3 (droits réservés ; non reproduites). Fiche EDB : Carlo Carletti / Antonio Enrico Felle, première version 2003.}
\ntext{\lnum{1}[- - -]ssima mi\{c\}hi Victorino\\
\lnum{2}[qui se vivo fecit s]⸢b⸣i et ipse\{i\} deposita III Kal(endas) Iunias\\
\lnum{3}[post(?)] conss(ulatum) Theodosi et Eugeni}
\nf{Témoins}{Export EDCS de 2020 : \emph{[3]ssima mi\{c\}hi Victorino / [qui se vivo fecit s]i et ipse\{i\} deposita III Kal(endas) Iunias / [post(?)] conss(ulatum) Theodosi et Eugeni}. EDB (EDB25703 ; écarts à examiner) : \emph{((monogramma Christi in crucem translatum cum alpha et omega)) / [--- dulci]ssima michi Victorino / [--- si]vi et ipsei deposita III kal(endas) iunias / [---] ((palma)) cons(ulatu) ((palma)) Theodosi et Eugeni}.}
\nf{Apparat}{⸢b⸣ : V gravé ; \{c\} : lettres superflues ; \{i\} : lettres superflues}
\nf{Traduction}{[…] très […] à moi, Victorinus, qui de son vivant a fait (ce monument) pour lui-même et elle ; déposée le 3e jour avant les calendes de juin, après le consulat de Théodose et d’Eugène.}
\nf{Langue}{L’éditeur signale une graphie corrigée : V gravé pour b (bêtacisme). Des lettres gravées sont jugées superflues par l’éditeur (entre accolades). 3 abréviations développées par l’éditeur (dont Kal(endas), post(?), conss(ulatum)).}
\nf{Chronologie}{Noms consulaires : Theodosius (lu) et Eugenius (lu). Datation par post-consulat. Dates du jour : déposition : III Kalendas Iunias = 30 mai 393 (lu).}
\nf{Bibliographie}{ICUR III 8429 ; ILCV 4556. Bibliographie complémentaire selon l’EDB : ILCV, 4556 ; CLRE, 657.}
\nf{Statut}{Lecture secondaire (EDCS) ; écarts avec l’EDB à arbitrer sur l’ICVR. L’EDCS signale une photographie (non consultée).}
\end{notice}

\begin{notice}{358}{393}{ICUR n.s. III, 8431 · EDCS-35802018 · EDB24939}
\label{not:358}
\nf{Datation}{393 — directe : nom consulaire lu}
\nf{Genre}{Épitaphe ; ensembles A, B, C ; 48 lettres transmises ; 65 \% restituées.}
\nf{Références}{ICUR III 8431.}
\nf{Provenance}{Rattachement selon l’EDB : \emph{Incertum locum in coem. Domitillae}. Conservation selon l’EDB : \emph{Musei Vaticani: Lapidario Cristiano}.}
\nf{Support}{Plaque de marbre. Inscription gravée. Fiche EDB : Carlo Carletti / Antonio Enrico Felle, première version 2003.}
\ntext{\lnum{1}]ri(a)e [- - -]\\
\lnum{2}[- - - con]sulibus [d(ominis) n(ostris) Theodosio]\\
\lnum{3}[III] et Eugen[io]}
\nf{Témoins}{EDB (EDB24939) : texte concordant aux conventions de notation près (indice de similarité 0,978).}
\nf{Traduction}{[…] sous les consuls nos seigneurs Théodose pour la 3e fois et Eugène.}
\nf{Langue}{3 abréviations développées par l’éditeur (dont ri(a)e, d(ominis), n(ostris)).}
\nf{Chronologie}{Noms consulaires : Theodosius (restitué) et Eugenius (lu).}
\nf{Bibliographie}{ICUR III 8431.}
\nf{Statut}{Lecture secondaire (EDCS) concordante avec l’EDB : lecture stabilisée entre deux bases, non collationnée sur l’ICVR ni sur le monument. Plus de la moitié des lettres sont restituées. L’EDCS signale une photographie (non consultée).}
\end{notice}

\begin{notice}{359}{393}{ICUR n.s. IV, 11123 · EDCS-37200406 · EDB39585}
\label{not:359}
\nf{Datation}{393 — directe : nom consulaire lu}
\nf{Genre}{Épitaphe ; ensembles A, B, C ; 89 lettres transmises ; 21 \% restituées.}
\nf{Références}{ICUR IV 11123 = ILCV 4414.}
\nf{Provenance}{Rattachement selon l’EDB : \emph{Coem. Callisti area subdialis et incerti loci}. Conservation selon l’EDB : \emph{Coem. Callisti area subdialis et incerti loci}.}
\nf{Support}{Plaque de marbre. 32 × 45 cm (h. × l.). Inscription gravée. Reproductions signalées par l’EDB : photographie PCAS ×2 (droits réservés ; non reproduites). Fiche EDB : Marida Pierno, première version 2016.}
\ntext{\lnum{1}]\\
\lnum{2}qui v[ix(it) ann(os) - - - e]\\
\lnum{3}t mens⸢e⸣s s[eptem - - -]\\
\lnum{4}sextu(m) Idus Dec⸢e⸣m[br(es) - - -]\\
\lnum{5}in pac⸢e⸣ consu⸢l⸣i(bu)s\\
\lnum{6}T(h)e⸢o⸣dosio ter et Eu⸢g⸣e\\
\lnum{7}nio consu⸢l⸣ibus}
\nf{Témoins}{Export EDCS de 2020 : \emph{] / qui v[ix(it) ann(os) 3 e]/t menss s[eptem 3] / sextu(m) Idus Decm[br(es) 3] / in pac consui(bu)s / T(h)edosio ter et Eue/nio consuibus}. EDB (EDB39585 ; écarts à examiner) : \emph{q̲u̲i̲ v̲[ixit --- e]/t̲ ṃensis s[---] / s̲e̲xtu idus ḍecim̲(bres) / i̲n̲ paci consu⌜l⌝is / Teudosio ter et Euce/nio consu⌜l⌝ibus}.}
\nf{Apparat}{⸢e⸣ : I gravé ; ⸢e⸣ : I gravé ; ⸢e⸣ : I gravé ; ⸢l⸣ : I gravé ; ⸢o⸣ : V gravé ; ⸢g⸣ : C gravé ; ⸢l⸣ : I gravé}
\nf{Traduction}{[…] qui vécut […] ans et sept mois […] ; le 6e jour avant les ides de décembre […] ; en paix ; sous les consuls Théodose pour la 3e fois et Eugène.}
\nf{Langue}{L’éditeur signale 7 graphies corrigées : I gravé pour e (confusion e/i) ×3 ; I gravé pour l ×2 ; V gravé pour o (confusion o/u) ; C gravé pour g (c pour g).}
\nf{Chronologie}{Noms consulaires : Theodosius (lu) et Eugenius (lu). Dates du jour : événement indéterminé : sextum Idus Decembres = 8 décembre 393 (partiellement restitué). Âge déclaré : 7 m. (partiellement restitué).}
\nf{Commentaire}{La formule in pace est conservée.}
\nf{Bibliographie}{ICUR IV 11123 ; ILCV 4414.}
\nf{Statut}{Lecture secondaire (EDCS) ; écarts avec l’EDB à arbitrer sur l’ICVR. L’EDCS signale une photographie (non consultée).}
\end{notice}

\begin{notice}{360}{393}{ICUR n.s. IV, 11124a · EDCS-37200407 · EDB746}
\label{not:360}
\nf{Datation}{393 — directe : lecture partielle d’un nom consulaire}
\nf{Genre}{Épitaphe ; ensembles A, B, C ; 15 lettres transmises ; 53 \% restituées.}
\nf{Références}{ICUR IV 11124a.}
\nf{Provenance}{Rattachement selon l’EDB : \emph{Coem. Callisti area subdialis et incerti loci}. Conservation selon l’EDB : \emph{Coem. Callisti area subdialis et incerti loci}.}
\nf{Support}{Plaque de marbre. Inscription gravée. Fiche EDB : Antonio Enrico Felle, première version 2015.}
\ntext{\lnum{1}]\\
\lnum{2}[- - - ἔ]τη ιʹ Θεο[δοσίῳ τό - - -]\\
\lnum{3}[}
\nf{Témoins}{Export EDCS de 2020 : \emph{texte grec absent}. EDB (EDB746 ; écarts à examiner) : \emph{[---] / [--- ἔ]τη ι´ Θεο[δοσίῳ τό ---] / [---]}.}
\nf{Traduction}{(grec) […] 10 ans ; (sous le consulat de) Théodose pour la […] fois […].}
\nf{Langue}{Le texte est rédigé en grec, en tout ou en partie ; l’export de 2020 ne le transmettait pas.}
\nf{Chronologie}{Noms consulaires : Theodosius (lu en partie).}
\nf{Bibliographie}{ICUR IV 11124a.}
\nf{Statut}{Lecture secondaire (EDCS) ; écarts avec l’EDB à arbitrer sur l’ICVR. Plus de la moitié des lettres sont restituées.}
\end{notice}

\begin{notice}{361}{393}{ICUR n.s. IV, 12422 · EDCS-37600946 · EDB41699}
\label{not:361}
\nf{Datation}{393 — directe : nom consulaire lu}
\nf{Genre}{Épitaphe ; ensembles A, B, C ; 92 lettres transmises ; 48 \% restituées.}
\nf{Références}{ICUR IV 12422.}
\nf{Provenance}{Rattachement selon l’EDB : \emph{inter vias Appiam et Ardeatinam}.}
\nf{Support}{Support non précisé par l’EDB. Fiche EDB : Fillippo Trotta, première version 2017.}
\ntext{\lnum{1}Victoriae ben[e merenti]\\
\lnum{2}in pace quae v[ixit annos]\\
\lnum{3}plus minus XIV [Theodosio]\\
\lnum{4}A(u)g(usto) III et Eugeni[o Aug(usto) cons(ulibus)]}
\nf{Témoins}{EDB (EDB41699) : texte concordant aux conventions de notation près (indice de similarité 0,975).}
\nf{Traduction}{À Victoria, bien méritante, en paix, qui vécut environ 14 ans ; sous le consulat de Théodose Auguste pour la 3e fois et d’Eugène Auguste.}
\nf{Langue}{4 abréviations développées par l’éditeur (dont A(u)g, (usto), Aug(usto), cons(ulibus)).}
\nf{Chronologie}{Noms consulaires : Theodosius (restitué) et Eugenius (lu). Âge déclaré : env. 14 a. (partiellement restitué).}
\nf{Commentaire}{La formule in pace est conservée.}
\nf{Bibliographie}{ICUR IV 12422.}
\nf{Statut}{Lecture secondaire (EDCS) concordante avec l’EDB : lecture stabilisée entre deux bases, non collationnée sur l’ICVR ni sur le monument.}
\end{notice}

\annee{394}

\begin{notice}{362}{394}{EDCS-13000396 · EDB43058}
\label{not:362}
\nf{Datation}{394 — directe : lecture partielle d’un nom consulaire}
\nf{Genre}{Épitaphe ; ensembles A, B, C ; 51 lettres transmises ; 59 \% restituées.}
\nf{Références}{CEACelio 311.}
\nf{Provenance}{Rattachement selon l’EDB : \emph{Incertae originis} (le volume I des ICVR réunit les inscriptions dont la provenance est perdue). Conservation selon l’EDB : \emph{Antiquarium del Celio}.}
\nf{Support}{La fiche EDB ne figure pas dans l’état 2019 des données EpiDoc ; support non documenté.}
\ntext{\lnum{1}deposit]us in pac[e - - -]\\
\lnum{2}[cons(ulatu) Nicomaci Fl]abiani(!) v(iri) [c(larissimi)]}
\nf{Témoins}{EDB (EDB43058 ; écarts à examiner) : \emph{------? / [---]N̲A̲S̲I̲A̲N̲[---] / [---? deposit]us in pac̣[e ---?] / [cons(ulatu) Nicomaci Fl]ạbiani v̅(iri) [c̅(larissimi)]}.}
\nf{Apparat}{(!) : forme signalée par l’éditeur}
\nf{Traduction}{Déposé en paix […] sous le consulat de Nicomachus Flavianus, clarissime.}
\nf{Langue}{L’éditeur signale au moins une forme anormale (sic). 4 abréviations développées par l’éditeur (dont cons(ulatu), Flabiani(!), v(iri), c(larissimi)).}
\nf{Chronologie}{Noms consulaires : Nicomachus Flavianus (absent).}
\nf{Commentaire}{La formule in pace est attestée seulement par restitution.}
\nf{Bibliographie}{CEACelio 311. Bibliographie complémentaire selon l’EDB : Gregori 2001, 324, n. 311, tav. LXXV, 2 (S. Lorenzi).}
\nf{Statut}{Lecture secondaire (EDCS) ; écarts avec l’EDB à arbitrer sur l’ICVR. Plus de la moitié des lettres sont restituées.}
\end{notice}

\begin{notice}{363}{394}{ICUR n.s. II, 4487 · EDCS-34100550 · EDB31779}
\label{not:363}
\nf{Datation}{394 — directe : lecture partielle d’un nom consulaire}
\nf{Genre}{Épitaphe ; ensembles A, B, C ; 71 lettres transmises ; 51 \% restituées.}
\nf{Références}{ICUR II 4487.}
\nf{Provenance}{Rattachement selon l’EDB : \emph{Coem. Calepodii}. Conservation selon l’EDB : \emph{Coem. Calepodii}.}
\nf{Support}{Plaque de marbre. 32 × 36 cm (h. × l.). Inscription gravée. Fiche EDB : Carlo Carletti, première version 2017.}
\ntext{\lnum{1}]us qui vixit ann[os - - -]\\
\lnum{2}[deposi]tus VII Idus Oct[obres - - -]\\
\lnum{3}[Arcadi]o III et Hon[norio II co(n)s(ulibus)]}
\nf{Témoins}{EDB (EDB31779) : texte concordant aux conventions de notation près (indice de similarité 0,936).}
\nf{Traduction}{[…] qui vécut […] ans ; déposé le 7e jour avant les ides d’octobre […] ; sous le consulat d’Arcadius pour la 3e fois et d’Honorius pour la 2e fois.}
\nf{Langue}{2 abréviations développées par l’éditeur (dont co(n)s, (ulibus)).}
\nf{Chronologie}{Noms consulaires : Arcadius (restitué). Dates du jour : déposition : VII Idus Octobres = 9 octobre 394 (partiellement restitué).}
\nf{Bibliographie}{ICUR II 4487. Bibliographie complémentaire selon l’EDB : Mazzoleni et alii 1999, 605.}
\nf{Statut}{Lecture secondaire (EDCS) concordante avec l’EDB : lecture stabilisée entre deux bases, non collationnée sur l’ICVR ni sur le monument. Plus de la moitié des lettres sont restituées.}
\end{notice}

\begin{notice}{364}{394}{ICUR n.s. II, 4503 · EDCS-34100566 · EDB20919}
\label{not:364}
\nf{Datation}{394 — directe : nom consulaire lu}
\nf{Genre}{Épitaphe ; ensembles A, B, C ; 137 lettres transmises ; 40 \% restituées.}
\nf{Références}{ICUR II 4503.}
\nf{Provenance}{Rattachement selon l’EDB : \emph{Coem. Pontiani}.}
\nf{Support}{Plaque de marbre. Inscription gravée. Fiche EDB : Antonio Enrico Felle, première version 2003.}
\ntext{\lnum{1}]dus\\
\lnum{2}[- - - qui vixit anno]s vit(a)e suae\\
\lnum{3}[- - -]ta et quinque et fec⸢i⸣t\\
\lnum{4}[cum coniuge s]ua Donata dulcissima\\
\lnum{5}[annos quinqu]e depositus in pa[c]e\\
\lnum{6}[- - -]as Nicomac(h)o [Fla]via\\
\lnum{7}[no v(iro) c(larissimo) c]onsule}
\nf{Témoins}{Export EDCS de 2020 : \emph{]dus / [3 qui vixit anno]s vit(a)e suae / [3]ta et quinque et fect / [cum coniuge s]ua Donata dulcissima / [annos quinqu]e depositus in pa[c]e / [3]as Nicomac(h)o [Fla]via/[no v(iro) c(larissimo) c]onsule}. EDB (EDB20919) : texte concordant aux conventions de notation près (indice de similarité 0,978).}
\nf{Apparat}{⸢i⸣ : E gravé}
\nf{Traduction}{[…] qui vécut, des années de sa vie, […]-cinq ans, et a vécu avec son épouse Donata, très douce, cinq ans ; déposé en paix […] sous le consulat de Nicomachus Flavianus, clarissime.}
\nf{Langue}{L’éditeur signale une graphie corrigée : E gravé pour i (confusion e/i).}
\nf{Chronologie}{Noms consulaires : Nicomachus Flavianus (lu).}
\nf{Commentaire}{La formule in pace est attestée seulement par restitution.}
\nf{Bibliographie}{ICUR II 4503.}
\nf{Statut}{Lecture secondaire (EDCS) concordante avec l’EDB : lecture stabilisée entre deux bases, non collationnée sur l’ICVR ni sur le monument.}
\end{notice}

\begin{notice}{365}{394}{ICUR n.s. II, 6460 · EDCS-34300812 · EDB20818}
\label{not:365}
\nf{Datation}{394 — directe : nom consulaire lu}
\nf{Genre}{Épitaphe ; ensembles A, B, C ; 147 lettres transmises.}
\nf{Références}{ICUR II 6460 = ILCV 3822.}
\nf{Provenance}{Rattachement selon l’EDB : \emph{via Ostiensis}. Conservation selon l’EDB : \emph{Musei Vaticani: Lapidario Cristiano}.}
\nf{Support}{Sarcophage de marbre. 60 × 45 cm (h. × l.). Inscription gravée. Reproductions signalées par l’EDB : reproduction, Perret 1851 (droits réservés ; non reproduites). Fiche EDB : Maria Cristina Grisanzio, première version 2003.}
\ntext{\lnum{1}Dep(ositus) III Idus Mai(as) Iosimus qui\\
\lnum{2}vixit ann⸢o⸣s XXVIII qui fecit\\
\lnum{3}cum co⸢m⸣pare sua ann⸢o⸣s septe(m)\\
\lnum{4}mens⸢e⸣s VIIII bene merenti in pace con\\
\lnum{5}sulatu Nicomac(h)i Fla⸢v⸣iani locum ar\\
\lnum{6}marari(i) quadrisomum}
\nf{Témoins}{Export EDCS de 2020 : \emph{Dep(ositus) III Idus Mai(as) Iosimus qui / vixit anns XXVIII qui fecit / cum copare sua anns septe(m) / menss VIIII bene merenti in pace con/sulatu Nicomac(h)i Flaiani locum ar/marari(i) quadrisomum}. EDB (EDB20818) : texte concordant aux conventions de notation près (indice de similarité 0,944).}
\nf{Apparat}{⸢o⸣ : V gravé ; ⸢m⸣ : N gravé ; ⸢o⸣ : V gravé ; ⸢e⸣ : I gravé ; ⸢v⸣ : B gravé}
\nf{Traduction}{Déposé le 3e jour avant les ides de mai, Iosimus (Zosimus), qui vécut 28 ans et a vécu avec sa compagne sept ans 9 mois ; bien méritant, en paix ; sous le consulat de Nicomachus Flavianus ; emplacement quadruple de l’armarius (?).}
\nf{Langue}{L’éditeur signale 5 graphies corrigées : V gravé pour o (confusion o/u) ×2 ; N gravé pour m (n pour m) ; I gravé pour e (confusion e/i) ; B gravé pour v (bêtacisme).}
\nf{Chronologie}{Noms consulaires : Nicomachus Flavianus (lu). Dates du jour : déposition : III Idus Maias = 13 mai 394 (lu). Âge déclaré : 28 a. (lu).}
\nf{Commentaire}{La formule in pace est conservée. Fonction ou métier mentionné : armararii.}
\nf{Bibliographie}{ICUR II 6460 ; ILCV 3822.}
\nf{Statut}{Lecture secondaire (EDCS) concordante avec l’EDB : lecture stabilisée entre deux bases, non collationnée sur l’ICVR ni sur le monument. L’EDCS signale une photographie (non consultée).}
\end{notice}

\begin{notice}{366}{394}{ICUR n.s. III, 8648 · EDCS-35900847 · EDB19494}
\label{not:366}
\nf{Datation}{394 — directe : nom consulaire lu}
\nf{Genre}{Épitaphe ; ensembles A, B, C ; 86 lettres transmises.}
\nf{Références}{ICUR III 8648.}
\nf{Provenance}{Rattachement selon l’EDB : \emph{Coem. Commodillae}.}
\nf{Support}{Plaque de marbre. Inscription gravée. Fiche EDB : Carlo Carletti, première version 2003.}
\ntext{\lnum{1}Quintiliano innocenti qui ⸢v⸣ixit\\
\lnum{2}ann(os) V mens⸢e⸣s VII depositus III Idus Au⸢g⸣(ustas)\\
\lnum{3}Fla⸢v⸣io Nicomaco Fla⸢v⸣iano cons(ule)}
\nf{Témoins}{Export EDCS de 2020 : \emph{Quintiliano innocenti qui ixit / ann(os) V menss VII depositus III Idus Au(ustas) / Flaio Nicomaco Flaiano cons(ule)}. EDB (EDB19494 ; écarts à examiner) : \emph{Quintiliano innocenti qui bixit / ann(os) V mensis VII depositus III idus aug(ustas) / Flabio Nicomaco Flabiano cons(ule) [texte tronqué dans la liste de résultats]}.}
\nf{Apparat}{⸢v⸣ : B gravé ; ⸢e⸣ : I gravé ; ⸢g⸣ : C gravé ; ⸢v⸣ : B gravé ; ⸢v⸣ : B gravé}
\nf{Traduction}{À Quintilianus, innocent, qui vécut 5 ans 7 mois ; déposé le 3e jour avant les ides d’août, sous le consulat de Flavius Nicomachus Flavianus.}
\nf{Langue}{L’éditeur signale 5 graphies corrigées : B gravé pour v (bêtacisme) ×3 ; I gravé pour e (confusion e/i) ; C gravé pour g (c pour g).}
\nf{Chronologie}{Noms consulaires : Nicomachus Flavianus (lu). Dates du jour : déposition : III Idus Augustas = 11 août 394 (lu). Âge déclaré : 5 a. 7 m. (lu).}
\nf{Bibliographie}{ICUR III 8648.}
\nf{Statut}{Lecture secondaire (EDCS) ; écarts avec l’EDB à arbitrer sur l’ICVR.}
\end{notice}

\begin{notice}{367}{394}{ICUR n.s. IV, 11125 · EDCS-37200411 · EDB39589}
\label{not:367}
\nf{Datation}{394 — directe : nom consulaire lu}
\nf{Genre}{Épitaphe ; ensembles A, B, C ; 109 lettres transmises ; 36 \% restituées.}
\nf{Références}{ICUR IV 11125 = ILCV +2971a.}
\nf{Provenance}{Rattachement selon l’EDB : \emph{Coem. Callisti area subdialis et incerti loci}. Conservation selon l’EDB : \emph{Coem. Callisti area subdialis et incerti loci}.}
\nf{Support}{Plaque de marbre. 18 × 27 × 5 cm (h. × l. × ép.). Inscription gravée. Reproductions signalées par l’EDB : photographie PCAS ×2 (droits réservés ; non reproduites). Fiche EDB : Marida Pierno, première version 2015.}
\ntext{\lnum{1}Alexandr(a)e bene [meren]\\
\lnum{2}ti qu(a)e vixit ann(os) I[- - -]\\
\lnum{3}d⸢e⸣posita in pace [- - -]\\
\lnum{4}Idus Dec(embres) conss(ulibus) [dd(ominis) nn(ostris)]\\
\lnum{5}Arc\{h\}adio Aug(usto) III [et Hono]\\
\lnum{6}rio I[I(?) conss(ulibus)]}
\nf{Témoins}{Export EDCS de 2020 : \emph{Alexandr(a)e bene [meren]/ti qu(a)e vixit ann(os) I[3] / dposita in pace [3] / Idus Dec(embres) conss(ulibus) [dd(ominis) nn(ostris)] / Arc\{h\}adio Aug(usto) III [et Hono]/rio I[I(?) conss(ulibus)]}. EDB (EDB39589) : texte concordant aux conventions de notation près (indice de similarité 0,922).}
\nf{Apparat}{⸢e⸣ : I gravé ; \{h\} : lettres superflues}
\nf{Traduction}{À Alexandra, bien méritante, qui vécut […] ans ; déposée en paix […] ides de décembre, sous le consulat de nos seigneurs Arcadius Auguste pour la 3e fois et Honorius pour la 2e fois.}
\nf{Langue}{L’éditeur signale une graphie corrigée : I gravé pour e (confusion e/i). Des lettres gravées sont jugées superflues par l’éditeur (entre accolades). 10 abréviations développées par l’éditeur (dont Alexandr(a)e, qu(a)e, ann(os), Dec(embres), conss(ulibus), dd(ominis)).}
\nf{Chronologie}{Noms consulaires : Arcadius (lu) et Honorius (lu en partie). Dates du jour : déposition : Idus Decembres = 13 décembre 394 (lu).}
\nf{Commentaire}{La formule in pace est conservée.}
\nf{Bibliographie}{ICUR IV 11125 ; ILCV +2971a.}
\nf{Statut}{Lecture secondaire (EDCS) concordante avec l’EDB : lecture stabilisée entre deux bases, non collationnée sur l’ICVR ni sur le monument. L’EDCS signale une photographie (non consultée).}
\end{notice}

\annee{395}

\begin{notice}{368}{395}{EDCS-23702281 · EDB41941}
\label{not:368}
\nf{Datation}{395 — directe : nom consulaire lu}
\nf{Genre}{Épitaphe ; ensembles A, B, C ; 119 lettres transmises.}
\nf{Références}{AE 2001, 526.}
\nf{Provenance}{Rattachement selon l’EDB : \emph{ex-vigna Chiaraviglio}. Conservation selon l’EDB : \emph{ex-vigna Chiaraviglio}.}
\nf{Support}{La fiche EDB ne figure pas dans l’état 2019 des données EpiDoc ; support non documenté.}
\ntext{\lnum{1}]\\
\lnum{2}fratres dulcissim\\
\lnum{3}o mages picto locum\\
\lnum{4}geminum fecerunt ubi nos in pace prae\\
\lnum{5}cess⸢i⸣t qui vix(it) an(n)o(s) XXX dep(ositus) IIII Non(as) Oct(obres) Anicis\\
\lnum{6}Olybr⸢i⸣o et Probino con(sulibus)}
\nf{Témoins}{Export EDCS de 2020 : \emph{] / fratres dulcissim/o mages picto locum / geminum fecerunt ubi nos in pace prae/cesst qui vix(it) an(n)o(s) XXX dep(ositus) IIII Non(as) Oct(obres) Anicis / Olybro et Probino con(sulibus)}. EDB (EDB41941) : texte concordant aux conventions de notation près (indice de similarité 0,986).}
\nf{Apparat}{⸢i⸣ : E gravé ; ⸢i⸣ : Y gravé}
\nf{Traduction}{[…] les frères au très doux Magès, peintre (?), ont fait un emplacement double, où il nous a précédés dans la paix ; il vécut 30 ans ; déposé le 4e jour avant les nones d’octobre, sous le consulat des Anicii Olybrius et Probinus.}
\nf{Langue}{L’éditeur signale 2 graphies corrigées : E gravé pour i (confusion e/i) ; Y gravé pour i.}
\nf{Chronologie}{Noms consulaires : Probinus (lu) et Olybrius (lu). Dates du jour : déposition : IIII Nonas Octobres = 4 octobre 395 (lu). Âge déclaré : 30 a. (lu).}
\nf{Commentaire}{La formule in pace est conservée.}
\nf{Bibliographie}{AE 2001, 526. Bibliographie complémentaire selon l’EDB : Giuliani 2001, 175-176, A.a. 1 ; AE, 2001, 526.}
\nf{Statut}{Lecture secondaire (EDCS) concordante avec l’EDB : lecture stabilisée entre deux bases, non collationnée sur l’ICVR ni sur le monument.}
\end{notice}

\begin{notice}{369}{395}{EDCS-44100241 · EDB42064}
\label{not:369}
\nf{Datation}{395 — directe : nom consulaire lu}
\nf{Genre}{Épitaphe ; ensembles A, B, C ; 37 lettres transmises.}
\nf{Références}{HEp 2006, 00435 = AE 2006, 240.}
\nf{Provenance}{Rattachement selon l’EDB : \emph{Incertae originis} (le volume I des ICVR réunit les inscriptions dont la provenance est perdue). Conservation selon l’EDB : \emph{Valencia   ; Prehistoric Museum of Valencia}.}
\nf{Support}{La fiche EDB ne figure pas dans l’état 2019 des données EpiDoc ; support non documenté.}
\ntext{\lnum{1}III Nonas Aug(ustas) An\{n\}iciis vv(iris) cc(larissimis)\\
\lnum{2}Bincenti(a)e in pace}
\nf{Témoins}{EDB (EDB42064) : texte concordant aux conventions de notation près (indice de similarité 0,938). L'EDB (EDB42064) distingue deux textes (a, la date consulaire des \emph{Anicii} ; b, \emph{Bincentie in pace}), suivis d'une croix monogrammatique avec alpha et oméga. Rattachement : origine incertaine ; conservation actuelle : Valence (Espagne), Museo de Prehistoria, selon l'EDB. Ce lieu de conservation n'a pas été contrôlé.}
\nf{Apparat}{\{n\} : lettres superflues}
\nf{Traduction}{Le 3e jour avant les nones d’août, sous le consulat des Anicii, clarissimes ; à Vincentia, en paix.}
\nf{Langue}{Des lettres gravées sont jugées superflues par l’éditeur (entre accolades). 4 abréviations développées par l’éditeur (dont Aug(ustas), vv(iris), cc(larissimis), Bincenti(a)e).}
\nf{Chronologie}{Noms consulaires : Anicii (lu). Dates du jour : événement indéterminé : III Nonas Augustas = 3 août 395 (lu).}
\nf{Commentaire}{La formule in pace est conservée. L’EDB transcrit les signes suivants : croix monogrammatique.}
\nf{Bibliographie}{HEp 2006, 00435 ; AE 2006, 240. Bibliographie complémentaire selon l’EDB : Mazzoleni 2006b, 429-430 ; AE, 2006, 240.}
\nf{Statut}{Lecture secondaire (EDCS) concordante avec l’EDB : lecture stabilisée entre deux bases, non collationnée sur l’ICVR ni sur le monument. L’EDCS signale une photographie (non consultée).}
\end{notice}

\begin{notice}{370}{395}{ICUR n.s. I, 1212 · EDCS-38700751 · EDB30159}
\label{not:370}
\nf{Datation}{395 — directe : nom consulaire lu}
\nf{Genre}{Épitaphe ; ensembles A, B, C ; 65 lettres transmises ; 58 \% restituées.}
\nf{Références}{ICUR I 1212.}
\nf{Provenance}{Rattachement selon l’EDB : \emph{Incertae originis} (le volume I des ICVR réunit les inscriptions dont la provenance est perdue). Conservation selon l’EDB : \emph{Antiquarium del Celio}.}
\nf{Support}{Plaque de marbre. 34 × 44 cm (h. × l.). Inscription gravée. Fiche EDB : Filippo Piazzolla, première version 2014.}
\ntext{\lnum{1}]\\
\lnum{2}Aur[- - -]\\
\lnum{3}vixi[t - - -]\\
\lnum{4}ann(os) [- - -]\\
\lnum{5}dep(ositus ?) VII[- - -]\\
\lnum{6}Olybri[o et Probino vv(iris) cc(larissimis) - - -]\\
\lnum{7}quies[cit in pace}
\nf{Témoins}{EDB (EDB30159) : texte concordant aux conventions de notation près (indice de similarité 0,95).}
\nf{Traduction}{[…] Aur[…] vécut […] ans ; déposé le 7e […] ; sous le consulat d’Olybrius et de Probinus, clarissimes […] ; il repose en paix.}
\nf{Langue}{4 abréviations développées par l’éditeur (dont ann(os), dep(ositus ?), vv(iris), cc(larissimis)).}
\nf{Chronologie}{Noms consulaires : Probinus (restitué) et Olybrius (lu).}
\nf{Commentaire}{La formule in pace est attestée seulement par restitution. Vocabulaire religieux : repos et sommeil de la mort — quiescit (partiellement restitué).}
\nf{Bibliographie}{ICUR I 1212. Bibliographie complémentaire selon l’EDB : Ferrua 1979b, 29, n. 1212.}
\nf{Statut}{Lecture secondaire (EDCS) concordante avec l’EDB : lecture stabilisée entre deux bases, non collationnée sur l’ICVR ni sur le monument. Plus de la moitié des lettres sont restituées.}
\end{notice}

\begin{notice}{371}{395}{ICUR n.s. I, 2105 · EDCS-39300324 · EDB32099}
\label{not:371}
\nf{Datation}{395 — directe : nom consulaire lu}
\nf{Genre}{Épitaphe ; ensembles A, B, C ; 84 lettres transmises ; 2 \% restituées.}
\nf{Références}{ICUR I 2105 = ILCV 3099a.}
\nf{Provenance}{Rattachement selon l’EDB : \emph{Incertae originis} (le volume I des ICVR réunit les inscriptions dont la provenance est perdue). Conservation selon l’EDB : \emph{Musei Vaticani: Galleria Lapidaria}.}
\nf{Support}{Plaque de marbre. 30 × 78 cm (h. × l.). Inscription gravée. Fiche EDB : Filippo Piazzolla, première version 2015.}
\ntext{\lnum{1}Anici(i)s Hermoginia[no]\\
\lnum{2}O⸢ly⸣b⸢ri⸣o et Probino v(iris) c(larissimis)\\
\lnum{3}XVII Kal(endas) Octobres d(eposita)\\
\lnum{4}Galla an(n)o⸢r⸣(um) II nata e⸢t⸣ m(ensium)\\
\lnum{5}IIII qu(i)es(cit) in pa\{a\}ce}
\nf{Témoins}{Export EDCS de 2020 : \emph{Anici(i)s Hermoginia[no] / Obo et Probino v(iris) c(larissimis) / XVII Kal(endas) Octobres d(eposita) / Galla an(n)o(um) II nata e m(ensium) / IIII qu(i)es(cit) in pa\{a\}ce}. EDB (EDB32099) : texte concordant aux conventions de notation près (indice de similarité 0,972).}
\nf{Apparat}{⸢ly⸣ : IV gravé ; ⸢ri⸣ : K gravé ; ⸢r⸣ : B gravé ; ⸢t⸣ : I gravé ; \{a\} : lettres superflues}
\nf{Traduction}{Sous le consulat des Anicii Hermogenianus Olybrius et Probinus, clarissimes, le 17e jour avant les calendes d’octobre, déposée Galla, âgée de 2 ans et 4 mois ; elle repose en paix.}
\nf{Langue}{L’éditeur signale 4 graphies corrigées : IV gravé pour ly ; K gravé pour ri ; B gravé pour r ; I gravé pour t. Des lettres gravées sont jugées superflues par l’éditeur (entre accolades). 10 abréviations développées par l’éditeur (dont Anici(i)s, v(iris), c(larissimis), Kal(endas), d(eposita), an(n)o).}
\nf{Chronologie}{Noms consulaires : Probinus (lu), Olybrius (lu) et Anicii (lu). Dates du jour : décès : XVII Kalendas Octobres = 15 septembre 395 (lu). Âge déclaré : 2 a. (lu).}
\nf{Commentaire}{La formule in pace est conservée. Vocabulaire religieux : repos et sommeil de la mort — quiescit (lu).}
\nf{Bibliographie}{ICUR I 2105 ; ILCV 3099a. Bibliographie complémentaire selon l’EDB : Lega 2000 ; ILCV, 3099A.}
\nf{Statut}{Lecture secondaire (EDCS) concordante avec l’EDB : lecture stabilisée entre deux bases, non collationnée sur l’ICVR ni sur le monument. L’EDCS signale une photographie (non consultée).}
\end{notice}

\begin{notice}{372}{395}{ICUR n.s. I, 3213 · EDCS-40300349 · EDB33161}
\label{not:372}
\nf{Datation}{395 — directe : nom consulaire lu}
\nf{Genre}{Épitaphe ; ensembles A, B, C ; 74 lettres transmises.}
\nf{Références}{ICUR I 3213.}
\nf{Provenance}{Rattachement selon l’EDB : \emph{Incertae originis} (le volume I des ICVR réunit les inscriptions dont la provenance est perdue).}
\nf{Support}{Support non précisé par l’EDB. Fiche EDB : Filippo Piazzolla, première version 2015.}
\ntext{\lnum{1}Gaudentius qui vixit annos XLX\\
\lnum{2}depositus VIII Kal(endas) Maias Aniciis Olybrio et\\
\lnum{3}Probino cons(ulibus)}
\nf{Témoins}{EDB (EDB33161) : texte concordant aux conventions de notation près (indice de similarité 0,989).}
\nf{Traduction}{Gaudentius, qui vécut 40 (+10 ?) ans [chiffre XLX anomal] ; déposé le 8e jour avant les calendes de mai, sous le consulat des Anicii Olybrius et Probinus.}
\nf{Langue}{2 abréviations développées par l’éditeur (dont Kal(endas), cons(ulibus)).}
\nf{Chronologie}{Noms consulaires : Probinus (lu), Olybrius (lu) et Anicii (lu). Dates du jour : déposition : VIII Kalendas Maias = 24 avril 395 (lu). Âge déclaré : 50 a. (lu).}
\nf{Bibliographie}{ICUR I 3213. Bibliographie complémentaire selon l’EDB : ILCV, 2946 adn.}
\nf{Statut}{Lecture secondaire (EDCS) concordante avec l’EDB : lecture stabilisée entre deux bases, non collationnée sur l’ICVR ni sur le monument.}
\end{notice}

\begin{notice}{373}{395}{ICUR n.s. II, 4510 · EDCS-34100573 · EDB18940}
\label{not:373}
\nf{Datation}{395 — directe : nom consulaire lu}
\nf{Genre}{Épitaphe ; ensembles A, B, C ; 155 lettres transmises ; 56 \% restituées.}
\nf{Références}{ICUR II 4510.}
\nf{Provenance}{Rattachement selon l’EDB : \emph{Coem. Pontiani}.}
\nf{Support}{Plaque de marbre. Inscription gravée. Reproductions signalées par l’EDB : photographie PCAS (droits réservés ; non reproduites). Fiche EDB : Antonio Enrico Felle, première version 2003.}
\ntext{\lnum{1}Aniciis Olybrio [et Probino vv(iris) cc(larissimis) conss(ulibus)]\\
\lnum{2}Victor ed D[- - -]\\
\lnum{3}Donata [quae vixit annos - - -]\\
\lnum{4}dies V et fecit cum [marito suo annos - - -]\\
\lnum{5}dies III deposita [- - -]\\
\lnum{6}conss(ulibus) Arca[dio - - - et Honorio - - - Augg(ustis)]}
\nf{Témoins}{EDB (EDB18940) : texte concordant aux conventions de notation près (indice de similarité 0,93).}
\nf{Traduction}{Sous le consulat des Anicii Olybrius et Probinus, clarissimes, Victor et D[…] ; Donata, qui vécut […] ans […] 5 jours, et a vécu avec son mari […] ans […] 3 jours ; déposée […] sous le consulat d’Arcadius […] et d’Honorius […], Augustes.}
\nf{Langue}{5 abréviations développées par l’éditeur (dont vv(iris), cc(larissimis), conss(ulibus), Augg(ustis)).}
\nf{Chronologie}{Noms consulaires : Arcadius (lu en partie), Honorius (restitué), Probinus (restitué), Olybrius (lu) et Anicii (lu). Âge déclaré : 5 j. (partiellement restitué).}
\nf{Bibliographie}{ICUR II 4510. Bibliographie complémentaire selon l’EDB : Ferrua 1988, 12.}
\nf{Statut}{Lecture secondaire (EDCS) concordante avec l’EDB : lecture stabilisée entre deux bases, non collationnée sur l’ICVR ni sur le monument. Plus de la moitié des lettres sont restituées. L’EDCS signale une photographie (non consultée).}
\end{notice}

\begin{notice}{374}{395}{ICUR n.s. II, 4829 · EDCS-36101309}
\label{not:374}
\nf{Datation}{395 — directe : nom consulaire lu}
\nf{Genre}{Épitaphe ; ensembles A, B, C ; 117 lettres transmises ; 14 \% restituées.}
\nf{Références}{ICUR II 4829 = ILCV 3800a.}
\nf{Provenance}{Rome (indication de l’EDCS) ; aucun rattachement cimétérial transmis.}
\nf{Support}{Non documenté dans les sources consultées.}
\ntext{\lnum{1}Petronius in pace XVII Kalen[das - - -]\\
\lnum{2}nis qui vixit ann⸢o⸣s LXV consulat[u Olybrio et]\\
\lnum{3}Probino vv(iris) cc(larissimis) hic requiescit in p[ace - - -]\\
\lnum{4}[s]e ⸢v⸣i⸢v⸣a fecit bisomus vacat}
\nf{Témoins}{Export EDCS de 2020 : \emph{Petronius in pace XVII Kalen[das 3]/nis qui vixit anns LXV consulat[u Olybrio et] / Probino vv(iris) cc(larissimis) hic requiescit in p[ace 3] / [s]e ia fecit bisomus vacat}. Aucune fiche EDB appariée (numéro ICUR absent de la liste publique de l’EDB).}
\nf{Apparat}{⸢o⸣ : V gravé ; ⸢v⸣ : B gravé ; ⸢v⸣ : B gravé}
\nf{Traduction}{Petronius, en paix, le 17e jour avant les calendes […], qui vécut 65 ans, sous le consulat d’Olybrius et de Probinus, clarissimes ; il repose ici en paix ; […] de son vivant a fait une tombe double ; [vacat].}
\nf{Langue}{L’éditeur signale 3 graphies corrigées : V gravé pour o (confusion o/u) ; B gravé pour v (bêtacisme) ×2.}
\nf{Chronologie}{Noms consulaires : Probinus (lu) et Olybrius (restitué). Âge déclaré : 65 a. (partiellement restitué).}
\nf{Commentaire}{La formule in pace est conservée. Vocabulaire religieux : repos et sommeil de la mort — requiescit (lu).}
\nf{Bibliographie}{ICUR II 4829 ; ILCV 3800a.}
\nf{Statut}{Lecture secondaire (EDCS) sans second témoin numérique.}
\end{notice}

\begin{notice}{375}{395}{ICUR n.s. II, 5861 · EDCS-38302567 · EDB22029}
\label{not:375}
\nf{Datation}{395 — directe : lecture partielle d’un nom consulaire}
\nf{Genre}{Épitaphe ; ensembles A, B, C ; 86 lettres transmises ; 65 \% restituées.}
\nf{Références}{ICUR II 5861.}
\nf{Provenance}{Rattachement selon l’EDB : \emph{Area Lucinae}. Conservation selon l’EDB : \emph{s. Paolo f.l.m.}.}
\nf{Support}{Plaque de marbre. 33 × 33 cm (h. × l.). Inscription gravée. Fiche EDB : Anita Rocco, première version 2007.}
\ntext{\lnum{1}depos]itus VI Kal[endas - - -]\\
\lnum{2}[- - - qui vixit an]n⸢o⸣s XIII me[nses - - -]\\
\lnum{3}[- - - Olybri]o et Prob[ino vv(iris) cc(larissimis) cons(ulibus)]}
\nf{Témoins}{Export EDCS de 2020 : \emph{depos]itus VI Kal[endas 3] / [3 qui vixit an]ns XIII me[nses 3] / [3 Olybri]o et Prob[ino vv(iris) cc(larissimis) cons(ulibus)]}. EDB (EDB22029) : texte concordant aux conventions de notation près (indice de similarité 0,933).}
\nf{Apparat}{⸢o⸣ : V gravé}
\nf{Traduction}{Déposé le 6e jour avant les calendes […] ; […] qui vécut 13 ans […] mois ; […] sous le consulat d’Olybrius et de Probinus, clarissimes.}
\nf{Langue}{L’éditeur signale une graphie corrigée : V gravé pour o (confusion o/u).}
\nf{Chronologie}{Noms consulaires : Probinus (lu en partie) et Olybrius (restitué). Âge déclaré : 13 a. (partiellement restitué).}
\nf{Bibliographie}{ICUR II 5861.}
\nf{Statut}{Lecture secondaire (EDCS) concordante avec l’EDB : lecture stabilisée entre deux bases, non collationnée sur l’ICVR ni sur le monument. Plus de la moitié des lettres sont restituées. L’EDCS signale une photographie (non consultée).}
\end{notice}

\begin{notice}{376}{395}{ICUR n.s. II, 6053 · EDCS-34200679 · EDB14893}
\label{not:376}
\nf{Datation}{395 — directe : nom consulaire lu}
\nf{Genre}{Épitaphe ; ensembles A, B, C ; 91 lettres transmises.}
\nf{Références}{ICUR II 6053 = ILCV 2606 (em) = AE 1906, 82.}
\nf{Provenance}{Rattachement selon l’EDB : \emph{Coem. Commodillae}.}
\nf{Support}{Plaque de marbre. Inscription gravée. Signes décrits par l’EDB : chrisme, chrisme avec alpha et oméga, croix monogrammatique et croix monogrammatique avec alpha et oméga. Reproductions signalées par l’EDB : photographie PCAS ×3 (droits réservés ; non reproduites). Fiche EDB : Carlo Carletti, première version 2015.}
\ntext{\lnum{1}Vitalin(a)e bene meren\\
\lnum{2}ti in pace quae vixit ann⸢o⸣s\\
\lnum{3}p(lus) m(inus) XXV deposita XI Kal(endas) Feb(ruarias)\\
\lnum{4}Olybrio et Probino vv(iris) cc(larissimis)\\
\lnum{5}conss(ulibus) in pace}
\nf{Témoins}{Export EDCS de 2020 : \emph{Vitalin(a)e bene meren/ti in pace quae vixit anns / p(lus) m(inus) XXV deposita XI Kal(endas) Feb(ruarias) / Olybrio et Probino vv(iris) cc(larissimis) / conss(ulibus) in pace}. EDB (EDB14893) : texte concordant aux conventions de notation près (indice de similarité 0,953).}
\nf{Apparat}{⸢o⸣ : V gravé}
\nf{Traduction}{À Vitalina, bien méritante, en paix, qui vécut environ 25 ans ; déposée le 11e jour avant les calendes de février, sous le consulat d’Olybrius et de Probinus, clarissimes, en paix.}
\nf{Langue}{L’éditeur signale une graphie corrigée : V gravé pour o (confusion o/u).}
\nf{Chronologie}{Noms consulaires : Probinus (lu) et Olybrius (lu). Dates du jour : déposition : XI Kalendas Februarias = 22 janvier 395 (lu). Âge déclaré : env. 25 a. (lu).}
\nf{Commentaire}{La formule in pace est conservée.}
\nf{Bibliographie}{ICUR II 6053 ; ILCV 2606 (em) ; AE 1906, 82.}
\nf{Statut}{Lecture secondaire (EDCS) concordante avec l’EDB : lecture stabilisée entre deux bases, non collationnée sur l’ICVR ni sur le monument. L’EDCS signale une photographie (non consultée).}
\end{notice}

\begin{notice}{377}{395}{ICUR n.s. III, 6509 · EDCS-34801309 · EDB22067}
\label{not:377}
\nf{Datation}{395 — directe : nom consulaire lu}
\nf{Genre}{Épitaphe ; ensembles A, B, C ; 23 lettres transmises ; 13 \% restituées.}
\nf{Références}{ICUR III 6509.}
\nf{Provenance}{Rattachement selon l’EDB : \emph{Coem. Domitillae pars inferior}. Conservation selon l’EDB : \emph{Coem. Domitillae pars inferior}.}
\nf{Support}{Plaque de marbre. 35 × 69 cm (h. × l.). Inscription gravée. Fiche EDB : Carlo Carletti / Antonio Enrico Felle, première version 2014.}
\ntext{\lnum{1}[- - -] F]laviis Anic[iis]\\
\lnum{2}vv(iris) cc(larissimis) conss(ulibus)}
\nf{Témoins}{EDB (EDB22067 ; écarts à examiner) : \emph{------ / [--- F]laviis Anic[iis] / vv(iris) cc(larissimis) cons(ulibu)s [texte tronqué dans la liste de résultats]}.}
\nf{Traduction}{[…] sous le consulat des Flavii Anicii, clarissimes.}
\nf{Langue}{3 abréviations développées par l’éditeur (dont vv(iris), cc(larissimis), conss(ulibus)).}
\nf{Chronologie}{Noms consulaires : Anicii (lu).}
\nf{Bibliographie}{ICUR III 6509.}
\nf{Statut}{Lecture secondaire (EDCS) ; écarts avec l’EDB à arbitrer sur l’ICVR. L’EDCS signale une photographie (non consultée).}
\end{notice}

\begin{notice}{378}{395}{ICUR n.s. III, 8161 · EDCS-35700190 · EDB25348}
\label{not:378}
\nf{Datation}{395 — directe : lecture partielle d’un nom consulaire}
\nf{Genre}{Épitaphe ; ensembles A, B, C ; 318 lettres transmises ; 26 \% restituées.}
\nf{Références}{ICUR III 8161 = ICUR III 8345 = ILCV 994 (em) = Mossong-2022, 00028.}
\nf{Provenance}{Rattachement selon l’EDB : \emph{Basilica ss. Nerei et Achillei in coem. Domitillae}. Conservation selon l’EDB : \emph{Basilica ss. Nerei et Achillei in coem. Domitillae}.}
\nf{Support}{Plaque de marbre. Inscription gravée. Reproductions signalées par l’EDB : photographie PCAS (droits réservés ; non reproduites). Fiche EDB : Carlo Carletti / Antonio Enrico Felle, première version 2003.}
\ntext{\lnum{1}Limenius v]ixit ann[os - - -]\\
\lnum{2}[- - -] Ade[od]a[- - -]\\
\lnum{3}[- - -]des\\
\lnum{4}[- - - prud]ens(?) ha[be]tur\\
\lnum{5}[- - - d]ulcis primo mihi\\
\lnum{6}[- - -] quiescit\\
\lnum{7}[- - -] mens(i)um coniu[gi - - -]\\
\lnum{8}[- - -]m de[p(ositus)] XI Kal(endas) Octob[res]\\
\lnum{9}[- - - Olybrio et Pro]bino\\
\lnum{10}[filia eius Vict]oria d(e)p(osita) Idi(bu)s Ian(uariis)\\
\lnum{11}[cum viro suo - - -]orio filiae ita est\\
\lnum{12}[genu]s Lim[enio p]atre [pr]esbytero\\
\lnum{13}nepo[s] Crescen[tii n]epos [- - -]\\
\lnum{14}diac(oni) reg(ionis) V [ne]pos [- - -]\\
\lnum{15}[p]ronep(os) Marcelli episc(opi) [ne]pos Cresce[ntii]\\
\lnum{16}[m]atrona presb(yteri) abnepos\\
\lnum{17}Basili socrus epis(copi) episcop(i)\\
\lnum{18}Ursini Albanens(is)}
\nf{Témoins}{EDB (EDB25348 ; divergence forte) : \emph{------ / [---]des / [--- prud?]ens ha[be]tur / [--- d]ulcis primo mịhi / [---] quiescit / [---] m̂ensum coniu[gi] / ------}.}
\nf{Traduction}{Limenius vécut […] ans […] Adeodatus (?) […] ; […] il est tenu pour sage […] ; […] doux d’abord pour moi […] ; […] il repose […] ; […] de mois, à son épouse […] ; déposé le 11e jour avant les calendes d’octobre […] sous (le consulat d’)Olybrius et Probinus ; sa fille Victoria, déposée aux ides de janvier avec son mari […]orius ; telle est la lignée de la fille : Limenius, son père, prêtre ; petit-fils de Crescentius […] ; […] diacre de la Ve région ; petit-fils de […] ; arrière-petit-fils de Marcellus, évêque ; petit-fils de Crescentius ; épouse d’un prêtre ; arrière-arrière-petit-fils de Basilius ; belle-mère de l’évêque […] ; Ursinus d’Albano (?) [généalogie à structure incertaine].}
\nf{Langue}{16 abréviations développées par l’éditeur (dont prudens(?), mens(i)um, dep(ositus), Kal(endas), d(e)p, (osita)).}
\nf{Chronologie}{Noms consulaires : Probinus (lu en partie) et Olybrius (restitué). Ne sont pas des consuls : Limenius v]ixit ann[os (défunt). La fiche EDB d'ICVR III 8345 (EDB25348) ne contient que le début du poème ; elle est à joindre à ICVR III 8274c (EDB25442) et 8387v (EDB42793). Aucune de ces fiches ne contient la formule \emph{[Olybrio et Pro]bino} ni la généalogie cléricale transmises par l'EDCS. L'EDB date l'ensemble « 390–424 ». L'année 395 dépend d'une partie du texte non retrouvée dans la seconde base : elle est conservée comme année de classement, sous réserve du contrôle de l'ICVR III 8345. Dates du jour : déposition : XI Kalendas Octobres = 21 septembre 395 (partiellement restitué) ; déposition : Idibus Ianuariis = 13 janvier 395 (lu).}
\nf{Commentaire}{Vocabulaire religieux : repos et sommeil de la mort — quiescit (lu). Fonction ou métier mentionné : presbytero, diaconi, episcopi, presbyteri. Commentaire développé : dossier p.~\pageref{dos:378}.}
\nf{Bibliographie}{ICUR III 8161 ; ICUR III 8345 ; ILCV 994 (em) ; Mossong-2022, 00028. Bibliographie complémentaire selon l’EDB : ICVR III, 8274c ; ICVR III, 8387v.}
\nf{Statut}{Lecture secondaire (EDCS) ; divergence forte avec l’EDB : le texte ne doit pas être exploité sans contrôle de l’édition imprimée. L’EDCS signale une photographie (non consultée).}
\end{notice}

\begin{notice}{379}{395}{ICUR n.s. III, 8162 · EDCS-35700191 · EDB24684}
\label{not:379}
\nf{Datation}{395 — directe : nom consulaire lu}
\nf{Genre}{Épitaphe ; ensembles A, B, C ; 52 lettres transmises ; 62 \% restituées.}
\nf{Références}{ICUR III 8162.}
\nf{Provenance}{Rattachement selon l’EDB : \emph{Basilica ss. Nerei et Achillei in coem. Domitillae}. Conservation selon l’EDB : \emph{Basilica ss. Nerei et Achillei in coem. Domitillae}.}
\nf{Support}{Plaque de marbre. 26 × 19 × 2 cm (h. × l. × ép.). Inscription gravée. Reproductions signalées par l’EDB : photographie PCAS ×2 (droits réservés ; non reproduites). Fiche EDB : Carlo Carletti / Antonio Enrico Felle, première version 2003.}
\ntext{\lnum{1}Locus [- - -]\\
\lnum{2}ann⸢o⸣s [- - -]\\
\lnum{3}Kalen[das - - - Aniciis]\\
\lnum{4}Olybr[io et Probino conss(ulibus)}
\nf{Témoins}{Export EDCS de 2020 : \emph{Locus [3] / anns [3] / Kalen[das 3 Aniciis] / Olybr[io et Probino conss(ulibus)}. EDB (EDB24684 ; écarts à examiner) : \emph{loc̲u̲s̲ [---] / annus [---] / kalen[d---] / Olybr[io et Probino cons(ulibu)s]}.}
\nf{Apparat}{⸢o⸣ : V gravé}
\nf{Traduction}{Emplacement de […] ans […] calendes […] ; sous le consulat des Anicii Olybrius et Probinus.}
\nf{Langue}{L’éditeur signale une graphie corrigée : V gravé pour o (confusion o/u).}
\nf{Chronologie}{Noms consulaires : Probinus (restitué), Olybrius (lu) et Anicii (lu).}
\nf{Bibliographie}{ICUR III 8162.}
\nf{Statut}{Lecture secondaire (EDCS) ; écarts avec l’EDB à arbitrer sur l’ICVR. Plus de la moitié des lettres sont restituées. L’EDCS signale une photographie (non consultée).}
\end{notice}

\begin{notice}{380}{395}{ICUR n.s. III, 8163 · EDCS-35700192 · EDB24685}
\label{not:380}
\nf{Datation}{395 — directe : nom consulaire lu}
\nf{Genre}{Épitaphe ; ensembles A, B, C ; 20 lettres transmises.}
\nf{Références}{ICUR III 8163.}
\nf{Provenance}{Rattachement selon l’EDB : \emph{Basilica ss. Nerei et Achillei in coem. Domitillae}. Conservation selon l’EDB : \emph{Basilica ss. Nerei et Achillei in coem. Domitillae}.}
\nf{Support}{Plaque de marbre. 26 × 34 × 5 cm (h. × l. × ép.). Inscription gravée. Fiche EDB : Carlo Carletti / Antonio Enrico Felle, première version 2003.}
\ntext{\lnum{1}Olybrio] et Pro⸢b⸣ino cons(ulibus)}
\nf{Témoins}{Export EDCS de 2020 : \emph{Olybrio] et Proino cons(ulibus)}. EDB (EDB24685 ; écarts à examiner) : \emph{------ / [--- Olybrio] et Provino cons(ulibu)[s]}.}
\nf{Apparat}{⸢b⸣ : V gravé}
\nf{Traduction}{Sous le consulat d’Olybrius et de Probinus.}
\nf{Langue}{L’éditeur signale une graphie corrigée : V gravé pour b (bêtacisme).}
\nf{Chronologie}{Noms consulaires : Probinus (lu) et Olybrius (restitué).}
\nf{Bibliographie}{ICUR III 8163.}
\nf{Statut}{Lecture secondaire (EDCS) ; écarts avec l’EDB à arbitrer sur l’ICVR. L’EDCS signale une photographie (non consultée).}
\end{notice}

\begin{notice}{381}{395}{ICUR n.s. III, 8164 · EDCS-35700193 · EDB24686}
\label{not:381}
\nf{Datation}{395 — directe : nom consulaire lu}
\nf{Genre}{Épitaphe ; ensembles A, B, C ; 132 lettres transmises.}
\nf{Références}{ICUR III 8164 = ILCV 2146 (em).}
\nf{Provenance}{Rattachement selon l’EDB : \emph{Basilica ss. Nerei et Achillei in coem. Domitillae}. Conservation selon l’EDB : \emph{Basilica ss. Nerei et Achillei in coem. Domitillae}.}
\nf{Support}{Plaque de marbre. 61 × 62 cm (h. × l.). Inscription gravée. Reproductions signalées par l’EDB : photographie PCAS (droits réservés ; non reproduites). Fiche EDB : Carlo Carletti / Antonio Enrico Felle, première version 2003.}
\ntext{\lnum{1}Beatus d⸢e⸣functus\\
\lnum{2}est III Idus Maias\\
\lnum{3}die \{s\} Saturni \{s\} an(norum) XXVIII\\
\lnum{4}Anicio Olybrio et\\
\lnum{5}Probino vv(iris) cc(larissimis) conss(ulibus)\\
\lnum{6}Vicentia d⸢e⸣functa\\
\lnum{7}est XII Kal(endas) Iunias\\
\lnum{8}die \{e\} Lun⸢ae⸣ ann(orum) XXVII\\
\lnum{9}in pace}
\nf{Témoins}{Export EDCS de 2020 : \emph{Beatus dfunctus / est III Idus Maias / die \{s\} Saturni \{s\} an(norum) XXVIII / Anicio Olybrio et / Probino vv(iris) cc(larissimis) conss(ulibus) / Vicentia dfuncta / est XII Kal(endas) Iunias / die \{e\} Lun ann(orum) XXVII / in pace}. EDB (EDB24686) : texte concordant aux conventions de notation près (indice de similarité 0,963).}
\nf{Apparat}{⸢e⸣ : I gravé ; ⸢e⸣ : I gravé ; ⸢ae⸣ : IS gravé ; \{s\} : lettres superflues ; \{s\} : lettres superflues ; \{e\} : lettres superflues}
\nf{Traduction}{Beatus est mort le 3e jour avant les ides de mai, un samedi, âgé de 28 ans, sous le consulat d’Anicius Olybrius et de Probinus, clarissimes ; Vincentia est morte le 12e jour avant les calendes de juin, un lundi, âgée de 27 ans ; en paix.}
\nf{Langue}{L’éditeur signale 3 graphies corrigées : I gravé pour e (confusion e/i) ×2 ; IS gravé pour ae. Des lettres gravées sont jugées superflues par l’éditeur (entre accolades). 6 abréviations développées par l’éditeur (dont an(norum), vv(iris), cc(larissimis), conss(ulibus), Kal(endas), ann(orum)).}
\nf{Chronologie}{Noms consulaires : Probinus (lu), Olybrius (lu) et Anicii (lu). Dates du jour : décès : III Idus Maias = 13 mai 395 (lu) ; décès : XII Kalendas Iunias = 21 mai 395 (lu). Âges déclarés : 28 a. (lu) ; 27 a. (lu).}
\nf{Commentaire}{La formule in pace est conservée.}
\nf{Bibliographie}{ICUR III 8164 ; ILCV 2146 (em).}
\nf{Statut}{Lecture secondaire (EDCS) concordante avec l’EDB : lecture stabilisée entre deux bases, non collationnée sur l’ICVR ni sur le monument. L’EDCS signale une photographie (non consultée).}
\end{notice}

\begin{notice}{382}{395}{ICUR n.s. III, 8729 · EDCS-31900337 · EDB20584}
\label{not:382}
\nf{Datation}{395 — directe : nom consulaire lu}
\nf{Genre}{Épitaphe ; ensembles A, B, C ; 308 lettres transmises ; 18 \% restituées.}
\nf{Références}{ICUR III 8729 = CLE 676 = ILCV 3424.}
\nf{Provenance}{Rattachement selon l’EDB : \emph{Tituli qui in coem. Callisti reperti traduntur}. Conservation selon l’EDB : \emph{Musei Vaticani: Galleria Lapidaria}.}
\nf{Support}{Plaque de marbre. Inscription gravée. Texte métrique. Fiche EDB : Carlo Carletti, première version 2003.}
\ntext{\lnum{1}[Pal]ma celestia regna t[enet - - -]\\
\lnum{2}[at] vero totius castitati[s - - -]\\
\lnum{3}nuper tante fides in te[- - -]\\
\lnum{4}tante pudicitiae ment[- - -]\\
\lnum{5}hanc vero grata anim[am - - -]\\
\lnum{6}acer⸢b⸣us natus quanta qu(a)e [- - -]\\
\lnum{7}generosa mi\{c\}hi co(n)iux et in[- - -]\\
\lnum{8}aurea sunt c(a)elestis qu(a)e tu c(a)elo [- - -]\\
\lnum{9}trium subolum mater in sin[u dei pausat]\\
\lnum{10}eius autem nomen capita ver[suum declarant]\\
\lnum{11}vixit autem ann(os) XXV [- - - dep(osita)]\\
\lnum{12}XVII Kal(endas) Sep(tembres) Olybrio et Probi[no cons(ulibus)]\\
\lnum{13}[Pa]nthagate}
\nf{Témoins}{Export EDCS de 2020 : \emph{[Pal]ma celestia regna t[enet 3] / [at] vero totius castitati[s 3] / nuper tante fides in te[3] / tante pudicitiae ment[3] / hanc vero grata anim[am 3] / acerus natus quanta qu(a)e [3] / generosa mi\{c\}hi co(n)iux et in[3] / aurea sunt c(a)elestis qu(a)e tu c(a)elo [3] / trium subolum mater in sin[u dei pausat] / eius autem nomen capita ver[suum declarant] / vixit autem ann(os) XXV [3 dep(osita)] /…}. EDB (EDB20584) : même texte ; les écarts relevés dans la liste de résultats tiennent à la notation (relecture de la fiche détaillée). L'EDB (EDB20584) donne le même texte ; l'EDCS ajoute \emph{// [Pa]nthagate}, absent de l'EDB. Le nom de la défunte, que les initiales des vers devaient former (\emph{eius autem nomen capita ver[suum declarant]}), ne peut être établi sur les seules lettres conservées. L'EDB qualifie le texte de métrique.}
\nf{Apparat}{⸢b⸣ : V gravé ; \{c\} : lettres superflues}
\nf{Traduction}{(vers acrostiches, fragmentaires) Avec la palme, elle tient les royaumes célestes […] ; mais de toute chasteté […] ; naguère d’une si grande foi en toi […] ; d’une si grande pudeur, l’esprit […] ; cette âme, reconnaissante […] ; […] né prématurément, combien […] ; épouse généreuse pour moi et […] ; dorées sont les (demeures ?) célestes que toi, au ciel […] ; mère de trois enfants, elle repose dans le sein de Dieu ; son nom, les initiales des vers le déclarent ; elle vécut 25 ans […] ; déposée le 17e jour avant les calendes de septembre, sous le consulat d’Olybrius et de Probinus // Panthagathè.}
\nf{Langue}{L’éditeur signale une graphie corrigée : V gravé pour b (bêtacisme). Des lettres gravées sont jugées superflues par l’éditeur (entre accolades). 10 abréviations développées par l’éditeur (dont qu(a)e, co(n)iux, c(a)elestis, c(a)elo, ann(os)).}
\nf{Chronologie}{Noms consulaires : Probinus (lu en partie) et Olybrius (lu). Dates du jour : déposition : XVII Kalendas Septembres = 16 août 395 (lu). Âge déclaré : 25 a. (partiellement restitué).}
\nf{Commentaire}{Vocabulaire religieux : désignation de Dieu — dei (entièrement restitué) ; initiation et foi — fides (lu) ; repos et sommeil de la mort — pausat (entièrement restitué) ; allusion interprétative (niveau 5) — caelestis (lu), caelo (lu). Commentaire développé : dossier p.~\pageref{dos:382}.}
\nf{Bibliographie}{ICUR III 8729 ; CLE 676 ; ILCV 3424.}
\nf{Statut}{Lecture secondaire (EDCS) concordante avec l’EDB : lecture stabilisée entre deux bases, non collationnée sur l’ICVR ni sur le monument.}
\end{notice}

\begin{notice}{383}{395}{ICUR n.s. IV, 9377 · EDCS-36100891 · EDB19574}
\label{not:383}
\nf{Datation}{395 — directe : nom consulaire lu}
\nf{Genre}{Épitaphe ; ensembles A, B, C ; 75 lettres transmises.}
\nf{Références}{ICUR IV 9377 = ILCV 3003a (em).}
\nf{Provenance}{Rattachement selon l’EDB : \emph{Crypta s.Cornelii}.}
\nf{Support}{Plaque de marbre. Inscription gravée. Fiche EDB : Antonio Enrico Felle, première version 2003.}
\ntext{\lnum{1}Dep(ositio) Amanti(a)es(!) XI Kal(endas) Martias q(uae) vixit\\
\lnum{2}ann(os) p(lus) m(inus) XLV Aniciis Olybrio et Probino vv(iris) cc(larissimis) conss(ulibus)\\
\lnum{3}in pace}
\nf{Témoins}{EDB (EDB19574) : texte concordant aux conventions de notation près (indice de similarité 0,954).}
\nf{Apparat}{(!) : forme signalée par l’éditeur}
\nf{Traduction}{Déposition d’Amantia le 11e jour avant les calendes de mars, qui vécut environ 45 ans, sous le consulat des Anicii Olybrius et Probinus, clarissimes ; en paix.}
\nf{Langue}{L’éditeur signale au moins une forme anormale (sic). 10 abréviations développées par l’éditeur (dont Dep(ositio), Amanti(a)es, Kal(endas), q(uae), ann(os), p(lus)).}
\nf{Chronologie}{Noms consulaires : Probinus (lu), Olybrius (lu) et Anicii (lu). Ne sont pas des consuls : Dep(ositio) Amanti(a)es(!) (défunte). Dates du jour : déposition : XI Kalendas Martias = 19 février 395 (lu). Âge déclaré : env. 45 a. (lu).}
\nf{Commentaire}{La formule in pace est conservée.}
\nf{Bibliographie}{ICUR IV 9377 ; ILCV 3003a (em).}
\nf{Statut}{Lecture secondaire (EDCS) concordante avec l’EDB : lecture stabilisée entre deux bases, non collationnée sur l’ICVR ni sur le monument.}
\end{notice}

\begin{notice}{384}{395}{ICUR n.s. IV, 11127 · EDCS-37200414 · EDB39591}
\label{not:384}
\nf{Datation}{395 — directe : consulat restitué}
\nf{Genre}{Épitaphe ; ensembles A, B, C ; 37 lettres transmises ; 70 \% restituées.}
\nf{Références}{ICUR IV 11127.}
\nf{Provenance}{Rattachement selon l’EDB : \emph{Coem. Callisti area subdialis et incerti loci}. Conservation selon l’EDB : \emph{Coem. Callisti area subdialis et incerti loci}.}
\nf{Support}{Plaque de marbre. 41 × 30 × 4 cm (h. × l. × ép.). Inscription gravée. Fiche EDB : Marida Pierno, première version 2015.}
\ntext{\lnum{1}fe]cit Do[- - -]\\
\lnum{2}[- - - O]ct(o)b(res) L[ybrio et Probino conss(ulibus)]}
\nf{Témoins}{EDB (EDB39591 ; écarts à examiner) : \emph{[--- fe]cit Do[---] / [o]ct(o)b(r---?) Ol[ybrio et Probino consulibus]}.}
\nf{Traduction}{[…] a fait […] ; […] octobre, sous le consulat d’Olybrius et de Probinus.}
\nf{Langue}{3 abréviations développées par l’éditeur (dont Oct(o)b, (res), conss(ulibus)).}
\nf{Chronologie}{Noms consulaires : Probinus (restitué).}
\nf{Bibliographie}{ICUR IV 11127.}
\nf{Statut}{Lecture secondaire (EDCS) ; écarts avec l’EDB à arbitrer sur l’ICVR. Plus de la moitié des lettres sont restituées.}
\end{notice}

\annee{396}

\begin{notice}{385}{396}{EDCS-75200321}
\label{not:385}
\nf{Datation}{396 — directe : nom consulaire lu}
\nf{Genre}{Épitaphe ; ensembles A, B, C ; 45 lettres transmises ; 24 \% restituées.}
\nf{Références}{AE 2016, 170.}
\nf{Provenance}{Rome (indication de l’EDCS) ; aucun rattachement cimétérial transmis.}
\nf{Support}{Non documenté dans les sources consultées.}
\ntext{\lnum{1}vir]go(?) vixit an(nos)\\
\lnum{2}[- - -] d(e)p(ositus ?) pri(di)e Idus\\
\lnum{3}[- - - A]rc(adio) IIII Hon(orio) III\\
\lnum{4}[co(n)s(ulibus ?)}
\nf{Témoins}{Aucune fiche EDB appariée (pas de numéro ICUR ; recherche par le texte infructueuse).}
\nf{Traduction}{[…] vierge, vécut […] ans ; […] déposé la veille des ides […] ; sous le consulat d’Arcadius pour la 4e fois et d’Honorius pour la 3e fois.}
\nf{Langue}{9 abréviations développées par l’éditeur (dont virgo(?), an(nos), d(e)p, (ositus ?), pri(di)e, Arc(adio)).}
\nf{Chronologie}{Noms consulaires : Arcadius (restitué) et Honorius (lu).}
\nf{Bibliographie}{AE 2016, 170.}
\nf{Statut}{Lecture secondaire (EDCS) sans second témoin numérique.}
\end{notice}

\begin{notice}{386}{396}{ICUR n.s. VII, 17510 · EDCS-03900102 · EDB17954}
\label{not:386}
\nf{Datation}{396 — directe : nom consulaire lu}
\nf{Genre}{Épitaphe ; ensembles A, B, C ; 69 lettres transmises.}
\nf{Références}{Anagni 151 = ICUR VII 17510 = ILCV +2946.}
\nf{Provenance}{Rattachement selon l’EDB : \emph{Coem. Cyriacae ad s.Laurentium}. Conservation selon l’EDB : \emph{Anagni (FR)   ; monastero delle suore Cistercensi della Carità}.}
\nf{Support}{Plaque de marbre. 14 × 37 cm (h. × l.). Inscription gravée. Reproductions signalées par l’EDB : reproduction, Perret 1851 et reproduction, Solin et alii 1996 (droits réservés ; non reproduites). Fiche EDB : Anita Rocco, première version 2015.}
\ntext{\lnum{1}Arcadio V et Honorio III\\
\lnum{2}deposita Staturina X Kal(endas)\\
\lnum{3}Apriles\\
\lnum{4}Sabbatius maritus fecit}
\nf{Témoins}{EDB (EDB17954) : texte concordant aux conventions de notation près (indice de similarité 1).}
\nf{Traduction}{Sous (le consulat d’)Arcadius pour la 5e (sic) fois et Honorius pour la 3e fois, déposée Staturina le 10e jour avant les calendes d’avril ; Sabbatius, son mari, a fait (ce monument).}
\nf{Langue}{1 abréviation développée par l’éditeur (dont Kal(endas)).}
\nf{Chronologie}{Noms consulaires : Arcadius (lu) et Honorius (lu). Dates du jour : déposition : X Kalendas Apriles = 23 mars 396 (lu).}
\nf{Bibliographie}{Anagni 151 ; ICUR VII 17510 ; ILCV +2946. Bibliographie complémentaire selon l’EDB : Solin et alii 1996, 94, n. 151 con foto.}
\nf{Statut}{Lecture secondaire (EDCS) concordante avec l’EDB : lecture stabilisée entre deux bases, non collationnée sur l’ICVR ni sur le monument. L’EDCS signale une photographie (non consultée).}
\end{notice}

\begin{notice}{387}{396}{EDCS-71500011}
\label{not:387}
\nf{Datation}{396 — directe : nom consulaire lu}
\nf{Genre}{Mixte (épitaphe et acte d’achat) ; ensembles A, B, C ; 156 lettres transmises ; 47 \% restituées.}
\nf{Références}{Carletti-2008, 00177.}
\nf{Provenance}{Rome (indication de l’EDCS) ; aucun rattachement cimétérial transmis.}
\nf{Support}{Non documenté dans les sources consultées.}
\ntext{\lnum{1}[- - -] deposita X[- - - d(omino) n(ostro) Honorio Aug(usto) III et]\\
\lnum{2}d(omino) n(ostro) Fl(avio) Ar(ca)dio Aug(usto) IIII conss(ulibus) M[- - -]\\
\lnum{3}Eros et Leo et Setema locu(m) [emerunt ad limina Petri et]\\
\lnum{4}Pauli sub episcopatu Si[rici a presbb(yteris) Urso]\\
\lnum{5}et Proclino}
\nf{Témoins}{Aucune fiche EDB appariée (pas de numéro ICUR ; recherche par le texte infructueuse).}
\nf{Traduction}{[…] déposée le 10 […] sous le consulat de notre seigneur Honorius Auguste pour la 3e fois et de notre seigneur Flavius Arcadius Auguste pour la 4e fois ; […] Eros, Leo et Setema ont acheté un emplacement près du seuil de Pierre et de Paul, sous l’épiscopat de Sirice, auprès des prêtres Ursus et Proclinus.}
\nf{Langue}{11 abréviations développées par l’éditeur (dont d(omino), n(ostro), Aug(usto), Fl(avio)).}
\nf{Chronologie}{Noms consulaires : Arcadius (lu) et Honorius (restitué). Ne sont pas des consuls : [ad limina Petri et] / Pauli (apôtre Paul).}
\nf{Commentaire}{Fonction ou métier mentionné : episcopatu, presbbyteris. Commentaire développé : dossier p.~\pageref{dos:387}.}
\nf{Bibliographie}{Carletti-2008, 00177.}
\nf{Statut}{Lecture secondaire (EDCS) sans second témoin numérique. L’EDCS signale une photographie (non consultée).}
\end{notice}

\begin{notice}{388}{396}{ICUR n.s. I, 1450 · EDCS-39000429 · EDB31239}
\label{not:388}
\nf{Datation}{396 — formule transmise par l’EDB (lue)}
\nf{Genre}{Épitaphe ; ensembles A, B, C ; 101 lettres transmises.}
\nf{Références}{ICUR I 1450 = ILCV 301 (add).}
\nf{Provenance}{Rattachement selon l’EDB : \emph{Incertae originis} (le volume I des ICVR réunit les inscriptions dont la provenance est perdue). Conservation selon l’EDB : \emph{Musei Vaticani: Lapidario Cristiano}.}
\nf{Support}{Plaque de marbre. 71 × 145 cm (h. × l.). Inscription gravée. Fiche EDB : Filippo Piazzolla, première version 2015.}
\ntext{\lnum{1}Hic requiesc⸢i⸣t Quodvul(t)deus ho\\
\lnum{2}nest(a)e recordation⸢i⸣s viro ⸢q⸣ui vi\{c\}\\
\lnum{3}xit annos LVII depositus in pace\\
\lnum{4}die V Idus Octobres conss(ulibus)}
\nf{Témoins}{Export EDCS de 2020 : \emph{Hic requiesct Quodvul(t)deus ho/nest(a)e recordations viro ui vi\{c\}/xit annos LVII depositus in pace / die V Idus Octobres conss(ulibus)}. EDB (EDB31239 ; écarts à examiner) : \emph{hic requiescet Quodvuldeus ho/neste recordationes vir ⌜q⌝ui vic/xit annos LVII depositus in pace / die V idus octobres cons̅(ulibu)s̅ dd(ominis) n̅n(ostris) / Arcadio Au̅g(usto) quater et Hono/rio Au̅g(usto) ter consulibus ((folium))}.}
\nf{Apparat}{⸢i⸣ : E gravé ; ⸢i⸣ : E gravé ; ⸢q⸣ : O gravé ; \{c\} : lettres superflues}
\nf{Traduction}{Ici repose Quodvultdeus, homme d’honorable mémoire, qui vécut 57 ans ; déposé en paix le 5e jour avant les ides d’octobre, sous le consulat […].}
\nf{Langue}{L’éditeur signale 3 graphies corrigées : E gravé pour i (confusion e/i) ×2 ; O gravé pour q (confusion de lettres). Des lettres gravées sont jugées superflues par l’éditeur (entre accolades). 3 abréviations développées par l’éditeur (dont Quodvul(t)deus, nest(a)e, conss(ulibus)).}
\nf{Chronologie}{Le texte EDCS s'arrête à \emph{conss(ulibus)}, sans nom. La fiche EDB (EDB31239) transmet la suite : \emph{dd(ominis) n̅n(ostris) / Arcadio Au̅g(usto) quater et Hono/rio Au̅g(usto) ter consulibus} (396). L'année 396 est donc fondée sur une formule consulaire lue. La plaque (71 × 145 cm) est conservée aux Musées du Vatican (Lapidario Cristiano) ; l'EDB décrit une feuille. Dates du jour : déposition : die V Idus Octobres = 11 octobre 396 (lu). Âge déclaré : 57 a. 5 j. (lu).}
\nf{Commentaire}{La formule in pace est conservée. Vocabulaire religieux : repos et sommeil de la mort — requiescit (lu) ; anthroponyme (niveau 6) — quodvultdeus (anthroponyme). Commentaire développé : dossier p.~\pageref{dos:388}.}
\nf{Bibliographie}{ICUR I 1450 ; ILCV 301 (add). Bibliographie complémentaire selon l’EDB : Ferrua 1964, 109 ; Ferrua 1979b, 38, n. 1450 ; Lega 2000, 141 ; ILCV, 301.}
\nf{Statut}{Lecture secondaire (EDCS) ; écarts avec l’EDB à arbitrer sur l’ICVR. L’EDCS signale une photographie (non consultée).}
\end{notice}

\begin{notice}{389}{396}{ICUR n.s. I, 2106 · EDCS-39300325 · EDB32100}
\label{not:389}
\nf{Datation}{396 — directe : nom consulaire lu}
\nf{Genre}{Épitaphe ; ensembles A, B, C ; 77 lettres transmises ; 29 \% restituées.}
\nf{Références}{ICUR I 2106.}
\nf{Provenance}{Rattachement selon l’EDB : \emph{Incertae originis} (le volume I des ICVR réunit les inscriptions dont la provenance est perdue). Conservation selon l’EDB : \emph{Musei Vaticani: Galleria Lapidaria}.}
\nf{Support}{Plaque de marbre. 36 × 40 cm (h. × l.). Inscription gravée. Fiche EDB : Filippo Piazzolla, première version 2015.}
\ntext{\lnum{1}vi]xit annis p(lus) m(inus) n(umero) XI\\
\lnum{2}[- - - O]ctobr⸢e⸣s dd(ominis) nn(ostris)\\
\lnum{3}[Arcadio] IIII et Honorio III\\
\lnum{4}[cons(ulibus) bene] merenti in pace}
\nf{Témoins}{Export EDCS de 2020 : \emph{vi]xit annis p(lus) m(inus) n(umero) XI / [3 O]ctobrs dd(ominis) nn(ostris) / [Arcadio] IIII et Honorio III / [cons(ulibus) bene] merenti in pace}. EDB (EDB32100) : texte concordant aux conventions de notation près (indice de similarité 0,973).}
\nf{Apparat}{⸢e⸣ : I gravé}
\nf{Traduction}{[…] vécut environ 11 ans ; […] octobre, sous le consulat de nos seigneurs Arcadius pour la 4e fois et Honorius pour la 3e fois ; bien méritant(e), en paix.}
\nf{Langue}{L’éditeur signale une graphie corrigée : I gravé pour e (confusion e/i).}
\nf{Chronologie}{Noms consulaires : Arcadius (restitué) et Honorius (lu). Âge déclaré : env. 11 a. (partiellement restitué).}
\nf{Commentaire}{La formule in pace est conservée.}
\nf{Bibliographie}{ICUR I 2106. Bibliographie complémentaire selon l’EDB : Lega 2000.}
\nf{Statut}{Lecture secondaire (EDCS) concordante avec l’EDB : lecture stabilisée entre deux bases, non collationnée sur l’ICVR ni sur le monument.}
\end{notice}

\begin{notice}{390}{396}{ICUR n.s. I, 2107 · EDCS-39300326 · EDB32101}
\label{not:390}
\nf{Datation}{396 — directe : nom consulaire lu}
\nf{Genre}{Épitaphe ; ensembles A, B, C ; 85 lettres transmises ; 42 \% restituées.}
\nf{Références}{ICUR I 2107.}
\nf{Provenance}{Rattachement selon l’EDB : \emph{Incertae originis} (le volume I des ICVR réunit les inscriptions dont la provenance est perdue). Conservation selon l’EDB : \emph{Musei Vaticani: Galleria Lapidaria}.}
\nf{Support}{Plaque de marbre. 17 × 61 cm (h. × l.). Inscription gravée. Fiche EDB : Filippo Piazzolla, première version 2015.}
\ntext{\lnum{1}Hic iac⸢e⸣t Rerinta [- - -]\\
\lnum{2}qu(a)e vixit an(n)is LIII d[iebus - - - dep(osita)]\\
\lnum{3}X Kal(endas) Aug(ustas) Arcadio III[I et Honorio III cons(ulibus)]}
\nf{Témoins}{Export EDCS de 2020 : \emph{Hic iacit Rerinta [3] / qu(a)e vixit an(n)is LIII d[iebus 3 dep(osita)] / X Kal(endas) Aug(ustas) Arcadio III[I et Honorio III cons(ulibus)]}. EDB (EDB32101 ; écarts à examiner) : \emph{hic iacit Rerinta [---] / que vixit anis LIII d(ies) [--- deposita] / X kal(endas) aug(ustas) Arcadio III[I et Honorio III cons(ulibus)]}.}
\nf{Apparat}{⸢e⸣ : I gravé}
\nf{Traduction}{Ici gît Rerinta (?) […] qui vécut 53 ans […] jours ; déposée le 10e jour avant les calendes d’août, sous le consulat d’Arcadius pour la 4e fois et d’Honorius pour la 3e fois.}
\nf{Langue}{L’éditeur signale une graphie corrigée : I gravé pour e (confusion e/i). L’export de 2020 diffère de l’état 2022 (voir apparat). 6 abréviations développées par l’éditeur (dont qu(a)e, an(n)is, dep(osita), Kal(endas), Aug(ustas), cons(ulibus)).}
\nf{Chronologie}{Noms consulaires : Arcadius (lu) et Honorius (restitué). Dates du jour : déposition : X Kalendas Augustas = 23 juillet 396 (lu). Âge déclaré : 53 a. (partiellement restitué).}
\nf{Bibliographie}{ICUR I 2107. Bibliographie complémentaire selon l’EDB : Ferrua 1979b, 63 ; ILCV, 3057A adn. ; Lega 2000.}
\nf{Statut}{Lecture secondaire (EDCS) ; écarts avec l’EDB à arbitrer sur l’ICVR.}
\end{notice}

\begin{notice}{391}{396}{ICUR n.s. I, 4073 · EDCS-40301037 · EDB34240}
\label{not:391}
\nf{Datation}{396 — directe : nom consulaire lu}
\nf{Genre}{Épitaphe ; ensembles A, B, C ; 209 lettres transmises ; 12 \% restituées.}
\nf{Références}{ICUR I 4073 = ILCV 1483 (add).}
\nf{Provenance}{Rattachement selon l’EDB : \emph{Incertae originis} (le volume I des ICVR réunit les inscriptions dont la provenance est perdue).}
\nf{Support}{Texte métrique. Fiche EDB : Filippo Piazzolla, première version 2015.}
\ntext{\lnum{1}Huic tumulo quiesc⸢i⸣t totius fi[dei - - - et]\\
\lnum{2}integritatis homo opinionis BONIAV[- - -]\\
\lnum{3}ratione fidelis mentis integrae [- - - amicus]\\
\lnum{4}amicorum ob hoc igitur neof⸢y⸣tus SPAC[- - -]\\
\lnum{5}Aufinianus qui vixit annis XLII depositus [- - -]\\
\lnum{6}dd(ominis) nn(ostris) Arcadio IIII et Honorio III Aug[g(ustis)]\\
\lnum{7}co(n)s[s(ulibus)]}
\nf{Témoins}{Export EDCS de 2020 : \emph{Huic tumulo quiesct totius fi[dei 3 et] / integritatis homo opinionis BONIAV[3] / ratione fidelis mentis integrae [3 amicus] / amicorum ob hoc igitur neoftus SPAC[3] / Aufinianus qui vixit annis XLII depositus [3] / dd(ominis) nn(ostris) Arcadio IIII et Honorio III Aug[g(ustis)] / co(n)s[s(ulibus)]}. EDB (EDB34240 ; écarts à examiner) : \emph{huic tumulo quiescet totius fi[dei --- et] / integritatis homo opinionis boni av[---] / ratione fidelis mentis integrae [--- amicus] / amicorum ob hoc igitur neofitus Spac[---] / Aufinianus qui vixit annis XLII depositus [---] / dd(ominis) nn(ostris) Arcadio IIII et Honorio III aug[g(ustis) ---] / co(n)s[(ulibu)s]}.}
\nf{Apparat}{⸢i⸣ : E gravé ; ⸢y⸣ : I gravé}
\nf{Traduction}{Dans ce tombeau repose un homme d’une foi et d’une intégrité entières […] d’une réputation […] ; fidèle par la raison, d’un esprit intègre […] ami des amis ; pour cette raison donc néophyte […] Aufinianus, qui vécut 42 ans ; déposé […] sous le consulat de nos seigneurs Arcadius pour la 4e fois et Honorius pour la 3e fois, Augustes.}
\nf{Langue}{L’éditeur signale 2 graphies corrigées : E gravé pour i (confusion e/i) ; I gravé pour y (graphie de y).}
\nf{Chronologie}{Noms consulaires : Arcadius (lu) et Honorius (lu). Âge déclaré : 42 a. (lu).}
\nf{Commentaire}{Vocabulaire religieux : initiation et foi — fidei (partiellement restitué), fidelis (lu), neofytus (lu) ; repos et sommeil de la mort — quiescit (lu).}
\nf{Bibliographie}{ICUR I 4073 ; ILCV 1483 (add). Bibliographie complémentaire selon l’EDB : ILCV, 1483.}
\nf{Statut}{Lecture secondaire (EDCS) ; écarts avec l’EDB à arbitrer sur l’ICVR.}
\end{notice}

\begin{notice}{392}{396}{ICUR n.s. III, 6510 · EDCS-34801310 · EDB22305}
\label{not:392}
\nf{Datation}{396 — directe : nom consulaire lu}
\nf{Genre}{Épitaphe ; ensembles A, B, C ; 50 lettres transmises ; 56 \% restituées.}
\nf{Références}{ICUR III 6510.}
\nf{Provenance}{Rattachement selon l’EDB : \emph{Coem. Domitillae pars inferior}. Conservation selon l’EDB : \emph{Basilica ss. Nerei et Achillei in coem. Domitillae}.}
\nf{Support}{Plaque de marbre. 27 × 43 cm (h. × l.). Inscription gravée. Reproductions signalées par l’EDB : photographie PCAS (droits réservés ; non reproduites). Fiche EDB : Carlo Carletti / Antonio Enrico Felle, première version 2013.}
\ntext{\lnum{1}Dep(ositus ?) pr[idie - - - con]\\
\lnum{2}sulatus Arc(a)d[io Aug(usto) quater et Hono]\\
\lnum{3}rio ter [}
\nf{Témoins}{EDB (EDB22305 ; écarts à examiner) : \emph{------ / ++[---] / dep(osit---) pr[idie --- con]/sulatus Arcad[io Aug(usto) quater et Hono]rio ter [texte tronqué dans la liste de résultats]}.}
\nf{Traduction}{Déposé la veille […] sous le consulat d’Arcadius Auguste pour la 4e fois et d’Honorius pour la 3e fois.}
\nf{Langue}{3 abréviations développées par l’éditeur (dont Dep(ositus ?), Arc(a)d, Aug(usto)).}
\nf{Chronologie}{Noms consulaires : Arcadius (lu) et Honorius (lu en partie).}
\nf{Bibliographie}{ICUR III 6510.}
\nf{Statut}{Lecture secondaire (EDCS) ; écarts avec l’EDB à arbitrer sur l’ICVR. Plus de la moitié des lettres sont restituées. L’EDCS signale une photographie (non consultée).}
\end{notice}

\begin{notice}{393}{396}{ICUR n.s. IV, 11126a · EDCS-37200412 · EDB39590}
\label{not:393}
\nf{Datation}{396 — directe : consulat restitué}
\nf{Genre}{Épitaphe ; ensembles A, B, C ; 41 lettres transmises ; 71 \% restituées.}
\nf{Références}{ICUR IV 11126a.}
\nf{Provenance}{Rattachement selon l’EDB : \emph{Coem. Callisti area subdialis et incerti loci}. Conservation selon l’EDB : \emph{Coem. Callisti area subdialis et incerti loci}.}
\nf{Support}{Plaque de marbre. 16 × 13 × 2 cm (h. × l. × ép.). Inscription gravée. Fiche EDB : Marida Pierno, première version 2015.}
\ntext{\lnum{1}Arcadio] IIII H[onorio III Augg(ustis) conss(ulibus)]}
\nf{Témoins}{EDB (EDB39590 ; divergence forte) : \emph{⟨:a⟩ ------ / [---]+[---] / [---]IIII H[---] / ------ /⟨:b⟩ [---]ADIO[---] / [---]o III Aug(usto) [---]}. L'EDB (EDB39590 = ICVR IV 11126) publie une fiche unique pour les fragments a et b, de très petite taille (16 × 13 × 2 cm). Elle ne lit ni \emph{Arc-} ni \emph{Hon-} : le texte lu se réduit à \emph{…IIII H…}, \emph{…ADIO…} et \emph{…o III Aug(usto)…}, compatible avec le consulat d'Arcadius IV et Honorius III (396). Les deux fragments sont réunis sous ce numéro ; l'ancien n° 394 y est joint. Un seul texte daté de 396 est compté. Absorbe la notice 394 (voir cette notice).}
\nf{Traduction}{Sous le consulat d’Arcadius pour la 4e fois et d’Honorius pour la 3e fois, Augustes.}
\nf{Langue}{2 abréviations développées par l’éditeur (dont Augg(ustis), conss(ulibus)).}
\nf{Chronologie}{Noms consulaires : Arcadius (restitué) et Honorius (restitué).}
\nf{Commentaire}{Commentaire développé : dossier p.~\pageref{dos:393}.}
\nf{Bibliographie}{ICUR IV 11126a.}
\nf{Statut}{Lecture secondaire (EDCS) ; divergence forte avec l’EDB : le texte ne doit pas être exploité sans contrôle de l’édition imprimée. Plus de la moitié des lettres sont restituées.}
\end{notice}

\begin{notice}{394}{396}{ICUR n.s. IV, 11126b · EDCS-37200413 · EDB39590}
\label{not:394}
\nf{Statut}{Fragment b de la même plaque que le n° 393 (EDB39590 = ICVR IV 11126). Joint au n° 393 ; un seul texte daté est compté.}
\nf{Références}{ICUR IV 11126b}
\ntext{\lnum{1}Arc]adio [Hon]\\
\lnum{2}[ori]o III Aug[g(ustis) conss(ulibus)]}
\nf{Traduction}{Sous le consulat d’Arcadius et d’Honorius pour la 3e fois, Augustes.}
\nf{Voir}{la notice 393 (p.~\pageref{not:393}).}
\end{notice}

\begin{notice}{395}{396}{ICUR n.s. IV, 12252a · EDCS-35700729 · EDB7067}
\label{not:395}
\nf{Datation}{396 — directe : nom consulaire lu}
\nf{Genre}{Épitaphe ; ensembles A, B, C ; 91 lettres transmises.}
\nf{Références}{ICUR IV 12252a.}
\nf{Provenance}{Rattachement selon l’EDB : \emph{Coem. anonymum ad viam Ardeatinam}. Conservation selon l’EDB : \emph{Coem. anonymum ad viam Ardeatinam}.}
\nf{Support}{Enduit de mortier de mortier. Inscription tracée dans le mortier frais. Fiche EDB : Marida Pierno, première version 2013.}
\ntext{\lnum{1}Lo\\
\lnum{2}cus tri\\
\lnum{3}so\\
\lnum{4}mus\\
\lnum{5}ubi positus Iustus qui depositus\\
\lnum{6}est VIIII Kalendas Octobr⸢e⸣s\\
\lnum{7}Fl(avio) Arcadio IIII et Honorio IIII}
\nf{Témoins}{Export EDCS de 2020 : \emph{Lo/cus tri/so/mus // ubi positus Iustus qui depositus // est VIIII Kalendas Octobrs // Fl(avio) Arcadio IIII et Honorio IIII}. EDB (EDB7067) : texte concordant aux conventions de notation près (indice de similarité 0,986).}
\nf{Apparat}{⸢e⸣ : I gravé}
\nf{Traduction}{Emplacement triple // où a été placé Iustus, qui a été déposé // le 9e jour avant les calendes d’octobre // sous (le consulat de) Flavius Arcadius pour la 4e fois et Honorius pour la 4e (sic) fois.}
\nf{Langue}{L’éditeur signale une graphie corrigée : I gravé pour e (confusion e/i).}
\nf{Chronologie}{Noms consulaires : Arcadius (lu) et Honorius (lu). Réserve : Itération d’Honorius (IIII transmis) à confronter au collège de 396 (Arcadius IV et Honorius III). Dates du jour : déposition : VIIII Kalendas Octobres = 23 septembre 396 (lu).}
\nf{Commentaire}{Le locus trisomus est un emplacement qualifié pour trois corps ; cette désignation ne prouve pas que trois défunts sont nommés ou que trois dépositions sont ici documentées.}
\nf{Bibliographie}{ICUR IV 12252a.}
\nf{Statut}{Lecture secondaire (EDCS) concordante avec l’EDB : lecture stabilisée entre deux bases, non collationnée sur l’ICVR ni sur le monument.}
\end{notice}

\annee{397}

\begin{notice}{396}{397}{EDCS-02700238 · EDB42676}
\label{not:396}
\nf{Datation}{397 — directe : nom consulaire lu}
\nf{Genre}{Épitaphe ; ensembles A, B, C ; 107 lettres transmises ; 27 \% restituées.}
\nf{Références}{CEACelio 466 = AE 1994, 163 = AE 2001, 219.}
\nf{Provenance}{Rattachement selon l’EDB : \emph{Incertae originis} (le volume I des ICVR réunit les inscriptions dont la provenance est perdue). Conservation selon l’EDB : \emph{Antiquarium del Celio}.}
\nf{Support}{La fiche EDB ne figure pas dans l’état 2019 des données EpiDoc ; support non documenté.}
\ntext{\lnum{1}] et Maximiano bene merenti in pace qui vixit\\
\lnum{2}[an(nos) p(lus) m(inus) L]XX dep(osito) die XV Kal(endas) Septe(mbres) Fl(avio) C(a)esari\\
\lnum{3}[o et Nonio Attico] v(iris) c(larissimis) cons(ulibus) hic e(s)t locus\\
\lnum{4}[}
\nf{Témoins}{Export EDCS de 2020 : \emph{] et Maximiano bene merenti in pace qui vixit / [an(nos) p(lus) m(inus) L]XX dep(osito) die XV Kal(endas) Septe(mbres) Fl(avio) C(a)esari / et Nonio Attico v(iris) c(larissimis) cons(ulibus) hic e(s)t locus / [}. EDB (EDB42676) : texte concordant aux conventions de notation près (indice de similarité 0,98).}
\nf{Traduction}{[…] et à Maximianus, bien méritant, en paix, qui vécut environ 70 ans ; déposé le 15e jour avant les calendes de septembre, sous le consulat de Flavius Caesarius et de Nonius Atticus, clarissimes ; ceci est l’emplacement […].}
\nf{Langue}{L’export de 2020 diffère de l’état 2022 (voir apparat). 12 abréviations développées par l’éditeur (dont an(nos), p(lus), m(inus), dep(osito), Kal(endas), Septe(mbres)).}
\nf{Chronologie}{Noms consulaires : Caesarius (lu) et Atticus (restitué). Dates du jour : déposition : die XV Kalendas Septembres = 18 août 397 (lu). Âge déclaré : env. 70 a. 15 j. (partiellement restitué).}
\nf{Commentaire}{La formule in pace est conservée.}
\nf{Bibliographie}{CEACelio 466 ; AE 1994, 163 ; AE 2001, 219. Bibliographie complémentaire selon l’EDB : Ricci et alii 1994, 187-188, n. 18 ; AE, 1994, 163 ; Gregori 2001, 394, n. 466 (M. Giocoli - M. Manganaro).}
\nf{Statut}{Lecture secondaire (EDCS) concordante avec l’EDB : lecture stabilisée entre deux bases, non collationnée sur l’ICVR ni sur le monument.}
\end{notice}

\begin{notice}{397}{397}{ICUR n.s. I, 3214 · EDCS-23102819 · EDB35690}
\label{not:397}
\nf{Datation}{397 — directe : nom consulaire lu}
\nf{Genre}{Épitaphe ; ensembles A, B, C ; 83 lettres transmises ; 35 \% restituées.}
\nf{Références}{CIL VI 39147a = CLE 2005 = ICUR I 3214 = ICUR X 27324 = ILCV 4768 = AE 1909, 27 = AE 1933, 158.}
\nf{Provenance}{Rattachement selon l’EDB : \emph{Coem. et basilica s.Valentini}.}
\nf{Support}{Plaque de marbre. 20,6 × 36,5 × 3 cm (h. × l. × ép.). Inscription gravée. Texte métrique. Fiche EDB : Filippo Piazzolla, première version 2015.}
\ntext{\lnum{1}]romi[- - -]\\
\lnum{2}[- - -] sed properata dies [- - -]\\
\lnum{3}[- - -]os bis denos iuvenis post [- - -]\\
\lnum{4}[- - -] Caesario et Att[ico v(iris) c(larissimis) cons(ulibus)}
\nf{Témoins}{EDB (EDB35690 ; divergence forte) : \emph{------ / Ṛọṃạṇ[---] / sed properata dies f[ato consumpsit acerbo?] /[---]os bis denos iuvenis post[quam compleverat annos?] / Caesario et Atti[co conss(ulibus)]}. L'EDB (EDB35690 = ICVR X 27324 = CIL VI 39147a) lit \emph{Ṛọṃạṇ[---]} au lieu de \emph{]romi[} et propose une restitution métrique (\emph{sed properata dies f[ato consumpsit acerbo?]}). Pierre perdue.}
\nf{Traduction}{(vers, fragmentaire) […] mais le jour (fatal) s’est hâté […] ; […] jeune homme, après deux fois dix (ans) […] ; sous le consulat de Caesarius et d’Atticus, clarissimes.}
\nf{Langue}{3 abréviations développées par l’éditeur (dont v(iris), c(larissimis), cons(ulibus)).}
\nf{Chronologie}{Noms consulaires : Caesarius (lu) et Atticus (lu en partie).}
\nf{Bibliographie}{CIL VI 39147a ; CLE 2005 ; ICUR I 3214 ; ICUR X 27324 ; ILCV 4768 ; AE 1909, 27 ; AE 1933, 158. Bibliographie complémentaire selon l’EDB : ICVR X, 27324 ; CIL VI, 39147.a.}
\nf{Statut}{Lecture secondaire (EDCS) ; divergence forte avec l’EDB : le texte ne doit pas être exploité sans contrôle de l’édition imprimée.}
\end{notice}

\begin{notice}{398}{397}{ICUR n.s. I, 471 · EDCS-38700879 · EDB28397}
\label{not:398}
\nf{Datation}{397 — directe : nom consulaire lu}
\nf{Genre}{Épitaphe ; ensembles A, B, C ; 50 lettres transmises ; 32 \% restituées.}
\nf{Références}{ICUR I 471.}
\nf{Provenance}{Rattachement selon l’EDB : \emph{Incertae originis} (le volume I des ICVR réunit les inscriptions dont la provenance est perdue). Conservation selon l’EDB : \emph{s. Maria in Aracoeli}.}
\nf{Support}{Plaque de marbre. 82 × 69 cm (h. × l.). Inscription gravée. Fiche EDB : Filippo Piazzolla, première version 2014.}
\ntext{\lnum{1}Paulus qui vixit ann(os) LIIII [- - -]\\
\lnum{2}Feb(ruarias) C(a)esario et [Attico cons(ulibus)]}
\nf{Témoins}{EDB (EDB28397 ; écarts à examiner) : \emph{Paulus qui vixit ann(is) LIIII [ men(sibus) --- d(iebus) --- dep(ositus) ---] / feb(ruar---) Cesario et [Attico cons(ulibu)s --- ]}.}
\nf{Traduction}{Paulus, qui vécut 54 ans […] ; […] février, sous le consulat de Caesarius et d’Atticus.}
\nf{Langue}{4 abréviations développées par l’éditeur (dont ann(os), Feb(ruarias), C(a)esario, cons(ulibus)).}
\nf{Chronologie}{Noms consulaires : Caesarius (lu) et Atticus (restitué). Ne sont pas des consuls : Paulus qui vixit ann(os) LIIII (défunt). Âge déclaré : 54 a. (lu).}
\nf{Bibliographie}{ICUR I 471.}
\nf{Statut}{Lecture secondaire (EDCS) ; écarts avec l’EDB à arbitrer sur l’ICVR.}
\end{notice}

\begin{notice}{399}{397}{ICUR n.s. I, 717 · EDCS-38701118 · EDB28944}
\label{not:399}
\nf{Datation}{397 — directe : nom consulaire lu}
\nf{Genre}{Épitaphe ; ensembles A, B, C ; 263 lettres transmises.}
\nf{Références}{ICUR I 717 = ILCV 4671.}
\nf{Provenance}{Rattachement selon l’EDB : \emph{Incertae originis} (le volume I des ICVR réunit les inscriptions dont la provenance est perdue).}
\nf{Support}{Support non précisé par l’EDB. Fiche EDB : Filippo Piazzolla, première version 2014.}
\ntext{\lnum{1}Sanctae coniugi Caelestinae quae vixit annis XX men\\
\lnum{2}sem(!) I dies XXII fecit autem cum Ampliato marito annis\\
\lnum{3}IIII dies XXI bene merenti in pace deposita VIII Kal(endas) Iunias\\
\lnum{4}Hilaritati filiae bene merenti in pace quae vixit annis\\
\lnum{5}XXIII m(ensibus) VIIII d(iebus) XXVIII fecit cum marito annos V d(ies) VI\\
\lnum{6}deposita IIII Non(as) Iul(ias) Caesario et Attico cons(ulibus)}
\nf{Témoins}{EDB (EDB28944) : texte concordant aux conventions de notation près (indice de similarité 0,981).}
\nf{Apparat}{(!) : forme signalée par l’éditeur}
\nf{Traduction}{À sa sainte épouse Caelestina, qui vécut 20 ans 1 mois 22 jours et a vécu avec Ampliatus, son mari, 4 ans 21 jours ; bien méritante, en paix ; déposée le 8e jour avant les calendes de juin. À Hilaritas, fille bien méritante, en paix, qui vécut 23 ans 9 mois 28 jours et a vécu avec son mari 5 ans 6 jours ; déposée le 4e jour avant les nones de juillet, sous le consulat de Caesarius et d’Atticus.}
\nf{Langue}{L’éditeur signale au moins une forme anormale (sic). 8 abréviations développées par l’éditeur (dont sem(!), Kal(endas), m(ensibus), d(iebus), d(ies), Non(as)).}
\nf{Chronologie}{Noms consulaires : Caesarius (lu) et Atticus (lu). Dates du jour : déposition : VIII Kalendas Iunias = 25 mai 397 (lu) ; déposition : IIII Nonas Iulias = 4 juillet 397 (lu). Âges déclarés : 20 a. 1 m. 22 j. (lu) ; 23 a. 9 m. 28 j. (lu).}
\nf{Commentaire}{La formule in pace est conservée. Vocabulaire religieux : épithète du défunt (non doctrinal) — Sanctae (lu) ; allusion interprétative (niveau 5) — Caelestinae (lu).}
\nf{Bibliographie}{ICUR I 717 ; ILCV 4671. Bibliographie complémentaire selon l’EDB : Ferrua 1979b, 21, n. 717 ; ILCV, 4671.}
\nf{Statut}{Lecture secondaire (EDCS) concordante avec l’EDB : lecture stabilisée entre deux bases, non collationnée sur l’ICVR ni sur le monument.}
\end{notice}

\begin{notice}{400}{397}{ICUR n.s. I, 941 · EDCS-38700455 · EDB29526}
\label{not:400}
\nf{Datation}{397 — directe : nom consulaire lu}
\nf{Genre}{Mixte (épitaphe et acte d’achat) ; ensembles A, B, C ; 188 lettres transmises.}
\nf{Références}{ICUR I 941 = ILCV +1002 = ILCV 3444 = Mossong-2022, 00030.}
\nf{Provenance}{Rattachement selon l’EDB : \emph{Incertae originis} (le volume I des ICVR réunit les inscriptions dont la provenance est perdue). Conservation selon l’EDB : \emph{s. Prassede}.}
\nf{Support}{Plaque de marbre. 44 × 175 cm (h. × l.). Inscription gravée. Signes décrits par l’EDB : chrisme et croix monogrammatique. Reproductions signalées par l’EDB : reproduction, Perret 1851 (droits réservés ; non reproduites). Fiche EDB : Filippo Piazzolla, première version 2014.}
\ntext{\lnum{1}Dulcis et innoce(n)s hic dormit Severianus in somn\{i\}o pacis\\
\lnum{2}qui vixit ann⸢o⸣s p(lus) m(inus) L cuius spiritus in luce domini suscep\\
\lnum{3}tus est die VIIIIX(!) K(a)l(endas) Mart(ias) dd(ominis) nn(ostris) Flavio Caesario et Nonio Attico\\
\lnum{4}vv(iris) cc(larissimis) (con)ss(ulibus) quem locum emit Viricunda ux\{x\}or Pascasio ep(is)c(opo)}
\nf{Témoins}{Export EDCS de 2020 : \emph{Dulcis et innoce(n)s hic dormit Severianus in somn\{i\}o pacis / qui vixit anns p(lus) m(inus) L cuius spiritus in luce domini suscep/tus est die VIIIIX(!) K(a)l(endas) Mart(ias) dd(ominis) nn(ostris) Flavio Caesario et Nonio Attico / vv(iris) cc(larissimis) (con)ss(ulibus) quem locum emit Viricunda ux\{x\}or Pascasio ep(is)c(opo)}. EDB (EDB29526) : même texte ; les écarts relevés dans la liste de résultats tiennent à la notation (relecture de la fiche détaillée).}
\nf{Apparat}{⸢o⸣ : V gravé ; \{i\} : lettres superflues ; \{x\} : lettres superflues ; (!) : forme signalée par l’éditeur}
\nf{Traduction}{Doux et innocent, ici dort Severianus dans le sommeil de la paix ; il vécut environ 50 ans ; son esprit a été accueilli dans la lumière du Seigneur le 19e (?) jour avant les calendes de mars, sous le consulat de nos seigneurs (sic) Flavius Caesarius et Nonius Atticus, clarissimes ; Viricunda, son épouse, a acheté cet emplacement à l’évêque (?) Paschasius.}
\nf{Langue}{L’éditeur signale une graphie corrigée : V gravé pour o (confusion o/u). Des lettres gravées sont jugées superflues par l’éditeur (entre accolades). L’éditeur signale au moins une forme anormale (sic). 15 abréviations développées par l’éditeur (dont innoce(n)s, p(lus), m(inus), VIIIIX(!), K(a)l, (endas)).}
\nf{Chronologie}{Noms consulaires : Caesarius (lu) et Atticus (lu). Dates du jour : événement indéterminé : die VIIIIX Kalendas Martias = 13 février 397 (lu). Âge déclaré : env. 50 a. (lu).}
\nf{Commentaire}{Vocabulaire religieux : désignation de Dieu — domini (lu) ; âme humaine (non doctrinal) — spiritus (lu) ; repos et sommeil de la mort — dormit (lu) ; allusion interprétative (niveau 5) — luce (lu). Fonction ou métier mentionné : episcopo. Spiritus se rapporte dans la construction à Severianus ; il n’est pas automatiquement une désignation de l’Esprit saint. Domini relève d’une qualification religieuse distincte. La présence d’un évêque est transmise, mais sa fonction exacte dans l’acquisition n’est pas résolue. L'EDB (EDB29526) décrit un chrisme, deux croix monogrammatiques (dont une dans une couronne) et un oiseau tenant un rameau d'olivier. Elle reprend la note de Diehl : \emph{VIIIIX kal. mart.} serait à corriger en \emph{prid. id. Febr.}, et aucun évêque de Rome ne porte le nom de Pascasius. La fonction de ce \emph{Pascasius ep(is)c(opus)} dans l'acquisition reste donc indéterminée. Désignation divine : retenue dans le profil strict. Commentaire développé : dossier p.~\pageref{dos:400}.}
\nf{Bibliographie}{ICUR I 941 ; ILCV +1002 ; ILCV 3444 ; Mossong-2022, 00030. Bibliographie complémentaire selon l’EDB : ILCV, 3444 ; CLRE, 659.}
\nf{Statut}{Lecture secondaire (EDCS) concordante avec l’EDB : lecture stabilisée entre deux bases, non collationnée sur l’ICVR ni sur le monument. L’EDCS signale une photographie (non consultée).}
\end{notice}

\begin{notice}{401}{397}{ICUR n.s. I, 1451 · EDCS-39000430 · EDB31240}
\label{not:401}
\nf{Datation}{397 — directe : lecture partielle d’un nom consulaire}
\nf{Genre}{Épitaphe ; ensembles A, B, C ; 45 lettres transmises ; 58 \% restituées.}
\nf{Références}{ICUR I 1451.}
\nf{Provenance}{Rattachement selon l’EDB : \emph{Incertae originis} (le volume I des ICVR réunit les inscriptions dont la provenance est perdue). Conservation selon l’EDB : \emph{Musei Vaticani: Lapidario Cristiano}.}
\nf{Support}{Plaque de marbre. 29 × 24 cm (h. × l.). Inscription gravée. Fiche EDB : Filippo Piazzolla, première version 2015.}
\ntext{\lnum{1}[- - -]bina in [pace - - -]\\
\lnum{2}[- - - Cae]sario et [Attico cons(ulibus) - - -]\\
\lnum{3}[- - - qua]e vixit [}
\nf{Témoins}{EDB (EDB31240 ; écarts à examiner) : \emph{[--- ?] ((avis)) [---?] / [--- Sa?]bina in p̲a̲c̲[e ---?] / [Cae]sario et [Attico cons(ulibu)s] / [--- qua]e vixit [ ---]}.}
\nf{Traduction}{[…]bina, en paix […] ; sous le consulat de Caesarius et d’Atticus […] ; […] qui vécut […].}
\nf{Langue}{1 abréviation développée par l’éditeur (dont cons(ulibus)).}
\nf{Chronologie}{Noms consulaires : Caesarius (lu en partie) et Atticus (restitué).}
\nf{Commentaire}{La formule in pace est attestée seulement par restitution.}
\nf{Bibliographie}{ICUR I 1451. Bibliographie complémentaire selon l’EDB : Lega 2000, 61.}
\nf{Statut}{Lecture secondaire (EDCS) ; écarts avec l’EDB à arbitrer sur l’ICVR. Plus de la moitié des lettres sont restituées. L’EDCS signale une photographie (non consultée).}
\end{notice}

\begin{notice}{402}{397}{ICUR n.s. I, 2108 · EDCS-39300327 · EDB32105}
\label{not:402}
\nf{Datation}{397 — directe : lecture partielle d’un nom consulaire}
\nf{Genre}{Épitaphe ; ensembles A, B, C ; 113 lettres transmises ; 44 \% restituées.}
\nf{Références}{ICUR I 2108.}
\nf{Provenance}{Rattachement selon l’EDB : \emph{Incertae originis} (le volume I des ICVR réunit les inscriptions dont la provenance est perdue). Conservation selon l’EDB : \emph{Musei Vaticani: Galleria Lapidaria}.}
\nf{Support}{Plaque de marbre. 50 × 58 cm (h. × l.). Inscription gravée. Signes décrits par l’EDB : croix monogrammatique. Fiche EDB : Filippo Piazzolla, première version 2015.}
\ntext{\lnum{1}[Caesario et A]⸢t⸣tic⸢o⸣ vv(iris) cc(larissimis) con[s(ulibus) - - -]\\
\lnum{2}[- - -]AICI se ⸢v⸣i⸢v⸣um fec[it]\\
\lnum{3}[cum fili]⸢o⸣ su⸢o⸣ Beneri⸢o⸣\\
\lnum{4}[qui vixit anno]s III m(enses) VII d(epositus ?) II Idus Mai[as]\\
\lnum{5}[bene meren]ti in pace\{s\}}
\nf{Témoins}{Export EDCS de 2020 : \emph{[Caesario et A]tic vv(iris) cc(larissimis) con[s(ulibus) 3] / [3]AICI se ium fec[it] / [cum fili] su Beneri / [qui vixit anno]s III m(enses) VII d(epositus ?) II Idus Mai[as] / [bene meren]ti in pace\{s\}}. EDB (EDB32105 ; écarts à examiner) : \emph{[--- Caesarii et A]dtici vv(irorum) cc(larissimorum) con[s(ulum)] / [---]+ici se bibum fec[it] / [cum fili]ụm suum Benerium [---] / [qui vixit anno]s III m(enses) VII d(ies) II idus mar[tias] / [benemeren]ti in pace ((monogramma Christi in crucem translatum))}.}
\nf{Apparat}{⸢t⸣ : D gravé ; ⸢o⸣ : I gravé ; ⸢v⸣ : B gravé ; ⸢v⸣ : B gravé ; ⸢o⸣ : VM gravé ; ⸢o⸣ : VM gravé ; ⸢o⸣ : VM gravé ; \{s\} : lettres superflues}
\nf{Traduction}{Sous le consulat de Caesarius et d’Atticus, clarissimes […] ; […] de son vivant a fait (ce monument) avec son fils Venerius (?), qui vécut 3 ans 7 mois ; déposé le 2e jour avant les ides de mai ; bien méritant, en paix.}
\nf{Langue}{L’éditeur signale 7 graphies corrigées : D gravé pour t ; I gravé pour o ; B gravé pour v (bêtacisme) ×2 ; VM gravé pour o ×3. Des lettres gravées sont jugées superflues par l’éditeur (entre accolades). 5 abréviations développées par l’éditeur (dont vv(iris), cc(larissimis), cons(ulibus), m(enses), d(epositus ?)).}
\nf{Chronologie}{Noms consulaires : Caesarius (restitué) et Atticus (lu en partie). Dates du jour : déposition : II Idus Maias = 14 mai 397 (partiellement restitué). Âge déclaré : 3 a. 7 m. (partiellement restitué).}
\nf{Commentaire}{La formule in pace est conservée.}
\nf{Bibliographie}{ICUR I 2108. Bibliographie complémentaire selon l’EDB : Lega 2000 ; ILCV, 4146E adn.}
\nf{Statut}{Lecture secondaire (EDCS) ; écarts avec l’EDB à arbitrer sur l’ICVR.}
\end{notice}

\begin{notice}{403}{397}{ICUR n.s. I, 2109 · EDCS-39300328 · EDB32106}
\label{not:403}
\nf{Datation}{397 — directe : nom consulaire lu}
\nf{Genre}{Épitaphe ; ensembles A, B, C ; 61 lettres transmises ; 3 \% restituées.}
\nf{Références}{ICUR I 2109.}
\nf{Provenance}{Rattachement selon l’EDB : \emph{Incertae originis} (le volume I des ICVR réunit les inscriptions dont la provenance est perdue). Conservation selon l’EDB : \emph{Musei Vaticani: Galleria Lapidaria}.}
\nf{Support}{Plaque de marbre. 39 × 58 cm (h. × l.). Inscription gravée. Fiche EDB : Filippo Piazzolla, première version 2015.}
\ntext{\lnum{1}Bene meren[ti] in Pace Satyro\\
\lnum{2}qui vixit annos\\
\lnum{3}p(lus) m(inus) XLV Caesario\\
\lnum{4}et Attico co(n)ss(ulibus)}
\nf{Témoins}{EDB (EDB32106) : texte concordant aux conventions de notation près (indice de similarité 0,974).}
\nf{Traduction}{Au bien méritant, en paix, Satyrus, qui vécut environ 45 ans ; sous le consulat de Caesarius et d’Atticus.}
\nf{Langue}{4 abréviations développées par l’éditeur (dont p(lus), m(inus), co(n)ss, (ulibus)).}
\nf{Chronologie}{Noms consulaires : Caesarius (lu) et Atticus (lu). Âge déclaré : env. 45 a. (lu).}
\nf{Commentaire}{La formule in pace est conservée.}
\nf{Bibliographie}{ICUR I 2109. Bibliographie complémentaire selon l’EDB : ILCV, 2609A adn.}
\nf{Statut}{Lecture secondaire (EDCS) concordante avec l’EDB : lecture stabilisée entre deux bases, non collationnée sur l’ICVR ni sur le monument.}
\end{notice}

\begin{notice}{404}{397}{ICUR n.s. I, 2772 · EDCS-39900321 · EDB32701}
\label{not:404}
\nf{Datation}{397 — directe : nom consulaire lu}
\nf{Genre}{Épitaphe ; ensembles A, B, C ; 92 lettres transmises.}
\nf{Références}{ICUR I 2772 = ILCV 3099b.}
\nf{Provenance}{Rattachement selon l’EDB : \emph{Incertae originis} (le volume I des ICVR réunit les inscriptions dont la provenance est perdue). Conservation selon l’EDB : \emph{Fossombrone (PU)   ; Lapidario Diocesano}.}
\nf{Support}{Plaque de marbre. 14,7 × 22,2 cm (h. × l.). Inscription gravée. Reproductions signalées par l’EDB : reproduction (droits réservés ; non reproduites). Fiche EDB : Filippo Piazzolla, première version 2018.}
\ntext{\lnum{1}Fl(avio) C(a)esario et Noni\\
\lnum{2}o Atti(c)o conss(ulibus) depo\\
\lnum{3}sita Alexandra III\\
\lnum{4}Idus Novembres\\
\lnum{5}ann(orum) V mens(i)um\\
\lnum{6}VIII die⸢r⸣(um) V quie\{n\}s\\
\lnum{7}⸢c⸣it in pace}
\nf{Témoins}{Export EDCS de 2020 : \emph{Fl(avio) C(a)esario et Noni/o Atti(c)o conss(ulibus) depo/sita Alexandra III / Idus Novembres / ann(orum) V mens(i)um / VIII die(um) V quie\{n\}s/it in pace}. EDB (EDB32701) : texte concordant aux conventions de notation près (indice de similarité 0,937).}
\nf{Apparat}{⸢r⸣ : S gravé ; ⸢c⸣ : QV gravé ; \{n\} : lettres superflues}
\nf{Traduction}{Sous le consulat de Flavius Caesarius et de Nonius Atticus, déposée Alexandra le 3e jour avant les ides de novembre, âgée de 5 ans 8 mois 5 jours ; elle repose en paix.}
\nf{Langue}{L’éditeur signale 2 graphies corrigées : S gravé pour r ; QV gravé pour c. Des lettres gravées sont jugées superflues par l’éditeur (entre accolades). 7 abréviations développées par l’éditeur (dont Fl(avio), C(a)esario, Atti(c)o, conss(ulibus), ann(orum), mens(i)um).}
\nf{Chronologie}{Noms consulaires : Caesarius (lu) et Atticus (lu). Dates du jour : déposition : Idus Novembres = 13 novembre 397 (lu). Âge déclaré : 5 a. (lu).}
\nf{Commentaire}{La formule in pace est conservée. Vocabulaire religieux : repos et sommeil de la mort — quiescit (lu).}
\nf{Bibliographie}{ICUR I 2772 ; ILCV 3099b. Bibliographie complémentaire selon l’EDB : ILCV, 3099B ; Billanovich 1967, 56-58.}
\nf{Statut}{Lecture secondaire (EDCS) concordante avec l’EDB : lecture stabilisée entre deux bases, non collationnée sur l’ICVR ni sur le monument. L’EDCS signale une photographie (non consultée).}
\end{notice}

\begin{notice}{405}{397}{ICUR n.s. I, 2812 · EDCS-11700578 · EDB32736}
\label{not:405}
\nf{Datation}{397 — directe : nom consulaire lu}
\nf{Genre}{Épitaphe ; ensembles A, B, C ; 82 lettres transmises.}
\nf{Références}{ICUR I 2812 = ILCV 1508 = ILMN-01, 00447.}
\nf{Provenance}{Rattachement selon l’EDB : \emph{Incertae originis} (le volume I des ICVR réunit les inscriptions dont la provenance est perdue). Conservation selon l’EDB : \emph{Napoli   ; Museo Archeologico Nazionale}.}
\nf{Support}{Plaque de marbre. 23,5 × 36 × 4 cm (h. × l. × ép.). Inscription gravée. Reproductions signalées par l’EDB : reproduction, Camodeca et alii 2000 (droits réservés ; non reproduites). Fiche EDB : Filippo Piazzolla, première version 2015.}
\ntext{\lnum{1}Bonifae bene qu(i)\{a\}es\\
\lnum{2}⸢c⸣enti catec(h)umi\\
\lnum{3}no qui vixit annos\\
\lnum{4}p(lus) m(inus) LX d(epositus) IIII Kal(endas)\\
\lnum{5}Iulias Caesario et Attico\\
\lnum{6}conss(ulibus)}
\nf{Témoins}{Export EDCS de 2020 : \emph{Bonifae bene qu(i)\{a\}es/enti catec(h)umi/no qui vixit annos / p(lus) m(inus) LX d(epositus) IIII Kal(endas) / Iulias Caesario et Attico / conss(ulibus)}. EDB (EDB32736 ; écarts à examiner) : \emph{Bonifae bene quaes/quenti catecumi/no qui vixit annos / p(lus) m(inus) LX[[---]] ((folium)) d(epositus) 〚IV ---〛 IIII ka⌜l⌝(endas) / iulias Caesario et Attico / cons(ulibu)s [texte tronqué dans la liste de résultats]}.}
\nf{Apparat}{⸢c⸣ : QV gravé ; \{a\} : lettres superflues}
\nf{Traduction}{À Bonifatius (?), qui repose bien, catéchumène, qui vécut environ 60 ans ; déposé le 4e jour avant les calendes de juillet, sous le consulat de Caesarius et d’Atticus.}
\nf{Langue}{L’éditeur signale une graphie corrigée : QV gravé pour c. Des lettres gravées sont jugées superflues par l’éditeur (entre accolades). 7 abréviations développées par l’éditeur (dont qu(i), catec(h)umi, p(lus), m(inus), d(epositus), Kal(endas)).}
\nf{Chronologie}{Noms consulaires : Caesarius (lu) et Atticus (lu). Dates du jour : déposition : IIII Kalendas Iulias = 28 juin 397 (lu). Âge déclaré : env. 60 a. (lu).}
\nf{Commentaire}{Vocabulaire religieux : initiation et foi — catechumino (lu) ; repos et sommeil de la mort — quiescenti (lu). Le mot catec(h)umino qualifie un statut d’initiation. Il n’est ni une occurrence de Deus ni une désignation du Christ ; son classement doit rester distinct des noms divins. Commentaire développé : dossier p.~\pageref{dos:405}.}
\nf{Bibliographie}{ICUR I 2812 ; ILCV 1508 ; ILMN-01, 00447. Bibliographie complémentaire selon l’EDB : ILCV, 1508 ; Camodeca et alii 2000, 153, n. 447 con foto.}
\nf{Statut}{Lecture secondaire (EDCS) ; écarts avec l’EDB à arbitrer sur l’ICVR. L’EDCS signale une photographie (non consultée).}
\end{notice}

\begin{notice}{406}{397}{ICUR n.s. II, 4830 · EDCS-36101310}
\label{not:406}
\nf{Datation}{397 — directe : nom consulaire lu}
\nf{Genre}{Mixte (épitaphe et acte d’achat) ; ensembles A, B, C ; 75 lettres transmises.}
\nf{Références}{ICUR II 4830 = ILCV 3747. C0 conserve 468 textes pour la comparaison historique. C1 inclut la notice 406 : 469 textes, dont 196 avec in pace (41,8 \%). Le statut final et l’événement daté exigent un témoin = aucun mot transmis n’est modifié.}
\nf{Provenance}{Rome (indication de l’EDCS) ; aucun rattachement cimétérial transmis.}
\nf{Support}{Non documenté dans les sources consultées.}
\ntext{\lnum{1}Sussanna compara\\
\lnum{2}vit sibi memoriam quie\\
\lnum{3}vit die VII Kal(endas) Augusta⸢s⸣\\
\lnum{4}conss(ulibus) Caesario et Attico}
\nf{Témoins}{Export EDCS de 2020 : \emph{Sussanna compara/vit sibi memoriam quie/vit die VII Kal(endas) Augusta / conss(ulibus) Caesario et Attico}. Aucune fiche EDB appariée (numéro ICUR absent de la liste publique de l’EDB).}
\nf{Apparat}{⸢s⸣ : R gravé}
\nf{Traduction}{Susanna s’est acquis une memoria ; elle a reposé le 7e jour avant les calendes d’août, sous le consulat de Caesarius et d’Atticus.}
\nf{Langue}{L’éditeur signale une graphie corrigée : R gravé pour s.}
\nf{Chronologie}{Noms consulaires : Caesarius (lu) et Atticus (lu). Le verbe \emph{quievit} précède la date (26 juillet 397) : celle-ci se rapporte au repos, c'est-à-dire au décès, de Susanna, non à l'achat. Le texte réunit une acquisition (\emph{comparavit sibi memoriam}) et la mention d'un décès ; il est classé comme texte mixte et inclus dans l'ensemble funéraire. Les mesures sont aussi données sans lui (chapitre sur les dénominateurs). Dates du jour : achat : die VII Kalendas Augustas = 26 juillet 397 (lu).}
\nf{Bibliographie}{ICUR II 4830 ; ILCV 3747. C0 conserve 468 textes pour la comparaison historique. C1 inclut la notice 406 : 469 textes, dont 196 avec in pace (41,8 \%). Le statut final et l’événement daté exigent un témoin ; aucun mot transmis n’est modifié.}
\nf{Statut}{Lecture secondaire (EDCS) sans second témoin numérique. L’EDCS signale une photographie (non consultée).}
\end{notice}

\begin{notice}{407}{397}{ICUR n.s. II, 4832 · EDCS-36101312}
\label{not:407}
\nf{Datation}{397 — directe : nom consulaire lu}
\nf{Genre}{Épitaphe ; ensembles A, B, C ; 51 lettres transmises ; 51 \% restituées.}
\nf{Références}{ICUR II 4832.}
\nf{Provenance}{Rome (indication de l’EDCS) ; aucun rattachement cimétérial transmis.}
\nf{Support}{Non documenté dans les sources consultées.}
\ntext{\lnum{1}S[- - -]\\
\lnum{2}Benata[- - -]\\
\lnum{3}se vi⸢v⸣i lo[cum - - -]\\
\lnum{4}Fl(avio) Caesario [et Nonio Attico cons(ulibus)]}
\nf{Témoins}{Export EDCS de 2020 : \emph{S[3] / Benata[3] / se vii lo[cum 3] / Fl(avio) Caesario [et Nonio Attico cons(ulibus)]}. Aucune fiche EDB appariée (numéro ICUR absent de la liste publique de l’EDB).}
\nf{Apparat}{⸢v⸣ : B gravé}
\nf{Traduction}{[…] Benata […] de leur vivant un emplacement […] ; sous le consulat de Flavius Caesarius et de Nonius Atticus.}
\nf{Langue}{L’éditeur signale une graphie corrigée : B gravé pour v (bêtacisme).}
\nf{Chronologie}{Noms consulaires : Caesarius (lu) et Atticus (restitué).}
\nf{Bibliographie}{ICUR II 4832.}
\nf{Statut}{Lecture secondaire (EDCS) sans second témoin numérique. Plus de la moitié des lettres sont restituées.}
\end{notice}

\begin{notice}{408}{397}{ICUR n.s. II, 6055 · EDCS-34200681 · EDB16486}
\label{not:408}
\nf{Datation}{397 — directe : nom consulaire lu}
\nf{Genre}{Épitaphe ; ensembles A, B, C ; 56 lettres transmises ; 21 \% restituées.}
\nf{Références}{ICUR II 6055 = AE 1905, 82.}
\nf{Provenance}{Rattachement selon l’EDB : \emph{Coem. Commodillae}. Conservation selon l’EDB : \emph{Coem. Commodillae}.}
\nf{Support}{Plaque de marbre. Inscription gravée. Fiche EDB : Carlo Carletti, première version 2003.}
\ntext{\lnum{1}Fl(avio) C(a)esar(i)o [et A]ttico vv(iris) (!) consu[libus - - -]\\
\lnum{2}[i]n pace [- - -]X Kal(endas) No⸢v⸣(embres) qu(a)e [- - -]\\
\lnum{3}[- - -]s depo[sit]us II[}
\nf{Témoins}{Export EDCS de 2020 : \emph{Fl(avio) C(a)esar(i)o [et A]ttico vv(iris) (!) consu[libus 3] / [i]n pace [3]X Kal(endas) No(embres) qu(a)e [3] / [3]s depo[sit]us II[}. EDB (EDB16486 ; écarts à examiner) : \emph{Fl(avio) Cesaro [et A]⌜tt⌝ico vv(iris) ⟨cc(larissimis)⟩ conṣ(ulibu)ṣ [---] / [i]n pace [---] X kal(endas) nob(embres) que [---] / [---] s depo[sit]us I+[---]}.}
\nf{Apparat}{⸢v⸣ : B gravé ; (!) : forme signalée par l’éditeur}
\nf{Traduction}{Sous le consulat de Flavius Caesarius et d’Atticus […] ; en paix, le 1x jour avant les calendes de novembre, qui […] ; […] déposé 2 (?).}
\nf{Langue}{L’éditeur signale une graphie corrigée : B gravé pour v (bêtacisme). L’éditeur signale au moins une forme anormale (sic). 7 abréviations développées par l’éditeur (dont Fl(avio), C(a)esar, (i)o, vv(iris), Kal(endas), (embres)).}
\nf{Chronologie}{Noms consulaires : Caesarius (lu) et Atticus (lu en partie). Dates du jour : événement indéterminé : …X Kalendas Novembres = 23 octobre 397 (partiellement restitué).}
\nf{Commentaire}{La formule in pace est attestée seulement par restitution.}
\nf{Bibliographie}{ICUR II 6055 ; AE 1905, 82. Bibliographie complémentaire selon l’EDB : CLRE, 659 ; ICVR III, 6438d ; Fratticci 2016, 51 ; 281 nr. 10.}
\nf{Statut}{Lecture secondaire (EDCS) ; écarts avec l’EDB à arbitrer sur l’ICVR. L’EDCS signale une photographie (non consultée).}
\end{notice}

\begin{notice}{409}{397}{ICUR n.s. II, 6056 · EDCS-34200682 · EDB16487}
\label{not:409}
\nf{Datation}{397 — directe : nom consulaire lu}
\nf{Genre}{Épitaphe ; ensembles A, B, C ; 159 lettres transmises.}
\nf{Références}{ICUR II 6056 = ILCV 4988 (em).}
\nf{Provenance}{Rattachement selon l’EDB : \emph{Coem. Commodillae}.}
\nf{Support}{Plaque de marbre. Inscription gravée. Fiche EDB : Carlo Carletti, première version 2003.}
\ntext{\lnum{1}Reveccae innocenti quae vixit\\
\lnum{2}annum unum mensem unum\\
\lnum{3}dies XVII bene merenti in pace\\
\lnum{4}deposita VIIII Kal(endas) Septembres\\
\lnum{5}Flavio Caesario et Nonio\\
\lnum{6}Attico vv(iris) cc(larissimis) cons(ulibu)s\\
\lnum{7}parentes bene merenti fecerunt}
\nf{Témoins}{EDB (EDB16487) : texte concordant aux conventions de notation près (indice de similarité 0,995).}
\nf{Traduction}{À Rebecca, innocente, qui vécut un an un mois 17 jours ; bien méritante, en paix ; déposée le 9e jour avant les calendes de septembre, sous le consulat de Flavius Caesarius et de Nonius Atticus, clarissimes ; ses parents ont fait (ce monument) pour elle, bien méritante.}
\nf{Langue}{4 abréviations développées par l’éditeur (dont Kal(endas), vv(iris), cc(larissimis), cons(ulibu)s).}
\nf{Chronologie}{Noms consulaires : Caesarius (lu) et Atticus (lu). Dates du jour : déposition : VIIII Kalendas Septembres = 24 août 397 (lu). Âge déclaré : 1 a. 1 m. 17 j. (lu).}
\nf{Commentaire}{La formule in pace est conservée.}
\nf{Bibliographie}{ICUR II 6056 ; ILCV 4988 (em).}
\nf{Statut}{Lecture secondaire (EDCS) concordante avec l’EDB : lecture stabilisée entre deux bases, non collationnée sur l’ICVR ni sur le monument. L’EDCS signale une photographie (non consultée).}
\end{notice}

\begin{notice}{410}{397}{ICUR n.s. II, 6057 · EDCS-34200683 · EDB16488}
\label{not:410}
\nf{Datation}{397 — directe : nom consulaire lu}
\nf{Genre}{Épitaphe ; ensembles A, B, C ; 69 lettres transmises ; 48 \% restituées.}
\nf{Références}{ICUR II 6057.}
\nf{Provenance}{Rattachement selon l’EDB : \emph{Coem. Commodillae}.}
\nf{Support}{Plaque de marbre. Inscription gravée. Fiche EDB : Carlo Carletti, première version 2003.}
\ntext{\lnum{1}[Bene merenti i]n pace Hilari\\
\lnum{2}[no qui vixit a]nn(os) V m(enses) V depo\\
\lnum{3}[situs - - - Caesa]rio et Attico cons(ulibu)s}
\nf{Témoins}{EDB (EDB16488) : texte concordant aux conventions de notation près (indice de similarité 0,983).}
\nf{Traduction}{Au bien méritant, en paix, Hilarinus, qui vécut 5 ans 5 mois ; déposé […] sous le consulat de Caesarius et d’Atticus.}
\nf{Langue}{3 abréviations développées par l’éditeur (dont ann(os), m(enses), cons(ulibu)s).}
\nf{Chronologie}{Noms consulaires : Caesarius (lu en partie) et Atticus (lu). Âge déclaré : 5 a. 5 m. (partiellement restitué).}
\nf{Commentaire}{La formule in pace est attestée seulement par restitution.}
\nf{Bibliographie}{ICUR II 6057.}
\nf{Statut}{Lecture secondaire (EDCS) concordante avec l’EDB : lecture stabilisée entre deux bases, non collationnée sur l’ICVR ni sur le monument.}
\end{notice}

\begin{notice}{411}{397}{ICUR n.s. II, 6058 · EDCS-34200684 · EDB16489}
\label{not:411}
\nf{Datation}{397 — directe : nom consulaire lu}
\nf{Genre}{Mixte (épitaphe et acte d’achat) ; ensembles A, B, C ; 90 lettres transmises.}
\nf{Références}{ICUR II 6058 = ILCV 3781 = Epigraphica-1985-109.}
\nf{Provenance}{Rattachement selon l’EDB : \emph{Coem. Commodillae}.}
\nf{Support}{Plaque de marbre. Inscription gravée. Fiche EDB : Carlo Carletti, première version 2003.}
\ntext{\lnum{1}Deposita puella nomine\\
\lnum{2}Ursa in pace annorum\\
\lnum{3}viginti duo conss(ulatu) Fl(avi) Caesari et Maximi\\
\lnum{4}vv(irorum) cc(larissimorum) emptu(m) locu(m) (t)r⸢e⸣s et s(olidis)\\
\lnum{5}XI}
\nf{Témoins}{Export EDCS de 2020 : \emph{Deposita puella nomine / Ursa in pace annorum / viginti duo conss(ulatu) Fl(avi) Caesari et Maximi / vv(irorum) cc(larissimorum) emptu(m) locu(m) (t)rs et s(olidis) / XI //}. EDB (EDB16489) : texte concordant aux conventions de notation près (indice de similarité 0,973).}
\nf{Apparat}{⸢e⸣ : I gravé}
\nf{Traduction}{Déposée, la jeune fille nommée Ursa, en paix, âgée de vingt-deux ans ; sous le consulat de Flavius Caesarius et de Maximus (sic), clarissimes ; emplacement acheté pour trois (personnes ?) et pour 11 sous (solidi).}
\nf{Langue}{L’éditeur signale une graphie corrigée : I gravé pour e (confusion e/i). L’export de 2020 diffère de l’état 2022 (voir apparat). 8 abréviations développées par l’éditeur (dont conss(ulatu), Fl(avi), vv(irorum), cc(larissimorum), emptu(m), locu(m)).}
\nf{Chronologie}{Noms consulaires : Caesarius (lu). Le collège transmis, \emph{Fl(avi) Caesari et Maximi}, ne correspond pas aux fastes de 397 (Caesarius et Atticus). L'année de classement est conservée ; le second nom consulaire doit être collationné. Réserve : Collège Caesarius et Maximus ; les fastes donnent pour 397 Caesarius et Atticus. Âge déclaré : 22 a. (lu).}
\nf{Commentaire}{La formule in pace est conservée. Le groupe relatif au prix doit conserver sa graphie et sa syntaxe transmises. La mention de solidi documente une transaction, mais le nombre de places acquises et la construction exacte du prix demandent une relecture ; aucun niveau social du ménage n’en est déduit. L'âge transmis est \emph{annorum viginti duo} : vingt-deux ans, et non vingt comme le portait la table des personnes antérieure ; la valeur est corrigée dans les calculs. Commentaire développé : dossier p.~\pageref{dos:411}.}
\nf{Bibliographie}{ICUR II 6058 ; ILCV 3781 ; Epigraphica-1985-109.}
\nf{Statut}{Lecture secondaire (EDCS) concordante avec l’EDB : lecture stabilisée entre deux bases, non collationnée sur l’ICVR ni sur le monument.}
\end{notice}

\begin{notice}{412}{397}{ICUR n.s. IV, 9581 · EDCS-36101106 · EDB20443}
\label{not:412}
\nf{Datation}{397 — directe : lecture partielle d’un nom consulaire}
\nf{Genre}{Épitaphe ; ensembles A, B, C ; 96 lettres transmises ; 58 \% restituées.}
\nf{Références}{ICUR IV 9581.}
\nf{Provenance}{Rattachement selon l’EDB : \emph{Coem. Callisti pars inferior}.}
\nf{Support}{Plaque de marbre. Inscription gravée. Fiche EDB : Antonio Enrico Felle, première version 2003.}
\ntext{\lnum{1}Victorio filio [- - -]\\
\lnum{2}qui vixit ann(os) II[- - -]\\
\lnum{3}dep(ositus) VIIII K(a)l(endas) Mai(as) c[onss(ulibus) Fl(avio)\\
\lnum{4}Caess[ario et Nonio Attico vv(iris) cc(larissimis)}
\nf{Témoins}{EDB (EDB20443) : texte concordant aux conventions de notation près (indice de similarité 0,947).}
\nf{Traduction}{À Victorius, son fils […] qui vécut 2 (?) ans […] ; déposé le 9e jour avant les calendes de mai, sous le consulat de Flavius Caesarius et de Nonius Atticus, clarissimes.}
\nf{Langue}{9 abréviations développées par l’éditeur (dont ann(os), dep(ositus), K(a)l, (endas), Mai(as), conss(ulibus)).}
\nf{Chronologie}{Noms consulaires : Atticus (restitué) et Victor (lu). Dates du jour : déposition : VIIII Kalendas Maias = 23 avril 397 (lu).}
\nf{Bibliographie}{ICUR IV 9581.}
\nf{Statut}{Lecture secondaire (EDCS) concordante avec l’EDB : lecture stabilisée entre deux bases, non collationnée sur l’ICVR ni sur le monument. Plus de la moitié des lettres sont restituées.}
\end{notice}

\begin{notice}{413}{397}{ICUR n.s. IV, 9582 · EDCS-36101107 · EDB19321}
\label{not:413}
\nf{Datation}{397 — directe : nom consulaire lu}
\nf{Genre}{Épitaphe ; ensembles A, B, C ; 33 lettres transmises ; 64 \% restituées.}
\nf{Références}{ICUR IV 9582.}
\nf{Provenance}{Rattachement selon l’EDB : \emph{Coem. Callisti pars inferior}.}
\nf{Support}{Plaque de marbre. Inscription gravée. Fiche EDB : Antonio Enrico Felle, première version 2003.}
\ntext{\lnum{1}[- - -]a dep(osita) re[- - -]\\
\lnum{2}[Caesario e]t Attic[o conss(ulibus)]}
\nf{Témoins}{EDB (EDB19321 ; écarts à examiner) : \emph{1[---]a dep(osita) re[---] / [Caesario e]t Attic[o conss(ulibus)] [texte tronqué dans la liste de résultats]}.}
\nf{Traduction}{[…]a déposée […] ; sous le consulat de Caesarius et d’Atticus.}
\nf{Langue}{2 abréviations développées par l’éditeur (dont dep(osita), conss(ulibus)).}
\nf{Chronologie}{Noms consulaires : Caesarius (restitué) et Atticus (lu).}
\nf{Bibliographie}{ICUR IV 9582.}
\nf{Statut}{Lecture secondaire (EDCS) ; écarts avec l’EDB à arbitrer sur l’ICVR. Plus de la moitié des lettres sont restituées.}
\end{notice}

\begin{notice}{414}{397}{ICUR n.s. IV, 11128 · EDCS-37200415 · EDB39592}
\label{not:414}
\nf{Datation}{397 — directe : lecture partielle d’un nom consulaire}
\nf{Genre}{Épitaphe ; ensembles A, B, C ; 50 lettres transmises ; 46 \% restituées.}
\nf{Références}{ICUR IV 11128 = ILCV +4193.}
\nf{Provenance}{Rattachement selon l’EDB : \emph{Coem. Callisti area subdialis et incerti loci}. Conservation selon l’EDB : \emph{Coem. Callisti area subdialis et incerti loci}.}
\nf{Support}{Plaque de marbre. 22 × 26 cm (h. × l.). Inscription gravée. Fiche EDB : Marida Pierno, première version 2015.}
\ntext{\lnum{1}]r dolor es[- - -]\\
\lnum{2}[- - - bis]omu(m ?) acerbu(m ?) dem[- - -]\\
\lnum{3}[- - - Cae]sario et [Attico conss(ulibus)]}
\nf{Témoins}{EDB (EDB39592 ; écarts à examiner) : \emph{[---]r dolor es[---] / [---]omu acerbu dem[---] / [--- Cae]sario et [Attico ---]}.}
\nf{Traduction}{[…] douleur […] ; […] tombe double amère […] ; sous le consulat de Caesarius et d’Atticus.}
\nf{Langue}{3 abréviations développées par l’éditeur (dont bisomu(m ?), acerbu(m ?), conss(ulibus)).}
\nf{Chronologie}{Noms consulaires : Caesarius (lu en partie) et Atticus (restitué).}
\nf{Bibliographie}{ICUR IV 11128 ; ILCV +4193.}
\nf{Statut}{Lecture secondaire (EDCS) ; écarts avec l’EDB à arbitrer sur l’ICVR.}
\end{notice}

\begin{notice}{415}{397}{ICUR n.s. IV, 11778 · EDCS-37400866 · EDB5167}
\label{not:415}
\nf{Datation}{397 — directe : consulat restitué}
\nf{Genre}{Épitaphe ; ensembles A, B, C ; 37 lettres transmises ; 78 \% restituées.}
\nf{Références}{ICUR IV 11778.}
\nf{Provenance}{Rattachement selon l’EDB : \emph{Coem. anonymum ad viam Appiam}. Conservation selon l’EDB : \emph{Coem. anonymum ad viam Appiam}.}
\nf{Support}{Plaque de marbre. Inscription gravée. Fiche EDB : Maria Cristina Grisanzio, première version 2003.}
\ntext{\lnum{1}M]ai(i)s Fl(avio) C(a)e[sario et Nonio Attico conss(ulibus)}
\nf{Témoins}{EDB (EDB5167 ; écarts à examiner) : \emph{[--- m]ais Fl(avio) Ce[sario et Nonio Attico cons(ulibu)s] [texte tronqué dans la liste de résultats]}.}
\nf{Traduction}{[…] mai, sous le consulat de Flavius Caesarius et de Nonius Atticus.}
\nf{Langue}{4 abréviations développées par l’éditeur (dont Mai(i)s, Fl(avio), C(a)e, conss(ulibus)).}
\nf{Chronologie}{Noms consulaires : Caesarius (restitué) et Atticus (restitué).}
\nf{Bibliographie}{ICUR IV 11778.}
\nf{Statut}{Lecture secondaire (EDCS) ; écarts avec l’EDB à arbitrer sur l’ICVR. Plus de la moitié des lettres sont restituées.}
\end{notice}

\begin{notice}{416}{397}{ICUR n.s. IV, 11779 · EDCS-37400867 · EDB5168}
\label{not:416}
\nf{Datation}{397 — directe : nom consulaire lu}
\nf{Genre}{Épitaphe ; ensembles A, B, C ; 137 lettres transmises ; 15 \% restituées.}
\nf{Références}{ICUR IV 11779 = ILCV 4400a = RAC-1925-14.}
\nf{Provenance}{Rattachement selon l’EDB : \emph{Coem. anonymum ad viam Appiam}. Conservation selon l’EDB : \emph{Coem. anonymum ad viam Appiam}.}
\nf{Support}{Plaque de marbre. Inscription gravée. Reproductions signalées par l’EDB : photographie PCAS (droits réservés ; non reproduites). Fiche EDB : Maria Cristina Grisanzio, première version 2003.}
\ntext{\lnum{1}[- - - defu]ncta puella\\
\lnum{2}[- - - nom]ine Petronia\\
\lnum{3}[ann]orum XIII die⸢rum⸣\\
\lnum{4}[- - - de]posita post ter\\
\lnum{5}[tiu]m Kalendas Iulia\\
\lnum{6}[s di]e ⸢V⸣eneris bene me\\
\lnum{7}[re]nti in pace conss(ulibus)\\
\lnum{8}[F]l(avio) Caesario et Nonio Atti\\
\lnum{9}co vv(iris) cc(larissimis)}
\nf{Témoins}{Export EDCS de 2020 : \emph{[3 defu]ncta puella / [3 nom]ine Petronia / [ann]orum XIII die / [3 de]posita post ter/[tiu]m Kalendas Iulia/[s di]e eneris bene me/[re]nti in pace conss(ulibus) / [F]l(avio) Caesario et Nonio Atti/co vv(iris) cc(larissimis)}. EDB (EDB5168) : texte concordant aux conventions de notation près (indice de similarité 0,929).}
\nf{Apparat}{⸢rum⸣ : S gravé ; ⸢V⸣ : B gravé}
\nf{Traduction}{[…] morte, la jeune fille […] nommée Petronia, âgée de 13 ans […] jours ; déposée après le 3e jour avant les calendes de juillet (?), un vendredi ; bien méritante, en paix ; sous le consulat de Flavius Caesarius et de Nonius Atticus, clarissimes.}
\nf{Langue}{L’éditeur signale 2 graphies corrigées : S gravé pour rum ; B gravé pour V (bêtacisme).}
\nf{Chronologie}{Noms consulaires : Caesarius (lu) et Atticus (lu). Dates du jour : déposition : tertium Kalendas Iulias = 29 juin 397 (partiellement restitué). Âge déclaré : 13 a. (partiellement restitué).}
\nf{Commentaire}{La formule in pace est conservée.}
\nf{Bibliographie}{ICUR IV 11779 ; ILCV 4400a ; RAC-1925-14.}
\nf{Statut}{Lecture secondaire (EDCS) concordante avec l’EDB : lecture stabilisée entre deux bases, non collationnée sur l’ICVR ni sur le monument. L’EDCS signale une photographie (non consultée).}
\end{notice}

\annee{398}

\begin{notice}{417}{398}{EDCS-08400008 · EDB38747}
\label{not:417}
\nf{Datation}{398 — directe : nom consulaire lu}
\nf{Genre}{Acte d’achat ou de réservation ; ensembles A, B ; 98 lettres transmises ; 21 \% restituées.}
\nf{Références}{AE 1984, 51.}
\nf{Provenance}{Rattachement selon l’EDB : \emph{Ipogeo della Torretta}. Conservation selon l’EDB : \emph{Ipogeo della Torretta}.}
\nf{Support}{Plaque de marbre. 62 × 133 × 2,4 cm (h. × l. × ép.). Inscription gravée. Fiche EDB : Anita Rocco, première version 2017.}
\ntext{\lnum{1}Dom[inis nos]tris Ho(no)rio quater et Eut⸢y⸣c(h)iano consulibus\\
\lnum{2}loc[um Felici]tati(s) et Crescentiaes(!) [- - -] em[erunt s]e ⸢v⸣i⸢v⸣i sibi et suis}
\nf{Témoins}{Export EDCS de 2020 : \emph{Dom[inis nos]tris Ho(no)rio quater et Euticiano consulibus / loc[um Felici]tati(s) et Crescentiaes(!) [3] em[erunt s]e ii sibi et suis}. EDB (EDB38747) : texte concordant aux conventions de notation près (indice de similarité 0,967).}
\nf{Apparat}{⸢y⸣ : I gravé ; ⸢v⸣ : B gravé ; ⸢v⸣ : B gravé ; (!) : forme signalée par l’éditeur}
\nf{Traduction}{Sous le consulat de nos seigneurs Honorius pour la 4e fois et Eutychianus, Felicitas et Crescentia ont acheté de leur vivant un emplacement pour eux et les leurs.}
\nf{Langue}{L’éditeur signale 3 graphies corrigées : I gravé pour y (graphie de y) ; B gravé pour v (bêtacisme) ×2. L’export de 2020 diffère de l’état 2022 (voir apparat). L’éditeur signale au moins une forme anormale (sic). 4 abréviations développées par l’éditeur (dont Ho(no)rio, c(h)iano, Felicitati(s), Crescentiaes(!)).}
\nf{Chronologie}{Noms consulaires : Honorius (lu) et Eutychianus (lu).}
\nf{Bibliographie}{AE 1984, 51. Bibliographie complémentaire selon l’EDB : Fasola 1984b, 30.}
\nf{Statut}{Lecture secondaire (EDCS) concordante avec l’EDB : lecture stabilisée entre deux bases, non collationnée sur l’ICVR ni sur le monument.}
\end{notice}

\begin{notice}{418}{398}{ICUR n.s. I, 1452 · EDCS-39000431 · EDB31242}
\label{not:418}
\nf{Datation}{398 — directe : nom consulaire lu}
\nf{Genre}{Épitaphe ; ensembles A, B, C ; 106 lettres transmises ; 20 \% restituées.}
\nf{Références}{ICUR I 1452 = ILCV +1128.}
\nf{Provenance}{Rattachement selon l’EDB : \emph{Incertae originis} (le volume I des ICVR réunit les inscriptions dont la provenance est perdue). Conservation selon l’EDB : \emph{Musei Vaticani: Lapidario Cristiano}.}
\nf{Support}{Plaque de marbre. 41 × 117 cm (h. × l.). Inscription gravée. Fiche EDB : Filippo Piazzolla, première version 2015.}
\ntext{\lnum{1}[Locus Ti]tiani et Epinicenis\\
\lnum{2}[Honorio] IIII et Fl(avio) Eut⸢y⸣c(h)iano vv(iris) cc(larissimis) co(n)s(ulibus)\\
\lnum{3}[- - -] XVI Kal(endas) Octob(res) depositus est\\
\lnum{4}[in somno] pacis ⸢q⸣ui vixit ann(os) p(lus) m(inus) LXV}
\nf{Témoins}{Export EDCS de 2020 : \emph{[Locus Ti]tiani et Epinicenis / [Honorio] IIII et Fl(avio) Euticiano vv(iris) cc(larissimis) co(n)s(ulibus) / [3] XVI Kal(endas) Octob(res) depositus est / [in somno] pacis ui vixit ann(os) p(lus) m(inus) LXV}. EDB (EDB31242) : texte concordant aux conventions de notation près (indice de similarité 0,989).}
\nf{Apparat}{⸢y⸣ : I gravé ; ⸢q⸣ : O gravé}
\nf{Traduction}{Emplacement de Titianus et d’Epinicè ; sous le consulat d’Honorius pour la 4e fois et de Flavius Eutychianus, clarissimes ; le 16e jour avant les calendes d’octobre il a été déposé dans le sommeil de la paix ; il vécut environ 65 ans.}
\nf{Langue}{L’éditeur signale 2 graphies corrigées : I gravé pour y (graphie de y) ; O gravé pour q (confusion de lettres). L’export de 2020 diffère de l’état 2022 (voir apparat). 11 abréviations développées par l’éditeur (dont Fl(avio), c(h)iano, vv(iris), cc(larissimis), co(n)s, (ulibus)).}
\nf{Chronologie}{Noms consulaires : Honorius (restitué) et Eutychianus (lu). Dates du jour : événement indéterminé : … XVI Kalendas Octobres = 16 septembre 398 (partiellement restitué). Âge déclaré : env. 65 a. (lu).}
\nf{Bibliographie}{ICUR I 1452 ; ILCV +1128. Bibliographie complémentaire selon l’EDB : Lega 2000, 284 ; Ferrua 1979b, 38, n. 1452 ; Ferrua 1964, 109.}
\nf{Statut}{Lecture secondaire (EDCS) concordante avec l’EDB : lecture stabilisée entre deux bases, non collationnée sur l’ICVR ni sur le monument. L’EDCS signale une photographie (non consultée).}
\end{notice}

\begin{notice}{419}{398}{ICUR n.s. I, 1941 · EDCS-39200385 · EDB31859}
\label{not:419}
\nf{Datation}{398 — directe : nom consulaire lu}
\nf{Genre}{Épitaphe ; ensembles A, B, C ; 103 lettres transmises.}
\nf{Références}{ICUR I 1941 = ILCV 2570.}
\nf{Provenance}{Rattachement selon l’EDB : \emph{Incertae originis} (le volume I des ICVR réunit les inscriptions dont la provenance est perdue). Conservation selon l’EDB : \emph{Museo Nazionale Romano}.}
\nf{Support}{Plaque de marbre. 69 × 78 × 3,8 cm (h. × l. × ép.). Inscription gravée. Fiche EDB : Filippo Piazzolla, première version 2015.}
\ntext{\lnum{1}Maxima in pac\{a\}e q(uae) vixit an\\
\lnum{2}n⸢o⸣s plus minus XXXV conss(ulibus)\\
\lnum{3}dominis nostris (H)onorio IIII c(on)s\{s\}(ule) et Eu\\
\lnum{4}t⸢y⸣c(h)iano cons(ule) pridi\{a\}e Nonas Septe⸢m⸣bres}
\nf{Témoins}{Export EDCS de 2020 : \emph{Maxima in pac\{a\}e q(uae) vixit an/ns plus minus XXXV conss(ulibus) / dominis nostris (H)onorio IIII c(on)s\{s\}(ule) et Eu/ticiano cons(ule) pridi\{a\}e Nonas Septebres}. EDB (EDB31859 ; écarts à examiner) : \emph{Maxima in pacae q⟨u⟩ae vixit an/nus plus minus XXXV cons(ulibu)s / dominis nostris Onorio IIII c(on)s\{s\}(ule) et Eu/ticiano cons(ule) pridiae nonas septenbres / ((vitis)) [texte tronqué dans la liste de résultats]}.}
\nf{Apparat}{⸢o⸣ : V gravé ; ⸢y⸣ : I gravé ; ⸢m⸣ : N gravé ; \{a\} : lettres superflues ; \{s\} : lettres superflues ; \{a\} : lettres superflues}
\nf{Traduction}{Maxima, en paix, qui vécut environ 35 ans ; sous le consulat de nos seigneurs (sic) Honorius pour la 4e fois et Eutychianus, la veille des nones de septembre.}
\nf{Langue}{L’éditeur signale 3 graphies corrigées : V gravé pour o (confusion o/u) ; I gravé pour y (graphie de y) ; N gravé pour m (n pour m). L’export de 2020 diffère de l’état 2022 (voir apparat). Des lettres gravées sont jugées superflues par l’éditeur (entre accolades). 7 abréviations développées par l’éditeur (dont q(uae), conss(ulibus), (H)onorio, c(on)s, (ule), c(h)iano).}
\nf{Chronologie}{Noms consulaires : Honorius (lu) et Eutychianus (lu). Dates du jour : événement indéterminé : pridie Nonas Septembres = 4 septembre 398 (lu). Âge déclaré : env. 35 a. (lu).}
\nf{Commentaire}{La formule in pace est conservée.}
\nf{Bibliographie}{ICUR I 1941 ; ILCV 2570. Bibliographie complémentaire selon l’EDB : ILCV, 2570.}
\nf{Statut}{Lecture secondaire (EDCS) ; écarts avec l’EDB à arbitrer sur l’ICVR.}
\end{notice}

\begin{notice}{420}{398}{ICUR n.s. II, 4505 · EDCS-34100568 · EDB19792}
\label{not:420}
\nf{Datation}{398 — directe : nom consulaire lu}
\nf{Genre}{Épitaphe ; ensembles A, B, C ; 99 lettres transmises ; 42 \% restituées.}
\nf{Références}{ICUR II 4505.}
\nf{Provenance}{Rattachement selon l’EDB : \emph{Coem. Pontiani}. Conservation selon l’EDB : \emph{Coem. Pontiani}.}
\nf{Support}{Plaque de marbre. Inscription gravée. Reproductions signalées par l’EDB : photographie PCAS (droits réservés ; non reproduites). Fiche EDB : Antonio Enrico Felle, première version 2003.}
\ntext{\lnum{1}[D(omino)] n(ostro) Honorio Aug(usto) [IIII et Fl(avio) Euty]\\
\lnum{2}c(h)iano v(iro) c(larissimo) cons(ulibus)\\
\lnum{3}d(e)p(ositus) Victor qui v[ixit ann(os) - - -]\\
\lnum{4}d(ies) VIIII qui fecit cu[m coniuge suo]\\
\lnum{5}d(ies) LXX in p(ace)}
\nf{Témoins}{EDB (EDB19792) : texte concordant aux conventions de notation près (indice de similarité 0,992).}
\nf{Traduction}{Sous le consulat de notre seigneur Honorius Auguste pour la 4e fois et de Flavius Eutychianus, clarissime, déposé Victor, qui vécut […] ans […] 9 jours, qui a vécu avec son épouse 70 jours ; en paix.}
\nf{Langue}{14 abréviations développées par l’éditeur (dont D(omino), n(ostro), Aug(usto), Fl(avio), c(h)iano, v(iro)).}
\nf{Chronologie}{Noms consulaires : Honorius (lu) et Eutychianus (lu en partie). Âge déclaré : 9 j. (partiellement restitué).}
\nf{Commentaire}{La formule in pace est transmise sous forme abrégée, sans restitution de lacune ; son développement reste éditorial.}
\nf{Bibliographie}{ICUR II 4505.}
\nf{Statut}{Lecture secondaire (EDCS) concordante avec l’EDB : lecture stabilisée entre deux bases, non collationnée sur l’ICVR ni sur le monument. L’EDCS signale une photographie (non consultée).}
\end{notice}

\begin{notice}{421}{398}{ICUR n.s. III, 8165 · EDCS-35700194 · EDB24687}
\label{not:421}
\nf{Datation}{398 — directe : nom consulaire lu}
\nf{Genre}{Épitaphe ; ensembles A, B, C ; 85 lettres transmises ; 31 \% restituées.}
\nf{Références}{ICUR III 8165 = ILCV +1269.}
\nf{Provenance}{Rattachement selon l’EDB : \emph{Basilica ss. Nerei et Achillei in coem. Domitillae}. Conservation selon l’EDB : \emph{Basilica ss. Nerei et Achillei in coem. Domitillae}.}
\nf{Support}{Plaque de marbre. 34 × 88 × 3,5 cm (h. × l. × ép.). Inscription gravée. Reproductions signalées par l’EDB : photographie PCAS ×2 (droits réservés ; non reproduites). Fiche EDB : Carlo Carletti / Antonio Enrico Felle, première version 2003.}
\ntext{\lnum{1}[Hic quies]cit Pascentius lector de Fasc[iola]\\
\lnum{2}[qui vixit an]n⸢o⸣s XXI depositus in pace\\
\lnum{3}[- - -] cons(ule) d(omino) n(ostro) [H]onorio I[III]}
\nf{Témoins}{Export EDCS de 2020 : \emph{[Hic quies]cit Pascentius lector de Fasc[iola] / [qui vixit an]ns XXI depositus in pace / [3] cons(ule) d(omino) n(ostro) [H]onorio I[III]}. EDB (EDB24687) : texte concordant aux conventions de notation près (indice de similarité 0,986).}
\nf{Apparat}{⸢o⸣ : V gravé}
\nf{Traduction}{Ici repose Pascentius, lecteur de (titre de) Fasciola, qui vécut 21 ans ; déposé en paix […] sous le consulat de notre seigneur Honorius pour la 4e fois.}
\nf{Langue}{L’éditeur signale une graphie corrigée : V gravé pour o (confusion o/u).}
\nf{Chronologie}{Noms consulaires : Honorius (lu). Âge déclaré : 21 a. (partiellement restitué).}
\nf{Commentaire}{La formule in pace est conservée. Vocabulaire religieux : repos et sommeil de la mort — quiescit (partiellement restitué). Fonction ou métier mentionné : lector.}
\nf{Bibliographie}{ICUR III 8165 ; ILCV +1269. Bibliographie complémentaire selon l’EDB : ILCV, 1269adn.}
\nf{Statut}{Lecture secondaire (EDCS) concordante avec l’EDB : lecture stabilisée entre deux bases, non collationnée sur l’ICVR ni sur le monument. L’EDCS signale une photographie (non consultée).}
\end{notice}

\begin{notice}{422}{398}{ICUR n.s. III, 8173b · EDCS-35700212 · EDB24704}
\label{not:422}
\nf{Datation}{398 — directe : lecture partielle d’un nom consulaire}
\nf{Genre}{Épitaphe ; ensembles A, B, C ; 54 lettres transmises ; 61 \% restituées.}
\nf{Références}{ICUR III 8173b.}
\nf{Provenance}{Rattachement selon l’EDB : \emph{Basilica ss. Nerei et Achillei in coem. Domitillae}.}
\nf{Support}{Plaque de marbre. Inscription gravée. Fiche EDB : Carlo Carletti / Antonio Enrico Felle, première version 2003.}
\ntext{\lnum{1}[- - - de]posi[t - - -]\\
\lnum{2}[- - - Hon]orio Aug(usto) III[I et Eutychiano conss(ulibus) - - -]\\
\lnum{3}[- - - Mer]curius X[}
\nf{Témoins}{EDB (EDB24704 ; écarts à examiner) : \emph{------ / [--- de]posi[t ---] / [--- Hon]o⌜r⌝i Aug(usti) III[I et Eutychiani ---] / [--- Mer]curius X[---] [texte tronqué dans la liste de résultats]}.}
\nf{Traduction}{[…] déposé(e) […] ; sous le consulat d’Honorius Auguste pour la 4e fois et d’Eutychianus […] ; […] Mercurius 10 (?).}
\nf{Langue}{2 abréviations développées par l’éditeur (dont Aug(usto), conss(ulibus)).}
\nf{Chronologie}{Noms consulaires : Honorius (lu en partie) et Eutychianus (restitué).}
\nf{Bibliographie}{ICUR III 8173b.}
\nf{Statut}{Lecture secondaire (EDCS) ; écarts avec l’EDB à arbitrer sur l’ICVR. Plus de la moitié des lettres sont restituées.}
\end{notice}

\begin{notice}{423}{398}{ICUR n.s. III, 8432 · EDCS-35802019 · EDB24940}
\label{not:423}
\nf{Datation}{398 — directe : nom consulaire lu}
\nf{Genre}{Épitaphe ; ensembles A, B, C ; 80 lettres transmises ; 32 \% restituées.}
\nf{Références}{ICUR III 8432.}
\nf{Provenance}{Rattachement selon l’EDB : \emph{Incertum locum in coem. Domitillae}. Conservation selon l’EDB : \emph{coem. Domitillae, aula supra coemeterium}.}
\nf{Support}{Plaque de marbre. Inscription gravée. Reproductions signalées par l’EDB : photographie PCAS ×3 (droits réservés ; non reproduites). Fiche EDB : Carlo Carletti / Antonio Enrico Felle, première version 2003.}
\ntext{\lnum{1}inimitabi]li(?) [- - - vi]xit\\
\lnum{2}[- - - mens]es III dies XVI\\
\lnum{3}[- - - Kalen]das Sept[e]mbr⸢e⸣s\\
\lnum{4}[Honorio IIII et F]l(avio) Euthyciano conss(ulibus)}
\nf{Témoins}{Export EDCS de 2020 : \emph{inimitabi]li(?) [3 vi]xit / [3 mens]es III dies XVI / [3 Kalen]das Sept[e]mbrs / [Honorio IIII et F]l(avio) Euthyciano conss(ulibus)}. EDB (EDB24940) : texte concordant aux conventions de notation près (indice de similarité 0,952).}
\nf{Apparat}{⸢e⸣ : I gravé}
\nf{Traduction}{[…] inimitable […] vécut […] 3 mois 16 jours ; […] calendes de septembre, sous le consulat d’Honorius pour la 4e fois et de Flavius Eutychianus.}
\nf{Langue}{L’éditeur signale une graphie corrigée : I gravé pour e (confusion e/i).}
\nf{Chronologie}{Noms consulaires : Honorius (restitué) et Eutychianus (lu). Dates du jour : événement indéterminé : … Kalendas Septembres (partiellement restitué). Âge déclaré : 3 m. 16 j. (partiellement restitué).}
\nf{Bibliographie}{ICUR III 8432.}
\nf{Statut}{Lecture secondaire (EDCS) concordante avec l’EDB : lecture stabilisée entre deux bases, non collationnée sur l’ICVR ni sur le monument. L’EDCS signale une photographie (non consultée).}
\end{notice}

\begin{notice}{424}{398}{ICUR n.s. III, 8647 · EDCS-35900846 · EDB19129}
\label{not:424}
\nf{Datation}{398 — directe : nom consulaire lu}
\nf{Genre}{Acte d’achat ou de réservation ; ensembles A, B ; 104 lettres transmises.}
\nf{Références}{ICUR III 8647.}
\nf{Provenance}{Rattachement selon l’EDB : \emph{Coem. Commodillae}.}
\nf{Support}{Plaque de marbre. Inscription gravée. Reproductions signalées par l’EDB : photographie PCAS (droits réservés ; non reproduites). Fiche EDB : Carlo Carletti, première version 2003.}
\ntext{\lnum{1}Locus Zosimi et Concordi(a)es(!)\\
\lnum{2}co⸢m⸣paratus a Preiecticio\\
\lnum{3}fossore\\
\lnum{4}consulatus domi(no) nostro Ho\\
\lnum{5}norio quater et Eut⸢y⸣c(h)iano consulib(us)}
\nf{Témoins}{Export EDCS de 2020 : \emph{Locus Zosimi et Concordi(a)es(!) / coparatus a Preiecticio / fossore / consulatus domi(no) nostro Ho/norio quater et Euticiano consulib(us)}. EDB (EDB19129) : texte concordant aux conventions de notation près (indice de similarité 0,948).}
\nf{Apparat}{⸢m⸣ : N gravé ; ⸢y⸣ : I gravé ; (!) : forme signalée par l’éditeur}
\nf{Traduction}{Emplacement de Zosimus et de Concordia, acquis auprès du fossoyeur Proiecticius, sous le consulat de notre seigneur Honorius pour la 4e fois et d’Eutychianus.}
\nf{Langue}{L’éditeur signale 2 graphies corrigées : N gravé pour m (n pour m) ; I gravé pour y (graphie de y). L’export de 2020 diffère de l’état 2022 (voir apparat). L’éditeur signale au moins une forme anormale (sic). 4 abréviations développées par l’éditeur (dont Concordi(a)es, domi(no), c(h)iano, consulib(us)).}
\nf{Chronologie}{Noms consulaires : Honorius (lu) et Eutychianus (lu).}
\nf{Commentaire}{Fonction ou métier mentionné : fossore.}
\nf{Bibliographie}{ICUR III 8647.}
\nf{Statut}{Lecture secondaire (EDCS) concordante avec l’EDB : lecture stabilisée entre deux bases, non collationnée sur l’ICVR ni sur le monument. L’EDCS signale une photographie (non consultée).}
\end{notice}

\begin{notice}{425}{398}{ICUR n.s. IV, 9583 · EDCS-36101108 · EDB19322}
\label{not:425}
\nf{Datation}{398 — directe : nom consulaire lu}
\nf{Genre}{Épitaphe ; ensembles A, B, C ; 91 lettres transmises ; 30 \% restituées.}
\nf{Références}{ICUR IV 9583 = ILCV +2946.}
\nf{Provenance}{Rattachement selon l’EDB : \emph{Coem. Callisti pars inferior}.}
\nf{Support}{Plaque de marbre. Inscription gravée. Fiche EDB : Antonio Enrico Felle, première version 2003.}
\ntext{\lnum{1}]\\
\lnum{2}[c]ons(ulatu) d(omini) n(ostri) Honori Aug(usti) IIII [et Euty]\\
\lnum{3}[ch]iani v(iri) c(larissimi) VI Idus Ian(uarias) depo[sita]\\
\lnum{4}[est] Crescentia virgo quae vi[xit]\\
\lnum{5}[ann(os)] VII men[ses}
\nf{Témoins}{EDB (EDB19322) : texte concordant aux conventions de notation près (indice de similarité 0,962).}
\nf{Traduction}{Sous le consulat de notre seigneur Honorius Auguste pour la 4e fois et d’Eutychianus, clarissime, le 6e jour avant les ides de janvier, a été déposée Crescentia, vierge, qui vécut 7 ans […] mois.}
\nf{Langue}{8 abréviations développées par l’éditeur (dont cons(ulatu), d(omini), n(ostri), Aug(usti), v(iri), c(larissimi)).}
\nf{Chronologie}{Noms consulaires : Honorius (lu) et Eutychianus (lu en partie). Dates du jour : événement indéterminé : VI Idus Ianuarias = 8 janvier 398 (lu). Âge déclaré : 7 a. (partiellement restitué).}
\nf{Bibliographie}{ICUR IV 9583 ; ILCV +2946.}
\nf{Statut}{Lecture secondaire (EDCS) concordante avec l’EDB : lecture stabilisée entre deux bases, non collationnée sur l’ICVR ni sur le monument. L’EDCS signale une photographie (non consultée).}
\end{notice}

\begin{notice}{426}{398}{ICUR n.s. IV, 11129 · EDCS-37200416 · EDB39593}
\label{not:426}
\nf{Datation}{398 — directe : consulat restitué}
\nf{Genre}{Épitaphe ; ensembles A, B, C ; 22 lettres transmises ; 27 \% restituées.}
\nf{Références}{ICUR IV 11129.}
\nf{Provenance}{Rattachement selon l’EDB : \emph{Coem. Callisti area subdialis et incerti loci}. Conservation selon l’EDB : \emph{Coem. Callisti area subdialis et incerti loci}.}
\nf{Support}{Plaque de marbre. 13 × 26 cm (h. × l.). Inscription gravée. Fiche EDB : Marida Pierno, première version 2015.}
\ntext{\lnum{1}depo]situs\\
\lnum{2}[- - - Honori]o IIII c(on)s(ule)}
\nf{Témoins}{EDB (EDB39593) : texte concordant aux conventions de notation près (indice de similarité 0,964).}
\nf{Traduction}{Déposé […] sous le consulat d’Honorius pour la 4e fois.}
\nf{Langue}{2 abréviations développées par l’éditeur (dont c(on)s, (ule)).}
\nf{Chronologie}{Noms consulaires : Honorius (restitué).}
\nf{Bibliographie}{ICUR IV 11129.}
\nf{Statut}{Lecture secondaire (EDCS) concordante avec l’EDB : lecture stabilisée entre deux bases, non collationnée sur l’ICVR ni sur le monument.}
\end{notice}

\begin{notice}{427}{398}{ICUR n.s. IV, 12424 · EDCS-37600947 · EDB41702}
\label{not:427}
\nf{Datation}{398 — directe : nom consulaire lu}
\nf{Genre}{Épitaphe ; ensembles A, B, C ; 54 lettres transmises ; 4 \% restituées.}
\nf{Références}{ICUR IV 12424.}
\nf{Provenance}{Rattachement selon l’EDB : \emph{inter vias Appiam et Ardeatinam}. Conservation selon l’EDB : \emph{Musei Capitolini}.}
\nf{Support}{Plaque de marbre. 24 × 31 × 4 cm (h. × l. × ép.). Inscription gravée. Fiche EDB : Domenico Raiola, première version 2017.}
\ntext{\lnum{1}[- - -]a in pace IIII Nonas\\
\lnum{2}[- - - q]ui vixit ann⸢o⸣s LXXX\\
\lnum{3}[- - - d]ie II Honorio IIII EBOSI(?)}
\nf{Témoins}{Export EDCS de 2020 : \emph{[3]a in pace IIII Nonas / [3 q]ui vixit anns LXXX / [3 d]ie II Honorio IIII EBOSI(?)}. EDB (EDB41702) : texte concordant aux conventions de notation près (indice de similarité 0,95).}
\nf{Apparat}{⸢o⸣ : V gravé}
\nf{Traduction}{[…] en paix, le 4e jour avant les nones […] ; […] qui vécut 80 ans ; […] le jour 2 (?) ; Honorius pour la 4e fois [fin obscure : EBOSI].}
\nf{Langue}{L’éditeur signale une graphie corrigée : V gravé pour o (confusion o/u).}
\nf{Chronologie}{Noms consulaires : Honorius (lu). Âge déclaré : 80 a. 2 j. (partiellement restitué).}
\nf{Commentaire}{La formule in pace est conservée.}
\nf{Bibliographie}{ICUR IV 12424.}
\nf{Statut}{Lecture secondaire (EDCS) concordante avec l’EDB : lecture stabilisée entre deux bases, non collationnée sur l’ICVR ni sur le monument. L’EDCS signale une photographie (non consultée).}
\end{notice}

\begin{notice}{428}{398}{ICUR n.s. IV, 12543 · EDCS-37801603 · EDB41825}
\label{not:428}
\nf{Datation}{398 — directe : lecture partielle d’un nom consulaire}
\nf{Genre}{Épitaphe ; ensembles A, B, C ; 190 lettres transmises ; 48 \% restituées.}
\nf{Références}{ICUR IV 12543 = ILCV 4539 = Epigraphica-2017-417 = AE 2017, 91 = Mossong-2022, 00133.}
\nf{Provenance}{Rome (indication de l’EDCS) ; aucun rattachement cimétérial transmis. Conservation selon l’EDB : \emph{Musei Vaticani: Lapidario Cristiano}.}
\nf{Support}{Plaque de marbre. 23 × 45 cm (h. × l.). Inscription gravée. Fiche EDB : Antonio Enrico Felle, première version 2018.}
\ntext{\lnum{1}Hic quie[scit in pace vir no]\\
\lnum{2}mine Laur[entius sanctus]\\
\lnum{3}[p]r(es)b(yter) [eccl(esiae) cathol(icae) dep(ositus) die]\\
\lnum{4}Satu[rni] XV Kal(endas) Ia[n(uarias) qui vixit]\\
\lnum{5}ann(os) p(lus) m(inus) [X]XII post con[s(ulatum) H]ono[ri]\\
\lnum{6}Aug(usti) IIII [et E]utychiani v(iri) c(larissimi) bene\\
\lnum{7}merenti matrona eius\\
\lnum{8}Remigia fecit}
\nf{Témoins}{EDB (EDB41825 ; écarts à examiner) : \emph{hic quie[scit --- no]/mine Laure[ntius ---] / [p]r(es)b(yter) [--- die] / Satu[rni] XV kal(endas) ia[n(uarias) ---] / ann(is) p(lus) m(inus) [X]XII post cons(ulatum) Hono[ri] / Aug(usti) IIII [et E]utychian̲i̲ v̲(iri) c̲(larissimi) bene/merent̲i̲ m̲atron̲a̲ e̲i̲u̲s / Remi̲gi̲a̲ f̲e̲c̲i̲t̲}.}
\nf{Traduction}{Ici repose en paix un homme nommé Laurentius, saint prêtre de l’Église catholique, déposé un samedi, le 15e jour avant les calendes de janvier ; il vécut environ 22 (?) ans ; après le consulat d’Honorius Auguste pour la 4e fois et d’Eutychianus, clarissime ; à lui, bien méritant, son épouse Remigia a fait (ce monument).}
\nf{Langue}{14 abréviations développées par l’éditeur (dont pr(es)b, (yter), eccl(esiae), cathol(icae), dep(ositus), Kal(endas)).}
\nf{Chronologie}{Noms consulaires : Honorius (lu en partie) et Eutychianus (lu en partie). Datation par post-consulat. Réserve : Post-consulat d’Honorius IV et Eutychianus (399 ?) ou consulat (398) ; le jour de la semaine (samedi) convient au 18 décembre 398. Dates du jour : déposition : XV Kalendas Ianuarias = 18 décembre 398 (partiellement restitué). Âge déclaré : env. 22 a. (partiellement restitué).}
\nf{Commentaire}{La formule in pace est attestée seulement par restitution. Vocabulaire religieux : Église (épithète) — sanctus (entièrement restitué) ; vocabulaire ecclésial — ecclesiae (entièrement restitué) ; repos et sommeil de la mort — quiescit (partiellement restitué). Fonction ou métier mentionné : presbyter. Le titre [p]r(es)b(yter) est partiellement restitué ; [eccl(esiae) cathol(icae)] est entièrement restitué. Ces groupes ne sont pas des attestations intégralement lues de l’institution ou de son qualificatif.}
\nf{Bibliographie}{ICUR IV 12543 ; ILCV 4539 ; Epigraphica-2017-417 ; AE 2017, 91 ; Mossong-2022, 00133. Bibliographie complémentaire selon l’EDB : Bultrighini 2017, 417-418, no.3 ; ILCV, 4539 ; CLRE, 662.}
\nf{Statut}{Lecture secondaire (EDCS) ; écarts avec l’EDB à arbitrer sur l’ICVR. L’EDCS signale une photographie (non consultée).}
\end{notice}

\annee{399}

\begin{notice}{429}{399}{EDCS-18200424}
\label{not:429}
\nf{Datation}{399 — directe : nom consulaire lu}
\nf{Genre}{Épitaphe ; ensembles A, B, C ; 144 lettres transmises ; 43 \% restituées.}
\nf{Références}{CIL VI 8405 = ILCV 120.}
\nf{Provenance}{Rome (indication de l’EDCS) ; aucun rattachement cimétérial transmis.}
\nf{Support}{Non documenté dans les sources consultées.}
\ntext{\lnum{1}[Hic requiescit] Martielis numerarius vicari(a)e [sedis]\\
\lnum{2}[urbis Romae] aeternae qui v(i)xit p(lus) m(inus) annis LXX[- - -]\\
\lnum{3}[dep(ositus) in pa]ce die V Kal(endas) Dec(embres) Fl(avio) Mallio Theodoro v(iro) [c(larissimo)]\\
\lnum{4}[cons(ule) in pa]ce}
\nf{Témoins}{Aucune fiche EDB appariée (pas de numéro ICUR ; recherche par le texte infructueuse).}
\nf{Traduction}{Ici repose Martielis (?), comptable (numerarius) du siège vicarial de la Ville de Rome éternelle, qui vécut environ 70 ans ; déposé en paix le 5e jour avant les calendes de décembre, sous le consulat de Flavius Mallius Theodorus, clarissime ; en paix.}
\nf{Langue}{11 abréviations développées par l’éditeur (dont vicari(a)e, v(i)xit, p(lus), m(inus), dep(ositus), Kal(endas)).}
\nf{Chronologie}{Noms consulaires : Theodorus (lu). Dates du jour : déposition : die V Kalendas Decembres = 27 novembre 399 (lu). Âge déclaré : env. 5 j. (partiellement restitué).}
\nf{Commentaire}{La formule in pace est attestée seulement par restitution. Vocabulaire religieux : repos et sommeil de la mort — requiescit (entièrement restitué). Fonction ou métier mentionné : numerarius.}
\nf{Bibliographie}{CIL VI 8405 ; ILCV 120.}
\nf{Statut}{Lecture secondaire (EDCS) sans second témoin numérique. L’EDCS signale une photographie (non consultée).}
\end{notice}

\begin{notice}{430}{399}{ICUR n.s. I, 987 · EDCS-19900450 · EDB29620}
\label{not:430}
\nf{Datation}{399 — directe : nom consulaire lu}
\nf{Genre}{Épitaphe ; ensembles A, B, C ; 62 lettres transmises ; 27 \% restituées.}
\nf{Références}{CIL VI 32069 (p 4809) = ICUR I 987 = ILCV +139.}
\nf{Provenance}{Rattachement selon l’EDB : \emph{Incertae originis} (le volume I des ICVR réunit les inscriptions dont la provenance est perdue). Conservation selon l’EDB : \emph{ss. Quattro Coronati}.}
\nf{Support}{Plaque de marbre. 51 × 45 × 7 cm (h. × l. × ép.). Inscription gravée. Signes décrits par l’EDB : croix monogrammatique. Reproductions signalées par l’EDB : reproduction (droits réservés ; non reproduites). Fiche EDB : Filippo Piazzolla, première version 2014.}
\ntext{\lnum{1}cons(ulis) Fl(avi) Ma]lli Theodori v(iri) c(larissimi) No[nis - - -]\\
\lnum{2}[L]epus(?) v(ir) c(larissimus) defunctus in [pace - - -]\\
\lnum{3}[- - - locus Roma]nill(a)e}
\nf{Témoins}{EDB (EDB29620) : texte concordant aux conventions de notation près (indice de similarité 0,94).}
\nf{Traduction}{Sous le consulat de Flavius Mallius Theodorus, clarissime, aux nones […] ; Lepus, clarissime, mort en paix […] ; […] emplacement de Romanilla.}
\nf{Langue}{8 abréviations développées par l’éditeur (dont cons(ulis), Fl(avi), v(iri), c(larissimi), Lepus(?), v(ir)).}
\nf{Chronologie}{Noms consulaires : Theodorus (lu).}
\nf{Commentaire}{La formule in pace est attestée seulement par restitution. Rang sénatorial d’une personne commémorée (clarissimi nonis … lepus vir clarissimus defu).}
\nf{Bibliographie}{CIL VI 32069 (p 4809) ; ICUR I 987 ; ILCV +139. Bibliographie complémentaire selon l’EDB : ILCV, 139 ; CIL VI, 32069, cfr. p. 4809.}
\nf{Statut}{Lecture secondaire (EDCS) concordante avec l’EDB : lecture stabilisée entre deux bases, non collationnée sur l’ICVR ni sur le monument. L’EDCS signale une photographie (non consultée).}
\end{notice}

\begin{notice}{431}{399}{ICUR n.s. I, 309 · EDCS-38600308 · EDB27849}
\label{not:431}
\nf{Datation}{399 — directe : nom consulaire lu ; datation contestée}
\nf{Genre}{Épitaphe ; ensembles A, B, C ; 156 lettres transmises ; 12 \% restituées.}
\nf{Références}{ICUR I 309 (p 490) = ILCV 4400b.}
\nf{Provenance}{Rattachement selon l’EDB : \emph{Incertae originis} (le volume I des ICVR réunit les inscriptions dont la provenance est perdue).}
\nf{Support}{Support non précisé par l’EDB. Signes décrits par l’EDB : chrisme et chrisme avec alpha et oméga. Fiche EDB : Filippo Piazzolla, première version 2014.}
\ntext{\lnum{1}Bene merenti in pace [- - -]\\
\lnum{2}quae vixit ann⸢o⸣s quat(t)u(o)r [menses - - - a]\\
\lnum{3}mor filiae dulcissimae deposita[e in pace]\\
\lnum{4}septimum Kalendas Septembr⸢e⸣s die Veneris [post]\\
\lnum{5}consulatu d(omini) n(ostri) \{H\}Arc\{h\}adi IIII et Euty\\
\lnum{6}chiani v(irorum) c(larissimorum)\\
\lnum{7}AC}
\nf{Témoins}{Export EDCS de 2020 : \emph{Bene merenti in pace [3] / quae vixit anns quat(t)u(o)r [menses 3 a]/mor filiae dulcissimae deposita[e in pace] / septimum Kalendas Septembrs die Veneris [post] / consulatu d(omini) n(ostri) \{H\}Arc\{h\}adi IIII et Euty/chiani v(irorum) c(larissimorum) / AC}. EDB (EDB27849 ; divergence forte) : \emph{((monogramma Christi cum alpha et omega)) benemerenti in pace [---] / quae vixit annus quatur [menses --- me?]mor filiae dulcissimae deposita[e in pace ] / septimum kalendas septembris die Veneris [post?] / consulatu d(omini) n(ostri) Harchadi IIII et Euty/chiani v(iri) c(larissimi) cons(ulibu)s / AC [texte tronqué dans la liste de résultats]}.}
\nf{Apparat}{⸢o⸣ : V gravé ; ⸢e⸣ : I gravé ; \{H\} : lettres superflues ; \{h\} : lettres superflues}
\nf{Traduction}{Bien méritante, en paix, […] qui vécut quatre ans […] mois ; amour d’une fille très douce, déposée en paix le 7e jour avant les calendes de septembre, un vendredi, après le consulat de notre seigneur Arcadius pour la 4e fois et d’Eutychianus (sic), clarissimes.}
\nf{Langue}{L’éditeur signale 2 graphies corrigées : V gravé pour o (confusion o/u) ; I gravé pour e (confusion e/i). Des lettres gravées sont jugées superflues par l’éditeur (entre accolades). 6 abréviations développées par l’éditeur (dont quat(t)u, (o)r, d(omini), n(ostri), v(irorum), c(larissimorum)).}
\nf{Chronologie}{Noms consulaires : Arcadius (lu) et Eutychianus (lu). Datation par post-consulat. La lecture retenue date le texte de 399 (post-consulat d'Honorius IV et Eutychianus). L'EDB (EDB27849), citant la CLRE (p. 661), l'attribue à 398 : le \emph{post} restitué serait trop long, et un post-consulat aussi tardif en 399 est sans exemple. Le lapicide a gravé \emph{Harchadi} pour \emph{Honorii} (erreur déjà relevée par Muratori et Silvagni). Datation contestée (398/399). Dates du jour : déposition : septimum Kalendas Septembres = 26 août 399 (lu). Âge déclaré : 4 a. (partiellement restitué).}
\nf{Commentaire}{La formule in pace est conservée.}
\nf{Bibliographie}{ICUR I 309 (p 490) ; ILCV 4400b. Bibliographie complémentaire selon l’EDB : ILCV, 4400B ; CLRE, 661.}
\nf{Statut}{Lecture secondaire (EDCS) ; divergence forte avec l’EDB : le texte ne doit pas être exploité sans contrôle de l’édition imprimée.}
\end{notice}

\begin{notice}{432}{399}{ICUR n.s. I, 524 · EDCS-38700928 · EDB28503}
\label{not:432}
\nf{Datation}{399 — directe : consulat restitué}
\nf{Genre}{Épitaphe ; ensembles A, B, C ; 21 lettres transmises ; 29 \% restituées.}
\nf{Références}{ICUR I 524 (p 490).}
\nf{Provenance}{Rattachement selon l’EDB : \emph{Incertae originis} (le volume I des ICVR réunit les inscriptions dont la provenance est perdue).}
\nf{Support}{Support non précisé par l’EDB. Fiche EDB : Filippo Piazzolla, première version 2014.}
\ntext{\lnum{1}Laur]entia(?)\\
\lnum{2}[- - - Theodo]ro cons(ule)}
\nf{Témoins}{EDB (EDB28503) : texte concordant aux conventions de notation près (indice de similarité 0,962).}
\nf{Traduction}{Laurentia […] ; sous le consulat de Theodorus.}
\nf{Langue}{2 abréviations développées par l’éditeur (dont Laurentia(?), cons(ule)).}
\nf{Chronologie}{Noms consulaires : Theodorus (restitué).}
\nf{Bibliographie}{ICUR I 524 (p 490).}
\nf{Statut}{Lecture secondaire (EDCS) concordante avec l’EDB : lecture stabilisée entre deux bases, non collationnée sur l’ICVR ni sur le monument.}
\end{notice}

\begin{notice}{433}{399}{ICUR n.s. I, 1453 · EDCS-35100468 · EDB31243}
\label{not:433}
\nf{Datation}{399 — directe : nom consulaire lu}
\nf{Genre}{Épitaphe ; ensembles A, B, C ; 314 lettres transmises ; 43 \% restituées.}
\nf{Références}{ICUR I 1453 = ILCV 3432 = CLE 679.}
\nf{Provenance}{Rattachement selon l’EDB : \emph{Incertae originis} (le volume I des ICVR réunit les inscriptions dont la provenance est perdue). Conservation selon l’EDB : \emph{Musei Vaticani: Lapidario Cristiano}.}
\nf{Support}{Plaque de marbre. 26 × 32 cm (h. × l.). Inscription gravée. Fiche EDB : Filippo Piazzolla, première version 2015.}
\ntext{\lnum{1}Hic iace⸢t⸣ extinctus pri[mo sub limine vitae]\\
\lnum{2}augusto multum dilec[tus tempore pupus]\\
\lnum{3}nec reor hunc lacrimis [fas esse urgere beatum]\\
\lnum{4}corporis exutus vinc(u)lis q[ui gaudet in astris]\\
\lnum{5}nec mala terreni sensit c[ontagia sensus]\\
\lnum{6}menses namque nove[m postquam super addidit anno]\\
\lnum{7}conplevitque dies vitae s[ex ivit ad astra]\\
\lnum{8}[dep(ositus)] L(ucio) Mallio The(o)doro v(iro) c(larissimo) co(n)s(ule) pri(die) Non[as - -]\\
\lnum{9}q(ui) vixit annum m(enses) VIIII [dies VI]}
\nf{Témoins}{Export EDCS de 2020 : \emph{Hic iace extinctus pri[mo sub limine vitae] / augusto multum dilec[tus tempore pupus] / nec reor hunc lacrimis [fas esse urgere beatum] / corporis exutus vinc(u)lis q[ui gaudet in astris] / nec mala terreni sensit c[ontagia sensus] / menses namque nove[m postquam super addidit anno] / conplevitque dies vitae s[ex ivit ad astra] / [dep(ositus)] L(ucio) Mallio The(o)doro v(iro) c(larissimo) co(n)s(u…}. EDB (EDB31243 ; écarts à examiner) : \emph{hic iacet ext̲i̲n̲c̲t̲u̲s̲ p̲r̲i̲[mo sub limine vitae] / angusto mul̲t̲u̲m̲ d̲i̲l̲e̲c̲t̲[us tempore pupus] / nec reor hunc̲ l̲a̲c̲r̲i̲̲m̲̲i̲s̲ [fas esse urgere beatum] / corporis exut̲u̲s̲ v̲i̲n̲c̲l̲i̲s̲ q̲[ui gaudet in astris] / nec mala terr̲e̲n̲i̲ s̲e̲n̲s̲i̲t̲ c̲[ontagia sensus] / menses namq̲u̲e̲ n̲o̲v̲e̲[m postquam super addidit anno] / co⌜m⌝plevitque d̲i̲e̲s̲ v̲i̲t̲a̲e̲ s̲e̲[x ivit ad astra] / [depositus F]l(avio) Mallio Thedoro v(iro) c(larissimo) c̲o̲(n)s̲(ule) p̲r̲i̲(die) n̲o̲n̲(as) [---] / q(ui) vixit an(num) un̲u̲m̲ m̲(enses) V̲I̲I̲I̲I̲ [d(ies) VI]}.}
\nf{Apparat}{⸢t⸣ : I gravé}
\nf{Traduction}{(vers) Ici gît, éteint au premier seuil de la vie, un petit enfant très aimé pendant un temps bien court ; et je ne pense pas qu’il soit permis d’accabler de larmes ce bienheureux qui, dépouillé des liens du corps, se réjouit parmi les astres, et n’a pas ressenti les contagions mauvaises des sens terrestres ; car après avoir ajouté neuf mois à une année et accompli six jours de vie, il s’en alla vers les astres. Déposé sous le consulat de Lucius (sic) Mallius Theodorus, clarissime, la veille des nones […] ; il vécut un an 9 mois 6 jours.}
\nf{Langue}{L’éditeur signale une graphie corrigée : I gravé pour t.}
\nf{Chronologie}{Noms consulaires : Theodorus (lu). Âge déclaré : 9 m. 6 j. (partiellement restitué).}
\nf{Commentaire}{Vocabulaire religieux : allusion interprétative (niveau 5) — astris (entièrement restitué), astra (entièrement restitué). Commentaire développé : dossier p.~\pageref{dos:433}.}
\nf{Bibliographie}{ICUR I 1453 ; ILCV 3432 ; CLE 679. Bibliographie complémentaire selon l’EDB : Lega 2000, 141 ; ILCV, 3432.}
\nf{Statut}{Lecture secondaire (EDCS) ; écarts avec l’EDB à arbitrer sur l’ICVR. L’EDCS signale une photographie (non consultée).}
\end{notice}

\annee{397}

\begin{notice}{434}{397}{ICUR n.s. II, 6476 · EDCS-34300830 · EDB34468}
\label{not:434}
\nf{Datation}{397 — directe : nom consulaire lu ; datation contestée}
\nf{Genre}{Épitaphe ; ensembles A, B, C ; 32 lettres transmises ; 22 \% restituées.}
\nf{Références}{ICUR II 6476.}
\nf{Provenance}{Rattachement selon l’EDB : \emph{Coem. Commodillae}. Conservation selon l’EDB : \emph{Coem. Commodillae}.}
\nf{Support}{Enduit de mortier de mortier. Inscription tracée dans le mortier frais. Signes décrits par l’EDB : chrisme et chrisme avec alpha et oméga. Fiche EDB : Filippo Piazzolla, première version 2014.}
\ntext{\lnum{1}] d(e)p(ositus ?) V Kal(endas) Dec(embres)\\
\lnum{2}[- - -]cam p(ost) cons\\
\lnum{3}[ulatu(?)] Arcadi I[II]I}
\nf{Témoins}{EDB (EDB34468 ; écarts à examiner) : \emph{d(e)p(osit---) V kal(endas) dec(embres) ((monogramma Christi cum alpha et omega)) / Cam(---) p(ost) cons(ulatum) / Arcadi I[II]I [texte tronqué dans la liste de résultats]}.}
\nf{Traduction}{Déposé le 5e jour avant les calendes de décembre […] après le consulat d’Arcadius pour la 4e fois.}
\nf{Langue}{6 abréviations développées par l’éditeur (dont d(e)p, (ositus ?), Kal(endas), Dec(embres), p(ost), ulatu(?)).}
\nf{Chronologie}{Noms consulaires : Arcadius (lu). Datation par post-consulat. \emph{p(ost) cons(ulatum) Arcadi I[II]I} désigne l'année qui suit le quatrième consulat d'Arcadius (396), soit 397 ; l'EDB (EDB34468) date de 399. L'année 397 est conservée ; la divergence est signalée. Dates du jour : déposition : V Kalendas Decembres = 27 novembre 397 (lu).}
\nf{Bibliographie}{ICUR II 6476.}
\nf{Statut}{Lecture secondaire (EDCS) ; écarts avec l’EDB à arbitrer sur l’ICVR.}
\end{notice}

\annee{399}

\begin{notice}{435}{399}{ICUR n.s. III, 6511 · EDCS-34801311 · EDB22306}
\label{not:435}
\nf{Datation}{399 — conjecturale}
\nf{Genre}{Épitaphe ; ensembles A, B, C ; 13 lettres transmises ; 46 \% restituées.}
\nf{Références}{ICUR III 6511.}
\nf{Provenance}{Rattachement selon l’EDB : \emph{Coem. Domitillae pars inferior}. Conservation selon l’EDB : \emph{Coem. Domitillae pars inferior}.}
\nf{Support}{Plaque de marbre. 16 × 16 × 2,3 cm (h. × l. × ép.). Inscription gravée. Fiche EDB : Carlo Carletti / Antonio Enrico Felle, première version 2003.}
\ntext{\lnum{1}[- - -]on[- - -]\\
\lnum{2}[- - -]ro v(iro) c(larissimo) c[onsule]}
\nf{Témoins}{EDB (EDB22306 ; divergence forte) : \emph{------ / [---]on[---] / [--- Mallio Theodo]ro v(iro) c(larissimo) c[onsule ---] [texte tronqué dans la liste de résultats]}.}
\nf{Traduction}{[…] sous le consulat de (Theodo)rus, clarissime.}
\nf{Langue}{2 abréviations développées par l’éditeur (dont v(iro), c(larissimo)).}
\nf{Chronologie}{Seul \emph{…ro v(iro) c(larissimo) c[onsule]} est lu. L'EDB (EDB22306) restitue \emph{[Mallio Theodo]ro} (399). Datation conjecturale.}
\nf{Bibliographie}{ICUR III 6511.}
\nf{Statut}{Lecture secondaire (EDCS) ; divergence forte avec l’EDB : le texte ne doit pas être exploité sans contrôle de l’édition imprimée. L’EDCS signale une photographie (non consultée).}
\end{notice}

\begin{notice}{436}{399}{ICUR n.s. III, 8166 · EDCS-35700195 · EDB24688}
\label{not:436}
\nf{Datation}{399 — directe : lecture partielle d’un nom consulaire}
\nf{Genre}{Épitaphe ; ensembles A, B, C ; 55 lettres transmises ; 56 \% restituées.}
\nf{Références}{ICUR III 8166.}
\nf{Provenance}{Rattachement selon l’EDB : \emph{Basilica ss. Nerei et Achillei in coem. Domitillae}. Conservation selon l’EDB : \emph{Coem. Domitillae pars inferior}.}
\nf{Support}{Plaque de marbre. 13 × 24 × 1,8 cm (h. × l. × ép.). Inscription gravée. Reproductions signalées par l’EDB : photographie PCAS (droits réservés ; non reproduites). Fiche EDB : Carlo Carletti / Antonio Enrico Felle, première version 2003.}
\ntext{\lnum{1}dep(ositus ?) q]uartu(m) Non[as - - -]\\
\lnum{2}[- - - Fl(avio)] Mallio The[odoro]\\
\lnum{3}[v(iro) c(larissimo) cons]ule\{m\}}
\nf{Témoins}{EDB (EDB24688 ; écarts à examiner) : \emph{[--- q]uartu non[as ---] / [--- Fl(avio)] Mallio The[odoro] / [v(iro) c(larissimo) cons]ulem}.}
\nf{Apparat}{\{m\} : lettres superflues}
\nf{Traduction}{Déposé le 4e jour avant les nones […] ; sous le consulat de Flavius Mallius Theodorus, clarissime.}
\nf{Langue}{Des lettres gravées sont jugées superflues par l’éditeur (entre accolades). 5 abréviations développées par l’éditeur (dont dep(ositus ?), quartu(m), Fl(avio), v(iro), c(larissimo)).}
\nf{Chronologie}{Noms consulaires : Theodorus (lu en partie).}
\nf{Bibliographie}{ICUR III 8166.}
\nf{Statut}{Lecture secondaire (EDCS) ; écarts avec l’EDB à arbitrer sur l’ICVR. Plus de la moitié des lettres sont restituées. L’EDCS signale une photographie (non consultée).}
\end{notice}

\begin{notice}{437}{399}{ICUR n.s. IV, 11124d · EDCS-37200410 · EDB39588}
\label{not:437}
\nf{Datation}{399 — directe : lecture partielle d’un nom consulaire}
\nf{Genre}{Épitaphe ; ensembles A, B, C ; 25 lettres transmises ; 64 \% restituées.}
\nf{Références}{ICUR IV 11124d.}
\nf{Provenance}{Rattachement selon l’EDB : \emph{Coem. Callisti area subdialis et incerti loci}. Conservation selon l’EDB : \emph{Coem. Callisti area subdialis et incerti loci}.}
\nf{Support}{Plaque de marbre. 22 × 30 cm (h. × l.). Inscription gravée. Fiche EDB : Marida Pierno, première version 2015.}
\ntext{\lnum{1}]XII m(enses) [- - -]\\
\lnum{2}[- - - cons(ulatu) Mall]i T(h)eod[ori]}
\nf{Témoins}{EDB (EDB39588 ; écarts à examiner) : \emph{------ / [---]XII m(ensibus) [---] / [consulatu Mall]i Teod[ori ---]}.}
\nf{Traduction}{[…] 12 mois […] ; sous le consulat de Mallius Theodorus.}
\nf{Langue}{3 abréviations développées par l’éditeur (dont m(enses), cons(ulatu), T(h)eod).}
\nf{Chronologie}{Noms consulaires : Theodorus (lu en partie).}
\nf{Bibliographie}{ICUR IV 11124d.}
\nf{Statut}{Lecture secondaire (EDCS) ; écarts avec l’EDB à arbitrer sur l’ICVR. Plus de la moitié des lettres sont restituées.}
\end{notice}

\begin{notice}{438}{399}{ICUR n.s. IV, 11130 · EDCS-37200417 · EDB39594}
\label{not:438}
\nf{Datation}{399 — directe : consulat restitué}
\nf{Genre}{Épitaphe ; ensembles A, B, C ; 45 lettres transmises ; 89 \% restituées.}
\nf{Références}{ICUR IV 11130.}
\nf{Provenance}{Rattachement selon l’EDB : \emph{Coem. Callisti area subdialis et incerti loci}. Conservation selon l’EDB : \emph{Coem. Callisti area subdialis et incerti loci}.}
\nf{Support}{Plaque de marbre. 18 × 13 × 7,5 cm (h. × l. × ép.). Inscription gravée. Fiche EDB : Marida Pierno, première version 2015.}
\ntext{\lnum{1}] ux[or - - -]\\
\lnum{2}[- - -] Fl(avio) M[allio Theodoro v(iro) c(larissimo) conss(ulibus)]}
\nf{Témoins}{EDB (EDB39594 ; écarts à examiner) : \emph{------ / [---] ux[or ---] / [---] Fl(avio) M[allio Theodoro viro clarissimo consule] / ------}.}
\nf{Traduction}{[…] épouse […] ; sous le consulat de Flavius Mallius Theodorus, clarissime.}
\nf{Langue}{4 abréviations développées par l’éditeur (dont Fl(avio), v(iro), c(larissimo), conss(ulibus)).}
\nf{Chronologie}{Noms consulaires : Theodorus (restitué).}
\nf{Bibliographie}{ICUR IV 11130.}
\nf{Statut}{Lecture secondaire (EDCS) ; écarts avec l’EDB à arbitrer sur l’ICVR. Plus de la moitié des lettres sont restituées.}
\end{notice}

\annee{400}

\begin{notice}{439}{400}{ICUR n.s. I, 3215 · EDCS-05801896 · EDB33162}
\label{not:439}
\nf{Datation}{400 — directe : nom consulaire lu}
\nf{Genre}{Épitaphe ; ensembles A, B, C ; 62 lettres transmises.}
\nf{Références}{CIL XIV 3896 = InscrIt-04-01, 00543 = ILCV 4181 (em) = ICUR I 3215.}
\nf{Provenance}{Provenance non établie ; Tibur probable d’après CIL XIV 3896 = InscrIt IV 1, 543 ; exclue des analyses topographiques propres à Rome. Rattachement selon l’EDB : \emph{Incertae originis} (le volume I des ICVR réunit les inscriptions dont la provenance est perdue).}
\nf{Support}{Support non précisé par l’EDB. Fiche EDB : Filippo Piazzolla, première version 2015.}
\ntext{\lnum{1}]\\
\lnum{2}parenti s(uo) posuerunt\\
\lnum{3}Petulum contra\\
\lnum{4}votu(m) et dolo(re) suo recedi\\
\lnum{5}t co(n)s(ulis) \{I\}Stil(i)c(h)onis}
\nf{Témoins}{EDB (EDB33162) : texte concordant aux conventions de notation près (indice de similarité 0,947). Provenance non établie. L'EDB (EDB33162 = ICVR I 3215) rattache le texte à la série d'origine incertaine et renvoie à ILCV 4181 et CIL XIV 3896 ; la concordance avec Inscriptiones Italiae IV 1, 543 désigne Tibur. Variantes entre les bases : \emph{parenti s(uo) posuerunt} (EDCS) / \emph{parentis posuerun[t ---]} (EDB) ; \emph{recedi/t} / \emph{rec[essi]/t} ; \emph{co(n)s(ulis) \{I\}Stil(i)c(h)onis} / \emph{c̅o(n)s̅(ule) Istiliconis}. La notice est conservée et exclue des analyses topographiques propres à Rome.}
\nf{Apparat}{\{I\} : lettres superflues}
\nf{Traduction}{[…] à leur parent, ils ont posé (ce monument) ; Petulus, contre leur vœu et à leur douleur, s’en va, sous le consulat de Stilicon.}
\nf{Langue}{Des lettres gravées sont jugées superflues par l’éditeur (entre accolades). 7 abréviations développées par l’éditeur (dont s(uo), votu(m), dolo(re), co(n)s, (ulis), Stil(i)c).}
\nf{Chronologie}{Noms consulaires : Stilicho (lu).}
\nf{Commentaire}{L’attribution romaine reste incertaine ; la référence CIL XIV / Inscriptiones Italiae justifie un contrôle de provenance, mais ne suffit pas ici à déplacer la pièce avec certitude vers Tibur.}
\nf{Bibliographie}{CIL XIV 3896 ; InscrIt-04-01, 00543 ; ILCV 4181 (em) ; ICUR I 3215. Bibliographie complémentaire selon l’EDB : ILCV, 4181 ; CIL XIV, 3896.}
\nf{Statut}{Lecture secondaire (EDCS) concordante avec l’EDB : lecture stabilisée entre deux bases, non collationnée sur l’ICVR ni sur le monument. L’EDCS signale une photographie (non consultée).}
\end{notice}

\begin{notice}{440}{400}{ICUR n.s. I, 1454 · EDCS-35200404 · EDB31244}
\label{not:440}
\nf{Datation}{400 — directe : nom consulaire lu}
\nf{Genre}{Épitaphe ; ensembles A, B, C ; 124 lettres transmises ; 24 \% restituées.}
\nf{Références}{CLE 680 = ICUR I 1454 = ILCV 317 (add).}
\nf{Provenance}{Rattachement selon l’EDB : \emph{Incertae originis} (le volume I des ICVR réunit les inscriptions dont la provenance est perdue). Conservation selon l’EDB : \emph{Musei Vaticani: Lapidario Cristiano}.}
\nf{Support}{Plaque de marbre. 28 × 61 cm (h. × l.). Inscription gravée. Texte métrique. Fiche EDB : Filippo Piazzolla, première version 2015.}
\ntext{\lnum{1}Fl(avia) Crescentina h(onestissima) f(emina) quae vixit ann[os - - -]\\
\lnum{2}Fl(avio) Stilichone v(iro) c(larissimo) consule A[- - -]\\
\lnum{3}tu mihi aeternum reno⸢v⸣as tu[mulata dolorem]\\
\lnum{4}Crescentina meum si fas si[mul accipe corpus]}
\nf{Témoins}{Export EDCS de 2020 : \emph{Fl(avia) Crescentina h(onestissima) f(emina) quae vixit ann[os 3] / Fl(avio) Stilichone v(iro) c(larissimo) consule A[3] / tu mihi aeternum renoas tu[mulata dolorem] / Crescentina meum si fas si[mul accipe corpus]}. EDB (EDB31244 ; écarts à examiner) : \emph{Fl(avia) Crescentina h(onestissima) f(emina) quae vixit ann(is) [---] / Fl(avio) Stilicone v(iro) c(larissimo) consule A[--- maritus fecit?] / tu mihi aeternum renobas tụ[mulata dolorem] / Crescentina meum si fas si[t ---]}.}
\nf{Apparat}{⸢v⸣ : B gravé}
\nf{Traduction}{Flavia Crescentina, femme très honorable, qui vécut […] ans ; sous le consulat de Flavius Stilicon, clarissime […]. (vers) Toi, ensevelie, tu renouvelles pour moi une douleur éternelle ; Crescentina, s’il est permis, reçois en même temps mon corps.}
\nf{Langue}{L’éditeur signale une graphie corrigée : B gravé pour v (bêtacisme).}
\nf{Chronologie}{Noms consulaires : Stilicho (lu).}
\nf{Bibliographie}{CLE 680 ; ICUR I 1454 ; ILCV 317 (add). Bibliographie complémentaire selon l’EDB : Lega 2000, 134 ; Ferrua 1979b, 38, n. 1454.}
\nf{Statut}{Lecture secondaire (EDCS) ; écarts avec l’EDB à arbitrer sur l’ICVR. L’EDCS signale une photographie (non consultée).}
\end{notice}

\begin{notice}{441}{400}{ICUR n.s. I, 528 · EDCS-38700931 · EDB28506}
\label{not:441}
\nf{Datation}{400 — directe : nom consulaire lu}
\nf{Genre}{Épitaphe ; ensembles A, B, C ; 76 lettres transmises.}
\nf{Références}{ICUR I 528.}
\nf{Provenance}{Rattachement selon l’EDB : \emph{Incertae originis} (le volume I des ICVR réunit les inscriptions dont la provenance est perdue). Conservation selon l’EDB : \emph{s. Maria in Trastevere}.}
\nf{Support}{Plaque de marbre. 36 × 92 cm (h. × l.). Inscription gravée. Reproductions signalées par l’EDB : reproduction (droits réservés ; non reproduites). Fiche EDB : Filippo Piazzolla, première version 2015.}
\ntext{\lnum{1}]RANNS[- - -]\\
\lnum{2}[- - -]rcus d⸢e⸣positus\\
\lnum{3}[- - -]V K(a)l(n)d(as) Maias Fl(avio) ⸢S⸣til\{l\}ic(h)on⸢e⸣ v(iro) c(larissimo) c(on)s(ule)\\
\lnum{4}[- - -] et Benedictus fratr⸢e⸣s se ⸢v⸣i⸢v⸣i\\
\lnum{5}[- - -]m fecerunt}
\nf{Témoins}{Export EDCS de 2020 : \emph{]RANNS[3] / [3]rcus dpositus / [3]V K(a)l(n)d(as) Maias Fl(avio) til\{l\}ic(h)on v(iro) c(larissimo) c(on)s(ule) / [3] et Benedictus fratrs se ii / [3]m fecerunt}. EDB (EDB28506 ; écarts à examiner) : \emph{------ / [---]ṇịạṇụṣq̣(ue) [---] / [--- Lup?]ẹrcus dipositus / [---] V k(a)l(en)d(as) maias Fl(avi) Estilliconi v(iro) c(larissimo) c(on)s(ul)⌜e⌝) / [---] et Benedictus fratr⌜e⌝s se bibi / [---]m fecerunt ((orans puer))}.}
\nf{Apparat}{⸢e⸣ : I gravé ; ⸢S⸣ : E gravé ; ⸢e⸣ : I gravé ; ⸢e⸣ : I gravé ; ⸢v⸣ : B gravé ; ⸢v⸣ : B gravé ; \{l\} : lettres superflues}
\nf{Traduction}{[…] (Ma)rcus déposé le 1x (?) jour avant les calendes de mai, sous le consulat de Flavius Stilicon, clarissime ; […] et Benedictus, frères, de leur vivant, ont fait […].}
\nf{Langue}{L’éditeur signale 6 graphies corrigées : I gravé pour e (confusion e/i) ×3 ; E gravé pour S ; B gravé pour v (bêtacisme) ×2. Des lettres gravées sont jugées superflues par l’éditeur (entre accolades). 9 abréviations développées par l’éditeur (dont K(a)l, (n)d, (as), Fl(avio), ic(h)on, v(iro)).}
\nf{Chronologie}{Noms consulaires : Stilicho (lu). Dates du jour : déposition : …V Kalndas Maias = 27 avril 400 (partiellement restitué).}
\nf{Bibliographie}{ICUR I 528. Bibliographie complémentaire selon l’EDB : ILCV, 4146F adn. ; Ferrua 1979b, 18, n. 528.}
\nf{Statut}{Lecture secondaire (EDCS) ; écarts avec l’EDB à arbitrer sur l’ICVR. L’EDCS signale une photographie (non consultée).}
\end{notice}

\begin{notice}{442}{400}{ICUR n.s. I, 2110 · EDCS-39300329 · EDB32110}
\label{not:442}
\nf{Datation}{400 — directe : lecture partielle d’un nom consulaire}
\nf{Genre}{Épitaphe ; ensembles A, B, C ; 65 lettres transmises ; 11 \% restituées.}
\nf{Références}{ICUR I 2110 = ILCV 3696.}
\nf{Provenance}{Rattachement selon l’EDB : \emph{Incertae originis} (le volume I des ICVR réunit les inscriptions dont la provenance est perdue). Conservation selon l’EDB : \emph{Musei Vaticani: Galleria Lapidaria}.}
\nf{Support}{Plaque de marbre. 14 × 21 cm (h. × l.). Inscription gravée. Fiche EDB : Filippo Piazzolla, première version 2015.}
\ntext{\lnum{1}]nnul() vix(it) Iuli\\
\lnum{2}[- - -] Denudatae(?) VN\\
\lnum{3}[- - -]ae in q(ua ?) vivae\\
\lnum{4}[- - -] casam perpetuam\\
\lnum{5}[- - - Stilich]onis anno s(ae)cul(ari)}
\nf{Témoins}{EDB (EDB32110 ; écarts à examiner) : \emph{------ / [--- s]trenue vix(---?) reli/[---] denudatae un/[---]tae inq(u-) vivae / [---] casam perpetuam / [cons(ulatu) Stilic]onis anno scul(psit?)}.}
\nf{Traduction}{[…] vécut […] Iuli[…] ; […] dans laquelle, de leur vivant […] une demeure perpétuelle […] ; sous (le consulat de) Stilicon, en l’année séculaire (?).}
\nf{Langue}{6 abréviations développées par l’éditeur (dont nnul, vix(it), Denudatae(?), q(ua ?), s(ae)cul, (ari)).}
\nf{Chronologie}{Noms consulaires : Stilicho (lu en partie).}
\nf{Bibliographie}{ICUR I 2110 ; ILCV 3696. Bibliographie complémentaire selon l’EDB : Lega 2000 ; ILCV, 3696.}
\nf{Statut}{Lecture secondaire (EDCS) ; écarts avec l’EDB à arbitrer sur l’ICVR.}
\end{notice}

\begin{notice}{443}{400}{ICUR n.s. II, 6059 · EDCS-34200685 · EDB16490}
\label{not:443}
\nf{Datation}{400 — directe : nom consulaire lu}
\nf{Genre}{Épitaphe ; ensembles A, B, C ; 68 lettres transmises ; 3 \% restituées.}
\nf{Références}{ICUR II 6059.}
\nf{Provenance}{Rattachement selon l’EDB : \emph{Coem. Commodillae}.}
\nf{Support}{Plaque de marbre. Inscription gravée. Fiche EDB : Carlo Carletti, première version 2003.}
\ntext{\lnum{1}Ursus qui vixit\\
\lnum{2}ann⸢o⸣s pl(us) minus XL\\
\lnum{3}depositus XVIII Kal(endas) A(u)gus\\
\lnum{4}[ta]s Fl(avio) St⸢i⸣l\{l\}ic(h)on⸢e⸣ v(iro) c(larissimo) cons(ule)}
\nf{Témoins}{Export EDCS de 2020 : \emph{Ursus qui vixit / anns pl(us) minus XL / depositus XVIII Kal(endas) A(u)gus/[ta]s Fl(avio) Stl\{l\}ic(h)on v(iro) c(larissimo) cons(ule)}. EDB (EDB16490) : texte concordant aux conventions de notation près (indice de similarité 0,959).}
\nf{Apparat}{⸢o⸣ : V gravé ; ⸢i⸣ : E gravé ; ⸢e⸣ : I gravé ; \{l\} : lettres superflues}
\nf{Traduction}{Ursus, qui vécut environ 40 ans ; déposé le 18e jour avant les calendes d’août, sous le consulat de Flavius Stilicon, clarissime.}
\nf{Langue}{L’éditeur signale 3 graphies corrigées : V gravé pour o (confusion o/u) ; E gravé pour i (confusion e/i) ; I gravé pour e (confusion e/i). Des lettres gravées sont jugées superflues par l’éditeur (entre accolades). 8 abréviations développées par l’éditeur (dont pl(us), Kal(endas), A(u)gus, Fl(avio), ic(h)on, v(iro)).}
\nf{Chronologie}{Noms consulaires : Stilicho (lu). Dates du jour : déposition : XVIII Kalendas Augustas = 15 juillet 400 (partiellement restitué). Âge déclaré : env. 40 a. (lu).}
\nf{Bibliographie}{ICUR II 6059.}
\nf{Statut}{Lecture secondaire (EDCS) concordante avec l’EDB : lecture stabilisée entre deux bases, non collationnée sur l’ICVR ni sur le monument.}
\end{notice}

\begin{notice}{444}{400}{ICUR n.s. II, 6060 · EDCS-34200686 · EDB16491}
\label{not:444}
\nf{Datation}{400 — directe : nom consulaire lu}
\nf{Genre}{Épitaphe ; ensembles A, B, C ; 64 lettres transmises ; 23 \% restituées.}
\nf{Références}{ICUR II 6060 = AIIRoma-05, p 111 = AE 1905, 83.}
\nf{Provenance}{Rattachement selon l’EDB : \emph{Coem. Commodillae}.}
\nf{Support}{Plaque de marbre. Inscription gravée. Reproductions signalées par l’EDB : photographie PCAS (droits réservés ; non reproduites). Fiche EDB : Carlo Carletti, première version 2003.}
\ntext{\lnum{1}dep]osita bene merenti in pace s[ub]\\
\lnum{2}[die - - - K]al(endas) No⸢v⸣(embres) Fl(avio) Stil\{l\}ichone v(iro) c(larissimo)\\
\lnum{3}[iter(um) con]sule}
\nf{Témoins}{Export EDCS de 2020 : \emph{dep]osita bene merenti in pace s[ub] / [die 3 K]al(endas) No(embres) Fl(avio) Stil\{l\}ichone v(iro) c(larissimo) / [iter(um) con]sule}. EDB (EDB16491 ; écarts à examiner) : \emph{[dep]osita benemerenti in pace / [--- k]al(end---) nob(embr---) Fl(avio) Stillichone v(iro) c(larissimo) / [con]sule}.}
\nf{Apparat}{⸢v⸣ : B gravé ; \{l\} : lettres superflues}
\nf{Traduction}{Déposée, bien méritante, en paix, le jour […] avant les calendes de novembre, sous le consulat de Flavius Stilicon, clarissime, pour la 2e fois [restitution « iterum » problématique pour 400].}
\nf{Langue}{L’éditeur signale une graphie corrigée : B gravé pour v (bêtacisme). Des lettres gravées sont jugées superflues par l’éditeur (entre accolades). 6 abréviations développées par l’éditeur (dont Kal(endas), (embres), Fl(avio), v(iro), c(larissimo), iter(um)).}
\nf{Chronologie}{Noms consulaires : Stilicho (lu). Réserve : Itération [iter(um)] de Stilichon restituée ; compatibilité avec 400 à contrôler. Dates du jour : déposition : … Kalendas Novembres (partiellement restitué).}
\nf{Commentaire}{La formule in pace est conservée. L’itération restituée appartient à la reconstruction éditoriale, non au niveau des lettres conservées ; toute analyse de chronologie doit la traiter séparément.}
\nf{Bibliographie}{ICUR II 6060 ; AIIRoma-05, p 111 ; AE 1905, 83.}
\nf{Statut}{Lecture secondaire (EDCS) ; écarts avec l’EDB à arbitrer sur l’ICVR. L’EDCS signale une photographie (non consultée).}
\end{notice}

\begin{notice}{445}{400}{ICUR n.s. III, 8168a · EDCS-35700197 · EDB24690}
\label{not:445}
\nf{Datation}{400 — formule transmise par l’EDB (lue, itération restituée)}
\nf{Genre}{Épitaphe ; ensembles A, B, C ; 45 lettres transmises ; 36 \% restituées.}
\nf{Références}{ICUR III 8168a.}
\nf{Provenance}{Rattachement selon l’EDB : \emph{Basilica ss. Nerei et Achillei in coem. Domitillae}. Conservation selon l’EDB : \emph{Basilica ss. Nerei et Achillei in coem. Domitillae}.}
\nf{Support}{Plaque de marbre. 44 × 58 × 2,5 cm (h. × l. × ép.). Inscription gravée. Reproductions signalées par l’EDB : photographie PCAS (droits réservés ; non reproduites). Fiche EDB : Carlo Carletti / Antonio Enrico Felle, première version 2003.}
\ntext{\lnum{1}[- - -]ino[- - -]\\
\lnum{2}[- - -]ev[- - -]hi[- - -]\\
\lnum{3}[- - - Prob]inianu[s(?) - - -]\\
\lnum{4}in p[ace - - -]\\
\lnum{5}[- - - pi]la secunda in ba[silica - - -]}
\nf{Témoins}{EDB (EDB24690 ; divergence forte) : \emph{[---]̣ḷ̲i̲n̲o̲ [---] / [---]ae VA[---] H̲I̲[---] / [--- ]inianu[s ---] / i̲n̲ p̲[ace ---] / [--- pi]la sẹc̣unda in b̲a̲[silica? ---] / cons(ulatu) Stellici v(iri) c(larissimi) [iterum? ---]}.}
\nf{Traduction}{[…] Probinianus […] en paix […] ; […] le deuxième pilier dans la basilique […].}
\nf{Langue}{1 abréviation développée par l’éditeur (dont Probinianus(?)).}
\nf{Chronologie}{Le texte EDCS ne conserve aucun nom consulaire. La fiche EDB (EDB24690) transmet une dernière ligne : \emph{cons(ulatu) Stellici v(iri) c(larissimi) [iterum? ---]}. Le nom de Stilichon est lu ; l'itération est restituée avec doute. Sans \emph{iterum}, l'année est 400 ; avec \emph{iterum}, 405, hors corpus. A. E. Felle note dans l'EDB que les lettres \emph{BA[silica]} sont visibles sur la photographie PCAS mais non sur l'image publiée par Marucchi en 1914.}
\nf{Commentaire}{La formule in pace est attestée seulement par restitution.}
\nf{Bibliographie}{ICUR III 8168a.}
\nf{Statut}{Lecture secondaire (EDCS) ; divergence forte avec l’EDB : le texte ne doit pas être exploité sans contrôle de l’édition imprimée. L’EDCS signale une photographie (non consultée).}
\end{notice}

\begin{notice}{446}{400}{ICUR n.s. III, 8730 · EDCS-35900973 · EDB20586}
\label{not:446}
\nf{Datation}{400 — directe : nom consulaire lu}
\nf{Genre}{Épitaphe ; ensembles A, B, C ; 55 lettres transmises.}
\nf{Références}{ICUR III 8730 = ILCV +2281.}
\nf{Provenance}{Rattachement selon l’EDB : \emph{Tituli qui in coem. Callisti reperti traduntur}. Conservation selon l’EDB : \emph{Musei Vaticani: Lapidario Cristiano}.}
\nf{Support}{Plaque de marbre. Inscription gravée. Fiche EDB : Carlo Carletti, première version 2018.}
\ntext{\lnum{1}Vincentius qui vi(xit)\\
\lnum{2}ann(os) p(lus) m(inus) LX\\
\lnum{3}d(epositus) p⸢r⸣id(ie) Id(us) Ian(uarias) Fl(avio)\\
\lnum{4}St\{h\}ilic(h)on⸢e⸣ v(iro) c(larissimo) co(n)s(ule)\\
\lnum{5}in pace q(uiescas)}
\nf{Témoins}{Export EDCS de 2020 : \emph{Vincentius qui vi(xit) / ann(os) p(lus) m(inus) LX / d(epositus) pid(ie) Id(us) Ian(uarias) Fl(avio) / St\{h\}ilic(h)on v(iro) c(larissimo) co(n)s(ule) / in pace q(uiescas)}. EDB (EDB20586) : texte concordant aux conventions de notation près (indice de similarité 0,966).}
\nf{Apparat}{⸢r⸣ : I gravé ; ⸢e⸣ : I gravé ; \{h\} : lettres superflues}
\nf{Traduction}{Vincentius, qui vécut environ 60 ans ; déposé la veille des ides de janvier, sous le consulat de Flavius Stilicon, clarissime ; puisses-tu reposer en paix.}
\nf{Langue}{L’éditeur signale 2 graphies corrigées : I gravé pour r ; I gravé pour e (confusion e/i). Des lettres gravées sont jugées superflues par l’éditeur (entre accolades). 15 abréviations développées par l’éditeur (dont vi(xit), ann(os), p(lus), m(inus), d(epositus), id(ie)).}
\nf{Chronologie}{Noms consulaires : Stilicho (lu). Ne sont pas des consuls : Vincentius qui vi(xit) (défunt). Dates du jour : déposition : pridie Idus Ianuarias = 12 janvier 400 (lu). Âge déclaré : env. 60 a. (lu).}
\nf{Commentaire}{La formule in pace est conservée. Vocabulaire religieux : repos et sommeil de la mort — quiescas (abréviation développée).}
\nf{Bibliographie}{ICUR III 8730 ; ILCV +2281.}
\nf{Statut}{Lecture secondaire (EDCS) concordante avec l’EDB : lecture stabilisée entre deux bases, non collationnée sur l’ICVR ni sur le monument. L’EDCS signale une photographie (non consultée).}
\end{notice}

\annee{401}

\begin{notice}{447}{401}{ICUR n.s. I, 1942 · EDCS-25402147 · EDB12500}
\label{not:447}
\nf{Datation}{401 — formule transmise par l’EDB (partiellement restituée)}
\nf{Genre}{Épitaphe ; ensembles A, B, C ; 122 lettres transmises ; 43 \% restituées.}
\nf{Références}{AE 1910, 205 = ICUR I 1942 = SEG 30, 1208.}
\nf{Provenance}{Rattachement selon l’EDB : \emph{Incertae originis} (le volume I des ICVR réunit les inscriptions dont la provenance est perdue). Conservation selon l’EDB : \emph{Museo Archeologico Nazionale}.}
\nf{Support}{Plaque de marbre. 34 × 26 × 3 cm (h. × l. × ép.). Inscription gravée. Fiche EDB : Antonio Enrico Felle, première version 2003.}
\ntext{\lnum{1}ὁρκίζομε δὲ ὑμᾶς [τὸν θ(εὸ)ν τὸν πα]\\
\lnum{2}ντοκράτορα εἴ τις το[λμήσει σκυ]\\
\lnum{3}λῖσε τὰ ὀστέα αὐτοῦ [ἢ ἄλλον τινὰ]\\
\lnum{4}[θῖνε ἔσ]τε ἔνοχος τῇ [κρίσει - - -]\\
\lnum{5}[- - -]αν ἡμαίραν τῆς [ἀναστάσεως(?)]}
\nf{Témoins}{Export EDCS de 2020 : \emph{texte grec absent}. EDB (EDB12500 ; divergence forte) : \emph{[ἐνθά]δε κεῖτε Εὐ[---] / [Θεο]δοσειοπολίτ[ης ἀμέμπτως ζήσ]/[ας] μικρῷ πλέω ἔ[τη --- ἐκοιμήθη τ]/ῇ πρὸ ιε´ καλ(ανδῶν) σεπτ(εμβρίων) [ὑπατίᾳ Βινκεντί]/ου καὶ Φραβίτου τῶ[ν λαμπροτάτων] / ὁρκίζομε δὲ ὑμᾶς [διὰ τὸν θεὸν τὸν πα]/ντοκράτορα εἴ τις το]λμήσειε σκυ]/λῖσε τὰ ᾡτέα αὐτοῦ [ἢ τάφον τοῦτο]/[ν βλάψε]ιε ἔνοχος τῇ [κολάσει τὴν φο]/[βερ]ὰν ἡμαίραν τῆς [κρίσεως ἔστω] [texte tronqué dans la liste de résultats]}. Les deux témoins diffèrent aussi dans les restitutions. Début (EDB seule) : \emph{[ἐνθά]δε κεῖτε Εὐ[---] / [Θεο]δοσειοπολίτ[ης]} : le défunt est originaire de Théodosiopolis. Fin : l'EDCS restitue \emph{ἔνοχος τῇ [κρίσει - - -] / [- - -]αν ἡμαίραν τῆς [ἀναστάσεως(?)]} (« passible du jugement … au jour de la résurrection ») ; l'EDB restitue \emph{ἔνοχος τῇ [κολάσει τὴν φο]/[βερ]ὰν ἡμαίραν τῆς [κρίσεως ἔστω]} (« qu'il soit passible du châtiment au jour terrible du jugement »). Seuls \emph{ἔνοχος τῇ} et \emph{ἡμαίραν τῆς} sont lus. La mention de la résurrection, déjà conjecturale, est donc contestée par la seconde base : elle ne peut entrer dans aucun compte.}
\nf{Traduction}{(grec, texte EDCS) Je vous adjure par [Dieu le Tout]-Puissant : si quelqu'un ose troubler ses os [ou (déposer) quelqu'un d'autre], qu'il soit passible du [jugement …] au jour de la [résurrection (?)]. — Pour le texte plus complet et les restitutions différentes de l'EDB, voir la confrontation des témoins.}
\nf{Langue}{Le texte est rédigé en grec, en tout ou en partie ; l’export de 2020 ne le transmettait pas. 1 abréviation développée par l’éditeur (dont (εὸ)).}
\nf{Chronologie}{Le texte EDCS ne transmet que les cinq dernières lignes. La fiche EDB (EDB12500 = ICVR I 1942), qui suit une autre édition, en donne dix, dont la date : \emph{τῇ πρὸ ιε´ καλ(ανδῶν) σεπτ(εμβρίων) [ὑπατίᾳ Βινκεντί]ου καὶ Φραβίτου} (18 août 401). Le nom de Fravitta est lu : l'année 401 est fondée sur une formule consulaire partiellement restituée.}
\nf{Commentaire}{Vocabulaire religieux (texte EDCS) : désignation de Dieu — \emph{θ(εὸ)ν} (restitué), \emph{-ντοκράτορα} (lu en partie) ; jugement — \emph{κρίσει} (restitué) ; résurrection — \emph{ἀναστάσεως} (restitué avec doute, contesté par l'EDB). L'adjuration (\emph{ὁρκίζομε δὲ ὑμᾶς}) et l'épithète \emph{-ντοκράτορα} sont lues ; le nom de Dieu est restitué. Le texte est une formule de protection de la tombe. L'idée d'un jugement (\emph{κρίσις}) apparaît dans les deux restitutions, à deux places différentes ; la résurrection, dans une seule. Commentaire développé : dossier p.~\pageref{dos:447}.}
\nf{Bibliographie}{AE 1910, 205 ; ICUR I 1942 ; SEG 30, 1208.}
\nf{Statut}{Lecture secondaire (EDCS) ; divergence forte avec l’EDB : le texte ne doit pas être exploité sans contrôle de l’édition imprimée.}
\end{notice}

\begin{notice}{448}{401}{EDCS-23702282 · EDB41942}
\label{not:448}
\nf{Datation}{401 — directe : nom consulaire lu}
\nf{Genre}{Épitaphe ; ensembles A, B, C ; 82 lettres transmises.}
\nf{Références}{AE 2001, 527.}
\nf{Provenance}{Rattachement selon l’EDB : \emph{ex-vigna Chiaraviglio}. Conservation selon l’EDB : \emph{ex-vigna Chiaraviglio}.}
\nf{Support}{La fiche EDB ne figure pas dans l’état 2019 des données EpiDoc ; support non documenté.}
\ntext{\lnum{1}Depositio Abundanti V Idus Octob(res)\\
\lnum{2}Fl(avio) Vincentio et Frav\{v\}it(t)⸢a⸣ vv(iris) cc(larissimis) conss(ulibus)\\
\lnum{3}qui vixit annos XXVIII in pace}
\nf{Témoins}{Export EDCS de 2020 : \emph{Depositio Abundanti V Idus Octob(res) / Fl(avio) Vincentio et Frav\{v\}it(t) vv(iris) cc(larissimis) conss(ulibus) / qui vixit annos XXVIII in pace}. EDB (EDB41942 ; écarts à examiner) : \emph{depositio Abundanti V idus octob(res) / Fl(avio) Vincentio et Fravitto vv(iris) cc(larissimis) cons(ulibu)s / qui vixit annos XΧVIΙΙ in pace [texte tronqué dans la liste de résultats]}.}
\nf{Apparat}{⸢a⸣ : O gravé ; \{v\} : lettres superflues}
\nf{Traduction}{Déposition d’Abundantius le 5e jour avant les ides d’octobre, sous le consulat de Flavius Vincentius et de Fravitta, clarissimes ; il vécut 28 ans, en paix.}
\nf{Langue}{L’éditeur signale une graphie corrigée : O gravé pour a. Des lettres gravées sont jugées superflues par l’éditeur (entre accolades). 6 abréviations développées par l’éditeur (dont Octob(res), Fl(avio), it(t), vv(iris), cc(larissimis), conss(ulibus)).}
\nf{Chronologie}{Noms consulaires : Vincentius (lu) et Fravitta (lu). Ne sont pas des consuls : Depositio Abundanti (défunt). Dates du jour : déposition : V Idus Octobres = 11 octobre 401 (lu). Âge déclaré : 28 a. (lu).}
\nf{Commentaire}{La formule in pace est conservée.}
\nf{Bibliographie}{AE 2001, 527. Bibliographie complémentaire selon l’EDB : AE, 2001, 527 ; Giuliani 2001, 176-178, A.a.2.}
\nf{Statut}{Lecture secondaire (EDCS) ; écarts avec l’EDB à arbitrer sur l’ICVR.}
\end{notice}

\begin{notice}{449}{401}{ICUR n.s. III, 9365 · EDCS-35900834 · EDB20435}
\label{not:449}
\nf{Datation}{401 — directe : nom consulaire lu}
\nf{Genre}{Épitaphe ; ensembles A, B, C ; 81 lettres transmises.}
\nf{Références}{CIL IX *630,3 = ICUR III 9365.}
\nf{Provenance}{Rattachement selon l’EDB : \emph{Tituli qui in coem. Callisti reperti traduntur}. Conservation selon l’EDB : \emph{Osimo (AN)   ; Palazzo Baldeschi-Balleani}.}
\nf{Support}{Plaque de marbre. Inscription gravée. Fiche EDB : Carlo Carletti, première version 2015.}
\ntext{\lnum{1}Pridie Non(as) Iul(ias)\\
\lnum{2}domina Volusiana Aufidia\\
\lnum{3}decessit et \{et\} ea\\
\lnum{4}deposita est\\
\lnum{5}Fl(aviis) Vincentio\\
\lnum{6}et Fravit(t)⸢a⸣ cons(ulibus)}
\nf{Témoins}{Export EDCS de 2020 : \emph{Pridie Non(as) Iul(ias) / domina Volusiana Aufidia / decessit et \{et\} ea / deposita est / Fl(aviis) Vincentio / et Fravit(t) cons(ulibus)}. EDB (EDB20435) : texte concordant aux conventions de notation près (indice de similarité 0,941).}
\nf{Apparat}{⸢a⸣ : O gravé ; \{et\} : lettres superflues}
\nf{Traduction}{La veille des nones de juillet, dame Volusiana Aufidia est morte et a été déposée, sous le consulat des Flavii Vincentius et Fravitta.}
\nf{Langue}{L’éditeur signale une graphie corrigée : O gravé pour a. Des lettres gravées sont jugées superflues par l’éditeur (entre accolades). 5 abréviations développées par l’éditeur (dont Non(as), Iul(ias), Fl(aviis), Fravit(t), cons(ulibus)).}
\nf{Chronologie}{Noms consulaires : Vincentius (lu) et Fravitta (lu). Dates du jour : événement indéterminé : Pridie Nonas Iulias = 6 juillet 401 (lu).}
\nf{Commentaire}{L’astérisque de la référence CIL est conservé comme signe transmis. Sa portée documentaire et la concordance avec ICUR doivent être vérifiées ; aucune conclusion d’authenticité n’est tirée du signe isolé.}
\nf{Bibliographie}{CIL IX *630,3 ; ICUR III 9365. Bibliographie complémentaire selon l’EDB : CIL XI, 630\textbackslash{}*, 03.}
\nf{Statut}{Lecture secondaire (EDCS) concordante avec l’EDB : lecture stabilisée entre deux bases, non collationnée sur l’ICVR ni sur le monument.}
\end{notice}

\begin{notice}{450}{401}{ICUR n.s. I, 349 · EDCS-38600338 · EDB28045}
\label{not:450}
\nf{Datation}{401 — directe : nom consulaire lu}
\nf{Genre}{Épitaphe ; ensembles A, B, C ; 78 lettres transmises ; 51 \% restituées.}
\nf{Références}{ICUR I 349.}
\nf{Provenance}{Rattachement selon l’EDB : \emph{Incertae originis} (le volume I des ICVR réunit les inscriptions dont la provenance est perdue).}
\nf{Support}{Support non précisé par l’EDB. Fiche EDB : Filippo Piazzolla, première version 2014.}
\ntext{\lnum{1}Hic quiescit in [pace - - -]\\
\lnum{2}qui vixit anno[s - - - consulatu]\\
\lnum{3}Vincenti et Fra[viti v(virorum) cc(larissimorum)]}
\nf{Témoins}{EDB (EDB28045 ; écarts à examiner) : \emph{hic quiescit in [pace ---] / qui vixit anno[s --- consulatu] / Vincenti et Fra[viti vv(irorum) cc(larissimorum duorum)] [texte tronqué dans la liste de résultats]}.}
\nf{Traduction}{Ici repose en paix […] qui vécut […] ans, sous le consulat de Vincentius et de Fravitta, clarissimes.}
\nf{Langue}{2 abréviations développées par l’éditeur (dont v(virorum), cc(larissimorum)).}
\nf{Chronologie}{Noms consulaires : Fravitta (lu en partie) et Vincentius (lu).}
\nf{Commentaire}{La formule in pace est attestée seulement par restitution. Vocabulaire religieux : repos et sommeil de la mort — quiescit (lu).}
\nf{Bibliographie}{ICUR I 349.}
\nf{Statut}{Lecture secondaire (EDCS) ; écarts avec l’EDB à arbitrer sur l’ICVR. Plus de la moitié des lettres sont restituées.}
\end{notice}

\begin{notice}{451}{401}{ICUR n.s. II, 4506 · EDCS-34100569 · EDB19793}
\label{not:451}
\nf{Datation}{401 — directe : nom consulaire lu}
\nf{Genre}{Épitaphe ; ensembles A, B, C ; 59 lettres transmises.}
\nf{Références}{ICUR II 4506 = ILCV 2974 (em).}
\nf{Provenance}{Rattachement selon l’EDB : \emph{Coem. Pontiani}. Conservation selon l’EDB : \emph{Urbino   ; Palazzo Ducale}.}
\nf{Support}{Plaque de marbre. 12,5 × 29 cm (h. × l.). Inscription gravée. Reproductions signalées par l’EDB : reproduction, Luni et alii 1986 (droits réservés ; non reproduites). Fiche EDB : Antonio Enrico Felle, première version 2015.}
\ntext{\lnum{1}Vincentio et Frav\{v\}i\\
\lnum{2}t⸢a⸣ consul(ibus) V Idus Septemb(res)\\
\lnum{3}deposita Gaudentia in\\
\lnum{4}pace}
\nf{Témoins}{Export EDCS de 2020 : \emph{Vincentio et Frav\{v\}i/t consul(ibus) V Idus Septemb(res) / deposita Gaudentia in / pace}. EDB (EDB19793) : texte concordant aux conventions de notation près (indice de similarité 0,979).}
\nf{Apparat}{⸢a⸣ : O gravé ; \{v\} : lettres superflues}
\nf{Traduction}{Sous le consulat de Vincentius et de Fravitta, le 5e jour avant les ides de septembre, déposée Gaudentia, en paix.}
\nf{Langue}{L’éditeur signale une graphie corrigée : O gravé pour a. Des lettres gravées sont jugées superflues par l’éditeur (entre accolades). 2 abréviations développées par l’éditeur (dont consul(ibus), Septemb(res)).}
\nf{Chronologie}{Noms consulaires : Vincentius (lu) et Fravitta (lu). Dates du jour : événement indéterminé : V Idus Septembres = 9 septembre 401 (lu).}
\nf{Commentaire}{La formule in pace est conservée (répartie sur deux lignes). La coupure in / pace ne crée pas deux unités lexicales indépendantes et ne justifie pas un faux négatif. La formule est présente sans restitution dans le champ transmis.}
\nf{Bibliographie}{ICUR II 4506 ; ILCV 2974 (em). Bibliographie complémentaire selon l’EDB : Luni et alii 1986, 67 fig.3.}
\nf{Statut}{Lecture secondaire (EDCS) concordante avec l’EDB : lecture stabilisée entre deux bases, non collationnée sur l’ICVR ni sur le monument. L’EDCS signale une photographie (non consultée).}
\end{notice}

\begin{notice}{452}{401}{ICUR n.s. II, 4836 · EDCS-36101314}
\label{not:452}
\nf{Datation}{401 — directe : consulat restitué}
\nf{Genre}{Épitaphe ; ensembles A, B, C ; 68 lettres transmises ; 57 \% restituées.}
\nf{Références}{ICUR II 4836.}
\nf{Provenance}{Rome (indication de l’EDCS) ; aucun rattachement cimétérial transmis.}
\nf{Support}{Non documenté dans les sources consultées.}
\ntext{\lnum{1}Theodosius v(ir) [c(larissimus) - - -]\\
\lnum{2}tina moribu[s - - -]\\
\lnum{3}M(a)rt(ias) FFll(avis) V[incentio et Fravita cons(ulibus)]}
\nf{Témoins}{Aucune fiche EDB appariée (numéro ICUR absent de la liste publique de l’EDB).}
\nf{Traduction}{Theodosius, clarissime, […]tina par ses mœurs […] ; […] mars, sous le consulat des Flavii Vincentius et Fravitta.}
\nf{Langue}{6 abréviations développées par l’éditeur (dont v(ir), c(larissimus), M(a)rt, (ias), FFll(avis), cons(ulibus)).}
\nf{Chronologie}{Noms consulaires : Vincentius (restitué) et Fravitta (restitué). Ne sont pas des consuls : Theodosius v(ir) [c(larissimus) (défunt ou dédicant, clarissime (non l’empereur)).}
\nf{Commentaire}{Rang sénatorial d’une personne commémorée (theodosius vir clarissimus …tin).}
\nf{Bibliographie}{ICUR II 4836.}
\nf{Statut}{Lecture secondaire (EDCS) sans second témoin numérique. Plus de la moitié des lettres sont restituées.}
\end{notice}

\begin{notice}{453}{401}{ICUR n.s. II, 6061 · EDCS-34200687 · EDB13893}
\label{not:453}
\nf{Datation}{401 — directe : nom consulaire lu}
\nf{Genre}{Épitaphe ; ensembles A, B, C ; 109 lettres transmises ; 44 \% restituées.}
\nf{Références}{ICUR II 6061.}
\nf{Provenance}{Rattachement selon l’EDB : \emph{Coem. Commodillae}.}
\nf{Support}{Plaque de marbre. Inscription gravée. Signes décrits par l’EDB : chrisme avec alpha et oméga. Fiche EDB : Carlo Carletti, première version 2003.}
\ntext{\lnum{1}Fl(avio) Vincentio v(iro) [c(larissimo) cons(ule)]\\
\lnum{2}XIII Kal(endas) Septemb(res) [in pace de]\\
\lnum{3}cessit nomine E[- - - vixit ann(os) plus]\\
\lnum{4}minus II et dies V[- - - bene meren]\\
\lnum{5}ti in pace}
\nf{Témoins}{EDB (EDB13893) : texte concordant aux conventions de notation près (indice de similarité 0,98).}
\nf{Traduction}{Sous le consulat de Flavius Vincentius, clarissime, le 13e jour avant les calendes de septembre, est mort(e) en paix, nommé(e) E[…], qui vécut environ 2 ans et 5 (?) jours ; bien méritant(e), en paix.}
\nf{Langue}{7 abréviations développées par l’éditeur (dont Fl(avio), v(iro), c(larissimo), cons(ule), Kal(endas), Septemb(res)).}
\nf{Chronologie}{Noms consulaires : Vincentius (lu). Dates du jour : événement indéterminé : XIII Kalendas Septembres = 20 août 401 (lu). Âge déclaré : env. 2 a. (partiellement restitué).}
\nf{Commentaire}{La formule in pace est conservée.}
\nf{Bibliographie}{ICUR II 6061.}
\nf{Statut}{Lecture secondaire (EDCS) concordante avec l’EDB : lecture stabilisée entre deux bases, non collationnée sur l’ICVR ni sur le monument.}
\end{notice}

\begin{notice}{454}{401}{ICUR n.s. II, 6062 · EDCS-34200688 · EDB13894}
\label{not:454}
\nf{Datation}{401 — directe : nom consulaire lu}
\nf{Genre}{Épitaphe ; ensembles A, B, C ; 83 lettres transmises ; 4 \% restituées.}
\nf{Références}{ICUR II 6062 = ILCV 3503 = AE 1905, 84.}
\nf{Provenance}{Rattachement selon l’EDB : \emph{Coem. Commodillae}.}
\nf{Support}{Plaque de marbre. Inscription gravée. Signes décrits par l’EDB : croix monogrammatique avec alpha et oméga. Fiche EDB : Carlo Carletti, première version 2003.}
\ntext{\lnum{1}In hunc locum posita est Marcia\\
\lnum{2}[n]a simul cum Martinum(!)\\
\lnum{3}[ma]ritum suum fecit con[- - -]\\
\lnum{4}[- - -]ii et Fravita v(iro) c(larissimo) consule}
\nf{Témoins}{EDB (EDB13894) : texte concordant aux conventions de notation près (indice de similarité 0,98).}
\nf{Apparat}{(!) : forme signalée par l’éditeur}
\nf{Traduction}{En ce lieu a été placée Marciana, avec Martinus son mari ; elle a fait […] ; sous le consulat de (Vincent)ius et de Fravitta, clarissime.}
\nf{Langue}{L’éditeur signale au moins une forme anormale (sic). 3 abréviations développées par l’éditeur (dont Martinum(!), v(iro), c(larissimo)).}
\nf{Chronologie}{Noms consulaires : Fravitta (lu).}
\nf{Bibliographie}{ICUR II 6062 ; ILCV 3503 ; AE 1905, 84.}
\nf{Statut}{Lecture secondaire (EDCS) concordante avec l’EDB : lecture stabilisée entre deux bases, non collationnée sur l’ICVR ni sur le monument.}
\end{notice}

\begin{notice}{455}{401}{ICUR n.s. II, 6063 · EDCS-34200689 · EDB13895}
\label{not:455}
\nf{Datation}{401 — directe : consulat restitué}
\nf{Genre}{Épitaphe ; ensembles A, B, C ; 115 lettres transmises ; 39 \% restituées.}
\nf{Références}{ICUR II 6063.}
\nf{Provenance}{Rattachement selon l’EDB : \emph{Coem. Commodillae}.}
\nf{Support}{Plaque de marbre. Inscription gravée. Fiche EDB : Carlo Carletti, première version 2003.}
\ntext{\lnum{1}[Toti]us innocentiae\\
\lnum{2}[mirae int]egritatis dulcis\\
\lnum{3}[simae co]niugi Constantiae\\
\lnum{4}[quae vixit a]nn(os) p(lus) m(inus) XXXIII m(enses) III\\
\lnum{5}[- - -] die V Kal(endas) I[- - -]\\
\lnum{6}[- - -]o et F[ravita cons(ulibu)s]}
\nf{Témoins}{EDB (EDB13895) : texte concordant aux conventions de notation près (indice de similarité 1).}
\nf{Traduction}{D’une entière innocence, d’une admirable intégrité, à son épouse très douce Constantia, qui vécut environ 33 ans 3 mois ; […] le 5e jour avant les calendes de […] ; sous le consulat de (Vincenti)us et de Fravitta.}
\nf{Langue}{6 abréviations développées par l’éditeur (dont ann(os), p(lus), m(inus), m(enses), Kal(endas), cons(ulibu)s).}
\nf{Chronologie}{Noms consulaires : Fravitta (restitué), Constantius (lu) et Constans (lu). Âge déclaré : env. 33 a. 3 m. 5 j. (partiellement restitué).}
\nf{Bibliographie}{ICUR II 6063.}
\nf{Statut}{Lecture secondaire (EDCS) concordante avec l’EDB : lecture stabilisée entre deux bases, non collationnée sur l’ICVR ni sur le monument.}
\end{notice}

\begin{notice}{456}{401}{ICUR n.s. III, 8167 · EDCS-35700196 · EDB24689}
\label{not:456}
\nf{Datation}{401 — directe : nom consulaire lu}
\nf{Genre}{Épitaphe ; ensembles A, B, C ; 33 lettres transmises ; 70 \% restituées.}
\nf{Références}{ICUR III 8167.}
\nf{Provenance}{Rattachement selon l’EDB : \emph{Basilica ss. Nerei et Achillei in coem. Domitillae}. Conservation selon l’EDB : \emph{Basilica ss. Nerei et Achillei in coem. Domitillae}.}
\nf{Support}{Plaque de marbre. 21 × 14 cm (h. × l.). Inscription gravée. Fiche EDB : Carlo Carletti / Antonio Enrico Felle, première version 2003.}
\ntext{\lnum{1}]XV d[- - -]\\
\lnum{2}[- - - Vincentio et] Fraudit[o conss(ulibus)}
\nf{Témoins}{EDB (EDB24689 ; écarts à examiner) : \emph{------ / [---]XV d[eposit---] / [--- Vincentio et] Fraudit[o cons(ulibu)s]}.}
\nf{Traduction}{[…] 15 […] ; sous le consulat de Vincentius et de Fravitta.}
\nf{Langue}{1 abréviation développée par l’éditeur (dont conss(ulibus)).}
\nf{Chronologie}{Noms consulaires : Vincentius (restitué) et Fravitta (lu).}
\nf{Bibliographie}{ICUR III 8167.}
\nf{Statut}{Lecture secondaire (EDCS) ; écarts avec l’EDB à arbitrer sur l’ICVR. Plus de la moitié des lettres sont restituées. L’EDCS signale une photographie (non consultée).}
\end{notice}

\begin{notice}{457}{401}{ICUR n.s. III, 8168b · EDCS-35700198 · EDB24691}
\label{not:457}
\nf{Datation}{401 — conjecturale}
\nf{Genre}{Épitaphe ; ensembles A, B, C ; 53 lettres transmises ; 36 \% restituées.}
\nf{Références}{ICUR III 8168b.}
\nf{Provenance}{Rattachement selon l’EDB : \emph{Basilica ss. Nerei et Achillei in coem. Domitillae}. Conservation selon l’EDB : \emph{Basilica ss. Nerei et Achillei in coem. Domitillae}.}
\nf{Support}{Plaque de marbre. 48 × 29 cm (h. × l.). Inscription gravée. Reproductions signalées par l’EDB : photographie PCAS ×2 (droits réservés ; non reproduites). Fiche EDB : Carlo Carletti / Antonio Enrico Felle, première version 2003.}
\ntext{\lnum{1}] sivi merit[- - -]\\
\lnum{2}[- - -]ur et nostro [- - -]\\
\lnum{3}[- - - ora]tionibus t[uis - - -]\\
\lnum{4}[- - - pro fil]io et fr[atribus]}
\nf{Témoins}{EDB (EDB24691 ; écarts à examiner) : \emph{------ / [---] sivi merit[---] / [---]ur et nostro [---] / [--- in ora]tionibus t[uis ---] / [--- Vincent]io et Ep̣[atio consulibus? ---]}.}
\nf{Traduction}{[…] pour soi (?) […] mérite […] ; […] et notre […] ; […] par tes prières […] ; […] pour son fils et ses frères.}
\nf{Chronologie}{Aucun nom consulaire n'est lu dans le texte EDCS. L'EDB (EDB24691) restitue une dernière ligne \emph{[--- Vincent]io et Ep̣[atio consulibus? ---]} et date le texte « 359 ; 401 ». Datation conjecturale ; l'année 401 est celle de la base.}
\nf{Commentaire}{Vocabulaire religieux : intercession — orationibus (partiellement restitué).}
\nf{Bibliographie}{ICUR III 8168b. Bibliographie complémentaire selon l’EDB : ILCV, 2350 adn.}
\nf{Statut}{Lecture secondaire (EDCS) ; écarts avec l’EDB à arbitrer sur l’ICVR. L’EDCS signale une photographie (non consultée).}
\end{notice}

\begin{notice}{458}{401}{ICUR n.s. IV, 11133 · EDCS-37200420 · EDB39597}
\label{not:458}
\nf{Datation}{401 — directe : nom consulaire lu}
\nf{Genre}{Épitaphe ; ensembles A, B, C ; 92 lettres transmises.}
\nf{Références}{ICUR IV 11133 = ILCV 1468.}
\nf{Provenance}{Rattachement selon l’EDB : \emph{Crypta s.Cornelii}. Conservation selon l’EDB : \emph{Coem. Callisti area subdialis et incerti loci}.}
\nf{Support}{Plaque de marbre. 19 × 23 × 2,5 cm (h. × l. × ép.). Inscription gravée. Reproductions signalées par l’EDB : photographie PCAS (droits réservés ; non reproduites). Fiche EDB : Marida Pierno, première version 2015.}
\ntext{\lnum{1}Pr(idie) Id(us) Iun(ias) pausa\\
\lnum{2}⸢vi⸣t Praetioso\\
\lnum{3}annorum puel(l)a\\
\lnum{4}virgo XII tantum\\
\lnum{5}ancilla dei et Chr(isti)\\
\lnum{6}Fl(avio) Vincentio et\\
\lnum{7}Fravito v(iris) c(larissimis) conss(ulibus)}
\nf{Témoins}{Export EDCS de 2020 : \emph{Pr(idie) Id(us) Iun(ias) pausa/t Praetioso / annorum puel(l)a / virgo XII tantum / ancilla dei et Chr(isti) / Fl(avio) Vincentio et / Fravito v(iris) c(larissimis) conss(ulibus)}. EDB (EDB39597) : texte concordant aux conventions de notation près (indice de similarité 0,965).}
\nf{Apparat}{⸢vi⸣ : BE gravé}
\nf{Traduction}{La veille des ides de juin a reposé Pretiosa, jeune fille âgée, vierge, de 12 ans seulement, servante de Dieu et du Christ, sous le consulat de Flavius Vincentius et de Fravitta, clarissimes.}
\nf{Langue}{L’éditeur signale une graphie corrigée : BE gravé pour vi.}
\nf{Chronologie}{Noms consulaires : Vincentius (lu) et Fravitta (lu). Dates du jour : événement indéterminé : Pridie Idus Iunias = 12 juin 401 (lu).}
\nf{Commentaire}{Vocabulaire religieux : désignation de Dieu — dei (lu), ancilla dei (lu) ; désignation du Christ — Christi (abréviation développée) ; repos et sommeil de la mort — pausavit (lu). Ancilla dei et Chr(isti) associe Dieu et le Christ dans une qualification personnelle. La formule ne nomme ni le Père ni le Fils comme tels et n’énumère pas les trois personnes de la Trinité. Désignation divine : retenue dans le profil strict. Commentaire développé : dossier p.~\pageref{dos:458}.}
\nf{Bibliographie}{ICUR IV 11133 ; ILCV 1468.}
\nf{Statut}{Lecture secondaire (EDCS) concordante avec l’EDB : lecture stabilisée entre deux bases, non collationnée sur l’ICVR ni sur le monument. L’EDCS signale une photographie (non consultée).}
\end{notice}

\begin{notice}{459}{401}{ICUR n.s. IV, 11134 · EDCS-37200421 · EDB39598}
\label{not:459}
\nf{Datation}{401 — directe : nom consulaire lu}
\nf{Genre}{Épitaphe ; ensembles A, B, C ; 30 lettres transmises ; 3 \% restituées.}
\nf{Références}{ICUR IV 11134.}
\nf{Provenance}{Rattachement selon l’EDB : \emph{Coem. Callisti area subdialis et incerti loci}. Conservation selon l’EDB : \emph{Coem. Callisti area subdialis et incerti loci}.}
\nf{Support}{Sarcophage de marbre. Inscription gravée. Fiche EDB : Marida Pierno, première version 2015.}
\ntext{\lnum{1}D(e)p(osita) Sirica VII[I] Id(us) Aug(ustas) Vincentio cons(ulibus)}
\nf{Témoins}{EDB (EDB39598) : texte concordant aux conventions de notation près (indice de similarité 1).}
\nf{Traduction}{Déposée, Sirica, le 8e jour avant les ides d’août, sous le consulat de Vincentius.}
\nf{Langue}{5 abréviations développées par l’éditeur (dont D(e)p, (osita), Id(us), Aug(ustas), cons(ulibus)).}
\nf{Chronologie}{Noms consulaires : Vincentius (lu). Dates du jour : déposition : VIII Idus Augustas = 6 août 401 (partiellement restitué).}
\nf{Bibliographie}{ICUR IV 11134.}
\nf{Statut}{Lecture secondaire (EDCS) concordante avec l’EDB : lecture stabilisée entre deux bases, non collationnée sur l’ICVR ni sur le monument.}
\end{notice}

\annee{402}

\begin{notice}{460}{402}{ICUR n.s. I, 718 · EDCS-38701119 · EDB28945}
\label{not:460}
\nf{Datation}{402 — directe : nom consulaire lu}
\nf{Genre}{Épitaphe ; ensembles A, B, C ; 109 lettres transmises ; 38 \% restituées.}
\nf{Références}{ICUR I 718 = ILCV 3142a.}
\nf{Provenance}{Rattachement selon l’EDB : \emph{Incertae originis} (le volume I des ICVR réunit les inscriptions dont la provenance est perdue).}
\nf{Support}{Support non précisé par l’EDB. Fiche EDB : Filippo Piazzolla, première version 2014.}
\ntext{\lnum{1}Hic requiescit Leontia in pace [quae vixit - - -]\\
\lnum{2}deposita est Kal(endis) Iuniis co[ns(ulatu) dd(ominorum) nn(ostrorum) Arcadi]\\
\lnum{3}Aug(usti) V et Honori Aug(usti) V con(sulum)}
\nf{Témoins}{EDB (EDB28945) : texte concordant aux conventions de notation près (indice de similarité 0,944).}
\nf{Traduction}{Ici repose Leontia, en paix, qui vécut […] ans ; elle a été déposée aux calendes de juin, sous le consulat de nos seigneurs Arcadius Auguste pour la 5e fois et Honorius Auguste pour la 5e fois.}
\nf{Langue}{7 abréviations développées par l’éditeur (dont Kal(endis), cons(ulatu), dd(ominorum), nn(ostrorum), Aug(usti)).}
\nf{Chronologie}{Noms consulaires : Arcadius (restitué) et Honorius (lu). Ne sont pas des consuls : requiescit Leontia (défunte). Dates du jour : déposition : Kalendis Iuniis = 1 juin 402 (lu).}
\nf{Commentaire}{La formule in pace est conservée. Vocabulaire religieux : repos et sommeil de la mort — requiescit (lu).}
\nf{Bibliographie}{ICUR I 718 ; ILCV 3142a. Bibliographie complémentaire selon l’EDB : ILCV, 3142A.}
\nf{Statut}{Lecture secondaire (EDCS) concordante avec l’EDB : lecture stabilisée entre deux bases, non collationnée sur l’ICVR ni sur le monument.}
\end{notice}

\begin{notice}{461}{402}{ICUR n.s. I, 1214 · EDCS-38700752 · EDB30165}
\label{not:461}
\nf{Datation}{402 — directe : nom consulaire lu}
\nf{Genre}{Épitaphe ; ensembles A, B, C ; 85 lettres transmises ; 47 \% restituées.}
\nf{Références}{ICUR I 1214.}
\nf{Provenance}{Rattachement selon l’EDB : \emph{Incertae originis} (le volume I des ICVR réunit les inscriptions dont la provenance est perdue). Conservation selon l’EDB : \emph{Musei Capitolini}.}
\nf{Support}{Plaque de marbre. 75 × 150 cm (h. × l.). Inscription gravée. Fiche EDB : Filippo Piazzolla, première version 2014.}
\ntext{\lnum{1}Matri dulcissim(a)e [- - - qu(a)e vixit]\\
\lnum{2}annos LXIII d(ies) XX de[p(osita) - - -]\\
\lnum{3}Mai(as) Arcadio et Hon[orio Augg(ustis) V conss(ulibus)]}
\nf{Témoins}{EDB (EDB30165) : texte concordant aux conventions de notation près (indice de similarité 0,95).}
\nf{Traduction}{À sa mère très douce […] qui vécut 63 ans 20 jours ; déposée […] mai, sous le consulat d’Arcadius et d’Honorius, Augustes pour la 5e fois.}
\nf{Langue}{7 abréviations développées par l’éditeur (dont dulcissim(a)e, qu(a)e, d(ies), dep(osita), Mai(as), Augg(ustis)).}
\nf{Chronologie}{Noms consulaires : Arcadius (lu) et Honorius (lu en partie). Âge déclaré : 63 a. 20 j. (partiellement restitué).}
\nf{Bibliographie}{ICUR I 1214. Bibliographie complémentaire selon l’EDB : Ferrua 1979b, 29, n. 1214.}
\nf{Statut}{Lecture secondaire (EDCS) concordante avec l’EDB : lecture stabilisée entre deux bases, non collationnée sur l’ICVR ni sur le monument.}
\end{notice}

\begin{notice}{462}{402}{ICUR n.s. I, 1256 · EDCS-39101966 · EDB35372}
\label{not:462}
\nf{Datation}{402 — directe : nom consulaire lu}
\nf{Genre}{Épitaphe ; ensembles A, B, C ; 37 lettres transmises ; 16 \% restituées.}
\nf{Références}{ICUR I 1256.}
\nf{Provenance}{Rattachement selon l’EDB : \emph{Coemeteria inter s.Felicitatem et viam Anapo}. Conservation selon l’EDB : \emph{Musei Capitolini}.}
\nf{Support}{Plaque de marbre. 48 × 36 × 2,5 cm (h. × l. × ép.). Inscription gravée. Fiche EDB : Filippo Piazzolla, première version 2014.}
\ntext{\lnum{1}[- - - S]aturnino Ru[- - -]\\
\lnum{2}[- - -]ore die Satu[rni - - -]\\
\lnum{3}[- - - co]ns(ulibus) Fl(avio) Arcadio [}
\nf{Témoins}{EDB (EDB35372 ; divergence forte) : \emph{S̲aturnino ⌜q⌝(ui) v̲i̲x̲i̲̲t̲ a̲n̲[nos ---] / [--- corp]ore die Satur̲n̲i̲s̲ v̲i̲[gilia? ---] / [co]ns(ulibus) Fl(avio) Arcadio e̲t̲ H̲o̲n̲o̲r̲[io] [texte tronqué dans la liste de résultats]}.}
\nf{Traduction}{[…] Saturninus […] ; […] un samedi […] ; […] sous le consulat de Flavius Arcadius […].}
\nf{Langue}{2 abréviations développées par l’éditeur (dont cons(ulibus), Fl(avio)).}
\nf{Chronologie}{Noms consulaires : Arcadius (lu). Ne sont pas des consuls : [S]aturnino q(ui) v(ixit) (EDB) (défunt). \emph{Saturnino} est le nom du défunt (\emph{S̲aturnino ⌜q⌝(ui) v̲i̲x̲i̲t̲}, EDB35372), non celui d'un consul. L'EDB lit \emph{[co]ns(ulibus) Fl(avio) Arcadio e̲t̲ H̲o̲n̲o̲r̲[io]} (lettres soulignées : vues par un éditeur antérieur), ce qui confirme 402.}
\nf{Bibliographie}{ICUR I 1256. Bibliographie complémentaire selon l’EDB : ICVR IX, 23763 ; Ferrua 1979b, 31, n.1256.}
\nf{Statut}{Lecture secondaire (EDCS) ; divergence forte avec l’EDB : le texte ne doit pas être exploité sans contrôle de l’édition imprimée.}
\end{notice}

\begin{notice}{463}{402}{ICUR n.s. I, 1943 · EDCS-39200386 · EDB31861}
\label{not:463}
\nf{Datation}{402 — directe : nom consulaire lu}
\nf{Genre}{Épitaphe ; ensembles A, B, C ; 81 lettres transmises.}
\nf{Références}{ICUR I 1943 = ILCV 2576 (add).}
\nf{Provenance}{Rattachement selon l’EDB : \emph{Incertae originis} (le volume I des ICVR réunit les inscriptions dont la provenance est perdue). Conservation selon l’EDB : \emph{Musei Vaticani: Lapidario Cristiano}.}
\nf{Support}{Plaque de marbre. 33 × 53 × 3,5 cm (h. × l. × ép.). Inscription gravée. Reproductions signalées par l’EDB : reproduction, Perret 1851 (droits réservés ; non reproduites). Fiche EDB : Filippo Piazzolla, première version 2015.}
\ntext{\lnum{1}Laurentius in pace\\
\lnum{2}qui vixit ann(os) IIII d(ies) XXXV\\
\lnum{3}depositus XIII Kal(endas) Octob(res)\\
\lnum{4}cons\{s\}(ulatu) Arcadi et Honori\\
\lnum{5}Augg(ustorum) VI}
\nf{Témoins}{EDB (EDB31861) : texte concordant aux conventions de notation près (indice de similarité 0,938).}
\nf{Apparat}{\{s\} : lettres superflues}
\nf{Traduction}{Laurentius, en paix, qui vécut 4 ans 35 jours ; déposé le 13e jour avant les calendes d’octobre, sous le consulat d’Arcadius et d’Honorius, Augustes, pour la 6e (sic) fois.}
\nf{Langue}{Des lettres gravées sont jugées superflues par l’éditeur (entre accolades). 6 abréviations développées par l’éditeur (dont ann(os), d(ies), Kal(endas), Octob(res), (ulatu), Augg(ustorum)).}
\nf{Chronologie}{Noms consulaires : Arcadius (lu) et Honorius (lu). Dates du jour : déposition : XIII Kalendas Octobres = 19 septembre 402 (lu). Âge déclaré : 4 a. 35 j. (lu).}
\nf{Commentaire}{La formule in pace est conservée. Le chiffre VI ne doit pas être attribué automatiquement aux deux empereurs ; la portée de l’itération reste à contrôler.}
\nf{Bibliographie}{ICUR I 1943 ; ILCV 2576 (add). Bibliographie complémentaire selon l’EDB : ILCV, 2576.}
\nf{Statut}{Lecture secondaire (EDCS) concordante avec l’EDB : lecture stabilisée entre deux bases, non collationnée sur l’ICVR ni sur le monument. L’EDCS signale une photographie (non consultée).}
\end{notice}

\begin{notice}{464}{402}{ICUR n.s. I, 3219 · EDCS-40300353 · EDB33166}
\label{not:464}
\nf{Datation}{402 — directe : nom consulaire lu}
\nf{Genre}{Épitaphe ; ensembles A, B, C ; 93 lettres transmises ; 62 \% restituées.}
\nf{Références}{ICUR I 3219.}
\nf{Provenance}{Rattachement selon l’EDB : \emph{Incertae originis} (le volume I des ICVR réunit les inscriptions dont la provenance est perdue).}
\nf{Support}{Support non précisé par l’EDB. Fiche EDB : Filippo Piazzolla, première version 2015.}
\ntext{\lnum{1}Locus ubi requi[escit - - -]\\
\lnum{2}annis LX pausa[vit die - - - conss(ulatu) dd(ominorum) nn(ostrorum) Arcadi et]\\
\lnum{3}Honori V do⸢r⸣[mit in pace}
\nf{Témoins}{Export EDCS de 2020 : \emph{Locus ubi requi[escit 3] / annis LX pausa[vit die 3 conss(ulatu) dd(ominorum) nn(ostrorum) Arcadi et] / Honori V do[mit in pace}. EDB (EDB33166 ; écarts à examiner) : \emph{locus ubi requi[escit ---] / annis L`X´ pausạ[vit die --- cons(ulibus) dd(ominis) nn(ostris) Arcadi et] / Honori V doṛ[mit ? ---]}.}
\nf{Apparat}{⸢r⸣ : I gravé}
\nf{Traduction}{Lieu où repose […] ; à 60 ans il a reposé, le jour […], sous le consulat de nos seigneurs Arcadius et Honorius pour la 5e fois ; il dort en paix.}
\nf{Langue}{L’éditeur signale une graphie corrigée : I gravé pour r.}
\nf{Chronologie}{Noms consulaires : Arcadius (restitué) et Honorius (lu).}
\nf{Commentaire}{La formule in pace est attestée seulement par restitution. Vocabulaire religieux : repos et sommeil de la mort — dormit (partiellement restitué), requiescit (partiellement restitué), pausavit (partiellement restitué).}
\nf{Bibliographie}{ICUR I 3219. Bibliographie complémentaire selon l’EDB : ILCV, 3239 adn.}
\nf{Statut}{Lecture secondaire (EDCS) ; écarts avec l’EDB à arbitrer sur l’ICVR. Plus de la moitié des lettres sont restituées.}
\end{notice}

\begin{notice}{465}{402}{ICUR n.s. I, 4054 · EDCS-40301026 · EDB34225}
\label{not:465}
\nf{Datation}{402 — directe : consulat restitué}
\nf{Genre}{Épitaphe ; ensembles A, B, C ; 32 lettres transmises ; 53 \% restituées.}
\nf{Références}{ICUR I 4054.}
\nf{Provenance}{Rattachement selon l’EDB : \emph{Incertae originis} (le volume I des ICVR réunit les inscriptions dont la provenance est perdue).}
\nf{Support}{Plaque de marbre. 19 × 30 cm (h. × l.). Inscription gravée. Signes décrits par l’EDB : chrisme. Fiche EDB : Filippo Piazzolla, première version 2015.}
\ntext{\lnum{1}]E\\
\lnum{2}[- - -]res\\
\lnum{3}[- - - Arcadio et Honorio A]Auugg(ustis) V conss(ulibus)}
\nf{Témoins}{EDB (EDB34225 ; écarts à examiner) : \emph{------ / [---]e ((monogramma Christi)) / [--- dep(osit---) ---]ḅṛes / [--- Arcadio et Honorio A]auug(ustis duobus) V cons(ulibu)s}.}
\nf{Traduction}{[…] ; sous le consulat d’Arcadius et d’Honorius, Augustes pour la 5e fois.}
\nf{Langue}{2 abréviations développées par l’éditeur (dont AAuugg(ustis), conss(ulibus)).}
\nf{Chronologie}{Noms consulaires : Arcadius (restitué) et Honorius (restitué).}
\nf{Bibliographie}{ICUR I 4054. Bibliographie complémentaire selon l’EDB : Ferrua 1979b, 121, n. 4054.}
\nf{Statut}{Lecture secondaire (EDCS) ; écarts avec l’EDB à arbitrer sur l’ICVR. Plus de la moitié des lettres sont restituées.}
\end{notice}

\begin{notice}{466}{402}{ICUR n.s. II, 6064 · EDCS-34200690 · EDB13896}
\label{not:466}
\nf{Datation}{402 — directe : nom consulaire lu}
\nf{Genre}{Épitaphe ; ensembles A, B, C ; 69 lettres transmises.}
\nf{Références}{ICUR II 6064 = ILCV 2927a (add, em) = ILCV 4415 (add, em) = AE 1905, 93.}
\nf{Provenance}{Rattachement selon l’EDB : \emph{Coem. Commodillae}. Conservation selon l’EDB : \emph{Coem. Commodillae}.}
\nf{Support}{Plaque de marbre. Inscription gravée. Reproductions signalées par l’EDB : photographie PCAS (droits réservés ; non reproduites). Fiche EDB : Carlo Carletti, première version 2003.}
\ntext{\lnum{1}(H)ic positus est\\
\lnum{2}Victor qui\\
\lnum{3}vixit ann⸢o⸣s\\
\lnum{4}LX m(enses) VI dep(ositus) XV\\
\lnum{5}K(a)l(e)n(das) I⸢a⸣nua\\
\lnum{6}r(i)as Arcadio\\
\lnum{7}et T(h)eodosio}
\nf{Témoins}{Export EDCS de 2020 : \emph{(H)ic positus est / Victor qui / vixit anns / LX m(enses) VI dep(ositus) XV / K(a)l(e)n(das) Inua/r(i)as Arcadio // et T(h)eodosio}. EDB (EDB13896) : texte concordant aux conventions de notation près (indice de similarité 0,977).}
\nf{Apparat}{⸢o⸣ : V gravé ; ⸢a⸣ : E gravé}
\nf{Traduction}{Ici a été placé Victor, qui vécut 60 ans 6 mois ; déposé le 15e jour avant les calendes de janvier, sous (le consulat d’)Arcadius // et Théodose (?).}
\nf{Langue}{L’éditeur signale 2 graphies corrigées : V gravé pour o (confusion o/u) ; E gravé pour a (e pour a).}
\nf{Chronologie}{Noms consulaires : Theodosius (lu) et Arcadius (lu). Réserve : Paire Arcadius–Théodose à contrôler sur les fastes de 402. Dates du jour : déposition : Kalendas Ianuarias = 1 janvier 402 (lu). Âge déclaré : 60 a. 6 m. (lu).}
\nf{Commentaire}{Aucun nom consulaire n’est remplacé dans le texte. L’alerte concerne la datation du commentaire et le statut de concordance, non une nouvelle lecture du monument.}
\nf{Bibliographie}{ICUR II 6064 ; ILCV 2927a (add, em) ; ILCV 4415 (add, em) ; AE 1905, 93.}
\nf{Statut}{Lecture secondaire (EDCS) concordante avec l’EDB : lecture stabilisée entre deux bases, non collationnée sur l’ICVR ni sur le monument. L’EDCS signale une photographie (non consultée).}
\end{notice}

\begin{notice}{467}{402}{ICUR n.s. III, 8433 · EDCS-35802020 · EDB24941}
\label{not:467}
\nf{Datation}{402 — directe : nom consulaire lu}
\nf{Genre}{Épitaphe ; ensembles A, B, C ; 83 lettres transmises ; 35 \% restituées.}
\nf{Références}{ICUR III 8433.}
\nf{Provenance}{Rattachement selon l’EDB : \emph{Incertum locum in coem. Domitillae}. Conservation selon l’EDB : \emph{coem. Domitillae, aula supra coemeterium}.}
\nf{Support}{Plaque de marbre. Inscription gravée. Reproductions signalées par l’EDB : photographie PCAS ×3 (droits réservés ; non reproduites). Fiche EDB : Carlo Carletti / Antonio Enrico Felle, première version 2003.}
\ntext{\lnum{1}[Hic] requiescit Pasc[- - -]\\
\lnum{2}qui vixit ann(os) VII\\
\lnum{3}[dep(ositus)] VII Idus Iulias d(ominis) n(ostris) [Arca]\\
\lnum{4}[dio] Aug(usto) V et Honor[io Aug(usto) V]}
\nf{Témoins}{EDB (EDB24941) : texte concordant aux conventions de notation près (indice de similarité 0,989).}
\nf{Traduction}{Ici repose Pasc[…], qui vécut 7 ans ; déposé le 7e jour avant les ides de juillet, sous (le consulat de) nos seigneurs Arcadius Auguste pour la 5e fois et Honorius Auguste pour la 5e fois.}
\nf{Langue}{6 abréviations développées par l’éditeur (dont ann(os), dep(ositus), d(ominis), n(ostris), Aug(usto)).}
\nf{Chronologie}{Noms consulaires : Arcadius (restitué) et Honorius (lu). Dates du jour : déposition : VII Idus Iulias = 9 juillet 402 (lu). Âge déclaré : 7 a. (partiellement restitué).}
\nf{Commentaire}{Vocabulaire religieux : repos et sommeil de la mort — requiescit (lu).}
\nf{Bibliographie}{ICUR III 8433.}
\nf{Statut}{Lecture secondaire (EDCS) concordante avec l’EDB : lecture stabilisée entre deux bases, non collationnée sur l’ICVR ni sur le monument. L’EDCS signale une photographie (non consultée).}
\end{notice}

\begin{notice}{468}{402}{ICUR n.s. IV, 12253 · EDCS-37600687 · EDB7068}
\label{not:468}
\nf{Datation}{402 — directe : consulat restitué}
\nf{Genre}{Épitaphe ; ensembles A, B, C ; 57 lettres transmises ; 49 \% restituées.}
\nf{Références}{ICUR IV 12253.}
\nf{Provenance}{Rattachement selon l’EDB : \emph{Coem. anonymum ad viam Ardeatinam}. Conservation selon l’EDB : \emph{Coem. anonymum ad viam Ardeatinam}.}
\nf{Support}{Plaque de marbre. Inscription gravée. Reproductions signalées par l’EDB : photographie PCAS (droits réservés ; non reproduites). Fiche EDB : Marida Pierno, première version 2003.}
\ntext{\lnum{1}] vixit annos XX\\
\lnum{2}[- - - d]uo in ⸢v⸣iduo\\
\lnum{3}[- - - Arcadio Aug]usto V et H[onorio V cons(ulibus)}
\nf{Témoins}{Export EDCS de 2020 : \emph{] vixit annos XX / [3 d]uo in iduo / [3 Arcadio Aug]usto V et H[onorio V cons(ulibus)}. EDB (EDB7068 ; écarts à examiner) : \emph{[---] vixit annos XX[---] / [--- d]uo in biduo / [Arcadio Aug]usto V et H[onorio Augusto V cons(ulibu)s]}.}
\nf{Apparat}{⸢v⸣ : B gravé}
\nf{Traduction}{[…] vécut 20 ans ; […] deux (ans) dans le veuvage ; sous le consulat d’Arcadius Auguste pour la 5e fois et d’Honorius pour la 5e fois.}
\nf{Langue}{L’éditeur signale une graphie corrigée : B gravé pour v (bêtacisme).}
\nf{Chronologie}{Noms consulaires : Arcadius (restitué) et Honorius (restitué). Âge déclaré : 20 a. (partiellement restitué).}
\nf{Bibliographie}{ICUR IV 12253. Bibliographie complémentaire selon l’EDB : CLRE, 663.}
\nf{Statut}{Lecture secondaire (EDCS) ; écarts avec l’EDB à arbitrer sur l’ICVR. L’EDCS signale une photographie (non consultée).}
\end{notice}

\annee{403}

\begin{notice}{469}{403}{ICUR n.s. VIII, 23434 · EDCS-23201556 · EDB21579}
\label{not:469}
\nf{Datation}{403 — directe : nom consulaire lu}
\nf{Genre}{Épitaphe ; ensembles A, B, C ; 89 lettres transmises ; 1 \% restituées.}
\nf{Références}{CIL VI 32955 = ICUR VIII 23434 = ILCV 732 (add).}
\nf{Provenance}{Rattachement selon l’EDB : \emph{Coem. s.Felicitatis}. Conservation selon l’EDB : \emph{Coem. s.Felicitatis}.}
\nf{Support}{Plaque de marbre. 43 × 79 × 3,5 cm (h. × l. × ép.). Inscription gravée. Fiche EDB : Ruggero Lombardi, première version 2003.}
\ntext{\lnum{1}V(ivit)\\
\lnum{2}[- - -]co sc(h)olastico qui vixit\\
\lnum{3}[- - -]III dep(ositus) IIII K(alendas) Feb(ruarias) in pace\\
\lnum{4}[- - -]inus \{i\}sc(h)olasticus sororis\\
\lnum{5}[- - - A]ug(usto) T(h)eodosio et Rumorido conss(ulibus)}
\nf{Témoins}{Export EDCS de 2020 : \emph{V(ivit) / [3]co sc(h)olastico qui vixit / [3]III dep(ositus) IIII K(alendas) Feb(ruarias) in pace / [3]inus \{i\}sc(h)olasticus sororis / [3] v(iris) c(larissimis ?) T(h)eodosio et Rumorido conss(ulibus)}. EDB (EDB21579) : texte concordant aux conventions de notation près (indice de similarité 0,955).}
\nf{Apparat}{\{i\} : lettres superflues}
\nf{Traduction}{Il vit […] ; au scholasticus […], qui vécut […]3 ans ; déposé le 4e jour avant les calendes de février, en paix ; […]inus, scholasticus, de sa sœur […] ; sous le consulat de […] Théodose Auguste et de Rumoridus.}
\nf{Langue}{L’export de 2020 diffère de l’état 2022 (voir apparat). Des lettres gravées sont jugées superflues par l’éditeur (entre accolades). 9 abréviations développées par l’éditeur (dont V(ivit), sc(h)olastico, dep(ositus), K(alendas), Feb(ruarias), sc(h)olasticus).}
\nf{Chronologie}{Noms consulaires : Theodosius (lu) et Rumoridus (lu). Dates du jour : déposition : IIII Kalendas Februarias = 29 janvier 403 (lu).}
\nf{Commentaire}{La formule in pace est conservée. Fonction ou métier mentionné : scholastico, scholasticus.}
\nf{Bibliographie}{CIL VI 32955 ; ICUR VIII 23434 ; ILCV 732 (add). Bibliographie complémentaire selon l’EDB : CIL VI, 32955.}
\nf{Statut}{Lecture secondaire (EDCS) concordante avec l’EDB : lecture stabilisée entre deux bases, non collationnée sur l’ICVR ni sur le monument.}
\end{notice}

\begin{notice}{470}{403}{ICUR n.s. I, 713 · EDCS-35100466 · EDB28939}
\label{not:470}
\nf{Datation}{403 — directe : nom consulaire lu}
\nf{Genre}{Épitaphe ; ensembles A, B, C ; 448 lettres transmises.}
\nf{Références}{ICUR I 713 = ILCV 4744 (em) = CLE 682.}
\nf{Provenance}{Rattachement selon l’EDB : \emph{Incertae originis} (le volume I des ICVR réunit les inscriptions dont la provenance est perdue).}
\nf{Support}{Sarcophage de marbre. Inscription gravée. Texte métrique. Fiche EDB : Filippo Piazzolla, première version 2014.}
\ntext{\lnum{1}Quod dulcis nati quod cara pignera praestant\\
\lnum{2}continet hic tumulus membra qui parva retentat\\
\lnum{3}dolorem sine fine dedit Felicitas isto\\
\lnum{4}clauditur infelix falso cognomine dicta\\
\lnum{5}nam cui prima dies in lucem protulit ann⸢o⸣s\\
\lnum{6}quattuor et male cum quintum sumeret ortum\\
\lnum{7}abstulit atra dies et funere mersit acer⸢b⸣o\\
\lnum{8}haec mater et genitor conscribunt carmina busto\\
\lnum{9}Quo legenti simul redeat sub corde figura\\
\lnum{10}et siccato saepe madescant lumina fletu\\
\lnum{11}sic medicatur amor nec curant carmina manes\\
\lnum{12}dep(osita) in p(ace) Kal(endis) Iuli⸢i⸣s\\
\lnum{13}d(omino) n(ostro) Theodosio Aug(usto) et Fl(avio) Rumordio\\
\lnum{14}v(iro) c(larissimo) conss(ulibus)}
\nf{Témoins}{Export EDCS de 2020 : \emph{Quod dulcis nati quod cara pignera praestant / continet hic tumulus membra qui parva retentat / dolorem sine fine dedit Felicitas isto / clauditur infelix falso cognomine dicta / nam cui prima dies in lucem protulit anns / quattuor et male cum quintum sumeret ortum / abstulit atra dies et funere mersit acero / haec mater et genitor conscribunt carmina busto // Quo legenti simul redeat sub corde fi…}. EDB (EDB28939) : même texte ; les écarts relevés dans la liste de résultats tiennent à la notation (relecture de la fiche détaillée).}
\nf{Apparat}{⸢o⸣ : V gravé ; ⸢b⸣ : V gravé ; ⸢i⸣ : A gravé}
\nf{Traduction}{(vers) Ce qu’offrent un doux enfant, ce qu’offrent de chers gages, ce tombeau le contient, qui retient de petits membres ; elle a causé une douleur sans fin, Felicitas, enfermée ici, malheureuse, appelée d’un surnom trompeur ; car à elle, que le premier jour avait portée à la lumière, quatre années (furent données), et, alors qu’elle prenait malheureusement son cinquième commencement, un jour noir l’emporta et la plongea dans une mort prématurée ; sa mère et son père composent ces vers pour le tombeau, // afin qu’en même temps, pour qui les lit, l’image revienne au cœur, et que les yeux, souvent séchés de pleurs, se mouillent (encore) ; ainsi l’amour se soigne, et les Mânes n’ont cure des vers. Déposée en paix aux calendes de juillet, sous le consulat de notre seigneur Théodose Auguste et de Flavius Rumoridus, clarissime.}
\nf{Langue}{L’éditeur signale 3 graphies corrigées : V gravé pour o (confusion o/u) ; V gravé pour b (bêtacisme) ; A gravé pour i.}
\nf{Chronologie}{Noms consulaires : Theodosius (lu) et Rumoridus (lu). Dates du jour : déposition : Kalendis Iuliis = 1 juillet 403 (lu).}
\nf{Commentaire}{La formule in pace est transmise sous forme abrégée, sans restitution de lacune ; son développement reste éditorial. La clausule nec curant carmina manes appartient au discours poétique sur les morts et sur l’efficacité du chant. Elle ne constitue pas à elle seule une négation de l’identité chrétienne. Une source littéraire précise ne peut être assignée sans parallèle effectivement contrôlé. Le v. 7, \emph{abstulit atra dies et funere mersit acerbo}, reprend un vers de Virgile (\emph{Énéide}, VI 429 ; repris en XI 28), comme l'EDB le signale. La clausule \emph{nec curant carmina manes} appartient au discours poétique sur les morts et ne constitue pas une négation de l'identité chrétienne. L'EDB décrit une orante et qualifie le texte de métrique. Le support est un sarcophage. Commentaire développé : dossier p.~\pageref{dos:470}.}
\nf{Bibliographie}{ICUR I 713 ; ILCV 4744 (em) ; CLE 682. Bibliographie complémentaire selon l’EDB : ILCV, 4744.}
\nf{Statut}{Lecture secondaire (EDCS) concordante avec l’EDB : lecture stabilisée entre deux bases, non collationnée sur l’ICVR ni sur le monument.}
\end{notice}

\begin{notice}{471}{403}{ICUR n.s. I, 1459 · EDCS-39000433 · EDB31246}
\label{not:471}
\nf{Datation}{403 — directe : lecture partielle d’un nom consulaire}
\nf{Genre}{Épitaphe ; ensembles A, B, C ; 67 lettres transmises ; 69 \% restituées.}
\nf{Références}{ICUR I 1459.}
\nf{Provenance}{Rattachement selon l’EDB : \emph{Incertae originis} (le volume I des ICVR réunit les inscriptions dont la provenance est perdue). Conservation selon l’EDB : \emph{Musei Vaticani: Lapidario Cristiano}.}
\nf{Support}{Plaque de marbre. 42 × 30 cm (h. × l.). Inscription gravée. Fiche EDB : Filippo Piazzolla, première version 2015.}
\ntext{\lnum{1}bene mere]nti in pa[ce - - -]\\
\lnum{2}[- - - Theodosio Aug(usto) et R]umorid[o v(iro) c(larissimo) cons(ulibus)}
\nf{Témoins}{EDB (EDB31246) : texte concordant aux conventions de notation près (indice de similarité 0,933).}
\nf{Traduction}{Bien méritant(e), en paix […] ; sous le consulat de Théodose Auguste et de Rumoridus, clarissime.}
\nf{Langue}{4 abréviations développées par l’éditeur (dont Aug(usto), v(iro), c(larissimo), cons(ulibus)).}
\nf{Chronologie}{Noms consulaires : Theodosius (restitué) et Rumoridus (lu en partie).}
\nf{Commentaire}{La formule in pace est attestée seulement par restitution.}
\nf{Bibliographie}{ICUR I 1459. Bibliographie complémentaire selon l’EDB : Ferrua 1964, 110 ; Lega 2000, 265 ; Ferrua 1979b, 39, n. 1459.}
\nf{Statut}{Lecture secondaire (EDCS) concordante avec l’EDB : lecture stabilisée entre deux bases, non collationnée sur l’ICVR ni sur le monument. Plus de la moitié des lettres sont restituées.}
\end{notice}

\begin{notice}{472}{403}{ICUR n.s. I, 1460 · EDCS-39000434 · EDB31247}
\label{not:472}
\nf{Datation}{403 — directe : nom consulaire lu}
\nf{Genre}{Épitaphe ; ensembles A, B, C ; 142 lettres transmises ; 30 \% restituées.}
\nf{Références}{ICUR I 1460 = ILCV 4703.}
\nf{Provenance}{Rattachement selon l’EDB : \emph{Incertae originis} (le volume I des ICVR réunit les inscriptions dont la provenance est perdue). Conservation selon l’EDB : \emph{Musei Vaticani: Lapidario Cristiano}.}
\nf{Support}{Plaque de marbre. 36 × 55 cm (h. × l.). Inscription gravée. Fiche EDB : Filippo Piazzolla, première version 2015.}
\ntext{\lnum{1}coniugi(?)] Venerand(a)e bo[nae(?) - - -]\\
\lnum{2}[- - -]DE innocuo BO[- - -]\\
\lnum{3}[- - -] Florentia dign[a - - -]\\
\lnum{4}qui [- - -]te nobis est ull[- - -]\\
\lnum{5}dign[a pi]a amabilis pu[dica - - -]\\
\lnum{6}qu(a)e vixit ann(os) p(lus) m(inus) [- - -]\\
\lnum{7}T(h)eodosio Aug(usto) et [Rumorido v(iro) c(larissimo) cons(ulibus)]}
\nf{Témoins}{EDB (EDB31247) : texte concordant aux conventions de notation près (indice de similarité 0,969).}
\nf{Traduction}{À son épouse Veneranda, bonne […] ; […] innocente […] ; […] Florentia, digne […] ; […] à nous […] ; digne, pieuse, aimable, chaste […] ; qui vécut environ […] ans ; sous le consulat de Théodose Auguste et de Rumoridus, clarissime.}
\nf{Langue}{12 abréviations développées par l’éditeur (dont coniugi(?), Venerand(a)e, bonae(?), qu(a)e, ann(os), p(lus)).}
\nf{Chronologie}{Noms consulaires : Theodosius (lu) et Rumoridus (restitué).}
\nf{Bibliographie}{ICUR I 1460 ; ILCV 4703. Bibliographie complémentaire selon l’EDB : Lega 2000, 293 ; ILCV, 4703.}
\nf{Statut}{Lecture secondaire (EDCS) concordante avec l’EDB : lecture stabilisée entre deux bases, non collationnée sur l’ICVR ni sur le monument. L’EDCS signale une photographie (non consultée).}
\end{notice}

\begin{notice}{473}{403}{ICUR n.s. I, 3220 · EDCS-40300354 · EDB33167}
\label{not:473}
\nf{Datation}{403 — directe : lecture partielle d’un nom consulaire}
\nf{Genre}{Épitaphe ; ensembles A, B, C ; 66 lettres transmises ; 67 \% restituées.}
\nf{Références}{ICUR I 3220.}
\nf{Provenance}{Rattachement selon l’EDB : \emph{Incertae originis} (le volume I des ICVR réunit les inscriptions dont la provenance est perdue).}
\nf{Support}{Support non précisé par l’EDB. Fiche EDB : Filippo Piazzolla, première version 2015.}
\ntext{\lnum{1}vixit an]n⸢o⸣s VIII[- - -]\\
\lnum{2}[- - - Theodosio Aug(usto)] ⸢e⸣t ⸢Fl⸣(avio) Rumo[rido v(iro) c(larissimo) cons(ulibus)}
\nf{Témoins}{Export EDCS de 2020 : \emph{vixit an]ns VIII[3] / [3 Theodosio Aug(usto)] t (avio) Rumo[rido v(iro) c(larissimo) cons(ulibus)}. EDB (EDB33167) : texte concordant aux conventions de notation près (indice de similarité 0,957).}
\nf{Apparat}{⸢o⸣ : V gravé ; ⸢e⸣ : F gravé ; ⸢Fl⸣ : TI gravé}
\nf{Traduction}{[…] vécut 8 ans […] ; sous le consulat de Théodose Auguste et de Flavius Rumoridus, clarissime.}
\nf{Langue}{L’éditeur signale 3 graphies corrigées : V gravé pour o (confusion o/u) ; F gravé pour e ; TI gravé pour Fl.}
\nf{Chronologie}{Noms consulaires : Theodosius (restitué) et Rumoridus (lu en partie).}
\nf{Bibliographie}{ICUR I 3220.}
\nf{Statut}{Lecture secondaire (EDCS) concordante avec l’EDB : lecture stabilisée entre deux bases, non collationnée sur l’ICVR ni sur le monument. Plus de la moitié des lettres sont restituées.}
\end{notice}

\begin{notice}{474}{403}{ICUR n.s. I, 3221 · EDCS-33900149 · EDB35329}
\label{not:474}
\nf{Datation}{403 — directe : nom consulaire lu}
\nf{Genre}{Épitaphe ; ensembles A, B, C ; 145 lettres transmises ; 22 \% restituées.}
\nf{Références}{ICUR I 3221 = ICUR II 4242 = ILCV 1605.}
\nf{Provenance}{Rattachement selon l’EDB : \emph{Coem. sub divo apud basilicam S.Petri apostoli}.}
\nf{Support}{Support non précisé par l’EDB. Fiche EDB : Filippo Piazzolla, première version 2015.}
\ntext{\lnum{1}Quintilianus homo dei [- - -]\\
\lnum{2}confirmans trinitatem [- - -]\\
\lnum{3}amans castitatem [- - -]\\
\lnum{4}respuens mundum [- - -]\\
\lnum{5}requiescet una cum [- - -]\\
\lnum{6}qui vixit ann(os) p(lus) m(inus) X[- - -]\\
\lnum{7}Theodosio Aug(usto) et R[umorido v(iro) c(larissimo) conss(ulibus)]}
\nf{Témoins}{EDB (EDB35329) : texte concordant aux conventions de notation près (indice de similarité 1). Texte identique dans l'EDCS et l'EDB (EDB35329 = ICVR I 3221 = ICVR II 4242). \emph{homo dei}, \emph{confirmans trinitatem}, \emph{amans castitatem}, \emph{respuens mundum} sont lus dans les deux bases ; chaque ligne est mutilée à droite. La date de 403 repose sur \emph{Theodosio Aug(usto)} (lu) et l'initiale \emph{R}, complétée en \emph{R[umorido]} : entre 340 et 403, seul Rumoridus (403) est collègue d'un Théodose Auguste avec un nom en R- (raisonnement à confirmer sur la CLRE). La pierre est perdue.}
\nf{Traduction}{Quintilianus, homme de Dieu, […] confessant la Trinité, […] aimant la chasteté, […] rejetant le monde, […] repose avec […] ; âge approximatif incomplet : X suivi d’une lacune ; sous Théodose Auguste et R[umoridus…].}
\nf{Langue}{7 abréviations développées par l’éditeur (dont ann(os), p(lus), m(inus), Aug(usto), v(iro), c(larissimo)).}
\nf{Chronologie}{Noms consulaires : Theodosius (lu) et Rumoridus (restitué).}
\nf{Commentaire}{Vocabulaire religieux : désignation de Dieu — dei (lu), homo dei (lu) ; mention de la Trinité — trinitatem (lu) ; repos et sommeil de la mort — requiescet (lu). Confirmans trinitatem est lu dans cet état transmis. Aucune énumération du Père, du Fils et de l’Esprit ne figure dans les lignes conservées ; elle n’est pas ajoutée aux lacunes. Désignation divine : retenue dans le profil strict. Commentaire développé : dossier p.~\pageref{dos:474}.}
\nf{Bibliographie}{ICUR I 3221 ; ICUR II 4242 ; ILCV 1605. Bibliographie complémentaire selon l’EDB : ICVR II, 4242.}
\nf{Statut}{Lecture secondaire (EDCS) concordante avec l’EDB : lecture stabilisée entre deux bases, non collationnée sur l’ICVR ni sur le monument.}
\end{notice}

\begin{notice}{475}{403}{ICUR n.s. II, 4272 · EDCS-33900178 · EDB11273}
\label{not:475}
\nf{Datation}{403 — directe : nom consulaire lu}
\nf{Genre}{Épitaphe ; ensembles A, B, C ; 71 lettres transmises ; 46 \% restituées.}
\nf{Références}{ICUR II 4272 = ILCV 2398 (addd) = BE 663.}
\nf{Provenance}{Rattachement selon l’EDB : \emph{Coem. Octavillae ad s.Pancratium}. Conservation selon l’EDB : \emph{Basilica s. Pancrazio}.}
\nf{Support}{Plaque de marbre. Inscription gravée. Reproductions signalées par l’EDB : photographie PCAS (droits réservés ; non reproduites). Fiche EDB : Antonio Enrico Felle, première version 2003.}
\ntext{\lnum{1}de]dit deus deus tulit\\
\lnum{2}[- - -] Kal(endas) No⸢v⸣(embres) T(h)e⸢o⸣dosio et Rumo\\
\lnum{3}[rido vv(iris) cc(larissimis) conss(ulibus)]}
\nf{Témoins}{Export EDCS de 2020 : \emph{de]dit deus deus tulit / [3] Kal(endas) No(embres) T(h)edosio et Rumo/[rido vv(iris) cc(larissimis) conss(ulibus)]}. EDB (EDB11273) : texte concordant aux conventions de notation près (indice de similarité 0,969).}
\nf{Apparat}{⸢v⸣ : B gravé ; ⸢o⸣ : V gravé}
\nf{Traduction}{Dieu a donné, Dieu a repris ; […] calendes de novembre, sous le consulat de Théodose et de Rumoridus, clarissimes.}
\nf{Langue}{L’éditeur signale 2 graphies corrigées : B gravé pour v (bêtacisme) ; V gravé pour o (confusion o/u).}
\nf{Chronologie}{Noms consulaires : Theodosius (lu) et Rumoridus (lu en partie). Dates du jour : événement indéterminé : … Kalendas Novembres (partiellement restitué).}
\nf{Commentaire}{Vocabulaire religieux : désignation de Dieu — deus (lu). Les deux formes deus sont conservées dans l’état reçu ; le début de dedit est lacunaire. Désignation divine : retenue dans le profil strict. Commentaire développé : dossier p.~\pageref{dos:475}.}
\nf{Bibliographie}{ICUR II 4272 ; ILCV 2398 (addd) ; BE 663. Bibliographie complémentaire selon l’EDB : ILCV, 2398 ; Felle 1995, 68-69, n. 4 ; Felle 2006, 314-315, n. 663.}
\nf{Statut}{Lecture secondaire (EDCS) concordante avec l’EDB : lecture stabilisée entre deux bases, non collationnée sur l’ICVR ni sur le monument. L’EDCS signale une photographie (non consultée).}
\end{notice}

\begin{notice}{476}{403}{ICUR n.s. II, 4507 · EDCS-34100570 · EDB19794}
\label{not:476}
\nf{Datation}{403 — directe : nom consulaire lu}
\nf{Genre}{Épitaphe ; ensembles A, B, C ; 38 lettres transmises ; 47 \% restituées.}
\nf{Références}{ICUR II 4507.}
\nf{Provenance}{Rattachement selon l’EDB : \emph{Coem. Pontiani}.}
\nf{Support}{Support non précisé par l’EDB. Fiche EDB : Antonio Enrico Felle, première version 2003.}
\ntext{\lnum{1}] XII Kal(endas) Iulias\\
\lnum{2}[Theodosio Aug(usto) et] Romodoro [}
\nf{Témoins}{EDB (EDB19794 ; écarts à examiner) : \emph{[---] XII kal(endas) iulias / [Theodosio Aug(usto) et] Romodoro [texte tronqué dans la liste de résultats]}.}
\nf{Traduction}{[…] le 12e jour avant les calendes de juillet, sous (le consulat de) Théodose Auguste et Rumoridus.}
\nf{Langue}{2 abréviations développées par l’éditeur (dont Kal(endas), Aug(usto)).}
\nf{Chronologie}{Noms consulaires : Theodosius (restitué) et Rumoridus (lu). Dates du jour : événement indéterminé : XII Kalendas Iulias = 20 juin 403 (lu).}
\nf{Bibliographie}{ICUR II 4507.}
\nf{Statut}{Lecture secondaire (EDCS) ; écarts avec l’EDB à arbitrer sur l’ICVR.}
\end{notice}

\begin{notice}{477}{403}{ICUR n.s. IV, 11136 · EDCS-37200423 · EDB39600}
\label{not:477}
\nf{Datation}{403 — directe : lecture partielle d’un nom consulaire}
\nf{Genre}{Épitaphe ; ensembles A, B, C ; 60 lettres transmises ; 75 \% restituées.}
\nf{Références}{ICUR IV 11136.}
\nf{Provenance}{Rattachement selon l’EDB : \emph{Coem. Callisti area subdialis et incerti loci}. Conservation selon l’EDB : \emph{Coem. Callisti area subdialis et incerti loci}.}
\nf{Support}{Plaque de marbre. 29 × 29 cm (h. × l.). Inscription gravée. Fiche EDB : Marida Pierno, première version 2015.}
\ntext{\lnum{1}]V Idus [- - -]\\
\lnum{2}[- - -]NIOI [- - - d(omino) n(ostro) Theodosio Aug(usto)]\\
\lnum{3}[et Rumo]rido v(iro) c(larissimo) [conss(ulibus)]}
\nf{Témoins}{EDB (EDB39600 ; écarts à examiner) : \emph{[---]u i[dus ---] / [---]NIOI[--- Theodosio Augusto] / [et Rumo]rido [v̅(iris) c̅(larissimis) consulibus]}.}
\nf{Traduction}{Le 5e jour avant les ides […] ; […] sous le consulat de notre seigneur Théodose Auguste et de Rumoridus, clarissime.}
\nf{Langue}{6 abréviations développées par l’éditeur (dont d(omino), n(ostro), Aug(usto), v(iro), c(larissimo), conss(ulibus)).}
\nf{Chronologie}{Noms consulaires : Theodosius (restitué) et Rumoridus (lu en partie).}
\nf{Bibliographie}{ICUR IV 11136.}
\nf{Statut}{Lecture secondaire (EDCS) ; écarts avec l’EDB à arbitrer sur l’ICVR. Plus de la moitié des lettres sont restituées.}
\end{notice}

\begin{notice}{478}{403}{ICUR n.s. IV, 12425 · EDCS-37600948 · EDB41721}
\label{not:478}
\nf{Datation}{403 — directe : nom consulaire lu}
\nf{Genre}{Mixte (épitaphe et acte d’achat) ; ensembles A, B, C ; 77 lettres transmises.}
\nf{Références}{ICUR IV 12425 = ILCV 3811a.}
\nf{Provenance}{Rattachement selon l’EDB : \emph{inter vias Appiam et Ardeatinam}. Conservation selon l’EDB : \emph{Firenze   ; Palazzo Rinuccini}.}
\nf{Support}{Plaque de marbre. 23 × 98 cm (h. × l.). Inscription gravée. Fiche EDB : Antonello Vilella, première version 2017.}
\ntext{\lnum{1}Rimorido c(larissimo) c(onsule) mense April(i) XVIII Kal(endas)\\
\lnum{2}Maias Barbarus defunctus lo\\
\lnum{3}cu(s) bisomu(s) em(p)tu(s) ab Urs⸢o⸣ fossore}
\nf{Témoins}{Export EDCS de 2020 : \emph{Rimorido c(larissimo) c(onsule) mense April(i) XVIII Kal(endas) / Maias Barbarus defunctus lo/cu(s) bisomu(s) em(p)tu(s) ab Urs fossore}. EDB (EDB41721) : texte concordant aux conventions de notation près (indice de similarité 0,953).}
\nf{Apparat}{⸢o⸣ : V gravé}
\nf{Traduction}{Sous le consulat de Rumoridus, clarissime, au mois d’avril, le 18e jour avant les calendes de mai, Barbarus est mort ; tombe double achetée au fossoyeur Ursus.}
\nf{Langue}{L’éditeur signale une graphie corrigée : V gravé pour o (confusion o/u).}
\nf{Chronologie}{Noms consulaires : Rumoridus (lu). Dates du jour : événement indéterminé : XVIII Kalendas Maias = 14 avril 403 (lu).}
\nf{Commentaire}{Fonction ou métier mentionné : fossore. Le texte doit être segmenté entre l’événement de décès et l’acquisition auprès du fossor. L’année ou le jour d’une opération ne sont pas reportés automatiquement sur l’autre.}
\nf{Bibliographie}{ICUR IV 12425 ; ILCV 3811a. Bibliographie complémentaire selon l’EDB : ILCV, 3811 A.}
\nf{Statut}{Lecture secondaire (EDCS) concordante avec l’EDB : lecture stabilisée entre deux bases, non collationnée sur l’ICVR ni sur le monument.}
\end{notice}
